\section{The other six-manifolds}
\subsection{The cyclic covers: a necessary correction}
\begin{theorem}[Counterexample: the proposed connected sum for cyclic covers]\label{maps:src:D.1.1}
\StatementRecord{maps:src:D.1.1}{proved}{Historical events: Claude Fable 5.2 (Anthropic), effort high, 16 November 2026: carried as background without fault found. GPT-6 Astra (OpenAI), ChatGPT Pro mode, effort Pro, S6\_09\_PAPERS, 17 November 2026: targeted proof check. Also checked in this writing. Pre-continuation writing events: GPT-6 Astra (OpenAI), ChatGPT Pro mode, effort Pro, S6\_10\_PAPERS, 18 November 2026: targeted proof check. Continuation metadata events: GPT-6 Astra (OpenAI), ChatGPT Pro mode, effort Pro, S6\_10\_PAPERS continuation, 19 November 2026: used and rechecked in this run; status-vocabulary reconciliation; metadata reconciliation only; no new mathematical check. \emph{Pre-continuation recorded proof scope.} GPT-6 Astra (OpenAI, ChatGPT Pro mode), effort Pro, 2026-11-17 UTC; Claude Fable 5.2 (Anthropic), effort high, 16 November 2026; also checked in this writing by GPT-6 Astra (OpenAI), ChatGPT Pro mode, effort Pro, 18 November 2026}{the hypotheses, notation and definitions stated here; the explicit constructions and calculations in the proof}


Put $d=|12\ell_0-4\ell_1-3\ell_2|$ for an admissible member $X_p$.
The universal cover $\widetilde X_p$ is a compact complex threefold of
Euler characteristic $2d$. Its Betti numbers necessarily satisfy
\[
 b_3(\widetilde X_p)=2+2b_2(\widetilde X_p)-2d,
 \qquad b_2(\widetilde X_p)\geq d-1.
\]
Consequently, for $d>1$ its second integral homology cannot vanish,
and it cannot be diffeomorphic to $\#^{d-1}(S^3\times S^3)$.
The proposed connected-sum classification under vanishing second homology has an impossible premise for $d>1$;
it is not a classification of the covers.
\end{theorem}


\begin{proof}
\begin{enumerate}[label=\arabic*.,leftmargin=*]
\item The covering has $d$ sheets, is unramified, and pulls back the complex
atlas, so is compact, complex, and oriented. Multiplicativity of Euler
characteristic for a finite covering follows by lifting a finite CW
structure: every cell has exactly $d$ lifts. Thus
$\chi(\widetilde X_p)=d\chi(X_p)=2d$.
Simple connectivity gives $b_1=b_5=0$, and oriented Poincare duality gives
$b_4=b_2$, $b_0=b_6=1$. Substitution gives the displayed identity. The
inequality follows from $b_3\geq0$. Finally a connected sum of $r$
copies of $S^3\times S^3$ has Euler characteristic $2-2r$; for $r=d-1$
this equals $4-2d$, which equals $2d$ only when $d=1$.
For a concrete admissible example take $(\ell_0,\ell_1,\ell_2)=(0,2,-1)$.
Then $d=5$, the cover has Euler characteristic $10$, whereas
$\#^4(S^3\times S^3)$ has Euler characteristic $-6$.
\end{enumerate}
\end{proof}

\begin{theorem}[the cusp of the cyclic cover]\label{maps:src:D.1.2}
\StatementRecord{maps:src:D.1.2}{proved}{Historical events: Claude Fable 5.2 (Anthropic), effort high, 16 November 2026: carried as background without fault found. GPT-6 Astra (OpenAI), ChatGPT Pro mode, effort Pro, S6\_09\_PAPERS, 17 November 2026: targeted proof check. Also checked in this writing. Pre-continuation writing events: GPT-6 Astra (OpenAI), ChatGPT Pro mode, effort Pro, S6\_10\_PAPERS, 18 November 2026: targeted proof check. Continuation metadata events: GPT-6 Astra (OpenAI), ChatGPT Pro mode, effort Pro, S6\_10\_PAPERS continuation, 19 November 2026: used and rechecked in this run; status-vocabulary reconciliation; metadata reconciliation only; no new mathematical check. \emph{Pre-continuation recorded proof scope.} GPT-6 Astra (OpenAI, ChatGPT Pro mode), effort Pro, 2026-11-17 UTC; Claude Fable 5.2 (Anthropic), effort high, 16 November 2026; also checked in this writing by GPT-6 Astra (OpenAI), ChatGPT Pro mode, effort Pro, 18 November 2026}{the hypotheses, notation and definitions stated here; \cite[Theorem~7.17]{AC26}}


The inverse image $\widetilde W_p$ of the cusp is connected. Its
normalization has $d$ components, each isomorphic to $dP_6$. Its
integral homology is torsion free, with ranks
\[
 (b_0,b_1,b_2,b_3,b_4)(\widetilde W_p)=(1,2,d+3,2,d).
\]
Here the covering of the dual torus is the sublattice
$\mathbb Z\widehat u\oplus d\mathbb Z\widehat\gamma$ of
$\overline\Lambda=\mathbb Z\widehat\gamma\oplus\mathbb Z\widehat u$.
\end{theorem}


\begin{proof}
\begin{enumerate}[label=\arabic*.,leftmargin=*]
\item The presentation of the fundamental group in
\cite[Theorem~7.17]{AC26} sends the regular-fibre lattice onto
$\mathbb Z/d$ by its primitive $\gamma$ coordinate: $12$ is a unit
modulo $d$, and $3x=\ell_1c$, $4y=\ell_2c$ express the other two
meridians as multiples of $c$. On the cusp $\pi_1(W)=\overline\Lambda$,
and the same map is surjective. Hence the inverse image is connected,
with the asserted lattice of deck periods. Since
$B_0(\widehat u)=e_1$ and $B_0(\widehat\gamma)=-e_2$, its toric
components are indexed by $\mathbb Z^2/(\mathbb Ze_1+d\mathbb Ze_2)$.

\item We compute the normalization resolution, retaining integral coefficients.
Let $R=\mathbb Z[t]/(t^d-1)$, where $t$ moves one component in the
$e_2$ direction. The degree-zero row of the resolution is the cellular
cochain complex of the dual torus with this $d$-sheeted marking; its
cohomology is $(\mathbb Z,\mathbb Z^2,\mathbb Z)$ in degrees $0,1,2$.
The degree-four row is $R$ in horizontal degree zero. In degree two,
write the six boundary curves on $dP_6$ in cyclic order as
\[
 E_1,\quad H-E_1-E_2,\quad E_2,\quad H-E_2-E_3,
 \quad E_3,\quad H-E_1-E_3.
\]
The restriction differential from the $d$ normalizations to the $3d$
double curves, in the basis $(H,E_1,E_2,E_3)$, is
\[
 B(t)=
 \begin{pmatrix}
 -1&-1&-1&-1\\
 1&1&1&t\\
 -t&-t&-1&-t
 \end{pmatrix}:R^4\longrightarrow R^3.
\]
Replacing the second row by its sum with the first and the third by
its difference with $t$ times the first gives rows with sole entries
$t-1$ in columns four and three. The first-row unit eliminates the
remaining first-row entries. Thus this matrix is equivalent over $R$
to a matrix with diagonal entries $1,t-1,t-1$ and one zero column.
Multiplication by $t-1$ on $R$ has kernel $\mathbb Z(1+t+\cdots+t^{d-1})$
and cokernel $\mathbb Z$. Therefore $\ker B$ is free abelian of rank
$d+2$, and $\operatorname{coker}B\simeq\mathbb Z^2$.
There can be no higher differential: all rows with odd vertical degree
are zero and the resolution has horizontal degrees only $0,1,2$.
The only extension in degree two is an extension of free abelian
groups, hence splits as an abelian group. This gives cohomology ranks
$(1,2,d+3,2,d)$ with no torsion. The universal coefficient theorem,
applied from the top degree downwards, gives the same assertion for
homology. The normalization resolution used here is exact on stalks:
at a point with $r$ local branches it is the augmented cochain complex
of the $(r-1)$-simplex, and hence computes ordinary cohomology even at
the self-identifications and the triple points.
\end{enumerate}
\end{proof}

\begin{theorem}[the rational homology of the cover]\label{maps:src:D.1.3}
\StatementRecord{maps:src:D.1.3}{proved}{Historical events: Claude Fable 5.2 (Anthropic), effort high, 16 November 2026: carried as background without fault found. GPT-6 Astra (OpenAI), ChatGPT Pro mode, effort Pro, S6\_09\_PAPERS, 17 November 2026: targeted proof check. Also checked in this writing. Pre-continuation writing events: GPT-6 Astra (OpenAI), ChatGPT Pro mode, effort Pro, S6\_10\_PAPERS, 18 November 2026: targeted proof check. Continuation metadata events: GPT-6 Astra (OpenAI), ChatGPT Pro mode, effort Pro, S6\_10\_PAPERS continuation, 19 November 2026: used and rechecked in this run; status-vocabulary reconciliation; metadata reconciliation only; no new mathematical check. \emph{Pre-continuation recorded proof scope.} GPT-6 Astra (OpenAI, ChatGPT Pro mode), effort Pro, 2026-11-17 UTC; Claude Fable 5.2 (Anthropic), effort high, 16 November 2026; also checked in this writing by GPT-6 Astra (OpenAI), ChatGPT Pro mode, effort Pro, 18 November 2026}{the hypotheses, notation and definitions stated here; the explicit constructions and calculations in the proof}


For every admissible $p$, the rational homology of $\widetilde X_p$ has
ranks
\[
 (1,0,d-1,0,d-1,0,1).
\]
The deck group on either $H^2(\widetilde X_p;\mathbb C)$ or
$H^4(\widetilde X_p;\mathbb C)$ is the direct sum of its nontrivial
characters, each with multiplicity one. This statement does not
assert the absence of integral torsion.
\end{theorem}


\begin{proof}
\begin{enumerate}[label=\arabic*.,leftmargin=*]
\item Let $\lambda\ne1$ be a $d$th root of unity and let
$\mathcal L_\lambda$ be the associated rank-one complex local system
on $X_p$. Its restriction to a regular fibre is nontrivial on
$\widehat\gamma$. The cellular cochain complex of a torus with this
local system is a Koszul complex containing the invertible scalar
$\lambda-1$ and is therefore contractible. On a reduced multiple
fibre the finite torus cover has the same property. Transfer, or
averaging over the finite deck group in characteristic zero, shows
that its cohomology also vanishes. The regular region and all three
boundary overlaps consequently have zero cohomology with these
coefficients, as follows either by their bundle spectral sequences
or by contraction of the fibre cochain complex. Mayer--Vietoris
therefore identifies $H^*(X_p;\mathcal L_\lambda)$ with
$H^*(W;\mathcal L_\lambda|_W)$.
In the preceding normalization calculation substitute $t=\lambda$.
The degree-zero row is acyclic, the matrix $B(\lambda)$ has rank three
and one-dimensional kernel, and the degree-four row is one-dimensional.
Thus this local system contributes one copy of $\mathbb C$ in degrees
two and four and nothing else. The trivial character contributes
$H^*(X_p;\mathbb C)$, which is concentrated in degrees zero and six by
\textup{[AC26, Theorem 7.22]}. The decomposition of the regular
representation of $\mathbb Z/d$ into its characters computes the
cohomology of the covering and proves the claim.
\end{enumerate}
\end{proof}

\begin{theorem}[the integral homology of the universal cover]\label{maps:src:D.1.4}
\StatementRecord{maps:src:D.1.4}{proved}{Historical events: Claude Fable 5.2 (Anthropic), effort high, 16 November 2026: carried as background without fault found. Claude Opus 5.5 (Anthropic), effort max, sweep 1, scope s6, 18 November 2026: computation or formal certificate. GPT-6 Astra (OpenAI), ChatGPT Pro mode, effort Pro, S6\_09\_PAPERS, 17 November 2026: targeted proof check. Also checked in this writing. Pre-continuation writing events: GPT-6 Astra (OpenAI), ChatGPT Pro mode, effort Pro, S6\_10\_PAPERS, 18 November 2026: targeted proof check. Continuation metadata events: GPT-6 Astra (OpenAI), ChatGPT Pro mode, effort Pro, S6\_10\_PAPERS continuation, 19 November 2026: used and rechecked in this run; status-vocabulary reconciliation; metadata reconciliation only; no new mathematical check. \emph{Pre-continuation recorded proof scope.} GPT-6 Astra (OpenAI, ChatGPT Pro mode), effort Pro, 2026-11-17 UTC; Claude Fable 5.2 (Anthropic), effort high, 16 November 2026; also checked in this writing by GPT-6 Astra (OpenAI), ChatGPT Pro mode, effort Pro, 18 November 2026}{the hypotheses, notation and definitions stated here; \cref{maps:src:D.1.3}}


For every admissible triple, with $d=|p|$, one has
\[
 H_i(\widetilde X_p;\mathbb Z)=
 \begin{cases}
 \mathbb Z,&i=0,6,\\
 \mathbb Z^{d-1}\oplus\mathbb Z/d,&i=2,\\
 \mathbb Z/d,&i=3,\\
 \mathbb Z^{d-1},&i=4,\\
 0,&i=1,5.
 \end{cases}
\]
Thus for $d>1$ the cover is also not a connected sum of copies of
$S^2\times S^4$. In particular a torsion-free classification theorem
cannot be applied to these covers.
\end{theorem}


\begin{proof}
\begin{enumerate}[label=\arabic*.,leftmargin=*]
\item We give the integral attachment calculation. Write
$\gamma'=\gamma/d$ on the covering lattice
$\Lambda_d=\mathbb Z(d\widehat\gamma)\oplus\mathbb Z\widehat u
\oplus\mathbb Z\widehat w\oplus\mathbb Z\widehat\delta$.
Choose integers $h_j$ with
$m_jh_j\equiv\ell_j\pmod d$. Replace each meridian lift by its
product with the translation $-h_j\widehat\gamma$. Its $m_j$th power
has translation vector
$v_j-h_j\mathcal N_j\widehat\gamma$, where
\[
 \mathcal N_1\widehat\gamma=3\varepsilon,
 \qquad \mathcal N_2\widehat\gamma=4\varepsilon'.
\]
The non-invariant part of an affine translation is removed by
conjugation by a real torus translation: averaging projects to the
invariant part, and its complementary part is in the image of $A_j-I$
over $\mathbb R$. This produces precisely the same smooth local
quotient with invariant translation vector. Its new integers are
\[
 a_j=(\ell_j-m_jh_j)/d,\qquad
 a_0=(\ell_0-h_1-h_2)/d.
\]
The numerator defining $a_0$ is divisible by $d$ by the original
meridian relation. They satisfy
\[
 p'=12a_0-4a_1-3a_2=p/d\in\{1,-1\}.
\]
In the basis of $\Lambda_d$, the two matrices are
\[
 A_1^{(d)}=\begin{pmatrix}1&0&0&0\\6d&0&1&0\\-6d&-1&-1&0\\-2d&1&0&1\end{pmatrix},
 \quad
 A_2^{(d)}=\begin{pmatrix}1&0&0&0\\0&0&-1&0\\-6d&1&0&0\\3d&0&1&1\end{pmatrix}.
\]
Their coinvariant lattices are the same saturated rank-two quotients
as for $d=1$. The invariant forms are
$\psi_1=2u+w+3\delta$, $\psi_2=u+w+2\delta$,
and their common invariant alternating form is
$q_d=u\wedge w+6d\gamma'\wedge\delta$.

\item Let $\widetilde X^\circ$ be the complement of the cusp tube. The
finite reduced fibres, denoted $S'_1,S'_2$, have torsion-free first
homology: their abelianized affine presentation is the free rank-two
coinvariant lattice with one meridian adjoined, modulo
$m_ja_j^{\rm mer}=[v'_j]$, a primitive relation since
$\gcd(a_j,m_j)=1$. Poincare duality, Euler characteristic zero and
the universal coefficient theorem give $H_2(S'_j)=\mathbb Z^2$.
The intersection form is even and unimodular: the canonical class is
torsion for a bielliptic surface and $H^2(S'_j)$ is torsion free,
so its mod-two reduction, the second Stiefel--Whitney class, is zero.

\item The pullback images in $H^2(F_d)$ have respective integral bases
\[
 \begin{aligned}
 L_1&=\langle\gamma'\psi_1,\ q_d-2d\gamma'\psi_1\rangle,\\
 L_2&=\langle2\gamma'\psi_2,\ q_d-3d\gamma'\psi_2\rangle.
 \end{aligned}
\]
Here is a verification of the sublattice assertion, including its
index. The invariant lattices are
$\langle\gamma'\psi_j,q_d\rangle$, with intersection determinants
$-9$ and $-4$. Pullback multiplies intersection pairings by $3$ and
$4$, so its indices are respectively one and two. At the second
fibre, $q_d$ cannot descend since its square divided by four is
$3d$, odd. The class $\gamma'\psi_2$ cannot descend either: on the
torus swept by $\delta$ and the affine meridian its integral would
be $-a_2/2$, which is not integral. Of the three index-two sublattices
of the invariant lattice, only the parity sublattice
$\{a\gamma'\psi_2+bq_d:a\equiv b\pmod2\}$ excludes both classes.
The displayed basis generates that sublattice because $d$ is odd.
Its pairing divided by four, and the first pairing divided by three,
are the hyperbolic plane.

\item In homology bases dual to these cohomology bases, the map from the
common regular fibre to the two finite pieces is
\[
 \alpha_2=
 \begin{pmatrix}
 2&1&3&0&0&0\\
 -4d&-2d&0&1&0&0\\
 -2&-2&-4&0&0&0\\
 3d&3d&0&-1&0&0
 \end{pmatrix}.
\]
Its rank is three. Among its maximal minors are $2$, $3d$, $4d$;
their greatest common divisor is one because $d$ is odd. Therefore
its cokernel is $\mathbb Z$, whose primitive dual is
\[
 \Phi=(4d,2,3d,2).
\]
The degree-one map has kernel
$\mathbb Z(-\widehat u-\widehat w+\widehat\delta)$, exactly as at
$d=1$. Mayer--Vietoris for the two finite pieces therefore gives
\[
 0\longrightarrow\mathbb Z\longrightarrow H_2(\widetilde X^\circ)
 \longrightarrow\mathbb Z\longrightarrow0.
\]
Furthermore $H_1(\widetilde X^\circ)=\mathbb Z$, with primitive dual
$\Theta$ restricting to $12\gamma'$ on $F_d$; the cusp meridian
has value $p'=\pm1$ under $\Theta$.

\item Write $N$ for the cusp tube and $M=\partial N$. The exterior-square
matrix of its monodromy has saturated image generated by
$\widehat w\wedge\widehat\delta$ and
$(d\widehat\gamma)\wedge\widehat w+d\widehat u\wedge\widehat\delta$.
Thus the Wang sequence makes $H_2(M)$ free of rank six. Sweeping the
two angular circles $\widehat w,\widehat\delta$ along the toric
meridian gives two classes $\xi_w,\xi_\delta$ which map under the
Wang projection to a basis of the invariant lattice. Both bound
integrally in $N$: sweep the corresponding circle along the extended
zero-section disc in the toric chart. Their span is primitive in
$H_2(M)$. On the other hand
$H_3(N,M)\simeq H^3(N)\simeq H^3(\widetilde W_p)=\mathbb Z^2$;
exactness bounds the kernel of $H_2(M)\to H_2(N)$ by rank two.
Consequently that kernel is exactly
$\mathbb Z\xi_w\oplus\mathbb Z\xi_\delta$.

\item Cutting $\xi_w$ into the two finite pieces gives connecting boundary
$\pm(-\widehat u-\widehat w+\widehat\delta)$, a primitive generator,
whereas $\xi_\delta$ has zero connecting boundary. The latter belongs
to $\operatorname{coker}\alpha_2$. If $\mathcal A$ is the class dual
to $\Phi$, then $\mathcal A|_{F_d}=2q_d$. Integration along the
free $\delta$-circle on $\widetilde X^\circ$ gives
\[
 (\iota_\delta\mathcal A)|_{F_d}=\pm12d\gamma',
 \qquad \iota_\delta\mathcal A=\pm d\Theta.
\]
The second equality follows from injectivity of restriction in degree
one. Evaluation on the meridian therefore gives
$\Phi(j_*\xi_\delta)=\pm d p'=\pm d$. Hence
\[
 H_2(\widetilde X^\circ)/j_*\ker(H_2(M)\to H_2(N))
 \simeq\mathbb Z/d.
\]
This evaluates the actual swept torus, rather than inferring its
coefficient from rational homology.

\item Finally, the exact sequence of the pair
$(\widetilde X_p,\widetilde X^\circ)$ and excision give
\[
 0\longrightarrow\mathbb Z/d\longrightarrow H_2(\widetilde X_p)
 \longrightarrow\ker\{H_2(N,M)\to H_1(\widetilde X^\circ)\}
 \longrightarrow0.
\]
Poincare--Lefschetz duality identifies $H_2(N,M)$ with
$H^4(\widetilde W_p)=\mathbb Z^d$. The extended section disc has
boundary the meridian, whose value is $p'=\pm1$, so the final map is
onto $\mathbb Z$ and its kernel is free of rank $d-1$. The sequence
splits as an abelian-group sequence. The rational calculation of
\ref{maps:src:D.1.3} gives $b_3=0$; duality and universal coefficients now yield
$H_3=\operatorname{Ext}(H_2,\mathbb Z)=\mathbb Z/d$,
$H_4=H^2=\mathbb Z^{d-1}$ and $H_5=0$.
\end{enumerate}
\end{proof}

\begin{proposition}[characteristic and rational ring constraints]\label{maps:src:D.1.5}
\StatementRecord{maps:src:D.1.5}{proved}{Historical events: Claude Fable 5.2 (Anthropic), effort high, 16 November 2026: carried as background without fault found. GPT-6 Astra (OpenAI), ChatGPT Pro mode, effort Pro, S6\_09\_PAPERS, 17 November 2026: targeted proof check. Also checked in this writing. Pre-continuation writing events: GPT-6 Astra (OpenAI), ChatGPT Pro mode, effort Pro, S6\_10\_PAPERS, 18 November 2026: targeted proof check. Continuation metadata events: GPT-6 Astra (OpenAI), ChatGPT Pro mode, effort Pro, S6\_10\_PAPERS continuation, 19 November 2026: used and rechecked in this run; status-vocabulary reconciliation; metadata reconciliation only; no new mathematical check. \emph{Pre-continuation recorded proof scope.} GPT-6 Astra (OpenAI, ChatGPT Pro mode), effort Pro, 2026-11-17 UTC; Claude Fable 5.2 (Anthropic), effort high, 16 November 2026; also checked in this writing by GPT-6 Astra (OpenAI), ChatGPT Pro mode, effort Pro, 18 November 2026}{the hypotheses, notation and definitions stated here; \cref{maps:src:D.1.3}}


For the complex orientation, $w_2(\widetilde X_p)=0$,
$c_1(\widetilde X_p)=0$, and $p_1(\widetilde X_p)$ is torsion of
order dividing $d$. The rational cup product
$H^2(\widetilde X_p;\mathbb Q)\otimes H^2(\widetilde X_p;\mathbb Q)
\to H^4(\widetilde X_p;\mathbb Q)$ is zero. The integral Pontryagin
class is not determined by these assertions.
\end{proposition}


\begin{proof}
\begin{enumerate}[label=\arabic*.,leftmargin=*]
\item The mod-two cohomology $H^2(X_p;\mathbb Z/2)$ vanishes, because all
the integral torsion in degrees one and two has odd order. Thus
$w_2(X_p)=0$ and the same holds on its cover. The class $c_1(X_p)$
is torsion and its pullback lies in the torsion-free group
$H^2(\widetilde X_p;\mathbb Z)$; hence the pullback is zero.
Similarly $H^4(X_p;\mathbb Z)=\mathbb Z/d$, so the pullback $p_1$
is annihilated by $d$.
For the ring statement, the restriction of every nontrivial-character
class in degree two to the normalization of a component of the cusp
lies in the kernel of $B(t)$ for $t\ne1$. This kernel is generated
by $-H+E_1$, whose square is zero in $H^4(dP_6)$.
The product of any two such restricted classes is therefore zero.
The restriction from degree-four cohomology of the covering to
$\bigoplus H^4(dP_6;\mathbb Q)$ is injective: it is the characterwise
isomorphism of \ref{maps:src:D.1.3} on the nontrivial summands, and the source has
no invariant summand. This proves that the original product is zero.
\end{enumerate}
\end{proof}

\begin{theorem}[the integral cohomology ring]\label{maps:src:D.1.6}
\StatementRecord{maps:src:D.1.6}{proved}{Historical events: Claude Fable 5.2 (Anthropic), effort high, 16 November 2026: carried as background without fault found. GPT-6 Astra (OpenAI), ChatGPT Pro mode, effort Pro, S6\_09\_PAPERS, 17 November 2026: targeted proof check. Also checked in this writing. Pre-continuation writing events: GPT-6 Astra (OpenAI), ChatGPT Pro mode, effort Pro, S6\_10\_PAPERS, 18 November 2026: targeted proof check. Continuation metadata events: GPT-6 Astra (OpenAI), ChatGPT Pro mode, effort Pro, S6\_10\_PAPERS continuation, 19 November 2026: used and rechecked in this run; status-vocabulary reconciliation; metadata reconciliation only; no new mathematical check. \emph{Pre-continuation recorded proof scope.} GPT-6 Astra (OpenAI, ChatGPT Pro mode), effort Pro, 2026-11-17 UTC; Claude Fable 5.2 (Anthropic), effort high, 16 November 2026; also checked in this writing by GPT-6 Astra (OpenAI), ChatGPT Pro mode, effort Pro, 18 November 2026}{the hypotheses, notation and definitions stated here; \cref{maps:src:D.1.4}; \cref{maps:src:D.1.3}; \cref{maps:src:D.1.2}}


Suppose $d>1$, so in this construction $d\geq5$. Number the cusp
components $D_j$, $j\in\mathbb Z/d$, by the positive $e_2$ deck
translation, and give all complex subvarieties their complex orientations.
Write $x_j=\operatorname{PD}[D_j]\in H^2(\widetilde X_p;\mathbb Z)$.
These classes generate, with the sole relation $\sum_jx_j=0$.
There is a distinguished generator $\kappa$ of the torsion subgroup of
$H^4(\widetilde X_p;\mathbb Z)$, namely the dual of the complex-oriented
ruling class $H-E_1$ in any normalized cusp component. The products are
\[
 x_i x_j=\begin{cases}
 -2\kappa,&i=j,\\
 \kappa,&i-j=1\text{ or }-1\pmod d,\\
 0,&\text{otherwise}.
 \end{cases}
\]
Together with the orientation pairing between the free parts of $H^2$
and $H^4$, these are every product in positive degree. For $d=1$ the
ring is the ring of $S^6$.
\end{theorem}


\begin{proof}
\begin{enumerate}[label=\arabic*.,leftmargin=*]
\item The relative cohomology group $H^2(N,M)=H_4(N)=\mathbb Z^d$ is
generated by the Thom classes of the irreducible Cartier divisors $D_j$.
This remains true for their self-identifications, since their fundamental
cycles generate $H_4(N)$ and Poincare--Lefschetz duality applies to the
smooth tube. The boundary homomorphism from $H^1(\widetilde X^\circ)$
takes its primitive generator to $(\pm1,\ldots,\pm1)$, with a common
sign. Indeed the toric meridians at adjacent height-one rays differ by
an angular $\widehat w$ or $\widehat\delta$ circle, and both have
zero image under $\Theta$; \ref{maps:src:D.1.4} evaluates any one of them as $p'=\pm1$.
Consequently $\mathbb Z^d/\mathbb Z(1,\ldots,1)$ injects into
$H^2(\widetilde X_p)$. The cokernel embeds in the torsion-free group
$H^2(\widetilde X^\circ)=\mathbb Z^2$ and has rank zero by \ref{maps:src:D.1.3}.
It is therefore zero. This proves the assertion about the generators
integrally, including saturation.

\item For clarity we compute the relevant curve relations before their
images in the whole manifold. The first opposite pair on component
$j$ identifies $E_{1j}$ with $H_j-E_{2j}-E_{3j}$. The other two
opposite pairs identify $E_{2j}$ with $E_{2,j+1}$ and $E_{3j}$ with
$E_{3,j+1}$, after using the first relation. These are precisely the
transpose relations of $B(t)$ in \ref{maps:src:D.1.2}. Thus the ruling classes
$f_j=H_j-E_{1j}=E_{2j}+E_{3j}$ all have the same image $f$ in
$H_2(N)$. Components $j$ and $j+1$ intersect along the two external
boundary curves, with multiplicity one and total class $f$; no
nonadjacent components meet. Hence $x_jx_{j+1}=\operatorname{PD}(f)$.
Multiplying $\sum_ix_i=0$ by $x_j$ gives $x_j^2=-2\operatorname{PD}(f)$.

\item It remains to verify that $f$ is a torsion generator, rather than a
proper multiple. The torus $\widehat u\wedge\widehat\delta$ specializes
to one of the two external boundary-curve classes, up to sign, and
$\widehat u\wedge\widehat w$ specializes to their signed sum. One
can check this in a single strip of the covering tiling: the base
circle has direction $B_0\widehat u=e_1$; it crosses the two edges
with stabilizer directions $e_2-e_1$ and $-e_2$, and its degree on
each surviving angular circle is the determinant with that direction.
These determinants are units. This is the integral annulus-collapse
calculation of AC26 (7.9), applied to the unchanged $e_1$ period.
In the first finite piece, $\widehat u\wedge\widehat\delta$ has
coordinates $(0,0)$ in the dual bases of \ref{maps:src:D.1.4}, whereas
$\widehat u\wedge\widehat w$ has coordinates $(0,1)$.
The functional $\Phi=(4d,2,3d,2)$ therefore evaluates their images
as $0$ and $2$, respectively. It follows that one external curve
has image zero and the other has image $\pm2$ in
$\mathbb Z/d\subset H_2(\widetilde X_p)$. Their complex-oriented
sum $f$ is consequently a generator, since $d$ is odd. Its sign
defines the stated $\kappa$, so the intersection formulas themselves
have no remaining sign ambiguity.

\item All products involving a torsion class and landing in $H^6=\mathbb Z$
vanish. The groups in degrees one and five vanish, and a product of
two degree-three classes is torsion and hence zero. Finally Poincare
duality gives the perfect integral pairing of the free degree-two
and degree-four groups. This exhausts all degree possibilities.
\end{enumerate}
\end{proof}

\begin{theorem}[dependence of the smooth type and the count]\label{maps:src:D.1.7}
\StatementRecord{maps:src:D.1.7}{proved}{Historical events: Claude Fable 5.2 (Anthropic), effort high, 16 November 2026: carried as background without fault found. GPT-6 Astra (OpenAI), ChatGPT Pro mode, effort Pro, S6\_09\_PAPERS, 17 November 2026: targeted proof check. Also checked in this writing. Pre-continuation writing events: GPT-6 Astra (OpenAI), ChatGPT Pro mode, effort Pro, S6\_10\_PAPERS, 18 November 2026: targeted proof check. Continuation metadata events: GPT-6 Astra (OpenAI), ChatGPT Pro mode, effort Pro, S6\_10\_PAPERS continuation, 19 November 2026: used and rechecked in this run; status-vocabulary reconciliation; metadata reconciliation only; no new mathematical check. \emph{Pre-continuation recorded proof scope.} GPT-6 Astra (OpenAI, ChatGPT Pro mode), effort Pro, 2026-11-17 UTC; Claude Fable 5.2 (Anthropic), effort high, 16 November 2026; also checked in this writing by GPT-6 Astra (OpenAI), ChatGPT Pro mode, effort Pro, 18 November 2026}{the hypotheses, notation and definitions stated here; \cref{maps:src:D.1.9}; \cref{maps:src:D.1.12}}


Within the multiplicity-$(3,4)$ construction, the oriented
diffeomorphism type of $X_p$ depends exactly on $d=|p|$.
For each $d$ coprime to six it carries a punctured-disc family of
pairwise non-biholomorphic complex structures, parametrized by
$W=e^{-\pi i(c+k_2)}$, $0<|W|<216$. At each fixed nonzero admissible
linearization scale $u=e^{-2\pi ic}$ it carries exactly two of these
biholomorphism classes. Its universal covers likewise have one
oriented diffeomorphism type for each $d$. This statement identifies
dependence on the construction data. The complete Zhubr classifying
invariant and a canonical cyclic action are supplied in \ref{maps:src:D.1.9}--\ref{maps:src:D.1.12}.
\end{theorem}


\begin{proof}
\begin{enumerate}[label=\arabic*.,leftmargin=*]
\item We use the all-twist comparison in \cref{fam:src:B.2.1}, whose sufficient step is the integral norm equation and analytic torsor gluing. Specifically, equality of $|p|$ and of $W$ gives
$p'=ap$, $a=\pm1$, and $c'-c\equiv k_2-k'_2\pmod2$.
The finite norm equations have integral solutions; the remaining
cusp defect is corrected in its entire rank-two quotient lattice,
and $H^1(\mathcal O\oplus\mathcal O(-1))=0$ glues the translations.
The cusp map preserves the prescribed fan. These are the steps that
give actual biholomorphisms, including every discrete twist, rather
than equality of abstract numerical data.

\item For fixed discrete data, varying $c$ along any compact path in the
admissible half-plane gives a proper smooth family. Here the
properness argument of \cref{fam:src:A.3.1} uses a bounded correction
$C(t,c)$ in $Z=sB_0+C$: on a compact parameter path the effective
logarithmic lattice is uniformly close to $B_0$, so finitely many
closed toric polydiscs provide compact representatives. The finite
free quotients and the ordinary torus parallelepipeds are uniform
as well. All charts are products with the parameter before their
free identifications, so projection is a submersion. The same proof
does not use $p=\pm1$, and therefore applies to the present data.
Integrating a horizontal vector field over the compact path gives
an orientation-preserving diffeomorphism of its end fibres.

\item Given two data with the same $d$, first change one parameter by an
integer to align its value of $c+k_2$ with that of the other
discrete datum; use the biholomorphism of \cref{fam:src:B.2.1} at this aligned parameter,
and connect the resulting parameter to the desired one inside the
admissible half-plane. This proves independence of every datum
except $d$. Distinct $d$ give nonisomorphic fundamental groups.
Every resulting diffeomorphism lifts to an oriented diffeomorphism
of universal covers; different $d$ are distinguished upstairs by
Euler characteristic $2d$.

\item By \cref{fam:src:B.2.1} and \cref{fam:src:C.2.3}, the map to $W$ is an exact biholomorphic classification.
Since $u=W^2$, its fibre at fixed $u\ne0$ is precisely the two
points $W$ and $-W$. No two of these are equal. The asserted count
and the punctured-disc family follow.
\end{enumerate}
\end{proof}

\begin{theorem}[the integral Pontryagin class]\label{maps:src:D.1.8}
\StatementRecord{maps:src:D.1.8}{proved}{Historical events: Claude Fable 5.2 (Anthropic), effort high, 16 November 2026: carried as background without fault found. GPT-6 Astra (OpenAI), ChatGPT Pro mode, effort Pro, S6\_09\_PAPERS, 17 November 2026: targeted proof check. Also checked in this writing. Pre-continuation writing events: GPT-6 Astra (OpenAI), ChatGPT Pro mode, effort Pro, S6\_10\_PAPERS, 18 November 2026: targeted proof check. Continuation metadata events: GPT-6 Astra (OpenAI), ChatGPT Pro mode, effort Pro, S6\_10\_PAPERS continuation, 19 November 2026: used and rechecked in this run; status-vocabulary reconciliation; metadata reconciliation only; no new mathematical check. \emph{Pre-continuation recorded proof scope.} GPT-6 Astra (OpenAI, ChatGPT Pro mode), effort Pro, 2026-11-17 UTC; Claude Fable 5.2 (Anthropic), effort high, 16 November 2026; also checked in this writing by GPT-6 Astra (OpenAI), ChatGPT Pro mode, effort Pro, 18 November 2026}{the hypotheses, notation and definitions stated here; \cref{maps:src:D.1.6}}


The pullback
\[
 \pi^*:H^4(X_p;\mathbb Z)\longrightarrow
 H^4(\widetilde X_p;\mathbb Z)
\]
is zero. In particular
$c_2(\widetilde X_p)=p_1(\widetilde X_p)=0$ integrally.
\end{theorem}


\begin{proof}
\begin{enumerate}[label=\arabic*.,leftmargin=*]
\item We apply the external-curve calculation in \ref{maps:src:D.1.6} first downstairs.
The two mixed regular-fibre tori specialize to one external curve
and a signed sum of the two; their values under the primitive
functional $\Phi=(4,2,3,2)$ are $0$ and $2$.
The first external curve therefore has zero class in
$H_2(X_p)=\mathbb Z/d$, and the other, which we call $C$, has
class $\pm2$ and generates because $d$ is odd.
The covering has exactly $d$ lifts $C_j$ of $C$, each mapped
orientation-preservingly to $C$, and the transpose relations of
$B(t)$ identify all $[C_j]$ in $H_2(\widetilde W_p)$.
Their image is torsion of order $d$ in $H_2(\widetilde X_p)$
by the same mixed-torus calculation. Homology transfer consequently
takes the generator $[C]$ to
\[
 \operatorname{tr}[C]=\sum_j[C_j]=d[C_0]=0.
\]
For a finite orientation-preserving cover, Poincare duality
intertwines cohomological pullback with homological transfer;
this follows by lifting a transverse cycle or from the cap-product
chain identity with the lifted fundamental class. Thus $\pi^*$
vanishes in degree four. Naturality of Chern and Pontryagin
classes, applied to $T\widetilde X_p\simeq\pi^*TX_p$, proves
both assertions. For $d=1$ the target degree-four group already
vanishes.
\end{enumerate}
\end{proof}

\begin{theorem}[the finite-coefficient cubic]\label{maps:src:D.1.9}
\StatementRecord{maps:src:D.1.9}{proved}{Historical events: Claude Fable 5.2 (Anthropic), effort high, 16 November 2026: carried as background without fault found. GPT-6 Astra (OpenAI), ChatGPT Pro mode, effort Pro, S6\_09\_PAPERS, 17 November 2026: targeted proof check. Also checked in this writing. Pre-continuation writing events: GPT-6 Astra (OpenAI), ChatGPT Pro mode, effort Pro, S6\_10\_PAPERS, 18 November 2026: targeted proof check. Continuation metadata events: GPT-6 Astra (OpenAI), ChatGPT Pro mode, effort Pro, S6\_10\_PAPERS continuation, 19 November 2026: used and rechecked in this run; status-vocabulary reconciliation; metadata reconciliation only; no new mathematical check. \emph{Pre-continuation recorded proof scope.} GPT-6 Astra (OpenAI, ChatGPT Pro mode), effort Pro, 2026-11-17 UTC; Claude Fable 5.2 (Anthropic), effort high, 16 November 2026; also checked in this writing by GPT-6 Astra (OpenAI), ChatGPT Pro mode, effort Pro, 18 November 2026}{the hypotheses, notation and definitions stated here; \cref{maps:src:D.1.6}}


Let $d>1$ and let $C_\delta$ be the complex-oriented fixed curve of
the $\delta$ action downstairs. There is a class
$v\in H^2(X_p;\mathbb Z/d)$ with $v(C_\delta)=1$.
For $u=\pi^*v$ and the reductions of the classes $x_j$ in \ref{maps:src:D.1.6},
the full degree-two cubic, evaluated on the complex fundamental
class of $\widetilde X_p$, is
\[
 \begin{aligned}
 \langle x_i x_j x_k,[\widetilde X_p]\rangle&=0,
 &\langle u x_i x_j,[\widetilde X_p]\rangle&=C_{ij},\\
 \langle u^2x_i,[\widetilde X_p]\rangle&=-12,
 &\langle u^3,[\widetilde X_p]\rangle&=0
 \quad\text{in }\mathbb Z/d.
 \end{aligned}
\]
where $C_{ii}=-2$, $C_{i,i\pm1}=1$ and all other entries are zero.
The $d-1$ classes $x_0,\ldots,x_{d-2}$ together with $u$ form a
$(\mathbb Z/d)$-basis of $H^2(\widetilde X_p;\mathbb Z/d)$.
In particular $u$ is deck invariant. No nontrivial extension of the
deck representation is needed to compute this cubic.
\end{theorem}


\begin{proof}
\begin{enumerate}[label=\arabic*.,leftmargin=*]
\item The fixed curve generates $H_2(X_p)=\mathbb Z/d$ by \ref{maps:src:D.1.6}'s
downstairs calculation. The universal coefficient epimorphism
\[
 H^2(X_p;\mathbb Z/d)\longrightarrow
 \operatorname{Hom}(H_2(X_p),\mathbb Z/d)
\]
therefore supplies $v$. Each lift of $C_\delta$ maps to it with
degree one, so $\langle u\kappa,[\widetilde X_p]\rangle=1$.
The diagonal-direction external curve has zero homology class,
and the ruling is its sum with the fixed curve. Thus the assertion
about this pairing uses the complex ruling orientation, not a
choice of a torsion generator up to sign. Universal coefficients
upstairs identify degree-two cohomology with
$\operatorname{Hom}(\mathbb Z^{d-1}\oplus\mathbb Z/d,\mathbb Z/d)$.
Evaluation on its torsion summand shows that $u$ completes the
reductions of the free integral basis to a basis. We choose free
lifts in $H_2$ annihilated by $u$, subtracting multiples of its
torsion generator if necessary; this fixes the splitting used in
the classifying coordinates below. The first two
cubic formulas follow from \ref{maps:src:D.1.6}.

\item We compute the remaining component restriction with its orientation.
In the toric model with lattice coordinates $(r,s)$ on its height-one
vertices, take the integral support function
\[
 \psi(r,s)=-\frac{r(r-1)}2+3s(s-1),
\]
interpolated linearly over each unimodular cone. It defines a toric
line bundle. Its changes under the two period translations are
\[
 \psi(r+1,s)-\psi(r,s)=-r,
 \qquad
 \psi(r,s-1)-\psi(r,s)=-6s+6.
\]
These affine functions give commuting character multipliers
$x_1^{-1}$ and $x_2^{-6}t^6$, so the bundle descends over the
period lattice generated by $e_1$ and $-e_2$. The holomorphic-unit
corrections in the actual cusp action can be deformed to one in a
sufficiently small tube, as in the checked cusp isotopy of AC26,
Section~7.2. This transports the integral topological class and
does not change any of the following pairings.

\item For a wall with endpoints $v,v'$ and the two adjacent third vertices
$w,w'$, the relation of height-one rays is $w+w'=v+v'$.
The two toric trivializations on the curve differ by a power of its
coordinate of exponent
$\psi(v)+\psi(v')-\psi(w)-\psi(w')$. This is its degree for the
complex orientation. Using respectively
\[
 \begin{array}{c|c|c}
 (v,v')&(w,w')&\deg\\\hline
 ((0,0),(1,0))&((0,1),(1,-1))&-6\\
 ((1,0),(0,1))&((0,0),(1,1))&0\\
 ((0,0),(0,1))&((1,0),(-1,1))&1
 \end{array}
\]
gives the internal, diagonal and $\delta$-fixed curve degrees
$(-6,0,1)$. The regular-fibre first Chern class of this bundle is
$-(u\wedge w+6\gamma\wedge\delta)$: the winding of the first
multiplier on the ordered $(\widehat u,\widehat w)$ rectangle is
$-1$, and that of the second on
$(\widehat\gamma,\widehat\delta)$ is $-6$.
There are no other windings. This also fixes the relative signs
that AC26 (7.9), when used with separately oriented curve generators,
leaves implicit.

\item The complement class dual to $\Phi=(4,2,3,2)$ restricts to twice
$q=u\wedge w+6\gamma\wedge\delta$. The specialization from the
regular fibre to the downstairs cusp is onto in degree two
(AC26, Proposition~7.12, whose integral annulus proof is used here).
Explicitly, for any of the three complex curves $C$, choose
$y\in H_2(F;\mathbb Z)$ specializing to $[C]$. The descended
toric line bundle evaluates on $y$ as $-q(y)=\deg_L(C)$,
while its image in the exterior cyclic generator $g$ has
coefficient $2q(y)$. Thus, integrally modulo $d$,
\[
 [C]=-2\deg_L(C)\,g\quad\text{in }H_2(X_p;\mathbb Z).
\]
Since $2$ is a unit modulo $d$, the images of the three
complex curves have precisely the ratios $-6:0:1$ after
normalizing by the complex $\delta$-fixed curve. On the normalized $dP_6$ the
internal curve is $E_1$, the zero external curve is $E_2$, and the
fixed external curve is $E_3$. The opposite-pair relation is
$H=E_1+E_2+E_3$ after mapping to the cusp. Thus the restriction of
$u$ to every normalized component is
\[
 \eta=-5H+6E_1-E_3=-5(H-E_1)+(E_1-E_3)
 \quad\text{in }H^2(dP_6;\mathbb Z/d).
\]
Indeed its evaluations on $E_1,E_2,E_3$ are $-6,0,1$, and on $H$
the evaluation is $-5$. Hence $\eta^2=25-36-1=-12$.
The projection formula for the component fundamental cycle gives
$\langle u^2x_i,[\widetilde X_p]\rangle=-12$, also when that
component has a self-identification: its normalization pushes its
fundamental class with degree one.

\item Finally $u$ is an actual pullback from downstairs. Therefore
\[
 \langle u^3,[\widetilde X_p]\rangle
 =\langle v^3,\pi_*[\widetilde X_p]\rangle
 =d\langle v^3,[X_p]\rangle=0\quad\text{in }\mathbb Z/d.
\]
This evaluates the otherwise undetermined finite-coefficient cube.
The exact support-function, wall-degree and intersection computations
are executed in \texttt{code/manifolds\_mu.py}.
\end{enumerate}
\end{proof}

\begin{lemma}[the cubic detects the complete fundamental-class invariant]\label{maps:src:D.1.10}
\StatementRecord{maps:src:D.1.10}{proved}{Historical events: Claude Fable 5.2 (Anthropic), effort high, 16 November 2026: carried as background without fault found. GPT-6 Astra (OpenAI), ChatGPT Pro mode, effort Pro, S6\_09\_PAPERS, 17 November 2026: targeted proof check. Also checked in this writing. Pre-continuation writing events: GPT-6 Astra (OpenAI), ChatGPT Pro mode, effort Pro, S6\_10\_PAPERS, 18 November 2026: targeted proof check. Continuation metadata events: GPT-6 Astra (OpenAI), ChatGPT Pro mode, effort Pro, S6\_10\_PAPERS continuation, 19 November 2026: used and rechecked in this run; status-vocabulary reconciliation; metadata reconciliation only; no new mathematical check. \emph{Pre-continuation recorded proof scope.} GPT-6 Astra (OpenAI, ChatGPT Pro mode), effort Pro, 2026-11-17 UTC; Claude Fable 5.2 (Anthropic), effort high, 16 November 2026; also checked in this writing by GPT-6 Astra (OpenAI), ChatGPT Pro mode, effort Pro, 18 November 2026}{the hypotheses, notation and definitions stated here; \cref{maps:src:D.1.9}}


Let $A=\mathbb Z^r\oplus\mathbb Z/d$ with $\gcd(d,6)=1$.
The integral homology group of $K(A,2)$ in degree six is
\[
 H_6(K(A,2);\mathbb Z)\simeq
 \mathbb Z^{\binom{r+2}{3}}\oplus
 (\mathbb Z/d)^{\binom{r+1}{2}+r+1}.
\]
Its free coordinates are detected by integral degree-two cubics;
its torsion coordinates are detected by the finite-coefficient
monomials $u x_i x_j$ ($i\leq j$), $u^2x_i$ and $u^3$, where
$u$ evaluates to one on the cyclic summand of $A$. Consequently
\ref{maps:src:D.1.9} and the integral cubic determine the complete class
$\mu(\widetilde X_p)$ in Zhubr's theorem, not merely a quotient of it.
\end{lemma}


\begin{proof}
\begin{enumerate}[label=\arabic*.,leftmargin=*]
\item We give the low-degree calculation to exclude invisible torsion.
Put $R=\mathbb Z[1/6]$. In the path-loop Serre spectral sequence of
$K(\mathbb Z,3)$ the fibre is $K(\mathbb Z,2)=\mathbb CP^\infty$.
If $t$ has degree two, its transgression is the degree-three
generator $z$. Multiplicativity gives
$d_3(t^2)=2tz$ and $d_3(t^3)=3t^2z$.
Since two and three are units and the total path space is
contractible, induction in total degree forces the base cohomology
through degree seven to be the exterior algebra on $z$:
there is no surviving fibre class that could kill an additional
base class in degrees four through seven. Also $z^2=0$ because
$z$ is odd and two is a unit. This proves exactly the low-degree
base calculation needed below.

\item The coefficient extension
$0\to\mathbb Z\xrightarrow{d}\mathbb Z\to\mathbb Z/d\to0$
gives the fibration
\[
 K(\mathbb Z,2)\longrightarrow K(\mathbb Z/d,2)
 \longrightarrow K(\mathbb Z,3)
\]
with transgression $d_3(t)=dz$. Its cohomology through degree seven
is therefore computed by $R[t]\otimes\Lambda(z)$ with this
 differential. The groups in degrees three, five and seven are
$R/(d)$, $R/(2d)$ and $R/(3d)$, all $\mathbb Z/d$, and the positive even
groups through degree six vanish. There is no higher differential
in this range, by the base calculation. Universal coefficients,
starting in degree two, give
\[
 H_2=H_4=H_6=\mathbb Z/d,\qquad H_3=H_5=0.
\]
With coefficients $\mathbb Z/d$ the indicated transgressions vanish,
and $t,t^2,t^3$ survive; in particular the cube evaluates as a
generator of the dual of $H_6$. The same reasoning applies after
localization at every prime dividing $d$. At all other primes,
$K(\mathbb Z/d,2)$ is homologically trivial after localization,
as follows from its single nonzero homotopy group and the localized
Hurewicz theorem. This rules out additional two- or three-primary
classes omitted by passing to $R$.

\item Finally we use the product decomposition
\[
 K(A,2)=(\mathbb CP^\infty)^r\times K(\mathbb Z/d,2).
\]
The first factor has torsion-free homology concentrated in even
degrees, so the integral Kunneth formula has no Tor term here.
Its degree-six free part has rank $\binom{r+2}{3}$; tensoring its
degree-four, degree-two and degree-zero groups with the three
cyclic groups just computed gives respectively
$\binom{r+1}{2}$, $r$ and one cyclic summands. The stated products
evaluate on these distinct summands, and each is a perfect
detector in its coordinate. This proves the lemma.
\end{enumerate}
\end{proof}

The complete credited statement and re-derivation are in \cref{maps:src:D.1.11}.
\begin{theorem}[a canonical cyclic action and its quotient]\label{maps:src:D.1.12}
\StatementRecord{maps:src:D.1.12}{proved}{Historical events: Claude Fable 5.2 (Anthropic), effort high, 16 November 2026: carried as background without fault found. GPT-6 Astra (OpenAI), ChatGPT Pro mode, effort Pro, S6\_09\_PAPERS, 17 November 2026: targeted proof check. Also checked in this writing. Pre-continuation writing events: GPT-6 Astra (OpenAI), ChatGPT Pro mode, effort Pro, S6\_10\_PAPERS, 18 November 2026: targeted proof check. Continuation metadata events: GPT-6 Astra (OpenAI), ChatGPT Pro mode, effort Pro, S6\_10\_PAPERS continuation, 19 November 2026: used and rechecked in this run; status-vocabulary reconciliation; metadata reconciliation only; no new mathematical check. \emph{Pre-continuation recorded proof scope.} GPT-6 Astra (OpenAI, ChatGPT Pro mode), effort Pro, 2026-11-17 UTC; Claude Fable 5.2 (Anthropic), effort high, 16 November 2026; also checked in this writing by GPT-6 Astra (OpenAI), ChatGPT Pro mode, effort Pro, 18 November 2026}{the hypotheses, notation and definitions stated here; \cref{maps:src:D.1.7}; \cref{maps:src:D.1.4}; \cref{maps:src:D.1.11}}


There is an explicit free action $\sigma_d$ of $\mathbb Z/d$ on
$M_d$, defined by the affine and toric data below, such that every
admissible multiplicity-$(3,4)$ member with $|p|=d$ is
orientation-preservingly diffeomorphic to $M_d/\langle\sigma_d\rangle$.
All cyclic actions arising from these members are equivariantly
diffeomorphic to this action, allowing an automorphism of the
abstract deck group. Their quotients carry exactly the complex
structures and fixed-$u$ count in \ref{maps:src:D.1.7}.
\end{theorem}


\begin{proof}
\begin{enumerate}[label=\arabic*.,leftmargin=*]
\item Choose the specific original datum
$(\ell_0,\ell_1,\ell_2)=(0,-d,d)$ and scalar period $c=-i$.
It is admissible and has $p=d$. Its cover has lattice $\Lambda_d$
and matrices $A_1^{(d)},A_2^{(d)}$ of \ref{maps:src:D.1.4}. Use the respective
finite invariant translation vectors
\[
 v'_1=(-1,-2d,4d,0),\qquad v'_2=(1,3d,-3d,0),
\]
divided by three and four in the affine finite actions, and use
the cusp period lattice $\langle e_1,d e_2\rangle$. The invariant
meridian value is $p'=1$. This specifies a single smooth model,
independently of any input twist ambiguity; \ref{maps:src:D.1.11} identifies its
underlying manifold with $M_d$.

\item On a regular torus define $\sigma_d$ by translation through the
original $\widehat\gamma$, whose coordinate in $\Lambda_d$ is
$(1/d,0,0,0)$. The matrices $A_j^{(d)}$ change this vector by an
integral vector, so this translation is globally defined on the
regular family and commutes with the finite quotient actions.
It has order $d$. On a finite reduced fibre a fixed point of its
$a$th power would require, in the invariant $\gamma'$ coordinate,
$a/d\equiv\pm k/m_j\pmod1$ for some $k$. Since $\gcd(d,m_j)=1$,
this forces $d\mid a$. On the cusp it is the affine period
translation by $-e_2$, including the prescribed holomorphic unit,
and cycles the $d$ components. No nonzero power below $d$ fixes a
stratum: after subtracting a covering period it would translate a
finite triangle, edge or vertex to itself by a nonzero vector.
Thus it extends over the entire toric tube and the action is free.
The quotient restores the original lattice, twists $(-d,d)$ and
period $c=-i$, so is exactly that fixed representative.

\item \ref{maps:src:D.1.7} gives an oriented diffeomorphism from any other constructed
quotient with the same $d$ to this one. Lift it to universal covers.
Uniqueness of path lifting makes the lift intertwine deck
transformations by the induced automorphism of $\pi_1=\mathbb Z/d$.
This proves the claimed equivariant classification within the
entire specified construction by a direct proof, without assuming
a classification of arbitrary free actions on $M_d$. Different
$d$ are distinguished downstairs by fundamental group and upstairs
by Euler characteristic.

\item The fixed-$u$ count is two for the present multiplicities. In the
coordinates of Viaclovsky's higher-multiplicity sphere family,
the dictionary in \cref{push:src:E.1.1} and the classification in \cref{fam:src:B.3.1} give $2k\ell$ at fixed $u$ for the pair $(3k,4\ell)$. The tangent computation in \cref{push:src:E.2.1} gives $(1,1,1,0)$, and \cref{push:src:E.2.2} determines the Kodaira dimension by the sign of
$k\ell-3k-4\ell$ for those sphere members; at $(k,\ell)=(1,1)$
this specializes to the count two and negative Kodaira dimension.
These sphere results concern $d=1$ and do not assert tangent
cohomology for the torsion covers $M_d$ with $d>1$.
\end{enumerate}
\end{proof}

\subsection{Rational elliptic starting configurations and exact computations}
The complete credited statement and re-derivation are in \cref{maps:src:D.2.1}.
\begin{computed}[as a bounded exact computation]\label{maps:src:D.2.2}
\StatementRecord{maps:src:D.2.2}{proved}{Historical events: Claude Fable 5.2 (Anthropic), effort high, 16 November 2026: carried as background without fault found. Claude Fable 5.2 (Anthropic), effort high, 17 November 2026: carried as background without fault found. GPT-6 Astra (OpenAI), ChatGPT Pro mode, effort Pro, S6\_09\_PAPERS, 17 November 2026: targeted proof check. Also checked in this writing. Pre-continuation writing events: GPT-6 Astra (OpenAI), ChatGPT Pro mode, effort Pro, S6\_10\_PAPERS, 18 November 2026: computation or formal certificate. Continuation metadata events: GPT-6 Astra (OpenAI), ChatGPT Pro mode, effort Pro, S6\_10\_PAPERS continuation, 19 November 2026: used and rechecked in this run; status-vocabulary reconciliation; metadata reconciliation only; no new mathematical check. \emph{Pre-continuation recorded proof scope.} GPT-6 Astra (OpenAI, ChatGPT Pro mode), effort Pro, 2026-11-17 UTC; Claude Fable 5.2, effort high, 2026-11-17; GPT-6 Astra (OpenAI, ChatGPT Pro mode), S6\_06\_BLIND, effort Pro, 2026-11-17; Claude Fable 5.2 (Anthropic), effort high, 16 November 2026; exact archived inputs and replies checked in this writing by GPT-6 Astra (OpenAI), ChatGPT Pro mode, effort Pro, 18 November 2026}{the hypotheses, notation and definitions stated here; \cite{manExplorer26}; \cref{maps:src:D.2.1}; \cref{mapsD:repair:collision-realization}}


Fifty specified inputs were evaluated in the public exact integral-cochain engine
\cite{manExplorer26}. In every configuration of \ref{maps:src:D.2.1} we take
$P=G$, $Q=nG$, one simple zero at the first $I_1$ in the stored
marking (or at the smallest multiplicative fibre if no $I_1$ exists),
and independently take at each additive fibre either its original
filling or the invariant log vector $(0,1/m)$, where $m$ is the
minimal semistable-reduction order. Within this exact slice the
engine recognizes 24 simply connected inputs and 26 non-simply-connected
inputs. Exactly three have the integral cohomology of $S^6$:
OS43 with $III^*+3I_1$ and the order-four modification;
OS43 with $III^*+II+I_1$ and only the $III^*$ modification;
OS49 with $IV^*+III+I_1$ and both modifications. The executable is
\path{code/manifolds_explorer_21.py}.

Every exact input, fundamental-group result, and group in degrees
zero through six is retained in
\texttt{records/manifolds\_explorer\_21.csv}, its JSON companion,
and the individual raw responses. The homology groups are the
cohomology groups in reverse order, by oriented Poincare duality.
The integer operations use Smith forms, not floating-point ranks.
The public implementation checks endpoint groups, Poincare and
torsion duality, and agreement of abelianized fundamental group
with the universal coefficient theorem. Re-running the script
with its recorded inputs produces the table; cached responses
retain their URL and API version~2.

The proposition here concerns the executed computation, with these
bounds. Its interpretation as topology of each smooth model uses
the author's supplied marked local-model database. The public scope
document explicitly leaves detailed geometric validation certificates
in the author's research workspace. We therefore do not promote this
run into an independent proof of all its geometric attachment maps,
nor into the unbounded classification asked for in D.2.

For the archived computation, all fifty exact request payloads and
saved API-v2 replies were regenerated and compared, giving the stated
$24$ simply connected outputs, $26$ other outputs and three outputs
with the integral cohomology of a sphere; the eight specified clutch
comparisons were checked as well.
This is an audit of the retained requests and responses.  The source audit records a timeout of its fresh remote request. In the present writing the fifty saved requests and API-v2 replies, and the eight integer-clutch comparisons, were regenerated and matched exactly by \path{code/III/audit_explorer_archive.py}. This is the stated bounded archive audit.  The independent geometric input for the two
collision profiles is \cref{mapsD:repair:collision-realization}.


\end{computed}


\begin{proposition}[Euler checks and parameter sensitivity]\label{maps:src:D.2.3}
\StatementRecord{maps:src:D.2.3}{proved}{Historical events: Claude Fable 5.2 (Anthropic), effort high, 16 November 2026: carried as background without fault found. GPT-6 Astra (OpenAI), ChatGPT Pro mode, effort Pro, S6\_09\_PAPERS, 17 November 2026: targeted proof check. Also checked in this writing. Pre-continuation writing events: GPT-6 Astra (OpenAI), ChatGPT Pro mode, effort Pro, S6\_10\_PAPERS, 18 November 2026: targeted proof check. Continuation metadata events: GPT-6 Astra (OpenAI), ChatGPT Pro mode, effort Pro, S6\_10\_PAPERS continuation, 19 November 2026: used and rechecked in this run; status-vocabulary reconciliation; metadata reconciliation only; no new mathematical check. \emph{Pre-continuation recorded proof scope.} GPT-6 Astra (OpenAI, ChatGPT Pro mode), effort Pro, 2026-11-17 UTC; Claude Fable 5.2 (Anthropic), effort high, 16 November 2026; also checked in this writing by GPT-6 Astra (OpenAI), ChatGPT Pro mode, effort Pro, 18 November 2026}{the hypotheses, notation and definitions stated here; the explicit constructions and calculations in the proof}


For any of the semistable triangular rank-two Mumford fillings, its
reduced central fibre has Euler number twice the number of vertices
in the dual torus triangulation. A rank-one wheel over an elliptic
curve, a smooth torus, and a free finite quotient of a torus have
Euler number zero. Consequently the total Euler number is the sum
of these rank-two contributions whenever all other fillings are of
these types. The fibre configuration and multiplicities alone do
not specify the input topology: integer clutching must also be kept.
\end{proposition}


\begin{proof}
\begin{enumerate}[label=\arabic*.,leftmargin=*]
\item Inclusion--exclusion over the toric strata counts the zero-dimensional
strata, since every positive-dimensional complex torus has Euler
number zero. These strata correspond to triangles. A torus
triangulation satisfies $3F=2E$ and $V-E+F=0$, hence $F=2V$.
A rank-one wheel is stratified by bundles over elliptic curves, and
each stratum has Euler number zero. A finite torus quotient has Euler
number zero by multiplicativity of Euler number for covers.
The regular locus of the total fibration also has Euler number zero;
Euler additivity gives the claim, applying the localization argument to the stated cells.

\item For the last assertion, even the finite computation needs the
full rational vector, not its class modulo the lattice. A smooth
weight-zero log vector $\theta$ of denominator $m$ has boundary
relation $\mathrm{meridian}^{m}=\mathrm{fibre}^{m\theta}$.
Replacing $\theta$ by $\theta+v$ for an integral $v$ changes that
relation and its marked cohomological boundary map. The executed
semistable examples OS45,47,55,56 have $\pi_1=\mathbb Z$ for the
zero clutch and trivial fundamental group for the recorded primitive
integer clutch, while the fibre multiplicity is one in both cases.
These four exact comparisons are saved individually in
\texttt{records/manifolds\_explorer\_runs/}.
\end{enumerate}
\end{proof}

The exact remaining classification questions and all established subcases are in \cref{maps:src:D.2.4}.
\begin{theorem}[arbitrary multiplicities in the original configuration]\label{maps:src:D.2.5}
\StatementRecord{maps:src:D.2.5}{proved}{Historical events: Claude Fable 5.2 (Anthropic), effort high, 16 November 2026: carried as background without fault found. GPT-6 Astra (OpenAI), ChatGPT Pro mode, effort Pro, S6\_09\_PAPERS, 17 November 2026: targeted proof check. Also checked in this writing. Pre-continuation writing events: GPT-6 Astra (OpenAI), ChatGPT Pro mode, effort Pro, S6\_10\_PAPERS, 18 November 2026: targeted proof check. Continuation metadata events: GPT-6 Astra (OpenAI), ChatGPT Pro mode, effort Pro, S6\_10\_PAPERS continuation, 19 November 2026: used and rechecked in this run; status-vocabulary reconciliation; metadata reconciliation only; no new mathematical check. \emph{Pre-continuation recorded proof scope.} GPT-6 Astra (OpenAI, ChatGPT Pro mode), effort Pro, 2026-11-17 UTC; Claude Fable 5.2 (Anthropic), effort high, 16 November 2026; also checked in this writing by GPT-6 Astra (OpenAI), ChatGPT Pro mode, effort Pro, 18 November 2026}{the hypotheses, notation and definitions stated here; \cite[\S2.1, (2.11)--(2.12)]{manViaclovskyLogs26}; \cite[Propositions 4.1, 4.3, 4.5]{manViaclovskyLogs26}; \cref{maps:src:D.3.2}}


Consider the $IV^*,III,I_1$ configuration with its rank-two Mumford
filling, and the logarithmic data of
\cite[\S2.1, (2.11)--(2.12)]{manViaclovskyLogs26}:
\[
 m_0=3k,\qquad m_1=4\ell,\qquad
 v_i=r_i\delta+n_i u_i,\qquad
 \gcd(n_i,m_i)=1,\qquad k,\ell>0.
\]
Here $r_i,n_i$ are arbitrary integers, not residue classes. Put
$P=m_1n_0+m_0n_1$. There is no upper bound on these parameters.
If $P\ne0$, the resulting smooth compact threefold has
\[
 \pi_1(X)=\mathbb Z/|P|,\qquad
 H_j(X;\mathbb Z)=
 \begin{cases}
 \mathbb Z&j=0,6,\\
 \mathbb Z/|P|&j=1,2,3,4,\\
 0&j=5.
 \end{cases}
\]
If $P=0$, then $\pi_1(X)=\mathbb Z$ and all integral homology
is free, with Betti numbers $(1,1,2,2,2,1,1)$.
In particular, in this admissible family $X$ is simply connected
exactly when $|P|=1$, and then it is diffeomorphic to $S^6$.
Such sphere data exist exactly when $\gcd(3k,4\ell)=1$.
This proves a parametric table for every multiplicity and every
integer twist in this family; it does not impose the determinant-one
hypothesis when computing the other manifolds.
\end{theorem}


\begin{proof}
\begin{enumerate}[label=\arabic*.,leftmargin=*]
\item We give the integral reduction explicitly, including the step that
extends beyond the simply connected hypothesis in
\cite[Propositions 4.1, 4.3, 4.5]{manViaclovskyLogs26}.
In the integral basis $(\delta,a,b,d)$, the three homological
monodromies are
\[
 A_0=\begin{pmatrix}1&-1&0&2\\0&-1&-1&6\\0&1&0&0\\0&0&0&1\end{pmatrix},\quad
 A_1=\begin{pmatrix}1&0&1&-3\\0&0&1&0\\0&-1&0&6\\0&0&0&1\end{pmatrix},\quad
 A_\infty=\begin{pmatrix}1&0&0&1\\0&1&1&0\\0&0&1&0\\0&0&0&1\end{pmatrix}.
\]
They satisfy $A_0A_1A_\infty=I$, $A_0^3=A_1^4=I$,
$u_0=d+2a+2b$, $u_1=d+3a+3b$ and
$\Lambda^{A_i}=\mathbb Z\delta\oplus\mathbb Zu_i$ for $i=0,1$.
The good reduction followed by the degree-$k$ or degree-$\ell$
base change remains smooth. The affine generator has translation
$v_i/m_i$ and rotates the base by $e^{2\pi i/m_i}$.
The invariant covector $d^\vee$ changes by $n_i/m_i$;
$\gcd(n_i,m_i)=1$ therefore proves that every nonidentity power
is free even on the central torus. Quotienting gives a smooth
filling retracting onto a bielliptic surface $F_i$.
The logarithm translation used in the punctured-disc gluing has
normal-circle relation $x_i^{m_i}=v_i$. These are precisely the
charts and gluing of Proposition 2.3 of the cited paper; the
freeness argument uses no relation between the two multiplicities.

\item The Mumford filling kills $\delta,a,x_\infty$ in the fundamental
group. The conjugation relation $x_1ax_1^{-1}=b^{-1}$ then kills
$b$, and $x_1=x_0^{-1}$. Thus the full group is already abelian,
generated by $d$ and $x=x_0$, with relation matrix
\[
 R=\begin{pmatrix}-n_0&m_0\\-n_1&-m_1\end{pmatrix}.
\]
Its entries have greatest common divisor one because
$\gcd(n_0,m_0)=1$, and $\det R=P$. Its Smith form is therefore
$(1,|P|)$ when $P\ne0$, or $(1,0)$ when $P=0$.
This is a computation of the group, not just its abelianization.

\item Let $V$ be the inverse image of the thrice-punctured sphere, and
$M_i=\partial N_i$ the boundary of each filling. All $H_1(M_i)$
are free of rank three: the individual coinvariants give
$b=-a,\delta=-3a$ at $0$, and $b=-a,\delta=-2a$ at $1$;
the cusp has primitive rank-two vanishing lattice.
The common coinvariants give $H_1(V)=\mathbb Z^3$.
Independently of $P$, the finite central surfaces have
\[
 H_1(F_0)=\mathbb Z\langle a,d,g_0\rangle/
   \langle3r_0a-n_0d+m_0g_0\rangle,
\quad
 H_1(F_1)=\mathbb Z\langle a,d,g_1\rangle/
   \langle2r_1a-n_1d+m_1g_1\rangle.
\]
These are free of rank two. Choose
$m_i\alpha_i+n_i\beta_i=1$ and put
$z_i=\alpha_i d+\beta_i g_i$; $(a,z_i)$ is an integral basis.
Also $H_1(F_\infty)=\mathbb Z^2$.
Consequently the degree-one Mayer--Vietoris map is a map between
free rank-nine groups with cokernel $H_1(X)$.
If $P\ne0$ its kernel is zero, even though its cokernel may be
nonzero. If $P=0$ its kernel is free of rank one.
This replaces the simply connected assumption in the cited
Proposition 4.1 by the actual rank computation.

\item The map $H_2(M_i)\to H_2(N_i)$ is onto for every filling.
Indeed its cokernel maps to
$H_2(N_i,M_i)\simeq H^4(F_i)=\mathbb Z$ by intersection
with $F_i$; the multiple $m_iF_i$ is principal, so that
intersection is zero. By the pair sequence, its kernel consists
of the normal-circle tori of classes in $H^3(F_i)$.
The homology Serre sequence over the graph underlying the base gives
\[
 0\longrightarrow (\bigwedge^2\Lambda)_{A_0,A_1}
 \longrightarrow H_2(V)\longrightarrow
 \ker\bigl[(A_0-I)\ (A_1-I)\bigr]\longrightarrow0.
\]
The left group is $\mathbb Z[a\wedge b]$: the exterior relations
give $\delta\wedge a=\delta\wedge b=a\wedge d=b\wedge d=0$
and $\delta\wedge d=-6a\wedge b$.
The invariant form
$\xi=a^\vee\wedge b^\vee+6d^\vee\wedge\delta^\vee$
evaluates to one on $a\wedge b$, proving primitivity.
The right group is free of rank five.
The two cusp normal tori are the sweeps of $\delta,a$.
In the graph kernel they are
$(-\delta,-\delta)$ and $(b,-a)$; quotienting by them leaves
the integral basis $(\delta,0),(u_0,0),(0,u_1)$.
For verification, the equation in the $A_0$ coinvariants forces
$b^\vee(v_1)=3d^\vee(v_1)$; subtracting the appropriate cusp
vector then puts $v_1$ in $\mathbb Zu_1$ and $v_0$ in
$\mathbb Z\delta\oplus\mathbb Zu_0$. Uniqueness follows
by comparing the second components.
Thus, after the cusp relations, the target for the remaining
normal tori is a free rank-four group with basis $(s,t,t_0,t_1)$,
where $s=[a\wedge b]$ and the last three classes are these sweeps.

\item For a loop $\lambda$ in the finite covering torus, the normal
torus cuts the reference fibre in
$\mathcal N_i\lambda$, where
$\mathcal N_i=\sum_{j=0}^{m_i-1}A_i^j$, and its $s$ coordinate
is $-\xi(\lambda,v_i)$. The normal torus of $g_i$ has
fibre cut $v_i$ and $s$ coordinate zero; both follow by applying
the homomorphisms to $m_i g_i=v_i$ and cancelling in free groups.
The norms are
\[
 \mathcal N_0a=-k\delta,\quad \mathcal N_0d=m_0u_0,
 \qquad
 \mathcal N_1a=-2\ell\delta,\quad \mathcal N_1d=m_1u_1.
\]
Moreover $\xi(a,u_0)=2$, $\xi(a,u_1)=3$ and
$\xi(d,\delta)=6$. Since $(0,\delta)=-(\delta,0)$ modulo
the cusp relations, the normal-torus map in the bases
$(a,z_0,a,z_1)$ and $(s,t,t_0,t_1)$ is
\[
 B=\begin{pmatrix}
 -2n_0&-6r_0\alpha_0&-3n_1&-6r_1\alpha_1\\
 -k&r_0\beta_0&2\ell&-r_1\beta_1\\
 0&1&0&0\\
 0&0&0&1
 \end{pmatrix}.
\]
Removing the two unit pivots reduces its cokernel to that of
\[
 C=\begin{pmatrix}-2n_0&-3n_1\\-k&2\ell\end{pmatrix},
 \qquad \det C=-P.
\]
The entries of $C$ have greatest common divisor one: any common
divisor divides $\gcd(k,2n_0)$ and hence divides $2$, whereas
$n_1$ is odd because $\gcd(n_1,4\ell)=1$.
Thus $\operatorname{coker}B=\mathbb Z/|P|$ for $P\ne0$
and $\mathbb Z$ for $P=0$. Mayer--Vietoris gives $H_2(X)$
equal to this cokernel for $P\ne0$; for $P=0$ it is an extension
of $\mathbb Z$ by $\mathbb Z$, hence $\mathbb Z^2$.

\item Every fibre except the rank-two cusp has Euler number zero, so
$\chi(X)=2$. Integral Poincare duality and the universal
coefficient theorem give $H_5\simeq\operatorname{Hom}(H_1,\mathbb Z)$,
$H_4\simeq\operatorname{Hom}(H_2,\mathbb Z)
\oplus\operatorname{Ext}(H_1,\mathbb Z)$, and
$\operatorname{Tor}H_3\simeq\operatorname{Tor}H_2$.
The Euler characteristic determines the remaining rank, giving
all stated groups. For $|P|=1$, sphere recognition follows as in
\ref{maps:src:D.3.2} from Hurewicz, Whitehead and $\Theta_6=0$.
Finally the Bezout equation $m_1n_0+m_0n_1=\pm1$ is solvable
exactly when $\gcd(m_0,m_1)=1$, and it automatically implies both
local admissibility conditions.

\item The unrestricted exact implementation is the script
\texttt{manifolds\_unbounded\_logs.py} in \texttt{code/}. It symbolically checks
the determinant and norm formulas and actually computes the Smith
forms on the stated finite validation grid; its all-parameter
correctness is the integral argument above. For example,
$(m_0,m_1;n_0,n_1)=(9,4;-2,1)$ gives $P=1$ and $S^6$;
$(9,4;1,1)$ gives all four torsion groups $\mathbb Z/13$;
$(12,12;1,-1)$ gives the stated infinite-cyclic case.
\end{enumerate}
\end{proof}

The logarithmic quotient and its marked local filling are those of \cite[Proposition~2.3]{manViaclovskyLogs26}; the original integral homology case is \cite[Proposition~4.1]{manViaclovskyLogs26}. The unrestricted determinant and torsion calculation above retains the complete integer boundary maps.

\begin{theorem}[infinitely many orders for $III^*+3I_1$]\label{maps:src:D.2.6}
\StatementRecord{maps:src:D.2.6}{proved}{Historical events: Claude Fable 5.2 (Anthropic), effort high, 16 November 2026: carried as background without fault found. GPT-6 Astra (OpenAI), ChatGPT Pro mode, effort Pro, S6\_09\_PAPERS, 17 November 2026: targeted proof check. Also checked in this writing. Pre-continuation writing events: GPT-6 Astra (OpenAI), ChatGPT Pro mode, effort Pro, S6\_10\_PAPERS, 18 November 2026: targeted proof check. Continuation metadata events: GPT-6 Astra (OpenAI), ChatGPT Pro mode, effort Pro, S6\_10\_PAPERS continuation, 19 November 2026: used and rechecked in this run; status-vocabulary reconciliation; metadata reconciliation only; no new mathematical check. \emph{Pre-continuation recorded proof scope.} GPT-6 Astra (OpenAI, ChatGPT Pro mode), effort Pro, 2026-11-17 UTC; Claude Fable 5.2 (Anthropic), effort high, 16 November 2026; also checked in this writing by GPT-6 Astra (OpenAI), ChatGPT Pro mode, effort Pro, 18 November 2026}{the hypotheses, notation and definitions stated here; \cite[Proposition 3.1, Lemma 3.2, Proposition 4.2]{manViaclovskyTwo26}; \cref{maps:src:D.2.5}}


Use the corrected $III^*,I_1,I_1,I_1$ construction of \cref{push:src:E.5.2} and \cref{push:src:E.5.3}, with
two primitive rank-one Mumford fillings and one rank-two hexagonal
filling. In the ordinary homology basis $(e_0,e_1,e_2,e_3)$ dual
to $(\psi,u,w,\delta)$, put $\ell=e_0-e_1-e_2$.
At $III^*$ one may make a logarithmic transformation with
\[
 m=4k,\qquad v=n\ell+r e_3,\qquad
 k>0,\quad n,r\in\mathbb Z,\quad \gcd(n,4k)=1.
\]
For every such datum the resulting compact complex threefold has
\[
 \pi_1(X)=\mathbb Z/|n|,\qquad
 H_0=H_6=\mathbb Z,\quad H_5=0,\quad
 H_1=H_2=H_3=H_4=\mathbb Z/|n|.
\]
Consequently $n=1$ or $n=-1$ gives a complex structure on $S^6$
for every $k>0$, with no coprimality restriction on $k$.
Thus this configuration answers Viaclovsky's Question 4
affirmatively. This assertion concerns the displayed analytic
fillings and integer twists, not all modifications of that surface.
\end{theorem}


\begin{proof}
\begin{enumerate}[label=\arabic*.,leftmargin=*]
\item The geometric input is the construction in \cref{push:src:E.5.2} and \cref{push:src:E.5.3}, with the original geometric statements in
\cite[Proposition 3.1, Lemma 3.2, Proposition 4.2]{manViaclovskyTwo26}.
The order-four good-reduction family is smooth and proper and its
punctured comparison preserves the distinguished section.
Pull it back by a degree-$k$ map of the covering disc, so that the
original base coordinate is $s^{4k}$. The linear deck action still
has order four on the torus lattice. It fixes $\ell$ and $e_3$.
Translation by $v/(4k)$, combined with rotation of $s$, defines
an order-$4k$ action. On the central torus, the invariant map
$\psi:T^4\to\mathbb R/\mathbb Z$ changes by $n/(4k)$;
coprimality proves freeness. The punctured comparison is changed
by the holomorphic translation
$[\log(s)v(s)/(2\pi i)]$, exactly as in \ref{maps:src:D.2.5}.
The quotient is therefore smooth and glues to the unchanged
three Mumford pieces. Properness follows from properness of the
four fillings and compactness of the base. Its meridian relation
is $g^m=v$. This constructs the claimed analytic threefold for
every displayed parameter, rather than extrapolating a finite table.

\item Here is the integral topology. The cohomological matrices, in the
stated marking, are
\[
 \begin{aligned}
 T_0&=\begin{pmatrix}1&0&-2&-1\\0&0&-1&0\\0&1&0&-1\\0&0&0&1\end{pmatrix},
 &T_1&=\begin{pmatrix}1&0&2&2\\0&1&1&1\\0&0&1&0\\0&0&0&1\end{pmatrix},\\
 T_2&=\begin{pmatrix}1&0&0&0\\0&1&0&0\\0&-1&1&0\\0&0&0&1\end{pmatrix},
 &T_3&=\begin{pmatrix}1&0&0&-1\\0&1&1&0\\0&0&1&0\\0&0&0&1\end{pmatrix}.
 \end{aligned}
\]
Let $A_i=T_i^{-t}$. The primitive vanishing subgroups of the
three Mumford pieces are respectively
$\mathbb Z(e_2+e_3)$, $\mathbb Ze_1$, and
$\mathbb Ze_2\oplus\mathbb Ze_3$.
They kill $e_1,e_2,e_3$. Their meridians are killed by the
extending section, and the sphere relation then kills the finite
meridian as well. The finite relation $g^m=v$ leaves exactly
$ne_0=0$. All boundary fundamental-group maps are onto, so
no generators have been omitted. This proves $\pi_1(X)=\mathbb Z/|n|$.

\item We justify the rank-one normal tori used in the degree-two
calculation. In a primitive rank-one Mumford filling, the
normalization is a flat $\mathbb P^1$-bundle over an elliptic
curve, with its two sections identified by a translation.
Topologically trivialize the flat line bundle and isotope the
translation to the identity, keeping the extending section fixed.
The collar pair is then the product of the primitive nodal
elliptic filling with $T^2$; its local equation is $xy=q$.
It follows that $H^3$ of its central fibre is free of rank three.
The normal tori are the sweeps of all three integral invariant
loops: the vanishing circle, whose rotation extends over $xy=q$,
and the two loops of the remaining elliptic factor. They are
boundaries of relative three-chains obtained from the section
disc and those rotations or translations. The unit-factor
deformation leaves the marked boundary and section comparison
unchanged up to isotopy. For the rank-two filling, the same
argument gives the two invariant angular loops and $H^3=\mathbb Z^2$.
Thus these are integral normal-torus classes; retaining only the
vanishing loop in a rank-one piece would omit two relations.

\item For completeness we give primitive-minor certificates for the
global lattice calculation. Set
\[
 D=\bigl[(A_0-I)\ (A_1-I)\ (A_2-I)\bigr],
 \qquad
 j_3(x)=-(A_1A_2x,A_2x,x).
\]
For $i=0,1,2$ let $j_i(x)$ place $x$ in block $i$ of
$\Lambda^3$. The fixed lattices at the three Mumford points have
the respective integral bases
\[
 (-e_0+2e_1,e_2,e_3),\qquad (e_0,e_1,e_3),\qquad(e_2,e_3).
\]
Form a $12\times9$ matrix $K$ whose first eight columns are
their images under $j_1,j_2,j_3$, in that order, and whose ninth
column is $j_0(\ell)$. Direct multiplication gives $DK=0$ and
$\operatorname{rank}D=3$. The minor of $K$ using rows
\[
 (1,2,4,5,7,8,9,10,12)
\]
is $-1$, with rows numbered from one. Hence its columns form
an integral basis of $\ker D$. Moreover
$j_0(e_3)=-K_3-K_6-K_8$, so the finite $e_3$ sweep vanishes
after the Mumford relations.

\item In the exterior basis
$(e_0e_1,e_0e_2,e_0e_3,e_1e_2,e_1e_3,e_2e_3)$,
the matrix $[\bigwedge^2A_i-I]_{i=0}^3$ has primitive rank five:
rows $(1,2,3,5,6)$ and columns $(1,2,4,5,7)$ have minor $-1$.
Its cokernel is detected by
$\xi=u\wedge w+2\psi\wedge\delta$, whose value on $e_1e_2$
is one. Thus the torus-bundle Serre sequence gives
$H_2(V)=\mathbb Z^{10}$, generated by this fibre class and
the nine sweep classes defining $K$. Each sweep has zero
integral against the flat form $\xi$; these classes consequently
fix an integral splitting, not only a rational one.
The eight Mumford normal tori span a primitive subgroup.
Their domain has rank eight, and the preceding local construction
realizes all eight columns, so their image is exactly that subgroup.
The quotient is
\[
 Q=\mathbb Z\langle s,t\rangle,
 \qquad s=[e_1\wedge e_2],\quad t=[T_{\ell,0}].
\]

\item Abelianization of the finite fibre group gives
$e_2=-e_1$, $e_3=2e_1$ and
$mg=ne_0+2r e_1$. Since $\gcd(n,m)=1$, it is free of rank
two with basis $e_1,z$, where
$z=\alpha e_0+\beta g$ and $m\alpha+n\beta=1$.
The normal-circle formulas proved in \ref{maps:src:D.2.5} apply. Here
\[
 \mathcal N_me_1=2k e_3,\qquad
 \mathcal N_me_0=m\ell,\qquad
 \xi(e_1,v)=-n,\qquad \xi(e_0,v)=2r.
\]
The first norm dies in $Q$. The normal tori of $e_1,e_0,g$
therefore have coordinates $(n,0)$, $(-2r,m)$ and $(0,n)$.
In the basis $e_1,z$, the remaining attachment matrix is
\[
 \begin{pmatrix}n&-2r\alpha\\0&1\end{pmatrix}.
\]
It has Smith form $(1,|n|)$ for every $k$ and $r$.

\item To check that this cokernel is actually $H_2(X)$, the boundary
$H_1$ groups have ranks $3,4,4,3$ and are free; the target
of their Mayer--Vietoris map has rank
$4+2+3+3+2=14$ and is also free. Its cokernel is the finite
group $H_1(X)=\mathbb Z/|n|$, so its kernel is zero.
Every central fibre is irreducible and a multiple of it is
principal; the pair-intersection argument in \ref{maps:src:D.2.5} makes each
$H_2(M_i)\to H_2(N_i)$ onto. Thus $H_2(X)$ is precisely
the displayed normal-torus cokernel. Euler localization gives
$\chi(X)=2$, and integral duality gives all remaining groups
as in \ref{maps:src:D.2.5}. Sphere recognition for $|n|=1$ is the same argument.

\item The script \texttt{code/manifolds\_iii\_star.py} executes all
stated lattice certificates, the norm identities, and 312 actual
$10\times10$ Smith forms for $1\leq k\leq12$,
$-15\leq n\leq15$, $\gcd(n,4k)=1$, and $r=0,3$.
The proof is independent of that finite validation bound.
\end{enumerate}
\end{proof}

\begin{theorem}[all orders $3k$ for Wang's primitive $I_2$ filling]\label{maps:src:D.2.7}
\StatementRecord{maps:src:D.2.7}{proved}{Historical events: Claude Fable 5.2 (Anthropic), effort high, 16 November 2026: carried as background without fault found. GPT-6 Astra (OpenAI), ChatGPT Pro mode, effort Pro, S6\_09\_PAPERS, 17 November 2026: targeted proof check. Also checked in this writing. Pre-continuation writing events: GPT-6 Astra (OpenAI), ChatGPT Pro mode, effort Pro, S6\_10\_PAPERS, 18 November 2026: targeted proof check. Continuation metadata events: GPT-6 Astra (OpenAI), ChatGPT Pro mode, effort Pro, S6\_10\_PAPERS continuation, 19 November 2026: used and rechecked in this run; status-vocabulary reconciliation; metadata reconciliation only; no new mathematical check. \emph{Pre-continuation recorded proof scope.} GPT-6 Astra (OpenAI, ChatGPT Pro mode), effort Pro, 2026-11-17 UTC; Claude Fable 5.2 (Anthropic), effort high, 16 November 2026; also checked in this writing by GPT-6 Astra (OpenAI), ChatGPT Pro mode, effort Pro, 18 November 2026}{the hypotheses, notation and definitions stated here; \cite{manWang26}; \cref{maps:src:D.2.2}; \cref{maps:src:D.2.6}; \cref{maps:src:D.2.5}}


Use the rational elliptic surface
\[
 y^2=x^3-3(1+t)x+(2+3t+2t^2),
 \qquad P=(1,\sqrt2\,t),\qquad Q=6P.
\]
Its fibres are $IV^*,I_2,I_1,I_1$, and the degree-one
linearization vanishes at the second $I_1$ point. Use Wang's
primitive one-component rank-one Mumford model at $I_2$, the
rank-one model at the other $I_1$, and the hexagonal rank-two
model at the zero of the linearization. In the homology marking
dual to $(\psi,u,w,\delta)$ put
\[
 \ell=e_0-2e_1-4e_2,\qquad
 m=3k,\qquad v=n\ell+r e_3,
 \qquad k>0,\quad \gcd(n,3k)=1.
\]
The order-$m$ logarithmic quotient at $IV^*$ is smooth for every
integer $r$, and the resulting threefold satisfies
\[
 \begin{gathered}
 \pi_1(X)=\mathbb Z/|n|,\qquad H_0=H_6=\mathbb Z,
 \qquad H_5=0,\\
 H_1=H_4=\mathbb Z/|n|,\qquad
 H_2=H_3=\mathbb Z/(2|n|).
 \end{gathered}
\]
In particular every $k>0$, with $n=\pm1$, gives a simply
connected complex rational homology six-sphere with two-torsion.
The case $(k,n,r)=(1,1,0)$ independently reproduces Wang's
homology theorem. The other orders are included by the proof,
without a finite parameter bound.
\end{theorem}


\begin{proof}
\begin{enumerate}[label=\arabic*.,leftmargin=*]
\item The source is \cite{manWang26}: Propositions 2.1 and
2.7--2.10 verify the surface and local models, and
Sections 4--5 fix their markings. In particular the $I_2$ period
exponent is $(2,1)$, not $(2,0)$. The unimodular change
$(x,y)=(v,v^2/z)$ makes it $(1,0)$ and gives a single
elliptic ruled normalization. This is the geometric input
missing from the zero-weight ``original $I_2$'' API run in \ref{maps:src:D.2.2}.
It is a different compactification, so the two computations
must not be equated.

\item Pull the smooth cubic good-reduction family back by degree $k$.
The order-three linear monodromy fixes $\ell,e_3$.
Combine base rotation of order $3k$ with translation by
$(n\ell+r e_3)/(3k)$, and use the section-preserving punctured
comparison and logarithmic clutching from \ref{maps:src:D.2.6}. The map $\psi$
changes by $n/(3k)$, so the action is free. This supplies the
analytic construction, properness and the marked relation
$g^{3k}=v$ for all the asserted data.

\item For the integral computation, use the fibre order
$(IV^*,I_2,I_1,I_1)$ and cohomological matrices
\[
 \begin{aligned}
 T_0&=\begin{pmatrix}1&0&-6&-2\\0&-1&-1&1\\0&1&0&-1\\0&0&0&1\end{pmatrix},
 &T_1&=\begin{pmatrix}1&0&6&3\\0&1&2&1\\0&0&1&0\\0&0&0&1\end{pmatrix},\\
 T_2&=\begin{pmatrix}1&0&0&0\\0&2&1&0\\0&-1&0&0\\0&0&0&1\end{pmatrix},
 &T_3&=\begin{pmatrix}1&0&0&-1\\0&1&1&0\\0&0&1&0\\0&0&0&1\end{pmatrix}.
 \end{aligned}
\]
Set $A_i=T_i^{-t}$. Their product is the identity and
$A_0^3=I$. The first two primitive vanishing directions are
$2e_2+e_3$ and $e_1+e_2$; the last filling kills $e_2,e_3$.
These generate $\langle e_1,e_2,e_3\rangle$. As in \ref{maps:src:D.2.6},
the extending sections kill all three Mumford meridians and
the sphere relation kills the fourth. The remaining finite
relation is $ne_0=0$, which proves the actual fundamental-group
assertion. Before the finite filling, the same calculation
gives $\pi_1(X_D)=\mathbb Z\langle e_0\rangle$ for the
inverse image $X_D$ of a disc containing the three Mumford points.

\item Use $D$ and the peripheral maps $j_i$ from \ref{maps:src:D.2.6}. The three
Mumford invariant lattices have bases
\[
 (-e_0+3e_1,e_2,e_3),\qquad
 (e_0,e_1+e_2,e_3),\qquad(e_2,e_3).
\]
Let $K$ consist of their eight peripheral columns, followed by
$j_0(\ell)$. The submatrix on rows
\[
 (1,2,4,5,7,8,9,10,12)
\]
has determinant $-1$. Since $DK=0$ and $\operatorname{rank}D=3$, these columns
are a basis of the entire integral kernel. Also
$j_0(e_3)=-K_3-K_6-K_8$.
For the exterior matrix $[\bigwedge^2 A_i-I]$, rows
$(1,2,3,5,6)$ and columns $(1,2,4,6,7)$ have determinant $1$.
Its primitive cokernel is $\mathbb Z$, detected by
\[
 \xi=u\wedge w+6\psi\wedge\delta.
\]
The normal-circle argument of \ref{maps:src:D.2.6} applies to these actual
primitive Mumford models and accounts for all eight columns.
Consequently
\[
 H_2(X_D)=Q=\mathbb Z\langle s,t\rangle,
 \quad s=[e_1\wedge e_2],\quad t=[T_{\ell,0}].
\]
Here the degree-one Mayer--Vietoris map has free domain of rank
$11$, free target of rank $12$, and cokernel $\mathbb Z$;
its kernel is zero. The stated generators therefore also
describe $H_2(X_D)$, rather than only a quotient of it.
The flat form $\xi$ has values $(1,0)$ on these generators.

\item In the finite linear coinvariants,
$e_1=-2e_2$ and $e_3=-3e_2$. Hence its reduced fibre has
\[
 H_1(F)=
 \mathbb Z\langle e_0,e_2,g\rangle/
 \langle mg-ne_0+3r e_2\rangle.
\]
This is free, with basis $e_2,z$, where
$z=\alpha e_0+\beta g$ and $m\alpha+n\beta=1$.
The exact norm and contraction identities are
\[
 \mathcal N_m e_2=-k e_3,\qquad
 \mathcal N_m e_0=m\ell,\qquad
 \xi(e_2,v)=2n,\qquad \xi(e_0,v)=6r.
\]
The first norm dies in $Q$. Thus the normal tori of
$e_2,e_0,g$ have coordinates $(-2n,0)$, $(-6r,m)$,
$(0,n)$. In the free basis $e_2,z$ the remaining matrix is
\[
 B_{k,n,r}=\begin{pmatrix}-2n&-6r\alpha\\0&1\end{pmatrix}.
\]
Its Smith form is $(1,2|n|)$.

\item All four boundary $H_1$ groups are free of total rank $14$;
the target has the same rank and cokernel $\mathbb Z/|n|$.
The degree-one kernel is zero. Each central fibre is irreducible
and a multiple is principal, so the boundary-to-filling map
in degree two is onto, by \ref{maps:src:D.2.5}. Therefore
$H_2(X)=\operatorname{coker}B_{k,n,r}$.
Euler localization gives $\chi(X)=2$; rational duality gives
$b_3=0$. Integral duality and the universal coefficient theorem
then give exactly the groups in the statement.

\item The program \texttt{code/manifolds\_wang.py} executes both
primitive-minor certificates, the norm identities and 336 full
$10\times10$ Smith forms with $1\leq k\leq12$,
$-15\leq n\leq15$, $\gcd(n,3k)=1$, and $r=0,3$.
Its input matrices and outputs are retained. The symbolic
reduction proves the formula for arbitrary displayed parameters.
\end{enumerate}
\end{proof}

The order-three geometric model and its original surgery calculation are given in \cite[Theorem~7.1 and Appendix~B]{manWang26}. The present comparison retains the primitive $I_2$ period $(2,1)$, which becomes $(1,0)$ under $(x,y)=(v,v^2/z)$; it is not replaced by the zero-weight $I_2$ archival API choice.

\begin{theorem}[the oriented type of all simply connected Wang members]\label{maps:src:D.2.8}
\StatementRecord{maps:src:D.2.8}{proved}{Historical events: Claude Fable 5.2 (Anthropic), effort high, 16 November 2026: carried as background without fault found. GPT-6 Astra (OpenAI), ChatGPT Pro mode, effort Pro, S6\_09\_PAPERS, 17 November 2026: targeted proof check. Also checked in this writing. Pre-continuation writing events: GPT-6 Astra (OpenAI), ChatGPT Pro mode, effort Pro, S6\_10\_PAPERS, 18 November 2026: targeted proof check. Continuation metadata events: GPT-6 Astra (OpenAI), ChatGPT Pro mode, effort Pro, S6\_10\_PAPERS continuation, 19 November 2026: used and rechecked in this run; status-vocabulary reconciliation; metadata reconciliation only; no new mathematical check. \emph{Pre-continuation recorded proof scope.} GPT-6 Astra (OpenAI, ChatGPT Pro mode), effort Pro, 2026-11-17 UTC; Claude Fable 5.2 (Anthropic), effort high, 16 November 2026; also checked in this writing by GPT-6 Astra (OpenAI), ChatGPT Pro mode, effort Pro, 18 November 2026}{the hypotheses, notation and definitions stated here; \cref{maps:src:D.2.7}; \cite[Theorem 7.1 and Appendix B]{manWang26}; \cite[Theorems 3, 2.10, pp.795,803, and \S3.12, p.807]{manZhubr75}}


Every member of \ref{maps:src:D.2.7} with $n=\pm1$, arbitrary $k>0$ and
arbitrary $r\in\mathbb Z$, is orientation-preservingly
diffeomorphic to
\[
 M_{\mathrm{conic}}=
 \bigl(\mathbb {CP}^3\setminus\operatorname{int}
       (S^2\times D^4)\bigr)
 \mathop{\cup}_{S^2\times S^3}(D^3\times S^3),
\]
where the removed sphere is the conic
$[s:t]\mapsto[s^2:st:t^2:0]$ with its oriented normal framing.
For the nonzero $a\in H^2(X;\mathbb F_2)$,
\[
 \langle a^3,[X]_2\rangle=1.
\]
This also supplies an independent check of the smooth
classification asserted for order three in
\cite[Theorem 7.1 and Appendix B]{manWang26}.
\end{theorem}


\begin{proof}
\begin{enumerate}[label=\arabic*.,leftmargin=*]
\item We first state precisely the classification being used, including
its two-primary information. Zhubr's Theorem 3 identifies the
oriented simply connected spin six-manifolds with fixed $H_2=G$
and $b_3=0$ with the corresponding spin-bordism classes over
$K(G,2)$, modulo automorphisms of $G$.
For $G=\mathbb Z/2$ his Theorem 2.10 and the calculation in
\S3.12 give
\[
 \Omega^{\mathrm{Spin}}_6(K(\mathbb Z/2,2))=\mathbb Z/2;
\]
the generator is the image of
$(\mathbb {CP}^3,\mathbb {CP}^3\hookrightarrow
K(\mathbb Z,2))$ under reduction $\mathbb Z\to\mathbb Z/2$
\cite[Theorems 3, 2.10, pp.795,803, and \S3.12, p.807]{manZhubr75}.
The homomorphism evaluating the cube of the universal
degree-two mod-two class is nonzero on this generator, and
therefore is an isomorphism to $\mathbb F_2$.
The automorphism group of $\mathbb Z/2$ is trivial. Hence there
are precisely two such oriented smooth types and their
distinguishing invariant is $\langle a^3,[X]_2\rangle$.
In particular, the vanishing of the ordinary integral $p_1$
does not need to be mistaken for a calculation of every refined
two-primary invariant.

\item The conic has normal bundle $\mathcal O(4)\oplus\mathcal O(2)$.
As an oriented real rank-four bundle on $S^2$ it is trivial:
its class is detected by $w_2$, which is zero. Its framing is
unique up to homotopy since $\pi_2SO(4)=0$.
Surgery preserves simple connectivity and spin, and its integral
homology sequence, with the attaching sphere representing twice
the generator, gives $H_2=H_3=\mathbb Z/2$ and all other reduced
groups zero except $H_6=\mathbb Z$.
The mod-two hyperplane class extends across the surgery trace,
because it evaluates to zero on the conic. Bordism invariance
therefore gives cube one on $M_{\mathrm{conic}}$.
It remains to compute this cube on the entire family.

\item Keep the common piece $X_D$ from \ref{maps:src:D.2.7}. The scalar involution
$\iota(z,v)=(z,-v)$ has exactly two fixed curves there: the
double elliptic curve $E$ of the primitive $I_2$ model, and the
rational curve $C$ corresponding to a vertical edge in the
rank-two triangulation. At $I_2$ the monomial coordinates
$(v,v^2/z)$ have involution $(-x,y)$, so both normal directions
at $E$ have sign $-1$. At the other rank-one point, the induced
translation of the elliptic factor is a nonzero two-torsion
point, with no fixed points. In a rank-two triangle the weights
are two odd normal weights and one even tangential weight.
There is one vertical edge orbit and its two endpoints are the
two distinct triple points, giving the smooth curve $C$.

\item The involution extends over the finite logarithmic quotient,
and introduces no additional fixed points there even when $k$
is even. Indeed, an equality $\iota x=g^j x$ on its central
covering torus would imply $jn/(3k)=0$ in $\mathbb R/\mathbb Z$
after applying $\psi$. Thus $3k\mid j$, since $n=\pm1$.
The equality would then assert that translation by $e_3/2$
fixes a point on a torus, which is impossible. Away from the
central torus the base rotation gives the same conclusion.
This replaces the odd-order shortcut in Wang's order-three
proof and establishes the fixed-set statement for every $k,r$.

\item We give the integral line-bundle calculation on the unchanged
piece $X_D$. In multiplicative coordinates let $h_1,h_2$ be
the elliptic period and the representative of $P$, and let $R$
be the multiplier of the contraction $F$. On an auxiliary line
with coordinate $\zeta$, use
\[
 \begin{aligned}
 h(z,v,\zeta)&=(h_1z,h_2v,-z^{-1}\zeta),\\
 F(z,v,\zeta)&=(h_2^{-6}z,Rv,h_2^{-3}R^3v^6\zeta).
 \end{aligned}
\]
The two lifts commute and their alternating Chern form is
$\xi=uw+6\psi\delta$. Their semicharacter modulo two is
$u+w+uw$, with zero value on the radical spanned by $e_0,e_3$;
direct substitution shows that all four monodromies preserve it.
The symmetry is checked, not inferred from vertical curvature.
Writing $B=h_2^{-3}R^3$, the lift
\[
 J(z,v,\zeta)=(z^{-1},v^{-1},-z^{-1}\zeta)
\]
satisfies $J^2=1$, $JhJ=h^{-1}$ and $JFJ=F^{-1}$;
the last identity is exactly $B^2=h_2^{-6}R^6$.
Multiplying $J$ by $-1$ normalizes it on the rigidifying section.
Thus these are symmetric rigidified Appell--Humbert data.
The displayed semicharacter is unchanged on overlap markings,
so the normalized fibrewise identifications patch.

\item Compare with the prequantum model on the flat real universal
cover, whose connection potential is $\pi i\xi(x,dx)$ and
whose lattice multipliers have that same semicharacter.
It is preserved by the linear monodromies. Its fibrewise
holomorphic structure is the Appell--Humbert structure for $\xi$.
Both families are symmetric and rigidified and have the same
alternating form. Their ratio is therefore a family of flat
line bundles of order dividing two. Its square is canonically
trivial on each fibre after rigidification: the trivializing
holomorphic section is unique with value one at the origin.
Those unique sections patch over the base. Hence the comparison
line is globally two-torsion on the regular family, rather than
merely having trivial fibrewise Chern class. Since $H^2(V;\mathbb Z)$
is free by the preceding integral graph calculation, its Chern
class is zero. The actual holomorphic line consequently has
exactly the prequantum Chern class on $V$, whose values on
fibre and swept cycles are respectively one and zero.
Finally $H_2(V)\to H_2(X_D)$ is onto by \ref{maps:src:D.2.7} and
$H_1(X_D)=\mathbb Z$, so restriction
$H^2(X_D;\mathbb Z)\to H^2(V;\mathbb Z)$ is injective.
The extension $\mathcal L$ below therefore has values $(1,0)$
on $(s,t)$, including the mixed sweep value.

\item These factors extend over the three Mumford pieces. At $I_2$
the three valuations are
\[
 (\operatorname{ord}h_1,\operatorname{ord}h_2,
   \operatorname{ord}R)=(2,1,-3).
\]
Replacing $F$ by $Fh^3$ gives line multipliers, up to units,
$-x^{-2}y$ and $y^3$. The support function in the degenerating
coordinate is $-j(j-1)$; all its slopes are even. Hence the
stabilizer character of the resulting line $\mathcal L$ is
trivial along $E$. At the rank-two point the valuations are
$(1,0,1)$ and a support function is
\[
 \varphi(a,b)=-\frac{a(a-1)}2+3b^2.
\]
Its period differences are $-a$ and $6b+3$, as required by
the line multipliers $z^{-1}$ and $q^3v^6$.
Its vertical slopes are odd, giving the nontrivial character
along $C$; the adjacent horizontal slopes differ by one, so
$\deg(\mathcal L|_C)$ is odd. At the remaining rank-one
point the same factors extend over its primitive fan.
Units contribute no winding or parity, and translation to the
section vertex changes the vertical slopes by a multiple of six.

\item We now check the global mod-two extension without assuming that
the degree-$3k$ transfer is invertible. \ref{maps:src:D.2.7} proved
$H_2(X_D)=\mathbb Z(s,t)$ and identified
$c_1(\mathcal L)$ with $(1,0)$ on that basis.
For $n=\pm1$ both first entries of
\[
 B_{k,n,r}=\begin{pmatrix}-2n&-6r\alpha\\0&1\end{pmatrix}
\]
are even. Its cokernel is generated by $s$ of order two, and
its second relation sends $t$ to an even multiple of $s$.
The nonzero global class $a$ therefore restricts exactly to
$\rho_2c_1(\mathcal L)$ on $X_D$. There is no universal
coefficient ambiguity on $X_D$ since $H_1(X_D)=\mathbb Z$.
This is the integral attachment step that makes the argument
valid for all $k$ and every $r$.

\item Finally $H^2(X;\mathbb Z)=0$ implies $c_1(X)=0$, hence spin.
Write $G=\langle\iota\rangle$ and
$\tau\in H^1(BG;\mathbb F_2)$ for the generator.
The class $a$ admits an equivariant lift $\widetilde a$:
the degree-two Borel obstruction with target
$H^2(BG;H^1(X;\mathbb F_2))$ vanishes, and a fixed point
splits the map to $BG$, killing the remaining differential
into $H^3(BG;\mathbb F_2)$.
The circle containing $\iota$ acts trivially on
$\pi_1(X_D)=\mathbb Z$, and a fixed point splits the extension,
so $\pi_1((X_D)_G)=\mathbb Z\times\mathbb Z/2$.
The difference between $\widetilde a$ and the equivariant
class of $\mathcal L$ vanishes in ordinary cohomology and thus
comes from this group's degree-two cohomology, by the
universal-cover low-degree sequence. Every fixed point induces
the same order-two element $(0,1)$ in this group. A single
adjustment by a multiple of $\tau^2$ consequently gives
\[
 \widetilde a|_C=a|_C+\tau^2,
 \qquad \widetilde a|_E=a|_E+\tau\eta,
 \quad\eta\in H^1(E;\mathbb F_2).
\]
The normal bundle of each fixed curve has $w_2=0$, since
$c_1(X)=0$ and $c_1(TC),c_1(TE)$ are even.
Both normal complex directions have sign $-1$, so their
equivariant Euler class is $\tau^4$.
The equivariant Thom isomorphism after inverting $\tau$ gives
\[
 \langle a^3,[X]_2\rangle
 =\int_C\frac{(a|_C+\tau^2)^3}{\tau^4}
  +\int_E\frac{(a|_E+\tau\eta)^3}{\tau^4}
 =\int_C a|_C=1.
\]
For the second equality, squares of degree-one classes on the
oriented elliptic curve vanish, and all other unwanted terms
have the wrong ordinary degree on a curve. The stated
classification now identifies $X$ with $M_{\mathrm{conic}}$.
\end{enumerate}
\end{proof}

\begin{theorem}[all orders $4k$ in the Liu--Rollenske configuration]\label{maps:src:D.2.9}
\StatementRecord{maps:src:D.2.9}{proved}{Historical events: Claude Fable 5.2 (Anthropic), effort high, 16 November 2026: carried as background without fault found. GPT-6 Astra (OpenAI), ChatGPT Pro mode, effort Pro, S6\_09\_PAPERS, 17 November 2026: targeted proof check. Also checked in this writing. Pre-continuation writing events: GPT-6 Astra (OpenAI), ChatGPT Pro mode, effort Pro, S6\_10\_PAPERS, 18 November 2026: targeted proof check. Continuation metadata events: GPT-6 Astra (OpenAI), ChatGPT Pro mode, effort Pro, S6\_10\_PAPERS continuation, 19 November 2026: used and rechecked in this run; status-vocabulary reconciliation; metadata reconciliation only; no new mathematical check. \emph{Pre-continuation recorded proof scope.} GPT-6 Astra (OpenAI, ChatGPT Pro mode), effort Pro, 2026-11-17 UTC; Claude Fable 5.2 (Anthropic), effort high, 16 November 2026; also checked in this writing by GPT-6 Astra (OpenAI), ChatGPT Pro mode, effort Pro, 18 November 2026}{the hypotheses, notation and definitions stated here; \cref{maps:src:D.2.6}; \cite[Lemma 3.10, Proposition 3.11]{manLR26}; \cite[Proposition 3.13]{manLR26}; \cref{maps:src:D.2.5}}


Use the $III^*,II,I_1$ surface and the unchanged type-$II$
filling of Liu--Rollenske, with the degree-one linearization zero
at $I_1$ and the primitive hexagonal filling there.
In their homology basis $(e_1,e_2,e_3,e_4)$ put
\[
 \delta=-e_1+2e_3,\qquad \ell=e_4,\qquad
 m=4k,\qquad v=n\ell+r\delta,
 \quad k>0,\quad\gcd(n,4k)=1.
\]
The logarithmic transformation at $III^*$ exists for every
such integer datum, and its total space satisfies
\[
 \pi_1(X)=\mathbb Z/|n|,\qquad
 H_0=H_6=\mathbb Z,\quad H_5=0,\quad
 H_1=H_2=H_3=H_4=\mathbb Z/|n|.
\]
Thus $n=\pm1$ gives $S^6$ for every $k>0$. In particular,
the Liu--Rollenske configuration also answers Viaclovsky's
Question 4 affirmatively. The type-$II$ fibre is unchanged
in this statement.
\end{theorem}


\begin{proof}
\begin{enumerate}[label=\arabic*.,leftmargin=*]
\item The source's good-reduction construction and section-preserving
comparison at $III^*$ give the analytic model. The degree-$k$
pullback and free affine quotient argument of \ref{maps:src:D.2.6} applies
with invariant covector $e_4^*$, which takes value $n$ on $v$.
It proves existence, smoothness, properness and the relation
$g^m=v$ without restricting $k$.

\item The extra local input, absent from \ref{maps:src:D.2.6}, is the unchanged
type-$II$ filling. By \cite[Lemma 3.10, Proposition 3.11]{manLR26},
its marked neighbourhood is homeomorphic to the type-$II$
elliptic neighbourhood times $T^2$, with the section preserved.
The cuspidal rational curve is homeomorphic to $S^2$.
The section evaluates as one on both this central two-cycle
and the nearby elliptic curve, so the degree-two specialization
is primitive. The product description makes its two normal
tori precisely the meridian swept by each of the two transverse
$T^2$ loops: each bounds the section disc times that loop.
K\"unneth, or the relative sequence of this product pair,
shows that these generate the entire rank-two kernel of
the boundary map in degree two. This justifies the integral
attachment used below without treating the cusp as a normal
crossing curve.

\item Here are all matrices and primitive certificates. The source
uses $T_\infty T_1T_0=I$ in homology. We order its matrices
as $(A_0,A_1,A_2)=(T_0,T_\infty,T_1)$, so their product is $I$:
\[
 \begin{gathered}
 A_0=\begin{pmatrix}-1&-1&-1&0\\2&1&1&0\\0&0&1&0\\0&0&0&1\end{pmatrix},\qquad
 A_1=\begin{pmatrix}1&0&0&-1\\-2&0&-1&1\\2&1&2&1\\0&0&0&1\end{pmatrix},\\
 A_2=\begin{pmatrix}1&1&0&1\\-2&0&-1&-1\\0&-1&1&-1\\0&0&0&1\end{pmatrix}.
 \end{gathered}
\]
These are \cite[Proposition 3.13]{manLR26} in its actual
integral marking. The cusp's primitive vanishing lattice is
$\langle a,\delta\rangle$, where $a=e_1-e_2-e_3$.
Its monodromy-normal closure contains
$A_2a=-e_2$, and hence all of
$\langle e_1,e_2,e_3\rangle$. The two unchanged meridians
are killed by the extending section. The sphere relation kills
the third, and $g^m=v$ then leaves only $ne_4=0$.
Surjectivity of the boundary maps on fundamental groups gives
$\pi_1(X)=\mathbb Z/|n|$.

\item Set $D=[A_0-I\ \ A_1-I]$ and, in $\Lambda^2$, define
\[
 j_0(x)=(x,0),\quad j_1(x)=(0,x),\quad
 j_2(x)=-(A_1x,x).
\]
The cusp fixed basis is $(a,\delta)$ and the type-$II$ fixed
basis is $(\delta,e_4-e_2)$. Form
\[
 K=[j_1(a)\ j_1(\delta)\ j_2(\delta)\
       j_2(e_4-e_2)\ j_0(\ell)].
\]
Then $DK=0$, $\operatorname{rank}D=3$, and the $5\times5$
minor on rows $(1,2,4,5,6)$ is $1$.
Moreover $j_0(\delta)=-K_2-K_3$.
The exterior matrix $[\bigwedge^2A_i-I]_{i=0}^2$ has rank
five and the minor on rows $(1,2,3,4,5)$ and columns
$(2,4,5,6,11)$ is $1$. Its primitive cokernel is detected by
\[
 \omega=2e_1^*e_2^*-e_2^*e_3^*-e_3^*e_4^*.
\]
Thus $H_2(V)$ is free of rank six, split by a fibre class
$s=-e_3e_4$, of $\omega$-value one, and the five zero-value
sweeps $K$. The four normal tori of the two unchanged fillings
are exactly its first four sweep columns, by the preceding
product argument and the rank-two argument in \ref{maps:src:D.2.6}.
After quotienting by them the group is
$Q=\mathbb Z\langle s,t\rangle$, $t=j_0(\ell)$.

\item The finite linear coinvariants satisfy $e_1=e_2=0$, so
$\delta=2e_3$. The finite reduced fibre therefore has free
$H_1$ with basis $e_3,z$, where
\[
 mg=ne_4+2r e_3,\qquad z=\alpha e_4+\beta g,
 \qquad m\alpha+n\beta=1.
\]
Direct calculation gives
\[
 \mathcal N_me_3=2k\delta,\quad
 \mathcal N_me_4=m\ell,\quad
 \omega(e_3,v)=-n,\quad \omega(e_4,v)=2r.
\]
The $\delta$ sweep dies in $Q$, and the normal-circle formula
therefore reduces the last map to
\[
 \begin{pmatrix}n&-2r\alpha\\0&1\end{pmatrix}.
\]
All boundary $H_1$ groups are free of total rank nine, and their
target is free of rank nine with finite cokernel
$\mathbb Z/|n|$; the degree-one kernel is zero.
Every central fibre is irreducible and a multiple is principal,
so the degree-two boundary maps are onto as in \ref{maps:src:D.2.5}.
It follows that $H_2(X)=\mathbb Z/|n|$.
The unchanged type-$II$ piece and the finite quotient have
Euler number zero, whereas the hexagonal fibre has Euler
number two. Duality yields all the other displayed groups.
Sphere recognition when $|n|=1$ is as in \ref{maps:src:D.2.6}.

\item The independent script
\texttt{code/manifolds\_liu\_rollenske.py} checks the two
unit minors, the norm identities, and 312 full $6\times6$
Smith forms for the same $(k,n,r)$ bounds as \ref{maps:src:D.2.6}.
The proof establishes the formula without those bounds.
\end{enumerate}
\end{proof}

The product type-$II$ cap and the original marked torus monodromy are established in \cite[Lemma~3.10 and Propositions~3.11, 3.13]{manLR26}. The proof above gives the full indicated $3k$ family in that marking.

\begin{theorem}[unrestricted orders in both Yang families]\label{maps:src:D.2.10}
\StatementRecord{maps:src:D.2.10}{proved}{Historical events: Claude Fable 5.2 (Anthropic), effort high, 16 November 2026: carried as background without fault found. GPT-6 Astra (OpenAI), ChatGPT Pro mode, effort Pro, S6\_09\_PAPERS, 17 November 2026: targeted proof check. Also checked in this writing. Pre-continuation writing events: GPT-6 Astra (OpenAI), ChatGPT Pro mode, effort Pro, S6\_10\_PAPERS, 18 November 2026: targeted proof check. Continuation metadata events: GPT-6 Astra (OpenAI), ChatGPT Pro mode, effort Pro, S6\_10\_PAPERS continuation, 19 November 2026: used and rechecked in this run; status-vocabulary reconciliation; metadata reconciliation only; no new mathematical check. \emph{Pre-continuation recorded proof scope.} GPT-6 Astra (OpenAI, ChatGPT Pro mode), effort Pro, 2026-11-17 UTC; Claude Fable 5.2 (Anthropic), effort high, 16 November 2026; also checked in this writing by GPT-6 Astra (OpenAI), ChatGPT Pro mode, effort Pro, 18 November 2026}{the hypotheses, notation and definitions stated here; \cite[Lemma 4.1, Lemma 4.3]{manYang26}; \cite[Proposition 4.4]{manYang26}; \cite[Theorem 7.8]{manYang26}; \cite[Lemma 7.10, Appendix B.4]{manYang26}; \cite[Theorem~7.13]{manYang26}}


Fix a positive cusp index $R$ in Yang's Type B or Type C period
family, with its canonical cusp collar. Use the following constants:
\[
 \begin{array}{c|ccc|cc}
 \text{type}&c&d&a=d/c&m_1&m_2\\\hline
 B&2&4&2&4k&6l\\
 C&3&3&1&3k&6l
 \end{array}
 \qquad k,l\geq1.
\]
In the marking $\Lambda=\mathbb Z\langle\gamma,u,w,\delta\rangle$,
put $\ell_1=\gamma-R(u+w)$ and choose the invariant finite
translation vectors
\[
 v_1=n_1\ell_1+r_1\delta,\qquad
 v_2=n_2\gamma+r_2\delta,\qquad
 \gcd(n_j,m_j)=1,
 \qquad n_j,r_j\in\mathbb Z.
\]
These data give compact smooth complex threefolds for every displayed
integer choice. Define
\[
 P=m_2n_1+m_1n_2,\qquad h=\gcd(R,ak,P)>0.
\]
If $P\ne0$, set
$G=\mathbb Z/h\oplus\mathbb Z/(R|P|/h)$, omitting trivial factors.
Then
\[
 \begin{array}{c|ccccccc}
 i&0&1&2&3&4&5&6\\\hline
 H_i(X;\mathbb Z)&\mathbb Z&\mathbb Z/|P|&
 \mathbb Z^{R-1}\oplus G&G&
 \mathbb Z^{R-1}\oplus\mathbb Z/|P|&0&\mathbb Z
 \end{array}
\]
and $\pi_1(X)=\mathbb Z/|P|$.
If $P=0$, necessarily $m_1=m_2$ and $n_2=-n_1$; in this case
$h=\gcd(R,ak)$, $\pi_1(X)=\mathbb Z$, and
\[
 \begin{array}{c|ccccccc}
 i&0&1&2&3&4&5&6\\\hline
 H_i(X;\mathbb Z)&\mathbb Z&\mathbb Z&
 \mathbb Z^{R+1}\oplus\mathbb Z/h&
 \mathbb Z^2\oplus\mathbb Z/h&\mathbb Z^{R+1}&
 \mathbb Z&\mathbb Z.
 \end{array}
\]
In both cases $\chi(X)=2R$; the displayed homology is independent
of $r_1,r_2$. No member is simply connected. More generally, with
these two finite multiplicities, this last obstruction holds for
every free affine choice of the two logarithmic fillings, even if
the sufficient coprimality condition above is dropped.

For $k=l=1$, $(n_1,n_2)=(1,-1)$ and $r_1=r_2=0$, this independently
recovers Yang's original groups: Type B has fundamental group
$\mathbb Z/2$ and middle torsion $\mathbb Z/R\oplus\mathbb Z/2$;
Type C has fundamental group $\mathbb Z/3$ and middle torsion
$\mathbb Z/(3R)$.
\end{theorem}


\begin{proof}
\begin{enumerate}[label=\arabic*.,leftmargin=*]
\item \emph{Analytic construction and exact marking.}
The actual dual deck matrices of the primary source are
\[
 A_1=\begin{pmatrix}
 1&0&0&0\\0&0&1&0\\-cR&-1&2-c&0\\-R&0&-1&1
 \end{pmatrix},\qquad
 A_2=\begin{pmatrix}
 1&0&0&0\\0&1&-1&0\\0&1&0&0\\0&0&0&1
 \end{pmatrix}.
\]
Their orders are $d,6$, and their integral fixed lattices are
$\langle\ell_1,\delta\rangle$ and
$\langle\gamma,\delta\rangle$, respectively.
Pull the good-reduction root discs back by degrees $k,l$.
On the resulting root family the affine action is
\[
 (s,x)\longmapsto(\zeta_{m_j}^{-1}s,A_jx+v_j/m_j).
\]
The holomorphic period formula and the logarithmic collar maps of
\cite[Lemma 4.1, Lemma 4.3]{manYang26} apply unchanged: in complex
coordinates the translation is the holomorphic period section
$\Pi(\zeta_{m_j}^{-1}s)v_j/m_j$.
The invariant functional $\chi=\gamma^*$ changes by $n_j/m_j$.
Applying it to a purported fixed point of a nontrivial power
forces $m_j\mid qn_j$, impossible for $0<q<m_j$.
The actions are therefore free, their quotients are proper, and
the central divisors have multiplicities $m_j$.
The collar isomorphism adds or subtracts
$(\log s)/(2\pi i)\,\Pi(s)v_j$; changing the logarithm adds an
integral period. This proves existence with the stated marking
and leaves the toroidal cusp unchanged. Compactness follows from
the same three compact tubes and compact base core as in the
source's global gluing argument.

\item Write $g_j$ for the transferred clockwise meridian. The finite
boundary map imposes exactly $g_j^{m_j}=v_j$, by
\cite[Proposition 4.4]{manYang26}; this can also be checked from
the path $q\mapsto qv_j/m_j$ in the central affine quotient.
The cusp kills $w,\delta$ and $g_1g_2$.
Since $A_1w=u+(2-c)w-\delta$, their normal closure kills $u$.
The remaining generators $\gamma,g_1$ commute and satisfy
\[
 m_1g_1=n_1\gamma,\qquad -m_2g_1=n_2\gamma.
\]
The entries of this relation matrix have gcd one, and its
determinant has absolute value $|P|$. This proves the asserted
fundamental groups, including the rank-one group when $P=0$.
For arbitrary free affine fillings one can instead kill the
entire fibre subgroup: the resulting quotient is
\[
 \langle g_1,g_2\mid g_1g_2=1,
                  g_1^{m_1}=g_2^{m_2}=1\rangle
 =\mathbb Z/\gcd(m_1,m_2).
\]
It is nontrivial, because its order is divisible by two in Type B
and by three in Type C. This proves the stronger obstruction.

\item \emph{Integral finite-quotient lattices.}
Before the cusp is attached, let $X^\circ$ be the union of the two
finite pieces. Each retracts to its reduced quotient surface $S_j$,
and their intersection retracts to a regular torus $F$.
The first linear coinvariants satisfy $w=-u$, $\delta=cu$;
the second kill $u,w$. Thus
\[
 \begin{split}
 H_1(S_1)&=\langle\gamma,u,g_1\mid
             m_1g_1=n_1\gamma+cr_1u\rangle,\\
 H_1(S_2)&=\langle\gamma,\delta,g_2\mid
             m_2g_2=n_2\gamma+r_2\delta\rangle.
 \end{split}
\]
Each relation is primitive by $\gcd(n_j,m_j)=1$; both groups are
$\mathbb Z^2$. Euler characteristic zero, duality and the universal
coefficient theorem give $H^2(S_j)=\mathbb Z^2$ with a unimodular
intersection pairing.

\item Use the ordered regular-fibre cohomology basis
\[
 (\gamma^*u^*,\gamma^*w^*,\gamma^*\delta^*,
            u^*w^*,u^*\delta^*,w^*\delta^*).
\]
Solving the integral fixed equations gives the first invariant
lattice with basis
\[
 \omega_1=-R\gamma^*u^*+R\gamma^*w^*+u^*w^*,\qquad
 \omega_2=\gamma^*u^*-\gamma^*w^*+c\gamma^*\delta^*.
\]
Both squares vanish and $\omega_1\omega_2=c\,\mathrm{vol}_F$.
The second invariant lattice has basis
$u^*w^*,\gamma^*\delta^*$ and pairing $U$.
If $L_j=p_j^*H^2(S_j)$, transfer gives injectivity and the
discriminant formula gives the exact indices
\[
 [H^2(F)^{A_1}:L_1]=m_1/c=ak,
 \qquad [H^2(F)^{A_2}:L_2]=m_2.
\]
Indeed, the quotient pairing is unimodular and pullback multiplies
it by the covering degree $m_j$.

\item The square $[q\delta+t v_j/m_j]$ closes to a two-torus in $S_j$.
For a pulled-back integral class $\alpha$ its evaluation is
$\alpha(\delta,v_j)/m_j\in\mathbb Z$.
For $\alpha=A\omega_1+B\omega_2$ this is
$-cn_1B/m_1$; coprimality forces $B\in ak\mathbb Z$.
At the second fibre it forces the coefficient of
$\gamma^*\delta^*$ to lie in $m_2\mathbb Z$.
These candidate sublattices have exactly the required indices,
so they are the entire pullback lattices, without a hidden
saturation factor. Consequently their integral basis matrices are
\[
 P_1=\begin{pmatrix}
 -R&ak\\R&-ak\\0&m_1\\1&0\\0&0\\0&0
 \end{pmatrix},\qquad
 P_2=\begin{pmatrix}
 0&0\\0&0\\0&-m_2\\1&0\\0&0\\0&0
 \end{pmatrix}.
\]
Choose the homology bases dual to these integral cohomology
bases. In particular, when $|n_j|>1$ the finite $\delta$ sweep is
not declared to be a basis vector. Its actual coordinates are
$(0,-n_1)$ and $(0,n_2)$ in the respective surfaces. This is the
integral distinction needed for arbitrary twists.
The degree-two Mayer--Vietoris map is therefore
\[
 \alpha_2=\begin{pmatrix}
 -R&R&0&1&0&0\\ak&-ak&m_1&0&0&0\\
 0&0&0&-1&0&0\\0&0&m_2&0&0&0
 \end{pmatrix}.
\]
The degree-one kernel is exactly the primitive line
$\ker\alpha_1=\mathbb Z(u+w)$: the second inclusion kills only
the $u,w$ directions of $\Lambda$, and the first then forces their
coefficients to agree. The extra relation involving $g_j$ cannot
create another kernel vector from $\Lambda$, because its
$g_j$ coefficient is the nonzero integer $m_j$.

\item \emph{The two cusp relations are primitive.}
Let $N$ be the unchanged cusp tube, $M=\partial N$, and
$P_0=A_1A_2$ its clockwise forward transport.
Direct exterior-square calculation gives rank two and a unit
$2\times2$ minor of $\bigwedge^2P_0-I$: use columns $(1,4)$
and rows $(2,6)$ in Type B or $(3,6)$ in Type C.
In both cases the determinant is $-1$, for every $R$.
Thus its cokernel is $\mathbb Z^4$. The Wang sequence then makes
$H_2(M)$ free, and the two swept tori $\xi_w,\xi_\delta$ project
to the integral basis $w,\delta$ of $\ker(P_0-I)$.
They generate a primitive direct summand.

\item For completeness, the normalization calculation used here is
integral. There are $R$ normalized $dP_6$ components, $3R$ double
curves and $2R$ triple points. With $\epsilon=1$ in Type B and
$\epsilon=2$ in Type C, component labels are
$j=a-\epsilon b\pmod R$. If a component class has coordinates
$(A_j,b_{1j},b_{2j},b_{3j})$, its three restriction differences are
\[
 \begin{split}
 X_j&=b_{1j}-A_{j+1}+b_{2,j+1}+b_{3,j+1},\\
 Y_j&=A_j-b_{1j}-b_{2j}-b_{3,j-\epsilon},\\
 Z_j&=b_{2j}-A_{j-\epsilon-1}
                  +b_{1,j-\epsilon-1}+b_{3,j-\epsilon-1}.
 \end{split}
\]
Its image is exactly the kernel of the surjective map
$(X,Y,Z)\mapsto(\sum(X+Y),\sum(Y+Z))$ to $\mathbb Z^2$.
To verify equality, put $T=S^{-1}$ for cyclic shift and use $Y$
to eliminate $A$. Solve
$Z+T^{\epsilon+1}Y=(I-T)q$, possible precisely when the second
sum vanishes; set $b_2=q$, $b_3=-T^{-\epsilon}q$.
The identity
$(I-T^{\epsilon+1})-T(I-T^\epsilon)=I-T$ solves that equation.
The other one is an integer cyclic difference equation for $b_1$
with zero coordinate sum, and is therefore soluble.
This proves saturation for all $R$, including self-identifications
when $R=1$. The normalization spectral sequence, whose other
rows are the dual two-torus and $\mathbb Z^R$ in degree four,
now gives $H^3(N)=\mathbb Z^2$, $H^4(N)=\mathbb Z^R$ and
$\chi(N)=2R$; this checks \cite[Theorem 7.8]{manYang26} directly.

\item The angular torus action on $N$ sweeps the extended zero-section
disc into integral solid-torus chains bounding $\xi_w,\xi_\delta$.
Poincar\'e--Lefschetz duality gives
$H_3(N,M)=H^3(N)=\mathbb Z^2$, so the kernel of
$H_2(M)\to H_2(N)$ has rank at most two. It contains the primitive
rank-two summand just constructed; a finite enlargement would
embed in the free quotient of $H_2(M)$ by that summand.
Hence
\[
 K_2=\ker(H_2(M)\to H_2(N))
       =\mathbb Z\xi_w\oplus\mathbb Z\xi_\delta.
\]
This establishes the complete kernel over $\mathbb Z$.

\item The two finite forward transports are $P_j=A_j^{-1}$.
They send $w$ first to $-u$ and then back to $w$.
Cutting the clockwise cusp sweep into the two cylinders gives
connecting class $-(u+w)$ in the convention
$\alpha=(i_{1*},-i_{2*})$, exactly as the explicit cylinder
boundaries in \cite[Lemma 7.10, Appendix B.4]{manYang26}.
It is a primitive generator of $\ker\alpha_1$, irrespective
of the finite translation vectors. Sweeping $\delta$ over the
pair of pants gives
$j_*\xi_\delta=-s_{1,\delta}-s_{2,\delta}$.
The evaluations already computed show that $-j_*\xi_\delta$
has the column $b=(0,-n_1,0,n_2)^{\mathsf T}$ in the integral
dual bases. Therefore
\[
 H_2(X^\circ)/j_*K_2
    =\operatorname{coker}[\alpha_2\mid b].
\]
Eliminate the unit relation from the $uw$ column. The remaining
three-generator presentation is
\[
 C=\begin{pmatrix}
 R&0&0\\-ak&m_1&-n_1\\0&m_2&n_2
 \end{pmatrix},\qquad \det C=RP.
\]
Its entry gcd is one. Its nonzero $2\times2$ minors, up to signs,
are
\[
 Rm_1,\ Rn_1,\ Rm_2,\ Rn_2,\ akm_2,\ akn_2,\ P.
\]
Their gcd is exactly $h=\gcd(R,ak,P)$, by the two coprimality
conditions. Thus for $P\ne0$ its Smith factors are
$(1,h,R|P|/h)$; for $P=0$ its rank is two and its factors are
$(1,h,0)$. The latter cokernel is $\mathbb Z\oplus\mathbb Z/h$.

\item \emph{The relative term, including $P=0$.}
Excision and duality give $H_2(N,M)=H^4(N)=\mathbb Z^R$.
Before the cusp is attached, the joint monodromy images kill
$\ker\chi$ in the abelianization. Hence
\[
 H_1(X^\circ)=\left\langle\gamma,g_1,g_2\ \middle|\
 m_1g_1=n_1\gamma,\quad m_2g_2=n_2\gamma\right\rangle
       =\mathbb Z\oplus\mathbb Z/d_0,
 \quad d_0=\gcd(m_1,m_2).
\]
The equality follows integrally: the relation matrix has entry
gcd one and maximal-minor gcd $d_0$.
The full image of the relative boundary map is
$\langle a_0\rangle$, where $a_0=-(g_1+g_2)$.
Indeed, the cusp boundary kernel is generated by its meridian
and its vanishing lattice; the latter is already zero in
$H_1(X^\circ)$, and the extended section disc supplies $a_0$.

\item When $P\ne0$, the primitive homomorphism from $H_1(X^\circ)$
to its free quotient is
\[
 \Phi(\gamma,g_1,g_2)
 =\frac1{d_0}(m_1m_2,\ m_2n_1,\ m_1n_2).
\]
The same gcd-of-minors calculation proves primitivity.
It sends $a_0$ to $-P/d_0\ne0$, so $\langle a_0\rangle$
is infinite cyclic. The kernel of the relative boundary map
is free of rank $R-1$. This does not assume that $a_0$ is
primitive in $H_1(X^\circ)$.

\item When $P=0$, equality of the reduced fractions
$n_1/m_1=-n_2/m_2$ forces $m_1=m_2=m$ and $n_2=-n_1$.
Use $a_0$ and $g_1$ as meridian generators: the relations are
$ma_0=0$ and $mg_1=n_1\gamma$. The second is primitive,
so $a_0$ has exact order $m$. The relative boundary image is
therefore the finite group $\mathbb Z/m$, and its kernel is
free of rank $R$, of index $m$ in $\mathbb Z^R$.
This is the extra free contribution in the degenerate case.

\item The exact sequence for $(X,X^\circ)$ now reads
\[
 0\longrightarrow\operatorname{coker}C
 \longrightarrow H_2(X)
 \longrightarrow\ker\bigl(H_2(N,M)\to H_1(X^\circ)\bigr)
 \longrightarrow0.
\]
Its right group is free, so it splits. This gives the displayed
degree-two groups in both cases. All pieces other than the cusp
have Euler characteristic zero, so $\chi(X)=2R$.
For $P\ne0$, duality gives $b_1=b_5=0$, $b_2=b_4=R-1$,
and Euler characteristic gives $b_3=0$.
For $P=0$, it gives $b_1=b_5=1$, $b_2=b_4=R+1$ and then
$b_3=2$. In either case the universal coefficient theorem and
duality identify the degree-three torsion with the degree-two
torsion and the degree-four torsion with the degree-one torsion.
This proves all the integral homology assertions separately
in the two cases.

\item The script \texttt{code/manifolds\_yang.py} verifies the symbolic
monodromy and primitive-minor identities, 32 full cusp
normalization Smith forms, 162 finite-quotient and sweep
checks, and 2592 full augmented $4\times7$ Smith forms,
including 72 with $P=0$. Its executable also accepts arbitrary
admissible integer inputs. These finite checks accompany the
unbounded integral proof; they impose no hypothesis on it.
The original cases are proved in \cite[Theorem~7.13]{manYang26}.
\end{enumerate}
\end{proof}

\begin{theorem}[all orders $3k$ for $IV^*+II+I_2$]\label{maps:src:D.2.11}
\StatementRecord{maps:src:D.2.11}{proved}{Historical events: Claude Fable 5.2 (Anthropic), effort high, 16 November 2026: carried as background without fault found. Claude Fable 5.2 (Anthropic), effort high, 17 November 2026: carried as background without fault found. Claude Opus 5.5 (Anthropic), effort max, sweep 1, scope s6, 18 November 2026: used and rechecked in this run. GPT-6 Astra (OpenAI), ChatGPT Pro mode, effort Pro, S6\_09\_PAPERS, 17 November 2026: targeted proof check. Also checked in this writing. Pre-continuation writing events: GPT-6 Astra (OpenAI), ChatGPT Pro mode, effort Pro, S6\_10\_PAPERS, 18 November 2026: targeted proof check. Continuation metadata events: GPT-6 Astra (OpenAI), ChatGPT Pro mode, effort Pro, S6\_10\_PAPERS continuation, 19 November 2026: used and rechecked in this run; status-vocabulary reconciliation; metadata reconciliation only; no new mathematical check. \emph{Pre-continuation recorded proof scope.} GPT-6 Astra (OpenAI, ChatGPT Pro mode), effort Pro, 2026-11-17 UTC; Claude Fable 5.2, effort high, 2026-11-17; GPT-6 Astra (OpenAI, ChatGPT Pro mode), S6\_06\_BLIND, effort Pro, 2026-11-17; Claude Fable 5.2 (Anthropic), effort high, 16 November 2026; also checked in this writing by GPT-6 Astra (OpenAI), ChatGPT Pro mode, effort Pro, 18 November 2026}{the hypotheses, notation and definitions stated here; \cref{maps:src:D.2.1}; \cite[Proposition 3.1]{manYang26}; \cite[Proposition 3.4]{manYang26}; \cite[Proposition 3.6]{manYang26}; \cite[Lemma 3.10]{manLR26}; \cite[Proposition 3.11]{manLR26}; \cref{maps:src:D.2.10}; \cref{maps:src:D.2.6}}


The rational elliptic surface
\[
 y^2=x^3-3t^3(t-1)x+t^4(t-1)(2t-1),
 \qquad P=(t^2,t^2),
 \tag{D.C1}\label{maps:eq:D.C1}
\]
has fibres $IV^*,II,I_2$ at $0,1,\infty$ and Mordell--Weil group
$\mathbb Z$ generated by $P$. Put $M=\mathcal O(P-O)$ and
$Q=3P$. Let $C_0$ be the identity component of the $I_2$ fibre.
There is a linearization with divisor $C_0$, and its torus
family has Yang's index-one Type C marking. Keep the type-$II$
filling unchanged and use the index-one toroidal filling at
$I_2$. At $IV^*$ choose any logarithmic order and invariant vector
\[
 m=3k,\qquad v=n(\gamma-u-w)+r\delta,
 \qquad k>0,\quad n,r\in\mathbb Z,
 \quad\gcd(n,3k)=1.
\]
The resulting compact smooth complex threefold has
\[
 \pi_1(X)=\mathbb Z/|n|,\qquad
 H_i(X;\mathbb Z)=
 \begin{cases}
 \mathbb Z&i=0,6,\\
 \mathbb Z/|n|&1\leq i\leq4,\\
 0&i=5.
 \end{cases}
\]
Consequently $n=\pm1$ gives the standard $S^6$ for every positive
$k$. This is an affirmative unbounded configuration for
Viaclovsky's Question 4. The linearization is supported on one
$I_2$ component, so the local-narrowness restriction of the public
engine does not include this input.
\end{theorem}


\begin{proof}
\begin{enumerate}[label=\arabic*.,leftmargin=*]
\item \emph{The explicit elliptic surface and linearization.}
For the coefficients $f,g$ of \eqref{maps:eq:D.C1}, direct calculation gives
\[
 4f^3+27g^2=27t^8(t-1)^2,\qquad
 j=6912t(1-t).
\]
The minimal valuations give $IV^*$ at zero and $II$ at one.
In the chart at infinity the discriminant has order two and
the cubic has a node, giving $I_2$. These fibres have Euler
sum twelve; the degree-four and degree-six Weierstrass
coefficients give the rational elliptic surface. Substitution
checks the section $P$, which is disjoint from $O$ and meets
the nonidentity components at $IV^*$ and $I_2$.
Its height is $2-4/3-1/2=1/6$. The Shioda--Tate rank is one,
and OS entry 49 has height lattice $\langle1/6\rangle$ and
no torsion, as checked in \ref{maps:src:D.2.1}; hence $P$ is a generator.

\item The section $Q=3P$ meets the identity component of $IV^*$
and the nonidentity component of $I_2$. Its height is $3/2$,
so the height formula gives $Q\cdot O=0$.
The mixed height $\langle P,Q\rangle=1/2$, with the sole
local contribution $1/2$ at $I_2$, gives $P\cdot Q=0$.
Let $J=M\otimes t_Q^*M^{-1}$. Its multidegree vanishes at
$IV^*$ and $II$. At $I_2$, translation by $Q$ exchanges the
two components, so its multidegrees are $(-2,2)$ on
$(C_0,C_1)$, exactly those of $\mathcal O(C_0)$.
Their difference is thus a pullback from the base. Restriction
to $O$ gives
\[
 J\cdot O=Q\cdot O-P\cdot Q+1=1=C_0\cdot O,
\]
so the base degree of that difference is zero. Therefore
$J\simeq\mathcal O(C_0)$, supplying the asserted effective
linearization. Its norm under the component exchange has divisor
$C_0+C_1=F_\infty$. This explains why a half full-fibre weight
appears in period coordinates while the geometric divisor is
integral.

\item \emph{Identification with the actual Type C periods.}
The elliptic $j$-function is the one in
\cite[Proposition 3.1]{manYang26}. Over $t=s^3$ the good
model has central curve $Y^2=X^3+1$, and $P$ specializes to
$(0,1)$, a nonzero fixed point of order three.
Over $t-1=s^6$, the same section specializes to the origin
of the good elliptic model. At $I_2$ it meets the nonidentity
component. Choose the sign of the elliptic marking so that its
Abel--Jacobi coordinate at the first point is $(\tau+1)/3$.
Its normalized lift therefore obeys the two Type C affine
laws and has $\mu-\tau/2$ regular at infinity.
The uniqueness proof of \cite[Proposition 3.4]{manYang26},
using $H^0(\mathcal O(-1))=H^1(\mathcal O(-1))=0$, identifies
it with $\mu_C$.
The extension line $M$ has this same Picard coordinate,
whereas the translation has coordinate $3\mu_C$.
The norm divisor just calculated fixes the cusp central
period. The remaining scalar period is the affine constant
of \cite[Proposition 3.6]{manYang26}. Consequently the open
two-torus family is the index-one Type C family, for a suitable
admissible constant, and its primitive toroidal cusp has one
normalized $dP_6$ component.

\item An exact integral check of the marking is also available.
In the order $(II,IV^*,I_2)$, the homology matrices are
\[
 \begin{split}
 M_{II}&=\begin{pmatrix}1&1&0&0\\-1&0&0&0\\0&0&1&0\\0&0&0&1\end{pmatrix},\\
 M_{IV}&=\begin{pmatrix}0&-1&0&0\\1&-1&0&-3\\0&1&1&1\\0&0&0&1\end{pmatrix},\\
 M_{I_2}&=\begin{pmatrix}1&2&0&3\\0&1&0&0\\0&-1&1&-1\\0&0&0&1\end{pmatrix}.
 \end{split}
\]
The first two follow from the locally narrow finite-point
section cocycles $(0,0)$ and $(0,-1)$ for $P$, multiplied by
three for $Q$. Their product determines the third matrix;
no local-narrowness hypothesis at $I_2$ is used.
Put $U=M_{IV}$, $V=U^{-1}M_{II}U$ and
\[
 S=\begin{pmatrix}
 0&0&-1&0\\0&-1&1&1\\0&0&0&1\\-1&0&3&0
 \end{pmatrix}.
\]
Then $\det S=-1$, $U^{-t}S=ST_1$ and $V^{-t}S=ST_2$,
where $T_1,T_2$ are precisely Yang's displayed Type C
cohomological matrices. This is the finite Hurwitz reorder
followed by an integral change of basis, not merely a rational
comparison of eigenvalues.

\item \emph{The unchanged type-$II$ filling with its integral marking.}
Here the product proof of \cite[Lemma 3.10]{manLR26} applies
with $3P$ in place of its translation section; we check all the
needed ingredients. A closed elliptic neighborhood $E$ of the
type-$II$ fibre retracts to its cuspidal cubic, homeomorphic to
$S^2$. The restriction of $M$ has degree zero and is therefore
differentiably trivial. In a smooth trivialization the lifted
contracting generator has form
\[
 (x,z)\longmapsto(t_{-Q}(x),\lambda(x)z),\qquad |\lambda|<1.
\]
As $E$ is simply connected, $\lambda$ has a smooth logarithm.
The relative exponential of the additive N\'eron group is an
isomorphism on the central smooth locus. After shrinking the
disc it provides a holomorphic logarithm of $-Q$ and hence a
fibre-preserving isotopy of $t_{-Q}$ to the identity.
The translations extend over the cuspidal cubic and fix its
cusp. Combining this isotopy with the logarithm of $\lambda$
trivializes the circle-bundle mapping torus of a fundamental
annulus. Thus this filling is homeomorphic over its base disc
to $E\times T^2$, with angular circle $\delta$ and a primitive
second circle of $\gamma^*$-value one.

\item Its regular specialization sends the elliptic area $uw$ to
the fundamental class of the cuspidal cubic and $\gamma\delta$
to the area of this second torus; it kills the mixed classes.
Equivalently the exact degree-two pullback lattice is
$\langle u^*w^*,\gamma^*\delta^*\rangle$.
There is no index-six factor: this is the unchanged filling,
not a free order-six affine quotient.
Both transverse normal tori bound products of a section disc
with a circle. In particular its $\delta$ meridian sweep is
zero. These assertions also follow by applying the marked
specialization proof of \cite[Proposition 3.11]{manLR26}
to the product just constructed.

\item \emph{All orders and all displayed twists.}
At $IV^*$ the order-$3k$ free quotient and its marked collar
are those of \ref{maps:src:D.2.10} with $R=1$, $n_1=n$, $r_1=r$.
The two other compact fillings have just been specified,
so the same finite-core gluing proves compactness and
smoothness. The unchanged type-$II$ piece kills $u,w$ and
its meridian; the cusp kills $\delta$ and the remaining
meridian. Van Kampen leaves exactly $n\gamma=0$, proving
$\pi_1(X)=\mathbb Z/|n|$.

\item The first finite pullback matrix is unchanged. For the second
piece the primitive product lattice replaces its coefficient
$m_2$ by one, and its vanishing $\delta$ sweep replaces $n_2$
by zero. Thus the complete augmented integral map is
\[
 D=\begin{pmatrix}
 -1&1&0&1&0&0&0\\
 k&-k&3k&0&0&0&-n\\
 0&0&0&-1&0&0&0\\
 0&0&1&0&0&0&0
 \end{pmatrix}.
\]
The primitive $w$ cusp sweep still kills the degree-one
Mayer--Vietoris kernel. Unit elimination reduces $D$ to
one unit factor and
\[
 \begin{pmatrix}1&0&0\\-k&3k&-n\\0&1&0\end{pmatrix}.
\]
Its Smith factors are $(1,1,|n|)$ for arbitrary allowed
integers, independently of $r$.
Before attaching the cusp, the remaining degree-one group is
$\langle\gamma,g\mid3kg=n\gamma\rangle=\mathbb Z$.
Its cusp meridian is $-g$, a nonzero multiple $-n$ of a
primitive generator. Hence the relative map
$H_2(N,M)=\mathbb Z\to H_1(X^\circ)=\mathbb Z$ is injective,
so it contributes no extra free degree-two class.
We conclude $H_2(X)=\mathbb Z/|n|$.
The cusp has Euler number two and the other pieces have
Euler number zero. Duality and the universal coefficient
theorem give the remaining groups. When $|n|=1$, sphere
recognition applies exactly as in \ref{maps:src:D.2.6}.

\item The independent executable
\texttt{code/manifolds\_type\_c\_ii.py} checks the Weierstrass
identity, section, $j$-function, integral monodromy conjugacy,
symbolic determinant, and 168 full augmented Smith forms.
The preceding proof has no upper bound on $k,n,r$.
\end{enumerate}
\end{proof}

The local Type-C geometric input is supplied by \cite[Propositions~3.1, 3.4, 3.6, 4.4 and Appendix~B.4]{manYang26}; the product type-$II$ marking is supplied by \cite[Lemma~3.10 and Propositions~3.11, 3.13]{manLR26}. The unrestricted $3P$ calculations use the explicitly displayed boundary maps.

\begin{theorem}[seven further unbounded one-log families]\label{maps:src:D.2.12}
\StatementRecord{maps:src:D.2.12}{proved}{Historical events: Claude Fable 5.2 (Anthropic), effort high, 16 November 2026: carried as background without fault found. GPT-6 Astra (OpenAI), ChatGPT Pro mode, effort Pro, S6\_09\_PAPERS, 17 November 2026: targeted proof check. Also checked in this writing. Pre-continuation writing events: GPT-6 Astra (OpenAI), ChatGPT Pro mode, effort Pro, S6\_10\_PAPERS, 18 November 2026: targeted proof check. Continuation metadata events: GPT-6 Astra (OpenAI), ChatGPT Pro mode, effort Pro, S6\_10\_PAPERS continuation, 19 November 2026: used and rechecked in this run; status-vocabulary reconciliation; metadata reconciliation only; no new mathematical check. \emph{Pre-continuation recorded proof scope.} GPT-6 Astra (OpenAI, ChatGPT Pro mode), effort Pro, 2026-11-17 UTC; Claude Fable 5.2 (Anthropic), effort high, 16 November 2026; also checked in this writing by GPT-6 Astra (OpenAI), ChatGPT Pro mode, effort Pro, 18 November 2026}{the hypotheses, notation and definitions stated here; \cref{maps:src:D.2.1}; \cref{maps:src:D.2.6}; \cite{manExplorer26}; \cref{maps:eq:D.R1}; \cref{maps:eq:D.R3}; \cref{maps:eq:D.R4}}


Let $G$ be the marked Mordell--Weil generator of the indicated
rational elliptic surface. In the table, the superscript $\bullet$
marks the fibre supporting the degree-one linearization, and $d$
is the linear order at the sole additive fibre.
\[
 \begin{array}{c|l|rr|r}
 \mathrm{OS}&\text{ordered fibre configuration}&P/G&Q/G&d\\\hline
 45&I_8,II,I_1^\bullet,I_1&8&1&6\\
 47&I_7,I_2,I_1^\bullet,II&14&1&6\\
 47&I_7,I_1^\bullet,I_1,III&7&2&4\\
 55&I_5,I_4,II,I_1^\bullet&20&1&6\\
 56&I_5,I_2,IV,I_1^\bullet&10&3&3\\
 56&I_5,I_3,III,I_1^\bullet&15&2&4\\
 49&I_2,IV^*,I_1^\bullet,I_1&2&3&3
 \end{array}
\]
Use $M=\mathcal O(P-O)$ and translation by $Q$.
Keep the two weight-zero multiplicative fillings and the
primitive toroidal filling at $I_1^\bullet$.
In the integral marking specified below, put
\[
 m=dk,\qquad v=n\ell+r\delta,
 \qquad k>0,\quad\gcd(n,dk)=1,
 \qquad n,r\in\mathbb Z.
\]
Every row admits these logarithmic fillings for every displayed
integer choice, and the resulting compact smooth threefold has
\[
 \begin{gathered}
 \pi_1(X)=\mathbb Z/|n|,\qquad
 H_i(X;\mathbb Z)=\mathbb Z/|n|\quad(1\leq i\leq4),\\
 H_5(X)=0,\qquad H_0(X)=H_6(X)=\mathbb Z.
 \end{gathered}
\]
In particular $n=\pm1$ gives $S^6$ for every $k>0$ in all seven
rows. Within the one-log construction, any simply connected
choice of affine vector must have $n=\pm1$, so the coprimality
hypothesis discards none of its sphere cases.
\end{theorem}


\begin{proof}
\begin{enumerate}[label=\arabic*.,leftmargin=*]
\item \emph{Geometric data and the actual multiplicative fillings.}
The configurations and height lattices are the primary
Oguiso--Shioda/Fukae data verified in \ref{maps:src:D.2.1}. Each displayed
pair has height pairing one. At every fibre at least one of
$P,Q$ is locally narrow, and at each weight-zero $I_j$ the
section $P$ is narrow while $Q$ generates its component group.
In the marked component groups the latter residues are
\[
 \begin{array}{c|l}
 \text{row}&Q\text{ on its higher multiplicative fibres}\\\hline
 1&3\pmod8\\
 2&5\pmod7,\quad1\pmod2\\
 3&3\pmod7\\
 4&2\pmod5,\quad3\pmod4\\
 5&3\pmod5,\quad1\pmod2\\
 6&2\pmod5,\quad1\pmod3\\
 7&1\pmod2.
 \end{array}
\]
All residues are units. At the additive fibre $Q$ is narrow;
its component residue is zero. The exact $P$ residues as well
as these $Q$ residues are retained with the matrices in the
accompanying \texttt{manifolds\_rankone\_families.json} record.

\item The line $J=M\otimes t_Q^*M^{-1}$ has multidegree zero on
every fibre: either $t_Q$ preserves its components, or $M$
already has multidegree zero there. Translation acts trivially on the degree-zero Picard group of
the generic elliptic fibre, so $J$ is generically trivial. Its
vertical divisor and zero component degrees then make it a
pullback from the base; all original fibres have multiplicity one.
Its degree is $\langle P,Q\rangle=1$, as all local mixed
height corrections vanish. Choose its section with its simple
zero at the marked $I_1^\bullet$. The contracting-quotient
construction and its properness estimate are the same as in
Engel's construction.

\item At a weight-zero $I_j$, the restriction of $M$ to each normalized
component $\mathbb P^1$ has degree zero and is trivial.
Translation by $Q$ cyclically permutes all $j$ components,
because its component residue is a unit. The quotient therefore
has one normalization component and one double curve.
After $j$ steps the first component returns; its normalization
is the quotient of $\mathbb P^1\times\mathbb C^*$ by a
translation on the second factor and a constant multiplicative
automorphism on the first. It is a $\mathbb P^1$ bundle over an
elliptic curve defined by a degree-zero line, hence is
topologically trivial. Its zero and infinity sections are
identified by a translation and flat units. These are isotopic
to the identity in the marked circle directions.
Each node has the original semistable chart $xy=t$ up to a
holomorphic unit: the infinite contracting action has no point
stabilizer and does not ramify this chart.
Thus the filled neighborhood and its normal-circle chains are
the primitive one-node model $N_{\rm node}\times T^2$ used
in \ref{maps:src:D.2.6}. In particular all three fixed one-cycles supply its
three normal tori, not just the vanishing circle.
This argument identifies the geometric local model, rather
than inferring it from a rational monodromy matrix.

\item At the marked $I_1$, the simple zero gives the usual one-component
hexagonal toroidal filling. Its two angular normal tori are
primitive. At the additive fibre, take the good-reduction root
model and then its degree-$k$ base change. An invariant vector
with primitive invariant-coordinate value $n$ defines a free
order-$dk$ affine action because $\gcd(n,dk)=1$.
The marked logarithmic collar and compact gluing are exactly
those proved in \ref{maps:src:D.2.6}. These steps prove existence and properness
for every allowed $k,n,r$.

\item \emph{Exact integral certificates, with no parameter bounds.}
Use the geometric open-family monodromy lifts from the section
cocycles and degree-one linearization, as in \ref{maps:src:D.2.1} and the
public primary monodromy derivation of \cite{manExplorer26}.
Cyclically reorder the four matrices so that the additive
one is first, and denote them by $A_0,A_1,A_2,A_3$, with
$A_0A_1A_2A_3=I$.
The marking here is $(e_1,e_2,\delta,\gamma)$, with
$\chi=\gamma^*$. The following integral vectors and constants
are obtained by exact unimodular reduction:
\[
 \begin{array}{c|r|r|l}
 \text{row}&d&q&\ell\\\hline
 1&6&1&(-1,0,0,1)\\
 2&6&1&(0,-1,0,1)\\
 3&4&2&(-3,-2,0,1)\\
 4&6&1&(0,0,0,1)\\
 5&3&3&(-1,3,0,1)\\
 6&4&2&(-1,3,0,1)\\
 7&3&3&(-3,-1,0,1).
 \end{array}
\]
In every row $A_0\ell=\ell$, $A_i\delta=\delta$ and
$\chi(\ell)=1$.
All four matrices, integral fixed bases and the integer
left inverses used next are written in the retained JSON;
the executable derives and verifies them using only the
MIT monodromy source, without a local-model database or API.

\item Let $D=[A_0-I\ A_1-I\ A_2-I]$ and set
\[
 j_0(x)=(x,0,0),\quad j_1(x)=(0,x,0),\quad
 j_2(x)=(0,0,x),\quad
 j_3(x)=-(A_1A_2x,A_2x,x).
\]
For each $i=1,2,3$, choose the full integral fixed basis
$F_i=\ker(A_i-I)$, of rank two at the cusp and rank three
at a one-node fibre. Form the $12\times9$ matrix
\[
 K=[j_1(F_1)\ j_2(F_2)\ j_3(F_3)\ j_0(\ell)].
\]
The exact certificates give
\[
 DK=0,\quad\operatorname{rank}D=3,\quad LK=I_9,\quad
 (Lj_0\delta)_9=0,
 \tag{D.R1}\label{maps:eq:D.R1}
\]
where $L$ is the explicitly retained integer left inverse.
Thus the nine actual sweep tori form a primitive integral
basis of the graph kernel. The exterior coinvariant matrix
$[\bigwedge^2A_i-I]_{i=0}^3$ has five nonzero Smith entries,
all one. Its primitive dual row is the invariant form
\[
 \omega=P_0\,e_1^*e_2^*-Q_0\,\delta^*\gamma^*,
 \qquad P=P_0G,\quad Q=Q_0G.
 \tag{D.R2}\label{maps:eq:D.R2}
\]
The regular-family $H_2$ is therefore free of rank ten:
one fibre coinvariant of $\omega$-value one and the nine
sweeps. The constant form $\omega$ has zero value on the
latter actual fixed-loop sweeps, so this is an integral
splitting. Quotienting by the eight normal tori of the
unchanged fillings gives
$\mathbb Z\langle s,t\rangle$, with $t=j_0(\ell)$.
Equation \eqref{maps:eq:D.R1} also proves $j_0(\delta)=0$ in this quotient.

\item In the vertical lattice $\ker\chi$, the finite coinvariants
are free of rank one. Let $b$ be a primitive representative;
the exact fixed-basis computation gives
\[
 [\delta]=q[b],\qquad
 \mathcal N_d b=(d/q)\delta,\qquad
 \mathcal N_d\ell=d\ell,\qquad
 -\omega(b,\ell)=1,\quad
 -\omega(\ell,\delta)=-q.
 \tag{D.R3}\label{maps:eq:D.R3}
\]
Here $\mathcal N_d=\sum_{i=0}^{d-1}A_0^i$; all are
identities of integral vectors, not rational averages.
The finite reduced fibre has free $H_1$, with generators
$b,z$, where
\[
 mg=n\ell+rq b,\qquad
 z=\alpha\ell+\beta g,\qquad m\alpha+n\beta=1.
\]
Apply the normal-circle torus formula of \ref{maps:src:D.2.6} to these two
cycles. The first column is the fibre area
$-\omega(b,v)$ and the sweep $j_0(\mathcal N_m b)$;
the second is $-\alpha\omega(\ell,v)$ and
$j_0(\alpha\mathcal N_m\ell+\beta v)$.
After the eight unchanged normal relations, \eqref{maps:eq:D.R1}--\eqref{maps:eq:D.R3}
give exactly
\[
 \begin{pmatrix}n&-qr\alpha\\0&1\end{pmatrix}.
 \tag{D.R4}\label{maps:eq:D.R4}
\]
Its Smith form is $(1,|n|)$ for all the asserted integers.

\item For van Kampen, the images of the three unchanged vanishing
lattices already generate the full primitive $\ker\chi$;
their combined matrix has Smith factors $(1,1,1)$, verified
in the same exact computation. Their meridians and the base
relation kill the remaining finite meridian. The only remaining
relation is $n\gamma=0$. This proves the nonabelian
fundamental-group assertion. It also shows, without assuming
coprimality, that any simply connected affine choice must have
$n=\pm1$.

\item The degree-one Mayer--Vietoris map has free source and target
of rank fourteen and finite cokernel $\mathbb Z/|n|$; hence
its kernel is zero. Every local degree-two boundary map is
onto, by the product models, the finite free quotient, and
the toroidal specialization proved above. Therefore \eqref{maps:eq:D.R4}
is the full remaining degree-two presentation and
$H_2(X)=\mathbb Z/|n|$.
Only the hexagonal cusp contributes Euler characteristic,
giving $\chi(X)=2$. Duality supplies all the other groups,
and homotopy-sphere recognition with $\Theta_6=0$ proves
the sphere statement.

\item The exact program \texttt{code/manifolds\_rankone\_families.py}
checks all constant matrices, integral bases and \eqref{maps:eq:D.R1}--\eqref{maps:eq:D.R3},
and also 234 full $10\times10$ Smith forms at displayed
sample parameters. Formula \eqref{maps:eq:D.R4} proves the result for
unrestricted orders and twists; the samples are corroboration.
\end{enumerate}
\end{proof}

\begin{theorem}[all smooth-fibre logs in four configurations]\label{maps:src:D.2.13}
\StatementRecord{maps:src:D.2.13}{proved}{Historical events: Claude Fable 5.2 (Anthropic), effort high, 16 November 2026: carried as background without fault found. GPT-6 Astra (OpenAI), ChatGPT Pro mode, effort Pro, S6\_09\_PAPERS, 17 November 2026: targeted proof check. Also checked in this writing. Pre-continuation writing events: GPT-6 Astra (OpenAI), ChatGPT Pro mode, effort Pro, S6\_10\_PAPERS, 18 November 2026: targeted proof check. Continuation metadata events: GPT-6 Astra (OpenAI), ChatGPT Pro mode, effort Pro, S6\_10\_PAPERS continuation, 19 November 2026: used and rechecked in this run; status-vocabulary reconciliation; metadata reconciliation only; no new mathematical check. \emph{Pre-continuation recorded proof scope.} GPT-6 Astra (OpenAI, ChatGPT Pro mode), effort Pro, 2026-11-17 UTC; Claude Fable 5.2 (Anthropic), effort high, 16 November 2026; also checked in this writing by GPT-6 Astra (OpenAI), ChatGPT Pro mode, effort Pro, 18 November 2026}{the hypotheses, notation and definitions stated here; \cref{maps:src:D.2.14}; \cref{maps:src:D.2.12}; \cite[Proposition 2.3]{manViaclovskyLogs26}; \cref{maps:eq:D.S6}; \cref{maps:eq:D.S5}; \cref{maps:eq:D.S8}}


Consider the following four marked rational elliptic surfaces,
with $G$ a Mordell--Weil generator:
\[
\begin{array}{c|l|r}
\mathrm{OS}&\text{ordered fibres}&p\\\hline
45&I_8,I_1,I_1,I_1,I_1^\bullet&8\\
47&I_7,I_2,I_1,I_1,I_1^\bullet&14\\
55&I_5,I_4,I_1,I_1,I_1^\bullet&20\\
56&I_5,I_3,I_2,I_1,I_1^\bullet&30.
\end{array}
\]
Use $P=pG$, $Q=G$, and the degree-one linearization with zero
at $I_1^\bullet$. Keep the original four primitive one-node
fillings and the one-component hexagonal cusp. At an additional
smooth fibre, perform the logarithmic transformation of order
$m>0$ and integral translation vector
\[
 v=a e_1+b e_2+c\delta+n\gamma,\qquad
 \gcd(m,a,b,c,n)=1.
 \tag{D.S1}\label{maps:eq:D.S1}
\]
This is the exact freeness condition; no coprimality condition
on $m,n$ alone is imposed. Put
\[
 t=\gcd(m,n,a,b)\,\gcd(m,n,p,c).
 \tag{D.S2}\label{maps:eq:D.S2}
\]
The two factors in \eqref{maps:eq:D.S2} are coprime, and $t\mid m$.
For $n\ne0$, also $t\mid n$, and the result has
\[
\begin{aligned}
 \pi_1(X)&=H_1(X)=\mathbb Z/|n|,\qquad H_5(X)=0,\\
 H_2(X)&=H_3(X)=G_{m,n,t},\qquad H_4(X)=\mathbb Z/|n|,\\
 G_{m,n,t}&=\mathbb Z/t\oplus\mathbb Z/m
               \oplus\mathbb Z/(m|n|/t).
\end{aligned}
\tag{D.S3}\label{maps:eq:D.S3}
\]
For $n=0$, the answer is instead
\[
\begin{aligned}
 \pi_1(X)&=H_1(X)=H_5(X)=\mathbb Z,\\
 H_2(X)&=H_3(X)=\mathbb Z^2\oplus\mathbb Z/t\oplus\mathbb Z/m,
 \qquad H_4(X)=\mathbb Z^2.
\end{aligned}
\tag{D.S4}\label{maps:eq:D.S4}
\]
In both cases $H_0=H_6=\mathbb Z$ and $\chi(X)=2$.
The simply connected members are precisely $n=\pm1$; they
have $H_2=H_3=(\mathbb Z/m)^2$. Consequently the sphere
members in these four families are precisely $m=1$, $n=\pm1$.
For $m>1$ the theorem gives simply connected complex rational
homology spheres with two cyclic factors of order $m$ in their
middle torsion. Their oriented smooth type is identified in
\ref{maps:src:D.2.14} below.
\end{theorem}


\begin{proof}
\begin{enumerate}[label=\arabic*.,leftmargin=*]
\item \emph{Geometric construction and complete affine scope.}
In these four configurations the height of $G$ is $1/p$.
The section $P$ is narrow at every fibre, whereas $Q$ has
unit component residue at each higher multiplicative fibre:
respectively $3$ modulo $8$; $(5,1)$ modulo $(7,2)$;
$(2,1)$ modulo $(5,4)$; and $(1,1,1)$ modulo $(5,3,2)$.
Thus $M=\mathcal O(P-O)$ has degree zero on each component,
and the quotient by $Q$ cycles through all components at each
weight-zero fibre. The geometric argument of \ref{maps:src:D.2.12} identifies
each such filling with the primitive one-node model, including
all three normal tori. The mixed height is one, so
$M\otimes t_Q^*M^{-1}$ is the pullback of $\mathcal O(1)$.
Its chosen simple zero gives the stated hexagonal cusp.

\item Over a small disc at the additional regular value, the torus
family has a holomorphic origin and trivial integral local
system. Pull it back by $z=s^m$ and quotient by
\[
 (x,s)\longmapsto (x+v(s)/m,e^{2\pi i/m}s).
\]
Here $v(s)$ is the holomorphic period corresponding to the
integral vector $v$. The action is free exactly when \eqref{maps:eq:D.S1}
holds: away from $s=0$ its nonidentity powers move the base,
and on the central torus translation by $v/m$ has exact
order $m$ precisely under \eqref{maps:eq:D.S1}. The usual logarithmic
collar, given by the translation
$(\log s)v(s)/(2\pi i)$, glues this proper quotient to the
original family; compare \cite[Proposition 2.3]{manViaclovskyLogs26}.
Changing $v$ by a multiple of $m$ retains its specified
integral collar clutch, rather than discarding that clutch.
The reduced central fibre is the torus with lattice
\[
 \Lambda'=\Lambda+\mathbb Z(v/m).
 \tag{D.S5}\label{maps:eq:D.S5}
\]
Its normal complex line is flat of finite order. Since the
integral $H^2$ of a torus is torsion-free, this line is
topologically trivial. Thus the boundary-to-fibre map is
onto on $H_2$, and its kernel consists of its four normal
circle tori.

\item \emph{Primitive global graph calculation.}
The marking is $(e_1,e_2,\delta,\gamma)$ and the invariant
coordinate is $\chi=\gamma^*$. Insert $A_0=I$ for the extra
regular fibre before the five original monodromies
$A_1,\ldots,A_5$. Their product is $I$.
All original vanishing lattices together generate precisely
the primitive $\ker\chi$; their combined matrix has three
unit Smith entries. Therefore van Kampen reduces the group
to the generator $\gamma$. All singular meridians vanish,
and the base relation kills the remaining regular meridian
$\xi$. The new relation $m\xi-v=0$ is then $n\gamma=0$.
This proves the assertion about the full fundamental group.

\item For completeness the integral degree-two calculation is
recorded explicitly. Set
\[
 D=[A_0-I\ A_1-I\ \cdots\ A_4-I]:\mathbb Z^{20}\to\mathbb Z^4.
\]
It has rank three and primitive image. As in \ref{maps:src:D.2.12}, denote
the five direct boundary embeddings into the graph chains by
$j_i$, and set
\[
 j_5(x)=-\bigl(A_1A_2A_3A_4x,\ A_2A_3A_4x,
                      \ A_3A_4x,\ A_4x,\ x\bigr).
\]
The unchanged fillings supply fourteen normal tori: three
at each of the four one-node fibres and two at the cusp.
If $B$ is their $20\times14$ graph matrix, then
\[
 K=[B\ j_0(e_1)\ j_0(e_2)\ j_0(\gamma)]
 \tag{D.S6}\label{maps:eq:D.S6}
\]
is a primitive integral basis of $\ker D$.
The retained integer matrix $L$ satisfies $LK=I_{17}$,
and the last three coordinates of $Lj_0(x)$ are
$(x_1,x_2,x_4)$. In particular the residual $\delta$ sweep
vanishes. These identities are verified for all four
constant marked monodromy tuples in the accompanying exact
program, which saves the matrices and their integer inverses.

\item The exterior coinvariants are free of rank one, detected by
\[
 \omega=p e_1^*e_2^*-\delta^*\gamma^*.
\]
Their presentation has five unit Smith entries. Every column
of \eqref{maps:eq:D.S6} is an actual fixed-loop sweep, on which the constant
vertical form $\omega$ has value zero. This specifies an
integral splitting of regular-family $H_2$.
After eliminating the fourteen unchanged normal tori,
the remaining group is freely generated by a fibre class
of $\omega$-value one and the three sweeps in \eqref{maps:eq:D.S6}.

\item \emph{The exact normal map.}
The lattice \eqref{maps:eq:D.S5} is generated by $e_1,e_2,\delta,\gamma,g=v/m$,
with its single primitive relation
$mg=a e_1+b e_2+c\delta+n\gamma$.
For an integral vector $x$, its normal torus has fibre
coordinate $-\omega(x,v)$ and sweep $j_0(mx)$.
The $g$ normal torus has fibre coordinate zero and sweep
$j_0(v)$. One can read these formulas directly on the
affine boundary torus: its normal circle is $m\xi-v$,
so the first contribution is $-x\wedge v$, and the second
is $mx$ swept around $\xi$. This gives the complete
$4\times5$ matrix
\[
 C=\begin{pmatrix}
 -pb&pa&n&-c&0\\
 m&0&0&0&a\\
 0&m&0&0&b\\
 0&0&0&m&n
 \end{pmatrix}.
 \tag{D.S7}\label{maps:eq:D.S7}
\]
Its columns satisfy exactly the stated lattice relation,
so the five-generator presentation introduces no additional
normal torus.

\item \emph{Smith reduction for every integer choice.}
The greatest common divisor of all entries of $C$ is one,
by \eqref{maps:eq:D.S1}. Its maximal minors, up to signs, are
\[
 m^3n,\quad m^2n^2,\quad m^2nc,\quad m^2nb,\quad m^2na.
\]
For $n\ne0$ their gcd is $m^2|n|$.
The lower determinantal divisors can be checked over each
local ring $\mathbb Z_{(\ell)}$, as follows.
If $m$ is a unit, three diagonal pivots leave just the
relation $n$.
If $n$ is a unit and $m$ is not, its pivot in the top row
and its pivot in the last column leave two copies of $m$.
In the remaining case $\ell\mid m,n$, condition \eqref{maps:eq:D.S1}
forces $c$ or one of $a,b$ to be a unit.
If $c$ is a unit, pivot on $-c$. Row operations give one
copy of $m$ together with the two-row presentation
\[
 \begin{pmatrix}m&0&h\\0&mn&n\end{pmatrix},
 \qquad (h)=(a,b)\text{ in }\mathbb Z_{(\ell)}.
\]
For $n\ne0$ the two nonzero determinantal divisors of this block
are $\gcd(m,n,a,b)$ and $m|n|$, up to local units.
For $n=0$ it has rank one and entry gcd $\gcd(m,a,b)$.
If instead $a$ or $b$ is a unit, pivot on that entry in
the last column. The remaining three-row presentation is
equivalent, up to units, to
\[
 \begin{pmatrix}p&n&-c\\m&0&0\\0&0&m\end{pmatrix}.
\]
Its first two determinantal divisors are
$\gcd(m,n,p,c)$ and $m\gcd(m,n,p,c)$.
The two gcd factors in \eqref{maps:eq:D.S2} have disjoint prime divisors:
a common divisor would divide $m,a,b,c,n$, contrary to
\eqref{maps:eq:D.S1}. These local calculations therefore give the global
Smith entries
\[
 (1,t,m,m|n|/t)\quad(n\ne0),\qquad
 (1,t,m,0)\quad(n=0).
 \tag{D.S8}\label{maps:eq:D.S8}
\]
The same unit-pivot argument proves the $n=0$ formula
directly, without treating a finite group order as a limit.

\item The degree-one Mayer--Vietoris map has free source and
target of rank twenty-four. Its cokernel is the already
computed $H_1(X)$. Thus its kernel is zero when $n\ne0$,
and is a free copy of $\mathbb Z$ when $n=0$.
All local degree-two boundary maps are onto, including
the smooth logarithmic piece by \eqref{maps:eq:D.S5}.
Consequently the full remaining degree-two cokernel is
$\operatorname{coker} C$, with precisely this extra free
relative term in the $n=0$ case. A free quotient splits
the resulting abelian-group extension. This proves $H_2$.
Only the hexagonal cusp contributes Euler characteristic,
giving two. Integral duality and the universal coefficient
theorem now give \eqref{maps:eq:D.S3}--\eqref{maps:eq:D.S4}; in particular $b_3=2$
when $n=0$, and torsion appears in both $H_2$ and $H_3$.
The sphere conclusion follows from $\Theta_6=0$.

\item The program \texttt{code/manifolds\_semistable.py} verifies
all four constant primitive graph and exterior calculations
and 2924 full $18\times19$ Smith presentations, including
388 cases with $n=0$. The proof of \eqref{maps:eq:D.S8} has no bounds
on any parameter. The source comparison is recorded in the verification index.
\end{enumerate}
\end{proof}

\begin{theorem}[the twisted Smale suspension]\label{maps:src:D.2.14}
\StatementRecord{maps:src:D.2.14}{proved}{Historical events: Claude Fable 5.2 (Anthropic), effort high, 16 November 2026: carried as background without fault found. GPT-6 Astra (OpenAI), ChatGPT Pro mode, effort Pro, S6\_09\_PAPERS, 17 November 2026: targeted proof check. Also checked in this writing. Pre-continuation writing events: GPT-6 Astra (OpenAI), ChatGPT Pro mode, effort Pro, S6\_10\_PAPERS, 18 November 2026: targeted proof check. Continuation metadata events: GPT-6 Astra (OpenAI), ChatGPT Pro mode, effort Pro, S6\_10\_PAPERS continuation, 19 November 2026: used and rechecked in this run; status-vocabulary reconciliation; metadata reconciliation only; no new mathematical check. \emph{Pre-continuation recorded proof scope.} GPT-6 Astra (OpenAI, ChatGPT Pro mode), effort Pro, 2026-11-17 UTC; Claude Fable 5.2 (Anthropic), effort high, 16 November 2026; also checked in this writing by GPT-6 Astra (OpenAI), ChatGPT Pro mode, effort Pro, 18 November 2026}{the hypotheses, notation and definitions stated here; \cite[Definition 5.1]{manGR25}; \cref{maps:src:D.2.13}; \cref{maps:src:D.2.12}; \cref{maps:eq:D.S7}; \cref{maps:eq:D.S11}; \cite[Theorem 2.9]{manGR25}; \cite[Theorems 2.2--2.3]{manBarden65}; \cite[Section 4.3]{manReiser24}; \cite[Theorem A]{manGR25}; \cref{maps:eq:D.S13}; \cite[Lemma 5.2]{manGR25}; \cite[Theorem 1]{manZhubr75}; \cref{maps:eq:D.S12}; \cref{maps:eq:D.S14}; \cite[Example 5.3(1)]{manGR25}}


Let $N_m$ be the unique closed simply connected spin five-manifold
with $H_2(N_m;\mathbb Z)=(\mathbb Z/m)^2$, so that $N_1=S^5$.
Define its twisted suspension by
\[
 \widetilde\Sigma_0 N_m
 =\bigl((N_m\setminus\mathring D^5)\times S^1\bigr)
       \cup_{\widetilde\alpha}(S^4\times D^2),
 \qquad \widetilde\alpha(x,z)=(\alpha_zx,z),
 \tag{D.S9}\label{maps:eq:D.S9}
\]
where $\alpha:S^1\to SO(5)$ represents its nontrivial
fundamental-group element. This specifies the framing in
\cite[Definition 5.1]{manGR25}; it is not replaced by an
unspecified choice between the two suspension framings.
Every simply connected member of the four families in
\ref{maps:src:D.2.13} is oriented diffeomorphic to
\[
                         \widetilde\Sigma_0 N_m.
 \tag{D.S10}\label{maps:eq:D.S10}
\]
In particular its oriented smooth type depends only on $m$,
and these types are distinct for distinct positive $m$.
It is spin and has $p_1=0$. The proof retains the full
refined invariant required when $m$ is even; vanishing of
$p_1$ alone is not used as a classification criterion.
\end{theorem}


\begin{proof}
\begin{enumerate}[label=\arabic*.,leftmargin=*]
\item \emph{The angular action and its fixed curve.}
Multiplication of the line-bundle coordinate by
$e^{i\theta}$ defines the global circle action of direction
$\delta$. It commutes with the contracting quotient,
the toroidal gluings and the logarithmic translations.
On the noncompact line-bundle core the action is free:
a nontrivial contracting iterate cannot agree with a unit
scalar, by the strict norm contraction. On the four
weight-zero one-node fillings it is free for the same
reason on the normalized $\mathbb P^1\times\mathbb C^*$
annulus, and it remains free on its double curve.

\item The primitive hexagonal cusp is different. In the height-one
toric lattice, use the cones on the two triangles of a unit
square, with the diagonal joining $(1,0)$ to $(0,1)$.
The cocharacter $\delta=(0,1,0)$ has weights $(-1,0,1)$
in each of their smooth affine charts, up to permutation.
Every nonzero weight is a unit, so the action is semifree
there. Its fixed curves correspond to vertical edges, which
form a single orbit under the two period translations.
They therefore give one embedded $C\simeq\mathbb P^1$
in the quotient. The adjacent wall relation is
\[
 (1,0,1)+(-1,1,1)=(0,0,1)+(0,1,1).
\]
It gives the two normal degrees $-1,-1$. With the circle
weights retained, the normal bundle is
\[
 \nu_C=\mathcal O(-1)_{+1}\oplus\mathcal O(-1)_{-1}.
 \tag{D.S11}\label{maps:eq:D.S11}
\]
Holomorphic unit factors in the actual annular charts
commute with the angular action and deform equivariantly
to the toric factors, so these are the actual smooth
normal data. This also gives a direct chart verification
of the fixed edge described in Engel's final remarks.

\item On the added smooth logarithmic fibre, the lattice is
$\Lambda'=\Lambda+\mathbb Z(v/m)$. If a scalar angle
$s\delta$ lies in this lattice, applying $\chi=\gamma^*$
shows that $kn/m$ is an integer for its $v/m$ coefficient
$k$. Since here $n=\pm1$, it follows that $m\mid k$,
and then $s$ is an integer. Thus there are no new finite
stabilizers at this fibre. The global circle action is
semifree with fixed locus exactly $C$.

\item \emph{Equivariant surgery and the spin quotient.}
The specialization $H_2(T^4)\to H_2$ of the cusp is onto,
by the integral annulus calculation used in \ref{maps:src:D.2.12}.
Consequently $[C]$ maps into the cyclic fibre-coinvariant
summand generated by $s$ in \eqref{maps:eq:D.S7}. Its $\delta$ normal
column imposes $ns=0$; hence $[C]=0$ when $n=\pm1$.
Since $X$ is simply connected, Hurewicz makes $C$
null-homotopic.

\item Conjugating the second complex coordinate in \eqref{maps:eq:D.S11}
identifies its equivariant real bundle with
$\mathcal O(-1)\oplus\mathcal O(1)$ carrying equal
scalar weights. Its transition loop is
$\operatorname{diag}(z^{-1},z)\in SU(2)$ and is
null-homotopic. Thus an equivariant tubular neighborhood
is $S^2\times D^4$, with trivial action on the first
factor and the free Hopf action on each punctured
normal sphere. Conjugating one normal coordinate reverses
normal orientation, so simultaneously reverse the model
$S^2$ orientation in this tubular identification; the
ambient orientation is then preserved.
Replace it equivariantly by $D^3\times S^3$.
The resulting $Y$ has a free circle action.

\item Two homotopic embeddings of $S^2$ in a six-manifold
are isotopic in this dimension; for instance apply
\cite[Theorem 2.9]{manGR25}, since $6>2\cdot2+1$.
Therefore the null-homotopic $C$ is isotopic to a small
unknotted sphere in a ball. Its framing is unique because
$\pi_2(SO(4))=0$. The underlying surgery is consequently
\[
                         Y\cong X\#(S^3\times S^3).
 \tag{D.S12}\label{maps:eq:D.S12}
\]
In particular $Y$ is simply connected.

\item The complex manifold $X$ is spin: $H^2(X;\mathbb Z)=0$
forces $c_1(X)=0$. Its spin structure is unique.
The circle action lifts with period $2\pi$ to this spin
structure. To see the period, lift the universal covering
action of $\mathbb R$ and inspect it at a point of $C$:
the parity of the isotropy weights is $1+(-1)=0$.
Thus the full $2\pi$ rotation acts trivially on the spin
lift, rather than by its nontrivial deck involution.
The equivariant trivialization of \eqref{maps:eq:D.S11} and its standard
Hopf replacement extend this same lift over the surgery.
The free action on $Y$ therefore has an equivariant spin
structure. Quotienting its horizontal spin frames shows
that the oriented orbit manifold $B=Y/S^1$ is spin.

\item The homotopy sequence gives $\pi_1(B)=0$ and the split
abelian-group sequence
\[
 0\longrightarrow H_2(Y)\longrightarrow H_2(B)
       \xrightarrow{\langle e,-\rangle}\mathbb Z
       \longrightarrow0,
\]
where $e$ is the Euler class of $Y\to B$. Thus
$H_2(B)=\mathbb Z\oplus(\mathbb Z/m)^2$ and $e$ is
primitive. The spin case of the oriented Smale--Barden
classification, with its realization of homology
isomorphisms, gives
\[
 (B,e)\cong\bigl((S^2\times S^3)\#N_m,(x,0)\bigr),
 \tag{D.S13}\label{maps:eq:D.S13}
\]
where the sign of the generator $x$ on the first summand
is chosen by the oriented circle-bundle convention.
The precise primary statements are
\cite[Theorems 2.2--2.3]{manBarden65}; the spin uniqueness
for $N_m$ is also recalled in \cite[Section 4.3]{manReiser24}.

\item \emph{The framing and cancellation.}
Apply \cite[Theorem A]{manGR25} to \eqref{maps:eq:D.S13}, with
$B_1=S^2\times S^3$ and $B_2=N_m$.
The Euler class on $B_1$ is primitive, its circle-bundle
total space is $S^3\times S^3$, and $B_1$ is spin.
The spin branch of that theorem specifies the twisted
framing, and yields
\[
          Y\cong(S^3\times S^3)\#\widetilde\Sigma_0 N_m.
 \tag{D.S14}\label{maps:eq:D.S14}
\]
By \cite[Lemma 5.2]{manGR25}, the latter suspension is
simply connected and spin, has the same $H_2$ as $N_m$,
and has $b_3=0$. Zhubr's complete oriented classification
\cite[Theorem 1]{manZhubr75} permits cancellation of the
$S^3\times S^3$ summand in \eqref{maps:eq:D.S12}--\eqref{maps:eq:D.S14}: connected
sum with it leaves $H_2$, the refined class $p$ and the
class $\mu$ unchanged, and adds two to $b_3$.
The diffeomorphism of the stabilized manifolds identifies
these full invariants; deleting the common increase of
$b_3$ identifies the classification data of the two
unstabilized manifolds. This proves \eqref{maps:eq:D.S10}, including
the two-primary data. No general cancellation theorem
outside this classified class is being assumed.

\item Finally $TY\simeq\pi^*TB\oplus\mathbb R$, and
$H^4(B;\mathbb Z)=H_1(B;\mathbb Z)=0$, so $p_1(Y)=0$.
Restriction across \eqref{maps:eq:D.S12} gives $p_1(X)=0$ even though
its target $H^4(X;\mathbb Z)=(\mathbb Z/m)^2$ can be
nonzero. Distinct $m$ have distinct integral torsion
orders. For $m=1$, \eqref{maps:eq:D.S9} is $S^6$, also directly by
\cite[Example 5.3(1)]{manGR25}.
The argument computes the smooth type rather than
deducing it from these additive groups or from $p_1$.
\end{enumerate}
\end{proof}

\begin{theorem}[all-order $S^2\times S^4$ outputs]\label{maps:src:D.2.15}
\StatementRecord{maps:src:D.2.15}{proved}{Historical events: Claude Fable 5.2 (Anthropic), effort high, 16 November 2026: carried as background without fault found. GPT-6 Astra (OpenAI), ChatGPT Pro mode, effort Pro, S6\_09\_PAPERS, 17 November 2026: targeted proof check. Also checked in this writing. Pre-continuation writing events: GPT-6 Astra (OpenAI), ChatGPT Pro mode, effort Pro, S6\_10\_PAPERS, 18 November 2026: targeted proof check. Continuation metadata events: GPT-6 Astra (OpenAI), ChatGPT Pro mode, effort Pro, S6\_10\_PAPERS continuation, 19 November 2026: used and rechecked in this run; status-vocabulary reconciliation; metadata reconciliation only; no new mathematical check. \emph{Pre-continuation recorded proof scope.} GPT-6 Astra (OpenAI, ChatGPT Pro mode), effort Pro, 2026-11-17 UTC; Claude Fable 5.2 (Anthropic), effort high, 16 November 2026; also checked in this writing by GPT-6 Astra (OpenAI), ChatGPT Pro mode, effort Pro, 18 November 2026}{the hypotheses, notation and definitions stated here; \cref{maps:src:D.2.12}; \cref{maps:eq:D.I1}; \cref{maps:src:D.1.2}; \cref{maps:eq:D.I2}; \cref{maps:eq:D.I3}; \cref{maps:eq:D.I4}; \cref{maps:eq:D.S11}; \cref{maps:src:D.2.14}; \cite[Theorem A and Theorem B(2)]{manGR25}}


Take the OS56 type-$II$ profile with ordered elliptic fibres
\[
                         I_2,\ II,\ I_5,\ I_3.
\]
Let $G$ be its marked Mordell--Weil generator, take
$P=30G$, $Q=G$, and put the degree-one linearization zero
at $I_2$. Use the two-component toroidal filling there,
the original primitive one-node fillings at $I_5,I_3$,
and the logarithmic filling at $II$ with
\[
 m=6k,\qquad v=n\gamma+r\delta,\qquad
 k>0,\quad\gcd(n,6k)=1,\quad r\in\mathbb Z.
\]
The resulting compact smooth threefold has
\[
\begin{aligned}
 \pi_1(X)&=H_1(X)=\mathbb Z/|n|,\qquad H_5(X)=0,\\
 H_2(X)&=H_4(X)=\mathbb Z\oplus\mathbb Z/|n|,
 \qquad H_3(X)=\mathbb Z/|n|,\\
 H_0(X)&=H_6(X)=\mathbb Z,\qquad \chi(X)=4.
\end{aligned}
\]
Every simply connected member, equivalently $n=\pm1$,
is oriented diffeomorphic to $S^2\times S^4$.
In particular this fixed configuration supplies that
smooth type for every logarithmic multiplicity $6k$.
\end{theorem}


\begin{proof}
\begin{enumerate}[label=\arabic*.,leftmargin=*]
\item \emph{The actual local geometries.}
The generator has height $1/30$, and $P$ is narrow at
every fibre. The component residues of $Q$ are $1$ in
each of $\mathbb Z/2$, $\mathbb Z/5$, $\mathbb Z/3$,
and zero at $II$. The mixed height is one, so the same
vertical-divisor argument as in \ref{maps:src:D.2.12} gives
$M\otimes t_Q^*M^{-1}=\pi^*\mathcal O(1)$.
Its section can have its simple zero on the whole $I_2$
fibre. The two weight-zero multiplicative fibres give
the primitive one-node models proved in \ref{maps:src:D.2.12}.
The type-$II$ good-reduction model has linear order six,
with invariant basis $\delta,\gamma$; its higher-order
free affine quotients are supplied by the same holomorphic
construction as in \ref{maps:src:D.2.12}.

\item At the zero of the linearization, the Tate parameter has
order two. Because $P$ is narrow, its line has order zero
in the angular Tate direction, while $Q$ has component
residue one. The two valuation periods of the torus
uniformization are therefore
\[
                 (2,0),\qquad(1,1).
 \tag{D.I1}\label{maps:eq:D.I1}
\]
The second coordinate one is the simple zero of the
linearization. Holomorphic units do not change this
integral lattice. Use the standard unit-triangle Mumford
fan and quotient by the lattice $\Gamma$ in \eqref{maps:eq:D.I1}.
It has index two, so the central fibre has two normalized
$dP_6$ components, six double curves and four triple
points. Its total space is smooth: all the fan cones are
unimodular, and the contracting period action is free
and properly discontinuous on the chosen tube.
This is the specified two-component filling, not a
primitive one-component replacement of its monodromy.

\item \emph{The integral cusp defect.}
Let $U$ be a compact tubular neighborhood of this cusp,
and $M=\partial U$. The fan symmetry
\[
 F=\begin{pmatrix}1&1\\0&-1\end{pmatrix},\qquad
 F\begin{pmatrix}2&1\\0&1\end{pmatrix}
     =\begin{pmatrix}2&2\\0&-1\end{pmatrix}
\]
preserves the triangular fan and identifies this lattice
with $2\mathbb Z\oplus\mathbb Z$.
Consequently the integral normalization resolution in
\ref{maps:src:D.1.2}, at $d=2$, applies directly and gives
\[
 \pi_1(U)=\mathbb Z^2,\qquad
 (H_0,H_1,H_2,H_3,H_4)(U)
   =(\mathbb Z,\mathbb Z^2,\mathbb Z^5,\mathbb Z^2,\mathbb Z^2).
 \tag{D.I2}\label{maps:eq:D.I2}
\]
This local normalization calculation does not require the
global covering degree to be coprime to six.

\item The cusp monodromy $A_3$ has $\operatorname{rank}(A_3-I)=2$
and Smith entries $(1,2)$, while
$\bigwedge^2A_3-I$ has two unit Smith entries.
The Wang sequences therefore give
\[
 H_1(M)=\mathbb Z^3\oplus\mathbb Z/2,
 \qquad H_2(M)=\mathbb Z^6.
\]
The cusp kills the full saturated angular lattice
$\ker(A_3-I)$, and its canonical meridian bounds the
extended section disc. Hence the map on $H_1$ is onto
$H_1(U)=\mathbb Z^2$, with kernel
$\mathbb Z\oplus\mathbb Z/2$.
Poincare--Lefschetz duality gives
$H_2(U,M)=H^4(U)=\mathbb Z^2$ and $H_1(U,M)=0$.
The relative sequence then shows that
\[
 H_2(U)/\operatorname{im}H_2(M)\cong\mathbb Z.
 \tag{D.I3}\label{maps:eq:D.I3}
\]
In particular this quotient has no hidden two-torsion.
The two angular normal tori bound in the toric model.
They project to the full primitive fixed lattice in the
Wang sequence for $H_2(M)$, so they form a primitive
rank-two subgroup. By \eqref{maps:eq:D.I2}--\eqref{maps:eq:D.I3} the kernel of
$H_2(M)\to H_2(U)$ has exactly rank two; these are
therefore its whole kernel.
Moreover the remaining image is exactly the regular-fibre
specialization image. Its cokernel is the free group
in \eqref{maps:eq:D.I3}.

\item \emph{The global integral map.}
Cyclically reorder the original four monodromies as
$A_0,A_1,A_2,A_3$ corresponding to $II,I_5,I_3,I_2$.
The full matrices are retained in the accompanying certificate
\texttt{manifolds\_i2\_cusp.json}.
They satisfy $A_0\cdots A_3=I$, $A_0^6=I$ and
fix $\delta$ and the covector $\gamma^*$.
Here $A_0\gamma=\gamma$ and
\[
 \omega=30e_1^*e_2^*-\delta^*\gamma^*
\]
is the primitive exterior coinvariant detector.
The exterior presentation has five unit Smith entries,
so its group is free of rank one.

\item There is one integral issue in degree one that cannot be
discarded. The combined matrix $[A_i-I]_{i=0}^3$ has
Smith entries $(1,1,2)$, so the regular family's fibre
coinvariants are $\mathbb Z\oplus\mathbb Z/2$.
The torsion in the cusp boundary is represented by
$-\delta$, and its image in this global $\mathbb Z/2$
has coefficient $-1$. Thus the torsion summands in
the degree-one Mayer--Vietoris map are mapped
isomorphically. Replacing the cusp image by its actual
saturated angular vanishing lattice gives Smith entries
$(1,1,1)$; it kills the entire primitive
$\ker\gamma^*$ in van Kampen. The base relation kills
the remaining finite meridian, leaving just
$n\gamma=0$. This proves $\pi_1(X)=\mathbb Z/|n|$.

\item Use the boundary maps $j_i$ of \ref{maps:src:D.2.12} and put
\[
 K=[j_1\ker(A_1-I)\ j_2\ker(A_2-I)
       \ j_3\ker(A_3-I)\ j_0(\gamma)].
\]
The exact integer certificate gives a left inverse
$LK=I_9$, and $(Lj_0\delta)_9=0$.
As before, these actual fixed-loop sweeps have zero
$\omega$-area. After the six one-node normal tori and
the two cusp normal tori, the regular-family $H_2$ is
therefore $\mathbb Z\langle s,t\rangle$ with
$t=j_0\gamma$.
The finite vertical coinvariants are generated by
$\delta$, and
$\mathcal N_m\delta=m\delta$,
$\mathcal N_m\gamma=m\gamma$.
Writing $m\alpha+n\beta=1$, its two normal tori give
\[
                        \begin{pmatrix}n&-r\alpha\\0&1\end{pmatrix}.
 \tag{D.I4}\label{maps:eq:D.I4}
\]
Their cokernel is $\mathbb Z/|n|$ for all parameters.
All these equalities are integral vector identities in
the retained matrices.

\item The degree-two boundary maps of the other fillings are
onto. At the cusp their cokernel is the free group
\eqref{maps:eq:D.I3}. Thus the full degree-two Mayer--Vietoris cokernel
is $\mathbb Z\oplus\mathbb Z/|n|$; the extension by
the free cusp defect splits. The degree-one source and
target both have free rank fourteen and one copy of
$\mathbb Z/2$. Their torsion map is the isomorphism
just checked, and their total cokernel is the finite
group $H_1(X)$. Their kernel is therefore zero.
This proves the asserted $H_2$. Four toric triple points
give Euler characteristic four; integral duality then
gives all the remaining groups.

\item \emph{Smooth recognition at $n=\pm1$.}
In this case $H_2(X)=\mathbb Z$, and the cusp map
$H_2(U)\to H_2(X)$ is onto, because \eqref{maps:eq:D.I4} kills the
entire regular-family contribution. The local Mumford
threefold has trivial canonical line. Hence $c_1(X)$
evaluates as zero on every $H_2$ class, so $c_1(X)=0$
and $X$ is spin.

\item The angular $\delta$ action is semifree with two fixed
rational curves. The one-node pieces are free as before.
At the finite logarithmic fibre, its invariant
$\gamma^*$ coordinate changes by $n/m$ under a deck
generator; for $n=\pm1$ this prevents any additional
stabilizer of the $\delta$ action. At the cusp there
are two vertical-edge orbits, each with the normal
weights and degrees \eqref{maps:eq:D.S11}.

\item It remains essential to check that these two fixed curves
have zero homology class; specialization is not onto in
\eqref{maps:eq:D.I3}. Colour a lattice vertex $(a,b)$ by $a+b$ modulo
two to label its normalized component. For a vertical
edge from $(a,b)$ to $(a,b+1)$, the two adjacent third
vertices are $(a+1,b)$ and $(a-1,b+1)$.
Each colour occurs once among the endpoints and once
among the third vertices. The toric wall formula thus
gives degree $-1+1=0$ for the line of either global
component on either fixed curve. Their images in
$H_2(U,M)$ are exactly these component intersection
numbers, so both images are zero. By the relative
sequence both curves lie in the boundary image, which
is precisely the regular-fibre specialization image.
Equation \eqref{maps:eq:D.I4}, with $n=\pm1$, now kills both classes
in $H_2(X)$.

\item The equivariant surgery and spin-lift argument of \ref{maps:src:D.2.14}
can therefore be performed on both disjoint curves.
Whitney isotopy applies also to their disjoint union,
and their normal framings are unique. It gives
\[
            Y\cong X\#_2(S^3\times S^3)
\]
with a free circle action and simply connected spin
quotient $B$. The homotopy sequence gives
$H_2(B)=\mathbb Z^2$ with primitive Euler class.
The oriented Smale--Barden theorem identifies the pair
with $((S^2\times S^3)\#(S^2\times S^3),(x,0))$.
Applying \cite[Theorem A and Theorem B(2)]{manGR25}
gives
\[
              Y\cong(S^2\times S^4)\#_2(S^3\times S^3).
\]
Zhubr cancellation as in \ref{maps:src:D.2.14} proves
$X\cong S^2\times S^4$ with its oriented smooth type.

\item The exact program \texttt{code/manifolds\_i2\_cusp.py}
verifies the primitive graph and exterior maps, the
nonzero boundary-torsion image, the actual saturated
vanishing lattice, the two fixed-curve degree vectors,
and 630 full $10\times10$ Smith forms.
The norm identities and \eqref{maps:eq:D.I4} prove the result without
any bounds on $k,n,r$.
\end{enumerate}
\end{proof}

\begin{theorem}[original odd-star fillings and all regular logs]\label{maps:src:D.2.16}
\StatementRecord{maps:src:D.2.16}{proved}{Historical events: Claude Fable 5.2 (Anthropic), effort high, 16 November 2026: carried as background without fault found. GPT-6 Astra (OpenAI), ChatGPT Pro mode, effort Pro, S6\_09\_PAPERS, 17 November 2026: targeted proof check. Also checked in this writing. Pre-continuation writing events: GPT-6 Astra (OpenAI), ChatGPT Pro mode, effort Pro, S6\_10\_PAPERS, 18 November 2026: targeted proof check. Continuation metadata events: GPT-6 Astra (OpenAI), ChatGPT Pro mode, effort Pro, S6\_10\_PAPERS continuation, 19 November 2026: used and rechecked in this run; status-vocabulary reconciliation; metadata reconciliation only; no new mathematical check. \emph{Pre-continuation recorded proof scope.} GPT-6 Astra (OpenAI, ChatGPT Pro mode), effort Pro, 2026-11-17 UTC; Claude Fable 5.2 (Anthropic), effort high, 16 November 2026; also checked in this writing by GPT-6 Astra (OpenAI), ChatGPT Pro mode, effort Pro, 18 November 2026}{the hypotheses, notation and definitions stated here; \cref{maps:src:D.2.14}; \cref{maps:src:D.2.12}; \cite[Section 2.4, (2.2)]{manUU21}; \cite[Sections 7.6--7.7]{manSS10}; \cite[Proposition 3.5(1),(3) and its proof]{manUU21}; \cref{maps:eq:D.T5}; \cref{maps:eq:D.T4}; \cref{maps:eq:D.T6}; \cref{maps:eq:D.T7}; \cref{maps:eq:D.T9}; \cref{maps:src:D.2.13}; \cref{maps:eq:D.T10}; \cite[Theorem A]{manGR25}; \cite[Theorem B(2)]{manGR25}}


Take one of the following marked configurations:
\[
\begin{array}{c|l|r|r}
 \mathrm{OS}&\text{ordered fibres}&p&b\\\hline
 46&I_3^*,I_1,I_1^\bullet,I_1&4&2\\
 50&I_1^*,I_3,I_1^\bullet,I_1&12&1.
\end{array}
\]
Set $P=pG$, $Q=G$, and place the degree-one linearization
zero at $I_1^\bullet$. Retain the original $I_j^*$ filling,
the primitive weight-zero nodal fillings and the hexagonal cusp.
At an additional regular fibre perform any free log of order
$m>0$ with vector $v=a e_1+b' e_2+c\delta+n\gamma$,
where $\gcd(m,a,b',c,n)=1$. Define
\[
 t=\gcd(m,n,a,b')\,\gcd(m,n,p,c),\qquad
 G=\mathbb Z/t\oplus\mathbb Z/m\oplus\mathbb Z/(m|n|/t)
 \quad(n\ne0).
\]
For $n\ne0$ the complete answer is
\[
\begin{aligned}
 \pi_1(X)&=H_1(X)=\mathbb Z/|n|, &H_5(X)&=0,\\
 H_2(X)&=\mathbb Z^b\oplus G,
 &H_3(X)&=\mathbb Z^{2b}\oplus G,\\
 H_4(X)&=\mathbb Z^b\oplus\mathbb Z/|n|.
\end{aligned}
\tag{D.T1}\label{maps:eq:D.T1}
\]
For $n=0$, put $T=\mathbb Z/t\oplus\mathbb Z/m$; then
\[
\begin{aligned}
 \pi_1(X)&=H_1(X)=H_5(X)=\mathbb Z,\\
 H_2(X)&=\mathbb Z^{b+2}\oplus T,
 &H_3(X)&=\mathbb Z^{2b+2}\oplus T,\\
 H_4(X)&=\mathbb Z^{b+2}.
\end{aligned}
\tag{D.T2}\label{maps:eq:D.T2}
\]
In both cases $H_0=H_6=\mathbb Z$ and $\chi(X)=2$.
All simply connected members, exactly those with $n=\pm1$,
have the oriented smooth type
\[
 X\cong \#^b(S^2\times S^4)\ \#\ \#^b(S^3\times S^3)
                   \ \#\ \widetilde\Sigma_0 N_m,
 \tag{D.T3}\label{maps:eq:D.T3}
\]
where the precisely framed twisted suspension and the spin
Smale manifold $N_m$ are as in \ref{maps:src:D.2.14}. In particular
$\widetilde\Sigma_0N_1=S^6$, so $m=1$ gives the corresponding
connected sums of sphere products.
\end{theorem}


\begin{proof}
\begin{enumerate}[label=\arabic*.,leftmargin=*]
\item \emph{The original filling, including its component action.}
The generator heights are $1/4$ and $1/12$ respectively.
In each case $P$ is narrow everywhere, $Q$ generates the
order-four component group at $I_j^*$, and in the second
case $Q$ also generates the component group of $I_3$.
Consequently the line $M=\mathcal O(P-O)$ has zero
multidegree at every fibre. As in \ref{maps:src:D.2.12} the line
$M\otimes t_Q^*M^{-1}$ is vertical and has zero multidegree,
hence is pulled back from the base; its degree is the mixed
height $\langle P,Q\rangle=1$. The specified linearization
exists. The $I_3$ quotient is the actual primitive nodal
model, by its unit component action. The simple zero on
$I_1^\bullet$ gives the primitive hexagonal cusp.

\item Here is an integral topological model of the original star
filling. Let $E\to\Delta$ be the minimal elliptic-surface
neighbourhood with fibre $I_j^*$, $j=3$ or $1$.
It retracts onto the affine $D_{j+4}$ tree of $j+5$ spheres.
A section in the identity component has a holomorphic local
logarithm. Explicitly the local exponential sequence is
$0\to\mathcal G\to\omega^{-1}\to E^{\rm id}\to0$,
and its stalk $\mathcal G_0$ is zero at additive reduction;
see \cite[Section 2.4, (2.2)]{manUU21}.
Thus choose $R$ in the identity component with $4R=4Q$.
Then $Q_0=Q-R$ has exact order four and the same component
as $Q$. The path $Q_0+sR$, $0\le s\le1$, is defined using
that logarithm and induces an isotopy of fibre-preserving
translations on the entire minimal surface. These
translations extend over all components, as also recalled in
\cite[Sections 7.6--7.7]{manSS10}. The order-four invariant
section can alternatively be read directly from
\cite[Proposition 3.5(1),(3) and its proof]{manUU21}.

\item Likewise the narrow section $P$ can be contracted to zero
through its local logarithm. Topologically this also follows
from the zero component degrees of $M$: $H^2(E;\mathbb Z)$
is free on those degrees, so $M$ is trivial. A unit in its
contraction lift has a logarithm on $E$, since $H^1(E;\mathbb Z)=0$.
Straightening this unit and the two section logarithms gives
a smooth family of contracting quotients with marked boundary.
A uniformly sufficiently small contraction parameter makes
the whole family proper and free. Ehresmann's theorem with
boundary therefore identifies the original filling with
\[
 U\cong S^1_\delta\times T(t_{Q_0}:E\to E),
 \tag{D.T4}\label{maps:eq:D.T4}
\]
where $T$ denotes the mapping torus. Integral fibre markings
are transported through this family, so this is a statement
about the actual original filling and its boundary map.
The line deformation is rigidified along $O$, and the unit
logarithm is chosen in that rigidification. Thus every member
has the same origin disc. On its regular fibres these are
deformations of group extensions and their translating period;
integral period loops are transported continuously, so their
integer monodromy is constant. The origin disc and fixed
period loops therefore give the transported meridian and sweep
marking, without an additional angular winding.

\item Translation by $Q_0$ cycles the four multiplicity-one end
components. The incidence tree then forces it to reflect the
chain of $j+1$ double components. All component maps preserve
complex orientation. Its permutation on $H_2(E)$ therefore
has
\[
                  \rho=1+(j+1)/2=b+1
\]
orbits. Invariants and coinvariants of a permutation lattice
are free on its orbits. The mapping-torus sequence and the
product with $S^1$ give, integrally,
\[
 \pi_1(U)=\mathbb Z^2,\qquad
 (H_0,H_1,H_2,H_3,H_4)(U)
  =(\mathbb Z,\mathbb Z^2,\mathbb Z^{\rho+1},
                         \mathbb Z^{2\rho},\mathbb Z^\rho).
 \tag{D.T5}\label{maps:eq:D.T5}
\]

\item \emph{The full boundary kernel and the free defect.}
In the global $(e_1,e_2,\delta,\gamma)$ marking, the star
monodromies are, respectively,
\[
 A_*^{46}=\begin{pmatrix}
 -1&-3&0&1\\0&-1&0&1\\-4&-8&1&3\\0&0&0&1
 \end{pmatrix},\qquad
 A_*^{50}=\begin{pmatrix}
 -1&-1&0&-1\\0&-1&0&1\\-12&-24&1&3\\0&0&0&1
 \end{pmatrix}.
 \tag{D.T6}\label{maps:eq:D.T6}
\]
In particular $(A_*^2-I)$ has rank one: these are not
finite-order monodromies. Both $A_*-I$ and
$\bigwedge^2A_*-I$ have only unit nonzero Smith entries,
with ranks two and four. For $M_*=\partial U$ the Wang
sequences give $H_1(M_*)=\mathbb Z^3$ and
$H_2(M_*)=\mathbb Z^4$. The map on $H_1$ is onto the
$\mathbb Z^2$ in \eqref{maps:eq:D.T5}, with primitive kernel generated
by the meridian, which bounds the extended origin disc.
Thus Poincare--Lefschetz duality and the relative sequence
imply
\[
 H_2(U)/\operatorname{im}H_2(M_*)=\mathbb Z^{\rho-1}
                                           =\mathbb Z^b.
 \tag{D.T7}\label{maps:eq:D.T7}
\]
The assertion is integral, since the map from
$H_2(U,M_*)=H^4(U)=\mathbb Z^\rho$ onto the meridian
kernel $\mathbb Z$ is surjective.

\item The two circle actions in \eqref{maps:eq:D.T4} give the entire normal-torus
kernel. One is $\delta$; the other is translation in the
mapping-torus coordinate with period four, using $Q_0^4=1$.
They sweep the extended origin disc, so their boundary tori
bound. Their orbit vectors on a regular fibre form the full
fixed lattice
\[
 \ker(A_*-I)=\mathbb Z\delta+\mathbb Z\ell,\qquad
 \ell=\begin{cases}(-1,2,0,4),&\mathrm{OS46},\\
                   (-3,2,0,4),&\mathrm{OS50}.
       \end{cases}
 \tag{D.T8}\label{maps:eq:D.T8}
\]
For an explicit check of the product marking, the integral
angular coordinate in \eqref{maps:eq:D.T4} is
$\delta^*-2e_1^*-e_2^*$ or
$\delta^*-6e_1^*-9e_2^*$ respectively. It is invariant
under \eqref{maps:eq:D.T6}, is one on $\delta$, and is zero on $\ell$.
The minimal nonzero $\gamma$-coordinate of a fixed vector
is four, exactly the period of the second action.
Consequently the two normal tori project isomorphically to
the full fixed lattice in the Wang sequence. They form a
primitive rank-two subgroup of $H_2(M_*)$. By \eqref{maps:eq:D.T5}--\eqref{maps:eq:D.T7}
the whole kernel has rank two, so they are its entire kernel.
It follows at once that the remaining image equals the
regular-fibre specialization image; its cokernel is \eqref{maps:eq:D.T7}.
This argument retains the star fibre's integral component
information. Replacing it by a product with an untwisted
elliptic surface, or by an order-two good-reduction filling,
would give a different boundary problem.

\item \emph{Global integral calculation for every affine vector.}
Insert $A_0=I$ before the four original monodromies. Set
$D=[A_0-I\ A_1-I\ A_2-I\ A_3-I]$, and use the direct
boundary maps $j_0,\ldots,j_3$ and
\[
 j_4(x)=-(A_1A_2A_3x,A_2A_3x,A_3x,x).
\]
The ten unchanged normal columns consist of two at the
star fibre, three at each of the two primitive nodal fibres,
and two at the cusp. Let their $16\times10$ matrix be $B$.
The exact constant integral certificates give
\[
 K=[B\ j_0(e_1)\ j_0(e_2)\ j_0(\gamma)],\qquad
 LK=I_{13},\qquad
 (Lj_0x)_{11,12,13}=(x_1,x_2,x_4).
 \tag{D.T9}\label{maps:eq:D.T9}
\]
The matrix $D$ has rank three with unit Smith entries;
$K$ is an integral basis of its kernel. The exterior
coinvariants are free of rank one, detected primitively by
$\omega=p e_1^*e_2^*-\delta^*\gamma^*$.
All ten normal tori are the marked fixed-loop sweeps proved
above and in \ref{maps:src:D.2.12}. The constant vertical form has zero
value on them. Thus \eqref{maps:eq:D.T9}, together with a fibre class of
$\omega$-value one, specifies the integral splitting,
without an unspecified fibre-area correction.

\item The original degree-one vanishing lattices generate the
primitive $\ker\gamma^*$. All original meridians bound.
Van Kampen therefore leaves only $\gamma$, and the regular
log imposes exactly $n\gamma=0$ as in \ref{maps:src:D.2.13}.
After eliminating the ten unchanged normals, the entire
remaining torsion presentation is
\[
 C=\begin{pmatrix}
 -pb'&pa&n&-c&0\\
 m&0&0&0&a\\
 0&m&0&0&b'\\
 0&0&0&m&n
 \end{pmatrix}.
 \tag{D.T10}\label{maps:eq:D.T10}
\]
Indeed the smooth-log lattice is
$\Lambda'=\Lambda+\mathbb Z(v/m)$, with primitive relation
$mg=v$; its four normal tori have coordinates
$(-\omega(x,v),j_0(mx))$, and the fifth lattice generator
has coordinates $(0,j_0v)$. This is the same integral
calculation as \ref{maps:src:D.2.13}, with $p=4$ or $12$.
Its all-parameter local-ring Smith proof gives entries
$(1,t,m,m|n|/t)$ for $n\ne0$, and $(1,t,m,0)$ for $n=0$.
The $b$ additional generators from \eqref{maps:eq:D.T7} are free and
split off, since their quotient is free.

\item The degree-one Mayer--Vietoris source and target are free
of the same rank nineteen. For $n\ne0$ their finite cokernel
has order $|n|$, so their kernel is zero. For $n=0$ the
cokernel and kernel each have free rank one; the latter
supplies the additional split free generator in $H_2$.
Thus $H_2$ is exactly as stated in \eqref{maps:eq:D.T1}--\eqref{maps:eq:D.T2}.
The original star filling has Euler zero by \eqref{maps:eq:D.T4}, and
the only nonzero Euler contribution is two from the cusp.
Integral duality, including its torsion terms, now gives
all the other displayed groups. In particular when $n=0$
one has $b_1=b_5=1$, $b_2=b_4=b+2$ and $b_3=2b+2$.

\item \emph{The exact smooth types.}
Assume $n=\pm1$. The free part of $H_2(X)$ is represented
in the original star filling. Its canonical line is trivial:
the canonical two-form on the minimal elliptic neighbourhood
is nonvanishing and translation-invariant, and its wedge
with $dw/w$ descends through the contraction. Thus $c_1$
evaluates trivially on every free generator. Since $X$ is
simply connected, $H^2(X;\mathbb Z)$ is torsion-free, so
$c_1(X)=0$ and $X$ is spin.

\item The global $\delta$ action is free on the original star
fill: in its contracting line model a nontrivial contraction
cannot equal a rotation. It is free on every other filling
except for the single cusp rational curve with normal
weights $(1,-1)$. This curve is in the cusp's fibre
specialization image, hence represents a multiple of the
regular class $s$. The normal relation $ns=0$ in \eqref{maps:eq:D.T10}
makes it zero. The equivariant nullsphere surgery of \ref{maps:src:D.2.14}
therefore gives a projectably spin free-circle manifold
$Y=X\#(S^3\times S^3)$.
Its simply connected spin orbit $B_5$ has
\[
 H_2(B_5)=\mathbb Z^{b+1}\oplus(\mathbb Z/m)^2,
\]
with primitive Euler class. Smale--Barden classifies the
oriented pair as
\[
 B_5=\#^{b+1}(S^2\times S^3)\#N_m,
 \qquad e=(x,0,\ldots,0).
\]
Apply \cite[Theorem A]{manGR25} successively, keeping the
primitive Euler class on the first $S^2\times S^3$ summand.
Its total space is $S^3\times S^3$. Each additional
$S^2\times S^3$ contributes its twisted suspension, which
is $(S^2\times S^4)\#(S^3\times S^3)$ by
\cite[Theorem B(2)]{manGR25}; $N_m$ contributes
$\widetilde\Sigma_0N_m$. Cancel the common $S^3\times S^3$
using the full Zhubr invariants exactly as in \ref{maps:src:D.2.14}.
This proves \eqref{maps:eq:D.T3}, including even $m$ and its framing.

\item The program \texttt{manifolds\_star\_smooth.py} retains all
constant marked matrices and integer left inverses,
checks \eqref{maps:eq:D.T6}--\eqref{maps:eq:D.T10}, and verifies 1,462 full $14\times15$
Smith forms, including 194 cases with $n=0$. The unrestricted
conclusion follows from the geometry and symbolic Smith
proof, not a bound on those parameters.
\end{enumerate}
\end{proof}

\begin{corollary}[two further original additive profiles]\label{maps:src:D.2.17}
\StatementRecord{maps:src:D.2.17}{proved}{Historical events: Claude Fable 5.2 (Anthropic), effort high, 16 November 2026: carried as background without fault found. GPT-6 Astra (OpenAI), ChatGPT Pro mode, effort Pro, S6\_09\_PAPERS, 17 November 2026: targeted proof check. Also checked in this writing. Pre-continuation writing events: GPT-6 Astra (OpenAI), ChatGPT Pro mode, effort Pro, S6\_10\_PAPERS, 18 November 2026: targeted proof check. Continuation metadata events: GPT-6 Astra (OpenAI), ChatGPT Pro mode, effort Pro, S6\_10\_PAPERS continuation, 19 November 2026: used and rechecked in this run; status-vocabulary reconciliation; metadata reconciliation only; no new mathematical check. \emph{Pre-continuation recorded proof scope.} GPT-6 Astra (OpenAI, ChatGPT Pro mode), effort Pro, 2026-11-17 UTC; Claude Fable 5.2 (Anthropic), effort high, 16 November 2026; also checked in this writing by GPT-6 Astra (OpenAI), ChatGPT Pro mode, effort Pro, 18 November 2026}{the hypotheses, notation and definitions stated here; \cref{maps:eq:D.T1}; \cref{maps:eq:D.T3}; \cref{maps:src:D.2.16}; \cref{maps:eq:D.T10}}


The complete formulas and oriented diffeomorphism classification
\eqref{maps:eq:D.T1}--\eqref{maps:eq:D.T3}, for every free smooth-fibre log, also hold for
\[
\begin{array}{c|c|l|r|r}
 \mathrm{OS}&\mathrm{profile}&\text{ordered fibres}&p&b\\\hline
 46&II&I_3^*,II,I_1^\bullet&4&2\\
 50&IV&I_1^*,IV,I_1^\bullet&12&1.
\end{array}
\]
Use $P=pG$, $Q=G$, the simple linearization zero shown,
and the original fillings at both additive fibres.
\end{corollary}


\begin{proof}
\begin{enumerate}[label=\arabic*.,leftmargin=*]
\item The section heights and component residues are the same at
the star fibre as in \ref{maps:src:D.2.16}. At $II$, both sections are
narrow. At $IV$, $P$ is narrow and $Q$ has residue two,
a generator of its order-three component group.
The mixed height is again one. The same zero-multidegree
argument gives the required degree-one linearization, with
no zero at either additive fibre.

\item For $II$, contract both section logarithms to zero. The
marked original filling is $E_{II}\times T^2$, where
$E_{II}$ retracts onto its cuspidal rational curve,
topologically $S^2$. For $IV$, divide $3Q$ in the additive
identity component as in \ref{maps:src:D.2.16}. The resulting section
$Q_0$ has order three and cycles all three components of
$IV$. Its central curve is three spheres meeting at one
point, so $H_2(E_{IV})=\mathbb Z^3$ and $H_1(E_{IV})=0$.
The marked fill is
$S^1_\delta\times T(t_{Q_0}:E_{IV}\to E_{IV})$.
The local logarithm, translation extension and boundary
rigidification arguments apply verbatim. In either case
there is one component orbit, and hence
\[
 \pi_1(U)=\mathbb Z^2,\qquad
 (H_0,H_1,H_2,H_3,H_4)(U)
       =(\mathbb Z,\mathbb Z^2,\mathbb Z^2,\mathbb Z^2,\mathbb Z).
\]

\item The two additive monodromies in the respective global
markings are
\[
 A_{II}=\begin{pmatrix}
 1&1&0&0\\-1&0&0&0\\0&0&1&0\\0&0&0&1
 \end{pmatrix},\qquad
 A_{IV}=\begin{pmatrix}
 -1&1&0&-2\\-1&0&0&0\\-24&0&1&-16\\0&0&0&1
 \end{pmatrix}.
\]
Both $A-I$ and $\bigwedge^2A-I$ have only unit nonzero
Smith entries, of ranks two and four. Thus their boundaries
have $H_1=\mathbb Z^3$, $H_2=\mathbb Z^4$. The two circle
actions sweep the origin disc and give the complete
primitive normal-torus kernel, with fixed vectors
\[
 \delta,\gamma\quad\text{at }II;\qquad
 \delta,(-2,2,0,3)\quad\text{at }IV.
\]
The second vector has minimal positive $\gamma$-coordinate
three, exactly the period of the order-three translation.
The relative sequence, with $H_2(U,\partial U)=\mathbb Z$
mapping onto the meridian kernel, shows that the boundary
specialization is onto $H_2(U)$. Thus these additional
original fillings contribute no free defect or torsion.

\item There are now four punctures including the regular log.
The graph differential has domain $\mathbb Z^{12}$ and
rank three with primitive image. The six unchanged normal
columns, together with $j_0(e_1),j_0(e_2),j_0(\gamma)$,
form a primitive basis of its rank-nine kernel. Its integer
left inverse has the same residual projection
$(x_1,x_2,x_4)$. Exterior coinvariants are still detected
primitively by $p e_1^*e_2^*-\delta^*\gamma^*$.
Consequently the entire residual normal matrix is \eqref{maps:eq:D.T10},
with the same free star defect $\mathbb Z^b$.
The degree-one Mayer--Vietoris source and target now have
rank fourteen; their kernel is zero for $n\ne0$ and is
$\mathbb Z$ for $n=0$. The full homology therefore follows
by the identical unrestricted Smith proof and duality.

\item Finally the global $\delta$ action remains free on both
original additive fillings. Its unique fixed curve is
still the null-homologous cusp curve for $n=\pm1$.
The free $H_2$ remains represented by the star neighbourhood,
where the canonical line is trivial. Thus the projectably
spin surgery, Smale--Barden orbit classification, repeated
Galaz-Garc\'ia--Reiser Theorem A and Zhubr cancellation in
\ref{maps:src:D.2.16} apply without modification and give \eqref{maps:eq:D.T3}.

\item The independent program
\texttt{manifolds\_star\_additive.py} records the full
constant integral certificates and 1,462 full $10\times11$
Smith checks, including 194 cases with $n=0$. The proof is
unrestricted in the affine vector and multiplicity.
\end{enumerate}
\end{proof}

\begin{theorem}[the remaining profile with a three-component cusp]\label{maps:src:D.2.18}
\StatementRecord{maps:src:D.2.18}{proved}{Historical events: Claude Fable 5.2 (Anthropic), effort high, 16 November 2026: carried as background without fault found. GPT-6 Astra (OpenAI), ChatGPT Pro mode, effort Pro, S6\_09\_PAPERS, 17 November 2026: targeted proof check. Also checked in this writing. Pre-continuation writing events: GPT-6 Astra (OpenAI), ChatGPT Pro mode, effort Pro, S6\_10\_PAPERS, 18 November 2026: targeted proof check. Continuation metadata events: GPT-6 Astra (OpenAI), ChatGPT Pro mode, effort Pro, S6\_10\_PAPERS continuation, 19 November 2026: used and rechecked in this run; status-vocabulary reconciliation; metadata reconciliation only; no new mathematical check. \emph{Pre-continuation recorded proof scope.} GPT-6 Astra (OpenAI, ChatGPT Pro mode), effort Pro, 2026-11-17 UTC; Claude Fable 5.2 (Anthropic), effort high, 16 November 2026; also checked in this writing by GPT-6 Astra (OpenAI), ChatGPT Pro mode, effort Pro, 18 November 2026}{the hypotheses, notation and definitions stated here; \cref{maps:src:D.2.16}; \cref{maps:src:D.2.15}; \cref{maps:src:D.1.2}; \cref{maps:src:D.2.17}; \cref{maps:eq:D.V5}; \cref{maps:eq:D.T10}; \cref{maps:eq:D.V4}; \cref{maps:src:D.2.14}}


For OS50 profile $II$, with ordered fibres $II,I_3,I_1^*$,
take $P=12G$, $Q=G$ and the degree-one linearization vanishing
on the whole $I_3$ fibre. Keep the original $II$ and $I_1^*$
fillings. At an additional regular fibre perform the free
order-$m$ log with vector $(a,b',c,n)$, where
$\gcd(m,a,b',c,n)=1$. Define $t$ and $G$ as in \ref{maps:src:D.2.16} with
$p=12$. For $n\ne0$ one has
\[
\begin{aligned}
 \pi_1(X)&=H_1(X)=\mathbb Z/|n|, &H_5(X)&=0,\\
 H_2(X)&=\mathbb Z^3\oplus G,
 &H_3(X)&=\mathbb Z^2\oplus G,\\
 H_4(X)&=\mathbb Z^3\oplus\mathbb Z/|n|.
\end{aligned}
\tag{D.V1}\label{maps:eq:D.V1}
\]
For $n=0$, put $T=\mathbb Z/t\oplus\mathbb Z/m$; the answer is
\[
\begin{aligned}
 \pi_1(X)&=H_1(X)=H_5(X)=\mathbb Z,\\
 H_2(X)&=\mathbb Z^5\oplus T,
 &H_3(X)&=\mathbb Z^4\oplus T,\\
 H_4(X)&=\mathbb Z^5.
\end{aligned}
\tag{D.V2}\label{maps:eq:D.V2}
\]
In both cases $H_0=H_6=\mathbb Z$ and $\chi(X)=6$.
Every simply connected member, exactly $n=\pm1$, has
\[
 X\cong \#^3(S^2\times S^4)\#(S^3\times S^3)
                          \#\widetilde\Sigma_0N_m
 \tag{D.V3}\label{maps:eq:D.V3}
\]
as an oriented smooth manifold.
\end{theorem}


\begin{proof}
\begin{enumerate}[label=\arabic*.,leftmargin=*]
\item \emph{The actual three-component cusp.}
The generator height is $1/12$. The section $P$ is narrow
at every fibre, while $Q$ has component residue two at
$I_3$ and is a generator at $I_1^*$. The mixed height is
one, so the same multidegree and vertical-line argument
as in \ref{maps:src:D.2.16} gives the required linearization.
The actual valuation periods at its simple zero are
\[
 \Gamma=\langle(3,0),(2,1)\rangle,
 \qquad
 \begin{pmatrix}1&1\\0&-1\end{pmatrix}
 \begin{pmatrix}3&2\\0&1\end{pmatrix}
       =\begin{pmatrix}3&3\\0&-1\end{pmatrix}.
 \tag{D.V4}\label{maps:eq:D.V4}
\]
Thus the same integral symmetry of the triangular fan as
in \ref{maps:src:D.2.15} identifies this lattice with
$3\mathbb Z\oplus\mathbb Z$. The smooth Mumford filling
has three normalized $dP_6$ components. The integral
normalization calculation of \ref{maps:src:D.1.2} gives
\[
 (H_0,H_1,H_2,H_3,H_4)(U_c)
  =(\mathbb Z,\mathbb Z^2,\mathbb Z^6,\mathbb Z^2,\mathbb Z^3),
 \qquad \pi_1(U_c)=\mathbb Z^2.
\]
The original $I_1^*$ filling has the order-four translation
model of \ref{maps:src:D.2.16}, with two component orbits and free
specialization defect one. The original $II$ filling has
the product model of \ref{maps:src:D.2.17} and no defect.

\item The cusp monodromy in the global marking is
\[
 A_c=\begin{pmatrix}
 7&12&0&-8\\-3&-5&0&4\\48&96&1&-63\\0&0&0&1
 \end{pmatrix}.
\]
The nonzero Smith entries of $A_c-I$ are $(1,3)$,
whereas those of $\bigwedge^2A_c-I$ are $(1,1)$.
Thus its boundary $M_c$ has
$H_1(M_c)=\mathbb Z^3\oplus\mathbb Z/3$ and
$H_2(M_c)=\mathbb Z^6$.
Its $H_1$ map kills precisely the meridian and the
$\mathbb Z/3$ angular torsion, and is onto $\mathbb Z^2$.
The relative sequence with
$H_2(U_c,M_c)=H^4(U_c)=\mathbb Z^3$ therefore gives
\[
 H_2(U_c)/\operatorname{im}H_2(M_c)=\mathbb Z^2.
 \tag{D.V5}\label{maps:eq:D.V5}
\]
Indeed this quotient is the kernel of a surjection from
$\mathbb Z^3$ to $\mathbb Z\oplus\mathbb Z/3$, hence is
free of rank two. The two toric angular normal tori bound
and project onto the full primitive fixed lattice of
$A_c$. The kernel of the degree-two boundary map has
rank two by \eqref{maps:eq:D.V5}, so these are its whole kernel.
The image is precisely the regular-fibre specialization
image. The total free local defect is consequently
$\mathbb Z^2\oplus\mathbb Z=\mathbb Z^3$.

\item \emph{The integral global map and its three-torsion.}
Insert $A_0=I$ before the three original monodromies.
The global fibre-coinvariant presentation has nonzero
Smith entries $(1,1,3)$. The cusp's boundary torsion maps
to this global $\mathbb Z/3$ by a unit; the exact integral
vector and its image are retained in the certificate.
Thus these torsion summands cancel isomorphically in
Mayer--Vietoris. They must not be discarded beforehand.
Replacing the cusp vanishing image by its actual saturated
angular lattice gives three unit Smith entries and kills
precisely $\ker\gamma^*$. All original meridians bound,
and van Kampen leaves exactly $n\gamma=0$.

\item The six unchanged normal tori consist of two from each
original filling. Their graph columns together with
$j_0(e_1),j_0(e_2),j_0(\gamma)$ form a primitive integral
basis of the rank-nine graph kernel. Its left inverse
has residual projection $(x_1,x_2,x_4)$, including zero
for the $\delta$ sweep. The exterior coinvariants are
free of rank one with primitive detector
$12e_1^*e_2^*-\delta^*\gamma^*$.
The actual fixed-loop sweeps specify the same integral
zero-area splitting as in \ref{maps:src:D.2.16}. Consequently after the
six unchanged normals are eliminated, the residual
regular-log presentation is exactly \eqref{maps:eq:D.T10} with $p=12$.
Its all-parameter Smith calculation supplies $G$ for
$n\ne0$, and $\mathbb Z\oplus T$ for $n=0$.
The free local defect \eqref{maps:eq:D.V5} and its star counterpart
split off. After cancelling the isomorphic degree-one
$\mathbb Z/3$ summands, the remaining source and target
have equal free rank fourteen. Their kernel is zero for
$n\ne0$ and is free of rank one for $n=0$, giving the
additional split generator in \eqref{maps:eq:D.V2}.
Only the three-component cusp contributes Euler
characteristic, namely six. Duality now proves all of
\eqref{maps:eq:D.V1}--\eqref{maps:eq:D.V2}, with the indicated distinct free ranks
in degree three.

\item \emph{Fixed-curve classes and smooth recognition.}
For $n=\pm1$, the free $H_2$ is generated by classes in
the cusp and star neighbourhoods. Both canonical lines
are locally trivial, by the Mumford volume form and
the translation-invariant form used in \ref{maps:src:D.2.16}.
Hence $c_1(X)=0$ and $X$ is spin.

\item There are three $\delta$-fixed rational curves, represented
by vertical edges of the triangular fan modulo $\Gamma$.
Their classes in the free cusp defect are zero, as the
following integral computation shows. Colour a vertex
$(a,b)$ by $a+b$ modulo three; this labels its global
component by \eqref{maps:eq:D.V4}. A vertical edge has endpoints
$(a,b),(a,b+1)$ and the two adjacent third vertices
$(a+1,b),(a-1,b+1)$. The colours of the endpoints are
$(c,c+1)$; those of the third vertices are $(c+1,c)$.
The toric wall formula therefore gives degree zero
against each global component divisor. Poincare--Lefschetz
identifies these degrees with the map
$H_2(U_c)\to H_2(U_c,M_c)=\mathbb Z^3$.
Each fixed curve is consequently in the boundary image,
which equals the regular-fibre image. It represents a
multiple of the regular class $s$, killed by the relation
$ns=0$. All three fixed curves are therefore null-homologous
in $X$, and hence null-homotopic. This checks their actual
integral classes, rather than inferring nullity from their
normal bundles.

\item Each normal bundle has the equivariant model
$\mathcal O(-1)_{+1}\oplus\mathcal O(-1)_{-1}$.
The three-component embedded link is isotopic to three
unknotted spheres: apply the dimension-six embedding
isotopy theorem of \ref{maps:src:D.2.14} to the disjoint union.
The same orientation-preserving equivariant surgeries
produce
$Y=X\#^3(S^3\times S^3)$ with a projectably spin free
circle action. Its simply connected spin orbit has
$H_2(B_5)=\mathbb Z^4\oplus(\mathbb Z/m)^2$ and primitive
Euler class. Thus Smale--Barden gives the oriented pair
\[
 (B_5,e)=\bigl(\#^4(S^2\times S^3)\#N_m,
                                      (x,0,0,0,0)\bigr).
\]
Repeated Galaz-Garc\'ia--Reiser Theorem A and Theorem B(2)
identify its total space with
$\#^3(S^2\times S^4)\#^4(S^3\times S^3)
\#\widetilde\Sigma_0N_m$.
Full Zhubr cancellation of three copies of $S^3\times S^3$
proves \eqref{maps:eq:D.V3}.

\item The program \texttt{manifolds\_star\_i3\_cusp.py} checks
the integral period lattice, all three fixed-curve degree
vectors, the boundary torsion map, saturated vanishing
lattice, full graph/exterior certificates, and 731 full
$10\times11$ Smith forms, including 97 cases with $n=0$.
The geometry and symbolic reduction prove the result
for every admissible affine vector and multiplicity.
\end{enumerate}
\end{proof}

\begin{proposition}[finite initial data for the whole-fibre problem]\label{maps:src:D.2.19}
\StatementRecord{maps:src:D.2.19}{proved}{Historical events: Claude Fable 5.2 (Anthropic), effort high, 16 November 2026: carried as background without fault found. GPT-6 Astra (OpenAI), ChatGPT Pro mode, effort Pro, S6\_09\_PAPERS, 17 November 2026: targeted proof check. Also checked in this writing. Pre-continuation writing events: GPT-6 Astra (OpenAI), ChatGPT Pro mode, effort Pro, S6\_10\_PAPERS, 18 November 2026: targeted proof check. Continuation metadata events: GPT-6 Astra (OpenAI), ChatGPT Pro mode, effort Pro, S6\_10\_PAPERS continuation, 19 November 2026: used and rechecked in this run; status-vocabulary reconciliation; metadata reconciliation only; no new mathematical check. \emph{Pre-continuation recorded proof scope.} GPT-6 Astra (OpenAI, ChatGPT Pro mode), effort Pro, 2026-11-17 UTC; Claude Fable 5.2 (Anthropic), effort high, 16 November 2026; also checked in this writing by GPT-6 Astra (OpenAI), ChatGPT Pro mode, effort Pro, 18 November 2026}{the hypotheses, notation and definitions stated here; \cref{maps:src:D.2.1}; \cref{maps:eq:D.F2}; \cite[Section 11.8]{manSS10}; \cref{maps:src:D.2.11}}


Restrict to the uncorrected line
$M=\mathcal O(P-O)$ and require its linearization line
$J=M\otimes t_Q^*M^{-1}$ to be $f^*\mathcal O_{\mathbb P^1}(1)$.
This is the case of a simple zero on one whole fibre.
Write $P=aG$, $Q=bG$ and $h(G)=1/d$.
Up to simultaneous sign reversal, the possible ordered
pairs are exactly the positive pairs $ab=d$ for which,
at every reducible fibre, $P$ or $Q$ is narrow there.
For the marked profiles of \ref{maps:src:D.2.1} they are
\[
\begin{array}{c|r|l}
 \mathrm{OS}&d& (a,b)\\\hline
 43&2&(1,2),(2,1)\\
 45&8&(1,8),(8,1)\\
 46&4&(1,4),(4,1)\\
 47&14&(1,14),(2,7),(7,2),(14,1)\\
 49&6&(1,6),(2,3),(3,2),(6,1)\\
 50&12&(1,12),(3,4),(4,3),(12,1)\\
 55&20&(1,20),(4,5),(5,4),(20,1)\\
 56&30&(1,30),(2,15),(3,10),(5,6),\\
 &&(6,5),(10,3),(15,2),(30,1).
\end{array}
\tag{D.F1}\label{maps:eq:D.F1}
\]
The same row applies to every profile of its OS entry.
Across all 21 marked profiles there are 104 positive
ordered factor-pair candidates, exactly 88 of which
satisfy the stated requirement. Including each marked
multiplicative fibre and one regular-fibre position for
the simple zero gives 348 finite initial data. All their
marked integral monodromy tuples are recorded exactly.
These counts classify initial data in the stated scope,
not diffeomorphism types or all logarithmic fillings.
\end{proposition}


\begin{proof}
\begin{enumerate}[label=\arabic*.,leftmargin=*]
\item On a smooth elliptic fibre, translation acts trivially
on $\operatorname{Pic}^0$, so $J$ is generically trivial
and hence is represented by a vertical divisor.
At a reducible fibre, label its multiplicity-one
components by the Neron component group $\Phi$.
The multidegree of $M$ on those components is
$e_{\bar P}-e_0$. Translation by $Q$ acts regularly
on this component group, and therefore its pullback
preserves that vector if and only if
\[
                         \bar P=0\quad\hbox{or}\quad\bar Q=0.
 \tag{D.F2}\label{maps:eq:D.F2}
\]
Indeed, if $\bar P\ne0$, its unique positive and unique
negative entries must both be fixed; the regular
translation fixes a component only when $\bar Q=0$.
If $P$ is narrow, every component degree is zero.
If $Q$ is narrow, its translation fixes every component
(the connected identity group cannot permute a finite
set). Thus \eqref{maps:eq:D.F2} is also sufficient for all multidegrees
of $J$ to vanish, including the multiple components.

\item A vertical divisor of zero multidegree is an integral
combination of whole fibres: each fibre intersection
matrix has radical generated by its multiplicity vector,
and that vector has an entry one because the surface has
a section. Hence precisely under \eqref{maps:eq:D.F2}, $J=f^*L$ for
a line $L$ on the base. Intersecting with $O$ gives
\[
 \deg L=P\cdot O+1+Q\cdot O-P\cdot Q
        =\langle P,Q\rangle=ab/d.
\]
For the first equality use $O^2=-1$ and translation
invariance of intersection, so
$(P-Q)\cdot O=P\cdot Q$ and $(-Q)\cdot O=Q\cdot O$.
For the second, every local correction in the height
pairing vanishes under \eqref{maps:eq:D.F2}; see
\cite[Section 11.8]{manSS10}.
Thus $J=f^*\mathcal O(1)$ is equivalent to \eqref{maps:eq:D.F2} and
$ab=d$. Inversion on the elliptic surface takes
$(P,Q,M,J)$ to $(-P,-Q,[-1]^*M,[-1]^*J)$ and preserves
complex orientation. We may consequently take $a,b>0$,
without identifying the two ordered factors.

\item The finite calculation now needs no parameter bound.
For each of the 21 profiles, list the positive divisors
$a\mid d$, put $b=d/a$, and evaluate both section
component maps. Applying \eqref{maps:eq:D.F2} gives exactly \eqref{maps:eq:D.F1}.
The 16 excluded marked pairs fail the requirement for
this particular uncorrected line; this is not a
nonexistence theorem for other line bundles.
For each admitted pair, place the single weight one
at each multiplicative site in turn and at one
additional regular site. The exact program
\texttt{manifolds\_degree\_one.py} gives 348 tuples,
checks every matrix is integral with determinant one,
and verifies every ordered monodromy product is $I$.
It retains all component data and all matrices, as well
as the additive sites and their finite-part orders.

\item There are precise boundaries to this reduction.
Component-supported linearizations are not covered:
\ref{maps:src:D.2.11} already has $\langle G,3G\rangle=1/2$ and
$J=\mathcal O(C_0)$, rather than a whole-fibre zero.
Adding an integral vertical divisor to $M$ changes its
multidegrees and can alter \eqref{maps:eq:D.F2}.
Nor does the finite list bound the logarithmic orders
or affine clutching vectors, choose a compactification
at an additive zero, or classify all Mumford subdivisions.
These remain geometric and integral-attachment choices
in the full D.2 problem. This proposition provides an
exhaustive finite set of starting data only under its
explicit whole-fibre and uncorrected-line hypotheses.
\end{enumerate}
\end{proof}

\begin{theorem}[six original-filled families, including the untwisted Engel case]\label{maps:src:D.2.20}
\StatementRecord{maps:src:D.2.20}{proved}{Historical events: Claude Fable 5.2 (Anthropic), effort high, 16 November 2026: carried as background without fault found. GPT-6 Astra (OpenAI), ChatGPT Pro mode, effort Pro, S6\_09\_PAPERS, 17 November 2026: targeted proof check. Also checked in this writing. Pre-continuation writing events: GPT-6 Astra (OpenAI), ChatGPT Pro mode, effort Pro, S6\_10\_PAPERS, 18 November 2026: targeted proof check. Continuation metadata events: GPT-6 Astra (OpenAI), ChatGPT Pro mode, effort Pro, S6\_10\_PAPERS continuation, 19 November 2026: used and rechecked in this run; status-vocabulary reconciliation; metadata reconciliation only; no new mathematical check. \emph{Pre-continuation recorded proof scope.} GPT-6 Astra (OpenAI, ChatGPT Pro mode), effort Pro, 2026-11-17 UTC; Claude Fable 5.2 (Anthropic), effort high, 16 November 2026; also checked in this writing by GPT-6 Astra (OpenAI), ChatGPT Pro mode, effort Pro, 18 November 2026}{the hypotheses, notation and definitions stated here; \cref{maps:src:D.2.19}; \cref{maps:src:D.2.16}; \cref{maps:eq:D.O4}; \cref{maps:eq:D.O5}; \cref{maps:eq:D.O8}; \cref{maps:src:D.3.3}; \cref{maps:eq:D.O7}; \cref{maps:src:D.2.14}; \cite[Theorem A and Theorem B(2)]{manGR25}}


Use the following configurations and section choices:
\[
\begin{array}{c|c|l|r}
 \mathrm{OS}&\mathrm{profile}&\text{fibres}&d\\\hline
 43&\mathrm{default}&III^*,3I_1&2\\
 43&II&III^*,II,I_1&2\\
 46&\mathrm{default}&I_3^*,3I_1&4\\
 46&II&I_3^*,II,I_1&4\\
 49&III&IV^*,III,I_1&6\\
 50&IV&I_1^*,IV,I_1&12.
\end{array}
\]
Take $P=G$, $Q=dG$, retain every original additive filling,
and put the simple linearization zero at the first marked
$I_1$. The other $I_1$ fillings are unchanged. At an
additional regular fibre perform any free log
\[
 m>0,\qquad v=a e_1+b e_2+c\delta+n\gamma,
 \qquad\gcd(m,a,b,c,n)=1.
\]
Set
\[
        g=\gcd(m,a,b,n),\qquad
        T_{m,g}=\mathbb Z/g\oplus(\mathbb Z/m)^2.
\]
For $n\ne0$, the answer is
\[
\begin{aligned}
 \pi_1(X)&=H_1(X)=\mathbb Z/|n|, &H_5(X)&=0,\\
 H_2(X)&=\mathbb Z^7\oplus T_{m,g},
 &H_3(X)&=\mathbb Z^{14}\oplus T_{m,g},\\
 H_4(X)&=\mathbb Z^7\oplus\mathbb Z/|n|.
\end{aligned}
\tag{D.O1}\label{maps:eq:D.O1}
\]
For $n=0$, it is
\[
\begin{aligned}
 \pi_1(X)&=H_1(X)=H_5(X)=\mathbb Z,\\
 H_2(X)&=\mathbb Z^8\oplus T_{m,g},
 &H_3(X)&=\mathbb Z^{14}\oplus T_{m,g},\\
 H_4(X)&=\mathbb Z^8.
\end{aligned}
\tag{D.O2}\label{maps:eq:D.O2}
\]
In all cases $H_0=H_6=\mathbb Z$ and $\chi(X)=2$.
The simply connected members are exactly $n=\pm1$, and
all of them have oriented diffeomorphism type
\[
 X\cong\#^7(S^2\times S^4)\#^7(S^3\times S^3)
                                    \#\widetilde\Sigma_0N_m.
 \tag{D.O3}\label{maps:eq:D.O3}
\]
In particular none is a six-sphere. For the Engel row,
the trivial regular log $(m,v)=(1,0)$ recovers the
untwisted groups of \cref{push:src:A.4.2} $(1,1,8,14,8,1,1)$ integrally,
with no torsion.
\end{theorem}


\begin{proof}
\begin{enumerate}[label=\arabic*.,leftmargin=*]
\item \emph{The original additive bundles.}
The section $Q$ is narrow everywhere. At every additive
fibre except $II$, $P$ meets a nonidentity component.
Proposition \ref{maps:src:D.2.19} supplies the degree-one linearization.
The zero on $I_1$ gives the same primitive hexagonal
cusp as before. This proof retains the original additive
bundles, including their nonzero component degrees.

\item Let $E\to\Delta$ be one such additive elliptic-surface
neighbourhood, with $r$ components. It is simply connected
and retracts onto a connected union of $r$ rational spheres
with $H_2(E)=\mathbb Z^r$ and no other positive homology.
This applies also to the tangent pair of type $III$ and
the three spheres of type $IV$, which each meet at one
point. The component degree vector of $M=\mathcal O(P-O)$
is the primitive vector
\[
                   c=e_{\bar P}-e_0\in H^2(E;\mathbb Z).
\]
Because $Q$ is narrow, its local logarithm contracts it
to zero through translations. Rigidify the line along
$O$ as in \ref{maps:src:D.2.16}. The unit in the lift has a logarithm
since $H^1(E)=0$, and can be straightened. The actual
original filling is therefore, with its transported
boundary marking, the principal $T^2$ bundle over $E$
whose two Chern classes are $(c,0)$. Its two circle
directions are denoted $\delta,\gamma$.

\item The integral bundle spectral sequence has only base
rows zero and two. Its nonzero differentials are
\[
 H_2(E)\longrightarrow H_1(T^2),\quad x\longmapsto c(x)\delta,
 \qquad
 H_2(E)\otimes H_1(T^2)\longrightarrow H_2(T^2),\quad
 x\otimes y\longmapsto c(x)\delta\wedge y.
\]
Both images are primitive because $c$ is primitive.
There are no higher differentials, and all surviving
groups and extensions are free. Consequently
\[
 \pi_1(U)=\mathbb Z\gamma,\qquad
 (H_0,H_1,H_2,H_3,H_4)(U)
  =(\mathbb Z,\mathbb Z,\mathbb Z^{r-1},
                             \mathbb Z^{2r-1},\mathbb Z^r).
 \tag{D.O4}\label{maps:eq:D.O4}
\]
In particular $H_2(U)\to H_2(E)$ identifies its image
with $\ker c$.

\item \emph{The marked fibre kernel.}
Write the upper-left elliptic monodromy block as $B_i$
and its $\gamma$ column as $q_i$. Local narrowness of $Q$
means that
\[
                 q_i=(B_i-I)z_i,\qquad z_i\in\mathbb Z^2.
\]
The logarithm deformation identifies the elliptic
projection of a regular four-torus by
$e_1\mapsto e_1$, $e_2\mapsto e_2$, $\delta\mapsto0$,
$\gamma\mapsto z_i$. Thus its area covector, in the
usual ordered exterior basis, is
\[
           w_i=(1,0,(z_i)_2,0,-(z_i)_1,0).
 \tag{D.O5}\label{maps:eq:D.O5}
\]
The specialization of the elliptic fibre in $H_2(E)$
is its component-multiplicity vector. This vector is
primitive, lies in $\ker c$, and generates the image
of $H_2(F)\to H_2(U)$. The latter map is exactly
$w_i$ times that primitive vector, by \eqref{maps:eq:D.O4}.
Therefore its kernel is the full integral $\ker w_i$.

\item The boundary monodromies in these six cases have two
unit nonzero Smith entries in $A_i-I$ and four in
$\bigwedge^2A_i-I$. Hence $H_1(\partial U)=\mathbb Z^3$
and $H_2(\partial U)=\mathbb Z^4$.
The two principal circle actions sweep the origin
disc, giving normal tori with orbit vectors
$\delta$ and $(-z_i,0,1)$. They form the full fixed
lattice of $A_i$. Thus these two sweeps split the
Wang projection integrally, and the remaining
boundary image equals the regular-fibre image just
computed. It follows that the local free defect is
\[
                H_2(U)/\operatorname{im}H_2(\partial U)
                         =\mathbb Z^{r-2}.
 \tag{D.O6}\label{maps:eq:D.O6}
\]
The full boundary kernel consists of the two actual
normal tori and the fibre kernel \eqref{maps:eq:D.O5}. This is why
using only the fixed-loop normal tori would lose
relations in this case.

\item The exterior coinvariants of the global regular family
are free of rank one, with primitive detector
$\omega=e_1^*e_2^*-d\delta^*\gamma^*$.
If $s$ has $\omega$-value one, the image of $\ker w_i$
in this group is $h_i\mathbb Zs$, where
\[
                  h_i=\gcd(d,(z_i)_1,(z_i)_2).
\]
Indeed an integral basis of $\ker w_i$ has $\omega$
values $0,-(z_i)_2,0,(z_i)_1,-d$.
Here all local inputs are explicit:
\[
\begin{array}{c|c|r|r|r}
 \mathrm{OS/profile}&\text{non-narrow-}P\text{ fibre}
       &r&z_i&h_i\\\hline
 43/\mathrm{default},43/II&III^*&8&(3,1)&1\\
 46/\mathrm{default},46/II&I_3^*&8&(1,-2)&1\\
 49/III&IV^*&7&(6,2)&2\\
       &III&2&(9,3)&3\\
 50/IV&I_1^*&6&(9,-6)&3\\
      &IV&3&(8,-8)&4.
\end{array}
\tag{D.O7}\label{maps:eq:D.O7}
\]
In every row of the theorem the $h_i$ together have gcd
one. Thus the actual original fillings kill $s$.
The original $II$ and ordinary $I_1$ pieces have the
product models already proved and add no defect.
The total free local defect $\beta$ is six in the first
four configurations and five in the last two.

\item \emph{The unrestricted global matrices.}
Insert $A_0=I$ for the regular log. If the total number
of punctures is $N$, the graph differential has domain
$\mathbb Z^{4(N-1)}$, rank three and primitive image.
Write $B$ for all unchanged fixed-loop normal columns,
including the two at the cusp. The exact constant
certificates give an integral basis
\[
 K=[B\ j_0(e_1)\ j_0(e_2)\ j_0(\gamma)]
\]
of the graph kernel, with left inverse $L$ whose last
three coordinates on $j_0(x)$ are $(x_1,x_2,x_4)$.
The actual circle sweeps above fix their zero fibre-area
coordinates. After these normal columns and the relations
$h_is=0$ are eliminated, the regular-family contribution
is freely generated by those three sweeps.

\item The smooth log has lattice $\Lambda+\mathbb Z(v/m)$,
with primitive relation $m(v/m)=v$.
Its normal tori have graph components $j_0(mx)$;
the additional lattice generator has component $j_0v$.
The fibre components no longer matter, since $s=0$.
The resulting presentation is exactly
\[
  \begin{pmatrix}
    m&0&0&a\\0&m&0&b\\0&0&m&n
  \end{pmatrix}.
 \tag{D.O8}\label{maps:eq:D.O8}
\]
Its determinantal divisors are $g,mg,m^2g$,
so its cokernel is $T_{m,g}$ for all integers in the
statement. Notice that $c$ enters the geometric
freeness condition even though it drops out of \eqref{maps:eq:D.O8}.

\item The original fillings kill all base meridians and
precisely the primitive $\ker\gamma^*$ in the fibre
fundamental group. Their complement over the punctured
base therefore has fundamental group $\mathbb Z\gamma$.
The new meridian relation is $m\xi-v=0$, and the base
relation kills $\xi$, leaving $n\gamma=0$. This proves
the asserted full fundamental group.

\item The degree-one Mayer--Vietoris map must now be counted
with the groups \eqref{maps:eq:D.O4}, not with the rank-two groups of
the narrow-$P$ filling. Its source and target are free.
For the first and third rows of the theorem their ranks
are $(19,18)$; for the second and fourth, $(14,13)$;
for the last two, $(14,12)$.
Its cokernel is the abelian group already computed by
van Kampen. Thus for $n\ne0$ its kernel rank is
$\varepsilon=1$ in the first four rows and two in the
last two. For $n=0$ this rank increases by one.
These free kernels split off in the degree-two sequence.
Since $\beta+\varepsilon=7$, this proves exactly the
stated $H_2$ groups, including their torsion.
All local pieces except the primitive cusp have Euler
zero, while the cusp contributes two. Integral duality
then gives \eqref{maps:eq:D.O1}--\eqref{maps:eq:D.O2} in every degree. In particular
the Engel row at $(m,v)=(1,0)$ independently reproduces
\cref{push:src:A.4.2}, including its integral freeness.

\item \emph{Canonical class, nullsphere surgery and smooth type.}
Let $F_m$ denote the reduced new multiple fibre.
The canonical-form computation in \ref{maps:src:D.3.3} applies to the
original pieces: before the regular log the canonical
line is $f^*\mathcal O(-1)$, and the primitive toroidal
cusp adds no divisor. A regular order-$m$ log adds
$(m-1)F_m$. Since $f^*(\mathrm{point})=mF_m$, one obtains
\[
                         K_X=\mathcal O_X(-F_m).
 \tag{D.O9}\label{maps:eq:D.O9}
\]
Here this formula also controls the topological class,
including the free classes supplied by the degree-one
Mayer--Vietoris kernel. The complement of $F_m$ has
$\pi_1=\mathbb Z\gamma$. Its primitive normal meridian
is $m\xi-v$, whose image is $-n\gamma$.
The residue sequence for the divisor therefore shows
that $[F_m]\in H^2(X;\mathbb Z)$ has exact order $|n|$
when $n\ne0$. For $n=0$ it is primitive of infinite
order: the quotient by its image embeds in the torsion-free
$H^2(X\setminus F_m)$, whose torsion vanishes by
$H_1(X\setminus F_m)=\mathbb Z$. For $n=\pm1$ it vanishes. Thus $c_1(X)=0$ in the
simply connected cases, and $X$ is spin; local canonical
triviality alone would not have sufficed for this step.

\item The angular $\delta$ action is free on every original
bundle and on the smooth log when $n=\pm1$. Its sole
fixed curve is the cusp rational curve with normal
weights $(1,-1)$. This curve belongs to the regular-fibre
specialization image and hence is a multiple of $s$.
By \eqref{maps:eq:D.O7}, $s=0$ already in the exterior after the
original additive fillings. The fixed curve is therefore
null-homologous, and then null-homotopic in a simply
connected member.
Apply the oriented equivariant nullsphere surgery and
projectable spin lift of \ref{maps:src:D.2.14}. It gives
$Y=X\#(S^3\times S^3)$ as a free circle bundle over a
simply connected spin five-manifold with
\[
            H_2(B_5)=\mathbb Z^8\oplus(\mathbb Z/m)^2
\]
and primitive Euler class. Smale--Barden identifies the
oriented pair as
\[
 (\#^8(S^2\times S^3)\#N_m,(x,0,\ldots,0)).
\]
Repeated \cite[Theorem A and Theorem B(2)]{manGR25}
give
\[
 Y\cong\#^7(S^2\times S^4)\#^8(S^3\times S^3)
                                        \#\widetilde\Sigma_0N_m.
\]
Cancellation by the complete Zhubr invariants proves
\eqref{maps:eq:D.O3}, including the specified twisted suspension
framing at even $m$.

\item The exact program \texttt{manifolds\_original\_additive.py}
retains all monodromies, fixed-loop bases, local covectors,
primitive fibre relations and full integral gluing maps.
It verifies 4,386 full degree-two Smith forms and 4,386
full degree-one Smith forms, including 582 cases with
$n=0$ in each degree. It separately checks the base-degree
map on bridge classes in both rank regimes. Formula
\eqref{maps:eq:D.O8} and the proved integral local maps supply the
unrestricted theorem.
\end{enumerate}
\end{proof}

\begin{theorem}\label{maps:src:D.2.21}
\StatementRecord{maps:src:D.2.21}{proved}{Historical events: Claude Fable 5.2 (Anthropic), effort high, 16 November 2026: carried as background without fault found. GPT-6 Astra (OpenAI), ChatGPT Pro mode, effort Pro, S6\_09\_PAPERS, 17 November 2026: targeted proof check. Also checked in this writing. Pre-continuation writing events: GPT-6 Astra (OpenAI), ChatGPT Pro mode, effort Pro, S6\_10\_PAPERS, 18 November 2026: targeted proof check. Continuation metadata events: GPT-6 Astra (OpenAI), ChatGPT Pro mode, effort Pro, S6\_10\_PAPERS continuation, 19 November 2026: used and rechecked in this run; status-vocabulary reconciliation; metadata reconciliation only; no new mathematical check. \emph{Pre-continuation recorded proof scope.} GPT-6 Astra (OpenAI, ChatGPT Pro mode), effort Pro, 2026-11-17 UTC; Claude Fable 5.2 (Anthropic), effort high, 16 November 2026; also checked in this writing by GPT-6 Astra (OpenAI), ChatGPT Pro mode, effort Pro, 18 November 2026}{the hypotheses, notation and definitions stated here; \cref{maps:src:D.2.20}; \cref{maps:src:D.2.5}; \cref{maps:src:D.2.12}; \cref{maps:src:D.2.14}; \cref{maps:src:D.2.16}}


Use the uncorrected degree-one data $P=G,Q=dG$ and the simple zero at the
marked $I_1$ in the following five rows. Keep the indicated original
additive filling and replace the other additive fibre by its finite-good
logarithmic model of order $m=qk$, with invariant power vector
$n\ell+r\delta$, where $k>0$, $r\in\mathbb Z$ and $\gcd(n,m)=1$.
\[
\begin{array}{c|c|c|r|r|r|r}
\mathrm{OS/profile}&d&\mathrm{finite\ log}\ (q)&\mathrm{original}&b&h&c\\\hline
49/III&6&IV^*\ (3)&III&1&3&2\\
49/III&6&III\ (4)&IV^*&6&2&3\\
50/IV&12&IV\ (3)&I_1^*&5&3&4\\
43/II&2&II\ (6)&III^*&7&1&2\\
46/II&4&II\ (6)&I_3^*&7&1&4
\end{array}
\]
Then $n\ne0$, $\pi_1(X)=H_1(X)=\mathbb Z/|n|$, and
\[
 H_2(X)=\mathbb Z^b,\qquad H_3(X)=\mathbb Z^{2b},\qquad
 H_4(X)=\mathbb Z^b\oplus\mathbb Z/|n|,\qquad H_5(X)=0.
\]
The endpoint groups are $\mathbb Z$, and $\chi(X)=2$.
Every member with $n=\pm1$ is oriented diffeomorphic to
\[
                 \#^b(S^2\times S^4)\ \#\ \#^b(S^3\times S^3).
\]
In particular none of these five simply connected families is $S^6$.

\end{theorem}
\begin{proof}
\begin{enumerate}[label=\arabic*.,leftmargin=*]
\item The component data, original local models and geometric
finite-good logarithmic replacements are those already proved in
\ref{maps:src:D.2.20} and \ref{maps:src:D.2.5}--\ref{maps:src:D.2.12}. We retain their markings rather than replacing
an original additive fibre by its invariant torus cohomology. Order the
three monodromies cyclically so that the finite-good one is first, and
write them $A_0,A_1,A_2$. Their product is the identity. The two remaining
full fixed lattices provide four primitive normal-torus vectors $B$.
If $\ell$ is the finite invariant vector of $\gamma$-coefficient one,
then $K=[B,j_0\ell]$ is a primitive integral basis of the rank-five
kernel of the graph differential. The exact integer inverses, monodromies,
component projections and complete matrices are saved in
\texttt{code/manifolds\_mixed\_additive.py} and its JSON certificate.

\item The regular degree-two coinvariants are freely generated by $s$, detected
by $\omega=e_1^*e_2^*-d\delta^*\gamma^*$. The original additive cap has
$r_0=b+1$ components and the actual principal $T^2$-bundle model of
\ref{maps:src:D.2.20}. Its local free defect is $r_0-2=b-1$. Its full fibre kernel maps
to $h\mathbb Zs$, with $h$ in the table. The one-component toric cusp
has no free defect. The finite cap has no free defect either.

\item Let $\beta$ generate the finite vertical coinvariant lattice and write
$\delta=q_\delta\beta$. Then
\[
 N_m\beta=(m/q_\delta)\delta,\qquad N_m\ell=m\ell,
 \qquad -\omega(\beta,\ell)=c=d/q_\delta,\quad
 -\omega(\ell,\delta)=-d.
\]
Choose $m\alpha+n\zeta=1$. The actual finite normal-circle tori are
\[
 \binom{-\omega(\beta,n\ell+r\delta)}{Lj_0(N_m\beta)},\qquad
 \binom{-\alpha\omega(\ell,n\ell+r\delta)}
 {Lj_0(\alpha N_m\ell+\zeta(n\ell+r\delta))}.
\]
After eliminating the four primitive columns $B$, the entire remaining
integral cokernel has the presentation
\[
                 \begin{pmatrix}h&cn&-dr\alpha\\0&0&1\end{pmatrix}.
\]
In each row $\gcd(h,cn)=1$: the nontrivial possibilities use respectively
$3\nmid n$, $2\nmid n$, $3\nmid n$; the last two have $h=1$.
Thus this cokernel vanishes for every allowed integer input, without an
order bound. The companion code verifies every constant primitive lattice
and $1,068$ full $6\times7$ Smith matrices as an independent certificate.

\item For completeness, the first-degree Mayer--Vietoris source has rank nine
and target rank eight: each boundary has rank three; the exterior has
rank three, while the finite, original and cusp caps have ranks two,
one and two. The original cap kills the elliptic and $\delta$ loops,
leaving $\gamma$. All original and cusp meridians vanish. The finite
power relation is therefore $n\gamma=0$. This proves
$\pi_1(X)=\mathbb Z/|n|$ directly and makes the rational first-degree
Mayer--Vietoris map onto. Its kernel is consequently free of rank one.
Together with the local defect $b-1$ and the zero cokernel above, this
proves $H_2(X)=\mathbb Z^b$. Integral duality, the universal coefficient
theorem and $\chi=2$ give the remaining groups as displayed.

\item When $n=\pm1$, the reduced finite fibre divisor is zero in integral
cohomology: the complement has fundamental group $\mathbb Z\gamma$,
and its primitive normal meridian maps to $\pm n\gamma$, so the divisor
residue sequence gives divisor order $|n|$. The canonical divisor differs
from a base pullback by a multiple of this reduced fibre; the full base
fibre is $m$ times the same divisor. Thus $c_1(X)=0$, including on the free
Mayer--Vietoris bridge class. The $\delta$ action is free on the original
principal bundle and on the finite quotient, since $\gcd(n,m)=1$.
Its only fixed component is the standard cusp sphere with normal weights
$+1,-1$. As in \ref{maps:src:D.2.20} this sphere lies in the image of the regular fibre,
hence is a multiple of $s$ and is nullhomologous by the unimodular
presentation. The equivariant surgery, spin lift and Smale--Barden
argument of \ref{maps:src:D.2.14} now give a free circle action on
$X\#(S^3\times S^3)$ with orbit
$\#^{b+1}(S^2\times S^3)$ and primitive Euler class.
Repeated Galaz-Garc\'ia--Reiser Theorem A and Theorem B(2), followed by
full Zhubr cancellation exactly as in \ref{maps:src:D.2.16}, give the asserted oriented
diffeomorphism type.\hfill$\square$
\end{enumerate}
\end{proof}


\begin{theorem}[finite initial data for sphere outputs]\label{maps:src:D.2.22}
\StatementRecord{maps:src:D.2.22}{proved}{Historical events: Claude Fable 5.2 (Anthropic), effort high, 16 November 2026: carried as background without fault found. GPT-6 Astra (OpenAI), ChatGPT Pro mode, effort Pro, S6\_09\_PAPERS, 17 November 2026: targeted proof check. Also checked in this writing. Pre-continuation writing events: GPT-6 Astra (OpenAI), ChatGPT Pro mode, effort Pro, S6\_10\_PAPERS, 18 November 2026: targeted proof check. Continuation metadata events: GPT-6 Astra (OpenAI), ChatGPT Pro mode, effort Pro, S6\_10\_PAPERS continuation, 19 November 2026: used and rechecked in this run; status-vocabulary reconciliation; metadata reconciliation only; no new mathematical check. \emph{Pre-continuation recorded proof scope.} GPT-6 Astra (OpenAI, ChatGPT Pro mode), effort Pro, 2026-11-17 UTC; Claude Fable 5.2 (Anthropic), effort high, 16 November 2026; also checked in this writing by GPT-6 Astra (OpenAI), ChatGPT Pro mode, effort Pro, 18 November 2026}{the hypotheses, notation and definitions stated here; \cref{maps:src:D.2.2}; \cref{maps:src:D.2.11}; \cref{maps:src:D.1.2}; \cite[Theorem 6.3 and Corollary 6.13]{manSS10}; \cref{maps:eq:D.FI3}; \cref{maps:eq:D.FI4}; \cref{maps:eq:D.FI5}}


Consider the standard section-line construction over any of the twenty-one
rational elliptic profiles with Mordell--Weil group $\mathbb Z$ in \ref{maps:src:D.2.2}.
Write $G$ for its chosen generator, $\langle G,G\rangle=1/d$,
$P=aG$, $Q=bG$, and $M=\mathcal O(P-O)$. Use a nonzero holomorphic
linearization with effective vertical divisor $D$, and the following local
fillings: a zero-free original quotient or a free affine logarithmic
replacement at each additive fibre; an original quotient or a standard
smooth Mumford filling at each multiplicative fibre; and, if desired,
free logarithmic replacements and rank-one Mumford fillings at regular
fibres. The latter occur at the regular zeros of the linearization.
These are the filling types specified in D.2; a new additive degeneration
at a zero of the linearization is not one of these local types.

If the completed manifold is a homotopy six-sphere, then $D$ is one
reduced component of one multiplicative fibre $I_n$. Translation by $Q$
acts transitively on its components and
\[
             \langle P,Q\rangle=\frac1n,
             \qquad ab=\frac dn.                                      \tag{D.FI1}\label{maps:eq:D.FI1}
\]
After simultaneously changing the signs of $P,Q$, take $a,b>0$.
There are exactly $172$ marked triples (profile, ordered section pair,
zero position) satisfying the resulting integral component criterion.
They are listed completely in
\texttt{records/manifolds\_sphere\_natural.json}.
This is a finite list of possible initial data, not a claim that all
its filling choices are spheres.

The same necessary reduction holds in the larger construction with
$M=\mathcal O(P-O+V)$ for an arbitrary integral vertical divisor $V$.
Modulo simultaneous translation of the chosen zero component, there
are $226$ marked (profile, ordered section pair, zero-fibre) candidates;
all $226$ component equations have integral solutions $V$.
Their remaining solution freedom is exactly the $Q$-invariant vertical
divisor lattice. This optional enlargement is distinguished from the
natural section-line construction in the preceding paragraph.

\end{theorem}


\begin{proof}
\begin{enumerate}[label=\arabic*.,leftmargin=*]
\item We give the local Euler identity, the rational rank obstruction, and the
finite integral enumeration separately. No logarithmic multiplicity is
bounded in any step.

\item \emph{1. The local line and the determinant.}
Near a multiplicative fibre choose a Tate coordinate with $q=t^n$;
holomorphic unit factors do not affect any valuation below. If $P=t^p$
times a unit, choose $0\le p<n$. The product defining the Tate theta
function gives
\[
 \theta(qz;q)=-z^{-1}\theta(z;q),\qquad
 \sigma(z)=\frac{\theta(z/P;q)}{\theta(z;q)},\qquad
 \sigma(qz)=P\sigma(z).
\]
Thus $\sigma$ is the meromorphic section with divisor $P-O$ of the
line with factor of automorphy $P$. At a generic point of component
$i$, represented by $z=t^i$ times a unit, its valuation in the flat
frame is
\[
 f_i=\min(0,i-p)\quad(0\le i<n),\qquad f_{i+n}=f_i+p.
\]
Indeed, in this range only the factor $1-z/P$ of the numerator can
have negative valuation; the other theta factors are units. The
meromorphic section has no vertical divisor in $\mathcal O(P-O)$.
Consequently the flat frame, expressed in an extending local frame,
has valuation $F_i=-f_i$. Adding $V=\sum v_iC_i$ changes this to
$F_i=-f_i+v_i$, where $v_{i+n}=v_i$. Put $A=-p$, so
$F_{i+n}=F_i+A$.

\item If translation by $Q$ shifts the component index by $B$, let $w$ be
the valuation of its flat linearization multiplier. Comparing the
extending frames in $M$ and $t_Q^*M$ gives, with no rational
intersection-theory ambiguity,
\[
 z_i=\operatorname{ord}_{C_i}(D)=w+F_i-F_{i+B}.
\]
For a quasi-periodic sequence with increment $A$, telescoping yields
\[
 \sum_{i=0}^{n-1}(F_{i+B}-F_i)=AB,
 \qquad \sum_i z_i=nw-AB.                                    \tag{D.FI2}\label{maps:eq:D.FI2}
\]
The identity holds for every integer $B$, by division by $n$ and
then telescoping the remaining finite interval. The actual valuation
vectors of the two long periods are $(n,-A)$ and $(B,-w)$: the latter
uses the inverse of the lift $M_{Qz}\to M_z$. Reversing the second
short-circle coordinate changes them to $(n,A),(B,w)$. Their index
is therefore $|nw-AB|$. For an effective nonzero local divisor this
equals $\sum_i z_i$. A standard unimodular triangular Mumford
cellulation has twice this many zero-dimensional torus strata, and
all its positive-dimensional strata have Euler number zero. Hence
\[
                  \chi(U_{I_n})=2\sum_i z_i.                  \tag{D.FI3}\label{maps:eq:D.FI3}
\]
This proves the local equality directly from the extending line
frames, rather than identifying an arbitrary support function with
the line. For example, $n=2,p=B=1,w=0$ gives
$F=(1,0)$, $A=-1$, $(z_0,z_1)=(1,0)$, and index one, precisely
the component-supported zero of \ref{maps:src:D.2.11}.

\item \emph{2. Local rational facts used in the rank obstruction.}
Let $F=T^4$, $A_i$ its monodromy at a deleted point, $L_i$ the
mapping-torus boundary, and $U_i$ a filling. Put
\[
 f_i=\dim H_1(F;\mathbb Q)_{A_i},\quad
 g_i=\dim H_2(F;\mathbb Q)_{A_i},\quad h_i=b_1(U_i).
\]
For every allowed local type, nearby-fibre specialization generates
$H_1(U_i;\mathbb Q)$; a multiple meridian is a fibre loop. Thus
$h_i\le f_i$. Here are the actual local models and the degree-two
fact needed below.

\item A rank-one Mumford cap is topologically a cycle of nodal elliptic
components times $T^2$: the degree-zero ruled line has topologically
trivial Chern class, and its translation gluing is isotopic to the
identity. The surviving circle of the nodal curve, its fibre-area
class, and the two transverse circles show that
$H^2(U_i;\mathbb Q)\to H^2(F;\mathbb Q)^{A_i}$ is onto.
For a rank-two triangular cap, take the universal toric model with
rays $(v,1)$, $v\in\mathbb Z^2$, and the standard periodic triangular
fan. Its period lattice $L\subset\mathbb Z^2$ acts by
\[
 (z_1,z_2,t)\longmapsto
 (u_{\ell,1}(t)t^{\ell_1}z_1,
  u_{\ell,2}(t)t^{\ell_2}z_2,t),\qquad \ell\in L.
\]
Here $u:L\to\mathcal O^*(\Delta)^2$ is a character. It extends to
$\mathbb Z^2$: put the inclusion in Smith form and take holomorphic
roots of its two nowhere-zero unit values on the disc. The extended
period action uses the same fan and shifts its dual complex by
$\mathbb Z^2$. Its quotient is a minimal hexagonal cap. The original
cap covers this quotient with deck group $\mathbb Z^2/L$. This
action is free even on the central fibre, because its action on
$\mathbb R^2/L$ is a nontrivial translation. This is an ordinary
finite covering, not the toric power map discussed in C.1.
Pull back the four independent specialization classes of \ref{maps:src:D.1.2}. On the regular torus
the pullback is an isomorphism over $\mathbb Q$ onto the invariant
space, of dimension four; hence it again gives surjectivity.
Subdividing the toric cellulation preserves these pulled-back classes.
For free finite affine quotients, rational transfer and averaging
give the same assertion. This includes higher logarithmic orders,
the good-reduction caps, and the semistable covers of odd $I_n^*$.

\item There is a separate check for an original cap; it cannot be replaced
without checking its component Chern class. Its central model is
the circle bundle of $M$ over the mapping torus of the component
translation of the Kodaira curve. Since the linearization is a unit
there, its multidegree vector $c$ is invariant under that translation.
If $c\ne0$, the rational circle-bundle transgression kills the
$\delta$ circle, so $h_i=f_i-1$. If $c=0$, let $\rho$ be the number
of component orbits. On the invariant component subspace, the
Kodaira intersection matrix is negative semidefinite with kernel
the whole fibre and rank $\rho-1$. The map
$H_2(U_i)\to H_2(U_i,L_i)$ has this rank: pair its image with the
component four-cycles (the component times the two circle factors).
Thus the cokernel of $H_2(L_i)\to H_2(U_i)$ has dimension at least
$\rho-1$. If $c=0$ and $\rho=1$, the invariant fibre-area class
and the circle classes give the same degree-two surjectivity as
above. These statements also follow directly for a multiplicative
curve from its cycle of spheres times $T^2$; for an additive curve
use its tree of spheres. Finite component permutations are handled
by averaging, which is legitimate over $\mathbb Q$.

\item \emph{3. The global rank budget.}
Delete $N$ discs containing every singular fibre, every zero of the
linearization, and every extra logarithmic site, and write $E$ for
the resulting torus bundle. Its fibre coinvariants in degree one
have dimension $r=1$. To verify this without assuming a special
section pair, observe that $\mathcal O(aG-O)$ differs from
$\mathcal O(G-O)^{\otimes a}$ by a vertical line, trivial on this
punctured base. The elliptic-plus-$\delta$ monodromy block therefore
scales its last row by $a$. For $a\ne0$ this does not change its
rational image rank. The exact twenty-one canonical blocks saved
by the program all have rank three. All full monodromy differences
lie in these first three coordinates, and hence have precisely
this image. Vertical modifications and the $Q$ column cannot
change the conclusion.

\item Also $v_2=\dim H_2(F;\mathbb Q)_{\pi_1}\ge1$. Explicitly,
in the marking $(e_1,e_2,\delta,\gamma)$ the actual monodromy
has the following blocks, where $z$ is the section cocycle of $G$:
\[
 A=\begin{pmatrix}
 B&0&bz\\ -az^tJB&1&c\\0&0&1
 \end{pmatrix},\qquad
 J=\begin{pmatrix}0&1\\-1&0\end{pmatrix},\qquad B\in SL_2(\mathbb Z).
\]
The same vertical-line trivialization gives the row multiplied by
$a$, while the Abel--Jacobi lift of $Q=bG$ gives the column $bz$.
The central entry $c$ is unrestricted. Direct block multiplication
shows that $A^t\Omega A=\Omega$, where
\[
 \Omega=\begin{pmatrix}aJ&0\\
 0&\begin{matrix}0&-b\\ b&0\end{matrix}\end{pmatrix}.
\]
In the only potentially nonzero cross term,
$B^tJz=(-z^tJB)^t$, so the two terms cancel. Thus
$\Omega=a e_1^*e_2^*+b\gamma^*\delta^*$ is a nonzero invariant,
independently of $c$ and of any vertical modification. The program
also verifies this identity symbolically. The case $a=0$ or $b=0$ cannot have an effective
nonempty $D$: norm over a component-preserving translation makes
such a divisor a base divisor of degree $\langle P,Q\rangle=0$.
It therefore cannot give Euler number two.

\item The graph-base homology sequence gives
\[
 b_1(E)=N,\qquad b_2(E)=v_2+4N-7,\qquad b_2(L_i)=g_i+f_i.
\]
If $H_1(X;\mathbb Q)=H_2(X;\mathbb Q)=0$, the first MV map is
onto and its kernel has dimension $\sum_i(f_i-h_i)$. This kernel
is a quotient of $H_2(X;\mathbb Q)$, so every $h_i=f_i$.
In particular every original cap has $c=0$. The degree-two map
from the boundary is also onto each cap; the component intersection
calculation then forces $\rho=1$ for every original cap.
All remaining caps consequently satisfy the degree-two invariant
surjectivity proved in step 2, and so $b_2(U_i)\ge g_i$.
After first mapping onto $\bigoplus H_2(U_i)$, the boundary kernel
has dimension at most $\sum f_i$. It must still surject onto
$H_2(E)$, since $H_2(X)=0$. Therefore
\[
 v_2+4N-7\le\sum_i f_i,
 \qquad\sum_i\operatorname{rank}(A_i-I)\le7-v_2\le6.          \tag{D.FI4}\label{maps:eq:D.FI4}
\]

\item Let $F_a,F_m$ count additive and multiplicative fibres of the
rational elliptic surface. The Shioda--Tate formula, $\rho(S)=10$,
and Mordell--Weil rank one give root rank seven. Comparing this
with $\chi(S)=12$ gives
\[
                            2F_a+F_m=5.
\]
This is also checked independently for every profile in the program;
see \cite[Theorem 6.3 and Corollary 6.13]{manSS10}.
Every additive monodromy difference has rank at least two, and
every multiplicative one rank at least one. A rank-two Mumford
zero at a multiplicative fibre raises its rank to two; a zero at
a regular fibre contributes an additional rank one. All original
and logarithmic caps have Euler number zero, as do rank-one caps.
Thus $\chi(X)=2$ and \eqref{maps:eq:D.FI3} force at least one rank-two zero,
already consuming the sole extra rank allowed by \eqref{maps:eq:D.FI4}.
There are no regular zeros and no second multiplicative zero.
Equation \eqref{maps:eq:D.FI3} now says the effective divisor at that fibre has
total coefficient one. It is one reduced component, as asserted.

\item \emph{4. Norm, height, and finite enumeration.}
Choose an integer $R$ such that $RQ$ meets the identity component
of every reducible fibre. Norm the line identity over the first
$R$ iterates of $t_Q$. The resulting divisor has zero component
multidegrees, and is therefore a sum of whole fibres. Since it is
effective and supported on the orbit of the single component,
$Q$ must act transitively on this $I_n$. Its norm is $(R/n)F$.
The normed line is
$M\otimes t_{RQ}^*M^{-1}$; the vertical correction cancels because
$RQ$ fixes every component. The height formula with a narrow second
section gives the base degree
$\langle P,RQ\rangle=R\langle P,Q\rangle$; all its local correction
terms vanish. Hence \eqref{maps:eq:D.FI1}. The divisor and height comparison is
integral: a vertical divisor of zero multidegree is a whole-fibre
multiple, by the corank-one component intersection matrix, and
$\operatorname{Pic}^0(S)=0$.

\item Here is the entire finite algorithm, with no implicit parameter
bound. For each multiplicative $I_n$ with $n\mid d$, list the
positive ordered divisors $ab=d/n$. Require that the component
of $Q$ generate $\mathbb Z/n$. For each Kodaira fibre let $I$ be
its integral component intersection matrix, $T$ the component
pullback permutation of $Q$, and $c_P=e_{P\text{-component}}-e_0$.
For a candidate divisor vector $z$, solve
\[
              (1-T)I V=Iz-(1-T)c_P.                         \tag{D.FI5}\label{maps:eq:D.FI5}
\]
At the chosen zero fibre, $z$ is one component; elsewhere $z=0$.
For the natural construction require $V=0$ and test every possible
zero component. For the larger construction solve \eqref{maps:eq:D.FI5} by
unimodular Smith operations. Since $I$ is invertible modulo the
invariant whole fibre, its solution kernel is exactly $\ker(1-T)$.
No additional numerical condition is omitted: \eqref{maps:eq:D.FI5} makes
$J\otimes\mathcal O(-D)$ a base pullback, and its norm has degree
zero by \eqref{maps:eq:D.FI1}, so the base pullback is trivial. Thus every retained
row actually has this effective linearization divisor.

\item The program checks the affine component diagrams, their primitive
multiplicity kernels and the translation permutations, then solves
every \eqref{maps:eq:D.FI5} integrally. The following counts retain marked fibre
positions; the first number uses $V=0$, the second allows $V$:
\[
\begin{array}{c|rr@{\qquad}c|rr}
43&6&6&43/II&2&2\\
45&8&17&45/II&4&9\\
46&6&9&46/II&2&3\\
47&14&16&47/II&6&8\\
47/III&8&10&49&10&10\\
49/II&2&2&49/III&4&4\\
50&10&15&50/II&2&3\\
50/IV&4&6&55&14&23\\
55/II&6&11&56&26&28\\
56/II&10&12&56/III&14&16\\
56/IV&14&16&\text{total}&172&226
\end{array}
\]
For $n>1$, the natural component equation in fact forces
$P$ and $Q$ to meet the same adjacent component, labelled $1$ or
$-1$, and the zero to be $C_0$; the code tests this directly,
including the doubled edges for $I_2$. The broader count normalizes
the zero component to $C_0$ by a power of $Q$; translating $M$
changes it only by an allowed vertical divisor. This explains
the different marking convention in the two counts.
\end{enumerate}
\end{proof}

\paragraph{Scope after \ref{maps:src:D.2.22}.}
The primary questions have different breadth. Viaclovsky's
Question 4 \cite[Section 1.1]{manViaclovskyLogs26} asks which starting configurations admit infinitely
many sphere multiplicities. Engel's Remark 8.5 \cite{manEngel26} expressly allows
an arbitrary fibrewise degree-zero line bundle, not only
$\mathcal O(P-O)$. We keep those scopes distinct. \ref{maps:src:D.2.22} removes
any need for an unexplained bound on section pairs or zero
positions when searching for spheres in the stated local sector.
The remaining task is the integral, all-order gluing calculation
for the finitely many initial rows and all allowed filling choices.
The $172$ and $226$ counts do not by themselves complete that task
or classify all nonsphere outputs.


\begin{theorem}[uniform all-order fundamental groups on the natural catalogue]\label{maps:src:D.2.23}
\StatementRecord{maps:src:D.2.23}{proved}{Historical events: Claude Fable 5.2 (Anthropic), effort high, 16 November 2026: carried as background without fault found. GPT-6 Astra (OpenAI), ChatGPT Pro mode, effort Pro, S6\_09\_PAPERS, 17 November 2026: targeted proof check. Also checked in this writing. Pre-continuation writing events: GPT-6 Astra (OpenAI), ChatGPT Pro mode, effort Pro, S6\_10\_PAPERS, 18 November 2026: targeted proof check. Continuation metadata events: GPT-6 Astra (OpenAI), ChatGPT Pro mode, effort Pro, S6\_10\_PAPERS continuation, 19 November 2026: used and rechecked in this run; status-vocabulary reconciliation; metadata reconciliation only; no new mathematical check. \emph{Pre-continuation recorded proof scope.} GPT-6 Astra (OpenAI, ChatGPT Pro mode), effort Pro, 2026-11-17 UTC; Claude Fable 5.2 (Anthropic), effort high, 16 November 2026; also checked in this writing by GPT-6 Astra (OpenAI), ChatGPT Pro mode, effort Pro, 18 November 2026}{the hypotheses, notation and definitions stated here; \cref{maps:src:D.2.22}; \cref{maps:src:D.2.19}}


For every one of the $172$ natural initial rows of \ref{maps:src:D.2.22}, retain any
of the stated original or Mumford fillings and perform any admissible
free affine logarithmic replacements. At an odd $I_n^*$, a logarithmic
replacement here means the free cyclic quotient of its rank-one
semistable Mumford cover. Regular logarithmic sites are allowed.
Let $m_i\ge1$ be their orders, $v_i$ their integral power vectors in
$H_1(T^4;\mathbb Z)$, and put $n_i=\gamma^*(v_i)$.
For a chosen system of peripheral lifts, let $e$ be the integer in
the global relation of their surviving $\gamma$ coordinates.
Then the complete fundamental group is
\[
 \pi_1(X)=\left\langle \gamma,c_1,\ldots,c_s\ \left|\
 [\gamma,c_i]=1,\ c_i^{m_i}=\gamma^{n_i},\
 c_1\cdots c_s=\gamma^e\right.\right\rangle .                 \tag{D.FG1}\label{maps:eq:D.FG1}
\]
The canonical lifts from the construction, obtained by unrolling the
translation quotient before filling, have $e=0$. Changing $c_i$ to
$c_i\gamma^{t_i}$ replaces $(n_i,e)$ by
$(n_i+m_it_i,e+\sum t_i)$ and leaves \eqref{maps:eq:D.FG1} unchanged.
Thus the formula also applies directly to other marked twist conventions.
The first integral homology is the Smith cokernel of the $(s+1)$ square
presentation matrix with rows
\[
 (-n_i,0,\ldots,m_i,\ldots,0)\quad(1\le i\le s),
 \qquad (-e,1,\ldots,1).                                    \tag{D.FG2}\label{maps:eq:D.FG2}
\]
There is no multiplicity bound and no restriction to invariant rational
representatives of the affine translations: their actual power vectors
are used.

After eliminating order-one sites by $c_i=\gamma^{n_i}$, replace $e$
by $e-\sum_{m_i=1}n_i$. If at least three sites of order greater than
one remain, the group is nontrivial. With zero, one, or two remaining
sites, it is abelian and is trivial precisely when, respectively,
\[
 |e|=1,\qquad |m_1e-n_1|=1,\qquad
 |m_2n_1+m_1n_2-em_1m_2|=1.                                \tag{D.FG3}\label{maps:eq:D.FG3}
\]
In the two-site case its invariant factors are the gcd of the entries
of the two-by-two relation matrix and its determinant divided by that
gcd; when the determinant vanishes use one free summand. In particular,
simple connectivity forces coprime remaining multiplicities and
$\gcd(m_i,n_i)=1$ at every remaining site.

\end{theorem}


\begin{proof}
\begin{enumerate}[label=\arabic*.,leftmargin=*]
\item We supply the integral vanishing comparison, since rational monodromy
alone would not prove this statement. The exact program
\texttt{manifolds\_sphere\_gluing.py} constructs all $172$ integral
monodromy tuples. Every zero-free point is locally narrow for at least
one of $P,Q$, so its lift is the one already proved in \ref{maps:src:D.2.19}.
The matrix at the one component-supported zero is independently fixed
by the global product relation. It agrees integrally with the normalized
weight $1/n$ derived in \ref{maps:src:D.2.22}. The program verifies every product,
all local Smith forms, and the preserved alternating form.

\item At a multiplicative Mumford cap the primitive vanishing sublattice is
$\operatorname{sat}\operatorname{im}(A_i-I)$: its individual primitive
circles bound the toric vanishing thimbles. The zero cap has rank two.
Every other multiplicative cap in all $172$ rows has rank one and
Smith width one; there is no nonprimitive necklace hidden in this
assertion. If instead the original zero-free quotient is retained,
these same circles are still nullhomotopic. When the line multidegree
is zero, the local line trivialization identifies the elliptic vanishing
thimble with this primitive circle. When it is nonzero, the natural
line has component degree vector $e_{P\text{-component}}-e_0$;
its entries have gcd one, so the component two-spheres kill the
$\delta$ circle integrally. The elliptic vanishing thimble then kills
its primitive elliptic coordinate as well, and hence kills the required
mixed circle. These comparisons do not replace an integral image by
its saturation without a geometric bounding disc.

\item Close these vanishing lattices under all monodromy matrices. Exact
unimodular operations show that their normal closure is
$\ker\gamma^*$ in $170$ rows. In the two rows
\[
 (46/II,a,b)=(46/II,1,4),\ (46/II,4,1)
\]
its quotient in $\ker\gamma^*$ is $\mathbb Z/3$. The $I_3^*$ cap
kills this last quotient. An original cap kills its elliptic vanishing
loops. For a semistable logarithmic cap, take its cyclic base cover.
The upstairs rank-one Mumford model has an actual disc bounding the
primitive circle in $\operatorname{im}(A_*^2-I)\otimes\mathbb Q$.
A selected regular fibre upstairs maps diffeomorphically to the
regular fibre downstairs: all nonidentity deck transformations move
its base root. Projecting this disc therefore bounds the same primitive
fibre circle downstairs, rather than a multiple. Appending this column
gives three unit Smith factors in each exceptional row. This proves
that all the fibre subgroup except the primitive $\gamma$ is killed,
for every allowed filling choice.

\item Over the punctured base the torus bundle has a section up to homotopy,
since the base retracts onto a graph. Its group is the corresponding
semidirect product. Modulo the vanishing normal closure just proved,
every monodromy acts trivially and its fibre group is exactly
$\mathbb Z\gamma$. A free affine log cap has a cyclic cover whose
fundamental group is its torus lattice, modulo the primitive vanishing
line in the semistable case. The covering extension has presentation
$c_i\lambda c_i^{-1}=A_i\lambda$, $c_i^{m_i}=v_i$; no further relation
is introduced. After the preceding quotient this becomes exactly
$[c_i,\gamma]=1$, $c_i^{m_i}=\gamma^{n_i}$.

\item For an original cap the translation quotient can be unrolled, and the
remaining group is $\mathbb Z\gamma$, with the base meridian trivial.
The same unrolling extends across a Mumford cap by unrolling its long
period in the infinite toric model. These are restrictions of the
original unrolled translation construction, so their global peripheral
relation has zero $\gamma$ coordinate. Equivalently, transition
functions of the line bundle contribute only the $\delta$ coordinate.
This proves $e=0$ in the canonical marking. Other meridian choices
produce exactly the stated change of $(n_i,e)$. Van Kampen now gives
\eqref{maps:eq:D.FG1}, including its completeness, and abelianization gives \eqref{maps:eq:D.FG2}.

\item Quotienting \eqref{maps:eq:D.FG1} by its central fibre yields the sphere-orbifold
presentation with cone orders $m_i>1$. With at least three such orders
it is nontrivial. One direct verification is the polygon-reflection
construction: a polygon with angles $\pi/m_i$ exists in the sphere,
plane or hyperbolic plane according to the angle sum, and products of adjacent side reflections give rotations of exact
orders $m_i$ whose cyclic product is the identity. They therefore
give a nontrivial representation of the orbifold presentation. The familiar finite spherical triangle
cases have the same nontrivial rotations. For at most two sites,
eliminate the product relation; the remaining generators commute.
The zero-, one-, and two-site integral relation matrices give \eqref{maps:eq:D.FG3}.
Reducing the last determinant modulo each multiplicity proves the
final coprimality assertions.
\end{enumerate}
\end{proof}

\begin{theorem}[both $I_3^*$ profiles are excluded from the natural sphere table]\label{maps:src:D.2.24}
\StatementRecord{maps:src:D.2.24}{proved}{Historical events: Claude Fable 5.2 (Anthropic), effort high, 16 November 2026: carried as background without fault found. GPT-6 Astra (OpenAI), ChatGPT Pro mode, effort Pro, S6\_09\_PAPERS, 17 November 2026: targeted proof check. Also checked in this writing. Pre-continuation writing events: GPT-6 Astra (OpenAI), ChatGPT Pro mode, effort Pro, S6\_10\_PAPERS, 18 November 2026: targeted proof check. Continuation metadata events: GPT-6 Astra (OpenAI), ChatGPT Pro mode, effort Pro, S6\_10\_PAPERS continuation, 19 November 2026: used and rechecked in this run; status-vocabulary reconciliation; metadata reconciliation only; no new mathematical check. \emph{Pre-continuation recorded proof scope.} GPT-6 Astra (OpenAI, ChatGPT Pro mode), effort Pro, 2026-11-17 UTC; Claude Fable 5.2 (Anthropic), effort high, 16 November 2026; also checked in this writing by GPT-6 Astra (OpenAI), ChatGPT Pro mode, effort Pro, 18 November 2026}{the hypotheses, notation and definitions stated here; \cref{maps:src:D.2.22}; \cref{maps:eq:D.IS2}; \cref{maps:eq:D.FI4}}


Neither OS profile $46$ nor $46/II$ can give a homotopy six-sphere
in the natural section-line construction and local filling sector
of \ref{maps:src:D.2.22}, for any admissible logarithmic orders or affine twists.
This includes arbitrary choices of original and semistable logarithmic
fillings at $I_3^*$.
\end{theorem}


\begin{proof}
\begin{enumerate}[label=\arabic*.,leftmargin=*]
\item \ref{maps:src:D.2.22} reduces these two profiles to the eight natural marked rows
in its finite catalogue. In every one of them, at $I_3^*$,
\[
 N=A_*^2-I,\qquad N^2=0,\qquad \operatorname{rank}N=1,
 \qquad \operatorname{Smith}(N)=(3).                         \tag{D.IS1}\label{maps:eq:D.IS1}
\]
The exact matrices also give
\[
 A_*N=-N,\qquad
 A_*=-1\text{ on }\Lambda/\ker N,\qquad
 A_*=1\text{ on }(\ker N/\operatorname{sat}\operatorname{im}N)
                         \otimes\mathbb Q.                 \tag{D.IS2}\label{maps:eq:D.IS2}
\]
These are the vanishing, dual-complex, and transverse factors of the
actual rank-one degeneration, respectively.

\item An original cap has more than one component orbit, so was already
excluded by the component intersection defect in \ref{maps:src:D.2.22}. For a log
cap of order $m=2k$, its semistable cover has monodromy
$A_*^{2k}=I+kN$. Thus its minimal rank-one Mumford cycle has exactly
$3k$ components. Its second rational homology has filtration with
component space $\mathbb Q^{3k}$, transverse area $\mathbb Q$, and
the two-dimensional product of the dual-complex circle with
$H_1(T^2;\mathbb Q)$. The finite deck group acts by permutation on
the first space, trivially on the area, and by minus one on the last
space, by \eqref{maps:eq:D.IS2}. Translation parts and unit gluing change neither
the component permutation nature nor these linear actions. Since the
group is finite, the filtration splits equivariantly over $\mathbb Q$;
rational transfer therefore gives
\[
                         b_2(U_*)=\rho+1,
\]
where $\rho$ is the number of component orbits. A group of order
$2k$ cannot act transitively on $3k$ objects: each orbit has size at
most $2k$, so $\rho\ge2$. Toric refinements can only add component
classes and do not remove this inequality. On the other hand
$g_*=\dim H_2(T^4;\mathbb Q)^{A_*}=2$, directly from the displayed
blocks. Hence $b_2(U_*)-g_*\ge1$.

\item The rank budget of \ref{maps:src:D.2.22} is saturated: its baseline is five, the
unique rank-two zero adds one, and $v_2\ge1$. Retaining an excess
$b_2(U_i)-g_i$ in that same MV calculation strengthens \eqref{maps:eq:D.FI4} to
\[
 \sum_i\operatorname{rank}(A_i-I)
 +\sum_i\bigl(b_2(U_i)-g_i\bigr)\le7-v_2.
\]
The left side is at least seven, while the right side is at most
six. This contradiction excludes a sphere. All comparisons took
place on the actual semistable cover and its free finite quotient;
no bound on $k$ or on the translation vector was used.
\end{enumerate}
\end{proof}


\begin{theorem}[all initial rows and all regular logarithms in the semistable sector]\label{maps:src:D.2.25}
\StatementRecord{maps:src:D.2.25}{proved}{Historical events: Claude Fable 5.2 (Anthropic), effort high, 16 November 2026: carried as background without fault found. GPT-6 Astra (OpenAI), ChatGPT Pro mode, effort Pro, S6\_09\_PAPERS, 17 November 2026: targeted proof check. Also checked in this writing. Pre-continuation writing events: GPT-6 Astra (OpenAI), ChatGPT Pro mode, effort Pro, S6\_10\_PAPERS, 18 November 2026: targeted proof check. Continuation metadata events: GPT-6 Astra (OpenAI), ChatGPT Pro mode, effort Pro, S6\_10\_PAPERS continuation, 19 November 2026: used and rechecked in this run; status-vocabulary reconciliation; metadata reconciliation only; no new mathematical check. \emph{Pre-continuation recorded proof scope.} GPT-6 Astra (OpenAI, ChatGPT Pro mode), effort Pro, 2026-11-17 UTC; Claude Fable 5.2 (Anthropic), effort high, 16 November 2026; also checked in this writing by GPT-6 Astra (OpenAI), ChatGPT Pro mode, effort Pro, 18 November 2026}{the hypotheses, notation and definitions stated here; \cref{maps:src:D.2.22}; \cref{maps:src:D.2.23}; \cref{maps:eq:D.FG2}; \cref{maps:eq:D.FG1}; \cref{maps:eq:D.FG3}; \cref{maps:src:D.2.14}; \cref{maps:src:D.2.12}; \cref{maps:src:D.2.13}; \cite[Theorem A]{manGR25}}


Consider any of the $62$ natural initial rows of \ref{maps:src:D.2.22} on the four
all-multiplicative profiles $45,47,55,56$. Write $P=aG,Q=bG$, with
$a,b>0$. Use the primitive Mumford model at each multiplicative fibre.
An original zero-free filling with zero line multidegree and one
component orbit may be used in its place; it has the same marked
primitive rank-one model. Perform any finite number of regular
logarithmic modifications, with data
\[
 m_i>0,\qquad v_i=(x_i,y_i,z_i,n_i)\in\mathbb Z^4,
 \qquad\gcd(m_i,x_i,y_i,z_i,n_i)=1.                           \tag{D.SC1}\label{maps:eq:D.SC1}
\]
The coordinates are $(e_1,e_2,\delta,\gamma)$. Put
$\omega=a e_1^*e_2^*+b\gamma^*\delta^*$. The initial rows have
$\gcd(a,b)=1$. Retain the peripheral integer $e$ of \ref{maps:src:D.2.23};
one may keep every order-one site and use the canonical value $e=0$.

Here is a complete integral homology algorithm without a bound on
orders, number of sites, or translation coordinates. Let $S$ be the
Seifert relation matrix \eqref{maps:eq:D.FG2}, and write
\[
 \operatorname{coker}S=\mathbb Z^r\oplus A_1,\qquad r\in\{0,1\}.
\]
Form an integral presentation $\mathcal C$ with generators
$s$ and $t_i(u)$, linear in $u\in\Lambda=\mathbb Z^4$, and relations
\[
 \begin{split}
 m_i t_i(u)&=\omega(u,v_i)s\quad(u=e_1,e_2,\delta,\gamma),\\
 t_i(v_i)&=0,\\
 \sum_i t_i(\delta)+be\,s&=0.
 \end{split}                                                \tag{D.SC2}\label{maps:eq:D.SC2}
\]
Write $T=\operatorname{Tor}\operatorname{coker}\mathcal C$.
Then $\operatorname{coker}\mathcal C$ has free rank $r$ and
\[
\begin{array}{c|c}
 j&H_j(X;\mathbb Z)\\\hline
 0,6&\mathbb Z\\
 1&\mathbb Z^r\oplus A_1\\
 2,3&\mathbb Z^{2r}\oplus T\\
 4&\mathbb Z^{2r}\oplus A_1\\
 5&\mathbb Z^r.
\end{array}                                                 \tag{D.SC3}\label{maps:eq:D.SC3}
\]
The fundamental group is exactly \eqref{maps:eq:D.FG1}, rather than merely the
abelian group in degree one. Smith reduction of the explicit matrix
\eqref{maps:eq:D.SC2} computes $T$, including all extensions.

If this group is trivial, let
\[
 M=\prod_i m_i,\quad
 X_0=\sum_i(M/m_i)x_i,\quad Y_0=\sum_i(M/m_i)y_i,\quad
 g=\gcd(b,M,X_0,Y_0).
\]
For no displayed sites use $M=1$ and $X_0=Y_0=0$.
Then
\[
 H_2(X)=H_3(X)=
 \mathbb Z/g\oplus\mathbb Z/M\oplus\mathbb Z/(bM/g).          \tag{D.SC4}\label{maps:eq:D.SC4}
\]
Consequently a sphere occurs precisely when the group criterion
\eqref{maps:eq:D.FG3} holds and $b=M=1$. In the canonical marking with all sites
retained, this says that every regular logarithmic order is one and
$|\sum_i n_i|=1$.

There are exactly $17$ permitted initial rows with $b=1$:
\[
\begin{array}{c|c|c|c}
 \text{profile}&\text{zero fibre}&(a,b)&\text{marked positions}\\\hline
 45&I_1&(8,1)&4\\
 47&I_2&(7,1)&1\\
 47&I_1&(14,1)&3\\
 55&I_4&(5,1)&1\\
 55&I_1&(20,1)&3\\
 56&I_5&(6,1)&1\\
 56&I_3&(10,1)&1\\
 56&I_2&(15,1)&1\\
 56&I_1&(30,1)&2.
\end{array}                                                 \tag{D.SC5}\label{maps:eq:D.SC5}
\]
This is the complete natural sphere table for these four profiles,
including original-versus-Mumford choices: any additional original
cap not covered by the first paragraph is excluded by the positive
local defect in \ref{maps:src:D.2.22}.

More generally, every simply connected output in these $17$ rows,
with arbitrary data \eqref{maps:eq:D.SC1}, is oriented diffeomorphic to
\[
                              \widetilde\Sigma_0N_M,         \tag{D.SC6}\label{maps:eq:D.SC6}
\]
with the twisted framing and the Smale five-manifold of \ref{maps:src:D.2.14}.
Thus its full smooth type depends only on $M$, including when $M$
is even. For $b>1$, \eqref{maps:eq:D.SC2}--\eqref{maps:eq:D.SC4} give the complete integral
homology and exclude a sphere; no smooth classification of all
those torsion types is asserted here.

\end{theorem}


\begin{proof}
\begin{enumerate}[label=\arabic*.,leftmargin=*]
\item \emph{1. The actual marked caps and the constant integral calculation.}
At the component zero the valuation lattice has index one by \ref{maps:src:D.2.22},
so its cap is the primitive hexagonal model. Every other
multiplicative monodromy in these $62$ rows has rank-one difference
and Smith width one, as checked in \ref{maps:src:D.2.23}. Its local Mumford cap is
the marked one-node model times $T^2$. The integral specialization
onto its $H_2$, and that of the hexagonal cap, are onto; their
boundary kernels are exactly the normal tori swept by
$\operatorname{Fix}(A_i)$. The toric angular coordinates give these
as pure swept classes, with zero fibre-area term. The same statement
for an original zero-degree, one-orbit cap follows by quotienting
its component cycle by $Q$, using the unit-factor isotopy of \ref{maps:src:D.2.12}.

\item These normal markings agree with the section-line construction.
Rigidify $M$ along $O$ on the punctured base. Put the divisor of this
rigidification at the chosen zero; its integral change there is
multiplication of the line coordinate by $t^h$. On the infinite
toric model this translates the fan by an integral lattice vector,
and hence extends over the cap. At all other points the
rigidification is a unit and the earlier relative unit isotopy
applies. Thus this normalization does not assume that equality of
vertical forms alone determines mixed swept classes.

\item The exact program \texttt{manifolds\_semistable\_catalogue.py}
performs the graph and exterior calculations for all $62$ rows.
The degree-two fibre coinvariants are freely generated by $s$,
with primitive detector $\omega$; all its Smith factors are one.
Adjoin one regular puncture, with monodromy the identity. In the
notation of \ref{maps:src:D.2.13}, the $14$ old normal columns form a primitive
rank-$14$ sublattice of the rank-$17$ graph kernel. Their quotient
is freely generated by the regular sweeps of $e_1,e_2,\gamma$,
while the sweep of $\delta$ is zero. All $62$ left inverses and
quotient matrices are saved, so this is an exact integral
certificate, not a numerical rank inference.

\item Adding another regular puncture adds four free graph coordinates.
The old normal lattice remains primitive by projection onto the
previous graph. Since every monodromy fixes $\delta$, the peripheral
Fox maps satisfy $\sum j_i(\delta)=0$. Induction therefore gives
all regular sweeps freely, with the sole relation
$\sum_i t_i(\delta)=0$. The degree-two exterior summand remains
$\mathbb Zs$. A change of global peripheral lift by $e\gamma$
adds $be\,s$ to this relation: the correction to a swept $\delta$
cycle is its wedge with that translation, and
$\omega(\delta,\gamma)=-b$. Equivalently, under
$v_i\mapsto v_i+m_i h_i\gamma$ the new sweep is
$t_i'(u)=t_i(u)+\omega(u,h_i\gamma)s$;
$e\mapsto e+\sum h_i$ makes \eqref{maps:eq:D.SC2} invariant. This checks its
sign against \eqref{maps:eq:D.FG1}.

\item \emph{2. The complete regular-log kernel.}
At a regular log point its reduced torus has lattice
$\Lambda'=\Lambda+\mathbb Z(v_i/m_i)$. The nearby fibre inclusion
and the base circle send a boundary bivector respectively to
$z$ and $u\wedge(v_i/m_i)$ in $\bigwedge^2\Lambda'$.
This map is onto. The condition in \eqref{maps:eq:D.SC1} means $v_i$ is primitive
modulo $m_i$; the integral normal lattice is exactly
$m_i\Lambda+\mathbb Z v_i$. This can be checked in a basis in which
$v_i$ modulo $m_i$ is the first basis vector, using the Euclidean
algorithm. Its generators give the full normal columns
\[
 (-\omega(u,v_i),\ m_it_i(u)),\qquad (0,t_i(v_i)).
\]
They are precisely the first two lines of \eqref{maps:eq:D.SC2}. The argument
identifies the full integral kernel, including its index, not only
its rational span. Since specialization is onto every cap in degree
two, Mayer--Vietoris reduces to this cokernel.

\item The remaining first MV kernel is free: all boundary $H_1$ groups
are free in these primitive toric models. Its rank is $r$, because
the domain and target of the first boundary map have equal rational
dimension and its cokernel is $H_1(X)$. Hence it contributes a
split $\mathbb Z^r$ to $H_2(X)$. Over $\mathbb Q$, eliminate the
sweeps in \eqref{maps:eq:D.SC2}; its sole residual relation is
\[
                 b\left(e-\sum_i n_i/m_i\right)s=0.
\]
Thus its free rank is the same $r$ as the Seifert matrix.
Only the primitive hexagonal cap contributes Euler characteristic,
so $\chi(X)=2$. Integral Poincare duality and the universal coefficient
theorem give \eqref{maps:eq:D.SC3}, including the identical torsion in $H_2,H_3$.
This also proves the formula when the rational Euler number is zero;
that case is not obtained by deleting a nonzero-order torsion factor.

\item For one site in canonical marking, quotienting the zero $\delta$
sweep displays the concrete four-by-five matrix
\[
 \begin{pmatrix}
 -ay&ax&bn&-bz&0\\
 m&0&0&0&x\\0&m&0&0&y\\0&0&0&m&n
 \end{pmatrix}.                                             \tag{D.SC7}\label{maps:eq:D.SC7}
\]
Its maximal minors are, up to signs,
$bm^2n\,(m,n,z,y,x)$. In particular, whenever $n\ne0$,
the finite degree-two cokernel has order $bm^2|n|$.
The implementation accepts the general multi-site matrix \eqref{maps:eq:D.SC2},
and verifies its comparison with all constant graph presentations.

\item \emph{3. The simply connected Smith form.}
Assume \eqref{maps:eq:D.FG1} is trivial. After removing order-one sites there are
at most two multiplicities greater than one, they are coprime, and
every $n_i$ is a unit modulo its $m_i$. Put
\[
                  D=M\left(\sum_i n_i/m_i-e\right).
\]
The group criterion gives $D=\pm1$. Work over $\mathbb Z_p$.
If $p$ does not divide $M$, all sweeps can be eliminated and the
cokernel is $\mathbb Z_p/(b)$. If $p\mid M$, it divides exactly
one multiplicity, say $m_i$. For every other site $m_j$ is a unit,
so its $\delta$ sweep equals $-bn_j/m_j$ times $s$. The global
relation then eliminates $t_i(\delta)$. Its own normal relation
becomes $bD/(M/m_i)$ times $s$, and hence is equivalent to $bs=0$.
Since $n_i$ is a unit, its relation $t_i(v_i)=0$ eliminates
$t_i(\gamma)$, whose remaining normal relation is redundant by
alternation of $\omega$. The three remaining generators have
presentation
\[
 \begin{pmatrix}
 b&0&0\\ -ay_i&m_i&0\\ ax_i&0&m_i
 \end{pmatrix}.                                             \tag{D.SC8}\label{maps:eq:D.SC8}
\]
Its determinantal divisors are $g_i,m_i g_i,bm_i^2$, where
$g_i=\gcd(b,m_i,x_i,y_i)$ in $\mathbb Z_p$; we used
$\gcd(a,b)=1$. The Chinese remainder theorem packages all these
primary invariant factors as \eqref{maps:eq:D.SC4}, because
$X_0,Y_0$ reduce at $p\mid m_i$ to unit multiples of $x_i,y_i$.
The invariant factors divide in the stated order, since $g$ divides
both $b$ and $M$. This proves the formula for every integer twist,
not by extrapolating from the finite checks.

\item The group in \eqref{maps:eq:D.SC4} vanishes exactly for $b=M=1$. Simple connectivity,
Euler number two, and integral duality then give a homotopy sphere;
$\Theta_6=0$ identifies it with the standard $S^6$. The seventeen
initial rows in \eqref{maps:eq:D.SC5} are the complete $b=1$ filter of \ref{maps:src:D.2.22}.
An excluded original cap has nonzero Chern degree or multiple
component orbits and would already force $b_2>0$, so cannot supply
additional sphere rows.

\item \emph{4. The smooth type when $b=1$.}
We verify the hypotheses of the geometric proof \ref{maps:src:D.2.14} rather than
inferring a smooth type from homology. At every zero-free
multiplicative point the component equation forces $P$ to be narrow
when $Q=G$; the component action of $Q$ is primitive and transitive. The code
checks these component statements and, more directly, that adjoining
$\delta$ to the actual primitive vanishing vector gives two unit
Smith factors at every such site. Thus $\delta$ stays primitive
modulo the vanishing line and its circle is free on the cap; width
one alone is not being used to infer this. At the zero, index one makes the valuation lattice the full original
short lattice, in which $\delta$ is the second coordinate. The code
checks its primitive inclusion in that lattice in all seventeen rows.
The actual hexagonal fan weights are therefore the same $(-1,0,1)$ as in \ref{maps:src:D.2.14}, with one fixed rational curve and
normal bundle $\mathcal O(-1)_{+1}\oplus\mathcal O(-1)_{-1}$.
For every regular log cap, $\gcd(m_i,n_i)=1$ implies that a scalar
angle lying in $\Lambda+\mathbb Z(v_i/m_i)$ is integral. Hence no
new circle stabilizers occur.

\item The fixed curve lies in the cyclic fibre summand generated by $s$.
The $\delta$ normal relations together with the global relation,
when $b=1$, are the unimodular Seifert matrix because $D=\pm1$.
They kill $s$ integrally. The fixed sphere is therefore nullhomologous
and nullhomotopic. Its equivariant surgery gives
$Y\cong X\#(S^3\times S^3)$ with a free, projectably spin circle
action, exactly as in \ref{maps:src:D.2.14}. The same orientation convention in its
normal-coordinate conjugation is retained. Its spin simply connected
quotient has $H_2=\mathbb Z\oplus(\mathbb Z/M)^2$, with primitive
Euler class. Smale--Barden and the spin branch of
\cite[Theorem A]{manGR25} give
$Y\cong(S^3\times S^3)\#\widetilde\Sigma_0N_M$.
Cancellation by the full Zhubr invariants, as justified in \ref{maps:src:D.2.14},
proves \eqref{maps:eq:D.SC6}, including the two-primary refined invariant.
\end{enumerate}
\end{proof}


\begin{theorem}[complete natural sphere table without $I_n^*$]\label{maps:src:D.2.26}
\StatementRecord{maps:src:D.2.26}{proved}{Historical events: Claude Fable 5.2 (Anthropic), effort high, 16 November 2026: carried as background without fault found. GPT-6 Astra (OpenAI), ChatGPT Pro mode, effort Pro, S6\_09\_PAPERS, 17 November 2026: targeted proof check. Also checked in this writing. Pre-continuation writing events: GPT-6 Astra (OpenAI), ChatGPT Pro mode, effort Pro, S6\_10\_PAPERS, 18 November 2026: targeted proof check. Continuation metadata events: GPT-6 Astra (OpenAI), ChatGPT Pro mode, effort Pro, S6\_10\_PAPERS continuation, 19 November 2026: used and rechecked in this run; status-vocabulary reconciliation; metadata reconciliation only; no new mathematical check. \emph{Pre-continuation recorded proof scope.} GPT-6 Astra (OpenAI, ChatGPT Pro mode), effort Pro, 2026-11-17 UTC; Claude Fable 5.2 (Anthropic), effort high, 16 November 2026; also checked in this writing by GPT-6 Astra (OpenAI), ChatGPT Pro mode, effort Pro, 18 November 2026}{the hypotheses, notation and definitions stated here; \cref{maps:src:D.2.22}; \cref{maps:eq:D.SC1}; \cref{maps:eq:D.SC3}; \cref{maps:eq:D.FG3}; \cref{maps:src:D.2.25}; \cref{maps:src:D.2.23}; \cite[Section 2.4, equation (2.2)]{manUU21}; \cref{maps:src:D.2.12}; \cref{maps:eq:D.HY6}; \cref{maps:eq:D.SC8}; \cref{maps:src:D.2.14}}


For the sixteen profiles in \ref{maps:src:D.2.22} containing no fibre $I_n^*$,
the natural section-line construction has exactly $59$ marked
initial/additive-filling patterns that can yield a sphere.
They lie on $50$ distinct initial rows. The complete table,
including the section pair, zero component, marked fibre positions,
and every additive original/log choice, is
\path{records/manifolds_nonstar_sphere_patterns.csv}.
At the zero the index-one Mumford filling is required; at every
other multiplicative fibre in these sphere patterns either the
original filling or the primitive Mumford filling is permitted.
The pattern counts are
\[
\begin{array}{c|r@{\qquad}c|r@{\qquad}c|r@{\qquad}c|r}
43&3&43/II&1&45&4&45/II&4\\
47&4&47/II&4&47/III&4&49&3\\
49/II&1&49/III&2&55&4&55/II&4\\
56&5&56/II&6&56/III&5&56/IV&5.
\end{array}                                                 \tag{D.HY1}\label{maps:eq:D.HY1}
\]
These are parameter families, not multiplicity samples. The exact
unrestricted parameter conditions for every row follow below.

More generally, take a natural initial row without $I_n^*$ and
retain only original caps of zero line multidegree with one
component orbit, using the primitive Mumford caps elsewhere.
A finite good-reduction log cap can occur in a simply connected
output only when its integral fixed lattice contains a vector
$\ell_i$ with $\gamma^*(\ell_i)=1$. Write its power vector
$v_i=n_i\ell_i+r_i\delta$. The following calculation applies
whenever $\gcd(m_i,n_i)=1$, a necessary condition for simple
connectivity. Set
\[
\begin{array}{c|ccccc}
 \text{finite type}&II&III&IV&III^*&IV^*\\\hline
 d_i\text{ (monodromy order)}&6&4&3&4&3\\
 q_i\text{ (vertical index)}&1&2&3&2&3.
\end{array}                                                 \tag{D.HY2}\label{maps:eq:D.HY2}
\]
Allow every order $m_i=d_i k_i$, $k_i>0$, every such power vector,
and every free affine lift with that power. Additional regular logs
have the unrestricted data \eqref{maps:eq:D.SC1}.

Here is their full integral degree-two presentation. For each finite
log retain a generator $d_i$; for each regular log retain its four
sweeps $t_j(\Lambda)$; retain the fibre generator $s$. Impose
\[
\begin{aligned}
 (m_i/q_i)d_i+(bn_i/q_i)s&=0 &&\text{(finite logs)},\\
 m_jt_j(u)-\omega(u,v_j)s&=0,\quad t_j(v_j)=0
                                      &&\text{(regular logs)},\\
 \sum_i d_i+\sum_jt_j(\delta)+be\,s&=0. &
\end{aligned}                                                 \tag{D.HY3}\label{maps:eq:D.HY3}
\]
All entries are integers. With $T$ the torsion of this cokernel and
$r,A_1$ obtained from the complete Seifert matrix, the full homology
is exactly \eqref{maps:eq:D.SC3}. Thus this is also an all-order integral algorithm
for the nonsimply-connected members satisfying the stated finite
coprimality condition. No conclusion for the omitted noncoprime
finite twists is silently inferred.

For a simply connected output put
\[
 Q_f=\prod_{i\text{ finite}}q_i,\quad c=b/Q_f,\quad
 R=\prod_{j\text{ regular}}m_j,
 \quad X_0=\sum_j(R/m_j)x_j,\quad Y_0=\sum_j(R/m_j)y_j.
\]
Empty products equal one. Then $c$ is a positive integer and
\[
 H_2(X)=H_3(X)=
 \mathbb Z/g\oplus\mathbb Z/R\oplus\mathbb Z/(cR/g),
 \qquad g=\gcd(c,R,X_0,Y_0).                                \tag{D.HY4}\label{maps:eq:D.HY4}
\]
Consequently the complete sphere criterion is
\[
 \pi_1(X)=1\text{ by \eqref{maps:eq:D.FG3}},\qquad b=Q_f,
 \qquad m_j=1\text{ at every regular log}.                  \tag{D.HY5}\label{maps:eq:D.HY5}
\]
For one finite log this includes every multiple of its order; for
two finite logs it requires their multiplicities to be coprime and
the displayed Seifert determinant to be $\pm1$. The only two-finite
sphere pattern is the original $IV^*,III,I_1$ pair. The finite
numerators $r_i$ are unrestricted. All twelve profiles in \eqref{maps:eq:D.HY1}
which contain additive fibres therefore admit infinitely many
allowable sphere multiplicities; the four all-multiplicative
profiles have the order-one conclusion of \ref{maps:src:D.2.25}.

Whenever $c=1$ and the output is simply connected, its exact oriented
smooth type is $\widetilde\Sigma_0N_R$. For $c>1$ the groups \eqref{maps:eq:D.HY4}
are exact, but a general smooth identification of these remaining
torsion types is not asserted by this theorem.

\end{theorem}


\begin{proof}
\begin{enumerate}[label=\arabic*.,leftmargin=*]
\item \emph{1. Necessary local choices and the integral graph lattice.}
The sphere reduction \ref{maps:src:D.2.22} already excludes every nonzero-degree
original cap and every original cap with multiple component orbits.
At a finite good-reduction point the image of $\gamma^*$ on the
integral fixed lattice is the component order of $Q$, namely one,
two or three in these profiles. If it is greater than one, that
prime divides both $m_i$ and every power-vector coordinate $n_i$;
\ref{maps:src:D.2.23} then excludes simple connectivity. These are exact fixed-lattice
calculations, included for all $172$ initial rows in
\texttt{manifolds\_sphere\_normal\_lattices.py}.

\item That program also computes the complete boundary fixed-sweep map
\[
 \bigoplus_i\operatorname{Fix}(A_i)
       \longrightarrow H_1(B^\circ;H_1(T^4;\mathbb Z)).
\]
For every row it is onto with unit Smith factors. Its kernel has
rank one and is precisely the primitive diagonal copy of $\delta$,
by the peripheral Fox identity. Thus the graph part is the direct
sum of these fixed lattices with the sole relation
$\sum t_i(\delta)=0$. Adding regular sites has the same integral
inductive comparison as \ref{maps:src:D.2.25}. The fibre coinvariants are
$\mathbb Zs$, detected primitively by $\omega$.

\item There is an integral check for the permitted original additive
caps, beyond their rational invariant cycles. A zero-degree original
cap is topologically $S^1_\delta$ times the mapping torus of the
component translation of its Kodaira tree. Here zero multidegree
trivializes the underlying line on the tree of rational curves.
In such a trivialization the remaining equivariant multiplier is a
unit. Its circle-valued part has a logarithm because the tree has
$H^1=0$; interpolation of that logarithm gives an isotopy to the
identity. The rigidification along the zero section fixes the
constant phase. This isotopy extends over the normal collar, and
its unit multiplier has zero meridian winding. Consequently it
preserves the marked normal tori used below and proves the stated
product description, rather than merely a rational homology
comparison. Its homology is free;
with one orbit, $H_2$ has rank two and the component intersection
map to $H_2(U,\partial U)$ is zero. Poincare--Lefschetz duality
therefore makes $H_2(\partial U)\to H_2(U)$ onto integrally.
The $\delta$ normal torus bounds the section disc times $S^1_\delta$.
For the other fixed loop, take the least power of $Q$ which is
narrow and its local Neron logarithm. The resulting family of loops
extends over the disc, so its normal torus also bounds. Local
logarithms at additive reduction are described in
\cite[Section 2.4, equation (2.2)]{manUU21}.
The linearization is a unit and the line has zero component degrees,
so the rigidified norm lift introduces no extra fibre term.
These two tori project onto the entire integral fixed lattice.
They are therefore the full free part of the kernel; any boundary
torsion maps to zero in both free targets and changes nothing in
MV. This proves that each unchanged minimal cap simply sets its
fixed sweeps to zero, with the asserted markings.

\item \emph{2. The actual finite log kernel.}
Let $\Lambda_0=\ker\gamma^*$. At every finite slot with primitive
$\gamma$, the exact vertical coinvariant calculation is
\[
 (\Lambda_0)_{A_i}=\mathbb Z\beta_i,
 \qquad [\delta]=q_i[\beta_i],
 \qquad \operatorname{Fix}(A_i)=\mathbb Z\delta\oplus\mathbb Z\ell_i.
                                                               \tag{D.HY6}\label{maps:eq:D.HY6}
\]
The table \eqref{maps:eq:D.HY2} follows either from the component group or the
saved primitive Smith bases. The code verifies all $68$ such slots.
For $m=m_i$, its norm satisfies
$N_m\beta_i=(m/q_i)\delta$ and $N_m\ell_i=m\ell_i$.
Invariance of $\omega$ then proves the exact, signed identity
\[
                     -\omega(\beta_i,\ell_i)=b/q_i.
\]
In particular $q_i\mid b$.

\item An arbitrary affine good-reduction lift with the prescribed integral
power is translation-conjugate, over the root disc, to the invariant
representative. To see the integral issue, the usual sequence of
period lattices and holomorphic vector sections identifies the
translation-conjugacy obstruction with
$\Lambda^{A_i}/N_m\Lambda$; it is exactly the power vector modulo
norms. The vector-section cohomology vanishes by finite averaging.
Thus equality of the power class gives an actual holomorphic
translation conjugacy extending over the smooth central torus.
This justifies the invariant representative for every affine lift
under discussion, rather than restricting the theorem to a chosen
constant lift.

\item The reduced finite fibre has first homology generated by
$\beta_i,\ell_i,c_i$ with
$m c_i=n_i\ell_i+r_iq_i\beta_i$. Choose
$\alpha m+\beta n_i=1$. Its first homology is free, with basis
$\beta_i$ and $\alpha\ell_i+\beta c_i$. Its flat normal complex
line has torsion Chern class, hence zero integral Chern class because
$H^2$ of this fibre is free. Consequently the boundary map in degree
two is onto and its kernel is exactly the two normal circle tori
of this basis. The actual affine collar maps give their images in
the exterior as
\[
\begin{split}
 &\bigl(-\omega(\beta_i,v_i),\ t_i(N_m\beta_i)\bigr),\\
 &\bigl(-\alpha\omega(\ell_i,v_i),\
               t_i(\alpha N_m\ell_i+\beta v_i)\bigr).
\end{split}                                                  \tag{D.HY7}\label{maps:eq:D.HY7}
\]
This is the marked rectangle-sweep computation of \ref{maps:src:D.2.12}; the basis
calculation here verifies that these are the full integral kernel,
not a sublattice of finite index. Substituting \eqref{maps:eq:D.HY6} gives
\[
 (bn_i/q_i)s+(m_i/q_i)t_i(\delta),\qquad
 -\alpha b r_i s+t_i(\ell_i)+\beta r_i t_i(\delta).
\]
The second relation eliminates $t_i(\ell_i)$ with coefficient one.
The first is exactly the finite relation in \eqref{maps:eq:D.HY3}. In particular
all $r_i$ disappear from the degree-two group for an integral reason.
The regular kernel was proved in \ref{maps:src:D.2.25}. Together with the one
global $\delta$ relation this proves \eqref{maps:eq:D.HY3}. All cap $H_1$ groups
in this coprime range are free, so the remaining MV kernel is the
split free group of rank $r$. Rational elimination and integral
duality give the whole table \eqref{maps:eq:D.SC3}.

\item \emph{3. All-primary Smith calculation and sphere enumeration.}
Suppose the Seifert group is trivial. Its nontrivial multiplicities
are pairwise coprime, and all numerators are units modulo them.
Thus the $q_i$ are pairwise coprime and $Q_f\mid b$.
Write $M$ for the product of all multiplicities, including the
finite ones, and $D=M(\sum n_i/m_i-e)=\pm1$.
Eliminate the $\delta$ sweeps over $\mathbb Z_p$ exactly as in
\ref{maps:src:D.2.25}. At a prime dividing a finite multiplicity $m_i$, its row
has been divided by $q_i$, so the residual fibre relation is
$(b/q_i)D/(M/m_i)$ times $s$. All other multiplicities are units
at that prime. At a prime dividing a regular multiplicity, the
residual relation is $b$ times a unit; the finite $q_i$ are units
there. At a prime dividing no multiplicity, it is again $b$ times
a unit. These are precisely the primary fibre orders of
$c=b/Q_f$.

\item The only extra generators at a prime dividing a regular multiplicity
are its two elliptic sweeps. Their three-generator matrix is
\eqref{maps:eq:D.SC8}, with the fibre order replaced by $c$. Because
$\gcd(a,c)=1$, its invariant factors are
$g_i,m_i,cm_i/g_i$. Chinese remainders yield \eqref{maps:eq:D.HY4}, with the
stated $X_0,Y_0$. This proves \eqref{maps:eq:D.HY5} for all integers.

\item The finite program now has no arbitrary search bound: list every
natural row, every additive original/log choice satisfying the
local conditions, and the vertical indices \eqref{maps:eq:D.HY2}. There are
$164$ minimal-cap patterns in this coprime finite range. Two have
an unavoidable orbifold obstruction from noncoprime base orders;
$162$ can satisfy simple connectivity. Exactly $59$ have $b=Q_f$.
Their complete data are saved in the CSV and JSON tables. Each
of them is realized: choose a numerator $\pm1$ for one finite
site, Bezout numerators for the coprime $(3k,4l)$ pair, or total
order-one regular numerator $\pm1$ if there is no finite site.
The free affine quotient exists by its primitive $\gamma$ value.
This also gives infinitely many orders for every listed pattern
with at least one finite site. Conversely every sphere was forced
into these choices by \ref{maps:src:D.2.22} and \ref{maps:src:D.2.23}, proving completeness in
all sixteen profiles.

\item \emph{4. The exact smooth type for $c=1$.}
The catalogue directly checks that at every zero-free multiplicative
site of a $b=Q_f$ pattern, $P$ is narrow, $Q$ acts transitively on
components, and adjoining $\delta$ to the actual vanishing vector
has two unit Smith factors. Thus the scalar circle is free there.
It is free on every permitted original additive cap by its trivial
component Chern class, and on each finite log cap because
$\gcd(m_i,n_i)=1$. At the zero its lattice is the full original
short lattice, retaining $\delta$ as the second coordinate, so its
single fixed sphere and normal weights are those of \ref{maps:src:D.2.14}.

\item The square subsystem on $s$ and the $\delta$ sweeps in \eqref{maps:eq:D.HY3}
has determinant $bD/Q_f=cD=\pm1$. Hence $s$ is zero integrally.
The fixed sphere is nullhomotopic, and the equivariant surgery,
projectable spin quotient, Smale--Barden classification, twisted
suspension formula and full Zhubr cancellation from \ref{maps:src:D.2.14} apply
with $H_2=(\mathbb Z/R)^2$. They give precisely
$X\cong\widetilde\Sigma_0N_R$. This supplies the asserted smooth
identification and all its refined spin data.
\end{enumerate}
\end{proof}


\begin{theorem}[the three $I_1^*$ profiles give no natural sphere]\label{maps:src:D.2.27}
\StatementRecord{maps:src:D.2.27}{proved}{Historical events: Claude Fable 5.2 (Anthropic), effort high, 16 November 2026: carried as background without fault found. GPT-6 Astra (OpenAI), ChatGPT Pro mode, effort Pro, S6\_09\_PAPERS, 17 November 2026: targeted proof check. Also checked in this writing. Pre-continuation writing events: GPT-6 Astra (OpenAI), ChatGPT Pro mode, effort Pro, S6\_10\_PAPERS, 18 November 2026: targeted proof check. Continuation metadata events: GPT-6 Astra (OpenAI), ChatGPT Pro mode, effort Pro, S6\_10\_PAPERS continuation, 19 November 2026: used and rechecked in this run; status-vocabulary reconciliation; metadata reconciliation only; no new mathematical check. \emph{Pre-continuation recorded proof scope.} GPT-6 Astra (OpenAI, ChatGPT Pro mode), effort Pro, 2026-11-17 UTC; GPT-6 Astra (OpenAI, ChatGPT Pro mode), S6\_06\_BLIND, effort Pro, 2026-11-17; Claude Fable 5.2 (Anthropic), effort high, 16 November 2026; also checked in this writing by GPT-6 Astra (OpenAI), ChatGPT Pro mode, effort Pro, 18 November 2026}{the hypotheses, notation and definitions stated here; \cref{maps:src:D.2.22}; \cref{maps:src:D.2.24}; \cite[Sections 7.6--7.7]{manSS10}; \cref{maps:src:D.2.23}; \cref{maps:eq:D.ST3}; \cref{maps:eq:D.ST4}; \cref{maps:eq:D.ST8}; \cref{maps:eq:D.HY3}; \cref{maps:eq:D.ST10}; \cref{maps:src:D.2.26}; \cref{maps:eq:D.ST11}; \cref{maps:eq:D.ST2}; \cref{maps:eq:D.ST12}}


None of the natural section-line inputs on
\[
 (I_1^*,I_3,I_1,I_1),\qquad (I_1^*,II,I_3),\qquad
 (I_1^*,IV,I_1)
\]
produces a homotopy six-sphere, for any permitted original,
logarithmic or Mumford filling, multiplicity or affine twist in
the sector of \ref{maps:src:D.2.22}. The obstruction is integral. In the only
remaining minimal star caps, the full normal-kernel indices are
\[
\begin{array}{c|c|c|c}
 m_*&\text{component action}&H_2(U_*;\mathbb Z)
       &\kappa_*\ \text{(normal projection index)}\\\hline
 2&\text{one component}&\mathbb Z^2&1\\
 4&\text{two components interchanged}&
                   \mathbb Z^2\oplus\mathbb Z/2&4.
\end{array}                                                   \tag{D.ST1}\label{maps:eq:D.ST1}
\]
If the complete output is simply connected and its other caps
have the minimal rational ranks, the order of $H_2$ has exactly
one factor of $2$ in the first row, and exactly two factors of $2$
in the second row. In particular $H_2\ne0$.

\end{theorem}


\begin{proof}
\begin{enumerate}[label=\arabic*.,leftmargin=*]
\item \emph{1. Exhausting the possible star caps.}
The sixteen natural initial rows on these three profiles are
listed in the certificate
\path{records/manifolds_i1star_obstruction.json}.
No original star cap has both zero component multidegree and one
component orbit, so \ref{maps:src:D.2.22} excludes every such original cap from
a sphere. A semistable root cover has order $m_*=2k$, and
$A_*^{2k}-1=kN$, where $N$ has rank one, square zero and primitive
image. Thus its minimal rank-one fibre has exactly $k$ components.
This is the integral version of the $I_n^*$ semistable reduction
used in \ref{maps:src:D.2.24}; see also \cite[Sections 7.6--7.7]{manSS10}.

\item On the dual-complex circle the residual action reverses
orientation. Its action on the $k$ components therefore has the
form $i\mapsto r-i$ and has at least $\lceil k/2\rceil$ orbits.
The transverse two-torus has trivial linear action after the
vanishing circle is collapsed. The component classes contribute
one class per orbit, the transverse area contributes one, and
the dual-circle cross terms are anti-invariant. Consequently
\[
 b_2(U_*)=\rho+1,\qquad
 \dim H_2(T^4;\mathbb Q)_{A_*}=2.
\]
The saturated rank budget of \ref{maps:src:D.2.22} forces $\rho=1$.
Hence $k=1$, or $k=2$ with the two components interchanged.
Further toric subdivisions add a component defect and cannot
restore this equality. This proves that \eqref{maps:eq:D.ST1} exhausts the
remaining cap types without an order bound imposed externally.

\item There is a second necessary integral condition. On the quotient
of the fibre lattice by the primitive vanishing circle, an
actual affine power vector is fixed by the induced involution.
In eight of the sixteen rows the image of this fixed lattice
under $\gamma^*$ is $2\mathbb Z$. These are precisely the ordered
section pairs
\[
 (a,b)=(4,1),\ (4,3),\ (12,1).
\]
Their star numerator is even, while their star multiplicity is
even. The Seifert presentation \ref{maps:src:D.2.23} therefore precludes simple
connectivity. This uses the fixed lattice \emph{after} collapsing
the vanishing circle, not the smaller fixed lattice before that
collapse.

\item For the other eight rows the pairs are
\[
 (a,b)=(1,4),\ (1,12),\ (3,4),                             \tag{D.ST2}\label{maps:eq:D.ST2}
\]
and there is an integral basis $(\alpha,\beta,\lambda,\ell)$
with $\gamma^*(\ell)=1$ and
\[
\begin{gathered}
 A\alpha=-\alpha,\quad A\beta=-\beta-\lambda,\quad
 A\lambda=\lambda-\alpha,\quad A\ell=\ell,\\
 \delta=-\alpha+2\lambda,\qquad
 A^2-1=\alpha\otimes\beta^*.
\end{gathered}                                               \tag{D.ST3}\label{maps:eq:D.ST3}
\]
Here $\alpha$ is primitive. The exact integral change of basis
for every marked row is in the certificate. For example, in a
standard marking put
\[
\begin{aligned}
 \alpha&=2e_1+a\delta,&
 \lambda&=e_1+\tfrac{a+1}{2}\delta,\\
 \beta&=e_2+\tfrac{3a-1}{2}\lambda,&
 \ell&=-\tfrac{3b}{4}e_1+\tfrac b2e_2+\gamma.
\end{aligned}
\]
The exceptional marking on $50/II$ is also checked integrally.
In this basis the degree-two fibre coinvariants are freely
generated by
\[
 x=[\beta\wedge\lambda],\qquad y=[\beta\wedge\ell],
\]
with relations
\[
 [\alpha\wedge\beta]=-2x,\quad
 [\lambda\wedge\ell]=-2y,\quad
 [\alpha\wedge\lambda]=[\alpha\wedge\ell]=0.                \tag{D.ST4}\label{maps:eq:D.ST4}
\]
In particular the boundary group $H_2(\partial U_*;\mathbb Z)$
is free of rank four; there is no boundary torsion that can be
discarded. The exterior fibre functional is still
$\omega=a e_1^*e_2^*+b\gamma^*\delta^*$, and
\[
 \omega(x)=-a,\qquad \omega(y)=b/4.                         \tag{D.ST5}\label{maps:eq:D.ST5}
\]

\item \emph{2. The normalization and its integral parity.}
We compute the retraction of $U_*$ to its central fibre by its
normalization and conductor. Since a simply connected output
has $\gcd(m_*,n_*)=1$, its star numerator is odd. All quotient
maps on the transverse elliptic curves below are therefore
unramified, with the displayed degrees.

\item For $m_*=2$ the normalization is an oriented $S^2$ bundle $N$
over a two-torus. Its conductor is a connected bisection $C$,
which maps with degree two to the double curve $D$. The relation
$A\lambda=\lambda-\alpha$ in \eqref{maps:eq:D.ST3} determines the parity of
this bundle. In a transverse circle coordinate $u$ its
pole-interchanging clutch can be represented by
\[
 z\longmapsto e^{-2\pi i u}/z.
\]
The corresponding $\mathrm{PGL}_2$ loop has representative
$\left(\begin{smallmatrix}0&e^{-2\pi i u}\\1&0\end{smallmatrix}\right)$.
Its determinant has odd winding, so its lift to the double cover
of $\mathrm{SO}(3)$ changes sign. Hence $w_2(N)$ evaluates to one
on the base. The remaining affine translations and unit
multipliers change the flat part of the clutch, not this winding;
the odd numerator changes the primitive base marking by an odd
index and preserves the parity. This connects the parity
calculation to every actual affine cap, not only a selected
linear model.

\item Let $f$ be the sphere fibre and choose a section class $\sigma$
with $\sigma^2=1$. Such a choice follows by adding multiples of
$f$ to any odd-square section. The bisection has zero
self-intersection, so
\[
 [C]=2\sigma-f.
\]
The normalization pushout gives
\[
 H_2(U_*)=
 \langle\sigma,f,D\mid 2\sigma-f-2D=0\rangle
          \cong\mathbb Z^2.                               \tag{D.ST6}\label{maps:eq:D.ST6}
\]
There is no additional $H_1$-kernel term: the conductor map to
$D$ is injective on $H_1$. The regular torus
$\alpha\wedge\beta$ maps to $f$, since it sweeps the single
semistable component with degree one. The regular
$\lambda\wedge\ell$ torus lies in the smooth part of the
normalization, covers its base twice and is disjoint from $C$.
Its class is consequently $2\sigma-f$, hence $2D$ in the
pushout. By \eqref{maps:eq:D.ST4},
\[
 x\longmapsto-\sigma+D,\qquad y\longmapsto-D.
\]
These are an integral basis. Thus the fibre-coinvariant map has
cokernel zero in this case.

\item For $m_*=4$ with the two components interchanged, choose one
component upstairs. Its stabilizer is generated by the square
of the deck generator. That square fixes
$\alpha,\lambda,\ell$. The open $\mathbb C^*$ bundle of its
normalization therefore has an abelian period lattice. More
explicitly it has generators $\alpha,\lambda,\ell,c^2$ and one
relation
\[
 2c^2=t\alpha+r\lambda+n_*\ell.
\]
Because $n_*$ is odd, $\alpha$ is primitive in this lattice.
The associated complex line over the quotient two-torus has
zero Chern class. Its two pole sections are thus homologous;
write their common class as $\sigma$. They map with degree two
to two different double curves $D_0,D_1$. Hence
\[
\begin{split}
 H_2(U_*)&=
 \langle f,\sigma,D_0,D_1\mid
                  \sigma-2D_0=\sigma-2D_1=0\rangle\\
 &\cong\mathbb Z f\oplus\mathbb Z D_0
                       \oplus\mathbb Z/2\,(D_1-D_0).
\end{split}                                                 \tag{D.ST7}\label{maps:eq:D.ST7}
\]
Again the conductor map on $H_1$ is injective into the sum of
its double-curve targets, so this is the complete group.

\item The regular $\alpha\wedge\beta$ torus sweeps both semistable
components, each with degree one onto the chosen component
fibre downstairs, and therefore maps to $2f$. The regular
$\lambda\wedge\ell$ torus covers the normalization base twice.
It avoids both pole sections, so its class is $2\sigma=4D_0$.
Thus, modulo possible multiples of the displayed order-two
class,
\[
 x\longmapsto-f,\qquad y\longmapsto-2D_0.
\]
The free part of this map has determinant two. Including the
order-two summand, its full cokernel has order four, independent
of the two possible torsion corrections. This is precisely the
integral effect missed by the rational vertical coinvariant
identity $[\delta]=4[\beta]$.

\item For completeness, the conductor has no third homology, and
its degree-two map into the direct sum of the double-curve
groups is injective: each conductor component maps with degree
two to its corresponding double curve. Thus the normalization
sequence identifies $H_3(U_*)$ with $H_3(N)=\mathbb Z^2$.
It also gives $H_4(U_*)=H_4(N)=\mathbb Z$, generated by the
reduced central component. A positive multiple of this divisor is principal,
being the divisor of the base parameter. Its intersection with
every two-cycle is therefore zero. Under Poincare--Lefschetz
duality this is the entire map
$H_2(U_*)\to H_2(U_*,\partial U_*)=H^4(U_*)=\mathbb Z$.
It follows that
\[
 H_2(\partial U_*)\longrightarrow H_2(U_*)
 \quad\hbox{is onto integrally}.                            \tag{D.ST8}\label{maps:eq:D.ST8}
\]

\item Let $C_*=H_2(T^4)_{A_*}$ and let $K_*$ be the full kernel in
\eqref{maps:eq:D.ST8}. The boundary Wang sequence gives
\[
 0\longrightarrow C_*\longrightarrow H_2(\partial U_*)
        \longrightarrow\operatorname{Fix}(A_*)
        \longrightarrow0.
\]
The fibre map $C_*\to H_2(U_*)$ is injective, and the preceding
calculations therefore give
\[
 [\operatorname{Fix}(A_*):\operatorname{proj}K_*]
 =\bigl|H_2(U_*)/\operatorname{im}C_*\bigr|
 =\begin{cases}1,&m_*=2,\\4,&m_*=4\text{ with swap}.
   \end{cases}                                             \tag{D.ST9}\label{maps:eq:D.ST9}
\]
This proves \eqref{maps:eq:D.ST1}, retaining the normalization torsion and
both full normal directions.

\item \emph{3. The full normal determinant and the global obstruction.}
Here is a useful integral determinant calculation. Suppose all
caps have the minimal rational ranks and their boundary maps
in degree two are onto. Put
\[
 \kappa_i=[\operatorname{Fix}(A_i):
                         \operatorname{proj}K_i],
\]
where $K_i$ is the full integral boundary kernel. The primitive
graph calculation for all $172$ rows gives the exterior group
as $\mathbb Zs$ plus the direct sum of the fixed sweep lattices,
with the sole relation $\sum_i t_i(\delta)+be\,s=0$.
Over $\mathbb Q$, each $K_i$ is the graph of a linear fibre
functional on its fixed lattice. Its value on $\delta$ is fixed
without choosing any of its other coefficients: the
$\delta$ normal torus bounds on the root cover, and transfer
with the actual affine power vector gives the exact relation
\[
 m_i t_i(\delta)+bn_i s=0.                                 \tag{D.ST10}\label{maps:eq:D.ST10}
\]
The sign is the same oriented convention as \eqref{maps:eq:D.HY3}.

\item Choose integral bases of the projected kernels. Their
covolumes are $\kappa_i$. Eliminate all independent sweep
coordinates over $\mathbb Q$ and then the single diagonal
$\delta$ relation. The determinant of the complete square
Mayer--Vietoris matrix is
\[
 \left(\prod_i\kappa_i\right)
                   b\left(e-\sum_i n_i/m_i\right).         \tag{D.ST11}\label{maps:eq:D.ST11}
\]
All longitudinal and mixed fibre coefficients cancel: only
the values \eqref{maps:eq:D.ST10} on the diagonal vector enter the last
Schur complement. Multiplying the rational calculation by the
integral covolumes is an equality of integer determinants,
not a saturation assumption. The certificate also verifies
this identity symbolically with arbitrary two-by-two normal
projection matrices and arbitrary longitudinal coefficients.

\item For an unchanged minimal cap or an index-one multiplicative
cap, $\kappa_i=1$. For a finite good-reduction cap \ref{maps:src:D.2.26} gives
$\kappa_i=m_i/q_i$. For a regular logarithm its projected
kernel is $m_i\Lambda+\mathbb Z v_i$, of index $m_i^3$;
this is the full lattice because the affine action is free.
These statements include every cap that could remain in a
sphere after the \ref{maps:src:D.2.22} rank obstruction.

\item If $X$ is simply connected, the Seifert determinant is $\pm1$,
so
\[
 \left|e-\sum_i n_i/m_i\right|=(\prod_i m_i)^{-1}.
\]
All degree-one groups in this Mayer--Vietoris map are free.
For the star boundary this follows from \eqref{maps:eq:D.ST3}; for the
remaining boundaries the local matrices have only unit
nonzero Smith factors, as checked in every retained row.
The star cap has the abelian presentation obtained from the
collapsed lattice and the relation $m_*c=v_*$; its relation
is primitive because $\gcd(m_*,n_*)=1$. The same primitive
relation proves freeness for the finite and regular log caps.
An unchanged minimal additive cap has the product/mapping-torus
model of \ref{maps:src:D.2.26} and free first homology. The exterior group is
free as well. The degree-one Mayer--Vietoris source and target
have equal rank, while their cokernel is $H_1(X)=0$.
Consequently that map is an isomorphism and has no additional
degree-one bridge kernel.
Thus \eqref{maps:eq:D.ST11} gives the exact order
\[
 |H_2(X)|=\frac{b}{q_*Q_o}\,R^2,\qquad
 q_*:=m_*/\kappa_*=
 \begin{cases}2,&m_*=2,\\1,&m_*=4\text{ with swap},\end{cases} \tag{D.ST12}\label{maps:eq:D.ST12}
\]
where $Q_o$ is the product of the indices $q_i$ of the other
finite caps and $R$ is the product of regular logarithmic
multiplicities.

\item Simple connectivity permits at most two multiple fibres and
forces their multiplicities to be coprime. Since $m_*$ is
even, every other multiplicity is odd. The only possible
additional finite index is $3$ from the $IV$ fibre of $50/IV$;
an order-six $II$ log cannot coexist with the star log in a
simply connected output. In every row of \eqref{maps:eq:D.ST2},
$b/Q_o$ has precisely two factors of $2$, and $R$ is odd.
Equation \eqref{maps:eq:D.ST12} now has two-primary order $2$ for $m_*=2$
and $4$ for $m_*=4$. Neither can be a homology sphere.
All excluded cap choices already have a positive rational
second Betti number or nontrivial fundamental group. This
completes the exhaustion.
\end{enumerate}
\end{proof}

\begin{corollary}[all twenty-one natural sphere profiles]\label{maps:src:D.2.28}
\StatementRecord{maps:src:D.2.28}{proved}{Historical events: Claude Fable 5.2 (Anthropic), effort high, 16 November 2026: carried as background without fault found. GPT-6 Astra (OpenAI), ChatGPT Pro mode, effort Pro, S6\_09\_PAPERS, 17 November 2026: targeted proof check. Also checked in this writing. Pre-continuation writing events: GPT-6 Astra (OpenAI), ChatGPT Pro mode, effort Pro, S6\_10\_PAPERS, 18 November 2026: targeted proof check. Continuation metadata events: GPT-6 Astra (OpenAI), ChatGPT Pro mode, effort Pro, S6\_10\_PAPERS continuation, 19 November 2026: used and rechecked in this run; status-vocabulary reconciliation; metadata reconciliation only; no new mathematical check. \emph{Pre-continuation recorded proof scope.} GPT-6 Astra (OpenAI, ChatGPT Pro mode), effort Pro, 2026-11-17 UTC; Claude Fable 5.2 (Anthropic), effort high, 16 November 2026; also checked in this writing by GPT-6 Astra (OpenAI), ChatGPT Pro mode, effort Pro, 18 November 2026}{the hypotheses, notation and definitions stated here; \cref{maps:src:D.2.22}; \cref{maps:eq:D.HY3}; \cref{maps:eq:D.HY5}; \cref{maps:src:D.2.23}; \cref{maps:src:D.2.24}; \cref{maps:src:D.2.27}; \cref{maps:src:D.2.25}; \cref{maps:src:D.3.2}}


For the natural line $M=\mathcal O(P-O)$, with all section pairs,
component-supported linearization zeros, permitted original and
Mumford fillings, additive and regular logarithms, arbitrary
multiplicities and affine twists in the sector of \ref{maps:src:D.2.22}, the
complete profile classification is as follows. A pattern means
a marked initial section/zero/additive-filling row, not an
individual complex structure or a sample of its multiplicities.
\[
\begin{array}{c|l|r|l}
\text{profile}&\text{singular fibres}&\text{sphere patterns}
                     &\text{sphere multiplicities}\\\hline
43&III^*,3I_1&3&\text{unbounded}\\
43/II&III^*,II,I_1&1&\text{unbounded}\\
45&I_8,4I_1&4&\text{all orders }1\\
45/II&I_8,II,2I_1&4&\text{unbounded}\\
46&I_3^*,3I_1&0&\text{none}\\
46/II&I_3^*,II,I_1&0&\text{none}\\
47&I_7,I_2,3I_1&4&\text{all orders }1\\
47/II&I_7,I_2,II,I_1&4&\text{unbounded}\\
47/III&I_7,III,2I_1&4&\text{unbounded}\\
49&IV^*,I_2,2I_1&3&\text{unbounded}\\
49/II&IV^*,II,I_2&1&\text{unbounded}\\
49/III&IV^*,III,I_1&2&\text{unbounded}\\
50&I_1^*,I_3,2I_1&0&\text{none}\\
50/II&I_1^*,II,I_3&0&\text{none}\\
50/IV&I_1^*,IV,I_1&0&\text{none}\\
55&I_5,I_4,3I_1&4&\text{all orders }1\\
55/II&I_5,I_4,II,I_1&4&\text{unbounded}\\
56&I_5,I_3,I_2,2I_1&5&\text{all orders }1\\
56/II&I_5,I_3,I_2,II&6&\text{unbounded}\\
56/III&I_5,I_3,III,I_1&5&\text{unbounded}\\
56/IV&I_5,I_2,IV,I_1&5&\text{unbounded}.
\end{array}                                                  \tag{D.ST13}\label{maps:eq:D.ST13}
\]
Thus twelve profiles admit unbounded sphere multiplicities,
four admit spheres only with all logarithmic orders one, and
five admit none. There are exactly $59$ marked sphere patterns
on $50$ distinct natural initial rows. Their full individual
conditions, including all affine parameters, are \eqref{maps:eq:D.HY3}--\eqref{maps:eq:D.HY5}
and the complete table
\path{records/manifolds_nonstar_sphere_patterns.csv}.
The profile summary is also saved as
\path{records/manifolds_complete_natural_sphere_profiles.csv}.
Every sphere in this classification is diffeomorphic to the
standard $S^6$.

\end{corollary}


\begin{proof}
\begin{enumerate}[label=\arabic*.,leftmargin=*]
\item \ref{maps:src:D.2.22} is the exhaustive finite initial reduction; \ref{maps:src:D.2.23} gives
the full fundamental group for every order. \ref{maps:src:D.2.24} eliminates
the two $I_3^*$ profiles, \ref{maps:src:D.2.27} eliminates the three $I_1^*$
profiles, and \ref{maps:src:D.2.25}--\ref{maps:src:D.2.26} compute every remaining sphere pattern
and all its unrestricted multiplicity conditions. Their
simply connected homology-sphere outputs are homotopy spheres
by Hurewicz and Whitehead, and $\Theta_6=0$ gives the stated
smooth type, as in \ref{maps:src:D.3.2}. The exact program
\path{code/III/manifolds_i1star_obstruction.py} checks all sixteen
star blocks, all eight integral canonical forms, both
normalization presentations, every possible torsion marking
in the order-four fibre map, the symbolic determinant identity,
and the twenty-one-profile tally.
\end{enumerate}
\end{proof}

The scope here is the natural section-line sphere question.
The optional vertically corrected catalogue of $226$ rows in
\ref{maps:src:D.2.22} has not been substituted for it. Formula \eqref{maps:eq:D.ST13} concerns the sphere outputs. The unrestricted
finite-log homology is supplied in \ref{maps:src:D.2.29}--\ref{maps:src:D.2.33}, and the entire specified
minimal simply connected smooth sector is completed by \ref{maps:src:D.2.35}--\ref{maps:src:D.2.36},
\ref{maps:src:D.2.38}--\ref{maps:src:D.2.39}, \ref{maps:src:D.2.42} and \ref{maps:src:D.2.44}--\ref{maps:src:D.2.45}. The remaining weighted and
additional-regular-log smooth classification is scoped in \ref{maps:src:D.2.4}.


\begin{theorem}[all free finite twists in the primitive-$\gamma$ sector]\label{maps:src:D.2.29}
\StatementRecord{maps:src:D.2.29}{proved}{Historical events: Claude Fable 5.2 (Anthropic), effort high, 16 November 2026: carried as background without fault found. GPT-6 Astra (OpenAI), ChatGPT Pro mode, effort Pro, S6\_09\_PAPERS, 17 November 2026: targeted proof check. Also checked in this writing. Pre-continuation writing events: GPT-6 Astra (OpenAI), ChatGPT Pro mode, effort Pro, S6\_10\_PAPERS, 18 November 2026: targeted proof check. Continuation metadata events: GPT-6 Astra (OpenAI), ChatGPT Pro mode, effort Pro, S6\_10\_PAPERS continuation, 19 November 2026: used and rechecked in this run; status-vocabulary reconciliation; metadata reconciliation only; no new mathematical check. \emph{Pre-continuation recorded proof scope.} GPT-6 Astra (OpenAI, ChatGPT Pro mode), effort Pro, 2026-11-17 UTC; Claude Fable 5.2 (Anthropic), effort high, 16 November 2026; also checked in this writing by GPT-6 Astra (OpenAI), ChatGPT Pro mode, effort Pro, 18 November 2026}{the hypotheses, notation and definitions stated here; \cref{maps:src:D.2.26}; \cref{maps:eq:D.HY2}; \cref{maps:eq:D.SC2}; \cref{maps:src:D.2.23}; \cref{maps:eq:D.UF6}; \cref{maps:eq:D.UF8}; \cref{maps:eq:D.HY3}; \cref{maps:eq:D.UF9}}


The restriction $\gcd(m_i,n_i)=1$ in the homology part of \ref{maps:src:D.2.26}
can be removed. The resulting exact formula retains torsion in
the finite cap itself. We keep the same natural initial rows,
minimal original and Mumford fillings, and the finite sites
whose integral fixed lattice contains $\ell_i$ with
$\gamma^*(\ell_i)=1$. These are the $68$ marked finite positions
certified in \ref{maps:src:D.2.26}; their orders and indices are $(d_i,q_i)$
from \eqref{maps:eq:D.HY2}.

At such a position write
\[
 m_i=d_i k_i,\qquad v_i=n_i\ell_i+r_i\delta,\qquad
 t_i=\gcd(m_i,n_i,q_i r_i),\qquad k_i>0.
\]
Every integer power vector and affine lift in this sector is
allowed exactly when
\[
 \gcd(m_i,n_i,r_i)=1,\qquad d_i\mid m_i/t_i.                 \tag{D.UF1}\label{maps:eq:D.UF1}
\]
This is an exact freeness criterion for all powers of the
cyclic action. Under it $t_i$ is either $1$ or $q_i$; the latter
case, when $q_i>1$, is permitted precisely when $q_i\mid k_i$.
The finite central quotient surface $F_i$ and its complex
normal line $L_i$ satisfy
\[
\begin{gathered}
 H_1(F_i)=\mathbb Z^2\oplus\mathbb Z/t_i,\qquad
 H_2(F_i)=\mathbb Z^2\oplus\mathbb Z/t_i,\\
 H_3(F_i)=\mathbb Z^2,\qquad H_0(F_i)=H_4(F_i)=\mathbb Z,\\
 \operatorname{ord}c_1(L_i)=t_i.
\end{gathered}                                               \tag{D.UF2}\label{maps:eq:D.UF2}
\]
Here order one means the zero class. Thus the normal line need
not be topologically trivial.

The full global degree-two presentation has a fibre generator
$s$, two fixed sweeps $D_i,L_i'$ for every finite log, and the
four sweeps $t_j(\Lambda)$ for every regular log. The finite
relations are
\[
\begin{aligned}
 (bn_i/q_i)s+(m_i/q_i)D_i&=0,\\
 -br_i s+m_iL_i'&=0,\\
 r_iD_i+n_iL_i'&=0.
\end{aligned}                                                \tag{D.UF3}\label{maps:eq:D.UF3}
\]
The regular-log relations remain those of \eqref{maps:eq:D.SC2}, and the
single global relation is
\[
 \sum_iD_i+\sum_jt_j(\delta)+be\,s=0.                       \tag{D.UF4}\label{maps:eq:D.UF4}
\]
There is no coprimality assumption on $m_i,n_i$ beyond the
actual freeness criterion \eqref{maps:eq:D.UF1}. The regular logs still have
arbitrary free affine vectors.

Let $T$ be the torsion of the cokernel of this matrix. Let
$\mathbb Z^r\oplus A_1$ be the Smith cokernel of the full
Seifert matrix of \ref{maps:src:D.2.23}, including these finite numerators
and all regular numerators. Then $r=0$ or $1$, and the complete
integral homology is
\[
\begin{array}{c|c}
 j&H_j(X;\mathbb Z)\\\hline
0,6&\mathbb Z\\
1&\mathbb Z^r\oplus A_1\\
2,3&\mathbb Z^{2r}\oplus T\\
4&\mathbb Z^{2r}\oplus A_1\\
5&\mathbb Z^r.
\end{array}                                                  \tag{D.UF5}\label{maps:eq:D.UF5}
\]
The complete fundamental group is still the Seifert
presentation, rather than its abelianization. In particular
this extends the non-star minimal-filling homology algorithm
to every free affine twist at each of its primitive-$\gamma$
finite sites.

\end{theorem}


\begin{proof}
\begin{enumerate}[label=\arabic*.,leftmargin=*]
\item \emph{1. Exact freeness and the cap groups.}
Let $A$ be the finite monodromy, of order $d$, and omit the
site index. Its nontrivial eigenvalues are the two eigenvalues
of the corresponding finite elliptic monodromy. Consequently
$A^j$ has the same rational fixed subspace as $A$ whenever
$d\nmid j$. Write $P$ for the rational projection to this fixed
subspace. The integral certificates of \ref{maps:src:D.2.26} give
\[
 \operatorname{Fix}(A)=\mathbb Z\delta\oplus\mathbb Z\ell,
 \qquad P\Lambda=\mathbb Z(\delta/q)\oplus\mathbb Z\ell,
 \qquad N_m\beta=(m/q)\delta,                               \tag{D.UF6}\label{maps:eq:D.UF6}
\]
where $N_m=1+A+\cdots+A^{m-1}$ and $\beta$ is the actual
primitive vertical coinvariant used there. The same certificates
give
$\Lambda/(A-1)\Lambda=\mathbb Z\beta\oplus\mathbb Z\ell$,
with $\delta=q\beta$ in that quotient.

\item Every affine action with the prescribed power is conjugate
on the good-reduction root disc to the invariant translation
form $h(x)=Ax+v/m$. Indeed its moving translation part has
zero finite-group average and is a translation coboundary;
the averaging solution is holomorphic on the disc. This is
the actual affine comparison of \ref{maps:src:D.2.26}, and it does not use
$\gcd(m,n)=1$. Constant unit choices extend over the cap and
preserve the marked normal comparison.

\item If $g=\gcd(m,n,r)>1$, then $h^{m/g}$ has integral translation
part and fixes the origin of its linear action. Freeness
therefore requires $g=1$. Assume this now. A power with
$d\mid j$ has identity linear part and nonintegral translation,
so is free. For $d\nmid j$, its moving linear map is surjective
on the moving torus. A fixed point exists precisely when its
fixed quotient translation is zero. By \eqref{maps:eq:D.UF6} this is
\[
 m\mid jn,\qquad m\mid jqr.
\]
The relevant powers are exactly the nonzero multiples of
$m/t$ below $m$, where $t=\gcd(m,n,qr)$. Since $g=1$ and
$q\in\{1,2,3\}$, $t$ divides $q$. Those powers all have
identity linear part exactly when $d\mid m/t$. This proves
\eqref{maps:eq:D.UF1} and the stated alternative for $t$. Freeness on the
central compact torus implies freeness on a sufficiently small
root neighborhood for this finite group, so the criterion is
also sufficient for the smooth holomorphic filling.

\item The quotient $F=T^4/\langle h\rangle$ has the affine fundamental
group presentation with the extra generator $c$ and the power
relation $mc=v$. On abelianizing, \eqref{maps:eq:D.UF6} gives
\[
 H_1(F)=\langle\beta,\ell,c\mid
                   mc-n\ell-qr\beta=0\rangle
        =\mathbb Z^2\oplus\mathbb Z/t.                     \tag{D.UF7}\label{maps:eq:D.UF7}
\]
It is an oriented four-manifold of Euler number zero and
rational first Betti number two. Integral Poincare duality and
the universal coefficient theorem give the other groups in
\eqref{maps:eq:D.UF2}, including the order-$t$ summand of $H_2$.

\item The normal line is the associated line for the root rotation
character. On $H_1(F)$ that character sends $c$ to $1/m$ in
$\mathbb R/\mathbb Z$ and the original fibre generators to
zero. The torsion generator
\[
 (m/t)c-(n/t)\ell-(qr/t)\beta
\]
has character $1/t$. The Bockstein of this character is
$c_1(L)$, so its order is exactly $t$. This proves the last
assertion of \eqref{maps:eq:D.UF2} without treating a torsion Chern class as
zero.

\item \emph{2. The full integral Thom kernel.}
The boundary is the marked mapping torus of $A$ and has
$H_1=\mathbb Z^3$, $H_2=\mathbb Z^4$. These statements follow
from the integral local Smith matrices, whose nonzero factors
are all units, in every one of the $68$ positions. The cap is
the disk bundle of $L\to F$. Since $c_1(L)$ is torsion, the
homology Gysin map $H_2(F)\to H_0(F)$ is zero. The pair sequence
therefore makes $H_2(\partial U)\to H_2(U)=H_2(F)$ onto,
including its torsion summand.

\item The full kernel is the image of the Thom map
\[
 H_3(U,\partial U)=H_1(F)\longrightarrow H_2(\partial U).
\]
The normal circle has the central group representative
$\mu=c^m v^{-1}$. For a lift of each of $\beta,\ell,c$, sweep
that lift by $\mu$. These three normal tori generate the
entire Thom image. On the boundary mapping cylinder, their
normal projections and their fibre components are respectively
\[
\begin{array}{c|c|c}
\text{generator}&\text{normal projection}&
                          \text{fibre component}\\\hline
\beta&N_m\beta&-\beta\wedge v\\
\ell&m\ell&-\ell\wedge v\\
c&v&0.
\end{array}                                                  \tag{D.UF8}\label{maps:eq:D.UF8}
\]
For the first row, subdividing the $m$-fold normal circle into
its $m$ mapping-cylinder intervals gives the sum $N_m\beta$;
the final power-vector edge supplies the displayed fibre
parallelogram. The second row is the same construction for an
invariant loop. The third is the torus of the two commuting
words $c$ and $c^m v^{-1}$, so its only normal sweep is $v$.
These are integral marked cycles, not invariant rational
forms.

\item Evaluation by the primitive exterior functional $\omega$,
using $-\omega(\beta,\ell)=b/q$ and
$\omega(\ell,\delta)=b$, turns \eqref{maps:eq:D.UF8} into the three columns
\[
 \begin{pmatrix}bn/q\\m/q\\0\end{pmatrix},\qquad
 \begin{pmatrix}-br\\0\\m\end{pmatrix},\qquad
 \begin{pmatrix}0\\r\\n\end{pmatrix}                       \tag{D.UF9}\label{maps:eq:D.UF9}
\]
on $(s,D,L')$. The target $H_2(\partial U)$ is free, so the
order-$t$ class of $H_1(F)$ maps to zero. Accordingly the sole
primitive relation among these columns is
\[
 (m/t)K_c-(n/t)K_\ell-(qr/t)K_\beta=0.
\]
No further saturation is allowed or needed: they are images
of generators of the whole Thom source. The projected
lattice has the exact index
\[
 \gcd\left(m^2/q,mn/q,mr\right)=(m/q)t.                    \tag{D.UF10}\label{maps:eq:D.UF10}
\]
In particular a torsion cap multiplies this index by $t$.
When $\gcd(m,n)=1$, the same columns reduce by Bezout to
\eqref{maps:eq:D.HY3}, recovering the previous theorem.

\item \emph{3. Global gluing and duality.}
Every degree-two boundary map is onto for these finite caps,
for the regular logs and for the retained minimal original and
Mumford caps. Eliminate their degree-two groups in
Mayer--Vietoris. The exterior fixed-sweep lattice is the
primitive direct sum with its single diagonal $\delta$
relation, as proved for all $172$ rows in \ref{maps:src:D.2.26}.
The full kernels \eqref{maps:eq:D.UF9}, the regular-log kernels and that
one diagonal relation give exactly \eqref{maps:eq:D.UF3}--\eqref{maps:eq:D.UF4}.
This calculation retains the cap torsion through its full
Thom image; it does not assume that the cap or its normal
line has torsion-free invariants.

\item The rational rank of this cokernel is $r$: after solving the
local sweep relations, the only remaining scalar equation is
$b(e-\sum n_i/m_i)s=0$, with the regular numerators included.
It has the same zero criterion as the Seifert matrix. The
kernel of the degree-one Mayer--Vietoris map is a subgroup
of the free boundary group and is therefore free. Its rank
is also $r$, because the source and target have equal free
rank and their cokernel has free rank $r$. Thus its extension
by the degree-two cokernel splits as an abelian group, giving
$H_2=\mathbb Z^{2r}\oplus T$. All caps still have the same
Euler contribution, so $\chi(X)=2$. Duality and the universal
coefficient theorem now give every remaining entry of \eqref{maps:eq:D.UF5}.
This includes $e-\sum n_i/m_i=0$ and does not suppress its
free groups or its torsion.
\end{enumerate}
\end{proof}

For example, take the $IV^*,II,I_2$ sphere initial row
$(a,b)=(1,3)$ of \ref{maps:src:D.2.11}, keep $II$ unchanged and use
$(q,m,n,r)=(3,9,3,1)$ at $IV^*$. The action is free, although
$\gcd(m,n)=3$. The result has
\[
\begin{gathered}
 \pi_1=\mathbb Z/3,\qquad H_1=H_4=\mathbb Z/3,\\
 H_2=H_3=(\mathbb Z/3)^2,\qquad H_5=0.
\end{gathered}
\]
Extrapolating the former primitive-numerator formula would
miss one order-three summand. For a $III^*$ sphere initial
row with $b=2$, the free data $(q,m,n,r)=(2,8,0,1)$ instead
give $\pi_1=\mathbb Z$, middle groups
$H_2=H_3=\mathbb Z^2\oplus\mathbb Z/2$, and
$H_4=\mathbb Z^2$, $H_5=\mathbb Z$.

The exact program
\path{code/III/manifolds_unrestricted_finite_logs.py} compares
\eqref{maps:eq:D.UF1} with every actual affine-power fixed-point test for
all $68$ marked matrices, $1\le k\le6$ and every pair
$(n,r)\bmod m$: $150059$ exact tests. It also compares $305$
full unreduced graph/normal Smith matrices with the reduced
presentation, retaining $99$ torsion-cap instances and $55$
zero-rational-Euler instances. These bounds describe the
independent checks; the proof and the matrix algorithm above
hold for every positive $k$ and every integer twist.


\begin{theorem}[every finite additive lattice and every free twist]\label{maps:src:D.2.30}
\StatementRecord{maps:src:D.2.30}{proved}{Historical events: Claude Fable 5.2 (Anthropic), effort high, 16 November 2026: carried as background without fault found. GPT-6 Astra (OpenAI), ChatGPT Pro mode, effort Pro, S6\_09\_PAPERS, 17 November 2026: targeted proof check. Also checked in this writing. Pre-continuation writing events: GPT-6 Astra (OpenAI), ChatGPT Pro mode, effort Pro, S6\_10\_PAPERS, 18 November 2026: targeted proof check. Continuation metadata events: GPT-6 Astra (OpenAI), ChatGPT Pro mode, effort Pro, S6\_10\_PAPERS continuation, 19 November 2026: used and rechecked in this run; status-vocabulary reconciliation; metadata reconciliation only; no new mathematical check. \emph{Pre-continuation recorded proof scope.} GPT-6 Astra (OpenAI, ChatGPT Pro mode), effort Pro, 2026-11-17 UTC; Claude Fable 5.2 (Anthropic), effort high, 16 November 2026; also checked in this writing by GPT-6 Astra (OpenAI), ChatGPT Pro mode, effort Pro, 18 November 2026}{the hypotheses, notation and definitions stated here; \cref{maps:src:D.2.29}; \cref{maps:eq:D.HY2}; \cref{maps:eq:D.UF2}; \cref{maps:src:D.2.25}; \cref{maps:eq:D.UF4}; \cref{maps:eq:D.UF5}; \cref{maps:src:D.2.23}; \cref{maps:src:D.2.28}; \cref{maps:src:D.2.26}; \cref{maps:eq:D.AF8}; \cref{maps:eq:D.UF8}}


The fixed-$\gamma$ restriction in \ref{maps:src:D.2.29} can also be removed.
All $102$ marked finite additive positions on the $172$ natural
candidate rows have the following integral description.
Let $d$ be the finite monodromy order, $q$ its index from
\eqref{maps:eq:D.HY2}, and put
\[
 \gamma^*(\operatorname{Fix}A)=g\mathbb Z,\qquad
 g\in\{1,q\},\qquad q_v=q/g.
                                                               \tag{D.AF1}\label{maps:eq:D.AF1}
\]
Here $g$ denotes the positive generator of that image.
There are $68$ positions with $g=1$ and $34$ with $g=q>1$.
Choose integral bases of the fixed and coinvariant lattices
so that
\[
\begin{gathered}
 \operatorname{Fix}A=\mathbb Z\delta\oplus\mathbb Z\ell,
 \qquad \gamma^*(\ell)=g,\\
 \Lambda/(A-1)\Lambda=\mathbb Z\beta\oplus\mathbb Z\eta,
 \qquad [\delta]=q_v\beta,\quad [\ell]=g\eta,\\
 N_m\beta=(m/q_v)\delta,
 \qquad N_m\eta=(m/g)\ell.
\end{gathered}                                                \tag{D.AF2}\label{maps:eq:D.AF2}
\]
The lifts $\beta,\eta$ are part of the retained integral
marking. In particular $\gamma^*(\eta)=1$.

For $m=dk>0$, every power vector is $v=r\delta+n\ell$.
Put
\[
 t=\gcd(m,q_vr,gn).
\]
The exact free-action criterion is
\[
 \gcd(m,n,r)=1,\qquad d\mid m/t.                            \tag{D.AF3}\label{maps:eq:D.AF3}
\]
The finite quotient and its normal line have precisely the
groups and the order-$t$ Chern class in \eqref{maps:eq:D.UF2}.
Their full normal Thom kernel, on the fibre generator $s$
and the sweeps $D,L'$ of $\delta,\ell$, has the three columns
\[
 \begin{pmatrix}bgn/q_v\\m/q_v\\0\end{pmatrix},\qquad
 \begin{pmatrix}-br\\0\\m/g\end{pmatrix},\qquad
 \begin{pmatrix}0\\r\\n\end{pmatrix}.                     \tag{D.AF4}\label{maps:eq:D.AF4}
\]
All entries are integers. The actual Seifert numerator is
$gn$, rather than $n$.

Consequently all $198$ minimal original/log additive-filling
patterns on all $148$ nonstar rows of the natural candidate
catalogue have a complete integral homology algorithm, for
every logarithmic multiplicity and every free affine twist.
Use \eqref{maps:eq:D.AF4} at each finite site, the full regular-log columns
of \ref{maps:src:D.2.25} at arbitrarily many regular sites, and the single
relation \eqref{maps:eq:D.UF4}. Take the torsion $T$ of the resulting Smith
cokernel, and take $r,A_1$ from the Seifert matrix with the
actual numerators $g_i n_i$. The answer is exactly \eqref{maps:eq:D.UF5},
including its zero-rational-Euler case. The complete
fundamental group remains \ref{maps:src:D.2.23}. The data of all $198$
patterns and all $102$ integral finite lattices are saved in
\path{records/manifolds_all_finite_logs.json}.

For one finite log and no regular logs there is also a closed
Smith formula, retaining an arbitrary integer clutch $e$.
Set
\[
 n'=n-(m/g)e,
 \qquad h=\gcd(b,m/g,n').                                  \tag{D.AF5}\label{maps:eq:D.AF5}
\]
If $n'\ne0$, then
\[
\begin{gathered}
 \pi_1(X)=\mathbb Z/(g|n'|),\qquad
 H_1(X)=H_4(X)=\mathbb Z/(g|n'|),\\
 H_2(X)=H_3(X)=
 \mathbb Z/h\oplus\mathbb Z/\bigl(bg|n'|t/(qh)\bigr),
 \qquad H_5(X)=0.
\end{gathered}                                               \tag{D.AF6}\label{maps:eq:D.AF6}
\]
When $n'=0$, instead
\[
\begin{gathered}
 \pi_1(X)=\mathbb Z,\qquad H_1(X)=H_5(X)=\mathbb Z,\\
 H_2(X)=H_3(X)=\mathbb Z^2\oplus
                         \mathbb Z/\gcd(b,m/g),\qquad
 H_4(X)=\mathbb Z^2.
\end{gathered}                                               \tag{D.AF7}\label{maps:eq:D.AF7}
\]
In both displays $H_0=H_6=\mathbb Z$, and a cyclic summand of
order one is omitted. Every case with $g>1$ is necessarily
nonsimply connected, so this extension does not alter \ref{maps:src:D.2.28}.

\end{theorem}


\begin{proof}
\begin{enumerate}[label=\arabic*.,leftmargin=*]
\item The $g=1$ basis and its full marked comparison were proved in
\ref{maps:src:D.2.26} and \ref{maps:src:D.2.29}. For the remaining positions the exact
matrices give
\[
 g=q,\qquad
 \Lambda/(A-1)\Lambda=\mathbb Z[\delta]\oplus\mathbb Z[\gamma],
 \qquad
 \ell=(q/d)N_d\gamma\in\Lambda.                            \tag{D.AF8}\label{maps:eq:D.AF8}
\]
Thus take $\beta=\delta$ and $\eta=\gamma$.
The vectors $\delta,\ell$ are an integral basis of the fixed
lattice, and their coinvariant images are $\delta,q\gamma$.
These assertions are checked by unimodular Smith operations
for all $34$ actual local matrices, not only by rational
conjugacy. They prove \eqref{maps:eq:D.AF2}. In all $102$ positions the
nonzero Smith factors of both $A-1$ and
$\bigwedge^2A-1$ are units. Hence the boundary groups used
in \ref{maps:src:D.2.29} are again free of the stated ranks.

\item The invariant exterior form satisfies
\[
 -\omega(\beta,\ell)=bg/q_v,\qquad
 \omega(\eta,\delta)=b,\qquad
 \omega(\eta,\ell)=0.
\]
For $g=q$ these follow directly from \eqref{maps:eq:D.AF8}, invariance of
$\omega$, and $\omega(\delta,\ell)=-bq$.
For $g=1$ they are the earlier identities. The required
divisibility $q_v\mid b$ is part of the same integral
catalogue, so the first coefficient of \eqref{maps:eq:D.AF4} is integral.

\item The fixed quotient torus has lattice $P\Lambda$, and the
coordinates of $v$ in it are $(q_vr,gn)$ by \eqref{maps:eq:D.AF2}.
Repeating the actual-power fixed-point argument of \ref{maps:src:D.2.29}
therefore gives \eqref{maps:eq:D.AF3}. The integer $t$ divides $q$ under
this freeness condition. The quotient surface has
\[
 H_1(F)=\langle\beta,\eta,c\mid
                   mc-q_vr\beta-gn\eta=0\rangle
       =\mathbb Z^2\oplus\mathbb Z/t.
\]
Its normal root character has value $1/t$ on the torsion
generator, exactly as before. Thus the normal line has Chern
class of order $t$, and duality gives all the cap groups.

\item The three full Thom generators in \eqref{maps:eq:D.UF8}, with $\eta$ in
place of the former $\ell$ coinvariant generator, now have
normal projections $(m/q_v)\delta,(m/g)\ell,v$.
Evaluation of their actual fibre parallelograms by $\omega$
gives \eqref{maps:eq:D.AF4}. Their primitive relation is
\[
 (m/t)K_c-(q_vr/t)K_\beta-(gn/t)K_\eta=0,
\]
and their normal-projection index is still $(m/q)t$.
The boundary-to-cap map in degree two remains onto, including
the torsion summand, because the normal Chern class is
torsion. Thus the full Mayer--Vietoris and duality argument of
\ref{maps:src:D.2.29} applies with the actual numerator $gn$ at every site.
This proves the all-order matrix algorithm.

\item To check its scope, each finite additive site can now be
log-transformed, including those with $g>1$; for example a
power along the primitive fixed quotient direction gives a
free action. An original cap is an additional choice exactly
when its line multidegree is zero and its component action
has one orbit. Those are precisely the original caps whose
full integral comparison is already proved in \ref{maps:src:D.2.26}.
Enumerating these choices gives $198$ patterns on all $148$
nonstar initial rows. It introduces no further multiplicity
restriction. Zero-free multiplicative sites and the unique
rank-two zero retain the marked minimal models already used
in that catalogue.

\item For the one-log formula eliminate $D=-be\,s$ in \eqref{maps:eq:D.AF4},
and subtract $e$ times the second column from the third.
On $(s,L')$ the matrix becomes
\[
 \begin{pmatrix}
  bgn'/q_v&-br&0\\
  0&m/g&n'
 \end{pmatrix}.                                            \tag{D.AF9}\label{maps:eq:D.AF9}
\]
Changing $n$ by $(m/g)e$ is exactly a peripheral lift change
by the coinvariant $\eta$. Thus \eqref{maps:eq:D.AF3} is preserved and
$t=\gcd(m,q_vr,gn')$. If $n'\ne0$, the gcd of the entries is
$h=\gcd(b,m/g,n')$: any prime dividing both $m/g$ and $n'$
cannot divide $r$ by freeness. The gcd of the two-by-two
minors is
\[
 bg|n'|t/q.
\]
These are the two determinantal divisors and prove the
middle groups in \eqref{maps:eq:D.AF6}. The Seifert presentation gives the
cyclic group of order $g|n'|$. If $n'=0$, the sole remaining
relation is $-brs+(m/g)L'=0$, whose torsion divisor is
$\gcd(b,m/g)$ since $r$ is prime to $m$ in this free case.
Its free rank one acquires the one free degree-one bridge
class, giving the rank two in \eqref{maps:eq:D.AF7}. Duality gives the
remaining groups in both cases.

\item Finally, if $g=q>1$, then $q$ divides both $m$ and its actual
numerator $gn$. Such a multiple fibre cannot occur in a
simply connected Seifert output by \ref{maps:src:D.2.23}. This verifies
that the new sector consists entirely of nonsimply-connected
members, without assuming that the group is abelian when
there are several logarithmic sites.
\end{enumerate}
\end{proof}

The exact program \path{code/III/manifolds_all_finite_logs.py}
checks all $34$ additional lattice bases, compares the
freeness criterion with $38675$ actual-matrix power tests
for $1\le k\le6$ and all residues, and compares $175$ full
unreduced normal/graph Smith matrices to \eqref{maps:eq:D.AF4}, including
$35$ torsion-cap instances. It also checks $725$ one-log
Smith forms with five different integer clutch values against
\eqref{maps:eq:D.AF6}--\eqref{maps:eq:D.AF7}. The finite checks certify the explicit
constant matrices and examples; the proof covers all orders
and all free integer twists. Initial weights outside the
finite sphere-candidate catalogue and nonminimal original
caps remain separate parts of the broader D.2 request.


\begin{theorem}[all zero-degree original caps and their free defects]\label{maps:src:D.2.31}
\StatementRecord{maps:src:D.2.31}{proved}{Historical events: Claude Fable 5.2 (Anthropic), effort high, 16 November 2026: carried as background without fault found. GPT-6 Astra (OpenAI), ChatGPT Pro mode, effort Pro, S6\_09\_PAPERS, 17 November 2026: targeted proof check. Also checked in this writing. Pre-continuation writing events: GPT-6 Astra (OpenAI), ChatGPT Pro mode, effort Pro, S6\_10\_PAPERS, 18 November 2026: targeted proof check. Continuation metadata events: GPT-6 Astra (OpenAI), ChatGPT Pro mode, effort Pro, S6\_10\_PAPERS continuation, 19 November 2026: used and rechecked in this run; status-vocabulary reconciliation; metadata reconciliation only; no new mathematical check. \emph{Pre-continuation recorded proof scope.} GPT-6 Astra (OpenAI, ChatGPT Pro mode), effort Pro, 2026-11-17 UTC; Claude Fable 5.2 (Anthropic), effort high, 16 November 2026; also checked in this writing by GPT-6 Astra (OpenAI), ChatGPT Pro mode, effort Pro, 18 November 2026}{the hypotheses, notation and definitions stated here; \cref{maps:src:D.2.30}; \cref{maps:src:D.2.23}; \cref{maps:src:D.2.14}; \cref{maps:src:D.2.26}; \cref{maps:src:D.2.16}; \cite[Section 2.4, equation (2.2)]{manUU21}; \cref{maps:eq:D.OD4}; \cref{maps:eq:D.OD5}; \cite[Theorems 2.2--2.3]{manBarden65}; \cite[Theorem A]{manGR25}}


In the natural candidate catalogue we may retain every original
additive cap at which $P$ is narrow, without requiring a single
component orbit. At $I_n^*$ this theorem uses the original cap;
finite good-reduction caps at the other additive sites have all
the choices and twists of \ref{maps:src:D.2.30}. The zero-free multiplicative
caps and the unique rank-two zero retain their preceding marked
models. These choices give $228$ additive-filling patterns on
$160$ natural initial rows. Exactly $30$ patterns have a positive
original component defect. The complete data are
\path{records/manifolds_original_zero_degree.json}.

For an original additive cap let $\rho_i$ be the number of
orbits of $Q$ on its components, and put
\[
 D=\sum_{i\text{ original additive}}(\rho_i-1).
\]
The only nonzero local defects in this catalogue are
\[
\begin{array}{c|rrrr}
\text{type}&III^*&IV^*&I_3^*&I_1^*\\\hline
\rho_i&5&3&3&2\\
\rho_i-1&4&2&2&1.
\end{array}                                                   \tag{D.OD1}\label{maps:eq:D.OD1}
\]
Every other permitted original cap has $\rho_i=1$.

Compute the Seifert group and its Smith cokernel
$\mathbb Z^r\oplus A_1$ as before, treating an unchanged cap
as an order-one meridian with zero numerator in the canonical
marking. Let $T$ be the torsion of the matrix in \ref{maps:src:D.2.30},
using its finite and regular columns and setting both fixed
sweeps of every original additive cap to zero. Then the full
integral groups are
\[
\begin{array}{c|c}
j&H_j(X;\mathbb Z)\\\hline
0,6&\mathbb Z\\
1&\mathbb Z^r\oplus A_1\\
2&\mathbb Z^{2r+D}\oplus T\\
3&\mathbb Z^{2r+2D}\oplus T\\
4&\mathbb Z^{2r+D}\oplus A_1\\
5&\mathbb Z^r.
\end{array}                                                   \tag{D.OD2}\label{maps:eq:D.OD2}
\]
The complete fundamental group is \ref{maps:src:D.2.23}, including every
order and affine twist. Thus the component defect adds free
homology, with no uncomputed torsion extension.

There is a full smooth identification in the $b=1$ sector.
For every simply connected member, let $M$ be the product of
its regular logarithmic multiplicities. Then
\[
 X\cong
 \mathop{\#}^{D}(S^2\times S^4)
 \ \#\ \mathop{\#}^{D}(S^3\times S^3)
 \ \#\ \widetilde\Sigma_0 N_M,                             \tag{D.OD3}\label{maps:eq:D.OD3}
\]
with the oriented twisted suspension and the spin five-manifold
$N_M$ defined in \ref{maps:src:D.2.14}. The allowed free defects in this sector
are $D=0,1,2,4$. The identification retains the full refined
spin data, including even $M$; it is not inferred from
additive homology alone.

\end{theorem}


\begin{proof}
\begin{enumerate}[label=\arabic*.,leftmargin=*]
\item \emph{1. The actual original cap.}
At an additive fibre the condition that $P$ is narrow means
that $\mathcal O(P-O)$ has degree zero on every component.
Its underlying line on the rational Kodaira tree is trivial.
The unit equivariant multiplier can be removed by its
rigidified logarithm isotopy, exactly as in \ref{maps:src:D.2.26}; the tree
has $H^1=0$, and the isotopy has zero meridian winding on the
normal collar. The cap consequently retracts, with its marked
angular circle, onto
\[
 S^1_\delta\times\operatorname{MT}_Q(K),
\]
where $K$ is the underlying Kodaira tree and $Q$ acts by its
actual component translation. The nontrivial contraction
modulus is retained as the mapping-torus direction; only its
unit circle part is straightened. This is the same relative
comparison proved for the odd-star caps in \ref{maps:src:D.2.16}, and uses
only narrowness of $P$, not the number of component orbits.

\item The tree has $H_1(K)=0$ and one free degree-two class per
component. The permutation invariants and coinvariants are
both free of rank $\rho$. The mapping-torus sequence and
Kunneth therefore give
\[
 H_1(U)=\mathbb Z^2,\quad H_2(U)=\mathbb Z^{\rho+1},\quad
 H_3(U)=\mathbb Z^{2\rho},\quad H_4(U)=\mathbb Z^\rho.
                                                               \tag{D.OD4}\label{maps:eq:D.OD4}
\]
There is no permutation torsion in these groups.

\item The two full normal tori are still primitive. The $\delta$
normal torus bounds the section disc times its circle. For
the other, take the least power $Q^q$ that is narrow and its
local Neron logarithm. Its marked loop extends across the
base disc; the rigidified norm of the unit linearization has
no meridian winding, so it introduces no fibre correction.
The local existence and extension of this logarithm are as
in \cite[Section 2.4, equation (2.2)]{manUU21} and \ref{maps:src:D.2.16}.
For each actual matrix the fixed lattice has basis
$\delta,\ell$, with $\gamma^*(\ell)=q$. Thus these two
normal tori project to the \emph{entire} integral fixed
lattice, not a finite-index sublattice. Their values under
the exterior fibre functional are zero in the stated
rigidified marking.

\item We also retain the component cokernel. Since \eqref{maps:eq:D.OD4} is free,
Poincare--Lefschetz duality gives
$H_2(U,\partial U)=H^4(U)=\mathbb Z^\rho$.
The map from $H_2(U)$ is the folded component-intersection
map. On the component classes its rank is $\rho-1$:
the Kodaira intersection matrix is negative semidefinite
with one-dimensional radical generated by the fibre
multiplicity vector, and averaging to permutation invariants
preserves that radical. The exact folded matrices are also
retained in the certificate. The principal fibre relation
bounds the full rank by $\rho-1$, so the additional
$\delta$--mapping-torus class cannot increase it.

\item It follows from the pair sequence that
\[
 \operatorname{im}\bigl(H_2(\partial U)\to H_2(U)\bigr)
 \text{ is primitive of rank }2,
 \qquad
 H_2(U)/\operatorname{im}H_2(\partial U)
                 \cong\mathbb Z^{\rho-1}.                 \tag{D.OD5}\label{maps:eq:D.OD5}
\]
Primitivity is integral: the quotient is the image of a map
into the free group $H_2(U,\partial U)$.
The boundary $H_2$ is free of rank four in every such marked
additive case. Its kernel has rank two, and the two primitive
fixed-loop normal tori just constructed therefore exhaust it.
Indeed their projection is already the whole fixed lattice,
and the fibre part is injective rationally. This proves that
the only change in the local Mayer--Vietoris data is the
free summand in \eqref{maps:eq:D.OD5}.

\item At a zero-free multiplicative site in this natural catalogue,
$P$ narrow implies that $Q$ is transitive on its components.
The direct component data and the primitive vanishing
matrices verify this at every occurrence. Hence an allowed
zero-degree original multiplicative cap has $\rho=1$ and
is already covered by the earlier primitive comparison.
The additive component permutations give exactly \eqref{maps:eq:D.OD1}.

\item \emph{2. Global integral homology.}
The subgroups in \eqref{maps:eq:D.OD5} are primitive, so choose integral
splittings in all the original cap groups. Eliminating their
boundary images in Mayer--Vietoris leaves their direct sum
$\mathbb Z^D$, together with exactly the torsion presentation
of \ref{maps:src:D.2.30}. The other cap boundary maps are onto as proved
there. The original caps have the same first Betti number as
their finite or minimal comparisons, so the degree-one bridge
kernel is still free of rank $r$. Consequently
$H_2=\mathbb Z^{2r+D}\oplus T$.
Every added original cap has Euler number zero, and
$\chi(X)=2$ remains unchanged. Duality gives all of \eqref{maps:eq:D.OD2},
including $b_3=2r+2D$ and the paired middle torsion.
This proves the formula without an implicit integral
spectral-sequence splitting.

\item There are $396$ exact local zero-degree original certificates.
Enumerating ``original'' precisely when $P$ is narrow,
and ``finite log'' at every finite good-reduction site,
with an original cap required at $I_n^*$ in the present
sector, gives the stated $228$ patterns on $160$ rows.
No logarithmic order is bounded in this enumeration.

\item \emph{3. Smooth classification when $b=1$.}
For a simply connected member all remaining finite logs have
primitive actual gamma numerator. Since $b=1$, their indices
$q_i$ are one; a site of fixed-gamma index greater than one
would contradict simple connectivity. \ref{maps:src:D.2.30} and the
regular-log calculation therefore give
\[
 H_2(X)=\mathbb Z^D\oplus(\mathbb Z/M)^2,
 \qquad H_3(X)=\mathbb Z^{2D}\oplus(\mathbb Z/M)^2.
\]
The scalar circle is free on every original cap, by the
product description in part 1, and free on every finite
or regular logarithmic cap by the primitive gamma numerator.
At each zero-free multiplicative cap, $Q=G$ is transitive
and the pair consisting of its actual primitive vanishing
vector and $\delta$ has unit Smith factors. The program
checks this pair at every such site in the relevant patterns.
The rank-two zero has the same full original short lattice
and the same unique fixed sphere as \ref{maps:src:D.2.14}, with normal
bundle $\mathcal O(-1)_{+1}\oplus\mathcal O(-1)_{-1}$.
Thus this is the actual semifree circle action on $X$.

\item The fixed sphere is nullhomotopic. Its specialization comes
from the common fibre class $s$; it has no contribution to
the extra component cokernels in \eqref{maps:eq:D.OD5}. The submatrix on
$s$ and the delta sweeps is the unimodular Seifert matrix
when $b=1$ and $\pi_1(X)=1$. It therefore kills $s$
integrally. Hurewicz gives the assertion about the sphere.

\item We check spin despite the new free second homology. On every
original cap the holomorphic volume form is the product of
a nonvanishing canonical two-form on the local minimal
elliptic surface with $dv/v$. Translation by $Q$ preserves
the first factor, and the change of $dv/v$ is a base-surface
one-form whose wedge with that two-form is zero. The volume
form descends, so the canonical line of that cap is trivial,
as in \ref{maps:src:D.2.16}. The free part of $H_2(X)$ is generated by the
component cokernels of these caps. Hence $c_1(X)$ evaluates
to zero on it. Since $X$ is simply connected,
$H^2(X;\mathbb Z)$ is free, so $c_1(X)=0$ and $X$ is spin.

\item The equivariant surgery of \ref{maps:src:D.2.14} now gives
$Y=X\#(S^3\times S^3)$ with a free circle action and a
projectable spin structure. Its orbit five-manifold is
simply connected and spin, and the homotopy/Gysin sequences
give
\[
 (B,e)\cong
 \left(\mathop{\#}^{D+1}(S^2\times S^3)\#N_M,
                        (x,0,\ldots,0)\right).
\]
Here $e$ is primitive; the spin Smale--Barden classification
and its realization theorem identify this pair, as in
\cite[Theorems 2.2--2.3]{manBarden65}.
Apply the spin branch of \cite[Theorem A]{manGR25}
successively, retaining the zero Euler class on the other
summands. Together with its suspension formula this gives
\[
 Y\cong\mathop{\#}^{D}(S^2\times S^4)
       \#\mathop{\#}^{D+1}(S^3\times S^3)
       \#\widetilde\Sigma_0N_M.
\]
The full Zhubr classification permits the same cancellation
of one $S^3\times S^3$ summand as in \ref{maps:src:D.2.14} and \ref{maps:src:D.2.16};
it retains the refined class $p$ and the full class
$\mu\in H_6(K(H_2,2);\mathbb Z)$, including all free cubic,
mixed and finite data. The added $S^3\times S^3$ maps
constantly to this Eilenberg--Mac Lane space and spin-bounds
$S^3\times D^4$. It therefore changes neither bordism
invariant and adds only two to $b_3$. The diffeomorphism of
the stabilized manifolds identifies these complete data,
which classify the unstabilized manifolds as well.
This proves \eqref{maps:eq:D.OD3}, including its oriented framing and
all two-primary data.
\end{enumerate}
\end{proof}

The exact program \path{code/III/manifolds_original_zero_degree.py}
checks all $396$ local original-cap lattices and folded
intersection ranks, the $228$ additive patterns, and $144$
actual primitive scalar-circle pairs in the $b=1$ patterns
that can be simply connected. It recovers the original
odd-star free defects of \ref{maps:src:D.2.16}--\ref{maps:src:D.2.17} and adds the same uniform
comparison for every permitted narrow-$P$ section pair.
Original caps with nonzero component Chern vector and higher
initial linearization weights remain outside this theorem.


\begin{theorem}[all original additive choices on the natural catalogue]\label{maps:src:D.2.32}
\StatementRecord{maps:src:D.2.32}{proved}{Historical events: Claude Fable 5.2 (Anthropic), effort high, 16 November 2026: carried as background without fault found. GPT-6 Astra (OpenAI), ChatGPT Pro mode, effort Pro, S6\_09\_PAPERS, 17 November 2026: targeted proof check. Also checked in this writing. Pre-continuation writing events: GPT-6 Astra (OpenAI), ChatGPT Pro mode, effort Pro, S6\_10\_PAPERS, 18 November 2026: targeted proof check. Continuation metadata events: GPT-6 Astra (OpenAI), ChatGPT Pro mode, effort Pro, S6\_10\_PAPERS continuation, 19 November 2026: used and rechecked in this run; status-vocabulary reconciliation; metadata reconciliation only; no new mathematical check. \emph{Pre-continuation recorded proof scope.} GPT-6 Astra (OpenAI, ChatGPT Pro mode), effort Pro, 2026-11-17 UTC; Claude Fable 5.2 (Anthropic), effort high, 16 November 2026; also checked in this writing by GPT-6 Astra (OpenAI), ChatGPT Pro mode, effort Pro, 18 November 2026}{the hypotheses, notation and definitions stated here; \cref{maps:src:D.2.30}; \cref{maps:src:D.2.31}; \cref{maps:src:D.2.23}; \cref{maps:src:D.2.20}; \cref{maps:eq:D.OC6}}


Every original additive cap on the $172$ natural candidate
rows can now be included. At $I_n^*$ we retain the original
cap; at every finite good-reduction additive site both the
original cap and every free logarithmic cap are allowed.
Keep the preceding marked Mumford models at multiplicative
sites, with the zero-degree original alternatives already
proved. These choices give exactly $282$ additive-filling
patterns on all $172$ initial rows. Of them, $54$ involve an
original additive cap with nonzero component Chern class.
The complete table is
\path{records/manifolds_original_chern.json}.

At a nonzero-Chern original cap let $s_i$ be the number of
Kodaira components, and let $q_i^c$ be its component-group
order. Its full boundary kernel adds the fibre relation
\[
 h_i s=0,\qquad h_i=b/q_i^c,                               \tag{D.OC1}\label{maps:eq:D.OC1}
\]
in addition to its two full fixed-loop normal relations.
Here $s$ is the common exterior fibre generator.
The local cokernel is free of rank $s_i-2$.

Let $L$ be the number of such nonzero-Chern original caps,
and put
\[
 \beta=
 \sum_{\text{zero-Chern originals}}(\rho_i-1)
 +\sum_{\text{nonzero-Chern originals}}(s_i-2).
                                                               \tag{D.OC2}\label{maps:eq:D.OC2}
\]
Take $T$ from the exact Smith matrix of \ref{maps:src:D.2.30}, adding
\eqref{maps:eq:D.OC1} for every nonzero-Chern original and setting both
fixed sweeps to zero at every unchanged additive cap.
Take $\mathbb Z^r\oplus A_1$ from the same complete Seifert
matrix, with the actual finite and regular numerators.
When $L=0$ the answer is \ref{maps:src:D.2.31}. When $L>0$ it is
\[
\begin{array}{c|c}
j&H_j(X;\mathbb Z)\\\hline
0,6&\mathbb Z\\
1&\mathbb Z^r\oplus A_1\\
2&\mathbb Z^{\beta+L+r}\oplus T\\
3&\mathbb Z^{2\beta+2L}\oplus T\\
4&\mathbb Z^{\beta+L+r}\oplus A_1\\
5&\mathbb Z^r.
\end{array}                                                   \tag{D.OC3}\label{maps:eq:D.OC3}
\]
All positive logarithmic orders and all free affine twists
are included, with arbitrary regular logarithms. The full
fundamental group is still \ref{maps:src:D.2.23}. In particular the extra
degree-one Mayer--Vietoris kernel has been retained, including
when the rational Seifert Euler number is zero.

\end{theorem}


\begin{proof}
\begin{enumerate}[label=\arabic*.,leftmargin=*]
\item \emph{The marked principal two-torus bundle.}
At an original additive site in the natural catalogue, either
$P$ or $Q$ is narrow. If $P$ is narrow, \ref{maps:src:D.2.31} applies.
Otherwise $Q$ is narrow, and the component Chern vector of
$M=\mathcal O(P-O)$ is
\[
 c=e_{\overline P}-e_0.
\]
It is primitive, and its pairing with the Kodaira
multiplicity vector is zero. There are exactly $46$ marked
occurrences of this case. The local Neron logarithm of $Q$
contracts its translation to the identity. Its lift to the
line bundle is straightened by the rigidified unit isotopy,
while the line bundle with its nonzero Chern vector is
retained. This is the actual relative construction of
D.2.20: the original cap retracts to the principal $T^2$
bundle over the Kodaira tree with Chern pair $(c,0)$, in
circle directions $\delta,\gamma$.

\item For clarity, the isotopy does not require the line bundle
to be trivial on the tree. Its Chern vector is fixed
throughout. A base isotopy lifts to this circle bundle,
and the difference between two lifts is a unit-valued
function. Since the tree has $H^1=0$, the unit has a
logarithm; rigidification fixes its constant phase and its
normal-meridian winding. Thus the transported boundary
marking is the one used in the following calculation.

\item The two nonzero transgressions in the integral bundle
spectral sequence are
\[
 x\longmapsto c(x)\delta,\qquad
 x\otimes y\longmapsto c(x)\delta\wedge y.
\]
Their images are primitive. As in \ref{maps:src:D.2.20}, all surviving
groups are free and
\[
 \pi_1(U)=\mathbb Z\gamma,\quad
 H_1(U)=\mathbb Z,\quad H_2(U)=\mathbb Z^{s_i-1},\quad
 H_3(U)=\mathbb Z^{2s_i-1},\quad H_4(U)=\mathbb Z^{s_i}.
                                                               \tag{D.OC4}\label{maps:eq:D.OC4}
\]
The projection $H_2(U)\to H_2(K)$ identifies it with
$\ker c$.

\item \emph{The entire fibre and normal kernel.}
Let $B_i$ be the upper-left elliptic block of $A_i$, and
let $Q_i$ be its elliptic gamma column. Narrowness of $Q$
gives the integral vector $z_i$ satisfying
\[
 Q_i=(B_i-I)z_i.
\]
The logarithm comparison identifies the elliptic projection
of a regular fibre by
$e_1\mapsto e_1$, $e_2\mapsto e_2$, $\delta\mapsto0$ and
$\gamma\mapsto z_i$. Its area covector is therefore
\[
 w_i=(1,0,(z_i)_2,0,-(z_i)_1,0).                            \tag{D.OC5}\label{maps:eq:D.OC5}
\]
The regular elliptic fibre specializes to the primitive
component-multiplicity vector in $\ker c$. Hence the full
map $H_2(T^4)\to H_2(U)$ is $w_i$ times that primitive
vector. Its kernel is exactly the integral $\ker w_i$.
Primitivity of the multiplicity vector inside $\ker c$
is checked directly, and also follows from its primitivity
in the ambient free component group.

\item The two principal circle actions sweep the origin disc.
Their normal vectors are
\[
 \delta,\qquad \ell_i=(-z_i,0,1).
\]
They are an integral basis of $\operatorname{Fix}(A_i)$.
Their actual fibre-area corrections are zero by the
rigidified logarithm comparison. Every relevant local
$A_i-I$ and $\bigwedge^2 A_i-I$ has only unit nonzero Smith
factors, so the boundary groups have the same free ranks
as in \ref{maps:src:D.2.20}. The two normal tori split the Wang projection
integrally. The remaining boundary image is the primitive
regular-fibre image above, and therefore
\[
 H_2(U)/\operatorname{im}H_2(\partial U)
                        =\mathbb Z^{s_i-2}.                \tag{D.OC6}\label{maps:eq:D.OC6}
\]
The full boundary kernel consists of both normal tori and
$\ker w_i$; its fibre part cannot be omitted.

\item Now evaluate $\ker w_i$ by the primitive global functional
$\omega=a e_1^*e_2^*-b\delta^*\gamma^*$.
An integral basis of the kernel has values
\[
 0,-a(z_i)_2,0,a(z_i)_1,-b.
\]
Thus its exact image in the exterior fibre group is
\[
 h_i\mathbb Zs,\qquad
 h_i=\gcd\bigl(b,a(z_i)_1,a(z_i)_2\bigr).
                                                               \tag{D.OC7}\label{maps:eq:D.OC7}
\]
All natural pairs in this catalogue have $\gcd(a,b)=1$,
so this is $\gcd(b,(z_i)_1,(z_i)_2)$.
The actual $46$ integral $z_i$ vectors give
$h_i=b/q_i^c$ in every case. The component orders are
$2$ for $III,III^*$, $3$ for $IV,IV^*$ and $4$ for
$I_1^*,I_3^*$. These are the finite constant lattice
calculations in the certificate; no assertion about
unbounded twists is obtained by enumerating them.
This proves the exact fibre relation \eqref{maps:eq:D.OC1}.

\item \emph{The global bridge and integral duality.}
All local cokernels in \eqref{maps:eq:D.OC6} and \ref{maps:src:D.2.31} are free,
so split them integrally. Their sum contributes
$\mathbb Z^\beta$ to the degree-two cokernel.
The remaining part is exactly the exterior fixed-sweep
matrix with the full finite/regular columns and the fibre
relations \eqref{maps:eq:D.OC1}. If $L>0$, at least one nonzero integer
multiple of $s$ is killed. Over $\mathbb Q$, all fixed
sweeps then vanish as well. This matrix consequently has
rational cokernel rank zero, even when the Seifert Euler
number is zero. Its integral torsion is precisely $T$.

\item Each nonzero-Chern original cap has first Betti number
one instead of the two of its boundary fixed lattice.
The degree-one Mayer--Vietoris source therefore has free
rank $L$ more than its target. Its cokernel is
$H_1(X)=\mathbb Z^r\oplus A_1$, so its kernel is free of
rank $L+r$. It is a subgroup of the free boundary group;
finite torsion in a logarithmic cap does not introduce a
hidden torsion kernel. The extension of this free bridge
by the degree-two cokernel splits as an abelian group.
It follows that
$H_2=\mathbb Z^{\beta+L+r}\oplus T$.
Every original cap has Euler number zero, so $\chi(X)=2$.
Duality gives the other entries of \eqref{maps:eq:D.OC3}, in particular
$b_3=2b_2-2b_1=2\beta+2L$.
This proves the complete formula, including the separate
role of the bridge when $r=1$.

\item Finally every finite additive site permits its two choices,
original or arbitrary free finite log, whereas $I_n^*$ is
kept original in this sector. Enumerating these choices
on all $172$ rows gives $282$ patterns. The $228$ patterns
with no nonzero-Chern original are exactly the preceding
zero-degree sector; the additional $54$ are those to which
\eqref{maps:eq:D.OC3} applies. This exhausts the stated additive choices.
\end{enumerate}
\end{proof}

For the untwisted original Engel input $P=G,Q=6G$, the
$IV^*$ and $III$ original caps have $s_i=7,2$ and
$h_i=2,3$. Hence $\beta=5$, $L=2$, $T=0$ and $r=1$.
Formula \eqref{maps:eq:D.OC3} gives the integral Betti numbers
$(1,1,8,14,8,1,1)$, independently recovering \cref{push:src:A.4.2} and
\ref{maps:src:D.2.20}. The one-original/one-finite cases similarly recover
the five all-order formulas of \ref{maps:src:D.2.21}. The theorem extends
these checks to every ordered section pair in the natural
candidate catalogue and every allowed additive choice.

The exact program \path{code/III/manifolds_original_chern.py}
retains all $46$ primitive Chern vectors, actual logarithm
vectors, full area kernels and primitive fibre images,
and the full $282$-pattern table. It also compares the
unreduced graph/normal/fibre-kernel Smith matrices with
the reduced presentation for every new pattern: $74$ full
Smith comparisons, retaining $9$ torsion-cap instances and
$54$ zero-rational-Euler instances. Original multiplicative caps with nonzero
component Chern class, logarithmic Mumford star caps and
initial weights outside this finite catalogue remain
separate parts of the broader request.


\begin{theorem}[all original multiplicative choices on the natural catalogue]\label{maps:src:D.2.33}
\StatementRecord{maps:src:D.2.33}{proved}{Historical events: Claude Fable 5.2 (Anthropic), effort high, 16 November 2026: carried as background without fault found. GPT-6 Astra (OpenAI), ChatGPT Pro mode, effort Pro, S6\_09\_PAPERS, 17 November 2026: targeted proof check. Also checked in this writing. Pre-continuation writing events: GPT-6 Astra (OpenAI), ChatGPT Pro mode, effort Pro, S6\_10\_PAPERS, 18 November 2026: targeted proof check. Continuation metadata events: GPT-6 Astra (OpenAI), ChatGPT Pro mode, effort Pro, S6\_10\_PAPERS continuation, 19 November 2026: used and rechecked in this run; status-vocabulary reconciliation; metadata reconciliation only; no new mathematical check. \emph{Pre-continuation recorded proof scope.} GPT-6 Astra (OpenAI, ChatGPT Pro mode), effort Pro, 2026-11-17 UTC; Claude Fable 5.2 (Anthropic), effort high, 16 November 2026; also checked in this writing by GPT-6 Astra (OpenAI), ChatGPT Pro mode, effort Pro, 18 November 2026}{the hypotheses, notation and definitions stated here; \cref{maps:src:D.2.32}; \cref{maps:src:D.2.23}; \cref{maps:eq:D.OC1}; \cref{maps:src:D.2.31}; \cref{maps:eq:D.OM5}; \cref{maps:eq:D.OM8}; \cref{maps:eq:D.OM7}; \cref{maps:eq:D.OM6}; \cref{maps:eq:D.OM10}; \cref{maps:src:D.2.30}; \cref{maps:eq:D.OC3}; \cref{maps:eq:D.OM3}}


At every zero-free multiplicative site in the $172$ natural
initial rows, either the original cap or the marked primitive
Mumford cap may be used. Together with all the original and
finite logarithmic additive choices of \ref{maps:src:D.2.32}, retaining
original caps at $I_n^*$, this gives $495$ distinct integral homology-presentation
patterns, after identifying the boundary-kernel agreements
for zero-degree original/Mumford choices proved below.
Every original/Mumford choice represented by such an
agreement is also listed explicitly in
\path{records/manifolds_original_multiplicative.json}.
Every finite and regular logarithmic order and every free
affine twist remains allowed.

The new original caps occur at $120$ marked zero-free
$I_n$ sites with nonzero component Chern vector. Each adds
\[
 h_i s=0,\qquad h_i=b/n,                                   \tag{D.OM1}\label{maps:eq:D.OM1}
\]
to the exact global matrix, has free local degree-two
cokernel of rank $n-2$, and adds one to the free
degree-one Mayer--Vietoris bridge. Consequently \ref{maps:src:D.2.32}
continues to hold verbatim if $L$ counts \emph{all}
nonzero-Chern original caps, additive or multiplicative,
and $\beta$ includes $n-2$ for each new multiplicative
original. In particular, when $L>0$,
\[
\begin{array}{c|c}
j&H_j(X;\mathbb Z)\\\hline
0,6&\mathbb Z\\
1&\mathbb Z^r\oplus A_1\\
2&\mathbb Z^{\beta+L+r}\oplus T\\
3&\mathbb Z^{2\beta+2L}\oplus T\\
4&\mathbb Z^{\beta+L+r}\oplus A_1\\
5&\mathbb Z^r.
\end{array}                                                \tag{D.OM2}\label{maps:eq:D.OM2}
\]
Here $\mathbb Z^r\oplus A_1$ is the exact Seifert
abelianization of \ref{maps:src:D.2.23}, and $T$ is the torsion of the
full finite/regular normal matrix with all the additive
relations \eqref{maps:eq:D.OC1} and multiplicative relations \eqref{maps:eq:D.OM1}.
The full fundamental group is still the Seifert group,
including when its rational Euler number is zero.

\end{theorem}


\begin{proof}
\begin{enumerate}[label=\arabic*.,leftmargin=*]
\item \emph{The actual bundle over the nodal cycle.}
At a zero-free $I_n$ site in these rows, either $P$ is
narrow or $Q$ is narrow. If $P$ is narrow, $Q$ acts
transitively on the components. Its original cap has
one component orbit and the same full integral normal
kernel and primitive fibre image as the marked Mumford
cap, by the multiplicative calculation in \ref{maps:src:D.2.31}.
This proves equality of the recorded integral homology
presentations for these zero-degree choices. It does not
assert a relative diffeomorphism between their caps.

\item Otherwise $n>1$, $Q$ is narrow and the component of $P$
is a unit modulo $n$. Write $K=I_n$ and
\[
 c=e_{\overline P}-e_0\in H^2(K;\mathbb Z)=\mathbb Z^n .
\]
This vector is primitive and its sum is zero. The local
Neron logarithm of $Q$ contracts its translation to the
identity. The lift preserves the line bundle with Chern
vector $c$; this bundle is not replaced by a trivial one.
There is one additional point compared with the tree
argument of D.2.32: the dual graph of $I_n$ has a circle.
On a cycle, translation in the identity component acts
on the generalized Jacobian of a multidegree-$c$ line
by the power given by $\sum_i c_i$. That sum is zero,
so the translation isotopy admits line-bundle lifts.
One can see this directly on the normalization: multiply
the coordinate on every rational component by the same
unit and use its integral powers in the line transition
functions; the final clutch factor is the unit to
$\sum_i c_i=0$.

\item The ratio between the resulting lift and the prescribed
holomorphic lift on $K$ is a global holomorphic unit on
the compact connected nodal cycle, and hence is constant.
It has no winding on the dual circle. Rigidification
removes this constant phase, and the nonvanishing
linearization on the base disc has a logarithm, fixing
the normal-meridian marking as well. Thus no additional
mixed Chern class is introduced. The actual cap retracts,
with its transported boundary marking, to
\[
 U\simeq Y_c\times S^1_\gamma,\qquad
 S^1_\delta\longrightarrow Y_c\longrightarrow K,           \tag{D.OM3}\label{maps:eq:D.OM3}
\]
where the circle bundle has Chern class $c$.

\item The integral Gysin sequence, or the bundle spectral
sequence over $K\simeq S^1\vee\bigvee_{i=1}^nS^2$, gives
\[
\begin{aligned}
 H_1(Y_c)&=\mathbb Z,&
 H_2(Y_c)&=\mathbb Z^n,&
 H_3(Y_c)&=\mathbb Z^n,\\
 H_1(U)&=\mathbb Z^2,&
 H_2(U)&=\mathbb Z^{n+1},&
 H_3(U)&=\mathbb Z^{2n},&
 H_4(U)&=\mathbb Z^n .
\end{aligned}                                             \tag{D.OM4}\label{maps:eq:D.OM4}
\]
All groups are free: the only transgression is the
primitive vector $c$. More precisely,
\[
 0\longrightarrow\mathbb Z[\delta\times\text{dual circle}]
 \longrightarrow H_2(Y_c)
 \longrightarrow\ker(c:\mathbb Z^n\to\mathbb Z)
 \longrightarrow0                                        \tag{D.OM5}\label{maps:eq:D.OM5}
\]
is an exact sequence of free groups. This also identifies
its first generator primitively.

\item \emph{The full marked kernel, including its normal correction.}
Let $u_1$ be the primitive elliptic vanishing circle and
choose $u_2$ with oriented determinant one. As in \ref{maps:src:D.2.32},
narrowness gives an integral vector $z$ satisfying
$(B-I)z=Q_{\rm col}$, and put $\ell=(-z,0,1)$.
In the integral basis $(u_1,u_2,\delta,\ell)$ the actual
monodromy has the form
\[
 A u_2=u_2+\nu u_1+p\delta,\qquad
 A u_1=u_1,\quad A\delta=\delta,\quad A\ell=\ell,
 \qquad |\nu|=n,\quad\gcd(\nu,p)=1.                        \tag{D.OM6}\label{maps:eq:D.OM6}
\]
These are equalities in the marked lattice, not only
a rational normal form. The certificate records the
unimodular basis for each of the $120$ occurrences.
In particular the full fixed lattice is
$\mathbb Zu_1\oplus\mathbb Z\delta\oplus\mathbb Z\ell$,
and both $A-I$ and $\bigwedge^2 A-I$ have only unit
nonzero Smith factors.

\item The regular fibre has primitive rank-three image in
$H_2(U)$. Its generators can be taken as
\[
 u_1\wedge u_2,\qquad u_2\wedge\delta,\qquad
 u_2\wedge\ell.                                           \tag{D.OM7}\label{maps:eq:D.OM7}
\]
The first projects to the all-one component vector in
$\ker c$, which is primitive; a possible addition of
the first term of \eqref{maps:eq:D.OM5} does not affect this assertion.
The second generates that first term primitively, and
the third is the dual-circle class times $\gamma$.
Hence their span is primitive. The full regular-fibre
kernel is exactly
\[
 K_{\rm fib}=
 \mathbb Z(u_1\wedge\delta)\oplus
 \mathbb Z(u_1\wedge\ell)\oplus
 \mathbb Z(\delta\wedge\ell).                              \tag{D.OM8}\label{maps:eq:D.OM8}
\]
For the first two generators, the vanishing disc and
its lifted products provide the required boundaries.
For the third, primitivity of $c$ kills the circle
$\delta$ integrally in $Y_c$. The image of
$\bigwedge^2 A-I$ is a primitive rank-two sublattice
of \eqref{maps:eq:D.OM8}.

\item The normal tori in the $\delta$ and $\ell$ directions
bound the sweeps of the origin disc under the two
principal circle actions. Their marked fibre corrections
are zero. A normal torus in the $u_1$ direction requires
slightly more care. Its projection to the original
elliptic neighbourhood is the usual vanishing-loop
normal torus, which bounds by the local vanishing-disc
construction. Its image in $H_2(U)$ therefore lies in
the first summand of \eqref{maps:eq:D.OM5}; its $\gamma$ coordinate
is zero. It is an integral multiple of the regular
torus $u_2\wedge\delta$. Subtracting that multiple
produces an actual boundary-kernel cycle with normal
projection $u_1$.

\item This correction cannot be discarded at the local level.
However the global primitive alternating functional
satisfies exactly
\[
 \omega(u_2,\delta)=0,\qquad
 \omega=a e_1^*e_2^*-b\delta^*\gamma^*.                    \tag{D.OM9}\label{maps:eq:D.OM9}
\]
Both $u_2$ and $\delta$ have zero $\gamma$ coordinate,
and $\delta$ has zero elliptic projection. Thus the
correction has zero image in the exterior fibre group.
The actual global image of the three normal kernels is
therefore the full marked fixed-sweep lattice with
zero fibre coefficient. This proves the needed marked
statement without guessing the local correction integer.

\item The boundary image is exactly the primitive rank-three
regular-fibre image \eqref{maps:eq:D.OM7}, and consequently
\[
 H_2(U)/\operatorname{im}H_2(\partial U)
                          =\mathbb Z^{n-2}.               \tag{D.OM10}\label{maps:eq:D.OM10}
\]
Evaluate \eqref{maps:eq:D.OM8} by $\omega$. Monodromy invariance gives
$\nu\,\omega(u_1,\ell)+p\,\omega(\delta,\ell)=0$,
where $\omega(\delta,\ell)=-b$ and
$\omega(u_1,\delta)=0$. Since $\gcd(\nu,p)=1$,
$n\mid b$ and the exact common divisor of these values
is $b/n$. This proves \eqref{maps:eq:D.OM1}, including its integral
coefficient.

\item \emph{Global gluing and exhaustion of the stated choices.}
The cap has first Betti number two, whereas the fixed
lattice in \eqref{maps:eq:D.OM6} has rank three. Thus it adds precisely
one free bridge, just as a nonzero-Chern additive original
does. The local degree-two cokernel \eqref{maps:eq:D.OM10} is free,
so it splits from the remaining global presentation.
The preceding kernel calculation shows that the latter
is exactly the old fixed-sweep matrix with the new
relation \eqref{maps:eq:D.OM1}. All finite torsion caps and their
nontrivial normal Chern classes are still handled by
the full three-column Thom matrices of \ref{maps:src:D.2.30}.

\item The proof of \eqref{maps:eq:D.OC3} now applies with the enlarged $L$
and $\beta$. In particular, a positive $L$ kills a
nonzero multiple of $s$, so the reduced degree-two
matrix has rational cokernel zero, even when $r=1$.
The degree-one bridge has rank $L+r$. Every original
cap in \eqref{maps:eq:D.OM3} has Euler number zero, and the unique
index-one toric zero still contributes two. Integral
duality therefore gives \eqref{maps:eq:D.OM2}.
At the group level the new original cap only adds
vertical relations, including the vanishing direction
and $\delta$, and adds no relation on $\gamma$.
The complete vertical lattice was already killed in
\ref{maps:src:D.2.23}, so its Seifert presentation is unchanged.

\item Finally, choose original or marked Mumford independently
at each of the nonzero-Chern multiplicative sites in
the $282$ additive patterns. This yields $495$ integral homology-presentation patterns
on all $172$ initial rows, with all the zero-degree
original/Mumford alternatives separately recorded. It
exhausts exactly the choices stated in the theorem.
\end{enumerate}
\end{proof}

The exact program \path{code/III/manifolds_original_multiplicative.py}
verifies all $120$ integral normal forms, primitive fibre
and monodromy kernels, actual values of $\omega$, and
primitive Chern and component-multiplicity vectors.
Their types are $I_8$ in $6$ occurrences, $I_7$ in $14$,
$I_2$ in $31$, $I_3$ in $23$, $I_5$ in $38$ and $I_4$
in $8$. It compares the full unreduced graph/normal/fibre
matrices with the stated reduced presentation for every
new pattern: $213$ full Smith comparisons, retaining
$18$ torsion-cap and $172$ zero-rational-Euler instances.
These are independent checks of an all-order proof.
The logarithmic Mumford choices at star fibres and
initial linearization weights outside the finite natural
catalogue remain separate parts of the broader request.


\begin{theorem}[every positive weight on the original Engel profile]\label{maps:src:D.2.34}
\StatementRecord{maps:src:D.2.34}{proved}{Historical events: Claude Fable 5.2 (Anthropic), effort high, 16 November 2026: carried as background without fault found. GPT-6 Astra (OpenAI), ChatGPT Pro mode, effort Pro, S6\_09\_PAPERS, 17 November 2026: targeted proof check. Also checked in this writing. Pre-continuation writing events: GPT-6 Astra (OpenAI), ChatGPT Pro mode, effort Pro, S6\_10\_PAPERS, 18 November 2026: targeted proof check. Continuation metadata events: GPT-6 Astra (OpenAI), ChatGPT Pro mode, effort Pro, S6\_10\_PAPERS continuation, 19 November 2026: used and rechecked in this run; status-vocabulary reconciliation; metadata reconciliation only; no new mathematical check. \emph{Pre-continuation recorded proof scope.} GPT-6 Astra (OpenAI, ChatGPT Pro mode), effort Pro, 2026-11-17 UTC; Claude Fable 5.2 (Anthropic), effort high, 16 November 2026; also checked in this writing by GPT-6 Astra (OpenAI), ChatGPT Pro mode, effort Pro, 18 November 2026}{the hypotheses, notation and definitions stated here; \cite[Lemma 2.1 and Remark 8.5]{manEngel26}; \cref{maps:src:D.2.5}; \cref{maps:src:D.2.29}; \cref{maps:src:D.1.2}; \cref{maps:eq:D.WE11}; \cref{maps:src:D.2.32}; \cref{maps:eq:D.WE9}; \cref{maps:eq:D.WE12}; \cref{maps:eq:D.WE13}}


On the $IV^*,III,I_1$ rational elliptic surface, let $G$
generate the Mordell--Weil group, retain
$M_S=\mathcal O(G-O)$, and use translation $6dG$ with
linearization zero $dI_1$, for any integer $d>0$.
Use the smooth $d$-period hexagonal Mumford cap at $I_1$.
At either additive site one may retain the original cap
or use any free finite logarithmic cap described below.
This gives an unbounded initial-weight family outside
the finite sphere-candidate catalogue.

For logarithmic caps put
\[
 m_0=3k,\quad m_1=4\ell,\qquad
 v_i=r_i\delta+n_i u_i,\qquad k,\ell>0.
\]
The exact freeness conditions are
\[
\begin{aligned}
 t_0&=\gcd(m_0,n_0,3r_0),&
 &\gcd(m_0,n_0,r_0)=1,\quad 3\mid m_0/t_0,\\
 t_1&=\gcd(m_1,n_1,2r_1),&
 &\gcd(m_1,n_1,r_1)=1,\quad 4\mid m_1/t_1 .
\end{aligned}                                             \tag{D.WE1}\label{maps:eq:D.WE1}
\]
There is no coprimality restriction on $d$.
All noninvariant parts of affine translations are
included through the marked translation conjugacy.
At an original site set its effective Seifert data to
$\overline m_i=1,\overline n_i=0$; otherwise set them
to $m_i,n_i$. Define
\[
 P=\overline m_1\overline n_0+
   \overline m_0\overline n_1,\qquad
 \mathbb Z^r\oplus A_1=
 \operatorname{coker}
 \begin{pmatrix}
  -\overline n_0&\overline m_0\\
  -\overline n_1&-\overline m_1
 \end{pmatrix}.                                           \tag{D.WE2}\label{maps:eq:D.WE2}
\]
The full fundamental group is this abelian group, and
$r=1$ precisely when $P=0$.

The exact degree-two matrix has target $(s,t,t_0',t_1')$.
For two logarithmic caps it is $B=(B_0\ B_1)$, where
\[
 B_0=
 \begin{pmatrix}
 -2dn_0&-6dr_0&0\\
 -k&0&r_0\\
 0&m_0&n_0\\
 0&0&0
 \end{pmatrix},
 \qquad
 B_1=
 \begin{pmatrix}
 -3dn_1&-6dr_1&0\\
 2\ell&0&-r_1\\
 0&0&0\\
 0&m_1&n_1
 \end{pmatrix}.                                           \tag{D.WE3}\label{maps:eq:D.WE3}
\]
For an original $IV^*$ or $III$ cap replace its three
columns, respectively, by
\[
 O_0=
 \begin{pmatrix}2d&0&0\\0&1&0\\0&0&1\\0&0&0\end{pmatrix},
 \qquad
 O_1=
 \begin{pmatrix}3d&0&0\\0&-1&0\\0&0&0\\0&0&1\end{pmatrix}.
                                                               \tag{D.WE4}\label{maps:eq:D.WE4}
\]
Let $T$ be the torsion of $\operatorname{coker}B$,
let $L$ count original additive caps, and put
$\beta=5$ if $IV^*$ is original, and $\beta=0$ otherwise.
Set
\[
 (b_2,b_3)=
 \begin{cases}
 (d-1+2r,\;2r),&L=0,\\
 (d-1+\beta+L+r,\;2\beta+2L),&L>0.
 \end{cases}                                               \tag{D.WE5}\label{maps:eq:D.WE5}
\]
Then all the integral homology is
\[
\begin{array}{c|c}
j&H_j(X;\mathbb Z)\\\hline
0,6&\mathbb Z\\
1&\mathbb Z^r\oplus A_1\\
2&\mathbb Z^{b_2}\oplus T\\
3&\mathbb Z^{b_3}\oplus T\\
4&\mathbb Z^{b_2}\oplus A_1\\
5&\mathbb Z^r .
\end{array}                                                \tag{D.WE6}\label{maps:eq:D.WE6}
\]
In particular $\chi(X)=2d$ for every choice.

There are useful closed forms. If both sites are
logarithmic and $\gcd(m_i,n_i)=1$, put
$g=\gcd(d,k,2\ell)$. Then
\[
 T=
 \begin{cases}
 \mathbb Z/g\oplus\mathbb Z/(d|P|/g),&P\ne0,\\
 \mathbb Z/g,&P=0.
 \end{cases}                                               \tag{D.WE7}\label{maps:eq:D.WE7}
\]
The fundamental group is $\mathbb Z/|P|$ for $P\ne0$
and $\mathbb Z$ for $P=0$. With at least one original
cap the torsion groups, for every free remaining twist,
are
\[
\begin{array}{c|c}
\text{original sites}&T\\\hline
IV^*\text{ only}&
 \mathbb Z/(d\gcd(2,n_1))\oplus\mathbb Z/\gcd(m_1,n_1)\\
III\text{ only}&
 \mathbb Z/(d\gcd(3,n_0))\oplus\mathbb Z/\gcd(m_0,n_0)\\
\text{both}&\mathbb Z/d .
\end{array}                                                \tag{D.WE8}\label{maps:eq:D.WE8}
\]
These are direct-sum descriptions, not necessarily
ordered invariant factors.

\end{theorem}


\begin{proof}
\begin{enumerate}[label=\arabic*.,leftmargin=*]
\item \emph{The weighted analytic construction.}
Engel's linearization identity for $M_S$ and $6G$
has base degree one
\cite[Lemma 2.1 and Remark 8.5]{manEngel26}.
Iterate that contracting generator $d$ times.
Translations preserve the base, so its linearization
is the product of $d$ copies and has divisor $dI_1$.
Equivalently, replace the long gamma period by $d$
times itself. On the punctured family this is the
ordinary torus cover defined by that sublattice.
The good-reduction models are preserved by it:
the integral matrices below have orders three and
four for every $d$. After the respective root base
changes one can therefore make the same finite
affine quotients as in \ref{maps:src:D.2.5} and \ref{maps:src:D.2.29}.

\item At the cusp retain the same unimodular triangular
fan, now modulo the sublattice generated by one
short-direction period and $d$ times the other.
Its quotient tube is smooth and proper by the same
finite-period toric construction as \ref{maps:src:D.1.2}. Its
normalization has $d$ copies of $dP_6$ and
\[
 H_*(N;\mathbb Z)
 =(\mathbb Z,\mathbb Z^2,\mathbb Z^{d+3},
                         \mathbb Z^2,\mathbb Z^d).         \tag{D.WE9}\label{maps:eq:D.WE9}
\]
The normalization proof of \ref{maps:src:D.1.2} uses
$\operatorname{diag}(1,t-1,t-1)$ over
$\mathbb Z[t]/(t^d-1)$ and is valid for every positive
$d$, including $2$ and $3$.
This constructs the weighted family directly.
It does not assume that the filled threefold is a
free cover of a weight-one finite quotient when $d$
shares a factor with its multiplicities.

\item \emph{Integral exterior and finite-cap data.}
Use the basis $(\delta,a,b,\gamma)$ of \ref{maps:src:D.2.5}, with
the new gamma period. The actual monodromies are
\[
\begin{aligned}
 A_0&=\begin{pmatrix}
 1&-1&0&2d\\0&-1&-1&6d\\0&1&0&0\\0&0&0&1
 \end{pmatrix},\\
 A_1&=\begin{pmatrix}
 1&0&1&-3d\\0&0&1&0\\0&-1&0&6d\\0&0&0&1
 \end{pmatrix},&
 A_\infty&=\begin{pmatrix}
 1&0&0&d\\0&1&1&0\\0&0&1&0\\0&0&0&1
 \end{pmatrix}.
\end{aligned}                                             \tag{D.WE10}\label{maps:eq:D.WE10}
\]
Their product is the identity. The two fixed lattices
are generated by $\delta$ and
\[
 u_0=\gamma+2d(a+b),\qquad
 u_1=\gamma+3d(a+b).
\]
The finite coinvariants still have
$b=-a,\delta=-3a$ and $b=-a,\delta=-2a$,
respectively, and $u_i$ maps to the primitive gamma
class. Thus their fixed-quotient indices and root
characters are exactly those of \ref{maps:src:D.2.29}, independently
of $d$. Its proof gives \eqref{maps:eq:D.WE1}, including every
nonprimitive free numerator, the torsion orders
$t_i$ of the caps and their normal lines, and the
three actual Thom generators.

\item Let $V$ be the regular bundle over a graph joining
the three peripheral circles. The common fibre
coinvariant in degree two is free of rank one:
its primitive functional is
\[
 \xi_d=a^*\wedge b^*+6d\gamma^*\wedge\delta^* .
\]
The integral relations are
$\delta a=\delta b=a\gamma=b\gamma=0$ and
$\delta\gamma=-6d\,ab$, so $s=[ab]$ is primitive.
For the graph part the five columns
\[
 (-\delta,-\delta),\quad(b,-a),\quad
 (\delta,0),\quad(u_0,0),\quad(0,u_1)                     \tag{D.WE11}\label{maps:eq:D.WE11}
\]
are an integral basis of
$\ker[(A_0-I)\ (A_1-I)]$.
They lie in the kernel and have a constant unit
$5$-row minor, using rows $0,2,3,4,7$ in zero-based
indexing. Their rational rank is five, so that unit
minor proves saturation for every $d$.
The common exterior-square relations likewise contain
a unit minor in the displayed primitive kernel of
$\xi_d$; the certificate gives its column indices
as a polynomial identity in $d$.

\item The first two columns of \eqref{maps:eq:D.WE11} are the actual
cusp normal tori in the angular directions
$\delta,a$. Their fibre coefficients are zero in
the rigidified cusp marking. Quotienting by these
two columns leaves the free target
$(s,t,t_0',t_1')$ used above.
The norms and contractions are
\[
\begin{aligned}
 N_0a&=-k\delta,&N_0\gamma&=m_0u_0,&
 \xi_d(a,u_0)&=2d,\\
 N_1a&=-2\ell\delta,&N_1\gamma&=m_1u_1,&
 \xi_d(a,u_1)&=3d,\\
 &&\xi_d(\gamma,\delta)&=6d .
\end{aligned}
\]
A covering-torus loop $\lambda$ has normal column
with graph part $N_i\lambda$ and fibre part
$-\xi_d(\lambda,v_i)$; the affine meridian has
graph part $v_i$ and fibre part zero. These are the
full Thom generators, including the torsion
character case of \ref{maps:src:D.2.29}. Since the delta sweep
at the second site is the negative of the first
modulo \eqref{maps:eq:D.WE11}, they give exactly \eqref{maps:eq:D.WE3}.

\item For an original additive cap, $Q=6dG$ is narrow and
the component Chern vector of $M_S$ is the same
primitive $e_G-e_O$ as at $d=1$. \ref{maps:src:D.2.32}'s actual
principal two-torus model applies with the marked
logarithm multiplied by $d$. The regular elliptic
area covectors, in the order
$(\delta a,\delta b,\delta\gamma,ab,a\gamma,b\gamma)$,
are
\[
 (0,0,0,1,-2d,2d),\qquad
 (0,0,0,1,-3d,3d).
\]
Their full kernels have $\xi_d$ images $2d\mathbb Z$
and $3d\mathbb Z$, and both complete fixed normal
lattices have zero fibre correction. Thus \eqref{maps:eq:D.WE4}
is the exact integral replacement. The local free
cokernels have ranks $7-2=5$ and $2-2=0$, and each
original loses one first-homology rank, giving
one additional free bridge. These statements retain
the primitive Chern class rather than making a
zero-degree approximation.

\item \emph{The cusp boundary torsion and the component meridians.}
Write $E=X\setminus\operatorname{int}N$ and $M=\partial N$.
For $d>1$, $H_1(M)$ has an actual $\mathbb Z/d$ summand:
the monodromy kills $a$ and $d\delta$, whereas the
cusp kills both primitive angular circles. It would
therefore be incorrect to reuse a free-boundary
Mayer--Vietoris splitting from weight one.
Instead use the exact sequence of the pair $(X,E)$.

\item The exterior-square cusp image is primitive of rank
two for every $d$, so $H_2(M)$ is free of rank six.
Its two angular normal tori bound by their toric
section-disc sweeps. They form a primitive sublattice
because their Wang projections are the primitive
fixed basis $\delta,a$.
Furthermore
$H_3(N,M)=H^3(N)=\mathbb Z^2$ by \eqref{maps:eq:D.WE9}.
Exactness bounds the kernel of
$H_2(M)\to H_2(N)$ by rank two. Hence those two
angular tori are its entire integral kernel.

\item Attaching the two finite or original additive caps
to $V$ first shows
\[
 H_2(E)/j_*\ker(H_2(M)\to H_2(N))
 \simeq\operatorname{coker}B\oplus\mathbb Z^{\beta+L}.
                                                               \tag{D.WE12}\label{maps:eq:D.WE12}
\]
Indeed every finite boundary-to-cap degree-two map
is onto, including when the cap and normal line
have torsion; its kernel is the full Thom image.
At an original cap split its free local cokernel.
The degree-one sequence then has a free kernel of
rank $L$, since the first Betti number of $E$ is
one and each original lowers the target rank by one.
The angular cusp classes come from $H_2(V)$, so
they have zero image in this new bridge. Thus
quotienting by them preserves its free splitting.
This proves \eqref{maps:eq:D.WE12} without dropping boundary
torsion or assuming a torsion extension splits.

\item Now $H_2(N,M)=H^4(N)=\mathbb Z^d$, with its basis
given by small positive transverse discs to the
$d$ components. Their boundary classes in $H_1(E)$
are \emph{identical}, not merely proportional.
In the smooth height-one fan a component ray
$(v,1)$ has its primitive meridian $(v,1)$ in
angular-plus-base coordinates. Two such meridians
differ by integral multiples of $a,\delta$.
Both angular classes already vanish in $H_1(E)$:
the unchanged vertical columns of $A_0-I,A_1-I$
generate the entire vertical lattice primitively.
Thus every disc has boundary the same class
$x_0+x_1$, with the common boundary orientation;
at an original site its $x_i$ is zero.

\item The abelian presentation of $H_1(E)$ has generators
$\gamma,x_0,x_1$ and relations
$\overline m_i x_i=\overline n_i\gamma$.
It has free rank one. Its element $x_0+x_1$ has
infinite order precisely when $P\ne0$.
Therefore the kernel of the actual map
\[
 \mathbb Z^d\longrightarrow H_1(E),\qquad
 (z_1,\ldots,z_d)\longmapsto
                 (z_1+\cdots+z_d)(x_0+x_1)
                                                               \tag{D.WE13}\label{maps:eq:D.WE13}
\]
is free of rank $d-1+r$. This remains true when
$H_1(E)$ or the indicated meridian has torsion.
The pair sequence and \eqref{maps:eq:D.WE12} now give a short
exact sequence with this free quotient, hence a
split sequence of abelian groups.

\item The matrix \eqref{maps:eq:D.WE3} has rational corank $r$.
For example a left-kernel vector has its other
three coordinates fixed by its $s$ coordinate;
the remaining equation is $dP$ times that coordinate.
If an original is present, \eqref{maps:eq:D.WE4} kills that
coordinate and the entire rational cokernel, so
the corank is zero. These facts and \eqref{maps:eq:D.WE13} give
the claimed $b_2$ and torsion $T$.

\item Van Kampen, separately, proves the group assertion.
The cusp kills the primitive $a,\delta$ loops;
finite monodromy then kills $b$. Only the central
$\gamma$ and the two meridians remain, with
$x_0x_1=1$ and the two power relations in \eqref{maps:eq:D.WE2}.
Thus the full group is abelian with that exact
integer presentation. Finally the cusp has Euler
number $2d$ by \eqref{maps:eq:D.WE9}, and every other cap has
Euler number zero. Integral duality and universal
coefficients give \eqref{maps:eq:D.WE5}--\eqref{maps:eq:D.WE6}.

\item \emph{Closed Smith forms.}
When both numerators are primitive, the same
Bezout reductions as \ref{maps:src:D.2.5} eliminate two unit
columns of \eqref{maps:eq:D.WE3} and leave
\[
 \begin{pmatrix}
 -2dn_0&-3dn_1\\-k&2\ell
 \end{pmatrix},\qquad \det=-dP .
\]
Its entries have gcd $\gcd(d,k,2\ell)$.
To check this prime by prime, a prime dividing
$k$ cannot divide $n_0$, one dividing $\ell$
cannot divide $n_1$, $3\nmid n_0$, and $2\nmid n_1$.
These observations give precisely the stated gcd
at $2,3$ and at every other prime. Its Smith
form gives \eqref{maps:eq:D.WE7}, including its rank-one case.
For an original $IV^*$, eliminate $t,t_0'$ using
\eqref{maps:eq:D.WE4}. The remaining relations are
$2ds=3dn_1s=0$ and
$m_1t_1'=n_1t_1'=0$, because $6dr_1s$ is already
zero. This is the first row of \eqref{maps:eq:D.WE8}.
The original $III$ case is identical with
$3d$ and $2dn_0$, and with both originals the
remaining fibre relation is $\gcd(2d,3d)s=ds=0$.
\end{enumerate}
\end{proof}

For two logarithmic caps, a simply connected member
has $P=\pm1$, which forces both local numerators
primitive and $\gcd(3k,4\ell)=1$. Thus $g=1$ and
\[
 H_2=\mathbb Z^{d-1}\oplus\mathbb Z/d,\qquad
 H_3=\mathbb Z/d,\qquad H_4=\mathbb Z^{d-1}.                \tag{D.WE14}\label{maps:eq:D.WE14}
\]
This includes even and three-divisible $d$, where the
free-cover reduction in D.1 is unavailable. At minimal
orders and $d$ prime to six it recovers D.1's groups.
The exact smooth classification beyond the cases
already settled there is not inferred from \eqref{maps:eq:D.WE14}.
With both original caps, \eqref{maps:eq:D.WE6} instead gives
\[
 \pi_1=\mathbb Z,\qquad
 H_2=\mathbb Z^{d+7}\oplus\mathbb Z/d,\qquad
 H_3=\mathbb Z^{14}\oplus\mathbb Z/d,\qquad
 H_4=\mathbb Z^{d+7},\quad H_5=\mathbb Z.
\]
At $d=1$ this is exactly \cref{push:src:A.4.2}; for every $d$ it is the
actual iterated-translation family.

The executable \path{code/III/manifolds_weighted_engel.py}
accepts arbitrary integer input, checks \eqref{maps:eq:D.WE1}, and
returns the full group and homology presentation.
It verifies polynomial identities for the monodromies,
primitive exterior-square quotient and graph basis,
and compares the unreduced graph and full area kernels
with \eqref{maps:eq:D.WE3}--\eqref{maps:eq:D.WE4}. Its retained output includes
nonprimitive free twists, nontrivial cap torsion,
every original/log choice and zero-Euler cases.


\begin{theorem}[the fifteen minimal $c=2$ patterns]\label{maps:src:D.2.35}
\StatementRecord{maps:src:D.2.35}{proved}{Historical events: Claude Fable 5.2 (Anthropic), effort high, 16 November 2026: carried as background without fault found. GPT-6 Astra (OpenAI), ChatGPT Pro mode, effort Pro, S6\_09\_PAPERS, 17 November 2026: targeted proof check. Also checked in this writing. Pre-continuation writing events: GPT-6 Astra (OpenAI), ChatGPT Pro mode, effort Pro, S6\_10\_PAPERS, 18 November 2026: targeted proof check. Continuation metadata events: GPT-6 Astra (OpenAI), ChatGPT Pro mode, effort Pro, S6\_10\_PAPERS continuation, 19 November 2026: used and rechecked in this run; status-vocabulary reconciliation; metadata reconciliation only; no new mathematical check. \emph{Pre-continuation recorded proof scope.} GPT-6 Astra (OpenAI, ChatGPT Pro mode), effort Pro, 2026-11-17 UTC; Claude Fable 5.2 (Anthropic), effort high, 16 November 2026; also checked in this writing by GPT-6 Astra (OpenAI), ChatGPT Pro mode, effort Pro, 18 November 2026}{the hypotheses, notation and definitions stated here; \cref{maps:eq:D.HY3}; \cref{maps:src:D.2.8}; \cref{maps:src:D.2.26}; \cref{maps:eq:D.C24}; \cref{maps:eq:D.C25}; \cref{maps:eq:D.C27}; \cref{maps:eq:D.C26}; \cref{maps:src:D.2.14}; \cref{maps:eq:D.HY7}; \cref{maps:eq:D.C215}; \cref{maps:eq:D.C212}; \cref{maps:eq:D.C213}; \cite[Theorems 3, 2.10 and Section 3.12]{manZhubr75}}


Take one of the fifteen patterns with $c=2$ in
\path{records/root_remaining_smooth_profiles.json}, use the primitive
Mumford filling at every multiplicative site, and make no additional
regular logarithmic transformations. Retain the additive original/log
choice of that pattern and the ordinary clutching $e\gamma$ of
\eqref{maps:eq:D.HY3}. Every simply connected output, for every allowed finite
multiplicity and affine twist, is orientation-preservingly diffeomorphic
to the conic surgery $M_{\mathrm{conic}}$ of \ref{maps:src:D.2.8}. Equivalently, if
$\alpha\ne0$ in $H^2(X;\mathbb F_2)$, then
\[
 H_2(X;\mathbb Z)=H_3(X;\mathbb Z)=\mathbb Z/2,
 \qquad \langle\alpha^3,[X]_2\rangle=1.                 \tag{D.C21}\label{maps:eq:D.C21}
\]
The marked patterns, with zero positions numbered from zero as in the
source tables, are
\[
\begin{array}{c|r|r|l|l}
\text{OS/profile}&a&b&\text{zero positions}&\text{additive choice}\\\hline
47&7&2&2,3,4&\text{none}\\
47/II&7&2&2&II\text{ original or log}\\
49&1&6&2,3&IV^*\text{ log}\\
56&5&2&1&\text{none}\\
56&15&2&3,4&\text{none}\\
56/II&5&2&3&II\text{ original or log}\\
56/IV&1&6&0&IV\text{ log}\\
56/IV&5&6&3&IV\text{ log}\\
56/IV&15&2&3&IV\text{ original}.
\end{array}                                                   \tag{D.C22}\label{maps:eq:D.C22}
\]
Thus these fifteen patterns occupy thirteen distinct marked initial
rows. The statement includes the original additive caps in this table;
it does not replace the stipulated multiplicative Mumford caps by
other caps, or include additional regular logarithmic transformations.

\end{theorem}


\begin{proof}
\begin{enumerate}[label=\arabic*.,leftmargin=*]
\item \emph{1. The common piece and its integral markings.}
Throughout, use the actual lattice basis $(u,w,\delta,\gamma)$ and
write
\[
 \Omega=a\,u^*\wedge w^*+b\,\gamma^*\wedge\delta^*.
                                                               \tag{D.C23}\label{maps:eq:D.C23}
\]
Here $a$ is odd, $b=2$ or $6$, and $\gcd(a,b)=1$.
Remove the additive caps and one ordinary disc carrying the clutching;
fill every multiplicative site by its stated Mumford cap, and denote
the resulting common piece by $W$. Denote its regular torus-bundle
part by $V$. In particular all finite affine twists are outside $W$,
and $V$ has the linear section used in the marked construction.

\item There is an exact integral fact about these thirteen rows. The columns
of their actual multiplicative angular vanishing lattices, without
saturation, generate precisely
\[
 \Lambda_0=\ker\gamma^*=\mathbb Zu\oplus\mathbb Zw
                    \oplus\mathbb Z\delta.                    \tag{D.C24}\label{maps:eq:D.C24}
\]
The same is true of all columns of $A_i-I$. Both assertions are checked
by integer Hermite forms, not ranks or component labels, in the
certificate below. Filling a Mumford cap kills its base meridian and
its angular lattice. The usual torus-bundle presentation therefore
gives
\[
 \pi_1(W)=\mathbb Z\gamma\times F_{r-1},
                                                               \tag{D.C25}\label{maps:eq:D.C25}
\]
where $r\ge1$ is the number of unfilled discs. Every remaining
monodromy is the identity on the quotient $\mathbb Z\gamma$.
The section makes the displayed extension split. In this table
$r\le2$, but only torsion-freeness is needed below.

\item We also need an integral cohomology comparison. Each primitive
Mumford cap has a surjective degree-two boundary map, as proved in
\ref{maps:src:D.2.26}. If there are $N$ original punctures, then $b_1(V)=N$ by the
second assertion in \eqref{maps:eq:D.C24}. A filled boundary has first Betti number
one greater than its cap. In the degree-one Mayer--Vietoris map the
target rank minus the source rank is consequently $r$, exactly
$b_1(W)$ by \eqref{maps:eq:D.C25}. Its kernel has rank zero and is a subgroup of a
free group, hence is zero. Eliminating the surjective cap degree-two
maps now shows that
\[
 H_2(V;\mathbb Z)\longrightarrow H_2(W;\mathbb Z)
                  \text{ is onto}.                            \tag{D.C26}\label{maps:eq:D.C26}
\]
Both $H_1(V)$ and $H_1(W)$ are free. Thus restriction from $W$ to $V$
is injective in degree-two integral cohomology and also with
$\mathbb F_2$ coefficients. These facts will prevent a swept-class
ambiguity in the auxiliary-line calculation.

\item \emph{2. The actual symmetric line and its Chern class.}
On a multiplicative chart let $h_1$ be the elliptic period, let $h_G$
represent the section $G$, and let $R$ be the contraction multiplier.
The section pair is $P=aG,Q=bG$. On an auxiliary line with coordinate
$\zeta$ put
\[
\begin{aligned}
 h(z,v,\zeta)&=(h_1z,h_G^av,(-1)^a z^{-a}\zeta),\\
 F(z,v,\zeta)&=(h_G^{-b}z,Rv,
                 h_G^{-ab/2}R^{b/2}v^b\zeta),\\
 J(z,v,\zeta)&=(z^{-1},v^{-1},z^{-a}\zeta).
\end{aligned}                                                   \tag{D.C27}\label{maps:eq:D.C27}
\]
All exponents are integers. Direct multiplication gives
$hF=Fh$, $J^2=1$, $JhJ=h^{-1}$ and $JFJ=F^{-1}$; the integral
logarithmic constant $a$ in the third identity exponentiates to one.
The displayed $J$ is normalized on the rigidifying origin.
The alternating Chern form is $\Omega$. The normalized half-period
signs are $(-1)^a$ on $u,w$ and $+1$ on $\delta,\gamma$; thus its
mod-two semicharacter is
\[
                    q=a(u+w+uw).                              \tag{D.C28}\label{maps:eq:D.C28}
\]
Every actual monodromy preserves this semicharacter. The certificate
checks this on all sixteen vectors modulo two at every marked site.
The symmetric rigidified Appell--Humbert identifications consequently
patch on overlap markings: an isomorphism preserving the rigidification
is unique, so their compositions satisfy the cocycle condition.
This constructs the holomorphic auxiliary line on $V$ from the actual
factors \eqref{maps:eq:D.C27}.
Since $b$ is even, the half-$\delta$ translation lifts in these factors
as $(z,v,\zeta)\mapsto(z,-v,\zeta)$. It commutes with $h,F$.
The overlap identifications preserve this lift: in the invariant
prequantum marking the lift depends only on $\Omega$ and $\delta/2$,
and each monodromy fixes $\delta$ and preserves $\Omega$ and $q$.
Thus the line has the global involution linearization used below.

\item Compare it with the flat prequantum family having connection potential
$\pi i\Omega(x,dx)$ and the same semicharacter, as in the fully marked
comparison in \ref{maps:src:D.2.8}. The monodromies preserve this family. Both lines
are symmetric and rigidified with alternating form $\Omega$; their
ratio is therefore fibrewise flat of order dividing two. The normalized
square-trivializations on the fibres are unique and patch, making the
ratio globally two-torsion. Since $H^2(V;\mathbb Z)$ is free, its Chern
class is zero. The actual line has precisely the pure fibre class
$\Omega$ and zero value on the marked swept cycles. This argument
controls the mixed values as well as the alternating fibre form.

\item These factors extend to a line $\mathcal L$ on $W$. Here is the
integral local extension, including the two parities that we need.
In a genuine $SL_2(\mathbb Z)$ elliptic marking the actual Tate
valuations are
\[
 (\operatorname{ord}h_1,\operatorname{ord}h_G,
       \operatorname{ord}R)=(n,j,\rho),\qquad
 P_{\rm val}=\begin{pmatrix}n&-bj\\aj&\rho\end{pmatrix}.
                                                               \tag{D.C29}\label{maps:eq:D.C29}
\]
The determinant is one at the zero and zero at every other
multiplicative site. These are raw valuations from the actual
monodromy; no residue of $j$ modulo $n$ is substituted.
At the zero, assign the following integer values to the lattice
vertices and interpolate linearly on the given unimodular fan:
\[
 \varphi(x,y)=
 \frac{-a\rho x^2-2abjxy+bn y^2+ax}{2}.                        \tag{D.C210}\label{maps:eq:D.C210}
\]
Indeed $n\rho+abj^2=1$, so $n,\rho$ are odd, and the expression is
integral at every integer vertex. Its two period differences are
\[
\begin{aligned}
 \varphi(x+n,y+aj)-\varphi(x,y)&=-ax,\\
 \varphi(x-bj,y+\rho)-\varphi(x,y)
          &=by+(b\rho-abj)/2.
\end{aligned}                                                   \tag{D.C211}\label{maps:eq:D.C211}
\]
They are exactly the valuations of the line factors in \eqref{maps:eq:D.C27}.
The usual monomial frames with these integral slopes therefore extend
the line and its commuting lattice identifications over every toric
chart. This uses an integral piecewise affine support function; no
positivity of the auxiliary line is required.

\item At a rank-one point write the two valuation periods as
$H=\alpha\nu$, $F=\beta\nu$, where
$\nu=(\nu_1,\nu_2)$ is primitive and
$\gcd(\alpha,\beta)=1$. Thus
$n=\alpha\nu_1$, $aj=\alpha\nu_2$,
$-bj=\beta\nu_1$ and $\rho=\beta\nu_2$.
The identity $-a\beta\nu_1=b\alpha\nu_2$ shows that
$k=a\nu_1/\alpha$ is an integer. On the valuation line use
\[
             \varphi_1(t)=-k t(t-\alpha)/2.
\]
Its two exact period differences are
$-a\nu_1t$ and
$b\nu_2t+(b\rho-abj)/2$, respectively, as required by
\eqref{maps:eq:D.C27}. At every odd-isotropy site $\nu_1$ is odd and
$\beta\nu_1=-bj$ is even, so $\beta$ is even and $\alpha$
is odd. The values of $\varphi_1$ at integer vertices are therefore
integral. At the unique even-isotropy site the valuation vector is
$(2,aj)$ with $aj$ odd, so $\alpha=1$ and the same integrality
holds. Integral monomial frames on the one-dimensional fan, retaining
the invertible second short coordinate, give the required extension.
The two displayed period identities prove that both actual deck
factors descend this extension; it is not enough to match only one
of them.
Holomorphic unit factors extend in the same frames and affect neither
integral slope nor stabilizer character. The local extensions agree
with the line on their annuli and hence glue to $\mathcal L$ on $W$.
By \eqref{maps:eq:D.C26}, its Chern class has exactly the pure fibre and zero swept
values already computed on $V$.

\item \emph{3. The fixed curves and the two line characters.}
Let $\iota$ be the scalar involution $(z,v)\mapsto(z,-v)$.
The exact primitive angular vectors show that its fixed locus on $W$
is the disjoint union of one elliptic curve $E$ at the unique $I_2$
site and the rational vertical-edge curve $C$ at the rank-two zero.
All other rank-one scalar stabilizers have odd order; this is checked
from the actual vanishing vectors in the certificate. The zero has
valuation lattice all of $\mathbb Z^2$, and $\delta=(0,1)$ in that
short lattice. Thus there is one vertical-edge orbit, with the same
normal model
$\mathcal O(-1)_{+1}\oplus\mathcal O(-1)_{-1}$ as in \ref{maps:src:D.2.14}.

\item At the zero the vertical slope difference is
\[
 \varphi(x,y+1)-\varphi(x,y)=-abjx+bny+bn/2.
\]
Its two nonconstant coefficients are even and its constant is odd,
since $b/2$ and $n$ are odd. The induced line character along $C$
is therefore nontrivial. The degree along its adjacent toric wall is,
up to the immaterial orientation convention for the divisor,
$-a\rho+abj$, which is odd. The other diagonal convention changes
this expression by an even integer, so
\[
 \operatorname{char}_{\iota}(\mathcal L|_C)=1,
          \qquad \deg(\mathcal L|_C)\equiv1\pmod2.            \tag{D.C212}\label{maps:eq:D.C212}
\]
At the $I_2$ site, $j$ is odd and $\rho=-abj^2/2$. The coordinates
\[
 x=z^{(1-aj)/2}v,\qquad y=z^{-aj}v^2
\]
are an integral unimodular change of the short torus coordinates;
they send the primitive vanishing vector $(2,aj)$ to $(1,0)$.
The scalar involution acts as $(-x,y)$, giving sign $-1$ on both
normal directions to $E$. The $h$ multiplier is, up to a unit,
$(-1)^a x^{-2a}y^a$. After replacing $F$ by $Fh^{bj/2}$, its
$x$ exponent is zero and its $y$ exponent is $b/2$. The support
function $-at(t-1)$ has all slopes even. Hence
\[
                 \operatorname{char}_{\iota}(\mathcal L|_E)=0.
                                                               \tag{D.C213}\label{maps:eq:D.C213}
\]

\item There are no further fixed points in the removed caps when they are
reattached. An original additive cap of this table has zero line
multidegree and a free scalar circle, by its actual product description
in \ref{maps:src:D.2.26}. On a finite cap, a scalar fixed-point equation on the covering
torus would imply $jn_i/m_i=0$ in $\mathbb R/\mathbb Z$ on applying
$\gamma^*$. Simple connectivity requires $\gcd(m_i,n_i)=1$, so the
root power is trivial; translation by $\delta/2$ cannot then fix a
torus point. Away from the central torus the base rotation gives the
same conclusion. The ordinary clutch disc also has free scalar action.
Thus $X^{\iota}=C\sqcup E$ for every multiplicity and twist in the
statement, including even finite multiplicities.

\item \emph{4. The global mod-two attachment.}
\ref{maps:src:D.2.26} gives $H_2(X)=H_3(X)=\mathbb Z/2$, generated in degree two by
the image of the primitive fibre class $s$; there are no additional
regular sweeps. We verify which cohomology class this is on $W$.
For every finite log in this table, $b/q_i=2$. Its complete two normal
columns \eqref{maps:eq:D.HY7} have pure fibre coefficients
\[
                 bn_i/q_i=2n_i,
                \qquad -\alpha_i b r_i,                       \tag{D.C214}\label{maps:eq:D.C214}
\]
where $\alpha_i m_i+\beta_i n_i=1$. Both coefficients are even,
for arbitrary allowed $m_i,n_i,r_i$. The original additive cap normal
columns have zero fibre coefficient in the rigidified marking.
The ordinary clutching $e\gamma$ has fibre coefficients divisible by
$b$, and its global relation is $be\,s$. Thus the functional that is
one on $s$ and zero on all swept generators annihilates every normal
column modulo two. It induces the nonzero class $\alpha$ on $X$.
This uses the full integral normal columns, not a transfer of even
degree. Its restriction to $V$ is exactly the reduction of the Chern
class just constructed, and \eqref{maps:eq:D.C26} now implies
\[
                  \alpha|_W=\rho_2c_1(\mathcal L).             \tag{D.C215}\label{maps:eq:D.C215}
\]
The absence of an $\operatorname{Ext}(H_1(W),\mathbb F_2)$ term is
part of \eqref{maps:eq:D.C25}, not an implicit universal-coefficient choice.

\item \emph{5. The full two-primary smooth invariant.}
Since $H^2(X;\mathbb Z)=0$, the complex manifold has $c_1=0$ and is
spin. Put $G=\langle\iota\rangle$ and let
$\tau\in H^1(BG;\mathbb F_2)$ be its generator. As in \ref{maps:src:D.2.8},
$\alpha$ has a Borel equivariant lift $\widetilde\alpha$:
$H^1(X;\mathbb F_2)=0$ kills the first possible obstruction, and the
fixed point splits the map to $BG$, killing the obstruction into
$H^3(BG;\mathbb F_2)$.

\item We check the relative fixed-point normalization on $W$. The scalar
circle acts trivially on its fundamental group, by an isotopy fixing
a point of $C$. That point splits the Borel extension, so
\[
                 \pi_1(W_G)=\pi_1(W)\times G.
\]
The first factor is torsion-free by \eqref{maps:eq:D.C25}. Consequently every fixed
point gives the same order-two element $(1,\iota)$ in this product.
By \eqref{maps:eq:D.C215}, the difference between $\widetilde\alpha|_{W_G}$ and
the equivariant Chern class of $\mathcal L$ vanishes in ordinary
cohomology. Pull it to the universal cover and use the degree-two
Cartan--Leray sequence: the difference comes from the group cohomology
of $\pi_1(W)\times G$. Its possible contributions are a common
multiple of $\tau^2$ and a mixed class $\tau\xi$ with
$\xi\in H^1(W;\mathbb F_2)$; its ordinary group contribution is
zero, since $H^2(B\pi_1(W))\to H^2(W)$ is injective. After one global
adjustment by $\tau^2$, \eqref{maps:eq:D.C212}--\eqref{maps:eq:D.C213} give
\[
 \widetilde\alpha|_C=\alpha|_C+\tau^2,
 \qquad \widetilde\alpha|_E=\alpha|_E+\tau\eta,
       \quad\eta\in H^1(E;\mathbb F_2).                       \tag{D.C216}\label{maps:eq:D.C216}
\]
In particular two unrelated fixed-point normalizations have not been
chosen.

\item Both complex normal directions have sign $-1$. Their normal first
Chern class is even because $c_1(X)=0$ and the tangent Chern class of
each fixed curve is even. Their equivariant Euler class is therefore
$\tau^4$. Equivariant Thom localization yields
\[
 \langle\alpha^3,[X]_2\rangle
 =\int_C\frac{(\alpha|_C+\tau^2)^3}{\tau^4}
  +\int_E\frac{(\alpha|_E+\tau\eta)^3}{\tau^4}
 =\int_C\alpha|_C=1.
\]
The elliptic contribution is zero since degree-one classes on an
oriented elliptic curve square to zero. Finally the exact primary
calculation
$\Omega^{\mathrm{Spin}}_6(K(\mathbb Z/2,2))=\mathbb Z/2$,
detected by this cube, and Zhubr's complete classification identify
$X$ with $M_{\mathrm{conic}}$; the classification and the conic's cube
were established in \ref{maps:src:D.2.8} directly from
\cite[Theorems 3, 2.10 and Section 3.12]{manZhubr75}.
This retains the two-primary refined invariant, not just ordinary
Pontryagin classes.
\end{enumerate}
\end{proof}

The exact extraction program \path{code/III/manifolds_tate_markings.py}
constructs the $SL_2(\mathbb Z)$ markings from the actual monodromies
and verifies \eqref{maps:eq:D.C29}. The independent certificate
\path{code/III/fourfolds_D2_c2_AH.py}
checks all $15$ patterns and $13$ marked rows, including $928$
semicharacter substitutions, both full integral Hermite lattices per
row, the exact conjugated monodromy, every toric character and degree
parity, and every finite normal-column parity. Its retained output is
\path{code/fourfolds_D2_c2_AH.out}; the detailed records are
\path{records/fourfolds_D2_c2_AH.json}. The symbolic identities prove
all displayed parameter formulas; no finite multiplicity or twist
bound is used in the theorem.


\begin{theorem}[the residual five- and seven-primary smooth types]\label{maps:src:D.2.36}
\StatementRecord{maps:src:D.2.36}{proved}{Historical events: Claude Fable 5.2 (Anthropic), effort high, 16 November 2026: carried as background without fault found. GPT-6 Astra (OpenAI), ChatGPT Pro mode, effort Pro, S6\_09\_PAPERS, 17 November 2026: targeted proof check. Also checked in this writing. Pre-continuation writing events: GPT-6 Astra (OpenAI), ChatGPT Pro mode, effort Pro, S6\_10\_PAPERS, 18 November 2026: targeted proof check. Continuation metadata events: GPT-6 Astra (OpenAI), ChatGPT Pro mode, effort Pro, S6\_10\_PAPERS continuation, 19 November 2026: used and rechecked in this run; status-vocabulary reconciliation; metadata reconciliation only; no new mathematical check. \emph{Pre-continuation recorded proof scope.} GPT-6 Astra (OpenAI, ChatGPT Pro mode), effort Pro, 2026-11-17 UTC; Claude Fable 5.2 (Anthropic), effort high, 16 November 2026; also checked in this writing by GPT-6 Astra (OpenAI), ChatGPT Pro mode, effort Pro, 18 November 2026}{the hypotheses, notation and definitions stated here; \cref{maps:src:D.2.26}; \cref{maps:src:D.1.11}; \cref{maps:eq:D.HY7}; \cref{maps:eq:D.OP3}; \cref{maps:eq:D.OP2}; \cref{maps:eq:D.OP4}; \cref{maps:eq:D.OP7}; \cref{maps:eq:D.OP8}; \cref{maps:src:D.1.10}; \cref{maps:eq:D.OP11}; \cref{maps:eq:D.OP10}; \cite{manZhubr75}}


Consider the $36$ marked additive-filling patterns on the $30$ initial
rows of \ref{maps:src:D.2.26} with residual $c=5$ or $7$ and no orbifold obstruction.
Use the primitive Mumford cap at every multiplicative position, the
specified minimal original or finite cap at each additive position,
and no additional regular logarithmic transformation. Allow every
positive multiple of the finite monodromy orders and every free
affine twist for which the resulting manifold is simply connected.
For $p=c$, every such manifold is spin and has complete Zhubr data
\[
 H_2=\mathbb Z/p,\qquad b_3=0,\qquad \mu=0,
 \qquad \mathfrak p\in\{0,1\}\quad\text{up to }\operatorname{Aut}(\mathbb Z/p).
 \tag{D.OP1}\label{maps:eq:D.OP1}
\]
Here $4\mathfrak p=\operatorname{PD}(p_1)$, as in \ref{maps:src:D.1.11}.
The following entirely local rule determines which value occurs.
There is exactly one zero-weight multiplicative position
whose double curve is fixed by the scalar subgroup $C_p$; it is of
type $I_p$. In its genuine Tate marking put
\[
 (\operatorname{ord}h_1,\operatorname{ord}h_G,
                  \operatorname{ord}R)=(p,r,\rho),\qquad
 w=(ar)^{-1}\pmod p.
\]
Then $\mathfrak p=0$ if and only if $w^2=1$; otherwise its canonical
value is $1$. Thus the $c=5$ patterns give exactly two oriented
diffeomorphism types and the $c=7$ patterns give exactly one.
These types are specified by the full published classifying
invariants \eqref{maps:eq:D.OP1}, rather than by homology alone.

\end{theorem}


\begin{proof}
\begin{enumerate}[label=\arabic*.,leftmargin=*]
\item \emph{Marked data and the class which must be extended.}
\ref{maps:src:D.2.26} gives $H_2=H_3=\mathbb Z/p$ and zero middle Betti number
under the stated simple-connectivity condition. Its primitive fibre
class $s$ generates $H_2$. Let $u\in H^2(X;\mathbb F_p)$ evaluate
to one on $s$. In the regular torus lattice the alternating form is
\[
 \Omega=a\,e_1^*e_2^*-b\,\delta^*\gamma^*,\qquad p\mid b,
                                      \quad p\nmid a.
 \tag{D.OP2}\label{maps:eq:D.OP2}
\]
Its values are one on $s$ and zero on the rigidified swept cycles.
The entire finite-cap normal image in \eqref{maps:eq:D.HY7} has fibre
coefficients divisible by $b/q_i$, and hence by $p$ in these
patterns. An allowed unchanged additive cap has zero fibre
coefficient, and the ordinary clutch contributes $be$.
Consequently this functional descends modulo $p$ through every
actual cap relation. This is the integral attachment calculation
which identifies its extension with $u$ for all finite orders
and twists; a finite covering transfer is not used.

\item Remove all additive caps and one ordinary clutch disc, retaining
all multiplicative Mumford caps, and denote the remaining space by
$W$. Its regular part is $V$. The actual multiplicative angular
vanishing columns generate $\ker\gamma^*$ primitively in all
$30$ rows: their Smith factors are $(1,1,1)$. The exact certificate
below uses the full marked columns, including the rank-two cusp.
Van Kampen therefore gives
\[
                  \pi_1(W)=\mathbb Z_\gamma\times F_r
 \tag{D.OP3}\label{maps:eq:D.OP3}
\]
for the free base group with the indicated caps removed. In
particular this group is torsion-free. Each primitive rank-one
Mumford cap and the one-component rank-two cap is generated in
degree two by its boundary. In degree one its sole additional
kernel is the base meridian: the angular lattice is already the
primitive image of the local monodromy minus the identity.
The unfilled ordinary disc makes these peripheral base classes
independent in the graph. Mayer--Vietoris therefore shows that
$H_2(V)\to H_2(W)$ is onto. Together with the free $H_1(W)$ of
\eqref{maps:eq:D.OP3}, this removes any universal-coefficient ambiguity in
comparing degree-two classes on $W$ from their actual values on $V$.

\item \emph{A doubled symmetric line.}
The use of $2\Omega$ makes its Appell--Humbert semicharacter
trivial modulo two, irrespective of the parity of $a,b$.
On an auxiliary holomorphic line use the multiplicative factors
\[
\begin{aligned}
 h(z,v,\zeta)&=(h_1z,h_G^a v,z^{-2a}\zeta),\\
 F(z,v,\zeta)&=(h_G^{-b}z,Rv,
                         h_G^{-ab}R^b v^{2b}\zeta),\\
 J(z,v,\zeta)&=(z^{-1},v^{-1},z^{-2a}\zeta).
\end{aligned}                                                    \tag{D.OP4}\label{maps:eq:D.OP4}
\]
Direct substitution gives $hF=Fh$, $J^2=1$, $JhJ=h^{-1}$
and $JFJ=F^{-1}$. The last identity uses precisely the square
of $h_G^{-ab}R^b$. The alternating Chern form is $2\Omega$.
These symmetric line bundles are rigidified at the origin.
Normalized identifications on overlap markings therefore patch:
an automorphism of a rigidified holomorphic line bundle on a
compact torus is the identity. Comparison with the flat real
prequantum family for $2\Omega$ has globally two-torsion Chern
class. Indeed its fibrewise square is canonically trivial after
rigidification, and the unique normalized trivializations patch.
Reduction modulo the odd prime $p$ kills this difference. Thus
half its Chern class modulo $p$ is exactly \eqref{maps:eq:D.OP2}, including
the zero values on every swept cycle.

\item We give the required extensions and scalar characters explicitly.
At the rank-two position write the two actual valuation periods as
\[
 H=(n,ar),\quad F=(-br,\rho),\qquad n\rho+ab r^2=1.
\]
An integral support function is
\[
 \psi(x,y)=-a\rho x^2-2ab rxy+bn y^2+ax.
 \tag{D.OP5}\label{maps:eq:D.OP5}
\]
Its differences under $H,F$ are respectively
$-2ax$ and $2by+b\rho-ab r$, the valuations of the
two factors in \eqref{maps:eq:D.OP4}. The same unimodular fan consequently
extends the line over the cusp. Its vertical slopes
$-2ab rx+bn(2y+1)$ are divisible by $p$.
The scalar $C_p$ character is therefore zero on its fixed
rational double curve $C$. The wall expression, using the two
opposite vertices $(1,0),(-1,1)$ and edge endpoints
$(0,0),(0,1)$, is
\[
 \psi(1,0)+\psi(-1,1)-\psi(0,0)-\psi(0,1)
                         =-2a(\rho-br).
 \tag{D.OP6}\label{maps:eq:D.OP6}
\]
The divisor convention can reverse the sign of this wall
expression; in either convention its half is nonzero modulo
$p$, since $\rho\equiv n^{-1}\pmod p$ and $p\nmid a$.
Thus the class initially normalized by $u(s)=1$ has
$u[C]\ne0$. For the remainder of the proof, rescale $u$
uniquely so that its evaluation on the complex-oriented
curve $C$ is $1$, and scale the auxiliary Chern class by
the same nonzero field element. This fixes the actual
cohomology-generator convention without a support-function
sign ambiguity. In particular $\alpha_C=1$ below.

\item At a rank-one position write its two periods as
$H=\alpha u,F=\beta u$, with $u=(u_1,u_2)$ primitive and
$\gcd(\alpha,\beta)=1$. Thus $n=\alpha u_1$,
$ar=\alpha u_2$, $-br=\beta u_1$, and $\rho=\beta u_2$.
The identity $a\beta u_1=-b\alpha u_2$ implies
$\alpha\mid au_1$. An explicit integral support is
\[
           \psi_1(j)=-\frac{au_1}{\alpha}\,j(j-\alpha).
 \tag{D.OP5a}\label{maps:eq:D.OP5a}
\]
Its differences under $\alpha,\beta$ are exactly
$-2au_1j$ and $2bu_2j+b\rho-ab r$.
The elliptic combination $\alpha F-\beta H$ has zero
exponent in the degenerating coordinate, by the same identity.
These facts extend both specified factors in \eqref{maps:eq:D.OP4} over the
primitive fan and retain their actual boundary marking.
At the unique $I_p$ position fixed by $C_p$, the short period
is already the primitive vector $(p,ar)$. Choose
integers $A,B$ with $Ap+B ar=1$. In the monomial coordinates
with first row $(A,B)$ and second row $(-ar,p)$, the scalar
action has normal weight $B=w$ and elliptic weight zero
modulo $p$. Its support function can be taken to be
$-ap\,j(j-1)$. All its slopes are divisible by $p$.
The scalar character of the auxiliary line on this elliptic
double curve $E$ is therefore also zero.

\item These are equivariant extensions. On the covering coordinates
the scalar subgroup lifts by $(z,v,\zeta)\mapsto
(z,\lambda v,\zeta)$ for $\lambda^p=1$; it commutes with
\eqref{maps:eq:D.OP4} because $p\mid b$. In the prequantum comparison the
same lift is preserved by every linear marking: each fixes
$\delta$ and $2\Omega$, and the semicharacter is trivial.
The symmetric two-torsion comparison has a unique compatible
odd-order lift. Holomorphic units on the local base discs
have logarithms and contribute no scalar character or wall
degree. Thus the above local lifts patch on $W$.

\item \emph{The fixed set and the global equivariant comparison.}
The primitive rank-one period at a multiplicative cap is the
column $(u_1,u_2)$ obtained by taking the Hermite basis of its
two period columns. Scalar isotropy along its double curve
has order $|u_1|$. The retained exact data show that precisely
one such order is divisible by $p$, and its type is $I_p$.
The rank-two fan has just the vertical rational fixed curve
$C$; its other primitive edge directions are not parallel to
$\delta$ modulo $p$. There are no further fixed strata.
An allowed original additive cap is a product with the free
scalar circle. At a finite cap, primitivity of the actual gamma
numerator in a simply connected member excludes scalar fixed
points: equality of a scalar translate with the $j$th affine
power would force $m\mid j$ from the gamma coordinate, and
then require a nontrivial scalar translation to fix a torus
point. This also covers orders divisible by $p$.
Consequently
\[
                         X^{C_p}=C\sqcup E.
 \tag{D.OP7}\label{maps:eq:D.OP7}
\]

\item The class $u$ admits a Borel equivariant lift. The obstruction
with $H^1(X;\mathbb F_p)$ vanishes by simple connectivity;
the remaining differential into $H^3(BC_p;\mathbb F_p)$
vanishes because any point of $C$ splits the base map.
The whole scalar circle fixes that point and acts trivially
on the based fundamental group of $W$. Hence
\[
 \pi_1(W_{C_p})=(\mathbb Z\times F_r)\times C_p.
\]
Every fixed point determines its same, unique order-$p$
subgroup, by torsion-freeness of the first factor.
The difference between the chosen lift of $u$ and the correspondingly scaled half of
the equivariant Chern class of the auxiliary line has zero
ordinary restriction on $W$. The universal-cover low-degree
sequence places this difference in group cohomology.
Its restrictions to the fixed-point copies of $BC_p$ are
therefore equal. Adjust the lift once by a base class.
Writing $t\in H^2(BC_p;\mathbb F_p)$ and
$\eta\in H^1(BC_p;\mathbb F_p)$, we obtain
\[
 \widetilde u|_C=\alpha_C[C]^*,\qquad
 \widetilde u|_E=\alpha_E[E]^*+\eta\zeta,
                  \quad\zeta\in H^1(E;\mathbb F_p).
 \tag{D.OP8}\label{maps:eq:D.OP8}
\]
In particular both restrictions have zero coefficient of $t$.

\item \emph{All the classifying invariants.}
Since $H^2(X;\mathbb Z)=0$, the complex first Chern class
vanishes and $X$ is spin. The two complex normal weights
at $C$ are $1,-1$, and those at $E$ are $w,-w$.
One can also check the sums from equivariant $c_1$:
it comes from $H^2(BC_p;\mathbb Z)$, and is zero at $C$,
hence zero everywhere. The equivariant Euler classes
have leading terms $-t^2$ and $-w^2t^2$.
Inverting $t$ in the equivariant Thom isomorphism gives
localization on \eqref{maps:eq:D.OP7}. Applying it to $u$, whose integral
has negative degree and is zero, gives
\[
                      \alpha_C+\alpha_E/w^2=0.
 \tag{D.OP9}\label{maps:eq:D.OP9}
\]
No mixed term in \eqref{maps:eq:D.OP8} contributes: it has ordinary
degree one and its product with an ordinary degree-two
Euler correction has degree three on a curve.

\item For each normal bundle with roots $wt+x_1,-wt+x_2$,
the square of their sum is zero on the fixed curve.
The tangent line's square is zero as well. Therefore its
equivariant Pontryagin class divided by the normal Euler
class is exactly $-2$. It follows that
\[
 \langle u\,p_1(X),[X]\rangle
   =-2(\alpha_C+\alpha_E)
   =-2(1-w^2).
 \tag{D.OP10}\label{maps:eq:D.OP10}
\]
Moreover the cube of each restriction in \eqref{maps:eq:D.OP8} is zero:
ordinary degree-two classes square to zero on a curve,
and $\eta^2=0$ at an odd prime. Thus
\[
                        \langle u^3,[X]\rangle=0.
 \tag{D.OP11}\label{maps:eq:D.OP11}
\]
Lemma \ref{maps:src:D.1.10} with $r=0,d=p$ identifies
$H_6(K(\mathbb Z/p,2);\mathbb Z)=\mathbb Z/p$
and proves that \eqref{maps:eq:D.OP11} detects its entire fundamental
class invariant. Consequently $\mu=0$.
Multiplication by four is invertible on $H_2$, so
\eqref{maps:eq:D.OP10} computes the entire refined Pontryagin class.
Because $\alpha_C\ne0$, it is zero exactly when $w^2=1$.
Every nonzero element of $\mathbb Z/p$ is carried to one
by an automorphism; these automorphisms leave $\mu=0$.
Zhubr's Theorem~1, with the definitions on pp.~794--795
\cite{manZhubr75}, now gives precisely \eqref{maps:eq:D.OP1} and the
claimed oriented diffeomorphism classification. Existence
comes from the actual constructed members, so no separate
realization assumption is needed.

\item The local weights in the retained table include both
$w^2=1$ and $w^2\ne1$ for $p=5$, and only $w^2\ne1$
for $p=7$. The code checks all $30$ actual marked rows;
the $36$ additive patterns and all allowed higher orders
use the same local data. This proves the final counts.
\end{enumerate}
\end{proof}

The complete Tate markings are reconstructed by
\path{code/III/manifolds_tate_markings.py}, and
\path{code/III/manifolds_odd_prime_AH.py} checks the doubled
line identities, support functions, primitive angular
lattices, scalar isotropy, unique fixed elliptic curves,
and the exact characteristic numbers for every listed row.
The refined Pontryagin coefficient in the curve-normalized
generator is exactly $(w^2-1)/2$ modulo $p$; its nonzero
orbit is subsequently normalized to one in \eqref{maps:eq:D.OP1}.
Its JSON output records both the signed support expression
and these curve-normalized characteristic coefficients.
No original
multiplicative replacement, extra regular log, or
higher initial weight is included in this theorem.


\begin{theorem}[a uniform weighted algorithm on all natural profiles]\label{maps:src:D.2.37}
\StatementRecord{maps:src:D.2.37}{proved}{Historical events: Claude Fable 5.2 (Anthropic), effort high, 16 November 2026: carried as background without fault found. GPT-6 Astra (OpenAI), ChatGPT Pro mode, effort Pro, S6\_09\_PAPERS, 17 November 2026: targeted proof check. Also checked in this writing. Pre-continuation writing events: GPT-6 Astra (OpenAI), ChatGPT Pro mode, effort Pro, S6\_10\_PAPERS, 18 November 2026: targeted proof check. Continuation metadata events: GPT-6 Astra (OpenAI), ChatGPT Pro mode, effort Pro, S6\_10\_PAPERS continuation, 19 November 2026: used and rechecked in this run; status-vocabulary reconciliation; metadata reconciliation only; no new mathematical check. \emph{Pre-continuation recorded proof scope.} GPT-6 Astra (OpenAI, ChatGPT Pro mode), effort Pro, 2026-11-17 UTC; Claude Fable 5.2 (Anthropic), effort high, 16 November 2026; also checked in this writing by GPT-6 Astra (OpenAI), ChatGPT Pro mode, effort Pro, 18 November 2026}{the hypotheses, notation and definitions stated here; \cref{maps:src:D.2.22}; \cref{maps:src:D.2.33}; \cref{maps:src:D.2.32}; \cref{maps:src:D.2.23}; \cref{maps:eq:D.WP10}; \cref{maps:src:D.2.29}; \cref{maps:src:D.2.34}; \cref{maps:eq:D.WP11}; \cref{maps:eq:D.WP12}}


Choose any of the $172$ natural initial rows of \ref{maps:src:D.2.22},
with $P=aG,Q=bG$, and any integer $d>0$ satisfying
\[
                         \gcd(a,d)=1.                    \tag{D.WP1}\label{maps:eq:D.WP1}
\]
Iterate its actual contracting translation $d$ times.
Thus the new translation is $dbG$ and its linearization
is the norm of the original one under the first $d$
translates. Keep the same unimodular Mumford fans,
quotiented by the new period lattice.

Every original/Mumford choice represented by \ref{maps:src:D.2.33},
every free finite-good additive logarithm, every number
of free regular logarithms and every integer peripheral
clutch is covered by the following exact integer
algorithm. Star caps are retained original. The
$495$ initial presentation patterns provide the finite
input; there is no upper bound on $d$, logarithmic
orders or affine coordinates.

Write $A_i$ for the recorded four-dimensional
monodromies in the basis $(e_1,e_2,\delta,\gamma)$.
Replace them by
\[
 A_i(d)=S_d^{-1}A_iS_d,\qquad
 S_d=\operatorname{diag}(1,1,1,d).
                                                               \tag{D.WP2}\label{maps:eq:D.WP2}
\]
Equivalently multiply their first three gamma-column
entries by $d$; every other entry is unchanged.
The primitive common alternating functional is
\[
 \omega_d=a e_1^*e_2^*-db\,\delta^*\gamma^*.                \tag{D.WP3}\label{maps:eq:D.WP3}
\]
Insert identity matrices for regular logarithms.
An additional ordinary order-one logarithm with power
$-e\gamma$ realizes the clutch integer $e$.

Let $N$ be the number of sites, with
$A_0(d)\cdots A_{N-1}(d)=I$, and compute the integral
kernel
\[
 K=\ker\bigl[(A_0(d)-I)\ \cdots\ (A_{N-2}(d)-I)\bigr].
                                                               \tag{D.WP4}\label{maps:eq:D.WP4}
\]
The target of the degree-two presentation is
$\mathbb Zs\oplus K$. For a fixed vector at site $i$,
let $J_i$ denote its usual oriented peripheral
insertion in $K$: the inclusion in block $i$ for
$i<N-1$, and, at the last site, the blocks
$-A_{j+1}(d)\cdots A_{N-2}(d)$.
All these maps are evaluated in a primitive integral
basis of \eqref{maps:eq:D.WP4}.

At a finite-good site of linear order $h$, compute
a primitive fixed basis $F:\mathbb Z^2\to\Lambda$
and the norm $N_h=I+\cdots+A^{h-1}$. If $H$ is an
integral basis matrix for the image of
$F^{-1}N_h:\Lambda\to\mathbb Z^2$, the projected lattice
is $F(H/h)\mathbb Z^2$. For any invariant power
$v=Fz$ and order $m=hk$, put
\[
 \begin{gathered}
 w=hH^{-1}z\in\mathbb Z^2,\qquad t=\gcd(m,w_1,w_2),\\
 \text{the action is free iff}\quad
 \gcd(m,v_1,v_2,v_3,v_4)=1,\quad h\mid m/t.
 \end{gathered}
                                                               \tag{D.WP5}\label{maps:eq:D.WP5}
\]
For every free such logarithm include the five columns
\[
 \bigl(-\omega_d(\lambda,v),\,J_iN_m\lambda\bigr)
       \quad(\lambda=e_1,e_2,\delta,\gamma),\qquad
                         (0,J_iv),                       \tag{D.WP6}\label{maps:eq:D.WP6}
\]
where $N_m=(m/h)N_h$. At a regular logarithm use the
same formula with $A=I,N_m=mI$ and the sole freeness
condition $\gcd(m,v_1,v_2,v_3,v_4)=1$.

At every unchanged original or Mumford cap include
the full marked fixed-normal columns
\[
                       (0,J_iF_i),\qquad
                   F_i\text{ a basis of }\ker(A_i(d)-I).
                                                               \tag{D.WP7}\label{maps:eq:D.WP7}
\]
For a nonzero-Chern original also include
$(dh_i,0)$, where $h_i$ is its recorded
$b/q_i^c$ or $b/n_i$ from \ref{maps:src:D.2.32}--\ref{maps:src:D.2.33}.
These are the entire global images of the boundary
kernels, including any torsion in a boundary.
Let $B$ be the resulting integer matrix and
$T=\operatorname{Tor}\operatorname{coker}B$.

For the component correction define $\beta$ as follows.
A zero-free rank-one site with $P$ narrow contributes
$\gcd(d,n_i)-1$; if $P$ is non-narrow, its marked
Mumford cap contributes zero. A zero-degree original
additive cap contributes $\rho_i(d)-1$, where
$\rho_i(d)$ is the number of component orbits under
$Q^d$, computed by raising the recorded permutation
to the $d$th power. A nonzero-Chern original contributes
its unchanged $s_i-2$ or $n_i-2$.
Let $L$ count these last originals. The weighted
rank-two zero has $d$ normalized components and is
not included in $\beta$.

The full fundamental group is the Seifert presentation
of \ref{maps:src:D.2.23} with the actual gamma coordinates of the
power vectors, the displayed orders and clutch.
Write its abelianization as $\mathbb Z^r\oplus A_1$.
Then
\[
 b_2=
 \begin{cases}
 d-1+\beta+2r,&L=0,\\
 d-1+\beta+L+r,&L>0,
 \end{cases}
 \qquad
 b_3=
 \begin{cases}
 2\beta+2r,&L=0,\\
 2\beta+2L,&L>0.
 \end{cases}                                               \tag{D.WP8}\label{maps:eq:D.WP8}
\]
The complete integral groups are
\[
\begin{array}{c|c}
j&H_j(X;\mathbb Z)\\\hline
0,6&\mathbb Z\\
1&\mathbb Z^r\oplus A_1\\
2&\mathbb Z^{b_2}\oplus T\\
3&\mathbb Z^{b_3}\oplus T\\
4&\mathbb Z^{b_2}\oplus A_1\\
5&\mathbb Z^r .
\end{array}                                                \tag{D.WP9}\label{maps:eq:D.WP9}
\]
In particular $\chi(X)=2d$. The executable
\path{code/III/manifolds_weighted_profiles.py} implements
all steps, including actual freeness, the full group
presentation and both Smith calculations.

\end{theorem}


\begin{proof}
\begin{enumerate}[label=\arabic*.,leftmargin=*]
\item \emph{The geometric meaning of the new weights.}
Iteration of the original lifted translation replaces
its line multiplier by the product of its first $d$
translates. If its initial zero is one component of
$I_n$, the new zero is its orbit sum, with repetitions.
It need not be a whole fibre. This is the prescribed
norm linearization, so no unsupported line-bundle
modification is being introduced.
Outside the finite replacements the new space is
the connected gamma cover defined by $S_d$ of the
old original or Mumford model. It follows directly
that the good-reduction monodromies are \eqref{maps:eq:D.WP2};
the same construction on their root discs permits
all the new free finite quotients.

\item At the unique rank-two zero the old valuation
lattice was the full short lattice. Multiplying
one long period by $d$ gives index $d$. The same
smooth triangular fan therefore has $d$ copies
of the old normalized $dP_6$ in its quotient.
At a zero-free rank-one site, write the old
monodromy as $I+\alpha\beta^*$ with $\alpha$
primitive and $\beta^*$ primitive. The cover
replaces $\beta^*$ by $\beta^*S_d$.
The exact factorization certificate retains all $436$
zero-free rank-one matrices as $\alpha\beta^*$, with
both vectors primitive. If $P$ is narrow, the content
of the first three entries of $\beta^*$ is exactly
$n_i$, and its gamma entry is prime to $n_i$. Thus
the new content is $\gcd(d,n_i)$. At every non-narrow
$P$ site that vertical content is exactly one, so
its new content is one for every $d$. These actual
vectors and entries are recorded in
\path{records/manifolds_weighted_rankone_certificates.json}.
Thus the respective component counts are precisely
those in the statement. This also follows directly
by lifting the component cycle through the gamma
cover. Zero-degree original multiplicative caps
have the same component count and the same global
kernel image; no relative diffeomorphism is asserted.

\item An original zero-degree additive cap is
$S^1_\delta$ times the mapping torus of the Kodaira
tree under $Q^d$. It has $\rho_i(d)$ fundamental
component cycles. At a nonzero-Chern original, $Q$
was already narrow and remains narrow; its
primitive Chern vector is unchanged. Its component
count and its principal-torus description from
\ref{maps:src:D.2.32} or \ref{maps:src:D.2.33} are therefore unchanged.
All these caps have torsion-free top homology and
Euler number zero. The rank-two cusp has Euler
number $2d$ by its normalization complex.

\item \emph{Why the regular integral lattice is free in the stated sector.}
For each original row, \ref{maps:src:D.2.23} gives
\[
 \Lambda_{\{A_i\}}=\mathbb Z\gamma,\qquad
 (\Lambda^2\Lambda)_{\{A_i\}}=\mathbb Zs,
\]
with primitive functional \eqref{maps:eq:D.WP3} at $d=1$.
There is a finite exact certificate for all other
weights. Set $d=0$ in the polynomial matrices
\eqref{maps:eq:D.WP2}, meaning that their first three gamma
entries are zero. For every one of the $172$ rows
the common degree-one relation matrix still has
rank three, the degree-two matrix has rank five,
and every prime in their nonunit Smith factors
divides $a$. The complete constant matrices and
Smith factors are retained in
\path{records/manifolds_weighted_zero_certificates.json}.

\item Here is why this finite certificate proves the
unbounded claim. At a prime not dividing $d$,
$S_d$ is an integral local isomorphism, so the
relations have the original saturated ranks.
At a prime dividing $d$, condition \eqref{maps:eq:D.WP1} makes
$a$ a unit, and reduction modulo that prime is
the just-certified $d=0$ matrix, still of full
rank three or five. Hence no prime divides a
nonunit Smith factor at the new weight.
The two common coinvariants are therefore free
of ranks one and one for every allowed $d$.
Moreover \eqref{maps:eq:D.WP3} is primitive since
$\gcd(a,db)=1$, so it is their actual integral
degree-two functional.

\item The graph Serre sequence consequently gives
\[
 H_1(V)=\mathbb Z\gamma\oplus\mathbb Z^{N-1},
 \qquad H_2(V)=\mathbb Zs\oplus K.                         \tag{D.WP10}\label{maps:eq:D.WP10}
\]
The splitting used in \eqref{maps:eq:D.WP10} is marked. On the
original regular piece the integral fibre class
has value zero on every fixed-loop peripheral
sweep, by the unrolled comparison and the single
diagonal relation of \ref{maps:src:D.2.23}--\ref{maps:src:D.2.26}. Pull it back
through the actual gamma cover. It restricts to
$\omega_d$ and still has value zero on every
new fixed sweep: that sweep maps to the old
sweep of $S_d\lambda$.
Primitivity of $\omega_d$ makes this a genuine
integral splitting. This is a pullback of an
already marked class, not an assumption that
an invariant differential form descends integrally.

\item \emph{All local boundary-kernel images.}
For a finite-good matrix, all nonidentity linear
powers have the same fixed real two-plane.
Its projected lattice is
$P\Lambda=(N_h/h)\Lambda=F(H/h)\mathbb Z^2$.
The translation $v/m$ on that quotient has exact
order $m/t$. A nonidentity linear power can have
a fixed point exactly when its exponent is
divisible by $m/t$. Requiring all such exponents
to have identity linear part is $h\mid m/t$.
For the identity-linear powers one must additionally
have the full translation of order $m$, which is
the primitive-vector condition in \eqref{maps:eq:D.WP5}.
This proves the stated freeness criterion for
arbitrary coordinates and orders.

\item The full normal Thom image is generated by the
four covering-torus loops and the affine meridian.
The cut of their swept normal tori into the regular
fibre is respectively $N_m\lambda$ and $v$.
Integration of the actual pulled-back marked
class from \eqref{maps:eq:D.WP10} gives respectively
$-\omega_d(\lambda,v)$ and zero, exactly as in
\ref{maps:src:D.2.29}. These are complete integral coordinates
in the free target \eqref{maps:eq:D.WP10}, so a local torsion
ambiguity cannot be lost by this computation.
In particular no claim of integral form descent
on the finite quotient is needed.
This proves \eqref{maps:eq:D.WP6}, including finite cap and
normal-line torsion.

\item For a Mumford or zero-degree original cap, every
new fixed normal torus bounds. Its image under
the covering is an old fixed normal torus; the
old bounding section-disc sweep lifts because
its loop is in the covering lattice. Equivalently
use the same toric or Neron-logarithm section
disc and its actual invariant circle actions.
These are all the normal projections, since
the Wang projection takes values in the full
fixed lattice.
After subtracting them, any further kernel
element comes from a fibre class. Rationally
its image is zero exactly on the monodromy-image
subspace. At the rank-two zero, the two primitive
angular normal tori exhaust the boundary kernel:
$H_2(\partial N)=\mathbb Z^6$ and
$H^3(N)=\mathbb Z^2$ by \ref{maps:src:D.2.34}. The regular-fibre
coinvariant of rank four therefore injects. At a
rank-one cap the nodal cycle times $T^2$ supplies
four independent fibre images (whole fibre, two
mixed surviving-circle tori, and transverse area).
At an original zero-degree additive cap the
Kodaira-tree permutation model supplies the whole
fibre and the transverse area, the full two
dimensional coinvariant. These explicit models
prove the required ranks for every weight.
The functional $\omega_d$ annihilates that
subspace. Since the global fibre coinvariant
in \eqref{maps:eq:D.WP10} is free, every such fibre-kernel
class has zero global image, even when it is
torsion in the boundary. This proves the entire
global image \eqref{maps:eq:D.WP7}.

\item For a nonzero-Chern original, the marked logarithm
vectors scale by $d$. The exact area-kernel
calculations of \ref{maps:src:D.2.32}--\ref{maps:src:D.2.33} therefore multiply
their fibre values by $d$, giving $dh_i\mathbb Zs$.
Their complete fixed normal projections still
bound, and the third multiplicative correction
is still annihilated by $\omega_d$.
This proves the remaining column, integrally.

\item \emph{The relative sequence and the free component contribution.}
First attach only the finite and regular logarithmic
caps to $V$, and call the result $E$. There is at
least one unfilled cap, the rank-two zero.
The relative degree-two group of each attached
disc bundle is $\mathbb Z$, whose boundary has
its nonzero multiplicity in its independent base
meridian coordinate. These relative maps into
$H_1(V)$ are injective. Thus $H_2(E)$ is exactly
the quotient of \eqref{maps:eq:D.WP10} by their full Thom
images, including when the cap has torsion.

\item Now attach all the remaining original and Mumford
caps $N_j$. Their relative degree-two groups are
\[
 H_2(N_j,\partial N_j)=H^4(N_j),
\]
free on their fundamental component cycles.
The degree-three boundary images are exactly
the kernels already computed above. Consequently
the pair sequence has the form
\[
 0\longrightarrow\operatorname{coker}B
 \longrightarrow H_2(X)
 \longrightarrow
 \ker\!\left\{\bigoplus_j H^4(N_j)\to H_1(E)\right\}
 \longrightarrow0.                                      \tag{D.WP11}\label{maps:eq:D.WP11}
\]
Its last group is a subgroup of a free group,
so the sequence splits as a sequence of abelian
groups. This avoids all potentially nonprimitive
degree-one boundary splittings.

\item If there are $s$ unfilled caps, $H_1(E)$ has free
rank $s$: the finite/regular power relations
remove their independent base directions.
The final rank is $r$, so the map in \eqref{maps:eq:D.WP11}
has rank $s-r$. Its domain has rank equal to
the sum of the component counts. Subtracting
one per cap gives
\[
 \operatorname{rank}\ker
   =d-1+\beta+L+r.                                        \tag{D.WP12}\label{maps:eq:D.WP12}
\]
This is an exact rank calculation in a free
relative group; torsion meridians can change
its index but cannot change its freeness or
introduce an extension in \eqref{maps:eq:D.WP11}.

\item Over the rationals \eqref{maps:eq:D.WP2} is conjugate to the
original monodromy. Thus the fixed-sweep graph
has the same single diagonal relation as in
\ref{maps:src:D.2.23}--\ref{maps:src:D.2.26}. With no nonzero-Chern original,
the remaining scalar equation is the actual
Seifert Euler equation, so
$\operatorname{corank}_{\mathbb Q}B=r$.
If $L>0$, a nonzero fibre multiple is killed,
and then every fixed sweep is killed rationally;
the corank is zero. Combining this with
\eqref{maps:eq:D.WP12} proves $b_2$ and its full torsion $T$.

\item The fundamental-group argument is also unchanged
integrally. The primitive angular vanishing
vectors are vertical, so they survive in the
cover without a multiple. Restrictions of
$A_i(d)$ to the vertical lattice are the
original restrictions. Their normal closures
therefore kill the same entire vertical
lattice as in \ref{maps:src:D.2.23}, including its original
star thimble when required. Each unchanged
cap has an actual transverse origin disc, so
adds no gamma twist beyond the recorded
peripheral clutch. The finite/regular power
relations then give exactly the full Seifert
group, not merely its abelianization.
Finally $\chi=2d$ and integral duality give
the remainder of \eqref{maps:eq:D.WP8}--\eqref{maps:eq:D.WP9}.
\end{enumerate}
\end{proof}

The program verifies all $172$ weight-zero
prime-support certificates and computes $1626$
full weighted graph presentations on the $495$
patterns, of which $1131$ have weight greater
than one. The all-parameter proof is the
prime-local argument and marked relative
calculation above, not a finite weight search.
Weights sharing a factor with $a$, additional
star logarithmic Mumford choices and arbitrary
initial pairs not obtained by these translation
iterates remain outside this theorem. Its
homology algorithm makes no new smooth
classification assertion.

The shared-even-factor case excluded from the stated formula is handled by the complete integral algorithm of \ref{maps:src:D.2.41}.

\begin{theorem}[the mixed ten- and fourteen-torsion smooth types]\label{maps:src:D.2.38}
\StatementRecord{maps:src:D.2.38}{proved}{Historical events: Claude Fable 5.2 (Anthropic), effort high, 16 November 2026: carried as background without fault found. GPT-6 Astra (OpenAI), ChatGPT Pro mode, effort Pro, S6\_09\_PAPERS, 17 November 2026: targeted proof check. Also checked in this writing. Pre-continuation writing events: GPT-6 Astra (OpenAI), ChatGPT Pro mode, effort Pro, S6\_10\_PAPERS, 18 November 2026: targeted proof check. Continuation metadata events: GPT-6 Astra (OpenAI), ChatGPT Pro mode, effort Pro, S6\_10\_PAPERS continuation, 19 November 2026: used and rechecked in this run; status-vocabulary reconciliation; metadata reconciliation only; no new mathematical check. \emph{Pre-continuation recorded proof scope.} GPT-6 Astra (OpenAI, ChatGPT Pro mode), effort Pro, 2026-11-17 UTC; Claude Fable 5.2 (Anthropic), effort high, 16 November 2026; also checked in this writing by GPT-6 Astra (OpenAI), ChatGPT Pro mode, effort Pro, 18 November 2026}{the hypotheses, notation and definitions stated here; \cref{maps:src:D.2.36}; \cref{maps:src:D.2.26}; \cref{maps:eq:D.C26}; \cref{maps:eq:D.C27}; \cref{maps:eq:D.MP3}; \cref{maps:eq:D.MP5}; \cref{maps:eq:D.C216}; \cref{maps:src:D.2.35}; \cref{maps:src:D.2.8}; \cref{maps:eq:D.OP4}; \cref{maps:eq:D.MP4}; \cref{maps:eq:D.OP8}; \cref{maps:src:D.1.10}; \cite[p.796]{manZhubr75}; \cref{maps:eq:D.MP9}; \cite[p.795]{manZhubr75}; \cref{maps:eq:D.MP7}}


Take one of the twelve minimal patterns in the following table,
with the primitive Mumford filling at every multiplicative position,
the indicated additive choice, and no additional regular logarithmic
transformation. Retain the ordinary clutch $e\gamma$ and allow every
finite multiplicity and affine twist for which the output $X$ is
simply connected. Zero positions use the zero-based source indexing.
Let $Z_{p,\epsilon}$ denote the oriented smooth type of \ref{maps:src:D.2.36} with
$H_2=\mathbb Z/p$, $b_3=0$, $\mu=0$, and refined Pontryagin
orbit $\mathfrak p=\epsilon$, for $\epsilon=0$ or $1$ when that
type occurs there. Then
\[
                  X\cong M_{\mathrm{conic}}\# Z_{p,\epsilon}
                     \quad\hbox{orientation-preservingly},
                                                        \tag{D.MP1}\label{maps:eq:D.MP1}
\]
where the parameters are
\[
\begin{array}{c|r|r|l|l|r|r}
\text{OS/profile}&a&b&\text{zero positions}&\text{additive choice}&p&\epsilon\\\hline
56&1&10&1&\text{none}&5&0\\
56&3&10&3,4&\text{none}&5&1\\
56/II&1&10&3&II\text{ original or log}&5&0\\
56/IV&1&30&3&IV\text{ log}&5&0\\
56/IV&3&10&3&IV\text{ original}&5&1\\
47&1&14&2,3,4&\text{none}&7&1\\
47/II&1&14&2&II\text{ original or log}&7&1.
\end{array}                                               \tag{D.MP2}\label{maps:eq:D.MP2}
\]
There are seven $c=10$ patterns and five $c=14$ patterns on ten
distinct marked initial rows. Thus they give exactly two oriented
smooth types for $c=10$ and one for $c=14$.

\end{theorem}


\begin{proof}
\begin{enumerate}[label=\arabic*.,leftmargin=*]
\item \emph{The actual classes and local extensions.}
\ref{maps:src:D.2.26} gives $H_2(X;\mathbb Z)=H_3(X;\mathbb Z)=\mathbb Z/(2p)$,
with $b_3=0$. In particular $H^2(X;\mathbb Z)=0$, so the complex
first Chern class is zero and $X$ is spin. Use the lattice basis
$(u,w,\delta,\gamma)$ and the primitive alternating form
\[
             \Omega=a\,u^*w^*+b\,\gamma^*\delta^*.
                                                        \tag{D.MP3}\label{maps:eq:D.MP3}
\]
Here $a$ is odd, $b/2$ is odd, and $p\mid b$. The coefficient
$b/q_i$ in every finite-cap normal relation is divisible by $2p$:
the only possibilities are $b/q_i=10$ or $14$, with the $IV$
logarithmic cap using $b/q_i=30/3=10$. The other coefficient is
a multiple of $b$. The allowed original additive caps have zero
fibre coefficient, and the ordinary clutch has coefficient $be$.
Thus the pure fibre functional, zero on the rigidified sweeps,
descends both modulo two and modulo $p$ through the actual
attachments for every allowed finite order and twist.

\item Let $W$ retain all multiplicative Mumford caps and omit the additive
caps and one ordinary disc, and let $V$ be its regular part. The
exact marked angular lattices in all ten rows have Hermite basis
$(u,w,\delta)$; so do the full images of $A_i-I$. Hence
$\pi_1(W)=\mathbb Z_\gamma\times F_{r-1}$, where $r$ is the number
of unfilled discs. The boundary-to-cap degree-two maps are onto.
As in the rank calculation \eqref{maps:eq:D.C26}, the independent peripheral
base classes, with the ordinary disc still unfilled, give
$H_2(V)\twoheadrightarrow H_2(W)$. Both first homology groups are
free. Therefore restriction from $W$ to $V$ is injective in degree
two with integral, mod-two, or mod-$p$ coefficients. This proves
the needed comparison with the globally descended classes.

\item Use the unsquared symmetric factors \eqref{maps:eq:D.C27} for $\Omega$.
All their exponents are integral since $b$ is even. Their
semicharacter $q=a(u+w+uw)$ is preserved by every actual
monodromy, as the exact certificate checks on all sixteen
vectors at each site. Thus the rigidified factors patch.
Their comparison with the marked prequantum family has globally
two-torsion Chern class; $H^2(V;\mathbb Z)$ is free, so that class
is zero. The resulting line has exactly \eqref{maps:eq:D.MP3} on $V$,
including its zero swept values.

\item For completeness the local extension uses the raw Tate periods
\[
 H=(n,ar),\qquad F=(-br,\rho).
\]
At the rank-two zero they have determinant one, and
\[
 \phi(x,y)=\frac{-a\rho x^2-2ab rxy+bn y^2+ax}{2}
                                                        \tag{D.MP4}\label{maps:eq:D.MP4}
\]
has period differences $-ax$ and $by+(b\rho-ab r)/2$.
Here $n,\rho$ are odd. Therefore $\phi$ is integral on every
lattice point. Its vertical slope has odd constant term
$bn/2$, whereas every coefficient of that slope is divisible
by $p$. The line's characters on the rational fixed double
curve $C$ are consequently one for the scalar involution and
zero for $C_p$. The wall degree is odd and a unit modulo $p$,
since its signed value is $-a(\rho-br)$.

\item At a rank-one site take its actual primitive period $u_0=(u_1,u_2)$
and write $H=\alpha u_0,F=\beta u_0$ with coprime
$\alpha,\beta$. The integral support is
\[
       \phi_1(j)=-\frac{a u_1}{2\alpha}j(j-\alpha).
                                                        \tag{D.MP5}\label{maps:eq:D.MP5}
\]
The extracted period identities give $\alpha\mid a u_1$;
either $\alpha$ is odd or $a u_1/\alpha$ is even. Hence \eqref{maps:eq:D.MP5}
is integral. Its two differences are precisely
$-a u_1j$ and $b u_2j+(b\rho-ab r)/2$. The elliptic
zero-period combination has zero exponent in the degenerating
coordinate. These are extensions of both actual factors, with
their given marking, over the primitive toric fans.

\item For the scalar involution there is exactly one further fixed
curve $E_2$, at an $I_2$ site. Its short period $(2,ar)$ is
primitive and \eqref{maps:eq:D.MP5} becomes $-a j(j-1)$, with even slopes.
For $C_p$ there is exactly one further fixed curve $E_p$,
at an $I_p$ site. Its short period $(p,ar)$ is primitive and
the doubled support becomes $-ap j(j-1)$, with slopes
divisible by $p$. The corresponding normal weight is
\[
                         w=(ar)^{-1}\pmod p.             \tag{D.MP6}\label{maps:eq:D.MP6}
\]
There are no other fixed strata for either subgroup. These
claims use the actual primitive two-dimensional period columns,
not merely the Kodaira types. They are checked for every row
in the retained certificate. The original additive caps have
free scalar action. In each simply connected finite cap the
actual gamma numerator is primitive modulo its full order;
applying gamma to an equality with a scalar translate forces
that affine power to be the identity. Thus these caps add no
fixed points even when their orders are divisible by $2$ or $p$.

\item \emph{The two complete primary invariants.}
The mod-two generator $v$ of $H^2(X;\mathbb F_2)$ has a Borel
lift whose characters on $C,E_2$ are respectively one and zero.
To justify a simultaneous choice, the based scalar action on
$W$ gives $\pi_1(W_{C_2})=\pi_1(W)\times C_2$. Its unique
order-two subgroup is the one seen at every fixed point.
The difference from the just-constructed line, being ordinary
zero on $W$, comes from group cohomology. Its restrictions
differ only by the same base class, and possibly a mixed
degree-one term on $E_2$. One base adjustment therefore gives
exactly the restrictions \eqref{maps:eq:D.C216}. The normal Euler class on
both curves is $\tau^4$, because both normal signs are negative
and their total Chern degree is even. The same Thom-localization
calculation as in \ref{maps:src:D.2.35} yields
\[
                         \langle v^3,[X]_2\rangle=1.     \tag{D.MP7}\label{maps:eq:D.MP7}
\]
Consequently the reduction of the spin-bordism class of $X$
to $K(\mathbb Z/2,2)$ is the conic-surgery class, by the full
calculation $\Omega^{\mathrm{Spin}}_6(K(\mathbb Z/2,2))
=\mathbb Z/2$ used in \ref{maps:src:D.2.8}.

\item For the odd reduction take the doubled line \eqref{maps:eq:D.OP4}; its
equivariant characters on $C,E_p$ are both zero. Normalize
$z\in H^2(X;\mathbb F_p)$ by $z[C]=1$. This is possible
because the wall degree in \eqref{maps:eq:D.MP4} is a unit. The torsion-free
group $\pi_1(W)$ again makes the base adjustment common to
both fixed curves. Thus the restrictions have exactly the
form \eqref{maps:eq:D.OP8}. The weights are $(1,-1)$ and $(w,-w)$.
Localization in degree two gives $z[E_p]=-w^2$, the ratio
of the equivariant $p_1$ to the normal Euler class is $-2$,
and the odd-degree mixed term squares to zero. Therefore
\[
 \langle z^3,[X]\rangle=0,\qquad
 \langle z\,p_1(X),[X]\rangle=-2(1-w^2).                 \tag{D.MP8}\label{maps:eq:D.MP8}
\]
\ref{maps:src:D.1.10} identifies the entire target
$H_6(K(\mathbb Z/p,2);\mathbb Z)=\mathbb Z/p$ with this cube
evaluation. Hence the odd fundamental-class invariant is zero,
and the refined Pontryagin coefficient in the stated generator
is $(w^2-1)/2$. Theorem~5 of \cite[p.796]{manZhubr75} says
that these two invariants inject the entire spin-bordism group;
it applies to an arbitrary map to $K(\mathbb Z/p,2)$, so no
isomorphism $H_2(X)\cong\mathbb Z/p$ is being assumed here.
Up to an automorphism its class is exactly that of
$Z_{p,0}$ if $w^2=1$, and $Z_{p,1}$ otherwise.

\item \emph{Why these primary computations identify the whole manifold.}
We supply the splitting, rather than assuming that ordinary
characteristic classes exclude a mixed refinement. Put
$K_2=K(\mathbb Z/2,2)$ and $K_p=K(\mathbb Z/p,2)$.
All their positive integral homology is respectively two-primary
and $p$-primary. One direct verification uses the path-loop
fibration: after inverting the relevant prime, the positive
homology of the fibre $K(\mathbb Z/q,1)$ vanishes by the
cyclic-group resolution. Its Serre sequence then identifies
the base homology with that of the contractible path space.
Thus every positive integral homology element of the base is
killed by a power of $q$, as claimed.

\item All tensor and Tor products between positive homology of
$K_2$ and $K_p$ vanish. The integral reduced Kunneth sequence
therefore makes $K_2\wedge K_p$ acyclic, and equivalently
the inclusion
\[
                K_2\vee K_p\longrightarrow K_2\times K_p
                                                        \tag{D.MP9}\label{maps:eq:D.MP9}
\]
an integral homology equivalence. Both spaces are simply
connected CW complexes. The relative Hurewicz theorem,
applied to the first possible nonzero relative homotopy group,
and then Whitehead's theorem make \eqref{maps:eq:D.MP9} a homotopy equivalence.
Since $\Omega^{\mathrm{Spin}}_6(\mathrm{pt})=0$, the wedge
axiom for reduced spin bordism consequently gives
\[
 \Omega^{\mathrm{Spin}}_6(K(\mathbb Z/(2p),2))
 \cong\Omega^{\mathrm{Spin}}_6(K_2)
       \oplus\Omega^{\mathrm{Spin}}_6(K_p).               \tag{D.MP10}\label{maps:eq:D.MP10}
\]
Here the comparison to the product is the actual Chinese
remainder isomorphism on $H_2$, and the two components are
its two projections. Its automorphism group is the product
of the corresponding primary automorphism groups.

\item The connected sum on the right of \eqref{maps:eq:D.MP1} maps to this product
by collapsing its separating sphere and using the canonical
classifying maps of its two summands. Its $H_2$ map is the
direct-sum isomorphism. A single oriented spin one-handle
bordism compares the connected sum with the disjoint union,
so its two bordism components are precisely the two just
computed for $X$. Both manifolds are simply connected, spin,
have this same $H_2$, and have $b_3=0$. Zhubr's Theorem~3
\cite[p.795]{manZhubr75} now gives the required oriented
diffeomorphism. In particular the possible two-primary
refinement has been determined by \eqref{maps:eq:D.MP7}, not discarded.

\item Finally the actual Tate weights are $w=4$ modulo five for
the three $\epsilon=0$ rows, $w=3$ modulo five for the three
$\epsilon=1$ rows, and $w=4$ modulo seven for all four
remaining rows. This gives exactly \eqref{maps:eq:D.MP2}. Distinct choices
of $\epsilon$ have distinct complete primary bordism
invariants, and both occur among the constructed members.
\end{enumerate}
\end{proof}

The exact certificate \path{code/III/manifolds_mixed_prime_AH.py}
retains all twelve patterns and ten rows, reconstructs their
Tate markings, and verifies the primitive angular lattices,
$736$ semicharacter substitutions, every period difference,
integrality and scalar character, both unique fixed elliptic
curves, and all odd normal weights. It imposes no finite
multiplicity or twist cutoff. The topological prime splitting
is proved above; it is not inferred from those finite checks.


\begin{theorem}[the full nine-valued invariant in the $c=3$ sector]\label{maps:src:D.2.39}
\StatementRecord{maps:src:D.2.39}{proved}{Historical events: Claude Fable 5.2 (Anthropic), effort high, 16 November 2026: carried as background without fault found. GPT-6 Astra (OpenAI), ChatGPT Pro mode, effort Pro, S6\_09\_PAPERS, 17 November 2026: targeted proof check. Also checked in this writing. Pre-continuation writing events: GPT-6 Astra (OpenAI), ChatGPT Pro mode, effort Pro, S6\_10\_PAPERS, 18 November 2026: targeted proof check. Continuation metadata events: GPT-6 Astra (OpenAI), ChatGPT Pro mode, effort Pro, S6\_10\_PAPERS continuation, 19 November 2026: used and rechecked in this run; status-vocabulary reconciliation; metadata reconciliation only; no new mathematical check. \emph{Pre-continuation recorded proof scope.} GPT-6 Astra (OpenAI, ChatGPT Pro mode), effort Pro, 2026-11-17 UTC; Claude Fable 5.2 (Anthropic), effort high, 16 November 2026; also checked in this writing by GPT-6 Astra (OpenAI), ChatGPT Pro mode, effort Pro, 18 November 2026}{the hypotheses, notation and definitions stated here; \cite[Sections 3.11--3.14 and Theorem 5]{manZhubr75}; \cite[equation (2.8), p.267]{ffThomas56}; \cref{maps:src:D.2.26}; \cref{maps:src:D.2.36}; \cref{maps:eq:D.C26}; \cref{maps:eq:D.OP4}; \cref{maps:eq:D.OP5}; \cref{maps:eq:D.OP5a}; \cref{maps:eq:D.T35}; \cite[equations (2.3), (2.5), (2.7), p.267]{ffThomas56}; \cref{maps:eq:D.T34}; \cref{maps:eq:D.T37}; \cref{maps:eq:D.T39}; \cite[p.795]{manZhubr75}}


Use the primitive Mumford filling at every multiplicative site in one
of the following eight minimal patterns, retain the indicated additive
choice and the ordinary clutch $e\gamma$, and make no additional
regular logarithmic transformations:
\[
\begin{array}{c|r|r|l|l}
\text{OS/profile}&a&b&\text{zero positions}&\text{additive choice}\\\hline
56&5&3&2&\text{none}\\
56&10&3&3,4&\text{none}\\
56/II&5&3&0&II\text{ original or log}\\
56/III&1&6&0&III\text{ log}\\
56/III&5&6&3&III\text{ log}\\
56/III&10&3&3&III\text{ original}.
\end{array}                                                   \tag{D.T31}\label{maps:eq:D.T31}
\]
Zero positions are numbered from zero. For every finite multiplicity
and affine twist producing a simply connected manifold, the exact
oriented diffeomorphism type is
\[
 M_{3,0}=\bigl((S^2\times S^4)\setminus
          \operatorname{int}(S^2\times D^4)\bigr)
               \mathop{\cup}_{S^2\times S^3}(D^3\times S^3),     \tag{D.T32}\label{maps:eq:D.T32}
\]
where the embedded sphere represents three times the generator of
$H_2(S^2\times S^4;\mathbb Z)$ and has its oriented normal framing.
For example it is the graph of the degree-three map
$z\mapsto z^3$ on $\mathbb {CP}^1=S^2$ and an equatorial embedding
$S^2\hookrightarrow S^4$.
The complete Zhubr data are
\[
 H_2=H_3=\mathbb Z/3,\qquad b_3=0,\qquad
 \mathfrak p=0,\qquad
 \mu=0\in H_6(K(\mathbb Z/3,2);\mathbb Z)=\mathbb Z/9.          \tag{D.T33}\label{maps:eq:D.T33}
\]
In particular the assertion about $\mu$ determines its order-three
subgroup as well as its ordinary mod-three cubic reduction.

\end{theorem}


\begin{proof}
\begin{enumerate}[label=\arabic*.,leftmargin=*]
\item \emph{1. The invariant which must be computed.}
We first identify its complete target from the primary classification
and the actual Pontryagin operation. Zhubr's epimorphism in Section~3.11
and his free calculation in Section~3.12 show that
$\Omega^{\mathrm{Spin}}_6(K(\mathbb Z/3,2))$ is generated by the
images of $\epsilon,\zeta$, where
\[
 (P,h)(\epsilon)=(6,0),\qquad
 \zeta=[\mathbb {CP}^3,\text{hyperplane}].
\]
The image of $\epsilon$ has both invariants zero: its first is six
modulo three, and its second is the image of zero. The injectivity
in Zhubr's Theorem~5 therefore kills that image.
Lemma~3.14 gives $N(\mathbb Z/3)=0$, so the fundamental-class map
identifies this cyclic bordism group with
$H_6(K(\mathbb Z/3,2);\mathbb Z)$; the mod-two obstruction target
is zero because all positive homology is three-primary.
These facts use \cite[Sections 3.11--3.14 and Theorem 5]{manZhubr75}.

\item The cyclic generator has order dividing nine. On
$\mathbb {CP}^1\times\mathbb {CP}^2$ take the integral class
$3x+y$, with $x,y$ the two hyperplane classes. Its mod-three
reduction factors through $\mathbb {CP}^2$, and hence sends the
six-dimensional fundamental class to zero. Before reduction its
cube evaluates to nine. Naturality therefore gives nine times the
universal generator equal to zero.
For the opposite bound, Thomas's Pontryagin third power is a natural
operation
\[
 \mathcal P_3:H^2(-;\mathbb Z/3)\longrightarrow
                              H^6(-;\mathbb Z/9).
\]
Its value on the reduced hyperplane class of $\mathbb {CP}^3$
evaluates to a generator of $\mathbb Z/9$
\cite[equation (2.8), p.267]{ffThomas56}. Thus the cyclic group has
order exactly nine. Evaluation of $\mathcal P_3$ detects its entire
generator, and hence the whole fundamental-class invariant. This is
stronger than evaluating the ordinary cube modulo three.

\item \emph{2. The common-piece comparison and its circle geometry.}
\ref{maps:src:D.2.26} gives $H_2(X)=H_3(X)=\mathbb Z/3$ and $b_3=0$ for every
simply connected member in \eqref{maps:eq:D.T31}. In particular $H^2(X;\mathbb Z)=0$,
so $c_1(X)=0$ and $X$ is spin. Let
$u\in H^2(X;\mathbb F_3)$ be the generator normalized below.
All the geometric hypotheses of the doubled-line calculation in
\ref{maps:src:D.2.36} hold here. We record the exact inputs to ensure that no
prime-greater-than-three conclusion is being imported.

\item Remove the additive caps and one ordinary clutch disc to form $W$.
In all seven marked initial rows the actual multiplicative angular
columns generate $\ker\gamma^*$ with three unit Smith factors.
Consequently $\pi_1(W)=\mathbb Z_\gamma\times F_r$ is torsion-free.
The integral degree-one bridge calculation and surjective degree-two
cap boundary maps give $H_2(V)\twoheadrightarrow H_2(W)$, as in
\eqref{maps:eq:D.C26}. The pure fibre functional is
$\Omega=a\,u_1^*u_2^*+b\,\gamma^*\delta^*$, with $3\mid b$
and $3\nmid a$. Its finite normal coefficients are divisible by
$b/q_i=3$, its original additive coefficients are zero, and the
ordinary clutch coefficient is $be$. Thus its mod-three reduction
is the actual global degree-two generator, with zero marked swept
values. The doubled factors \eqref{maps:eq:D.OP4} and the exact integral supports
\eqref{maps:eq:D.OP5}, \eqref{maps:eq:D.OP5a} produce the corresponding line on $W$.

\item The scalar circle has precisely two strata whose stabilizer order is
divisible by three: the rational curve $C$ fixed by the whole circle,
and one elliptic double curve $E$ with stabilizer exactly $C_3$.
The latter is at the unique $I_3$ position. On its complement all
stabilizers are finite of order prime to three. This includes the
finite logarithmic caps: primitivity of their actual gamma numerators
makes the scalar circle free there, for every allowed order.
The original additive caps have free scalar action as well.
The seven exact Tate markings in the certificate verify the remaining
multiplicative assertions, including the primitive short period
$(3,aj)$ at $E$ and the two normal characters
$w,-w$, where $w=(aj)^{-1}=\pm1$ in $\mathbb F_3$.

\item The doubled auxiliary line has zero $C_3$ character at both $C$ and
$E$. Its degree on $C$ is a unit modulo three. Normalize $u$ by
$\langle u,[C]\rangle=1$. The torsion-free group of $W$ gives the
same simultaneous fixed-point comparison as in D.2.36: once an
equivariant lift has zero point character on $C$, it also has zero
point character on $E$. Indeed the two fixed-point copies of $BC_3$
represent the same order-three subgroup in
$\pi_1(W_{C_3})=\pi_1(W)\times C_3$. The group-cohomology difference
from half the auxiliary-line class has a common degree-two base
coefficient; the possible mixed degree-one term does not change
that coefficient. This comparison uses the complete normal columns,
including arbitrary finite twists.

\item \emph{3. The full-circle lift and the elliptic Borel space.}
The generator $u$ admits a lift
\[
                 \widehat u\in H^2_{S^1}(X;\mathbb F_3).
\]
In the Borel spectral sequence the only possible differentials out
of ordinary degree two have targets involving $H^1(X;\mathbb F_3)$
or $H^3(BS^1;\mathbb F_3)$; both are zero. Adding a multiple of the
base generator normalizes its restriction on the fixed curve to
\[
                      \widehat u|_{C_{S^1}}=[C]^* .           \tag{D.T34}\label{maps:eq:D.T34}
\]
Restriction $H^2(BS^1;\mathbb F_3)\to H^2(BC_3;\mathbb F_3)$
is nonzero. The point-character comparison in part~2 consequently
makes the coefficient of the degree-two $BC_3$ class on $E$ zero
as well.

\item The effective action $S^1/C_3$ on the elliptic curve $E$ is free.
Its quotient is a circle. Every principal circle bundle over a
circle is trivial, so equivariantly $E$ is the product of a free
$S^1/C_3$ orbit and a circle. Therefore
\[
 E_{S^1}\simeq BC_3\times S^1,
 \qquad
 \widehat u|_{E_{S^1}}=\eta\times\zeta,                    \tag{D.T35}\label{maps:eq:D.T35}
\]
where $\eta\in H^1(BC_3;\mathbb F_3)$ and
$\zeta\in H^1(S^1;\mathbb F_3)$, allowing a scalar in either
factor. Formula \eqref{maps:eq:D.T35} retains the possible nonzero mixed class;
it does not assert that $u$ descends to the coarse quotient $X/C_3$.
In fact the mixed term is generally essential.

\item For odd primes Thomas proves exact external-product multiplicativity
and vanishing on odd-dimensional classes, and independently proves
vanishing on suspensions
\cite[equations (2.3), (2.5), (2.7), p.267]{ffThomas56}.
These are statements with the full target coefficients, so
\[
 \mathcal P_3(\widehat u)|_{E_{S^1}}
   =\mathcal P_3(\eta)\times\mathcal P_3(\zeta)=0
                       \quad\text{in }H^6(E_{S^1};\mathbb Z/9).
                                                               \tag{D.T36}\label{maps:eq:D.T36}
\]
There is also no contribution on $C_{S^1}$: \eqref{maps:eq:D.T34} is pulled back
from the two-dimensional curve, whose degree-six cohomology is zero.
The familiar nonvanishing phenomena for the Pontryagin square of
suspended classes are irrelevant here; the cited formulas explicitly
assume an odd prime.

\item \emph{4. Why the full order-nine class vanishes globally.}
Put $R=\mathbb Z/9$. Choose disjoint invariant tubular neighborhoods
$N$ of $C\sqcup E$ and let $Y=X\setminus\operatorname{int}N$.
The circle acts locally freely on $Y$, with stabilizer orders prime
to three. The quotient $Y/S^1$ is a compact five-dimensional orbifold
with boundary. The Borel projection to this quotient, for both $Y$
and its boundary, has local cohomology fibres $BH$ with $H$ finite
of order invertible in $R$. Finite averaging gives
$H^q(BH;R)=0$ for $q>0$. The Leray comparison and the long exact
sequences of the pairs therefore give
\[
 H^k_{S^1}(Y,\partial Y;R)
             \cong H^k(Y/S^1,\partial(Y/S^1);R)=0
                            \qquad(k>5).                     \tag{D.T37}\label{maps:eq:D.T37}
\]
The dimension bound can be read from an orbifold triangulation,
or from a finite good orbifold atlas and its relative cellular
cohomology. No division by three occurs in this comparison.

\item Equivariant excision identifies the left side of \eqref{maps:eq:D.T37} with
$H^k_{S^1}(X,N;R)$. At $k=6$ it is zero. Thus restriction
\[
 H^6_{S^1}(X;R)\longrightarrow H^6_{S^1}(C\sqcup E;R)
                         \text{ is injective}.               \tag{D.T38}\label{maps:eq:D.T38}
\]
The vanishing of both restrictions proved in part~3 now implies
$\mathcal P_3(\widehat u)=0$ in the full integral-coefficient group
with coefficients $R$. Restricting to the ordinary fibre $X$ gives
\[
                        \langle\mathcal P_3(u),[X]\rangle=0
                                     \quad\text{in }\mathbb Z/9.
                                                               \tag{D.T39}\label{maps:eq:D.T39}
\]
Part~1 proves that this is the entire invariant $\mu$. Since
$N(\mathbb Z/3)=0$, the spin-bordism class itself is zero, and its
refined Pontryagin class is zero as well. This also agrees with
ordinary localization, which yields $p_1=0$ because $w^2=1$;
that latter congruence alone would not have proved \eqref{maps:eq:D.T39}.

\item \emph{5. The stated concrete smooth model.}
The graph embedding in \eqref{maps:eq:D.T32} represents three times the generator.
Its oriented real rank-four normal bundle is trivial: the ambient
manifold is spin and $w_2(TS^2)=0$, and $w_2$ classifies such bundles
on $S^2$. The framing is unique up to homotopy because
$\pi_2SO(4)=0$. The framed surgery preserves simple connectivity
and spin. Its integral sequence and duality give
$H_2=H_3=\mathbb Z/3$, with $b_3=0$.
The mod-three class from the first $S^2$ extends across the surgery
trace, since it evaluates to zero on the attaching sphere.
The original pair $(S^2\times S^4,u)$ spin-bounds
$S^2\times D^5$ over $K(\mathbb Z/3,2)$.
Thus the resulting model has zero spin-bordism class over this
Eilenberg--Mac Lane space. Zhubr's Theorem~3
\cite[p.795]{manZhubr75} identifies it, orientation-preservingly,
with every manifold in \eqref{maps:eq:D.T31}, since both classifying maps induce
isomorphisms on $H_2$ and both have $b_3=0$.
\end{enumerate}
\end{proof}

The program \path{code/III/manifolds_tate_markings.py} was executed with
residual three; the exact extraction output is
\path{code/fourfolds_D2_c3_tate.out}.
The independent certificate \path{code/III/fourfolds_D2_c3_AH.py}
checks all eight patterns on seven marked rows, all 27 adapted
monodromy identities, 14 full four-row angular and global-image
Hermite forms, all three finite normal rows with both arbitrary-twist
factors, and both original additive choices. It also checks every
doubled-line period identity, the unique three-primary isotropy
curve and its normal weights. Its retained output and full
records are \path{code/fourfolds_D2_c3_AH.out} and
\path{records/fourfolds_D2_c3_AH.json}. The order-nine computation
uses the exact primary operation and the relative Borel argument
above; it is not extrapolated from a finite mod-three test.


\begin{theorem}[shared odd factors in every weighted profile]\label{maps:src:D.2.40}
\StatementRecord{maps:src:D.2.40}{proved}{Historical events: Claude Fable 5.2 (Anthropic), effort high, 16 November 2026: carried as background without fault found. GPT-6 Astra (OpenAI), ChatGPT Pro mode, effort Pro, S6\_09\_PAPERS, 17 November 2026: targeted proof check. Also checked in this writing. Pre-continuation writing events: GPT-6 Astra (OpenAI), ChatGPT Pro mode, effort Pro, S6\_10\_PAPERS, 18 November 2026: targeted proof check. Continuation metadata events: GPT-6 Astra (OpenAI), ChatGPT Pro mode, effort Pro, S6\_10\_PAPERS continuation, 19 November 2026: used and rechecked in this run; status-vocabulary reconciliation; metadata reconciliation only; no new mathematical check. \emph{Pre-continuation recorded proof scope.} GPT-6 Astra (OpenAI, ChatGPT Pro mode), effort Pro, 2026-11-17 UTC; Claude Fable 5.2 (Anthropic), effort high, 16 November 2026; also checked in this writing by GPT-6 Astra (OpenAI), ChatGPT Pro mode, effort Pro, 18 November 2026}{the hypotheses, notation and definitions stated here; \cref{maps:src:D.2.37}; \cref{maps:eq:D.WP8}; \cref{maps:eq:D.WP9}; \cref{maps:eq:D.WP6}; \cref{maps:eq:D.OW5}; \cref{maps:src:D.2.29}; \cref{maps:eq:D.WP5}; \cref{maps:src:D.2.32}; \cref{maps:src:D.2.34}; \cref{maps:src:D.2.33}; \cref{maps:eq:D.WP11}; \cref{maps:eq:D.OW6}}


The uniform weighted algorithm of \ref{maps:src:D.2.37} extends from
$\gcd(a,d)=1$ to
\[
                     \gcd(a,d)\text{ odd}.                \tag{D.OW1}\label{maps:eq:D.OW1}
\]
In particular it covers \emph{every} positive translation
iterate of every row with $a$ odd. All its original,
Mumford, finite-good and regular logarithmic choices,
freeness tests and full Seifert groups remain valid.
Star caps are still original. The free ranks and Euler
number are still exactly \eqref{maps:eq:D.WP8}--\eqref{maps:eq:D.WP9}, with
$\chi=2d$.

The torsion must now be computed with the full fibre
presentation rather than one alternating functional.
Use the target
\[
                       \Lambda^2\Lambda\oplus K
\]
and include all columns
\[
 ((\Lambda^2A_i(d)-I)\eta,0)\qquad
                  (\eta\text{ a standard bivector}).     \tag{D.OW2}\label{maps:eq:D.OW2}
\]
For a finite or regular logarithm replace \eqref{maps:eq:D.WP6} by
\[
 (-\lambda\wedge v,\;J_iN_m\lambda)
       \quad(\lambda=e_1,e_2,\delta,\gamma),\qquad
                         (0,J_iv).                       \tag{D.OW3}\label{maps:eq:D.OW3}
\]
At every other cap include its full fixed-normal
columns $(0,J_iF_i)$ and the complete fibre kernel
\[
 {\cal K}_i(d)=
 \bigl\{\eta\in\Lambda^2\Lambda:
          (\Lambda^2S_d)\eta\in
          {\cal K}_i(1)\otimes\mathbb Q\bigr\}.             \tag{D.OW4}\label{maps:eq:D.OW4}
\]
Here $\mathcal K_i(1)$ is the recorded actual fibre kernel
for a nonzero-Chern original. In every other case its
rational span is
$\operatorname{im}(\Lambda^2A_i-I)$.
Formula \eqref{maps:eq:D.OW4} is evaluated as an ordinary saturated
integer kernel, not by discarding the nonunit Smith
factors of \eqref{maps:eq:D.OW2}.

Let $T$ be the torsion cokernel of this full matrix.
Then \eqref{maps:eq:D.WP9}, with this $T$, is the complete integral
homology for every datum in \eqref{maps:eq:D.OW1}.
The executable \path{code/III/manifolds_odd_weighted_profiles.py}
implements this extension and retains all shared-factor
torsion.

\end{theorem}


\begin{proof}
\begin{enumerate}[label=\arabic*.,leftmargin=*]
\item \emph{The integral involution splitting.}
Put $g=\gcd(a,d)$. The prime-local proof in \ref{maps:src:D.2.37}
shows, without assuming $g=1$, that every torsion
prime of
\[
 C=(\Lambda^2\Lambda)_{\{A_i(d)\}}
\]
divides $g$. Its free rank is one. Under \eqref{maps:eq:D.OW1}
we may write $C\simeq\mathbb Z\oplus C_{\rm tor}$,
with $C_{\rm tor}$ of odd order. The regular graph
sequence is an exact sequence of abelian groups
\[
                 0\longrightarrow C\longrightarrow
           H_2(V)\longrightarrow K\longrightarrow0.      \tag{D.OW5}\label{maps:eq:D.OW5}
\]
Fibre inversion is an actual involution of this
linear torus bundle. It acts by $+1$ on $C$ and
by $-1$ on $K$.

\item Let $A:H_2(V)\to\mathbb Z$ be evaluation by the
pullback of the marked weight-one class used in
\ref{maps:src:D.2.37}. It is invariant under inversion: this is
true on the old fibre and on all its generating
fixed sweeps, and therefore remains true on the
cover. Its restriction to $C$ is
\[
 g\,\omega',\qquad
 \omega'=(a/g)e_1^*e_2^*-(db/g)\delta^*\gamma^*,
\]
where $\omega'$ is the primitive free functional.
For any $x\in H_2(V)$, $x+\iota x$ belongs to $C$.
Hence
$2A(x)=A(x+\iota x)$ is divisible by $g$.
Since $g$ is odd, $A/g$ is an integral homomorphism.
Call it $B$.

\item We spell out the resulting splitting to avoid a
formality or invariant-descent assumption. Given
$k\in K$, choose a lift $x$ with $B(x)=0$, subtracting
a fibre generator if necessary. Then
$x+\iota x\in C_{\rm tor}$. Multiplication by two
is an automorphism of this finite group, so
\[
                  y=x-\tfrac12(x+\iota x)                 \tag{D.OW6}\label{maps:eq:D.OW6}
\]
is a well-defined integral class with
$\iota y=-y$ and image $k$. If the lift $x$ is
changed by $c\in C_{\rm tor}$, the half term
changes by $c$, so $y$ is unchanged.
This constructs an additive section of \eqref{maps:eq:D.OW5}.
It is the unique section with $B=0$ and
anti-invariant image. Thus
$H_2(V)=C\oplus K$ with an actual integral marking,
including the odd torsion of $C$.

\item Every literal fixed-loop peripheral torus is
anti-invariant, since inversion reverses its
fibre circle and fixes the base circle. It
therefore has exactly its graph coordinate and
zero fibre coordinate in this splitting.
This remains true for the norm sweep of any
covering-torus loop: take the pushforward of
$\lambda\times S^1$ from the $m$-fold mapping-torus
cover, where $A^m=I$. It is anti-invariant and
has graph coordinate $N_m\lambda$.
This is the step that controls integral torsion
which a vertical alternating form alone cannot see.

\item \emph{The full finite and regular normal columns.}
On the $m$-fold boundary cover the normal meridian
is $m\theta-v$. The shear
$(x,t)\mapsto(x-tv,t)$ sends the product torus
of $\lambda$ and the covering base circle to
that norm sweep minus the regular torus
$\lambda\wedge v$. Pushing it to the ordinary
mapping torus and then to $V$, the preceding
anti-invariant description gives exactly the
first four columns of \eqref{maps:eq:D.OW3}.
There is no unknown torsion component in the
norm sweep.

\item For the affine meridian, use its commuting lifts
$\theta$ and $\mu=\theta^m v^{-1}$ in the boundary
group. Their torus is the fixed torus of $v$
and $\theta$, with the normal orientation
convention of \ref{maps:src:D.2.29}. Hence its image is
$(0,J_iv)$, the last column in \eqref{maps:eq:D.OW3}.
These loops generate the complete normal Thom
image, including torsion loops in the central
finite quotient. The proof of the general
freeness test \eqref{maps:eq:D.WP5} is unchanged.

\item \emph{Every nonfinite fibre kernel and its marked normal.}
All the unchanged caps in this construction have
torsion-free $H_2$. For the rank-two cusp this is
the explicit normalization calculation. For a
rank-one gamma cover the normalized components
are $S^2$ bundles of degree zero over a covered
two-torus, and consecutive double curves are
identified by translations; straightening those
translations gives the same free homology as a
cycle of spheres times $T^2$. Original
zero-degree additive caps have the permutation
mapping-torus homology of a Kodaira tree times
$S^1$, which is free. The nonzero-Chern originals
have the primitive Gysin presentations of
\ref{maps:src:D.2.32}--\ref{maps:src:D.2.33}, also free.

\item The covering square on regular fibres multiplies
bivectors by $\Lambda^2S_d$. The needed ranks come
from the actual caps, as follows. At the rank-two
zero the short-period map has Smith factors $(1,d)$.
In that basis its exterior-square difference has
primitive image of rank two, so the boundary has
$H_2=\mathbb Z^6$. The normalization calculation of
\ref{maps:src:D.2.34} gives $H^3(N)=\mathbb Z^2$. The two angular
normal tori are a primitive sublattice of boundary
homology and bound by their section-disc sweeps.
The pair sequence bounds its whole kernel by rank
two; hence they are exactly that kernel. In
particular the regular-fibre coinvariant of rank
four injects into the cap. This proves that the
rational fibre kernel is precisely the monodromy
image, without a local invariant-cycle assumption.

\item For the rank-one cap, the explicit cycle of
$\rho$ spheres times $T^2$ has four independent
nearby-fibre classes: the sum of the component
spheres, the surviving nodal-circle class times
each transverse circle, and the transverse
$T^2$ area. This is the full four-dimensional
fibre coinvariant, so again its kernel is exactly
the monodromy image over $\mathbb Q$. For an
original zero-degree additive cap, the permutation
mapping torus of the Kodaira tree gives the two
independent images: the nonzero multiplicity-weighted
whole-fibre class in the component-orbit space,
and the transverse $\delta\gamma$ area. These
are the full two-dimensional fibre coinvariant,
also for the original odd-star fibres. The
zero-degree multiplicative original has the same
four explicit classes as the nodal-cycle model.

\item Thus in every zero-degree case the rational
fibre kernel is exactly the inverse image of
the old one. For a nonzero-Chern original the
actual area and principal-bundle formulas give
the same statement directly. Since the target
$H_2$ is torsion free, the integral kernel is the
saturation of that rational kernel. This is
precisely \eqref{maps:eq:D.OW4}.

\item The full fixed-normal projections lift the old
bounding section-disc sweeps, as in \ref{maps:src:D.2.37}.
There is one previously isolated correction:
at a nonzero-Chern original multiplicative cap,
the vanishing-loop normal can differ from a
kernel cycle by an integral multiple of
$u_2\wedge\delta$. In \ref{maps:src:D.2.33} only its zero
alternating-form value was needed. Here it
vanishes in the full global fibre coinvariant.
Indeed, for all $172$ rows the elliptic matrices
$B_i$ satisfy the exact unit Smith certificate
\[
       \mathbb Z^2/\sum_i(B_i-I)\mathbb Z^2=0.             \tag{D.OW7}\label{maps:eq:D.OW7}
\]
These $172$ matrices and their unit factors
are retained in
\path{records/manifolds_elliptic_coinvariants.json}.
Since every $A_i(d)$ fixes $\delta$,
\[
 (\Lambda^2A_i(d)-I)(\delta\wedge z)
                 =\delta\wedge(B_i-I)z
       \quad(z\in\mathbb Z^2).
\]
Consequently $\delta\wedge e_1$ and
$\delta\wedge e_2$ are zero integrally in $C$,
for every weight. This kills the correction
without any assumption on its integer value.

\item Subtracting these complete marked normals
from a boundary-kernel element leaves a
regular-fibre element. Its full image is
exactly \eqref{maps:eq:D.OW4}. Thus \eqref{maps:eq:D.OW2}--\eqref{maps:eq:D.OW4}
give the entire integral global kernel;
in particular a primitive fibre relation
that was only a torsion class in the boundary
has not been omitted.

\item \emph{Global homology and the remaining scope.}
The relative pair argument \eqref{maps:eq:D.WP11} does not
require the regular coinvariant to be free.
Finite and regular caps still have free
relative degree-two groups, and their
nonzero base-meridian coefficients make the
first relative map injective. Their full
degree-three images are now \eqref{maps:eq:D.OW3}.
The remaining caps still have free
$H^4$, so the final kernel in \eqref{maps:eq:D.WP11} is
free and the abelian-group extension splits.
The top component counts and rational ranks
are unchanged. Therefore \eqref{maps:eq:D.WP8} and the
entire duality table hold, with exactly the
full torsion cokernel computed here.

\item The vertical normal closures are also unchanged,
so the same all-order Seifert presentation
computes the full fundamental group.
The argument has not used a bound on weights
or logarithmic parameters. It uses oddness
only in dividing $g$ in the invariant functional
and in the unique half of \eqref{maps:eq:D.OW6}. When both
$a$ and $d$ are even, a two-primary ambiguity
can remain in that marking; this theorem
does not discard it.
\end{enumerate}
\end{proof}

The exact implementation computes the saturated
preimages \eqref{maps:eq:D.OW4} before forming any Smith quotient.
Its overlap checks compare this full fibre-and-graph
matrix with \ref{maps:src:D.2.37}'s independent scalar-functional
matrix whenever $g=1$. Shared odd-factor cases
retain the nontrivial regular fibre coinvariant
and all its relations. The printed certificates
and explicit sample matrices are retained in
\path{records/manifolds_odd_weighted_profiles.json}.

The shared-even-factor case outside the odd-iterate statement is included in \ref{maps:src:D.2.41}.

\begin{theorem}[every positive weight by integral bar chains]\label{maps:src:D.2.41}
\StatementRecord{maps:src:D.2.41}{proved}{Historical events: Claude Fable 5.2 (Anthropic), effort high, 16 November 2026: carried as background without fault found. GPT-6 Astra (OpenAI), ChatGPT Pro mode, effort Pro, S6\_09\_PAPERS, 17 November 2026: targeted proof check. Also checked in this writing. Pre-continuation writing events: GPT-6 Astra (OpenAI), ChatGPT Pro mode, effort Pro, S6\_10\_PAPERS, 18 November 2026: targeted proof check. Continuation metadata events: GPT-6 Astra (OpenAI), ChatGPT Pro mode, effort Pro, S6\_10\_PAPERS continuation, 19 November 2026: used and rechecked in this run; status-vocabulary reconciliation; metadata reconciliation only; no new mathematical check. \emph{Pre-continuation recorded proof scope.} GPT-6 Astra (OpenAI, ChatGPT Pro mode), effort Pro, 2026-11-17 UTC; Claude Fable 5.2 (Anthropic), effort high, 16 November 2026; also checked in this writing by GPT-6 Astra (OpenAI), ChatGPT Pro mode, effort Pro, 18 November 2026}{the hypotheses, notation and definitions stated here; \cref{maps:src:D.2.37}; \cref{maps:src:D.2.40}; \cref{maps:eq:D.WP8}; \cref{maps:eq:D.WP9}; \cref{maps:eq:D.OW2}; \cref{maps:eq:D.OW4}; \cref{maps:eq:D.WP5}; \cite[Chapter I, Section 5]{manBrown82}; \cref{maps:eq:D.BC9}; \cref{maps:eq:D.OW3}; \cref{maps:eq:D.BC8}; \cref{maps:eq:D.BC10}; \cref{maps:eq:D.BC14}; \cref{maps:eq:D.BC13}; \cref{maps:src:D.2.29}; \cref{maps:eq:D.OW7}; \cref{maps:eq:D.WP11}}


The restriction on $\gcd(a,d)$ in \ref{maps:src:D.2.37} and \ref{maps:src:D.2.40} can
be removed. Every positive iterate $Q\mapsto dQ$ of
every one of the $172$ natural initial rows has an
exact integral homology algorithm, for all $495$
original/Mumford presentation patterns, every free
finite-good affine logarithm, and arbitrary regular
logarithms and clutching integer. Star caps remain
original. The full Seifert fundamental group and
the free ranks \eqref{maps:eq:D.WP8}--\eqref{maps:eq:D.WP9} are unchanged.

Here is the required all-parity correction to the
full matrix \eqref{maps:eq:D.OW2}--\eqref{maps:eq:D.OW4}. Index the boundary
matrices as $A_0,\ldots,A_r$, with
$A_0\cdots A_r=I$, and use the first $r$ base loops
as independent generators. All $A_i$ below mean
the actual weighted matrices $A_i(d)$.
For $u,v\in\Lambda=\mathbb Z^4$ set
\[
 p(u,v)=\sum_{\alpha<\beta}u_\alpha v_\beta\,
                  e_\alpha\wedge e_\beta,\qquad
 D_A(u,v)=p(Au,Av)-(\Lambda^2A)p(u,v).
                                                        \tag{D.BC1}\label{maps:eq:D.BC1}
\]
The bilinear map $D_A$ is symmetric. Define its
integral quadratic primitive by
\[
 f_A(u)=-\sum_\alpha D_A(e_\alpha,e_\alpha)
                       \binom{u_\alpha}{2}
        -\sum_{\alpha<\beta}D_A(e_\alpha,e_\beta)
                       u_\alpha u_\beta .
                                                        \tag{D.BC2}\label{maps:eq:D.BC2}
\]
Negative coordinates use the same integral polynomial
$\binom n2=n(n-1)/2$.
Define the peripheral correction
\[
 \begin{split}
 c_i(u)&=f_{A_i}(u) \quad (i<r),\\
 c_r(u)&=-\sum_{j<r}
       f_{A_j}(A_{j+1}\cdots A_{r-1}A_r u).
 \end{split}                                            \tag{D.BC3}\label{maps:eq:D.BC3}
\]
The transported reversed word in the second line
is part of the formula.

Keep every fibre relation \eqref{maps:eq:D.OW2} and every full
saturated cap kernel \eqref{maps:eq:D.OW4}. Replace the five
finite or regular logarithmic normal columns by
\[
 \begin{gathered}
 \left(\sum_{\nu=0}^{m-1}c_i(A_i^\nu\lambda)
                  -\lambda\wedge v,\quad J_iN_m\lambda\right)
       \quad(\lambda=e_1,e_2,\delta,\gamma),\\
                       (c_i(v),J_iv).
 \end{gathered}                                          \tag{D.BC4}\label{maps:eq:D.BC4}
\]
At every other cap the fixed-normal columns are
\[
                         (c_i(v),J_iv),
                         \qquad v\text{ in a basis of }\operatorname{Fix} A_i.
                                                        \tag{D.BC5}\label{maps:eq:D.BC5}
\]
The finite freeness test is still exactly \eqref{maps:eq:D.WP5}.
If $A_i$ has order $h\mid m$, the sum in \eqref{maps:eq:D.BC4}
is $(m/h)\sum_{\nu=0}^{h-1}c_i(A_i^\nu\lambda)$;
thus this is a finite algorithm for unbounded orders.
Let $T$ be the torsion cokernel of this full matrix.
Then the complete integral homology is \eqref{maps:eq:D.WP9}
with this $T$, for every $d>0$, including all
shared even factors.

\end{theorem}


\begin{proof}
\begin{enumerate}[label=\arabic*.,leftmargin=*]
\item \emph{An integral chain comparison with an explicit marking.}
Use the normalized bar chains of $\Lambda$ with
trivial coefficients. Their low-dimensional
boundaries are
\[
 \partial[u]=0,\qquad
 \partial[u|v]=[v]-[u+v]+[u].
                                                        \tag{D.BC6}\label{maps:eq:D.BC6}
\]
The bar complex is a functorial chain model of
$B\Lambda$; see the standard resolution
\cite[Chapter I, Section 5]{manBrown82}.
We give all the comparison formulas used here.
The maps
\[
             p_1([u])=u,\qquad p_2([u|v])=p(u,v)
                                                        \tag{D.BC7}\label{maps:eq:D.BC7}
\]
annihilate boundaries in the appropriate degree.
In degree two the cycle $[u|v]-[v|u]$ maps to
$u\wedge v$. The coordinate subtori give an
integral basis, so $p_2$ induces exactly
$H_2(T^4;\mathbb Z)\simeq\Lambda^2\Lambda$.

\item Antisymmetrizing \eqref{maps:eq:D.BC1} gives zero, proving
the symmetry of $D_A$. The binomial identity
$\binom{x+y}{2}=\binom{x}{2}+\binom{y}{2}+xy$
then gives the exact integral identity
\[
 f_A(u)+f_A(v)-f_A(u+v)
             =p(Au,Av)-(\Lambda^2A)p(u,v).
                                                        \tag{D.BC8}\label{maps:eq:D.BC8}
\]
Equivalently, extending $f_A([u])=f_A(u)$
linearly on bar one-chains,
$f_A\partial=p_2A-(\Lambda^2A)p_2$ on bar
two-chains.

\item The regular piece retracts to its torus bundle
over a rose with $r$ petals. Its integral bar
mapping-cone complex has
\[
 {\cal C}_n=B_n(\Lambda)\oplus
                   \bigoplus_{i<r}B_{n-1}(\Lambda),
 \qquad
 \partial(c;g_i)=
 \left(\partial c+\sum_{i<r}(A_i-I)g_i;\
                                     -\partial g_i\right).
                                                        \tag{D.BC9}\label{maps:eq:D.BC9}
\]
This is a model for the actual linear torus bundle,
not an assertion of equivariant formality. Its
universal cover is a torus-universal-cover bundle
over a tree, hence contractible. The bar model
has the same semidirect-product fundamental group,
and their maps on every peripheral torus agree
when their marked loops agree.

\item The sign in \eqref{maps:eq:D.BC9} uses the existing inverse
endpoint convention
$(x,1)\sim(A_i^{-1}x,0)$. Label a cylinder by its
terminal fibre chain $g$; its initial label is
$A_i g$. With fibre-first product orientation
its boundary is $(A_i-I)g$. This is the same cut
orientation used in \eqref{maps:eq:D.OW3}.

\item In degree two define
\[
 P_2(c;g_i)=
 \left(p_2c+\sum_{i<r}f_{A_i}(g_i);\
                                  (p_1g_i)_{i<r}\right).
                                                        \tag{D.BC10}\label{maps:eq:D.BC10}
\]
The target is the small complex whose degree-two
homology is
\[
 C\oplus K,\qquad
 C=\Lambda^2\Lambda/
          \sum_{i<r}(\Lambda^2A_i-I)\Lambda^2\Lambda,
 \quad
 K=\ker[A_0-I\ \cdots\ A_{r-1}-I].
                                                        \tag{D.BC11}\label{maps:eq:D.BC11}
\]
No saturation or torsion quotient is taken in $C$.
Formula \eqref{maps:eq:D.BC8} proves that $P_2$ carries a
degree-three boundary to a small-complex boundary:
its fibre coordinate is
\[
 \sum_{i<r}\bigl(p_2(A_i-I)h_i-f_{A_i}\partial h_i\bigr)
       =\sum_{i<r}(\Lambda^2A_i-I)p_2h_i .
                                                        \tag{D.BC12}\label{maps:eq:D.BC12}
\]
The contribution $p_2\partial c_3$ is zero.
Similarly $P_2$ commutes with the differential
to degree one. Thus it induces a homomorphism
on $H_2$.

\item Both cone homologies have the short exact
sequence with left term the degree-two fibre
coinvariant and right term $K$. On the left
$P_2$ is the integral isomorphism induced by
$p_2$, and on the right it is the identity
induced by $p_1$. The short five lemma proves
that \eqref{maps:eq:D.BC10} is an integral isomorphism
$H_2(V)\simeq C\oplus K$. This supplies the
missing two-primary marking directly. Neither
division by two nor descent of an invariant
differential form is used.

\item \emph{Fixed sweeps, norm sweeps, and the last boundary.}
If $A_i v=v$ and $i<r$, the literal fixed-loop
torus is the bar cycle with $g_i=[v]$ and
all other entries zero. Formula \eqref{maps:eq:D.BC10} sends
it to $(f_{A_i}(v),J_iv)$.
For an arbitrary loop $\lambda$ with
$A_i^m\lambda=\lambda$, the pushforward of
its torus in the $m$-fold mapping-torus cover
has edge chain
\[
                        \sum_{\nu=0}^{m-1}[A_i^\nu\lambda].
                                                        \tag{D.BC13}\label{maps:eq:D.BC13}
\]
Its boundary telescopes to zero. Hence its
coordinates are the norm term in \eqref{maps:eq:D.BC4}
before subtracting $\lambda\wedge v$.

\item For the last boundary the base word is
$t_r=t_{r-1}^{-1}\cdots t_0^{-1}$.
The local bar one-chain $[u]$ maps to the
independent edge chains
\[
         -[A_{j+1}\cdots A_{r-1}A_ru],
                            \qquad 0\le j<r .
                                                        \tag{D.BC14}\label{maps:eq:D.BC14}
\]
This identity is at the chain level: applying
$\sum(A_j-I)$ makes the bars telescope to
$[A_ru]-[u]$, with no leftover fibre two-chain.
For a fixed vector, its linear part is exactly
the previous marked $J_rv$. Applying
\eqref{maps:eq:D.BC10} gives the second line of \eqref{maps:eq:D.BC3}.
Applying \eqref{maps:eq:D.BC14} to every summand of
\eqref{maps:eq:D.BC13} gives the full transported last-boundary
norm correction. In particular no reordering
of the last peripheral segments has been
silently discarded.

\item \emph{The actual cap maps.}
For a finite or regular logarithm the boundary
group has its normal meridian
$\mu=t_i^m v^{-1}$, with $A_i^m=I$ and
$A_i v=v$. The loops $\lambda,t_i^m,v$ lie
in one abelian subgroup. Bilinearity of its
torus classes therefore gives the exact
integral relation
\[
              [\lambda,\mu]=[\lambda,t_i^m]-[\lambda,v].
                                                        \tag{D.BC15}\label{maps:eq:D.BC15}
\]
The first term is \eqref{maps:eq:D.BC13}, and the second is
the fibre bivector $\lambda\wedge v$.
This proves the first four columns of \eqref{maps:eq:D.BC4}
even when their images have two-primary
torsion. The affine meridian generator gives
$[t_i,\mu]=[v,t_i]$, the fixed sweep already
computed, proving the fifth column. These
are exactly the full normal Thom generators
from \ref{maps:src:D.2.29}--\ref{maps:src:D.2.30}, including the torsion
normal-line cases.

\item Every fixed-normal torus at an unchanged cap
is still the same literal peripheral sweep,
so its image is \eqref{maps:eq:D.BC5}. The fibre kernels
are exactly the saturated lattices proved
from the explicit local models in \ref{maps:src:D.2.40}.
The possible original multiplicative
correction $\delta\wedge u_2$ is zero in
the full global fibre coinvariant by the
$172$ unit elliptic certificates \eqref{maps:eq:D.OW7}.
Thus the listed columns give the entire
marked kernel image at every allowed cap.

\item Finally the relative pair sequence \eqref{maps:eq:D.WP11}
has free final quotient for every weight;
its proof did not use the parity of $g$.
Its ranks depend only on the actual component
counts already proved in \ref{maps:src:D.2.37}. It follows
that the corrected full matrix computes
exactly the middle torsion, while the Seifert
group, free ranks and duality table stay
unchanged. All formulas are integral
polynomials, binomial polynomials, saturated
integer kernels and finite Smith operations.
They impose no upper bound on $d$, the
logarithmic orders or affine coordinates.
\end{enumerate}
\end{proof}

The executable
\path{code/III/manifolds_all_weighted_profiles.py}
implements \eqref{maps:eq:D.BC1}--\eqref{maps:eq:D.BC5}, records the full
quadratic coefficient vectors, and compares
its results with the independent inversion
method whenever $\gcd(a,d)$ is odd.
This theorem concerns translation iterates
of the finite natural initial catalogue.
It does not assert that all broader initial
linearizations or star logarithmic caps are
such iterates, and it does not infer an
uncomputed smooth classification from the
homology.


\begin{theorem}[the mixed six-, fifteen-, and thirty-torsion types]\label{maps:src:D.2.42}
\StatementRecord{maps:src:D.2.42}{proved}{Historical events: Claude Fable 5.2 (Anthropic), effort high, 16 November 2026: carried as background without fault found. GPT-6 Astra (OpenAI), ChatGPT Pro mode, effort Pro, S6\_09\_PAPERS, 17 November 2026: targeted proof check. Also checked in this writing. Pre-continuation writing events: GPT-6 Astra (OpenAI), ChatGPT Pro mode, effort Pro, S6\_10\_PAPERS, 18 November 2026: targeted proof check. Continuation metadata events: GPT-6 Astra (OpenAI), ChatGPT Pro mode, effort Pro, S6\_10\_PAPERS continuation, 19 November 2026: used and rechecked in this run; status-vocabulary reconciliation; metadata reconciliation only; no new mathematical check. \emph{Pre-continuation recorded proof scope.} GPT-6 Astra (OpenAI, ChatGPT Pro mode), effort Pro, 2026-11-17 UTC; Claude Fable 5.2 (Anthropic), effort high, 16 November 2026; also checked in this writing by GPT-6 Astra (OpenAI), ChatGPT Pro mode, effort Pro, 18 November 2026}{the hypotheses, notation and definitions stated here; \cref{maps:src:D.2.39}; \cref{maps:src:D.2.36}; \cref{maps:src:D.2.26}; \cref{maps:src:D.2.35}; \cref{maps:eq:D.OP4}; \cref{maps:eq:D.MT2}; \cite[equations (2.3), (2.5), p.267]{ffThomas56}; \cref{maps:eq:D.T37}; \cref{maps:eq:D.T38}; \cref{maps:eq:D.C27}; \cref{maps:src:D.2.8}; \cref{maps:src:D.1.10}; \cref{maps:eq:D.MP9}; \cref{maps:eq:D.MP10}; \cite[p.795]{manZhubr75}}


Consider the fourteen minimal patterns in \eqref{maps:eq:D.MT1}, with the primitive
Mumford cap at every multiplicative site, the indicated additive
choice, and no additional regular logarithmic transformation.
Allow every finite multiplicity, affine twist and ordinary clutch
$e\gamma$ producing a simply connected manifold. Let
$M_{3,0}$ be the explicitly framed surgery of \ref{maps:src:D.2.39}, and
$Z_{5,\epsilon}$ the oriented type of \ref{maps:src:D.2.36}. Define
\[
\begin{aligned}
 Y_6&=M_{\mathrm{conic}}\# M_{3,0},\\
 Y_{15,\epsilon}&=M_{3,0}\# Z_{5,\epsilon},\\
 Y_{30}&=M_{\mathrm{conic}}\# M_{3,0}\# Z_{5,0}.
\end{aligned}
\]
The orientation-preserving diffeomorphism type is given by the
following table; zero positions have the source's zero-based indexing.
\[
\begin{array}{c|r|r|l|l|l}
\text{OS/profile}&a&b&\text{zero positions}&\text{additive choice}&\text{type}\\\hline
56&1&6&0&\text{none}&Y_6\\
56&5&6&3,4&\text{none}&Y_6\\
56/II&1&6&2&II\text{ original or log}&Y_6\\
56&1&15&2&\text{none}&Y_{15,0}\\
56&2&15&3,4&\text{none}&Y_{15,1}\\
56/II&1&15&0&II\text{ original or log}&Y_{15,0}\\
56/III&1&30&3&III\text{ log}&Y_{15,0}\\
56/III&2&15&3&III\text{ original}&Y_{15,1}\\
56&1&30&3,4&\text{none}&Y_{30}.
\end{array}                                               \tag{D.MT1}\label{maps:eq:D.MT1}
\]
Thus the five $c=6$ patterns give one oriented type, the seven
$c=15$ patterns give two, and the two $c=30$ patterns give one.

\end{theorem}


\begin{proof}
\begin{enumerate}[label=\arabic*.,leftmargin=*]
\item \emph{The geometric input for every primary calculation.}
\ref{maps:src:D.2.26} gives $H_2=H_3=\mathbb Z/c$ and $b_3=0$ in these
simply connected cases. As $H^2(X;\mathbb Z)=0$, each is spin.
There are twelve distinct marked initial rows. Write
$\Omega=a\,u^*w^*+b\,\gamma^*\delta^*$ in their actual markings.
The exact certificate below verifies its invariance and verifies
that both the multiplicative angular columns and the full images
of $A_i-I$ have integral Hermite lattice $\ker\gamma^*$.
Thus, after omitting the additive caps and one ordinary disc,
the common piece $W$ has torsion-free group
$\mathbb Z_\gamma\times F_r$. The primitive Mumford boundary
surjections and the independent remaining peripheral loops give
$H_2(V)\twoheadrightarrow H_2(W)$, with free first homology.
Consequently the degree-two comparisons on $W$ used in
\ref{maps:src:D.2.35}--\ref{maps:src:D.2.36} apply to the actual marked classes here as well.

\item The three possible finite patterns have normal fibre coefficients
$6$, $15$, and $30/2=15$, respectively. All are divisible by
the stated residual $c$, as is the second normal coefficient,
a multiple of $b$. Every allowed original additive choice has
zero coefficient, and the clutch coefficient is $be$.
Hence the primitive fibre functional with zero swept values
descends modulo every prime dividing $c$, independently of
the allowed finite order and affine twist.

\item We record the local inputs to this comparison. The actual Tate
periods are $H=(n,ar),F=(-br,\rho)$; at the zero
$n\rho+ab r^2=1$. The doubled symmetric line \eqref{maps:eq:D.OP4} has
integral support
\[
          \psi=-a\rho x^2-2ab rxy+bn y^2+ax,
                                                        \tag{D.MT2}\label{maps:eq:D.MT2}
\]
with the two exact period differences $-2ax$ and
$2by+b\rho-ab r$. For an odd prime $p\mid c$, all its
vertical slopes are divisible by $p$, and half its signed
wall degree is the unit $-a(\rho-br)$ modulo $p$.
At a rank-one site write $H=\alpha u_0,F=\beta u_0$, with
$u_0=(u_1,u_2)$ primitive and $\gcd(\alpha,\beta)=1$.
Its integral doubled support is
$-(a u_1/\alpha)j(j-\alpha)$; the extracted identities prove
$\alpha\mid a u_1$ and give both prescribed period differences.
Thus the auxiliary line extends on all the specified caps,
with the actual boundary marking.

\item For each $p\mid c$ there is exactly one additional
$C_p$-fixed elliptic double curve $E_p$, at an $I_p$ site.
Its primitive short period is $(p,ar)$ and its doubled support
is $-ap j(j-1)$. The line therefore has zero $C_p$ character
on both the rational fixed curve $C$ and $E_p$, for odd $p$.
The normal weights at $E_p$ are $w,-w$, where
$w=(ar)^{-1}\pmod p$. All these statements, including the
uniqueness of each fixed elliptic curve, are verified using
the actual primitive period columns in the twelve rows.
The finite caps have free scalar action by primitivity of
their gamma numerators; the original additive caps have
free scalar action. No extra prime-isotropy stratum is omitted.

\item \emph{The full three-primary class.}
For every row take $u_3\in H^2(X;\mathbb F_3)$ normalized
to have degree one on $C$. It lifts to $S^1$-equivariant
cohomology because $H^1(X;\mathbb F_3)=0$ and
$H^3(BS^1;\mathbb F_3)=0$. Normalize its point character
on $C$ to zero. The torsion-free group of $W$ and the
zero characters of \eqref{maps:eq:D.MT2} then make its $C_3$ point
character on $E_3$ zero as well, exactly as in \ref{maps:src:D.2.39}.
The restrictions of this lift are a class pulled back from
$C$ and an external product of degree-one classes on
$BC_3\times S^1\simeq(E_3)_{S^1}$.

\item Thomas's exact odd-prime external-product formula and
odd-degree vanishing make the Pontryagin third power of
both restrictions zero in coefficients $\mathbb Z/9$
\cite[equations (2.3), (2.5), p.267]{ffThomas56}.
The complementary locally free circle action has all
stabilizers of order prime to three. Its Borel pair has
the cohomology of a five-dimensional orbifold pair with
these coefficients, so its relative degree-six group
is zero. Thus restriction to $C\sqcup E_3$ is injective
in equivariant degree six, as proved in \eqref{maps:eq:D.T37}--\eqref{maps:eq:D.T38}.
The full Pontryagin third power on $X$ is consequently zero.

\item The target calculation in part~1 of \ref{maps:src:D.2.39} and Zhubr's
Theorem~5 now show that the entire spin-bordism image of
$X\to K(\mathbb Z/3,2)$ is zero. This use concerns an
arbitrary classifying map: it does not assume that
$H_2(X)$ itself equals $\mathbb Z/3$. It determines the
hidden subgroup of order three in $\mathbb Z/9$, in
addition to the ordinary mod-three cube. The primary
class is therefore exactly the class of $M_{3,0}$.

\item \emph{The two- and five-primary classes.}
When $c$ is even, $a$ is odd and $b/2$ is odd in every
row of \eqref{maps:eq:D.MT1}. The unsquared line \eqref{maps:eq:D.C27} has the
monodromy-invariant semicharacter $a(u+w+uw)$.
Half of \eqref{maps:eq:D.MT2} and half of its rank-one support are
integral. Their scalar involution characters are one
on $C$ and zero on $E_2$, and their degree on $C$ is odd.
The common-piece comparison and the same simultaneous
Borel normalization as in \ref{maps:src:D.2.35} thus give
$\langle u_2^3,[X]_2\rangle=1$. Its entire two-primary
bordism class is the conic-surgery class of \ref{maps:src:D.2.8}.

\item For $c=15$ or $30$, the five-primary fixed curves have
the zero characters supplied by \eqref{maps:eq:D.MT2}. Normalize
$u_5[C]=1$. The actual localization of \ref{maps:src:D.2.36} gives
\[
 \langle u_5^3,[X]\rangle=0,\qquad
 \mathfrak p_5=(w^2-1)/2\pmod5.                           \tag{D.MT3}\label{maps:eq:D.MT3}
\]
The first equality detects the entire five-primary
fundamental class by \ref{maps:src:D.1.10}, and the second is the
complete refined Pontryagin class since four is a
unit modulo five. The actual weight is $w=4$ for
all $a=1$ rows, and $w=2$ for the $a=2$ rows.
Thus their complete bordism orbits are respectively
those of $Z_{5,0}$ and $Z_{5,1}$, by Zhubr's Theorem~5.

\item \emph{The complete oriented comparison.}
Iterate the prime splitting \eqref{maps:eq:D.MP9}--\eqref{maps:eq:D.MP10}.
Positive integral homology of the distinct primary
Eilenberg--Mac Lane factors has zero cross tensor and
Tor products. The wedge-to-product inclusion is therefore
a simply connected homology equivalence and hence a
homotopy equivalence. It identifies the full spin-bordism
class over $K(\mathbb Z/c,2)$ with the primary classes
just computed, through the actual Chinese remainder
projections. The automorphism group splits in the same way.

\item The stated connected sums have exactly these classes:
their natural maps obtained by collapsing the separating
spheres induce the direct-sum isomorphism on $H_2$, and
the spin one-handle bordism identifies connected sum
with addition. They are simply connected and spin, have
$H_2=\mathbb Z/c$ and $b_3=0$. The uniqueness assertion
of Zhubr's Theorem~3 \cite[p.795]{manZhubr75} consequently
gives each diffeomorphism in \eqref{maps:eq:D.MT1}. The two $c=15$
types are distinct because their five-primary refined
classes have different automorphism orbits.
\end{enumerate}
\end{proof}

The extraction program was run at residuals $6,15,30$.
The independent exact local certificate is
\path{code/III/manifolds_three_mixed_AH.py}; its output records
all fourteen patterns, twelve marked rows, $464$
semicharacter substitutions and all three finite normal
coefficients. The full three-primary value is justified
by the primary operation and the Borel comparison above,
not by those finite algebraic checks alone.


\begin{theorem}[a general triangular cusp and its full boundary kernel]\label{maps:src:D.2.43}
\StatementRecord{maps:src:D.2.43}{proved}{Historical events: Claude Fable 5.2 (Anthropic), effort high, 16 November 2026: carried as background without fault found. GPT-6 Astra (OpenAI), ChatGPT Pro mode, effort Pro, S6\_09\_PAPERS, 17 November 2026: targeted proof check. Also checked in this writing. Pre-continuation writing events: GPT-6 Astra (OpenAI), ChatGPT Pro mode, effort Pro, S6\_10\_PAPERS, 18 November 2026: targeted proof check. Continuation metadata events: GPT-6 Astra (OpenAI), ChatGPT Pro mode, effort Pro, S6\_10\_PAPERS continuation, 19 November 2026: used and rechecked in this run; status-vocabulary reconciliation; metadata reconciliation only; no new mathematical check. \emph{Pre-continuation recorded proof scope.} GPT-6 Astra (OpenAI, ChatGPT Pro mode), effort Pro, 2026-11-17 UTC; Claude Fable 5.2 (Anthropic), effort high, 16 November 2026; also checked in this writing by GPT-6 Astra (OpenAI), ChatGPT Pro mode, effort Pro, 18 November 2026}{the hypotheses, notation and definitions stated here; \cref{maps:src:D.2.22}; \cref{maps:src:E.2.19}; \cref{maps:eq:D.GC9}; \cref{maps:eq:D.GC11}; \cref{maps:eq:D.GC10}; \cref{maps:eq:D.GC8}; \cref{maps:eq:D.GC13}}


Let $N_L$ be the smooth same-fan triangular Mumford cusp
whose short-period lattice is a full sublattice
$L\subset\mathbb Z^2$. Write its Smith factors as
$s\mid t$, and put $D=st$. Holomorphic unit factors
in the two periods are allowed. Let $M_L=\partial N_L$,
and let $\Lambda=\mathbb Z^4$ be the regular torus
lattice. Its unipotent monodromy is $A=I+N$,
where $N^2=0$ and the map from the long lattice to
the primitive short lattice has Smith factors $(s,t)$.
Put $F=\ker(A-I)$, the primitive short lattice.

Then $\pi_1(N_L)=L\simeq\mathbb Z^2$, and its ordinary
integral homology is
\[
 \begin{array}{c|ccccc}
 j&0&1&2&3&4\\ \hline
 H_j(N_L)&\mathbb Z&\mathbb Z^2&
       \mathbb Z^{D+3}\oplus\mathbb Z/s&
       \mathbb Z^2&\mathbb Z^D .
 \end{array}                                             \tag{D.GC1}\label{maps:eq:D.GC1}
\]
All higher groups vanish. In particular $\chi(N_L)=2D$.
The full regular-fibre kernel is the generally
non-saturated lattice
\[
 \ker\{H_2(T^4)\to H_2(N_L)\}
       =(\Lambda^2A-I)\Lambda^2\Lambda+\Lambda^2F .
                                                        \tag{D.GC2}\label{maps:eq:D.GC2}
\]
Its image is $\mathbb Z^4\oplus\mathbb Z/s$, with
free cokernel of rank $D-1$ in $H_2(N_L)$.
The two actual fixed angular normal tori bound
and give the entire fixed-normal projection.
Together with \eqref{maps:eq:D.GC2} they specify the full
boundary kernel. More explicitly,
\[
 \begin{split}
 H_1(M_L)&=\mathbb Z^3\oplus\mathbb Z/s\oplus\mathbb Z/t,\\
 H_2(M_L)&=\mathbb Z^6\oplus(\mathbb Z/s)^2,\\
 \ker\{H_2(M_L)\to H_2(N_L)\}
             &\simeq\mathbb Z^2\oplus\mathbb Z/s .
 \end{split}                                             \tag{D.GC3}\label{maps:eq:D.GC3}
\]
The torsion killed by this map is the primitive
angular area modulo the monodromy image. The other
boundary torsion summand maps isomorphically to
the actual cap torsion.

There is also an exact component-meridian marking.
Use the positive $q$-meridian $\theta$ and index
the $D$ normalized components by $v\in\mathbb Z^2/L$.
The positive transverse disc to that component has
boundary
\[
                     \theta+\bar v
           \quad\hbox{in }\mathbb Z\theta\oplus\mathbb Z^2/L
                     \subset H_1(M_L).                  \tag{D.GC4}\label{maps:eq:D.GC4}
\]
These are the complete columns of the relative
map $H_2(N_L,M_L)=H^4(N_L)=\mathbb Z^D\to H_1(M_L)$.
In particular the component differences need not
vanish before global gluing.

\end{theorem}


\begin{proof}
\begin{enumerate}[label=\arabic*.,leftmargin=*]
\item \emph{The actual covering and its two-variable normalization.}
The extension of period units and the free same-fan
cover construction in \ref{maps:src:D.2.22} identify $N_L$ with
the finite cover of the minimal hexagonal cusp
associated to $L$. The units can also be removed
through the relative smooth family of \ref{maps:src:E.2.19}.
The normalized components are $dP_6$ surfaces,
the double curves are spheres, and the triple
strata are points. Their dual complex is
$\mathbb R^2/L$. Van Kampen gives
$\pi_1(N_L)=L$, because each component and every
intersection component is simply connected.

\item View the common universal cover with the full deck
action of the minimal cusp, and put
$R=\mathbb Z[x^{\pm1},y^{\pm1}]$ and
$I=(x-1,y-1)$. In the cyclic toric order
$(E_1,H-E_1-E_2,E_2,H-E_2-E_3,E_3,H-E_1-E_3)$,
the three opposite-pair translations have
weights $x^{-1},y,xy$. Indeed the primitive
fan rays are $(1,0),(0,1),(-1,1)$, and the
opposite pairing translates by their negatives;
choose the second period generator with the
corresponding sign. Simultaneously inverting
the deck variables changes neither the ideal
$I^2$ nor the assertions below.

\item With intersection form $\operatorname{diag}(1,-1,-1,-1)$,
the degree-two normalization matrix is
\[
 B=\begin{pmatrix}
 -x^{-1}&-1&-x^{-1}&-x^{-1}\\
 1&1&1&y\\
 -xy&-xy&-1&-xy
 \end{pmatrix}.                                         \tag{D.GC5}\label{maps:eq:D.GC5}
\]
The homology differential is its transpose with
the corresponding deck involution. One unit
row/column cancellation reduces its transpose to
\[
 F=\begin{pmatrix}
 1-x&0\\0&xy-1\\y-1&xy(y-1)
 \end{pmatrix}.                                         \tag{D.GC6}\label{maps:eq:D.GC6}
\]
For clarity there is an explicit Laurent-unimodular
comparison, not just equality of ranks or minors.
Let
\[
 T=\begin{pmatrix}
 0&0&-x^{-1}\\
 -x^{-1}y^{-2}&x^{-1}y^{-1}&-x^{-2}y^{-1}\\
 -y^{-1}&0&0
 \end{pmatrix},\qquad
 S=\begin{pmatrix}y&xy^2\\x&0\end{pmatrix}.
\]
Their determinants are Laurent units, and
\[
 T^{-1}F=
 \begin{pmatrix}-(y-1)&0\\x-1&-(y-1)\\0&x-1\end{pmatrix}S.
                                                        \tag{D.GC7}\label{maps:eq:D.GC7}
\]
The middle matrix is the exact presentation
of $I^2$ by the generators
$(x-1)^2,(x-1)(y-1),(y-1)^2$.
Its exactness follows directly from the regular
sequence $x-1,y-1$: the two displayed relations
are all the syzygies among those three monomials.
Thus
\[
                H_2(\widetilde N_L;\mathbb Z)\simeq I^2
                         \quad\hbox{as an }R\hbox{-module}.
                                                        \tag{D.GC8}\label{maps:eq:D.GC8}
\]
Here the universal cover is the same for every $L$.
The degree-zero normalization row is the cellular
complex of the plane, and the degree-four row is
$R$ in normalization degree zero. The degree-two
differential is injective over the domain $R$.
No higher differential is possible, because
normalization degree is at most two and the odd
homology of every stratum vanishes.

\item \emph{The integral finite-cover groups.}
Write $G=\mathbb Z^2/L$ and $R_G=\mathbb Z[G]$.
The degree-two normalization cokernel is
$I^2\otimes_R R_G$. We compute its torsion exactly.
The two short exact sequences
\[
 0\to I\to R\to\mathbb Z\to0,\qquad
 0\to I^2\to I\to\mathbb Z^2\to0
                                                        \tag{D.GC9}\label{maps:eq:D.GC9}
\]
use the trivial $R$-action on $I/I^2=\mathbb Z^2$.
The Koszul resolution for $x-1,y-1$ identifies
$\operatorname{Tor}^R_j(\mathbb Z,R_G)$ with
$H_j(L;\mathbb Z)$.
Consequently $I\otimes_R R_G$ is free of rank $D+1$:
it is an extension of the free augmentation ideal
of $\mathbb Z[G]$, of rank $D-1$, by $L\simeq\mathbb Z^2$.
Its map to $\mathbb Z^2$ in \eqref{maps:eq:D.GC9} is onto,
since $x-1$ and $y-1$ map to the two standard
generators. Its kernel is free of rank $D-1$.

\item The relevant Tor segment from the second sequence is
\[
 \mathbb Z\ \xrightarrow{\;\alpha_L\;}\ \mathbb Z^4
 \longrightarrow I^2\otimes_RR_G
 \longrightarrow I\otimes_RR_G
 \longrightarrow\mathbb Z^2\longrightarrow0 .
                                                        \tag{D.GC10}\label{maps:eq:D.GC10}
\]
The source is $H_2(L)$ and the target is
$H_1(L)\otimes\mathbb Z^2$.
If $u,v$ are an oriented basis of $L$, the map is
\[
                  \alpha_L(1)=u\otimes v-v\otimes u,
                                                        \tag{D.GC11}\label{maps:eq:D.GC11}
\]
where the second tensor factor uses the inclusion
$L\hookrightarrow\mathbb Z^2$.
One can check \eqref{maps:eq:D.GC11} directly in the Koszul
complex. In Smith coordinates
$L=\langle(s,0),(0,t)\rangle$, let
$N_x=1+\cdots+x^{s-1}$ and
$N_y=1+\cdots+y^{t-1}$.
The map on resolutions sends the relation
$(-(y-1),x-1)$ of $I$ to
$-e_y\otimes e_x+e_x\otimes e_y$.
The invariant cycle $N_xN_y$ therefore maps,
in $H_1(L)\otimes\mathbb Z^2$, to a vector with
nonzero coordinates $-s,t$.
Its Smith factor is exactly $s=\gcd(s,t)$.
The augmentation ideal and its square are invariant
under changes of basis of $\mathbb Z^2$, so this
also proves \eqref{maps:eq:D.GC11} for an arbitrary $L$.

\item Thus $\alpha_L$ is injective and its cokernel is
$\mathbb Z^3\oplus\mathbb Z/s$.
The free last kernel in \eqref{maps:eq:D.GC10} gives
\[
            I^2\otimes_RR_G
                    \simeq\mathbb Z^{D+2}\oplus\mathbb Z/s.
                                                        \tag{D.GC12}\label{maps:eq:D.GC12}
\]
The preceding Tor segment gives
$\operatorname{Tor}_1^R(I^2,R_G)=\mathbb Z^2$.
The degree-zero normalization row now has the
homology of the dual torus, adding one free
summand to total degree two and two to degree
one. The degree-four row contributes $\mathbb Z^D$.
These observations prove \eqref{maps:eq:D.GC1}; the extension
by the dual torus's degree-two $\mathbb Z$ splits
as an abelian-group extension.

\item \emph{The actual boundary map and the nonsaturated kernel.}
Choose bases with
$A b_1=b_1+s a_1$, $A b_2=b_2+t a_2$ and
$A a_i=a_i$. The exterior-square difference has
image generated by
\[
 s\,a_1\wedge a_2,\qquad
 s\,a_1\wedge b_2-t\,a_2\wedge b_1
                         +st\,a_1\wedge a_2 .
                                                        \tag{D.GC13}\label{maps:eq:D.GC13}
\]
Its two Smith factors are $(s,s)$.
The Wang sequence gives the first two boundary
groups in \eqref{maps:eq:D.GC3}.

\item The minimal cusp's two fixed angular normal
tori have actual section-disc bounding sweeps.
In the unitless toric chart at the dense locus
of the component indexed by zero, take the
section with its two nonvanishing torus
coordinates equal to one and its normal
coordinate equal to $q$. Acting on this
section by either short circle gives a map
$S^1\times D^2\to N$ whose boundary is the
literal fixed angular normal torus. Its
fundamental circle lies in a simply connected
normalized component, so its image in the
cusp fundamental group is zero. The whole
map therefore lifts to every period cover
$N_L$, with its marked boundary. Transport
through the relative unit-removing family
preserves these boundary classes. Thus the
two primitive fixed-normal projections still
bound, and their actual classes split off
the full Wang projection $H_2(M_L)\to F$.

\item The primitive angular area also maps to zero.
In the universal cover the lifted regular
fibre is $T^2\times\mathbb R^2$. Its degree-two
group is a trivial $R$-module. Its image in
the module $I^2$ of \eqref{maps:eq:D.GC8} is zero, since
$I^2$ has no nonzero element fixed by both
$x$ and $y$. This nullhomology projects to
every finite period cover.

\item Poincare--Lefschetz duality gives
\[
 H_3(N_L,M_L)=H^3(N_L)=\mathbb Z^2\oplus\mathbb Z/s,
 \qquad H_2(N_L,M_L)=H^4(N_L)=\mathbb Z^D .
                                                        \tag{D.GC14}\label{maps:eq:D.GC14}
\]
The already bounding angular normal tori exhaust
the rational rank of the boundary kernel.
Hence the four-dimensional rational fibre
coinvariant injects into the cap.
Moreover the cokernel of
$H_2(M_L)\to H_2(N_L)$ injects into the free
group $H^4(N_L)$, so it is free.
Because the fixed-normal projections bound,
the boundary image equals the regular-fibre image.
It therefore contains all of the cap torsion.
There is no hidden free-to-torsion extension
here: if a fibre-coinvariant element maps to
a torsion class, rational injectivity forces
that element itself to be torsion. Hence
the cap torsion is the image of the boundary
fibre torsion subgroup, not of a free class
whose rational image has been overlooked.

\item The two boundary fibre torsion summands in
\eqref{maps:eq:D.GC13} have order $s$. The angular-area
one is killed. The other must map onto
$\operatorname{Tor}H_2(N_L)=\mathbb Z/s$ and
hence maps isomorphically. No free fibre
class can enter the kernel, since the rational
fibre coinvariant injects. This proves exactly
\eqref{maps:eq:D.GC2}, not its saturation, and gives all
the remaining assertions in \eqref{maps:eq:D.GC3}.
The free cokernel has rank $(D+3)-4=D-1$.

\item Finally the smooth fan has primitive rays
$(v,1)$. A positive small normal disc to
the associated component has this primitive
angular-plus-base meridian. Reducing the
angular coordinate modulo the period lattice
gives \eqref{maps:eq:D.GC4}, with the complex orientation
fixing the common sign. Every coset of $L$
occurs once. Their differences generate
$\mathbb Z^2/L$, and the zero coset gives
$\theta$; thus the displayed columns give
the entire kernel of $H_1(M_L)\to H_1(N_L)$.
\end{enumerate}
\end{proof}

The complete Laurent-unimodular operations,
finite HNF-lattice normalization matrices,
boundary exterior-square matrices and all
component-meridian columns are implemented in
\path{code/III/manifolds_general_cusp.py}.
When $s=1$, the fibre kernel is saturated
and the result recovers the local input used
in \ref{maps:src:D.2.34}--\ref{maps:src:D.2.41}. When $s>1$, replacing \eqref{maps:eq:D.GC2}
by its saturation would incorrectly erase
the actual cap torsion $\mathbb Z/s$.


\begin{theorem}[the seven semistable residual-four and residual-eight smooth types]\label{maps:src:D.2.44}
\StatementRecord{maps:src:D.2.44}{proved}{Historical events: Claude Fable 5.2 (Anthropic), effort high, 16 November 2026: carried as background without fault found. GPT-6 Astra (OpenAI), ChatGPT Pro mode, effort Pro, S6\_09\_PAPERS, 17 November 2026: targeted proof check. Also checked in this writing. Pre-continuation writing events: GPT-6 Astra (OpenAI), ChatGPT Pro mode, effort Pro, S6\_10\_PAPERS, 18 November 2026: targeted proof check. Continuation metadata events: GPT-6 Astra (OpenAI), ChatGPT Pro mode, effort Pro, S6\_10\_PAPERS continuation, 19 November 2026: used and rechecked in this run; status-vocabulary reconciliation; metadata reconciliation only; no new mathematical check. \emph{Pre-continuation recorded proof scope.} GPT-6 Astra (OpenAI, ChatGPT Pro mode), effort Pro, 2026-11-17 UTC; Claude Fable 5.2 (Anthropic), effort high, 16 November 2026; also checked in this writing by GPT-6 Astra (OpenAI), ChatGPT Pro mode, effort Pro, 18 November 2026}{the hypotheses, notation and definitions stated here; \cref{maps:src:D.2.8}; \cref{maps:src:D.2.26}; \cref{maps:src:D.2.41}; \cref{maps:src:D.2.35}; \cite[Sections 3.11--3.14]{manZhubr75}; \cite[Theorem 5 and Sections 3.11--3.14]{manZhubr75}; \cite{ffZhubr13}; \cite[pp.369--372]{ffBrown72}; \cite[Section 5]{ffZhubr13}; \cite[Definition 1.7, Remark 1.8, Theorem 3.1]{mapsKaplan79}; \cite[Theorem 5]{manZhubr75}; \cite[Theorem 3]{manZhubr75}}


Consider the default semistable rows OS55 with $(a,b)=(5,4)$ and zero
index $2,3,4$, and OS45 with $(a,b)=(1,8)$ and zero index $1,2,3,4$.
Use the primitive Mumford cap at every multiplicative site, no additional
regular logarithmic transforms, and either simply connected ordinary
clutch $e\gamma$, $e=\pm1$. Write $c=b$.
The resulting oriented six-manifolds have
\[
 H_2(X;\mathbb Z)=H_3(X;\mathbb Z)=\mathbb Z/c,
 \qquad \lambda(X)=0,\qquad \mu(X)=0,
 \qquad
 P(X)=\begin{cases}2,&\mathrm{OS55},\\0,&\mathrm{OS45}.
 \end{cases}                                                   \tag{D44.1}\label{maps:eq:D44.1}
\]
Here $P$ is Zhubr's refined quarter-Pontryagin class, transported to
$\mathbb Z/c$ by any generator; both displayed values are independent
of that choice. In particular all three OS55 manifolds are oriented
diffeomorphic to $Z_{4,1}$, and all four OS45 manifolds are oriented
diffeomorphic to $Z_{8,0}$, with the following concrete surgery models.
The manifold $Z_{c,0}$ is surgery on an embedded two-sphere representing
$c$ times the generator of $H_2(S^2\times S^4)$.
The manifold $Z_{c,1}$ is surgery in
$M_{\rm conic}\#(S^2\times S^4)$ on an embedded two-sphere representing
$(t,-c/2)\in\mathbb Z/2\oplus\mathbb Z$, where $t$ is the generator
of the conic-surgery manifold of \ref{maps:src:D.2.8}. These are spin surgeries with
their unique normal framings. This theorem makes no assertion about
the remaining additive profiles or additional regular logarithms.

\end{theorem}
\begin{proof}
\begin{enumerate}[label=\arabic*.,leftmargin=*]
\item The marked homology calculation \ref{maps:src:D.2.26} and its all-weight verification
\ref{maps:src:D.2.41} give the stated integral groups, with no other reduced homology
except $H_6=\mathbb Z$. The full integral support-function and
semicharacter certificate
\path{code/III/fourfolds_D2_2power_AH.py}
verifies the integral line-descent and actual-angle hypotheses
used in \ref{maps:src:D.2.35} for these rows. Here are the complete remaining
localization steps. The subgroup $G=C_c$ of the scalar circle has
fixed locus the rational curve $C$ and the one exceptional
elliptic double curve $E$ at $I_c$; it acts freely off them.
The actual normal characters are $(1,-1)$ on $C$ and $(w,-w)$
on $E$, with $w^2=1$ in $\mathbb Z/c$. The two normal lines on
$C$ have degree $-1$, and those on $E$ have degree zero.
The symmetric auxiliary line of \ref{maps:src:D.2.35} has character $c/2$ on
$C$, zero on $E$, and degree coprime to $c$ on $C$.
Its identification with the global mod-$c$ class is the same
integral common-piece argument as there: the boundary-to-cap
surjections and primitive degree-one map make $H_2(W)\to H_2(X)$
onto; the line has zero swept values and kills the ordinary
clutch relation modulo $c$. Moreover $\pi_1(W_G)$ has the
unique torsion subgroup $G$. The difference between equivariant
lifts therefore has the same point character on $C$ and $E$.
After one base-character adjustment, the stated auxiliary
characters are the characters of a global equivariant lift.

\item Normalize its ordinary class $\alpha\in H^2(X;\mathbb Z/c)$
to have degree one on $C$. Put $t\in H^2(BG;\mathbb Z/c)$,
$x=[C]^*$ and $y=[E]^*$. Its restrictions have the form
\[
 \alpha|_C=x+(c/2)t,\qquad
 \alpha|_E=A y+\eta\zeta,\qquad
 e_G(N_C)=-t^2,\qquad e_G(N_E)=-w^2t^2,                       \tag{D44.0a}\label{maps:eq:D44.0a}
\]
where $\eta$ has degree one on $BG$ and $\zeta$ lifts from
integral degree-one cohomology of $E$. Localization over
$\mathbb Z/c$ is legitimate: the complement is free, so its
cohomology vanishes after inverting $t$. The degree-one
localization identity gives $A=-w^2$. The elliptic restriction
has cube zero, because $y\zeta=0$ and $\zeta^2=0$; the rational
curve's contribution to the cube is $-3(c/2)^2=0\pmod c$.
Thus $\langle\alpha^3,[X]\rangle=0$ in the full coefficient
ring, not merely modulo two.
The equivariant complex determinant is trivial: its ordinary
class is zero and its total character on $C$ is zero.
Hence $\lambda_G=-c_{2,G}$. On either fixed curve,
$c_{2,G}(TX)=e_G(N)$, since the square of the tangent-line
class vanishes. Therefore
\[
 \langle\alpha\lambda,[X]\rangle=-(1+A)=-(1-w^2)=0.           \tag{D44.0b}\label{maps:eq:D44.0b}
\]
The normalized class detects the cyclic degree-four group, so
$\lambda=0$. Zhubr's Sections 3.11--3.14 identify the complete
group $H_6(K(\mathbb Z/c,2);\mathbb Z)=\mathbb Z/c$ by the cube,
giving $\mu=0$ \cite[Sections 3.11--3.14]{manZhubr75}.
The free-coefficient calculation
and coefficient surjectivity in Zhubr's Sections 3.11--3.14 give
$2P=\operatorname{PD}(\lambda)$
\cite[Theorem 5 and Sections 3.11--3.14]{manZhubr75}.
Thus the two remaining possibilities
are exactly $P=0,c/2$. The calculation below distinguishes them by
the actual tangent framing, without halving a congruence or using
the conjectural change-of-three-frame formula of \cite{ffZhubr13}.

\item \emph{A framing parity that depends only on the three-cycle.}
We first record the precise primary construction used to compare a
regular three-torus with a torsion three-sphere.
Let a closed oriented six-manifold $Y$ have a stable tangent framing.
For an embedded oriented three-manifold $N\subset Y$ with a chosen
normal framing, denote by $[N]_{\rm fr}\in\Omega^{\rm fr}_3
\cong\mathbb Z/24$ its induced stable framed-bordism class. Then
\[
 [N]_{\rm fr}\pmod2
   =q_Y\bigl(\operatorname{PD}_{\mathbb F_2}[N]\bigr)             \tag{D44.2}\label{maps:eq:D44.2}
\]
for a function of the indicated mod-two cohomology class alone.
This statement includes torsion classes in integral homology.
Here is a derivation directly from Brown's construction
\cite[pp.369--372]{ffBrown72}. Set $K_3=K(\mathbb F_2,3)$ and let
$j:S^3\to K_3$ represent its fundamental class. Brown's $F(u,0)$
is literally the stabilization of the map $u:Y\to K_3$.
The framed Pontryagin--Thom collapse $Y\to S^3$, followed by $j$,
therefore represents $F(\operatorname{PD}[N],0)$.
Apply the framed duality map $W_*A_\alpha$ of his Theorem 1.7 and
Corollary 1.11. Its value is the image of $[N]_{\rm fr}$ under
\[
 j_*:\pi_3^{\rm S}\longrightarrow\pi_6^{\rm S}K_3
          \cong\mathbb Z/2.                                  \tag{D44.3}\label{maps:eq:D44.3}
\]
The target is Brown's Lemma 1.2. The remark immediately following
his Corollary 1.13 identifies $j_*\nu$ with its nonzero generator:
the coefficient is the mod-two Hopf invariant, and the quaternionic
Hopf element $\nu$ has Hopf invariant one. Dividing Brown's
$\mathbb Z/4$-valued function, whose values here are $0,2$, by two
gives \eqref{maps:eq:D44.2}. This is a Pontryagin--Thom argument about the specified
normal framing; no compression of an arbitrary framed torus to a
sphere is being assumed.

\item For clarity we also fix the sphere convention. The homomorphism
$\pi_3SO(3)\to\pi_3SO(6)$ has image $2\mathbb Z$.
Consequently changing the normal framing of an embedded $S^3$
does not change the parity of its induced stable framed class.
For the standard sphere bounding a ball, the induced stable framed
class is zero and its tangent three-frame represents the nonzero
class in $\pi_3V_{6,3}=\mathbb Z/2$ relative to a constant ambient
three-frame. Equivalently, the tangent Lie framing on $S^3$ has
class $\nu$, and a change of ambient frame by degree $d$ changes
that class by $d\nu$, while the relative three-frame class changes
by $d\bmod2$. Thus Zhubr's required nontrivial relative
three-frame is equivalent to \emph{even} induced stable framed
class on the sphere. These conventions agree with the quaternionic
calculation in \cite[Section 5]{ffZhubr13}.

\item Since $\lambda(X)=0$, $TX$ is stably trivial: the only obstruction
for a spin bundle on a six-dimensional complex is its class in
$H^4(-;\pi_3\mathrm{Spin})$, namely $\lambda$.
Moreover $TX$ is actually trivial on $X^\circ=X\setminus\{p\}$.
For example its complex bundle has $c_1=c_2=0$ and its last
obstruction lies in $H^6(X^\circ;\mathbb Z)=0$.
Any framing there extends stably across $p$, since
$\pi_5SO=0$. All such full framings have the same relative
three-frame parity on an embedded torsion $S^3$: their difference
lifts to $\mathrm{Spin}(6)$ and its degree-three class is a class
of $H^3(X^\circ;\mathbb Z)$, which evaluates to zero on a torsion
three-cycle.
An embedded sphere in the integral class of the regular torus
$T=\langle e_1,e_2,\gamma\rangle$ used below can also be
constructed explicitly, so its existence is not an additional
representability assumption. Apply Kaplan's construction to
the spin structure induced on the normally framed $T$:
it gives a four-dimensional filling with one zero-handle and
only two-handles, preserving that almost framing
\cite[Definition 1.7, Remark 1.8, Theorem 3.1]{mapsKaplan79}.
Delete the zero-handle and reverse the cobordism. Realize the
resulting two-handle attachments as ambient surgeries in
$X^\circ$. Their attaching circles bound embedded disks
with interiors disjoint from $T$ and from one another, since
$\pi_1(X^\circ)=0$, $2+3<6$ and $2+2<6$.
Prescribe a boundary collar along the first normal vector.
The two tangent framing directions of each attachment and
the other two normal directions give a four-frame around
the boundary of its disk. Its extension obstruction in
$\pi_1SO(4)=\mathbb Z/2$ is exactly the spin obstruction
of that two-handle, and vanishes. Thickening these disks
therefore performs the prescribed surgeries inside $X^\circ$.
The resulting embedded $S^3$ has the same integral homology
class as $T$. Brown's identity already supplies its parity;
this construction additionally verifies its geometric
realization without changing the prescribed spin structure.
Zhubr's Theorem 1 now implies
\[
 P(X)=0\quad\Longleftrightarrow\quad
 \text{a full ambient framing is admissible on the torsion }S^3.
                                                               \tag{D44.4}\label{maps:eq:D44.4}
\]
Indeed the forward implication is not an assumption about arbitrary
three-frames. Take the admissible three-frame supplied by that
theorem. If $P=0$, its spin rank-three complement has vanishing
primary obstruction $c_2$; the only possible further obstruction
on the punctured six-manifold is in $H^5(X^\circ;\mathbb Z/2)=0$.
The complement is therefore trivial, and the three-frame extends
to a full framing. Conversely an admissible full framing has
trivial complement and hence refined class zero. Independence
of the full framing proves \eqref{maps:eq:D44.4} for the particular framing
used below.

\item \emph{The reference frame and the actual critical-cycle defect.}
Write $f:X\to\mathbb P^1$ for the torus fibration and
$\mathcal V=\ker(df)$ for its holomorphic differential kernel.
This is a locally free rank-two sheaf, including at the critical
curves. In the nodal chart $f=xy$ it has basis
$x\partial_x-y\partial_y,\partial_z$; in the triangular cusp
chart $f=xyz$ it has basis
$x\partial_x-y\partial_y,y\partial_y-z\partial_z$.
The scalar $\delta$-action defines a nowhere-zero section of this
\emph{abstract} vector bundle, although its image as a vector
field vanishes on the fixed curve. In a toric chart its two
coefficients are those of the nonzero primitive scalar lattice
vector, so they cannot both vanish. This description is unchanged
by the translation clutch.
As $H^2(X;\mathbb Z)=0$, the quotient line $\mathcal V/\mathcal O\delta$
is topologically trivial; hence $\mathcal V$ has a global smooth
frame whose first vector is $\delta$. The line
$\mathcal L=f^*T\mathbb P^1$ is also topologically trivial.

\item Let $Z$ denote the scheme defined by the coefficients of $df$.
There is an exact sequence
\[
 0\longrightarrow\mathcal V\longrightarrow TX
  \xrightarrow{df}\mathcal L
  \longrightarrow\mathcal L\otimes\mathcal O_Z
  \longrightarrow0.                                         \tag{D44.5}\label{maps:eq:D44.5}
\]
Its codimension-two cycle is the literal sum $C_{\rm crit}$ of
the complex double curves, each with coefficient one. This
includes the three axes at a triple point: the ideal
$(xy,xz,yz)$ has multiplicity one on each axis. Since
$\mathcal V$ and $\mathcal L$ are topologically trivial,
\eqref{maps:eq:D44.5} gives $c_2(TX)=[C_{\rm crit}]$. Thus this cycle is
integrally nullhomologous.

\item Off $Z$, splitting \eqref{maps:eq:D44.5} identifies $TX$ with the framed
bundle $\mathcal V\oplus\mathcal L$. Its comparison with a
global frame on $X^\circ$ has an integral degree-three class
$a$ whose local residue on every smooth double curve is one,
up to the same overall sign. In the model $xy$ this is the
explicit Bott matrix
\[
 \frac1{\sqrt{|x|^2+|y|^2}}
 \begin{pmatrix}x&\overline y\\-y&\overline x\end{pmatrix}
       \in SU(2).                                            \tag{D44.6}\label{maps:eq:D44.6}
\]
It has degree one on the normal three-sphere; adding the
tangent line to the critical curve sends this to the stable
generator, not twice that generator. The latter factor two
belongs to the different map $SO(3)\to SO$, as is also seen
from $\lambda=-c_2$ for this complex rank-two normal calculation.
Removing the finitely many triple points does not affect this
degree-three residue calculation. The integral relative
cohomology sequence, or equivalently the Thom intersection
construction, therefore gives
\[
 \langle a,[T]\rangle
      =\pm\operatorname{lk}(T,C_{\rm crit})                   \tag{D44.7}\label{maps:eq:D44.7}
\]
for any regular three-torus $T$ disjoint from the critical
cycle. The ambiguity of adding a class from $H^3(X^\circ;\mathbb Z)$
evaluates to zero on our torsion three-cycle. Thus the integer
in \eqref{maps:eq:D44.7}, not merely a residue modulo $c$, is well defined.

\item Take specifically $T=\langle e_1,e_2,\gamma\rangle$ in a regular
fibre inside the ordinary clutch disc. The reference frame
$\mathcal V\oplus\mathcal L$ has zero degree-three difference
from the flat translation frame there. Indeed the vertical
frames have the same first vector $\delta$; their change lies
in the stabilizer $\mathbb C\rtimes\mathbb C^*$, which retracts
to a circle. The base line contributes another $\mathbb C^*$,
and changing the horizontal splitting contributes a contractible
space. None contributes a degree-three class. The induced flat
framing of $T^3$ is its product Lie framing, with stable class
$\eta^3=12\nu$, an even element of $\mathbb Z/24$.
A possible determinant-circle transition cannot change this
parity. A map $T^3\to S^1$ is homotopic to a character and,
after an integral change of torus basis, depends on just one
circle. Its real rotation can be placed stably in that circle
and one normal direction. It changes just that circle's framed
class, so the resulting product-framed torus has class either
$0$ or $\eta^3$, both even. After splitting off the determinant,
the remaining $SU(3)$ transition is a degree-$k$ map on a
three-dimensional complex and changes the framed class by
$\pm k\nu$. Thus its exact parity is the integer in \eqref{maps:eq:D44.7}.
Equations \eqref{maps:eq:D44.2} and \eqref{maps:eq:D44.7} consequently give
\[
 q_X(\operatorname{PD}[T])
       =\operatorname{lk}(T,C_{\rm crit})\pmod2.              \tag{D44.8}\label{maps:eq:D44.8}
\]
The torus represents an odd generator of $H_3(X)=\mathbb Z/c$
by the explicit linking calculation below. Its reduction is
therefore the same as that of an embedded torsion generator
$S^3$. Combine \eqref{maps:eq:D44.4} and \eqref{maps:eq:D44.8} to obtain the actual decision
\[
 P(X)=\frac c2
       \bigl(\operatorname{lk}(T,C_{\rm crit})\bmod2\bigr).
                                                               \tag{D44.9}\label{maps:eq:D44.9}
\]

\item \emph{The integral linking calculation.}
All local preimages and bounding coefficients here are those
of the actual marked caps, not an invariant-form calculation
on an unmarked torus. At each site choose the recorded
determinant-one elliptic basis $(u_1,u_2)$ and use the short
and long planes $(u_1,\delta)$ and $(u_2,\gamma)$. The exact
valuation matrix is
\[
 L=\begin{pmatrix}n&-bH\\aH&R\end{pmatrix}.                  \tag{D44.10}\label{maps:eq:D44.10}
\]
At a rank-one site write $L=\alpha r$ with $\alpha$ primitive.
Choose a short complement $\sigma$ with
$\det(\alpha,\sigma)=1$ and long vectors $(\beta,\ell)$ with
$r\beta=\rho>0$, $r\ell=0$, and $\det(\beta,\ell)=1$.
The literal transverse complex elliptic curve has regular-fibre
preimage $\sigma\wedge\ell$. The product nodal chart supplies
the specialization chain with degree one. These signs follow
by orienting the vanishing elliptic factor as $(\alpha,\beta)$
and the transverse factor as $(\sigma,\ell)$.

\item At the unique rank-two cusp $\det L=1$. Put
$(b_1,b_2)=(u_2,\gamma)L^{-1}$ and $(a_1,a_2)=(u_1,\delta)$,
so $Nb_j=a_j$. The three integral support functions
$r(r-1)/2,rs,s(s-1)/2$ have degree matrix
\[
 \begin{pmatrix}0&-1&0\\1&1&-1\\-1&0&0\end{pmatrix}          \tag{D44.11}\label{maps:eq:D44.11}
\]
on the three complex opposite-pair double curves in cyclic
ray order $(1,0),(0,1),(-1,1)$. This unimodular matrix gives
degree vector $(-1,1,-1)$ for their sum. The fibre vector
\[
 a_1\wedge b_1+b_1\wedge a_2+a_2\wedge b_2                  \tag{D44.12}\label{maps:eq:D44.12}
\]
has precisely those degrees and zero long--long component.
The latter is the degree-two class of the dual torus, so these
four coordinates identify its specialization with the actual
sum of the three double curves integrally. In particular no
undetected fourth homology component is being suppressed.

\item Transport the five preimages to the common regular fibre along
the fixed base paths. Their sum is a representative of
$C_{\rm crit}$ in the exact marked presentation of \ref{maps:src:D.2.41}.
All transport and cap relations can be realized by chains
outside the ordinary disc. In that disc the four normal
Thom chains are the products of a lattice loop with the
disc, and the redundant fifth chain has loop vector
$(0,0,0,-e)$. The intersection of a bounding relation with
$T=e_1e_2\gamma$ is consequently its clutch $\delta$
coefficient. Integral Smith solutions give
\[
 \begin{array}{c|c|c|c}
  \mathrm{OS}&c&\text{zero indices}&
       \operatorname{lk}(T,C_{\rm crit})\text{ for }e=1\\\hline
  45&8&1,2,3,4&4,4,4,4\\
  55&4&2,3,4&37,7,27.
 \end{array}                                                  \tag{D44.13}\label{maps:eq:D44.13}
\]
For example, in OS55 at zero index $2$, the five cap vectors
in ordered bivector coordinates
$(e_1e_2,e_1\delta,e_1\gamma,e_2\delta,e_2\gamma,\delta\gamma)$
are
\[
 \begin{gathered}
 (0,0,0,-12,0,5),\quad(3,0,1,-12,0,4),\\
 (-19,-23,-10,0,5,6),\quad(0,0,0,0,0,1),\quad(0,4,0,0,0,1).
 \end{gathered}                                               \tag{D44.14}\label{maps:eq:D44.14}
\]
Their sum has an integral bounding relation with clutch loop
$(0,0,37,0)$. The same Smith calculation for $c$ times a
two-dimensional torsion generator gives a clutch coefficient
coprime to $c$ in every row. Nondegeneracy of the torsion
linking pairing proves that $T$ is an odd generator of
$H_3(X)$. Every integral relation among presentation columns
has zero clutch $\delta$ coefficient, so the displayed
linking does not depend on a selected Smith solution.
Changing the stated short complements or the cusp mixed
preimage by an actual cap-kernel class also changes that
coefficient by zero. For $e=-1$ the exact same calculation
with its actual clutch vector reverses the displayed integer
signs and retains the unit torsion-linking test.

\item The executable
\path{code/III/manifolds_even_critical_cycles.py}
retains all fourteen rows, the full integral coefficient
vectors, local bases, alternative preimages and both clutch
signs in
\path{records/manifolds_even_critical_cycles.json}.
Its integral computation is independent of the differential
and framing arguments above. Formula \eqref{maps:eq:D44.9} and table
\eqref{maps:eq:D44.13} prove \eqref{maps:eq:D44.1}.
The separate executable
\path{code/III/fourfolds_D2_frame_certificate.py}
checks the local kernel and Bott-matrix identities, the
unimodular cusp degree matrix and, for all fourteen marked
rows, the independent degree identity
$\Omega(C_{\rm crit})=-be\,\operatorname{lk}(T,C_{\rm crit})$
for the actual pure fibre functional
$\Omega=a e_1^*e_2^*+b\gamma^*\delta^*$.
It retains its inputs' hash and every local degree in
\path{records/fourfolds_D2_frame_certificate.json}.

\item \emph{Recognition of the surgery models.}
The normal rank-four bundle of any embedded sphere in these
spin six-manifolds is trivial, and its framing is unique
since $\pi_2SO(4)=0$. The stated surgeries have $H_2=H_3=\mathbb Z/c$
and zero free third Betti number. The classifying map to
$K(\mathbb Z/c,2)$ extends across each surgery trace.
The model $Z_{c,0}$ has refined class and fundamental-class
invariant both zero. For $Z_{c,1}$, the conic summand's refined
class maps to $c/2$, and the complete cubic maps to zero for
$c=4,8$, since $(c/2)^3=0\pmod c$. Thus their full Zhubr
invariants are respectively $(P,\mu)=(0,0)$ and $(c/2,0)$,
using naturality in \cite[Theorem 5]{manZhubr75} and the full
conic invariant established in \ref{maps:src:D.2.8}.
The classification theorem \cite[Theorem 3]{manZhubr75}, including
its finite invariants, gives the asserted oriented diffeomorphisms. No conclusion
here rests on ordinary homology alone.\hfill$\square$
\end{enumerate}
\end{proof}


\begin{theorem}[the remaining minimal residual-four, eight and twenty types]\label{maps:src:D.2.45}
\StatementRecord{maps:src:D.2.45}{proved}{Historical events: Claude Fable 5.2 (Anthropic), effort high, 16 November 2026: carried as background without fault found. GPT-6 Astra (OpenAI), ChatGPT Pro mode, effort Pro, S6\_09\_PAPERS, 17 November 2026: targeted proof check. Also checked in this writing. Pre-continuation writing events: GPT-6 Astra (OpenAI), ChatGPT Pro mode, effort Pro, S6\_10\_PAPERS, 18 November 2026: targeted proof check. Continuation metadata events: GPT-6 Astra (OpenAI), ChatGPT Pro mode, effort Pro, S6\_10\_PAPERS continuation, 19 November 2026: used and rechecked in this run; status-vocabulary reconciliation; metadata reconciliation only; no new mathematical check. \emph{Pre-continuation recorded proof scope.} GPT-6 Astra (OpenAI, ChatGPT Pro mode), effort Pro, 2026-11-17 UTC; Claude Fable 5.2 (Anthropic), effort high, 16 November 2026; also checked in this writing by GPT-6 Astra (OpenAI), ChatGPT Pro mode, effort Pro, 18 November 2026}{the hypotheses, notation and definitions stated here; \cref{maps:src:D.2.30}; \cref{maps:src:D.2.44}; \cref{maps:src:D.2.36}; \cref{maps:src:D.2.26}; \cref{maps:src:D.2.41}; \cite[Definition 1.7, Remark 1.8 and Theorem 3.1]{mapsKaplan79}; \cite{ffZhubr13}; \cite[Section 5]{ffZhubr13}; \cite[Sections 3.11--3.14]{manZhubr75}; \cref{maps:src:D.2.38}; \cite{manZhubr75}}


Retain the primitive Mumford filling at every multiplicative site,
no additional regular logarithmic transformations, and the indicated
choice at the type-$II$ site in the following eleven patterns:
\[
\begin{array}{c|c|c|c}
 \mathrm{OS/profile}&(a,b)&\text{zero indices}&\text{type-}II\text{ choice}\\\hline
 45/II&(1,8)&2,3&\text{original or good log}\\
 55/II&(5,4)&3&\text{original or good log}\\
 55/\mathrm{default}&(1,20)&2,3,4&\text{none}\\
 55/II&(1,20)&3&\text{original or good log}.
\end{array}                                                    \tag{D45.1}\label{maps:eq:D45.1}
\]
Indices are zero-based. In an original-only pattern the ordinary
clutch is $e\gamma$ with $e=\pm1$. At a good $II$ logarithm allow
every order $m=6k$, $k>0$, every integer ordinary clutch $e$, and
every power vector
\[
 v=r\delta+n\ell,\qquad r\in\mathbb Z,
 \qquad n-me=\sigma\in\{1,-1\},                             \tag{D45.2}\label{maps:eq:D45.2}
\]
where $\ell$ is the actual primitive fixed lift with
$\gamma^*(\ell)=1$. These are exactly the simply connected data
in this one-log sector; they are free by \ref{maps:src:D.2.30}.
All the resulting manifolds have $H_2=H_3=\mathbb Z/b$ and
zero free third Betti number. Their full Zhubr invariants, in
the primitive fibre coordinate specified below, are
\[
 \begin{array}{c|c|c|c|c}
 b&P&\lambda&\mu&\text{oriented diffeomorphism type}\\\hline
 4&2&0&0&Z_{4,1}\\
 8&0&0&0&Z_{8,0}\\
 20&14&8&0&Z_{4,1}\#Z_{5,1}.
 \end{array}                                                 \tag{D45.3}\label{maps:eq:D45.3}
\]
Here $Z_{4,1},Z_{8,0}$ are the explicit framed two-sphere surgery
models of \ref{maps:src:D.2.44}, and $Z_{5,1}$ is the nonzero refined-class,
zero fundamental-class model of \ref{maps:src:D.2.36}. In particular \ref{maps:src:D.2.44}--\ref{maps:src:D.2.45}
identify all eighteen minimal patterns of residual order
$4,8,20$, including every simply connected finite multiplicity
and twist in \eqref{maps:eq:D45.2}. No assertion about additional regular
logs or higher translation weights is included.

\end{theorem}
\begin{proof}
\begin{enumerate}[label=\arabic*.,leftmargin=*]
\item The homology and fundamental-group statements are \ref{maps:src:D.2.26},
\ref{maps:src:D.2.30} and \ref{maps:src:D.2.41}. In particular $H^2(X;\mathbb Z)=0$, so
$c_1(X)=0$ and $X$ is spin. The differential construction of
\ref{maps:src:D.2.44} extends to the stated caps with two precise changes.
At an original type-$II$ elliptic cusp, the smooth total
surface has base map
\[
 q=y^2-x^3,\qquad
 \ker(df)=\mathcal O(2y,3x^2,0)\oplus\mathcal O(0,0,1).
                                                               \tag{D45.4}\label{maps:eq:D45.4}
\]
To see that this is the map, not just the singular zero set,
solve the Weierstrass equation for the base using its nonzero
linear base coefficient, split off its nondegenerate $y$
direction, and take the cube root of the remaining holomorphic
unit. The extra scalar coordinate is a local product coordinate.
The critical ideal $(x^2,y)$ has transverse length two.
Its literal critical cycle is therefore $2E_{II}$.
The normal comparison is the normalized unitary matrix
\[
 \frac1{\sqrt{4|y|^2+9|x|^4}}
 \begin{pmatrix}2y&-3\overline{x}^{2}\\3x^2&2\overline y\end{pmatrix}.
                                                               \tag{D45.5}\label{maps:eq:D45.5}
\]
Its degree is two: the holomorphic map $(x,y)\mapsto(2y,3x^2)$
has two positive local preimages of a small regular value.
Thus the relative obstruction, as well as the scheme cycle,
has coefficient two. The scalar $\delta$ is nonzero in the
second abstract kernel summand, including on $E_{II}$.

\item At a good logarithmic cap the total cyclic quotient is free,
so its good cover is etale over the smooth quotient total
space. On that cover $q=s^m$. If $D$ is the reduced multiple
fibre, remove the common divisorial factor of the differential
and use the target line
\[
 \mathcal L=f^*T\mathbb P^1\otimes
                   \mathcal O(-(m-1)D).                     \tag{D45.6}\label{maps:eq:D45.6}
\]
The corrected differential pulls back to $m\,ds$ and is
surjective. The division is deck-equivariant, and a unit
change $q=u s^m$ gives $mu\,ds+s\,du$, still surjective.
Consequently a good logarithmic cap has no corrected critical
cycle. Counting the uncorrected divisorial zero would give
the wrong relative bundle. Its vertical kernel is a subbundle
and has the descended nowhere-zero scalar section.

\item Write $C$ for the complete corrected critical cycle: coefficient
one on each nodal elliptic curve and each triangular double
curve, coefficient two on an original $II$ elliptic curve,
and no contribution at a good logarithm. The locally free
kernel $\mathcal V$ has the global nonzero abstract section
$\delta$. Both $\mathcal V/\mathcal O\delta$ and $\mathcal L$
are topologically trivial because $H^2(X;\mathbb Z)=0$.
The framed bundle $\mathcal V\oplus\mathcal L$ therefore
provides a specified full reference frame $\rho$ of $TX$
away from $C$. The local matrices in \ref{maps:src:D.2.44} and \eqref{maps:eq:D45.5} give
the exact relative class
\[
                  \lambda(TX,\rho)=-\operatorname{PD}(C).
                                                               \tag{D45.7}\label{maps:eq:D45.7}
\]
This specifies local coefficients, not merely the absolute
identity $\lambda=-c_2$. In particular we do not assume
that $C$ is nullhomologous; it has order five when $b=20$.

\item \emph{A relative-characteristic formula retaining the rational link.}
We prove the needed general fact. Let $N\subset X\setminus C$
be an oriented normally framed three-manifold generating
$H_3(X)=\mathbb Z/b$, and suppose that the stable tangent
frame induced by $\rho$ on $N$ is an even multiple of
$\nu\in\Omega^{\rm fr}_3=\mathbb Z/24$.
Choose a four-chain $B$ with $\partial B=bN$, transverse to
the curve strata and avoiding their finitely many singular
points. Set
\[
                     \mathcal L(N,C)=\frac{B\cdot C}{b}
                             \in\mathbb Q.                  \tag{D45.8}\label{maps:eq:D45.8}
\]
This is the actual rational number for the marked cycles,
not a residue modulo integers. It is independent of $B$,
since $[C]$ is torsion. Then
\[
       \operatorname{lk}(N,\operatorname{PD}P)
                          =-\frac12\mathcal L(N,C)
                             \quad\text{in }\mathbb Q/\mathbb Z.
                                                               \tag{D45.9}\label{maps:eq:D45.9}
\]

\item Here is a proof that includes the missing parity. Remove a
point away from all chains. Write $b=2^a q$ with $q$ odd.
Join $q$ parallel copies of $N$ by normally framed tubes.
Kaplan's prescribed-spin construction, followed by deletion
of its zero-handle and reversal, gives a cobordism to a
sphere using only two-handles
\cite[Definition 1.7, Remark 1.8 and Theorem 3.1]{mapsKaplan79}.
As in \ref{maps:src:D.2.44}, its attaching disks and normal four-frames can
be realized in the punctured simply connected six-manifold:
general position gives disjoint disks and the spin condition
is exactly their $\pi_1SO(4)$ extension obstruction.
Thus the outgoing sphere is actually embedded and represents
$q[N]$, the two-primary generator.
Now choose Zhubr's admissible three-frame $v$ for this
particular embedded sphere, as permitted by his Theorem 1
\cite{ffZhubr13}, and put $\xi=v^\perp$. Its spin framing
obstruction is $o_4(\xi)=P$, and
$\lambda(\xi)=2o_4(\xi)$ with his normalization.
Choose a spin frame $\tau$ of $\xi|_N$.
It extends over the tube and two-handle trace: the relative
obstruction groups are $\pi_0\mathrm{Spin}(3)$ and
$\pi_1\mathrm{Spin}(3)$, both zero. Hence the framings
induced by $v+\tau$ on the incoming and outgoing manifolds
are stably bordant. On the outgoing sphere, admissibility
means even framed parity, as normalized in \ref{maps:src:D.2.44} and
\cite[Section 5]{ffZhubr13}. Since $q$ is odd,
$[N]_{v+\tau}$ is even as well.

\item The complement of the real codimension-four curve strata
is simply connected, by general position applied to disks.
Both $\rho|_N$ and $v+\tau$ therefore induce the same spin
structure. Their difference lifts to $\mathrm{Spin}(6)$.
Let $a$ be its normalized degree-three integral class and
$k=\langle a,[N]\rangle$. The stable $J$ image of the
$\pi_3\mathrm{Spin}(6)$ generator is $\nu$, so the two
even framed parities imply $k$ is even. The integral
transgression of $\lambda$ gives, with consistent
transition-map direction, the identity of relative classes
\[
 \lambda(TX,\rho|_N)=2o_4(\xi,\tau)+\delta a
                         \quad\text{in }H^4(X\setminus\{p\},N;\mathbb Z).
                                                               \tag{D45.10}\label{maps:eq:D45.10}
\]
If $Q=\langle o_4(\xi,\tau),B\rangle$, evaluation gives
the identity of integers
\[
                       -B\cdot C=2Q+bk.                    \tag{D45.11}\label{maps:eq:D45.11}
\]
By the integral definition of the torsion pairing,
$Q/b\bmod\mathbb Z=\operatorname{lk}(N,\operatorname{PD}P)$.
Dividing \eqref{maps:eq:D45.11} and using the evenness of $k$ proves
\eqref{maps:eq:D45.9}. This is not division of a class in an even cyclic
group. Its integer correction was fixed by actual spin
frames before division. The proof uses Zhubr's proved
Theorem 1, and does not use his conjectural general
change-of-three-frame formula. It also permits a nonzero
odd-primary $\lambda$.

\item \emph{The finite-core generator and its reference parity.}
For a single good $II$ logarithm the regular three-torus
can be an even class in $H_3(X)$. We use the following
actual core three-manifold instead. The central four-manifold
is an elliptic-torus bundle over the rank-two lattice
\[
 \langle\delta,\ell,g\mid mg=r\delta+n\ell\rangle.
                                                               \tag{D45.12}\label{maps:eq:D45.12}
\]
Its elliptic monodromy is $B$ on $g$, trivial on the other
two generators, where $B$ has order six and
$\det(B-I)=1$. With $\sigma=n-me$, define the primitive
integer character
\[
 \chi(\delta)=1,\qquad \chi(\ell)=-\sigma r,
                       \chi(g)=-\sigma re.                 \tag{D45.13}\label{maps:eq:D45.13}
\]
The relation in \eqref{maps:eq:D45.12} is respected. Let $N$ be the
connected level set of the resulting map of the core to
$S^1$, cooriented by $\chi$. Put
$h=g+\sigma re\delta$ and $k_0=\ell+\sigma r\delta$.
These lie in $\ker\chi$ and satisfy $mh=nk_0$.
Since $\gcd(m,n)=1$, their primitive generator has elliptic
monodromy $B^\sigma$: an integer inverse of $n$ modulo
$m$ is congruent to $\sigma$ modulo six. Thus $N$ is an
actual affine elliptic mapping torus, with
\[
                        H_1(N;\mathbb Z)=\mathbb Z.        \tag{D45.14}\label{maps:eq:D45.14}
\]
Translations do not change this Wang sequence.

\item For completeness, any actual tangent parallelization of
a closed oriented three-manifold $M$ has framed parity
\[
        [M,\tau]_{\rm fr}\bmod2
                    =1+\dim H_1(M;\mathbb F_2).             \tag{D45.15}\label{maps:eq:D45.15}
\]
Use Kaplan's spin-compatible filling $W$ with one zero-handle
and $j$ two-handles. Its tangent bundle is trivial, since
its spine is two-dimensional and $w_2=0$. On the boundary
use the outward normal followed by $\tau$.
The relative stable framing obstruction is $\lambda$;
its $J$ image is $[M,\tau]$. Modulo two it is the relative
$w_4$, hence the relative Euler class, whose evaluation is
$\chi(W)=1+j$. The mod-two intersection matrix of the
spin filling is alternating, so has even rank. Its
cokernel is $H_1(M;\mathbb F_2)$; hence
$j\equiv\dim H_1(M;\mathbb F_2)$. This proves \eqref{maps:eq:D45.15}
while retaining the prescribed spin structure.

\item The reference frame on the finite-core $N$ is the
stabilization of an actual tangent frame. Indeed $\delta$
is transverse to $N$ by $\chi(\delta)=1$. For the global
complex vertical frame $(\delta,v_2)$, project
$J\delta,\operatorname{Re}v_2,\operatorname{Im}v_2$ to $TN$
along $\delta$. They form a tangent frame. The normal
frame is $\delta$ together with the two real directions
of the corrected target line $\mathcal L$. Projection
and horizontal splitting are upper triangular homotopies.
Equation \eqref{maps:eq:D45.14}, then \eqref{maps:eq:D45.15}, therefore gives even
reference parity. This conclusion is not inferred from
the even-degree torus cover of $N$.

\item \emph{The actual marked cycles and the unbounded chain identity.}
Use the same integral regular-fibre coordinate as D.2.41:
\[
                \Omega=a e_1^*e_2^*-b\delta^*\gamma^*.
                                                               \tag{D45.16}\label{maps:eq:D45.16}
\]
The full degree-two monodromy cokernel is $\mathbb Z$,
primitively measured by $\Omega$. Choose an integral
fibre two-cycle $s$ with $\Omega(s)=1$; its image generates
$H_2(X)=\mathbb Z/b$. The nodal and triangular preimages
are those proved in \ref{maps:src:D.2.44}. At an original $II$ site put
\[
 z=(B-I)^{-1}q_0\in\mathbb Z^2,\qquad
 \ell=(-z,0,1),\qquad [E_{II}]=\delta\wedge\ell,             \tag{D45.17}\label{maps:eq:D45.17}
\]
where $q_0$ is the full elliptic long-column vector.
The full matrix fixes $\ell$, not just its quotient.
The principal $\delta,\ell$ torus through the elliptic
cusp is the reduced critical elliptic curve; its regular
torus specializes with degree one. The comparison takes
place within that cap. Thus $2\delta\wedge\ell$ is the
literal critical preimage, with the complex torus orientation.
For OS45/II, $\ell=(-8,0,0,1)$; for the OS55/II rows,
$\ell=(0,0,0,1)$.

\item Transport all these preimages along the fixed base paths
and let $K=\Omega(C)$. The exact integers are
\[
\begin{array}{c|c|c|r|c}
 \text{profile}&(a,b)&\text{zero and }II\text{ choice}&K&
                                  \operatorname{ord}[C]\\\hline
 45/II&(1,8)&2,3;\ \mathrm{original}&-32&1\\
 45/II&(1,8)&2,3;\ \mathrm{log}&-16&1\\
 55/II&(5,4)&3;\ \mathrm{original}&-108&1\\
 55/II&(5,4)&3;\ \mathrm{log}&-100&1\\
 55/II&(1,20)&3;\ \mathrm{original}&-68&5\\
 55/II&(1,20)&3;\ \mathrm{log}&-28&5\\
 55/\mathrm{default}&(1,20)&2&-108&5\\
 55/\mathrm{default}&(1,20)&3,4&-68&5.
\end{array}                                                    \tag{D45.18}\label{maps:eq:D45.18}
\]
The difference $C-Ks$ bounds a chain in the regular
complement, after the specified local specialization
chains are included. This uses the full integral
monodromy image, whose five nonzero Smith factors are
units; it does not replace an image by its saturation.
For a logarithmic pattern all these chains avoid its
finite core, since that cap has no critical curve.

\item Here is the exact all-parameter relation that fixes
both the generator and the rational linking.
Every monodromy fixes $\delta$. Its quadratic bar
correction on that single loop is zero, including
the transported last peripheral word, and the
peripheral graph maps satisfy $\sum_iJ_i\delta=0$.
Take the finite normal Thom column for $\delta$,
add $m$ times the ordinary-clutch $\delta$ column,
and add $m$ times the other filled-cap $\delta$ sweeps.
The graph terms cancel. Its actual fibre part is
\[
 -\delta\wedge(r\delta+n\ell)+me\delta\wedge\gamma,
 \qquad \Omega=b(n-me)=b\sigma.                            \tag{D45.19}\label{maps:eq:D45.19}
\]
The remaining difference from $b\sigma s$ is an
integral monodromy relation, supported away from $N$.
All other cap sweeps also bound away from $N$.
The finite normal Thom chain is the disc bundle
over its actual central loop. Its intersection
with the level set $N$ is its character value,
which here is $\chi(\delta)=1$.
Thus \eqref{maps:eq:D45.19} gives the actual pairings
\[
 \operatorname{lk}(N,s)=\frac{\sigma}{b}\pmod{\mathbb Z},
 \qquad \mathcal L(N,C)=\frac{\sigma K}{b}.                 \tag{D45.20}\label{maps:eq:D45.20}
\]
The first is a unit pairing, proving that $N$ is a
full $H_3$ generator, including when $b=20$.
The fifth affine-meridian column, if used in another
Smith bounding chain, pairs by its actual value
$\chi(g)=-\sigma re$, not zero. This retained term
makes the full-chain comparison agree with \eqref{maps:eq:D45.20}.
There is no additional factor of $m$: the covering
order is already present in the boundary norm column.

\item For an original-only pattern use the ordinary
$T^3=\langle e_1,e_2,\gamma\rangle$ and $e=\pm1$.
The same peripheral relation, now using its ordinary
normal Thom chain, has fibre value $-be$ and
intersection one with $T^3$. Therefore
\[
 \operatorname{lk}(T^3,s)=-\frac e b\pmod{\mathbb Z},
 \qquad \mathcal L(T^3,C)=-\frac{eK}{b}.                    \tag{D45.21}\label{maps:eq:D45.21}
\]
It is again a full generator with even reference parity,
as proved in \ref{maps:src:D.2.44}. If $b=20$, these equations can
equivalently be computed by bounding $5C$ and dividing
its literal intersection by five. No nullhomology
of $C$ is asserted.

\item Substitute \eqref{maps:eq:D45.20} or \eqref{maps:eq:D45.21} into the already
proved relative formula \eqref{maps:eq:D45.9}. Since $K$ is an
even integer in every row of \eqref{maps:eq:D45.18}, cancellation
of the displayed unit pairing gives
\[
                       P=-K/2\quad\text{in }\mathbb Z/b.
                                                               \tag{D45.22}\label{maps:eq:D45.22}
\]
Here the division was justified by \eqref{maps:eq:D45.10}--\eqref{maps:eq:D45.11}
using the actual cycles and framings. It is not
division of the congruence $2P=\lambda$.
The table gives $P=2,0,14$ as in \eqref{maps:eq:D45.3}, and the
absolute identity $\lambda=-[C]$ gives $0,0,8$.

\item \emph{The complete fundamental-class invariant and recognition.}
The global auxiliary-line comparison and localization
of \ref{maps:src:D.2.44} apply over $\mathbb Z/4$ and $\mathbb Z/8$,
respectively. All actual original-$II$ coefficients
are zero and all finite-$II$ normal coefficients
are $bn$ and $-br$; hence the same mod-$b$ class
extends for every parameter in \eqref{maps:eq:D45.2}.
For $b=20$ apply that calculation separately over
$\mathbb Z/4$ and $\mathbb Z/5$. The $C_4$ action
has only the rational fixed curve and the $I_4$
elliptic fixed curve, with characters $2,0$.
The $C_5$ action similarly has the rational and
$I_5$ curves, with characters $0,0$; all remaining
stabilizers are prime to the relevant coefficient
order. The actual Tate normal weights have square
one modulo four, and weight two modulo five.
The cube formula is zero in each coefficient ring:
on the rational curve it is $-3$ times the square
of its point character, and on the elliptic curve
the possible mixed degree-one term has square zero.
Thus the full fundamental-class invariant is zero,
since $H_6(K(\mathbb Z/b,2);\mathbb Z)=\mathbb Z/b$
is detected by the cube for $b=4,8,20$
\cite[Sections 3.11--3.14]{manZhubr75}.

\item The finite data therefore agree with the first two
surgery models of \ref{maps:src:D.2.44}. For $b=20$, the two
primary parts of $P=14$ are $2\in\mathbb Z/4$
and a nonzero element of $\mathbb Z/5$.
The inclusion
$K(\mathbb Z/4,2)\vee K(\mathbb Z/5,2)
\to K(\mathbb Z/4,2)\times K(\mathbb Z/5,2)$
is a homology equivalence: all cross tensor and
Tor terms vanish because their positive-degree
homology has coprime primary orders. Both spaces
are simply connected, so it is a homotopy equivalence.
The reduced spin-bordism invariant consequently
splits into these two summands, as in \ref{maps:src:D.2.38}.
The complete invariants agree with
$Z_{4,1}\#Z_{5,1}$, and the free third Betti numbers
also agree. Zhubr's Theorems 3 and 5
\cite{manZhubr75} give all the oriented
diffeomorphisms in \eqref{maps:eq:D45.3}.

\item \emph{Exact computations and their bounds.}
The local differential identities are checked in
\path{code/III/maps_D2_II_frame_checks.py}.
The line/support, full angular-image, actual
adapted-matrix and arbitrary finite-normal
coefficients for all eighteen patterns are checked in
\path{code/III/fourfolds_D2_remaining_AH.py}.
This comprises fourteen initial rows, $1056$ exact
semicharacter checks, $62$ full adapted matrices,
four symbolic arbitrary-twist finite normal rows
and four original-$II$ choices.
The actual curve vectors, complete Smith solutions,
finite root-loop characters and literal intersections
are retained by
\path{code/III/manifolds_extended_critical_cycles.py}.
Its $158$ full Smith replays comprise the two original
clutch signs and, for each of the four logarithmic
patterns, $k=1,2,3$, $r=0,1$, $e=-1,0,1$ and
$\sigma=\pm1$. This explicitly bounded replay is
supplemented by the eleven constant certificates in
\path{records/manifolds_extended_critical_symbolic.json}:
integral preimages of $-\delta\ell-bs$,
$\delta\gamma+bs$, and $C-Ks$, together with the
zero bar corrections and exact peripheral graph
identity. Their linear combinations prove \eqref{maps:eq:D45.19}
for every integer parameter in \eqref{maps:eq:D45.2}, with no
unproved finite search bound.\hfill$\square$
\end{enumerate}
\end{proof}


\begin{theorem}[arbitrary whole-fibre initial pairs and weights]\label{maps:src:D.2.46}
\StatementRecord{maps:src:D.2.46}{proved}{Historical events: Claude Fable 5.2 (Anthropic), effort high, 16 November 2026: carried as background without fault found. GPT-6 Astra (OpenAI), ChatGPT Pro mode, effort Pro, S6\_09\_PAPERS, 17 November 2026: targeted proof check. Also checked in this writing. Pre-continuation writing events: GPT-6 Astra (OpenAI), ChatGPT Pro mode, effort Pro, S6\_10\_PAPERS, 18 November 2026: targeted proof check. Continuation metadata events: GPT-6 Astra (OpenAI), ChatGPT Pro mode, effort Pro, S6\_10\_PAPERS continuation, 19 November 2026: used and rechecked in this run; status-vocabulary reconciliation; metadata reconciliation only; no new mathematical check. \emph{Pre-continuation recorded proof scope.} GPT-6 Astra (OpenAI, ChatGPT Pro mode), effort Pro, 2026-11-17 UTC; Claude Fable 5.2 (Anthropic), effort high, 16 November 2026; also checked in this writing by GPT-6 Astra (OpenAI), ChatGPT Pro mode, effort Pro, 18 November 2026}{the hypotheses, notation and definitions stated here; \cref{maps:src:D.2.19}; \cref{maps:eq:D.BC1}; \cref{maps:eq:D.BC14}; \cref{maps:eq:D.BC4}; \cref{maps:src:D.2.29}; \cref{maps:eq:D.BC5}; \cref{maps:src:D.2.22}; \cref{maps:src:D.2.23}; \cite[Section 11.8]{manSS10}; \cref{maps:src:D.2.37}; \cref{maps:src:D.2.43}; \cref{maps:src:D.2.31}; \cref{maps:src:D.2.32}; \cref{maps:src:D.2.40}; \cref{maps:eq:D.GI9}; \cref{maps:src:D.2.41}}


Fix any of the $21$ marked rational elliptic profiles with
$\operatorname{MW}=\mathbb ZG$ in \ref{maps:src:D.2.19}, and write
$\langle G,G\rangle=1/d_0$. Let
\[
 P=aG,\quad Q=bG,\quad a,b>0,
 \qquad \overline P_i=0\ \hbox{or}\ \overline Q_i=0
                         \quad\hbox{at every fibre}.       \tag{D.GI1}\label{maps:eq:D.GI1}
\]
Choose any nonnegative integral whole-fibre linearization
weights $w_i$, of total $ab/d_0$, supported on multiplicative
or regular fibres. Distinct regular support points are allowed.
At a positive weight use the marked periodic Mumford filling.
At a zero-free multiplicative fibre choose either its original
cap or its marked Mumford cap. At a finite good-reduction
additive fibre choose either its original cap or any free affine
logarithmic cap. Retain the original cap at $I_n^*$.
Arbitrary regular logarithms and an arbitrary ordinary clutching
integer are allowed.

There is a finite exact integral algorithm for the fundamental
group and every integral homology group of every such threefold.
It includes pairs not obtained by iterating any of the $172$
initial rows, several zero fibres, noncyclic toric period
lattices, and all their boundary torsion. Its implementation is
\path{code/III/manifolds_general_initial.py}; no bound on
$a,b,w_i$, logarithmic orders or affine numerators occurs in the
algorithm or proof.

Here are the complete data of the algorithm. Use the lattice
$\Lambda=\mathbb Z\langle e_1,e_2,\delta,\gamma\rangle$.
The recorded elliptic monodromies are $B_i$, and the section
$G$ has integral cocycle $g_i$. Put $p_i=ag_i$, $q_i=bg_i$,
choose any rational $z_i$ with $(B_i-I)z_i=q_i$, and set
\[
 r_i=((p_i)_2,-(p_i)_1)B_i,\qquad
 A_i=\begin{pmatrix}
 B_i&0&q_i\\ r_i&1&r_i z_i+w_i\\0&0&1
 \end{pmatrix}.                                           \tag{D.GI2}\label{maps:eq:D.GI2}
\]
The block sizes are $2,1,1$. At a regular zero use
$A_i=I+w_i\delta\gamma^*$. At an ordinary or regular
logarithm the boundary matrix is $I$. These are integral
matrices, their ordered product is $I$, and they are the
actual marked monodromies.

Apply the integral bar comparison \eqref{maps:eq:D.BC1}--\eqref{maps:eq:D.BC14} to these
matrices. Thus the regular-piece degree-two group is
$C\oplus K$, with its full unsaturated fibre presentation
and its transported last-boundary marking. At each finite or
regular logarithm use all five columns \eqref{maps:eq:D.BC4}, with its
actual order $m$ and fixed power $v$, and the exact freeness
test of \ref{maps:src:D.2.29}--\ref{maps:src:D.2.30}. At every other cap use the full fixed-normal
columns \eqref{maps:eq:D.BC5} and the following fibre kernel $F_i$ in
$\Lambda^2\Lambda$. Here $R_i=\Lambda^2A_i-I$ and
$\operatorname{sat}$ denotes saturation only where it is
explicitly written.
\[
\begin{array}{c|c|c}
\hbox{cap}&F_i& h_i=\operatorname{rank}H^4(U_i)\\ \hline
w_i>0\hbox{ at }I_n, n>0&
 \operatorname{im}R_i+\Lambda^2\operatorname{Fix}A_i&nw_i\\
\hbox{rank-one Mumford}&\operatorname{sat}(\operatorname{im}R_i)&\rho_i\\
\hbox{original, }\overline P_i=0&
 \operatorname{sat}(\operatorname{im}R_i)&\rho_i\\
\hbox{original additive, }\overline P_i\ne0&\ker\omega_i&s_i\\
\hbox{original }I_n,\ \overline P_i\ne0&
 \langle u_1\delta,u_1\ell_i,\delta\ell_i\rangle&n .
\end{array}                                                \tag{D.GI3}\label{maps:eq:D.GI3}
\]
Products in the last row are exterior products.
For a rank-one Mumford cap, $\rho_i$ is the unique nonzero
Smith factor of $A_i-I$. For a zero-Chern original cap it is
the number of orbits of translation by $Q$ on the Kodaira
components. The actual marked permutations were given in
\ref{maps:src:D.2.22}. In the last two rows, $Q$ is narrow, so $z_i$ may
and must be integral, and
\[
 \ell_i=(-z_i,0,1),\qquad
 \omega_i=(1,0,(z_i)_2,0,-(z_i)_1,0).
                                                               \tag{D.GI4}\label{maps:eq:D.GI4}
\]
Here $u_1$ is the primitive elliptic vanishing circle and
$s_i$ is the number of Kodaira components. At a positive
regular zero the second row applies, with $\rho_i=w_i$.

Let $\mathcal B$ be the resulting full integral degree-two
matrix, in $\Lambda^2\Lambda\oplus K$, and write
$\operatorname{coker}\mathcal B=\mathbb Z^g\oplus T$.
Let $E$ be the regular piece with all finite and regular
logarithmic caps already filled. Compute $H_1(E)$ from its
literal graph-bundle presentation and the relations
$\tau_i^{m_i}=v_i$. Compute $H_1(X)=\mathbb Z^r\oplus A_1$
by additionally imposing, at each other cap, $\tau_i=1$
and the following fibre relations:
\[
\begin{array}{c|c}
\hbox{cap}&\hbox{killed subgroup of }\Lambda\\\hline
\hbox{positive rank-two Mumford}&\operatorname{Fix}A_i\\
\hbox{rank-one Mumford or zero-Chern original}&
                  \operatorname{sat}(\operatorname{im}(A_i-I))\\
\hbox{nonzero-Chern additive original}&
                  \mathbb Z\langle e_1,e_2,\delta\rangle\\
\hbox{nonzero-Chern multiplicative original}&
                  \mathbb Z\langle u_1,\delta\rangle .
\end{array}                                                \tag{D.GI5}\label{maps:eq:D.GI5}
\]
Set
\[
 q=\sum_{\text{nonlog caps}}h_i-b_1(E)+r,\qquad
 b_2=g+q,\qquad
 \chi=2\sum_{I_n,\ n>0}nw_i,\qquad
 b_3=2+2b_2-2r-\chi .                                    \tag{D.GI6}\label{maps:eq:D.GI6}
\]
Then $q\ge0$ and the complete answer is
\[
\begin{array}{c|c}
j&H_j(X;\mathbb Z)\\\hline
0,6&\mathbb Z\\
1&\mathbb Z^r\oplus A_1\\
2&\mathbb Z^{b_2}\oplus T\\
3&\mathbb Z^{b_3}\oplus T\\
4&\mathbb Z^{b_2}\oplus A_1\\
5&\mathbb Z^r .
\end{array}                                                \tag{D.GI7}\label{maps:eq:D.GI7}
\]
For the nonabelian fundamental group, retain four commuting
fibre generators, all peripheral generators $\tau_i$, the
relations $\tau_i x\tau_i^{-1}=A_i x$ and
$\prod_i\tau_i=1$, all finite-power relations, and
\eqref{maps:eq:D.GI5}. This is a complete finite presentation, not just
its abelianization.

Except in the case described next, the entire vertical
subgroup $\langle e_1,e_2,\delta\rangle$ is killed and this
presentation reduces to the Seifert group of \ref{maps:src:D.2.23}, with
the actual gamma numerators. The exceptional possibility is
\[
\begin{gathered}
 \mathrm{OS49/II}=(II,IV^*,I_2),\quad a\text{ odd},\quad
 b\text{ even},\quad\text{both additive caps logarithmic},\\
 w_{I_2}=0,\quad I_2\text{ filled by Mumford},\quad
 \text{all positive weights at regular fibres}.
\end{gathered}                                             \tag{D.GI8}\label{maps:eq:D.GI8}
\]
The compatibility condition \eqref{maps:eq:D.GI1} remains required; it
can impose an additional factor $3$ on $a$ or $b$.
Before the logarithmic power relations the remaining
vertical group in \eqref{maps:eq:D.GI8} is $(\mathbb Z/2)^2$. The two
additive monodromies act by inverse order-three
automorphisms. With no additional regular logarithms,
the group is the split semidirect product of this Klein
four group by the two-fibre Seifert group. With arbitrary
regular logarithms use the full presentation above.
Every group in \eqref{maps:eq:D.GI8} surjects onto $\mathbb Z/3$, so none
is simply connected.

\end{theorem}


\begin{proof}
\begin{enumerate}[label=\arabic*.,leftmargin=*]
\item \emph{The initial data are actual section-line data.}
Let $M=\mathcal O(P-O)$ and
$J=M\otimes(T_Q^*M)^{-1}$. On a smooth elliptic fibre,
translation acts trivially on $\operatorname{Pic}^0$, so
$J$ is a vertical line bundle. At a singular fibre its
component multidegree is the difference of that of $M$
and its $Q$-translate. It is zero if $P$ is narrow, and
also if $Q$ is narrow. Thus \eqref{maps:eq:D.GI1} makes every local
multidegree zero. The intersection matrix of each
Kodaira fibre has kernel generated by its full
multiplicity vector; hence $J$ is pulled back from the
base. The height calculation in \ref{maps:src:D.2.19}, equivalently
the height pairing with all local mixed corrections
zero \cite[Section 11.8]{manSS10}, gives
\[
                         J\simeq f^*\mathcal O(ab/d_0).
\]
This proves integrality of $ab/d_0$ in the compatible
sector. A base section with precisely the chosen
nonnegative divisor gives the required linearization.
This does not introduce a component-supported vertical
correction to $M$.

\item For completeness, the monodromy formula is also an
integral geometric formula. The rigidified Poincare
bundle has elliptic section cocycles $p_i,q_i$; its
commutator pairing gives the row
$r_i=((p_i)_2,-(p_i)_1)B_i$. A local logarithm of $Q$
gives the central entry $r_i z_i$, and a zero of the
linearization adds $w_i\delta\gamma^*$. This is the
same marked local calculation used in \ref{maps:src:D.2.19} and
\ref{maps:src:D.2.37}, now before restricting either section coefficient.
Changing $z_i$ by a fixed vector changes no entry:
$r_i$ annihilates $\ker(B_i-I)$.

\item If $Q$ is narrow, $z_i$ is integral. If instead $P$
is narrow, write $p_i=(B_i-I)y_i$ with $y_i$ integral.
For $J_2=\left(\begin{smallmatrix}0&1\\-1&0\end{smallmatrix}\right)$,
symplectic invariance gives
\[
 r_i z_i=-y_i^t(B_i-I)^tJ_2B_i z_i
                         =y_i^tJ_2q_i\in\mathbb Z.
\]
Thus \eqref{maps:eq:D.GI2} is integral in both cases. At zero weights
its ordered product is the central matrix with entry
$-ab/d_0$; the local Poincare commutators are bilinear
in $a,b$. This identity is checked already for the
single generator pair on each of the $21$ recorded
elliptic tuples and follows for all coefficients by
bilinearity. Adding the weights changes that entry by
$\sum w_i=ab/d_0$, proving the product relation.
The implementation reconstructs \eqref{maps:eq:D.GI2} directly from
the saved cocycles and verifies integrality and product,
rather than treating a numerical monodromy output as
a topology theorem.

\item \emph{All local kernels, with their integral markings.}
At a positive $I_n$ weight, use a primitive elliptic
vanishing/long basis, followed by $\delta,\gamma$.
The nilpotent part has image in the primitive short
plane $\langle u_1,\delta\rangle$, squares to zero,
and has determinant $nw_i$ on long-to-short lattices.
Its Smith factors $s\mid t$ satisfy $st=nw_i$.
The actual linearization changes only the corresponding
integral long period; the same triangular periodic fan
is preserved by this period sublattice. The unit factors
are removed by the root-and-unit comparison of \ref{maps:src:D.2.22}.
Consequently this is exactly the cap of \ref{maps:src:D.2.43}, with
\[
 H_2(U_i)=\mathbb Z^{nw_i+3}\oplus\mathbb Z/s,
 \qquad H_3(U_i)=\mathbb Z^2,\qquad
 H_4(U_i)=\mathbb Z^{nw_i}.
\]
Its full fibre kernel is the first row of \eqref{maps:eq:D.GI3},
\emph{without saturation}. The two primitive short
normal tori bound their actual orbit discs. In
particular its angular-area torsion relation, and
the other boundary torsion class which maps onto
$\operatorname{Tor}H_2(U_i)$, remain distinguished.
The component meridians are the literal rays $(v,1)$
of \ref{maps:src:D.2.43}; thus no change from the original period
marking to a saturated lattice has occurred.

\item At rank one write $A-I=\rho\alpha\beta^*$ with
$\alpha$ and $\beta^*$ primitive. The local toric
model is the $\rho$-component nodal-cycle model
times the transverse torus, with the same unit
comparison and primitive vanishing disc as in
\ref{maps:src:D.2.37}. Its homology is free, its fourth rank is
$\rho$, and its full fibre kernel is precisely
the saturation of $\operatorname{im}R_i$.
The full fixed lattice supplies the normal tori.
This argument includes a regular zero, for which
$\alpha=\delta$, $\beta^*=\gamma^*$ and $\rho=w_i$.

\item For an original zero-Chern cap, the line bundle is
trivial on the Kodaira normalization and the prescribed
unit lift has no residual meridian winding. Its
retraction is $S^1_\delta$ times the mapping torus
of the actual component permutation by $Q$.
For an additive fibre the curve configuration is
a tree; for $I_n$ its dual graph is a circle.
The latter has no extra line-bundle winding: the
component degrees are all zero, the translation
lift on the compact connected cycle differs from
the rigidified lift by a constant unit, and the
nonvanishing base linearization has a logarithm.
The component permutation therefore contributes
one free fourth class per orbit. The explicit
normalization or tree-chain complex has free
homology and gives exactly the third row of
\eqref{maps:eq:D.GI3}, with all fixed normal tori. This is the
permutation calculation of \ref{maps:src:D.2.31}; it uses the
actual residues of $bG$, not primitivity of $b$.

\item Suppose next that the original cap has
$\overline P_i\ne0$. Then $Q$ is narrow, and the
Chern vector of $M$ on the components is
\[
                         c=e_{\overline P_i}-e_0.
\]
It is primitive even if the component of $P$
is not a generator of the component group.
Its pairing with the Kodaira multiplicity vector
is zero. The integral logarithm of $Q$ and the
rigidified unit comparison contract the translation
while retaining this Chern class. On a tree the
unit comparison has a logarithm because $H^1=0$.
On a nodal cycle its final clutch exponent is
$\sum c_i=0$, so the same comparison has no
dual-circle winding. Thus the actual retract is
the principal two-torus bundle with Chern pair
$(c,0)$, or $Y_c\times S^1_\gamma$ in the cycle
case, precisely as in \ref{maps:src:D.2.32}--\ref{maps:src:D.2.33}.

\item The primitivity of $c$, not the stronger and
unneeded assumption that $\overline P_i$ generates
the component group, makes all these local homology
groups free. In the tree case the regular elliptic
fibre maps to the primitive multiplicity vector in
$\ker c$. Its exact area functional is $\omega_i$,
giving the fourth row of \eqref{maps:eq:D.GI3}. In the cycle case,
the three regular fibre images
\[
                       u_1u_2,\quad u_2\delta,\quad u_2\ell_i
\]
remain primitive: respectively they project to the
primitive all-one component vector, the fibre-circle
times the dual-circle generator, and the dual circle
times the second principal circle. This proves the
last row of \eqref{maps:eq:D.GI3} even when $A_i-I$ has nonunit
Smith factors. Its fixed lattice is exactly
$\langle u_1,\delta,\ell_i\rangle$.

\item The $\delta$ and $\ell_i$ normal tori bound with
zero fibre correction. The $u_1$ normal torus in
the cycle case can differ from a bounding normal
torus by an integral multiple of $u_2\delta$.
As in \ref{maps:src:D.2.40}, this correction is zero in the
\emph{full} global fibre coinvariant, including its
torsion. Indeed the relations on $\delta\wedge\Lambda$
are the actual elliptic columns $(B_i-I)\mathbb Z^2$,
and their common Smith factors are $(1,1)$ for
each of the $21$ profiles. This assertion is
independent of $a,b,w_i$ and is retained as an
explicit constant certificate. Hence all the normal
columns in the theorem have their actual global
marking; no rational replacement is involved.

\item The same local retractions give the fundamental
group kernels \eqref{maps:eq:D.GI5}. All meridians of these
nonlogarithmic caps bound their normal discs.
At a finite logarithm the principal meridian is
$\tau_i^m v^{-1}$, with its full integral normal
line torsion and the five Thom columns already
proved in \ref{maps:src:D.2.29}--\ref{maps:src:D.2.30}. These local statements prove
the claimed van Kampen presentation.

\item \emph{The normal closure and the exceptional profile.}
There is at least one positive weight. Its primitive
vanishing lattice contains $\delta$, so $\delta=1$
in the global fundamental group. Modulo $\delta$,
a zero-free rank-one $I_n$ Mumford cap kills
\[
                   \frac{n}{\gcd(n,a)}\,u_i,              \tag{D.GI9}\label{maps:eq:D.GI9}
\]
where $u_i$ is its primitive elliptic vanishing vector.
To see the coefficient, in its Tate basis the
vanishing vector before primitivization is
$nu_i+aH_i\delta$, with
$\gcd(H_i,n)=1$. The generator component of $G$
has that unit residue on every recorded $I_n$.
If $P$ is narrow the projection in \eqref{maps:eq:D.GI9} is
already primitive. If $P$ is non-narrow, the
compatibility condition makes $Q$ narrow; changing
its integral logarithm changes the long coordinate
but not this projected vanishing line.

\item The elliptic normal closure of these lines under
the finitely many $B_i$ depends only on
$a\bmod d_0$. The original star caps additionally
kill the full elliptic lattice after $\delta$ is
killed: a nonzero Chern vector kills it directly,
while a narrow $P$ supplies its integral graph
trivialization. Thus there are precisely $284$
constant residue cases to check. Their recorded
Hermite matrices are full $\mathbb Z^2$ in every
case except OS49/II with $a$ odd, where the matrix
is $2I_2$. The certificate also computes every
one of the $912$ actual local projected vanishing
lines and verifies \eqref{maps:eq:D.GI9} directly; this is an
exhaustive residue calculation, not a bounded
search over unbounded section pairs.

\item In that sole residual case, a positive $I_2$
weight adds its primitive short circle, and an
original $I_2$ cap adds the primitive elliptic
vanishing circle because its Chern vector is
nonzero. Either one and its $B_i$ translates
generate $\mathbb Z^2$. An original additive
cap does the same. Hence only \eqref{maps:eq:D.GI8} remains.
The two additive matrices modulo $2$ have order
three, are inverse, and $B_i-I$ is invertible
modulo $2$. Since $b$ is even, their gamma
columns vanish modulo $2$. Their fixed integral
power vectors consequently have even elliptic
coordinates. With no extra regular logarithms,
all power and base-product relations thus lie
in the Seifert factor, proving the asserted
split semidirect product. Additional regular
affine powers can have nonzero Klein-four
coordinates, which is why the full presentation
is retained in that case. In all cases set
the entire fibre group and additional regular
peripherals to $1$, and send the two additive
peripherals to opposite generators of $\mathbb Z/3$.
Their orders are multiples of $6$ and $3$.
This is a well-defined surjection and proves
the non-simply-connected assertion.

\item \emph{Integral relative gluing, including cap torsion.}
The bar comparison of \ref{maps:src:D.2.41} applies to any integral
tuple \eqref{maps:eq:D.GI2}; its proof made no scalar-iterate
assumption. It gives $H_2$ of the regular piece
with all fibre torsion and every fixed, norm and
last-peripheral correction. Filling the finite
caps quotients by precisely their five columns,
so it gives the stated $H_2(E)$ presentation.

\item Every remaining cap retracts to a complex of
dimension at most four and has free $H^4$ of
rank $h_i$. For the general cusp this follows
from \ref{maps:src:D.2.43}, which gives free $H_3$ as well as
free $H_4$; for the other caps it follows from
the explicit free groups above. Excision and
Poincare--Lefschetz duality give
\[
 H_2(X,E)=\bigoplus_iH^4(U_i)=\mathbb Z^{\sum h_i},
                    \qquad H_1(X,E)=0 .                  \tag{D.GI10}\label{maps:eq:D.GI10}
\]
The group $H_3(U_i,\partial U_i)=H^3(U_i)$ need
\emph{not} be free: it contains the torsion dual
to $H_2(U_i)$ at a general cusp. Its boundary
image is nevertheless exactly
$\ker[H_2(\partial U_i)\to H_2(U_i)]$.
\ref{maps:src:D.2.43} identifies that image with the two literal
normal sweeps and the nonsaturated first fibre
kernel in \eqref{maps:eq:D.GI3}. In particular the angular-area
$\mathbb Z/s$ is included, while the other
$\mathbb Z/s$ maps isomorphically to cap torsion.
This accounts for the torsion relative boundary
classes with their actual marking; freeness of
cap $H_2$ is not assumed.

\item It follows that the image of $H_3(X,E)$ in
$H_2(E)$ is exactly the remaining columns of
$\mathcal B$. The pair sequence therefore has
the exact form
\[
 0\longrightarrow\operatorname{coker}\mathcal B
 \longrightarrow H_2(X)
 \longrightarrow\ker[\mathbb Z^{\sum h_i}\to H_1(E)]
 \longrightarrow0.                                      \tag{D.GI11}\label{maps:eq:D.GI11}
\]
The final term is a free subgroup of a free
group; its rank is $q$ because
$H_1(E)\to H_1(X)$ is surjective. The sequence
splits as abelian groups, proving the degree-two
formula with precisely the torsion $T$.
All caps except positive rank-two cusps have
Euler number zero; a cusp of period index $nw_i$
has Euler number $2nw_i$. Hence \eqref{maps:eq:D.GI6} gives
the Euler number and the middle Betti number.
Integral duality and the universal coefficient
theorem give every remaining group in \eqref{maps:eq:D.GI7}.
This includes zero rational Euler number and
arbitrary first Betti number, without a Seifert
rank assumption in the exceptional sector.
\end{enumerate}
\end{proof}

The replay \path{code/III/manifolds_general_initial_audit.py}
stores all $284$ normal-closure cases, $912$ local
rank-one projections, $21$ generator height products
and elliptic unit lattices, $90$ complete overlaps
with \ref{maps:src:D.2.41}, and $148$ full Smith examples in the
larger parameter sector. For example, OS45/default
with $P=2G$, $Q=8G$, weight $2$ at $I_8$, all other
weights zero, all Mumford caps and ordinary clutch
$e=1$ is simply connected and has
\[
 H_2=\mathbb Z^{15}\oplus\mathbb Z/8,\qquad
 H_3=\mathbb Z/8,\qquad H_4=\mathbb Z^{15},\qquad\chi=32.
\]
Its rank-two cusp has period index $16$, so it is
outside the scalar iterations of the natural
index-one initial catalogue. This theorem concerns
the full stated whole-fibre sector. Component-supported
linearizations outside \eqref{maps:eq:D.GI1}, arbitrary extra
vertical modifications, logarithmic star Mumford
caps, and general smooth classification of its
nonsphere outputs are not asserted here.


\begin{theorem}[the general triangular-cusp component intersection form]\label{maps:src:D.2.47}
\StatementRecord{maps:src:D.2.47}{proved}{Historical events: Claude Fable 5.2 (Anthropic), effort high, 16 November 2026: carried as background without fault found. GPT-6 Astra (OpenAI), ChatGPT Pro mode, effort Pro, S6\_09\_PAPERS, 17 November 2026: targeted proof check. Also checked in this writing. Pre-continuation writing events: GPT-6 Astra (OpenAI), ChatGPT Pro mode, effort Pro, S6\_10\_PAPERS, 18 November 2026: targeted proof check. Continuation metadata events: GPT-6 Astra (OpenAI), ChatGPT Pro mode, effort Pro, S6\_10\_PAPERS continuation, 19 November 2026: used and rechecked in this run; status-vocabulary reconciliation; metadata reconciliation only; no new mathematical check. \emph{Pre-continuation recorded proof scope.} GPT-6 Astra (OpenAI, ChatGPT Pro mode), effort Pro, 2026-11-17 UTC; Claude Fable 5.2 (Anthropic), effort high, 16 November 2026; also checked in this writing by GPT-6 Astra (OpenAI), ChatGPT Pro mode, effort Pro, 18 November 2026}{the hypotheses, notation and definitions stated here; \cref{maps:src:D.2.43}}


Let $N_L$ be the smooth triangular Mumford cap of \ref{maps:src:D.2.43},
with period sublattice $L\subset\mathbb Z^2$ of finite index
$D>0$. Write $V=\mathbb Z^2/L$ for its vertices and
$E_i$, $i\in V$, for its reduced irreducible central divisors.
Their normalizations $\nu_i:P_i\to E_i\subset N_L$ are the
degree-six toric del Pezzo surface, even when a divisor has
self-gluings. Use the cyclic ray order
\[
 (r_0,\ldots,r_5)=
       (e_1,e_2,e_2-e_1,-e_1,-e_2,e_1-e_2).                 \tag{D47.1}\label{maps:eq:D47.1}
\]
For a coefficient vector $u=(u_i)_{i\in V}$ put
$E(u)=\sum_i u_iE_i$ and
$a_{i,r}(u)=u_{i+r}-u_i$. The entire integral trilinear
form on these compact divisor classes is
\[
\begin{split}
 I_L(u,v,w)=\sum_{i\in V}u_i\Bigl(
 &-\sum_{r=0}^5 a_{i,r}(v)a_{i,r}(w)\\
 &+\sum_{r=0}^5\bigl(a_{i,r}(v)a_{i,r+1}(w)
                    +a_{i,r+1}(v)a_{i,r}(w)\bigr)\Bigr),
\end{split}                                                   \tag{D47.2}\label{maps:eq:D47.2}
\]
where ray indices are cyclic. It is symmetric in all three
arguments, and the constant vector annihilates it.
Equivalently, the cubic is
\[
\begin{split}
 I_L(u,u,u)={}&6\sum_{i\in V}u_i^3
       -3\sum_{\{i,j\}\in\mathcal E}
                      (u_i^2u_j+u_iu_j^2)\\
       &+6\sum_{(i,j,k)\in\mathcal T}u_iu_ju_k.
\end{split}                                                   \tag{D47.3}\label{maps:eq:D47.3}
\]
Here $\mathcal E$ is the set of the $3D$ quotient edge
\emph{cells}, indexed by $i$ and the first three directions
in \eqref{maps:eq:D47.1}, and $\mathcal T$ consists of the $2D$ triangle
cells. Loops, multiple edges and repeated vertex labels in
a triangle are retained with their actual cell multiplicities.
They are not replaced by a simple graph or a simplicial
complex with distinct labels.
The Chern and Pontryagin evaluations on every component are
\[
 \left\langle c_2(TN_L),\nu_{i*}[P_i]\right\rangle=0,
 \qquad
 \left\langle p_1(TN_L),\nu_{i*}[P_i]\right\rangle=0.         \tag{D47.4}\label{maps:eq:D47.4}
\]
All these formulas remain the same when the cap is included
in a closed smooth complex threefold.

For $L=\mathbb Ze_1\oplus d\mathbb Ze_2$, the entire
component trilinear form is zero for every $d>0$.
It is not zero for arbitrary period lattices: for
$L=2\mathbb Z^2$ and its four vertices,
\[
 E_i^3=6,\qquad E_i^2E_j=-2\ (i\ne j),\qquad
 E_iE_jE_k=2\ (i,j,k\text{ distinct}).                      \tag{D47.5}\label{maps:eq:D47.5}
\]
This is the component part of the cohomology ring and
the indicated characteristic pairings. It is not asserted
to determine torsion products, the full fundamental-class
invariant, or a global diffeomorphism type.

\end{theorem}
\begin{proof}
\begin{enumerate}[label=\arabic*.,leftmargin=*]
\item Denote the six toric boundary curves of $P_i$ by $C_{i,r}$.
The fan \eqref{maps:eq:D47.1} is the fan of $\mathbb P^2$ subdivided
at its three torus-fixed points. Consecutive rays have
determinant one and $r_{j-1}+r_{j+1}=r_j$. Consequently
\[
 C_{i,r}^2=-1,\qquad
 C_{i,r}C_{i,r+1}=1,\qquad
 C_{i,r}C_{i,s}=0\quad(s\not\equiv r,r\pm1).               \tag{D47.6}\label{maps:eq:D47.6}
\]
These also follow directly from the monomial coordinates
on the six smooth affine toric charts.

\item For $j\ne i$, the pullback of the Cartier divisor $E_j$
to $P_i$ is precisely the sum of those boundary curves
whose adjacent label is $j$. The total reduced central
divisor is $\sum_jE_j=\operatorname{div}(q)$, so its
absolute line bundle on the cap is trivial. Restricting
this identity determines the remaining pullback and gives
the equality of line bundles, or their divisor classes,
\[
 \nu_i^*\mathcal O(E(u))
       =\mathcal O_{P_i}\left(
                   \sum_{r=0}^5(u_{i+r}-u_i)C_{i,r}\right).
                                                               \tag{D47.7}\label{maps:eq:D47.7}
\]
When an edge returns to the same component, its coefficient
is zero. In local normal-crossing coordinates, the same
irreducible Cartier divisor then contains both branches;
the second branch contributes exactly the correction to
the normal line of the first branch. Thus \eqref{maps:eq:D47.7} does
not apply embedded-divisor adjunction incorrectly to a
self-glued component.

\item The fundamental cycle of $E_i$ is $\nu_{i*}[P_i]$ with
degree one. Intersection with the compact divisor $E(u)$,
then the projection formula, therefore gives
\[
 E(u)E(v)E(w)=\sum_i u_i
           \int_{P_i}c_1\nu_i^*\mathcal O(E(v))
                       c_1\nu_i^*\mathcal O(E(w)).          \tag{D47.8}\label{maps:eq:D47.8}
\]
Equations \eqref{maps:eq:D47.6}--\eqref{maps:eq:D47.7} give \eqref{maps:eq:D47.2} integrally.
It is symmetric because these are products of actual
Cartier classes on the smooth threefold. A constant
coefficient vector has trivial restriction on every
$P_i$, so annihilates the form.

\item There is a relative-cohomology distinction here. The
divisors have compact support and define classes in
$H^2(N_L,\partial N_L;\mathbb Z)$; their triple product
is evaluated in relative degree six. Although
$\sum_iE_i$ has a trivial absolute line bundle on the
cap, its compact-support class need not be zero: the
boundary value of $q$ winds around the origin.
The assertion is that the constant vector annihilates
the component intersection form, not that the sum of
the component fundamental cycles vanishes in homology.
For an inclusion into a closed threefold, the same
compact cycles and local divisor restrictions compute
the same products.

\item \emph{The quotient cell formula and examples.}
Set $v=w=u$ in \eqref{maps:eq:D47.2} and expand. A vertex contributes
$6u_i^3$. Pairing the two directed incidences of every
edge gives $-3(u_i^2u_j+u_iu_j^2)$. Each triangle has
three corner incidences and contributes $6u_iu_ju_k$.
This proves \eqref{maps:eq:D47.3}, with all incidence multiplicities
retained, even when labels coincide.
One explicit set of triangle cells is
\[
 (i,i+e_1,i+e_2),\qquad
 (i+e_1+e_2,i+e_1,i+e_2)\quad(i\in V).                    \tag{D47.9}\label{maps:eq:D47.9}
\]
The integral trilinear form is already specified by
\eqref{maps:eq:D47.2}; no divisibility assumption based on a rational
polarization is needed.

\item If $e_1\in L$ and $V$ is the cyclic sequence of length
$d$ in the $e_2$ direction, the $e_1$ loop terms cancel
the vertex terms. The other two edge families have the
same adjacent labels and cancel the two triangle
families in \eqref{maps:eq:D47.9}. Hence the cubic, and by \eqref{maps:eq:D47.2}
the whole trilinear form, vanish for every $d$, including
$d=1,2$. For $L=2\mathbb Z^2$, each vertex sees the
other three labels twice in its hexagon. Substitution
in \eqref{maps:eq:D47.2} gives \eqref{maps:eq:D47.5}, so a vanishing conclusion
cannot be extrapolated from the cyclic period case.

\item \emph{The canonical volume and characteristic pairings.}
On the torus cover away from $q=0$ consider
\[
                    \omega=d\log X\wedge d\log Y\wedge dq.
                                                               \tag{D47.10}\label{maps:eq:D47.10}
\]
The period translations multiply $X,Y$ by integral
powers of $q$ and by nonvanishing functions of $q$.
The added logarithmic differential terms are multiples
of $dq$, so $\omega$ is invariant and descends.
In every unimodular triangular chart one can choose
coordinates with $q=xyz$, $X=x$, $Y=y$ up to the
unimodular monomial change. Then
\[
                    d\log x\wedge d\log y\wedge d(xyz)
                                 =dx\wedge dy\wedge dz.
                                                               \tag{D47.11}\label{maps:eq:D47.11}
\]
Thus the descended form extends without a zero or pole
over the central fibre, including self-intersections.
The cap has a trivial canonical bundle near all its
compact components.

\item The normalization $\nu_i$ is a holomorphic immersion:
each preimage of a crossing lies on one smooth local
coordinate branch. Even when its image has self-gluings,
the immersion normal line is therefore
$\mathcal N_{\nu_i}=K_{P_i}$, by the adjunction sequence
and the nonvanishing volume form. This normal line
must not be confused with $\nu_i^*\mathcal O(E_i)$,
which includes the other self-intersection branches.
The exact sequence
\[
       0\longrightarrow TP_i\longrightarrow\nu_i^*TN_L
                        \longrightarrow K_{P_i}\longrightarrow0
                                                               \tag{D47.12}\label{maps:eq:D47.12}
\]
now gives
\[
 \nu_i^*c_2(TN_L)=c_2(TP_i)-c_1(TP_i)^2,\qquad
 \nu_i^*p_1(TN_L)=2c_1(TP_i)^2-2c_2(TP_i).                 \tag{D47.13}\label{maps:eq:D47.13}
\]
The hexagonal surface has Euler number six (one for
each torus-fixed point) and $c_1^2=6$, also obtained
by squaring $\sum_rC_{i,r}$ using \eqref{maps:eq:D47.6}.
Both evaluations in \eqref{maps:eq:D47.13} are therefore zero,
proving \eqref{maps:eq:D47.4}. The calculation uses the normal line
of the immersion and is valid for every quotient
lattice, not only for embedded central components.

\item \emph{Exact verification and bounds.}
The program
\path{code/III/fourfolds_D2_component_ring.py}
enumerates all $127$ column Hermite forms
$\left(\begin{smallmatrix}a&b\\0&d\end{smallmatrix}\right)$
with $ad\le12$ and $0\le b<a$. For every one it
computes the integral triple products directly from
the six boundary restrictions, verifies symmetry
in all three indices and the constant annihilator,
and compares every cubic coefficient with the
quotient cell formula while retaining loops and
repeated labels. It also verifies the toric ray,
intersection and canonical-volume identities and
the examples \eqref{maps:eq:D47.5}. All matrices, cells and
nonzero integral products are retained in
\path{records/fourfolds_D2_component_ring.json}.
The proof above applies to arbitrary finite index;
the finite computation tests the conventions and
does not supply an unstated global index bound.
\hfill$\square$
\end{enumerate}
\end{proof}


\begin{theorem}[complete smooth classification of the simply connected weighted Engel family]\label{maps:src:D.2.48}
\StatementRecord{maps:src:D.2.48}{proved}{Historical events: Claude Fable 5.2 (Anthropic), effort high, 16 November 2026: carried as background without fault found. GPT-6 Astra (OpenAI), ChatGPT Pro mode, effort Pro, S6\_09\_PAPERS, 17 November 2026: targeted proof check. Also checked in this writing. Pre-continuation writing events: GPT-6 Astra (OpenAI), ChatGPT Pro mode, effort Pro, S6\_10\_PAPERS, 18 November 2026: targeted proof check. Continuation metadata events: GPT-6 Astra (OpenAI), ChatGPT Pro mode, effort Pro, S6\_10\_PAPERS continuation, 19 November 2026: used and rechecked in this run; status-vocabulary reconciliation; metadata reconciliation only; no new mathematical check. \emph{Pre-continuation recorded proof scope.} GPT-6 Astra (OpenAI, ChatGPT Pro mode), effort Pro, 2026-11-17 UTC; Claude Fable 5.2 (Anthropic), effort high, 16 November 2026; also checked in this writing by GPT-6 Astra (OpenAI), ChatGPT Pro mode, effort Pro, 18 November 2026}{the hypotheses, notation and definitions stated here; \cref{maps:src:D.2.34}; \cref{maps:src:D.1.11}; \cref{maps:eq:D.WE3}; \cref{maps:eq:D.WI5}; \cref{maps:eq:D.WI4}; \cref{maps:src:D.2.44}; \cref{maps:src:D.2.45}; \cref{maps:eq:D.WI7}; \cite[Definition 1.7 and Theorem 3.1, pp.239,243--246]{mapsKaplan79}; \cref{maps:eq:D.WI9}; \cref{maps:eq:D.WI11}; \cref{maps:eq:D.WI13}; \cref{maps:eq:D.WI10}; \cref{maps:eq:D.WI12}}


Take the weight-$d$ Engel construction of \ref{maps:src:D.2.34}, for any $d>0$,
with both additive sites replaced by good logarithmic caps.  Retain
all the actual marked data
\[
 m_0=3k,\qquad m_1=4\ell,\qquad
 v_i=r_i\delta+n_i u_i,\qquad
 P=m_1n_0+m_0n_1=\pm1 .
 \tag{D.WI1}\label{maps:eq:D.WI1}
\]
The manifold is simply connected.  Its complex first and second
Chern classes, spin characteristic class, and refined class satisfy
\[
 c_1(X)=c_2(X)=\lambda(X)=0,\qquad P_{\mathrm{ref}}(X)=0.
 \tag{D.WI2}\label{maps:eq:D.WI2}
\]
Here $P_{\mathrm{ref}}$ denotes Zhubr's class, to distinguish it from
the Seifert determinant $P$ in \eqref{maps:eq:D.WI1}.

The classes $x_i=\operatorname{PD}(D_i)$ of the $d$ cusp components,
indexed cyclically, generate $H^2(X;\mathbb Z)$ with sole relation
$\sum_i x_i=0$.  The complex ruling $f=H-E_1$ has image
$s=[a\wedge b]$, a generator of $\operatorname{Tor}H_2(X)=\mathbb Z/d$.
Put $\kappa=\operatorname{PD}(f)$.  With repetitions retained when
$d=2$, the complete degree-two product is
\[
 x_i x_j=
 \bigl(\mathbf1_{j=i+1}+\mathbf1_{j=i-1}-2\mathbf1_{j=i}\bigr)\kappa.
 \tag{D.WI3}\label{maps:eq:D.WI3}
\]
For $d=1$ the sole component class is zero.  For $d=2$ the two
off-diagonal incidences both occur, and every product in \eqref{maps:eq:D.WI3}
is zero since $2\kappa=0$.  The integral cubic is zero for every $d$.
Here is the complete fundamental-class invariant, including its
two- and three-primary coordinates. Put
$r=d-1$, $D_2=d\gcd(d,2)$, and $D_3=d\gcd(d,3)$.
For $d>1$ there is a class $u\in H^2(X;\mathbb Z/d)$
which, together with $x_0,\ldots,x_{d-2}$, fixes the splitting
$H_2(X)=\mathbb Z^r\oplus\mathbb Z/d$ described below. In
\[
 H_6(K(\mathbb Z^r\oplus\mathbb Z/d,2);\mathbb Z)
 \simeq
 \mathbb Z^{\binom{r+2}{3}}\oplus
 (\mathbb Z/d)^{\binom{r+1}{2}}\oplus
 (\mathbb Z/D_2)^r\oplus\mathbb Z/D_3
 \tag{D.WI14}\label{maps:eq:D.WI14}
\]
the complete class $\mu$ is determined by
\[
 \begin{aligned}
 \langle x_i x_j x_k,[X]\rangle&=0,\\
 \langle u x_i x_j,[X]\rangle
    &=\mathbf1_{j=i+1}+\mathbf1_{j=i-1}-2\mathbf1_{j=i}
                                                     &&\pmod d,\\
 \langle \mathcal Q_2(u)x_i,[X]\rangle&=-12             &&\pmod {D_2},\\
 \langle \mathcal Q_3(u),[X]\rangle&=0                 &&\pmod {D_3}.
 \end{aligned}
 \tag{D.WI15}\label{maps:eq:D.WI15}
\]
The indices in the second line are cyclic modulo $d$; those
used as independent coordinates run from zero to $d-2$.
Here $\mathcal Q_2$ is the ordinary square if $d$ is odd and
the Pontryagin square with values modulo $2d$ if $d$ is even.
Likewise $\mathcal Q_3$ is the ordinary cube if $3\nmid d$
and Thomas's Pontryagin third power with values modulo $3d$
if $3\mid d$.

Consequently, for every $d>0$ all the simply connected data
\eqref{maps:eq:D.WI1} give one oriented diffeomorphism type $M_d$, namely
the unique type with
\[
 H_2=\mathbb Z^{d-1}\oplus\mathbb Z/d,\qquad b_3=0,\qquad
 P_{\rm ref}=0,\qquad \mu\ \text{given by \eqref{maps:eq:D.WI15}}.
 \tag{D.WI16}\label{maps:eq:D.WI16}
\]
One has $M_1=S^6$, and for $d$ prime to six this is the
type already identified in \ref{maps:src:D.1.11}. The types for different
$d$ are distinct, since their Euler numbers are $2d$.

\end{theorem}


\begin{proof}
\begin{enumerate}[label=\arabic*.,leftmargin=*]
\item \emph{Primitive component classes.}
Write $E$ for the complement of the interior of the cusp tube.
The actual finite-cap presentation before adjoining its cusp meridian is
\[
 H_1(E)=
 \langle\gamma,x_0',x_1'\mid m_0x_0'=n_0\gamma,\ 
                                   m_1x_1'=n_1\gamma\rangle .
\]
The determinant condition implies
$\gcd(m_0,m_1)=\gcd(m_i,n_i)=1$ and $H_1(E)=\mathbb Z$.
The cusp meridian is $x_0'+x_1'$ up to the common orientation sign.
It is primitive because adjoining its relation gives $H_1(X)=0$.
Angular differences of adjacent toric meridians are killed in
$H_1(E)$.  Thus the relative Thom map is
\[
 H^1(E)\longrightarrow H^2(X,E)=\mathbb Z^d,\qquad
                   1\longmapsto\pm(1,\ldots,1).
 \tag{D.WI4}\label{maps:eq:D.WI4}
\]
The cokernel of the resulting injection
$\mathbb Z^d/\mathbb Z(1,\ldots,1)\hookrightarrow H^2(X)$
embeds in $H^2(E)$, which is torsion-free by the integral universal
coefficient theorem.  It has rank zero by \ref{maps:src:D.2.34}, so it vanishes.
In particular the entire fibre class is
$f_X^*[\mathrm{point}]=\sum_i x_i=0$.
For a reduced multiple fibre $D$ one has
$m[D]=f_X^*[\mathrm{point}]=0$; the group $H^2(X)$ is free,
so $[D]=0$ as well.

\item The common ruling in the $d$-period cusp is the specialization
of $a\wedge b$, with its complex orientation.  This is an integral
unit-period annulus calculation: that torus remains in one gamma
component and limits to the two adjacent external curves
$E_2+E_3=H-E_1$.  The three support functions
$r(r-1)/2$, $rs$, and $s(s-1)/2$ evaluate on that ruling as
$(-1,0,0)$, exactly their actual multiplier windings on
$a\wedge b$; its dual-torus coordinate is zero.
There is no division by $d$ or saturation at this step.
The primitive functional of \ref{maps:src:D.2.34} evaluates it as
$\xi_d(a\wedge b)=1$.  The matrix \eqref{maps:eq:D.WE3} shows that $s$ has
exact order $d$ in the simply connected sector.  The actual
external intersections and the principal sum of components
therefore give \eqref{maps:eq:D.WI3}.  The adjacent-pair formula counts both
edges when they have the same pair of labels.  Since every
product is torsion, its further product with $H^2(X)$ vanishes
in $H^6(X)=\mathbb Z$.

\item \emph{The differential and its literal critical cycle.}
The normalization immersion $j_i:dP_6\longrightarrow X$ has
normal line $K_{dP_6}$.  Indeed the principal local central
divisor is the product of the normal coordinates to all its
branches; on a fixed normalized branch its normal transition
is the inverse of the six toric boundary transitions.  This
includes both conductor branches of each self-identification.
Consequently
\[
 j_i^*c_1(TX)=c_1(TdP_6)+c_1(K_{dP_6})=0,\qquad
 \langle j_i^*c_2(TX),[dP_6]\rangle=6-6=0 .
 \tag{D.WI5}\label{maps:eq:D.WI5}
\]
Restriction of $H^2(X)$ to the normalized components is injective.
For a component combination with coefficients $u_i$, its
restriction to component $i$ is
$\sum_r(u_{i+r}-u_i)C_r$.  Pairing with the anticanonical ample
class, which has degree one on each of the six boundary curves,
gives the component-graph Laplacian.  The graph is connected, so
zero restrictions force the $u_i$ to be constant.  That is
precisely the already removed principal relation.  Equation
\eqref{maps:eq:D.WI5} now gives $c_1(TX)=0$ integrally.

\item Remove the divisorial multiple-fibre factors from $df_X$.  Its
corrected target is the line
\[
 L=f_X^*T\mathbb P^1\otimes
                      \mathcal O\bigl(-\sum (m_i-1)D_i^{\rm fin}\bigr).
\]
It is topologically trivial by \eqref{maps:eq:D.WI4}.  The abstract rank-two
kernel $V$ is locally free, including at $xy$ and $xyz$ charts,
and has the nowhere-zero abstract section given by $\delta$.
At a node the kernel columns are $(x,-y,0)$ and $(0,0,1)$;
at a triple point they are $(x,-y,0)$ and $(x,0,-z)$.
The constant scalar-circle weight vector is a nonzero constant
combination in these abstract bases.  At the good finite caps
the corrected differential is a submersion and $\delta$ remains
a nonzero tangent vector.  Since
$c_1(V)=c_1(TX)-c_1(L)=0$, its quotient by this section is a
trivial complex line.  Thus $V$ and $L$ admit full smooth frames.
They give the reference frame $\rho$ on the complement of the
literal critical curve cycle $C$.

\item The local Bott clutch at $xy$ has coefficient one:
its two normal columns are
$\bigl(\begin{smallmatrix}x&\overline y\\-y&\overline x\end{smallmatrix}\bigr)$.
The $xyz$ Hilbert--Burch minors give multiplicity one on each
of its three axes and no further codimension-two contribution.
The good finite logarithmic caps have no residual critical
cycle after correcting the target.  Hence, with the complex
orientation and the spin normalization of \ref{maps:src:D.2.44}--\ref{maps:src:D.2.45},
\[
                 \lambda(TX,\rho)=-\operatorname{PD}(C).
 \tag{D.WI6}\label{maps:eq:D.WI6}
\]

\item Only a local covering is used to compute $[C]$.
In the positive period basis
$N b=a$, $N\gamma_1=\delta$, the weight-one cusp's literal
three-curve sum has regular-fibre preimage
$a\wedge b+b\wedge\delta+\delta\wedge\gamma_1$.
The weight-$d$ cusp is its degree-$d$ etale period cover, with
$\gamma_d=d\gamma_1$.  Transfer therefore gives the actual
preimage of its full $3d$-curve sum as
\[
 d\,a\wedge b+d\,b\wedge\delta+\delta\wedge\gamma_d.
 \tag{D.WI7}\label{maps:eq:D.WI7}
\]
The angular vector $\delta$ here is fixed by
$N\gamma_d=d\delta$; this convention removes the possible
opposite choice for the original scalar-circle parameter.
Since $\xi_d$ has values $1,0,-6d$ on these three bivectors,
\eqref{maps:eq:D.WI7} has image $-5d\,s=0$.  Thus $C$ is integrally
nullhomologous and $c_2(TX)=\lambda(TX)=0$.
No covering of either finite logarithmic cap is assumed.

\item \emph{A finite-cap detector with its actual reference parity.}
If $d$ is odd, $2P_{\rm ref}=\lambda=0$ already gives
$P_{\rm ref}=0$.  Suppose $d$ is even.  In the order-four
finite cap, put $\beta=[a]$ and $\eta=[u_1]$ in its two
central base directions.  Its marked base lattice is
\[
 \langle\beta,\eta,g\mid m_1g=-2r_1\beta+n_1\eta\rangle .
 \tag{D.WI8}\label{maps:eq:D.WI8}
\]
Choose integers $B,C$ with $m_1C=-2r_1+n_1B$ and define the
primitive character by $\chi(\beta)=1$, $\chi(\eta)=B$,
$\chi(g)=C$.  Such integers exist because $\gcd(m_1,n_1)=1$.
Its connected fibre $N$ in the central four-manifold is an
elliptic-torus mapping torus.  In its base kernel
$h=g-C\beta$ and $q=\eta-B\beta$ obey $m_1h=n_1q$.
A primitive generator has the root-loop exponent inverse to
$n_1$ modulo $m_1$, hence odd modulo four.  The elliptic
monodromy of $N$ is therefore the order-four matrix or its
inverse.  Its Wang sequence gives
\[
             H_1(N;\mathbb Z)=\mathbb Z\oplus\mathbb Z/2.
 \tag{D.WI9}\label{maps:eq:D.WI9}
\]
Moreover $\chi(\delta)=-2$, so the actual scalar tangent
vector is transverse to $N$ in the reduced finite fibre.
Projecting $J\delta,\operatorname{Re}v_2,\operatorname{Im}v_2$
to $TN$ along $\delta$ identifies the stable tangent frame
induced by $\rho$ with an actual tangent parallelization.
The normal three-frame is $\delta$ followed by the two real
directions of $L$.

\item For an actual parallelization of a closed oriented
three-manifold, its stable framed coefficient modulo two is
$1+\dim H_1(N;\mathbb F_2)$.  The spin-compatible Kaplan
zero-/two-handle proof was given in the finite-II frame lemma:
the relative $\lambda$ reduces to the relative Euler number,
and the even intersection matrix has even mod-two rank.
It applies to every prescribed spin structure
\cite[Definition 1.7 and Theorem 3.1, pp.239,243--246]{mapsKaplan79}.
Equation \eqref{maps:eq:D.WI9} therefore gives reference parity
\[
                         \epsilon=[N]_\rho=1\pmod2 .
 \tag{D.WI10}\label{maps:eq:D.WI10}
\]

\item Retain the actual Thom chains, not only their matrix classes.
Let $K_{a0},K_{a1}$ be the first columns of the two blocks in
\eqref{maps:eq:D.WE3}.  Directly in its four target coordinates,
\[
                    d s=-2\ell P\,K_{a0}-kP\,K_{a1}.
 \tag{D.WI11}\label{maps:eq:D.WI11}
\]
The chain at cap zero misses $N$, whereas the $a$-Thom
chain in cap one pairs with $N$ by $\chi(a)=1$.
Thus $\operatorname{lk}(N,s)=-kP/d$.
The determinant equation implies that $k$ is odd, so $N$
generates the two-primary summand of $H_3(X)$.
The specialization chains for \eqref{maps:eq:D.WI7} are supported in the
cusp and miss $N$.  Equations \eqref{maps:eq:D.WI7}--\eqref{maps:eq:D.WI11} then give
the actual integer link
\[
                          \ell(N,C)=5kP,
 \tag{D.WI12}\label{maps:eq:D.WI12}
\]
up to the common harmless orientation sign.  It is odd.

\item Finally the relative refined-class computation extends with
reference parity $\epsilon$ in the form
\[
 \operatorname{lk}(N,\operatorname{PD}P_{\rm ref})
                 =-\frac{\ell(N,C)}2-\frac{\epsilon}2
                                      \quad\text{in }\mathbb Q/\mathbb Z .
 \tag{D.WI13}\label{maps:eq:D.WI13}
\]
Indeed, choose Zhubr's admissible three-frame on a genuinely
embedded torsion sphere obtained from an odd multiple of $N$
by spin-compatible ambient two-handle surgery.  Its complement
spin three-frame extends over the trace.  The induced stable
framed class on $N$ is even.  The comparison with $\rho$ has
Spin(6) degree $q_0\equiv\epsilon\pmod2$.
For a bounding four-chain $B_4$ with boundary
$\operatorname{ord}(N)N$, the exact relative transgression is
$-B_4\cdot C=2q+\operatorname{ord}(N)q_0$.
Division of this integer identity gives \eqref{maps:eq:D.WI13}, retaining
the actual rational link before reduction modulo one.
This is the proof of the relative lemma in \ref{maps:src:D.2.45} with its
reference-parity term restored, not an application of the
unproved general change-of-frame conjecture.
Free $H_2$ causes no change: $C$ is nullhomologous, and
$2P_{\rm ref}=\lambda=0$ already kills the free part.

\item The two odd halves in \eqref{maps:eq:D.WI13} cancel by \eqref{maps:eq:D.WI10} and
\eqref{maps:eq:D.WI12}.  The two-primary part of $P_{\rm ref}$ is zero,
and its odd-primary part was already zero from
$2P_{\rm ref}=0$.  This proves \eqref{maps:eq:D.WI2}.
\end{enumerate}
\end{proof}

\begin{proof}[Proof of the complete invariant and classification]
\begin{enumerate}[label=\arabic*.,leftmargin=*]
\item \emph{A global finite-coefficient class, retaining boundary torsion.}
At weight one the regular bundle $V_1$ over the peripheral
graph has free $H_1$ and $H_2$. Its degree-two group has the
primitive fibre quotient generated by $s$, together with the
five marked graph generators in \eqref{maps:eq:D.WE11}. Choose the integral
cohomology class $\alpha_1$ which evaluates as $\xi_1$ on the
fibre part and zero on every graph generator.

\item On the weight-one cusp use the actual toric line with support
\[
                 \psi(r,s)=-\frac{r(r-1)}2+3s(s-1).
 \tag{D.WI17}\label{maps:eq:D.WI17}
\]
Its multiplier windings give $\xi_1$ in the positive period
marking used in \eqref{maps:eq:D.WI7}. Its boundary value agrees integrally
with $\alpha_1$: their fibre restrictions agree, and both
vanish on the two angular Wang sweeps, which bound in the
cusp. These are all the remaining generators of $H_2(M_1)$.
Also $H_1(M_1)$ is free, so the integral universal coefficient
theorem leaves no hidden $\operatorname{Ext}$ difference.

\item Before the finite caps are attached, the core and cusp have
the genuine degree-$d$ period covering defined by
$\gamma_d=d\gamma_1$. Pull back both classes along that
cover. They then agree on the entire boundary $M_d$, not
only on its regular-fibre restriction. In particular this
argument retains the $\mathbb Z/d$ term in $H_1(M_d)$ and
the resulting torsion in $H^2(M_d)$. The core pullback
$\alpha_d$ has fibre functional $\xi_d$ and zero values on
the five graph generators: the covering fixes the first
three such generators and multiplies the two $u_i$ generators
by $d$.

\item At a finite cap let $\pi:Y_i\longrightarrow Q_i$ be its
actual normal circle bundle. The obstruction for
$\alpha_d\bmod d$ to extend across the disk bundle is its
Gysin image
\[
                  \pi_!(\alpha_d\bmod d)\in
                         H^1(Q_i;\mathbb Z/d).
 \tag{D.WI18}\label{maps:eq:D.WI18}
\]
Its values on the complete marked generators
$a,\gamma,g$ of $H_1(Q_i)$ are precisely the first-row
coefficients of \eqref{maps:eq:D.WE3}:
\[
 \begin{array}{c|ccc}
       &a&\gamma&g\\ \hline
 i=0&-2dn_0&-6dr_0&0\\
 i=1&-3dn_1&-6dr_1&0.
 \end{array}
 \tag{D.WI19}\label{maps:eq:D.WI19}
\]
The root-loop column $g$ is retained. This is the Thom
definition of fibre integration: a loop in $Q_i$ sweeps its
normal circle, and pairing with $\alpha_d$ gives the displayed
actual normal-column coefficient. Every coefficient is zero
modulo $d$. Since $H^1(Q_i;\mathbb Z/d)$ is
$\operatorname{Hom}(H_1(Q_i),\mathbb Z/d)$, \eqref{maps:eq:D.WI18} vanishes,
even when the integral normal line has torsion. The circle
Gysin sequence therefore supplies an extension over each
finite cap. Mayer--Vietoris now glues them with the pulled-back
cusp line to a class $u\in H^2(X;\mathbb Z/d)$.
Only the core and cusp were covered; neither finite fill was
assumed to be a cover of a weight-one fill.

\item On every normalized cusp component the integral representative
of its restriction is
\[
              \eta=-5H+6E_1-E_3,\qquad
       \eta(E_1,E_2,E_3)=(-6,0,1),\qquad
       \eta(f)=1,\qquad \eta^2=-12.
 \tag{D.WI20}\label{maps:eq:D.WI20}
\]
The three degrees follow directly by taking the differences
of \eqref{maps:eq:D.WI17} across the internal, diagonal, and scalar-fixed
walls, as in \ref{maps:src:D.1.9}. In particular $u(s)=1$.
Choose free lifts in $H_2(X)$ annihilated by $u$, subtracting
appropriate multiples of $s$. This specifies the splitting
in \eqref{maps:eq:D.WI14}--\eqref{maps:eq:D.WI16}. The integral classes $x_i$ form its
dual free basis, and $u$ is the standard cyclic coordinate.

\item \emph{The full mixed coordinates.}
The first mixed value in \eqref{maps:eq:D.WI15} follows from \eqref{maps:eq:D.WI3} and
$\langle u\kappa,[X]\rangle=1$. The Pontryagin square of the
reduction of an integral degree-two class is its ordinary
square with the full target modulus $2d$; this is the
$p=2$ case of the actual cochain construction
\cite[equations (2.1), (2.8), and (3.4), pp.267--268]{ffThomas56}.
By naturality and the projection formula for the normalized
component fundamental class,
\[
 \langle \mathcal Q_2(u)x_i,[X]\rangle
       =\langle\mathcal Q_2(j_i^*u),[dP_6]\rangle
       =\eta^2=-12\pmod {D_2}.
 \tag{D.WI21}\label{maps:eq:D.WI21}
\]
This computes the coefficient modulo $2d$ when $d$ is even;
reducing the ordinary square merely modulo $d$ would lose
precisely the extra order-two information. The relation
$\sum_i x_i=0$ is consistent with \eqref{maps:eq:D.WI21}, since
$12d$ is divisible by $D_2$.

\item \emph{The pure coordinate.}
The following actual scalar-circle calculation evaluates
the full Pontryagin third power. It uses the class just
constructed and its specified component restrictions.
\subsubsection*{The full Pontryagin cube of the weighted Engel torsion class}
We use the actual class $u\in H^2(X_d;\mathbb Z/d)$ whose
restrictions to the normalized cusp components are specified
in \ref{maps:src:D.2.48}. In particular its degree on every scalar-fixed
rational curve $C_j$, $j\in\mathbb Z/d$, is one, and its
component restriction is represented by the integral toric
support
\[
                  \psi(r,s)=-\frac{r(r-1)}2+3s(s-1).
                                                               \tag{D.TP1}\label{maps:eq:D.TP1}
\]
The calculation below uses these actual restrictions. Merely
choosing an arbitrary degree-two class that is a unit on
torsion would not fix its pure coordinate in the presence
of the free summand.

\item The scalar circle action is semifree. On a finite cap of
order $m_i$ its gamma translation has primitive numerator
$n_i$, with $\gcd(m_i,n_i)=1$. A nontrivial root power cannot
be compensated by the scalar direction, because that would
require $j n_i/m_i$ integral. Thus these caps are scalar-free.
The cusp has exactly the $d$ fixed curves $C_j$; their normal
bundles are $\mathcal O(-1)\oplus\mathcal O(-1)$ with weights
$(1,-1)$, and the other cusp orbits are free. These are the
marked geometric inputs of the weighted family.

\item Because $X_d$ is simply connected, $u$ has an equivariant
lift $\widehat u\in H^2_{S^1}(X_d;\mathbb Z/d)$.
In the Borel spectral sequence the only possible
differentials out of ordinary degree two have targets
involving $H^1(X_d;\mathbb Z/d)$ or $H^3(BS^1;\mathbb Z/d)$,
and both vanish. Write
\[
         \widehat u|_{C_j}=x+\lambda_j t,
         \qquad x=[C_j]^*,\quad t\in H^2(BS^1;\mathbb Z/d).
                                                               \tag{D.TP2}\label{maps:eq:D.TP2}
\]
The difference of any two equivariant lifts of a fixed
ordinary class on the simply connected normalized
$dP_6$ is a multiple of $t$. Hence the differences of
the characters on its two scalar-fixed curves can be
read from any toric equivariant lift of \eqref{maps:eq:D.TP1}.
The fixed edges are the translates of the edge from
$(0,j)$ to $(0,j+1)$; the support slope in their scalar
direction is
\[
               \psi(0,j+1)-\psi(0,j)=6j.
                                                               \tag{D.TP3}\label{maps:eq:D.TP3}
\]
The fibre character of the associated monomial
trivialization is the negative of this slope, up to
one global character convention. Equivalently the
intervening weight-one orbit curve has degree $-6$.
Thus, for one common sign $\epsilon=\pm1$,
\[
                    \lambda_j=\lambda_0+6\epsilon j
                                 \quad\text{in }\mathbb Z/d.
                                                               \tag{D.TP4}\label{maps:eq:D.TP4}
\]
This equation is compatible with the period $j\mapsto j+d$.
It does not assert that the integral support line itself
descends with an integral global circle linearization:
the needed character equality is modulo $d$ and follows
from the existing global equivariant cohomology lift.

\item Suppose now that $3\mid d$. Thomas's operation is defined
with precisely the required full coefficients:
\[
 \mathcal P_3:H^2(-;\mathbb Z/d)
                    \longrightarrow H^6(-;\mathbb Z/(3d)).
                                                               \tag{D.TP5}\label{maps:eq:D.TP5}
\]
This is \cite[equation (2.1), p.267]{ffThomas56} with
$p=3$ and $m=d/3$. On the reduction of an integral
class $z$, it equals the reduction of $z^3$ modulo
$3d$. Indeed the correction term in Thomas's actual
cochain formula (3.4) is zero for an integral cocycle,
and its first term is the ordinary threefold diagonal
product \cite[equations (3.4)--(3.6), p.268]{ffThomas56}.
Each restriction \eqref{maps:eq:D.TP2} has such an integral lift.

\item Apply equivariant localization with coefficients
$\mathbb Z/(3d)$ to $\mathcal P_3(\widehat u)$.
No field assumption is needed here: the complement
is scalar-free, and each normal Euler class becomes
a unit after inverting $t$. On a fixed curve it is
\[
                    (-x+t)(-x-t)=-t^2,
\]
because $x^2=0$. Integrating the term of ordinary
curve degree two gives
\[
 \left\langle\mathcal P_3(u),[X_d]\right\rangle
                      =-3\sum_{j=0}^{d-1}\lambda_j^2
                                  \quad\text{in }\mathbb Z/(3d).
                                                               \tag{D.TP6}\label{maps:eq:D.TP6}
\]
The right side is independent of the integer
representatives of the characters: $3\mid d$ makes
the cube of an integral lift well defined modulo
$3d$. The base cohomology polynomial ring injects
into its localization at $t$, so the degree-zero
identity is an identity in the original coefficient
ring, not just a localized statement.

\item For arbitrary integers $d>0,\lambda_0$ and either
sign in \eqref{maps:eq:D.TP4}, the exact identity is
\[
 \sum_{j=0}^{d-1}(\lambda_0+6\epsilon j)^2
   =d\left(\lambda_0^2+6\epsilon\lambda_0(d-1)
                         +6(d-1)(2d-1)\right).              \tag{D.TP7}\label{maps:eq:D.TP7}
\]
Substitution into \eqref{maps:eq:D.TP6} proves
\[
                 \left\langle\mathcal P_3(u),[X_d]\right\rangle=0
                                  \quad\text{in }\mathbb Z/(3d).
                                                               \tag{D.TP8}\label{maps:eq:D.TP8}
\]
Thus the complete pure cyclic fundamental-class coordinate
is zero, including the order-three subgroup invisible to
the ordinary mod-$d$ cube. For $3\nmid d$ the same
localization calculation with the ordinary cube gives
zero modulo $d$. Mixed coordinates are separate from
this calculation; in particular it does not discard
the order-$2d$ mixed Pontryagin-square terms when $d$
is even.

\item The universal integer identities, support slope,
single-curve Euler/cube calculation and an explicitly
bounded supplementary check are executed in
\path{code/III/fourfolds_D48_thomas.py}.
The supplementary range is $d=3,6,\ldots,120$,
$\lambda_0=0,\ldots,d-1$ and both signs, giving
$4920$ exact checks. The arbitrary-$d$ conclusion
comes from \eqref{maps:eq:D.TP7}, not that finite range.

\item \emph{Why these detectors are complete.}
We include the low-degree integral target calculation to
exclude an undetected finite subgroup. For a cyclic group,
\[
 \begin{aligned}
 H_2(K(\mathbb Z/d,2);\mathbb Z)&=\mathbb Z/d,\\
 H_4(K(\mathbb Z/d,2);\mathbb Z)&=\mathbb Z/D_2,\\
 H_6(K(\mathbb Z/d,2);\mathbb Z)&=\mathbb Z/D_3 .
 \end{aligned}
 \tag{D.WI22}\label{maps:eq:D.WI22}
\]
The degree-four assertion follows from Whitehead's exact
sequence: on $K(A,2)$ its two adjacent homotopy groups
vanish, so $H_4=\Gamma(A)$
\cite[Theorem 1 and Section 2]{mapsWhitehead50}.
For $A=\mathbb Z/d$, the universal quadratic generator
$q(1)$ has relations $d^2q(1)=2dq(1)=0$.
Conversely $q(j)=j^2$ modulo $\gcd(d^2,2d)=D_2$
is well defined and quadratic. Thus these are its complete
relations and the image of $\mathbb {CP}^2$ generates.
The ordinary or Pontryagin square evaluates on that image
as one modulo $D_2$, so it detects the entire group.

\item For degree six, Zhubr's epimorphism in Section 3.11 and
his two free bordism generators in Section 3.12 show that
the spin-realizable part for a cyclic coefficient group
is generated in homology by the image of
$[\mathbb {CP}^3]$: his other generator has zero
fundamental-class image. For odd-primary cyclic groups
this is all of $H_6$, since the only obstruction in his
Lemma 3.7 has mod-two target. For a two-power cyclic
group Section 3.12 explicitly gives
$\ker\alpha=H_6=\mathbb Z/2^a$
\cite[Lemmas 3.7 and 3.11, and Section 3.12,
pp.805--807]{manZhubr75}.
The decomposition of $\mathbb Z/d$ into its relatively
prime primary factors gives the corresponding product
of Eilenberg--MacLane spaces. All cross tensor and
$\operatorname{Tor}$ terms vanish, since their positive
homology has relatively prime orders. Hence $H_6$ is
cyclic and is generated by the image $z$ of
$[\mathbb {CP}^3]$ for every $d$.

\item The order of $z$ divides both $d^3$ and $3d$.
Indeed the class $d h$ on $\mathbb {CP}^3$ becomes zero
modulo $d$, while its integral cube is $d^3$.
On $\mathbb {CP}^1\times\mathbb {CP}^2$ the class
$d x+y$ reduces to a class factoring through
$\mathbb {CP}^2$, whose degree-six image is zero,
while its integral cube evaluates to $3d$.
Thus the order divides
$\gcd(d^3,3d)=d\gcd(d,3)=D_3$.
The opposite bound is detected by the ordinary cube when
$3\nmid d$, or by Thomas's operation when $3\mid d$:
on the reduced hyperplane class of $\mathbb {CP}^3$
the full target value is one
\cite[equation (2.8), p.267]{ffThomas56}.
This proves the last assertion of \eqref{maps:eq:D.WI22} with its
stated perfect detector.

\item Finally
$K(\mathbb Z^r\oplus\mathbb Z/d,2)
 =(\mathbb {CP}^{\infty})^r\times K(\mathbb Z/d,2)$.
The first factor has torsion-free homology in even
degrees and zero homology in odd degrees. The integral
Kunneth formula has no $\operatorname{Tor}$ term and
gives exactly \eqref{maps:eq:D.WI14}. Integral cubic monomials,
$u x_i x_j$, $\mathcal Q_2(u)x_i$, and
$\mathcal Q_3(u)$ are the respective perfect coordinate
detectors. Therefore \eqref{maps:eq:D.WI15} computes all of $\mu$,
including the extra factors two and three.

\item \emph{Application of smooth classification.}
The manifold is simply connected and spin, has
$H_2=\mathbb Z^{d-1}\oplus\mathbb Z/d$ and middle Betti
number zero by \ref{maps:src:D.2.34}, refined class zero by the first
proof, and the complete fundamental-class invariant
\eqref{maps:eq:D.WI15}. These are exactly the four invariants in
Zhubr's Theorem 1, with the definitions on pp.794--795
\cite[Theorem 1]{manZhubr75}. Their values are independent
of every marked parameter in \eqref{maps:eq:D.WI1}.
Existence is supplied by the constructed manifolds,
and the injective classification gives one oriented
diffeomorphism type for each $d$. This proves \eqref{maps:eq:D.WI16}
by an applied complete classification, rather than by
ordinary cohomology or an asserted connected-sum model.
\end{enumerate}
\end{proof}



\begin{theorem}[unbounded canonical star logs with the full integral boundary map]\label{maps:src:D.2.49}
\StatementRecord{maps:src:D.2.49}{proved}{Historical events: Claude Fable 5.2 (Anthropic), effort high, 16 November 2026: carried as background without fault found. GPT-6 Astra (OpenAI), ChatGPT Pro mode, effort Pro, S6\_09\_PAPERS, 17 November 2026: targeted proof check. Also checked in this writing. Pre-continuation writing events: GPT-6 Astra (OpenAI), ChatGPT Pro mode, effort Pro, S6\_10\_PAPERS, 18 November 2026: targeted proof check. Continuation metadata events: GPT-6 Astra (OpenAI), ChatGPT Pro mode, effort Pro, S6\_10\_PAPERS continuation, 19 November 2026: used and rechecked in this run; status-vocabulary reconciliation; metadata reconciliation only; no new mathematical check. \emph{Pre-continuation recorded proof scope.} GPT-6 Astra (OpenAI, ChatGPT Pro mode), effort Pro, 2026-11-17 UTC; GPT-6 Astra (OpenAI, ChatGPT Pro mode), S6\_06\_BLIND, effort Pro, 2026-11-17; Claude Fable 5.2 (Anthropic), effort high, 16 November 2026; also checked in this writing by GPT-6 Astra (OpenAI), ChatGPT Pro mode, effort Pro, 18 November 2026}{the hypotheses, notation and definitions stated here; \cref{maps:src:D.2.27}; \cref{maps:src:D.2.41}; \cref{maps:src:D.2.46}; \cref{maps:eq:D.GI10}; \cref{maps:src:D.2.26}; \cref{maps:src:D.2.24}; \cref{maps:eq:D.SL11}; \cref{maps:eq:D.SL10}; \cref{maps:eq:D.SL12}; \cref{maps:eq:D.SL13}; \cref{maps:eq:D.SL15}; \cref{maps:eq:D.SL14}; \cref{maps:eq:D.SL19}; \cref{maps:eq:D.SL17}; \cref{maps:eq:D.SL18}; \cref{maps:eq:D.SL16}; \cref{maps:eq:D.SL20}; \cref{maps:eq:D.SL22}; \cref{maps:eq:D.SL24}}


Let a marked zero-free star boundary have the integral basis
$(\alpha,\beta,\lambda,\ell)$ and monodromy
\[
\begin{gathered}
 A\alpha=-\alpha,\quad A\beta=-\beta-n\lambda,\quad
 A\lambda=\lambda-\alpha,\quad A\ell=\ell,\\
 \delta=-\alpha+2\lambda,\qquad\gamma^*(\ell)=1,
 \qquad n\in\{1,3\}.
\end{gathered}                                               \tag{D.SL1}\label{maps:eq:D.SL1}
\]
There are twelve actual marked star positions of this type
in the natural catalogue: the eight primitive-gamma $I_1^*$
positions of \ref{maps:src:D.2.27} and the four primitive-gamma $I_3^*$
positions. Their unimodular bases are retained in the
certificate below. The theorem also applies to any further
input for which the displayed integral marked basis is supplied.

Take a root of order $m=2k$, with $k>0$. Its rank-one
semistable Mumford cap has $q=nk$ components. Let its affine
deck action on their cyclic labels be
$i\mapsto j-i$, with a marked integer lift $j\in\mathbb Z$.
The component labels use $j$ modulo $q$; its integer lift is
retained in the peripheral marking. Write its collapsed
power as
\[
 v_{\rm col}=R\lambda+N\ell,\qquad R=2r-j,
                \qquad \gcd(N,2k)=1 .                     \tag{D.SL2}\label{maps:eq:D.SL2}
\]
Every $j,r$ and every such $N$ is realized by an actual free
affine root action. All its holomorphic unit factors and
integer affine twists are retained. Put
\[
 f=\#\{i:2i=j\pmod q\},\qquad \rho=(q+f)/2.
\]
Thus $f=1$ for odd $k$, $f=2$ for even $k$ and even $j$,
and $f=0$ for even $k$ and odd $j$.
For the quotient cap $U$, the complete integral groups are
\[
\begin{array}{c|c}
a&H_a(U;\mathbb Z)\\\hline
0&\mathbb Z\\
1&\mathbb Z^2\\
2&\mathbb Z^{\rho+1}\oplus
                    (\mathbb Z/2\text{ if }f=0)\\
3&\mathbb Z^{2\rho}\\
4&\mathbb Z^\rho,
\end{array}                                                \tag{D.SL3}\label{maps:eq:D.SL3}
\]
and higher groups vanish.

The boundary has $H_2(\partial U)=\mathbb Z^4$. Its
degree-two fibre coinvariant has the primitive basis
\[
 x=[\beta\lambda],\qquad
 w=[\beta\ell]+\frac{n-1}{2}[\lambda\ell],\qquad
 [\alpha\beta]=-2x,\quad [\beta\ell]=nw,\quad
 [\lambda\ell]=-2w.                                     \tag{D.SL4}\label{maps:eq:D.SL4}
\]
Let $D,L$ be its literal fixed-normal tori in the
$\delta,\ell$ directions. In coordinates $(x,w,D,L)$,
the \emph{full} boundary kernel is as follows:
\[
\begin{array}{c|c}
f&\ker[H_2(\partial U)\longrightarrow H_2(U)]\\\hline
1&\{(0,b,c,d):kb=2Nc-Rd\}\\
2&\{(0,b,c,d):(k/2)b=Nc-(R/2)d\}\\
0&\{(0,b,c,d):kb=2Nc-Rd,\ b\equiv0\pmod2\}.
\end{array}                                                \tag{D.SL5}\label{maps:eq:D.SL5}
\]
These are unsaturated marked kernels where appropriate.
Their projected indices in the full fixed lattice are
respectively $k,k/2,2k$. The boundary cokernel is free of
rank $\rho-1$. In particular the last row retains the
actual cap $\mathbb Z/2$ and its nonzero fibre correction;
it is not a conclusion from rational ranks or determinant
orders alone.

These columns can be inserted directly into the integral
bar algorithm of \ref{maps:src:D.2.41} or \ref{maps:src:D.2.46}. A vector $(0,b,c,d)$
contributes the actual regular-fibre representative of
$bw$, together with $c$ times the marked $\delta$ sweep
and $d$ times the marked $\ell$ sweep. The cap contributes
$\rho$ to the free relative rank in \eqref{maps:eq:D.GI10} and zero to
the Euler number. Thus all other already proved cap
choices and full finite Thom columns can be retained.

There is also a closed all-order result when this is the
sole logarithmic cap and every other cap has the minimal
natural marking of \ref{maps:src:D.2.26}, with no extra regular logs.
Write the initial coefficient as $b_0=4B$, and let $e$
be the ordinary clutching integer. Put
\[
                         s=N-2ke .                       \tag{D.SL6}\label{maps:eq:D.SL6}
\]
It is an odd nonzero integer, prime to $2k$. Then
\[
 \pi_1(X)=\mathbb Z/|s|,\qquad
 T=\operatorname{Tor}H_2(X)=
 \begin{cases}
 \mathbb Z/(2B|s|),&f=1,\\
 \mathbb Z/(B|s|),&f=2,\\
 \mathbb Z/2\oplus\mathbb Z/(2B|s|),&f=0.
 \end{cases}                                             \tag{D.SL7}\label{maps:eq:D.SL7}
\]
The other groups are
\[
\begin{gathered}
 H_2=\mathbb Z^{\rho-1}\oplus T,\qquad
 H_3=\mathbb Z^{2\rho-2}\oplus T,\\
 H_1=\mathbb Z/|s|,\qquad
 H_4=\mathbb Z^{\rho-1}\oplus\mathbb Z/|s|,\qquad H_5=0,
 \qquad H_0=H_6=\mathbb Z .
\end{gathered}                                             \tag{D.SL8}\label{maps:eq:D.SL8}
\]
Cyclic summands of order one are omitted. For example,
the six default natural rows with $b_0=4$ admit the
simply connected choice $s=\pm1$, $k=2h$, $j$ even.
It has
\[
          H_2=H_4=\mathbb Z^{nh},\qquad
                         H_3=\mathbb Z^{2nh}.              \tag{D.SL9}\label{maps:eq:D.SL9}
\]
Here $n=1$ on OS50/default and $n=3$ on OS46/default.
These are nonminimal star-log outputs, and no smooth
classification is inferred from their homology.

\end{theorem}


\begin{proof}
\begin{enumerate}[label=\arabic*.,leftmargin=*]
\item \emph{The actual unbounded affine actions.}
From \eqref{maps:eq:D.SL1}, $A^{2k}-I=q\alpha\beta^*$, so the
minimal semistable root model has exactly $q$ components.
This is the same integral Tate model used in \ref{maps:src:D.2.24}
and \ref{maps:src:D.2.27}, now at arbitrary root order. Its canonical
deck action reverses the component circle and fixes
the chosen zero component. Flat sections of the root
torus which are relevant to the affine action are
\[
                t=\frac{j}{q}\beta+\frac r k\lambda
                                      +\frac{N}{2k}\ell.
                                                               \tag{D.SL10}\label{maps:eq:D.SL10}
\]
They are single-valued because
$(A^{2k}-I)t=j\alpha$ is integral. The first term
translates the Mumford fan by $j$ components and
therefore extends over the toric cap. The other
terms are ordinary unit translations in the
transverse periods. Their full norm is
\[
         \sum_{a=0}^{2k-1}A^a t
                     =-r\alpha+(2r-j)\lambda+N\ell.
                                                               \tag{D.SL11}\label{maps:eq:D.SL11}
\]
It is integral, so composing the canonical deck
map with this translation still has order $2k$.
Using the inverse endpoint convention changes
only the coefficient of $\alpha$ in \eqref{maps:eq:D.SL11};
its collapsed vector \eqref{maps:eq:D.SL2} is unchanged, and
the corresponding exterior differences vanish
in the fibre coinvariant \eqref{maps:eq:D.SL4}.

\item The gamma displacement of a nonidentity deck
power is nonintegral because $N$ is prime to
$2k$. Hence the action is free on every stratum,
including fixed components and fixed double
curves. Conversely, in this canonical primitive-
gamma sector the flat translation conditions
give $q t_\beta\in\mathbb Z$, $k t_\lambda\in\mathbb Z$
and $2k t_\ell\in\mathbb Z$; a remaining angular
translation is removed by conjugation since
$A\alpha=-\alpha$. This is the marking \eqref{maps:eq:D.SL10}.
The usual root-and-unit isotopy keeps these
integral characters and all component meridians.
Thus the following topological calculation
applies to each actual action, not merely to
an abstract permutation of components.

\item \emph{The normalization and conductor at every order.}
The cap retracts equivariantly to its central
rank-one toric fibre. The quotient of its
$q$-gon dual graph by the reflection is an
interval. There are $\rho$ vertex orbits and
$\epsilon$ edge orbits, with
$\rho+\epsilon=q+1$; the number of fixed vertices
plus fixed edges is two. The normalization
and double-curve incidence graph is the
barycentric subdivision of this interval,
with exactly $q$ conductor incidences.

\item A paired vertex has stabilizer generated by
the square of the deck map. Its open
$\mathbb C^*$ bundle has an abelian period
lattice. The angular circle is primitive
because $\gcd(N,k)=1$. Thus its normalization
is a ruled surface over a torus with zero
Chern class, and its two pole sections have
the same class $\sigma$.
A fixed vertex instead has a pole-interchanging
clutch. Around the transverse $\lambda$ circle
the relation $A\lambda=\lambda-\alpha$ gives
the projective loop
$z\mapsto e^{-2\pi i u}/z$.
Its determinant winding is odd. Passing to
the primitive quotient gamma loop uses an
odd power of this map, since $N$ is odd, and
does not change the parity. Its normalization
is therefore the odd oriented $S^2$ bundle
over a torus. Choose a section $\sigma$ with
$\sigma^2=1$ and denote its complex sphere
fibre by $F$. Its connected conductor
bisection has class $2\sigma-F$; its
self-intersection is zero because the pole
normal line becomes trivial on the cover.

\item Every conductor flag has the even subgroup
as stabilizer. Its map to a normalization
base has degree two at a fixed vertex and
one at a paired vertex. Its map to the
double curve has degree two at a fixed edge
and one at a paired edge. Consequently
$H_2(U)$ has the following \emph{integral}
presentation. For each vertex orbit use
$F_v,\sigma_v$, and for each edge orbit use
its complex double-curve class $D_e$.
For each conductor incidence impose
\[
 \begin{cases}
 2\sigma_v-F_v-d_eD_e=0,&v\text{ fixed},\\
 \sigma_v-d_eD_e=0,&v\text{ paired},
 \end{cases}
 \qquad d_e=\begin{cases}2,&e\text{ fixed},\\1,&e\text{ paired}.
 \end{cases}                                             \tag{D.SL12}\label{maps:eq:D.SL12}
\]
There is no hidden degree-one bridge in this
normalization sequence: the conductor map
on $H_1$ is injective, by leaf elimination
on the incidence tree, since every local
isogeny is injective on $H_1$. The map on
$H_2$ is likewise injective. Hence
$H_3(U)$ is the direct sum of the
$\mathbb Z^2$ groups of the normalizations,
and $H_4(U)$ is their direct sum of top
classes.

\item All conductor section classes in \eqref{maps:eq:D.SL12}
are a single common class $S$. Internal
paired vertices have $\sigma_v=S$;
fixed vertices satisfy $2\sigma_v-F_v=S$;
a fixed endpoint edge satisfies $2D_e=S$.
Leaf elimination now gives free rank
$\rho+1$. If a vertex is fixed its sphere
coefficient is a unit and no torsion
survives. If no vertex is fixed, the two
fixed endpoint edges give
\[
 S=2D_0=2D_1,\qquad
       \tau=D_1-D_0\text{ of exact order two}.             \tag{D.SL13}\label{maps:eq:D.SL13}
\]
This proves the degree-two part of \eqref{maps:eq:D.SL3}.

\item For clarity, the fundamental group of the
cap follows directly from the free-cover
description after collapsing $\alpha$:
\[
\begin{gathered}
 \Gamma=\langle\beta,\lambda,\ell,g\mid
 [\beta,\lambda]=[\beta,\ell]=[\lambda,\ell]=1,\\
 g\beta g^{-1}=-\beta-n\lambda,\quad
 [g,\lambda]=[g,\ell]=1,\quad
                    2kg=R\lambda+N\ell\rangle .
\end{gathered}                                             \tag{D.SL14}\label{maps:eq:D.SL14}
\]
Additive notation in the last relation is
within the abelian subgroup generated by
$g,\lambda,\ell$. Its abelianization is free
of rank two: the relation $2\beta+n\lambda=0$
is primitive because $n$ is odd, and the
remaining relation is primitive because
$N$ is prime to $2k$. This proves $H_1(U)$
and finishes the local groups. As an
independent calculation, choose a primitive
quotient gamma circle. The cap is a mapping
torus of $I_q\times S^1_\lambda$; on its
component spheres the return is the
reflection, while on the actual specialized
torus $t=[\beta\lambda]$ naturality of the
exterior-square action gives
$t\mapsto-t-\sum_iF_i$. Its integral Wang
sequence gives the same groups, including
the order-two summand in \eqref{maps:eq:D.SL13}.

\item \emph{The regular-fibre map and its missing parity.}
The coinvariant relations \eqref{maps:eq:D.SL4} follow
by direct integral exterior-square reduction
of \eqref{maps:eq:D.SL1}; all four nonzero Smith factors
are units. The regular torus $\alpha\beta$
specializes to the sum of the $q$ complex
component spheres, each with degree one.
The transverse torus $\lambda\ell$ maps
to $kS$. Indeed its chosen component has
stabilizer of order $k$ in the paired case,
so its base cover has degree $k$ and its
class is $k\sigma$. In the fixed case its
base cover has degree $2k$; it is disjoint
from the bisection and has class
$k(2\sigma-F)$. Both give $kS$ in \eqref{maps:eq:D.SL12}.
Thus
\[
        2x=-\sum_v|v|F_v,\qquad 2w=-kS.                  \tag{D.SL15}\label{maps:eq:D.SL15}
\]
When $f>0$ the cap group is free, so these
equations determine the images uniquely.
Put
$T_0=\sum_{v\,\mathrm{paired}}F_v+
       \sum_{v\,\mathrm{fixed}}\sigma_v$.
For $f=1$ let $D_0$ be the sole fixed-edge
class. For $f=2$ retain the common $S$.
Then the images are
\[
\begin{array}{c|c|c}
 f&x&w\\\hline
1&-T_0+D_0&-kD_0\\
2&-T_0+S&-(k/2)S.
\end{array}                                               \tag{D.SL16}\label{maps:eq:D.SL16}
\]
For $f=0$, equations \eqref{maps:eq:D.SL15} alone leave
two genuine order-two ambiguities. We
resolve them using actual conductor loops.

\item Choose integers $a_0,b_0$ with
$a_0N+b_0(2k)=1$ and put
$\kappa=a_0g+b_0\ell$. The classifying
group \eqref{maps:eq:D.SL14} becomes the mapping-torus
group of the two-torus $(\beta,\lambda)$
with return
\[
           B=\begin{pmatrix}-1&0\\-n&1\end{pmatrix}.
\]
In fact
$\ell=2k\kappa-a_0R\lambda$ and
$g=N\kappa+b_0R\lambda$.
Here $a_0$ is odd. Its integral degree-two
group is
\[
 H_2(\Gamma)=\mathbb Zz\oplus\mathbb Z/2\,t,
       \qquad z=[\lambda\kappa],\quad t=[\beta\lambda].
                                                               \tag{D.SL17}\label{maps:eq:D.SL17}
\]
This group is computed by the same integral
bar comparison as \ref{maps:src:D.2.41}. For the two-dimensional
return $B$, its quadratic correction is
$f_B(u)=-n\binom{u_\beta}{2}t$.
Thus a $2k$-fold $\beta$ norm torus has
free part $-knz$ and torsion part $-knt$.
The actual regular tori therefore map as
\[
 [\beta\ell]\longmapsto-knz-(kn+a_0R)t,
 \qquad [\lambda\ell]\longmapsto2kz.                     \tag{D.SL18}\label{maps:eq:D.SL18}
\]

\item First choose $0\le j<q$ and label the two fixed endpoint edges so that
their stabilizer loops are $g$ and $\beta g$.
These are the actual two reflected edge
lifts, separated by half the dual circle.
Their primitive transverse periods are
$\kappa$ and
$\beta\kappa-\frac{n(a_0-1)}2\lambda$, respectively.
Their complex torus classes consequently
map to
\[
                        D_0\mapsto z,\qquad D_1\mapsto z-t.
                                                               \tag{D.SL19}\label{maps:eq:D.SL19}
\]
Both use the positive quotient gamma period,
so their orientations agree. It follows
that $\tau$ maps to the nonzero $t$.
After quotienting $H_2(U)$ by its component
spheres, \eqref{maps:eq:D.SL19} is an isomorphism onto
\eqref{maps:eq:D.SL17}. In particular it detects precisely
the torsion ambiguity left in \eqref{maps:eq:D.SL15}.

\item Since $k$ is even and $j,R,a_0$ are odd
in this case, \eqref{maps:eq:D.SL18} sends the primitive
$w$ of \eqref{maps:eq:D.SL4} to $-kz+t$; it sends $x$
to $t$. Hence the exact, marked images are
\[
       x=-\sum_vF_v+\tau,\qquad
                         w=-kD_0+\tau.                   \tag{D.SL20}\label{maps:eq:D.SL20}
\]
Both order-two terms are forced by the
actual classifying map. In particular the
fibre cokernel has cyclic torsion
$\mathbb Z/(2k)$ in this case; replacing
the second image by $-kD_0$ would give a
different marked map.

\item \emph{The literal fixed-normal tori.}
The root meridian is $g^{2k}v^{-1}$,
with \eqref{maps:eq:D.SL11}. Sweeping its bounding
normal disc by $\delta$ and $\ell$ gives
the exact relations in $H_2(U)$
\[
                 2kD+4Nw=0,\qquad
                             2kL-2Rw=0.                  \tag{D.SL21}\label{maps:eq:D.SL21}
\]
The omitted $\alpha$ terms vanish by
\eqref{maps:eq:D.SL4}. These relations determine the
free parts of both images. When $f>0$
there is no torsion ambiguity. When $f=0$,
the classifying map fixes it as well:
the cap loops are $\delta=2\lambda$,
$g=N\kappa+b_0R\lambda$ and
$\ell=2k\kappa-a_0R\lambda$, so their
actual commuting tori map to
\[
                 D\mapsto2Nz,\qquad L\mapsto-Rz
                            \quad\text{in }H_2(\Gamma).
\]
There is no $t$ term. In the same notation
as \eqref{maps:eq:D.SL16}--\eqref{maps:eq:D.SL20}, the full images are
\[
\begin{array}{c|c|c}
f&D&L\\\hline
1&2ND_0&-RD_0\\
2&NS&-(R/2)S\\
0&2ND_0&-RD_0.
\end{array}                                               \tag{D.SL22}\label{maps:eq:D.SL22}
\]
Together with the fibre images these
prove \eqref{maps:eq:D.SL5} by solving an integral
linear relation. For $f=0$ the coefficient
of $x$ must vanish in the independent
component-sphere direction; the coefficient
of $w$ must then be even by \eqref{maps:eq:D.SL20}.
The free coefficient gives the displayed
equation. This explicitly retains the
relative torsion information rather than
saturating the normal lattice.

\item The cokernel of the full boundary map is
free of rank $\rho-1$, as one sees directly
from \eqref{maps:eq:D.SL12} and \eqref{maps:eq:D.SL16}--\eqref{maps:eq:D.SL22}.
Equivalently, Poincare--Lefschetz duality
puts it into the free group $H^4(U)$.
The projected-kernel indices are $k,k/2,2k$;
the computations for $m=2$ and the swapped
$m=4$ in \ref{maps:src:D.2.27} are recovered exactly.
The full local kernel can therefore be
used in the global bar presentation,
with all its existing endpoint corrections.

\item The integer lift of $j$ must also be retained.
Write $j=j_0+hq$, $0\le j_0<q$. The corresponding
peripheral generator changes from $g_0$ to
$g=\beta^h g_0$, and $R=R_0-hq$.
Its literal normal tori satisfy
\[
                    D=D_0^{\rm norm},\qquad
                    L=L_0-hn w.                           \tag{D.SL24}\label{maps:eq:D.SL24}
\]
Here the first symbol denotes the old normal
torus, not a double-curve class. Indeed
$[\delta\beta]=0$ and $[\beta\ell]=nw$
in \eqref{maps:eq:D.SL4}. Substituting \eqref{maps:eq:D.SL24} into
the kernel formulas replaces $R_0$ by
$R$ exactly. In the $f=0$ case it changes
the displayed image of $L$ in \eqref{maps:eq:D.SL22} to
$-RD_0+h\tau$, but the equation in
\eqref{maps:eq:D.SL5} already forces its coefficient
$d$ to be even; the kernel parity is
therefore still precisely $b$ even.
Thus integer affine lifts are included
without suppressing their actual torsion
correction.

\item \emph{The closed one-star Smith calculation.}
For the natural minimal exterior, the
fibre functional has
$\Omega(x)=-a$ and $\Omega(w)=b_0/4=B$.
The primitive graph calculation has the
single diagonal relation $D=-b_0e\,s_0$,
where $s_0$ is its regular-fibre generator.
Because every vector in \eqref{maps:eq:D.SL5} has zero
$x$ coordinate, its global relation on
$(s_0,L)$ is
\[
                         (B(b-4ec),d).                   \tag{D.SL23}\label{maps:eq:D.SL23}
\]
Set $u=b-4ec$ and $s=N-2ke$.
For $f=1$ its defining condition becomes
$ku+Rd=2sc$. Since $k$ is odd and prime
to $2s$, $d$ is unrestricted and
$u$ is a fixed residue modulo $2s$.
The Smith cokernel is cyclic of order
$2B|s|$.
For $f=2$, the condition is
$(k/2)u+(R/2)d=sc$.
Here $k/2$ is prime to $s$, so the same
argument gives a cyclic cokernel of
order $B|s|$.
For $f=0$, $u$ and $d$ are even; write
$u=2u_0,d=2d_0$. The condition becomes
$ku_0+Rd_0=sc$. Its relation matrix has
Smith factors $2$ and $2B|s|$.
This proves the three torsion groups
in \eqref{maps:eq:D.SL7} for every parameter.

\item The one-log Seifert presentation gives
$\pi_1=\mathbb Z/|s|$. Its rational
Euler number cannot vanish because
$N$ is prime to $2k$. The free local
boundary defect is $\rho-1$ and there
is no additional rational degree-one
bridge. All other caps in the stated
corollary have their minimal ranks.
Thus $b_2=\rho-1$. The unique primitive
rank-two zero still contributes Euler
number two, while this star cap has
Euler number zero. Integral duality
gives \eqref{maps:eq:D.SL8}, and the simply connected
specialization gives \eqref{maps:eq:D.SL9}.
\end{enumerate}
\end{proof}

The exact script \path{code/III/manifolds_star_log_caps.py}
retains the twelve actual canonical bases,
the full normalization matrices, the
literal fibre and normal maps, and their
complete integral kernels. It checks
the conductor calculation independently
against the mapping-torus sequence and
compares the global Smith formula on
unbounded symbolic lattices and a finite
set of exact instances. The proof above,
including the order-two marking, is
independent of the finite parameter
bounds used in those corroborating runs.
Root independently read the full proof and replayed all
$12$ actual marked bases, $95$ normalization comparisons
and $1638$ global Smith calculations successfully.
An independent adversarial review by the fourfolds
agent also checked \eqref{maps:eq:D.SL17}--\eqref{maps:eq:D.SL24}, including the
nonzero order-two correction and the integer phase lift.
Root actions with nonprimitive gamma
power, other integral star blocks and
positive-weight star degenerations are
outside this theorem's stated scope.


\begin{theorem}[all vertical corrections and new simply connected fillings]\label{maps:src:D.2.50}
\StatementRecord{maps:src:D.2.50}{proved}{Historical events: Claude Fable 5.2 (Anthropic), effort high, 16 November 2026: carried as background without fault found. GPT-6 Astra (OpenAI), ChatGPT Pro mode, effort Pro, S6\_09\_PAPERS, 17 November 2026: targeted proof check. Also checked in this writing. Pre-continuation writing events: GPT-6 Astra (OpenAI), ChatGPT Pro mode, effort Pro, S6\_10\_PAPERS, 18 November 2026: targeted proof check. Continuation metadata events: GPT-6 Astra (OpenAI), ChatGPT Pro mode, effort Pro, S6\_10\_PAPERS continuation, 19 November 2026: used and rechecked in this run; status-vocabulary reconciliation; metadata reconciliation only; no new mathematical check. \emph{Pre-continuation recorded proof scope.} GPT-6 Astra (OpenAI, ChatGPT Pro mode), effort Pro, 2026-11-17 UTC; GPT-6 Astra (OpenAI, ChatGPT Pro mode), S6\_06\_BLIND, effort Pro, 2026-11-17; Claude Fable 5.2 (Anthropic), effort high, 16 November 2026; also checked in this writing by GPT-6 Astra (OpenAI), ChatGPT Pro mode, effort Pro, 18 November 2026}{the hypotheses, notation and definitions stated here; \cref{maps:eq:D.GI2}; \cref{maps:eq:D.VC3}; \cref{maps:src:D.2.19}; \cref{maps:src:D.2.46}; \cref{maps:eq:D.VC8}; \cref{maps:eq:D.VC6}; \cref{maps:eq:D.VC2}; \cref{maps:src:D.2.43}; \cref{maps:eq:D.GI3}; \cref{maps:eq:D.GI5}; \cref{maps:src:D.2.41}; \cref{maps:eq:D.GI6}}


Let $f:S\to\mathbb P^1$ be one of the twenty-one recorded
rational elliptic surfaces with Mordell--Weil group $\mathbb ZG$,
zero section $O$, and height $\langle G,G\rangle=1/d_0$.
For integers $a,b$, put $P=aG$, $Q=bG$ and
\[
 M_D=\mathcal O_S(P-O+D),\qquad
 J_D=M_D\otimes T_Q^*M_D^{-1},
\]
where $D$ is an integral divisor supported on fibre components.
Here $T_Q:S\to S$ is the actual section-translation automorphism.
Some such $D$ makes
$J_D$ the pullback of a base line bundle if and only if
$d_0\mid ab$. For every such correction the exact identity is
\[
                  J_D\simeq f^*\mathcal O_{\mathbb P^1}(ab/d_0).
                                                        \tag{D.VC1}\label{maps:eq:D.VC1}
\]
Thus the compatible initial line bundles are not exhausted
by requiring $P$ or $Q$ narrow at every fibre.
All corrections are determined by the integral local equations
below, with precisely the invariant component orbit sums as ambiguity.
For $a,b>0$ and any nonnegative whole-fibre weights of total $ab/d_0$,
the corrected punctured family has exactly the marked monodromies
\eqref{maps:eq:D.GI2}. This does not identify its newly permitted original caps.

There are additional explicitly marked simply connected examples.
On OS45/default, choose $(a,b)=(2,4)$ or $(4,2)$, a simple
linearization zero at any of its four $I_1$ fibres, multiplicative
Mumford fillings at every singular fibre, no extra affine boundary
translations, and ordinary gamma clutch $e=\pm1$. Use the actual
vertical corrections specified in the proof. All sixteen choices have
\[
 \begin{gathered}
 \pi_1(X)=1,\qquad H_0(X)=H_6(X)=\mathbb Z,\\
 H_2(X)=H_3(X)=\mathbb Z/b,\qquad
 H_1(X)=H_4(X)=H_5(X)=0 .
 \end{gathered}                                             \tag{D.VC7}\label{maps:eq:D.VC7}
\]
These are new initial and marked filling data beyond the earlier
narrowness condition. Their smooth diffeomorphism classes are not
inferred from these groups.

\end{theorem}


\begin{proof}
\begin{enumerate}[label=\arabic*.,leftmargin=*]
\item \emph{The integral local equation.}
Index the components of a singular fibre by $C_0,\ldots,C_{v-1}$,
with $C_0$ the component meeting $O$. Let $m$ be its positive
primitive multiplicity vector, $J=(C_iC_j)$ its intersection
matrix and $\sigma$ the permutation matrix induced by $Q$.
If $P$ meets component $C_p$, the actual multidegree of
$M_D$ on this fibre is
\[
                    c_D=e_p-e_0+Jz,
                                                        \tag{D.VC2}\label{maps:eq:D.VC2}
\]
where $z$ is the vector of coefficients of $D$ there.
All section components have multiplicity one, so $m^tc_D=0$.
The vanishing of the component degrees of $J_D$ is exactly
\[
                 (\sigma-I)Jz=-(\sigma-I)(e_p-e_0).
                                                        \tag{D.VC3}\label{maps:eq:D.VC3}
\]
Using the inverse permutation for pullback gives the same
invariance condition and the same solution set.
The equation is integral: it may not be replaced by its
rational solution or by the saturation of its image.

\item For completeness, its entire ambiguity also has an exact
description. If $z,z'$ solve \eqref{maps:eq:D.VC3}, then
$J(\sigma-I)(z-z')=0$, since $J$ commutes with the permutation.
The integral kernel of $J$ is $\mathbb Zm$, so
$(\sigma-I)(z-z')=h m$ for some integer $h$.
Apply $I+\sigma+\cdots+\sigma^{q-1}$, where $\sigma^q=I$.
Because $\sigma m=m$, the result is $0=qh m$, hence $h=0$.
Conversely every invariant vector can be added to a solution.
Thus all corrections are one particular solution plus
\[
                         \ker(\sigma-I)\subset\mathbb Z^v.
                                                        \tag{D.VC4}\label{maps:eq:D.VC4}
\]
In particular the ambiguity consists of integral orbit sums
of components, including the whole-fibre ambiguity.

\item \emph{Why vanishing component degrees is sufficient globally.}
On the generic elliptic fibre, translation acts trivially
on $\operatorname{Pic}^0$, so $J_D$ is generically trivial.
It therefore has a rational section whose divisor is vertical:
one extends a generic trivialization as a rational section
on the smooth surface. If \eqref{maps:eq:D.VC3} holds at every fibre,
this vertical divisor has intersection zero with every
component. On each connected fibre the negative-semidefinite
intersection matrix has kernel exactly its multiplicity
vector. The divisor on that fibre is an integral multiple
of the whole fibre, since the multiplicities have gcd one.
Consequently $J_D=f^*\mathcal O(t)$ for an integer $t$.
This proves sufficiency for the actual line bundle, not
only for its rational first Chern class. Necessity follows
by restricting a pullback to the components.

\item We determine $t$ without suppressing a local correction.
Choose $h>0$ with $hQ$ narrow at every singular fibre.
Then $T_{hQ}$ fixes each component divisor. Telescoping
the $h$ translated copies of $J_D$ gives
\[
 \bigotimes_{j=0}^{h-1} T_Q^{j*}J_D
       =M_D\otimes T_{hQ}^*M_D^{-1}
       =\mathcal O(P-O)\otimes
                    T_{hQ}^*\mathcal O(P-O)^{-1}.          \tag{D.VC5}\label{maps:eq:D.VC5}
\]
The left side is $f^*\mathcal O(ht)$, since $T_Q$ preserves
the base. The last expression is in the already proved
narrow-translation case of \ref{maps:src:D.2.19} and \ref{maps:src:D.2.46}. Its degree
is the integral height pairing
$\langle P,hQ\rangle=hab/d_0$.
Pullback on base Picard groups is injective, by restriction
to the zero section. Hence $ht=hab/d_0$, proving \eqref{maps:eq:D.VC1}
and the necessary divisibility $d_0\mid ab$.

\item \emph{The actual punctured monodromies.}
Let the chosen base linearization have whole-fibre weights $w_i$.
Choose $h$ as above and iterate the actual contracting generator
$h$ times. The canonical nonzero section of $\mathcal O(D)$
off the singular fibres telescopes in this iteration, because
$T_{hQ}D=D$ as an actual divisor. Thus the iterated punctured
family is the uncorrected narrow-translation construction with
sections $P,hQ$ and weights $h w_i$. A constant nonzero scalar
in this identification changes no monodromy.

\item In the actual period lattice its long generator is
$\gamma_h=h\gamma$; the elliptic generators and $\delta$ are
unchanged. If $A_i$ denotes the original monodromy and $A_{i,h}$
the iterated one, put $S=\operatorname{diag}(1,1,1,h)$. Then
\[
                         A_iS=SA_{i,h}.                    \tag{D.VC8}\label{maps:eq:D.VC8}
\]
The narrow case \eqref{maps:eq:D.GI2} has $q_{i,h}=h q_i$,
$z_{i,h}=h z_i$, and central entry $r_i(hz_i)+h w_i$.
Solving \eqref{maps:eq:D.VC8} gives exactly the unscaled matrix \eqref{maps:eq:D.GI2}.
This proves its integrality in every corrected case as well.
The argument concerns the punctured family and its actual long
period. It does not assume a covering of a filled original or
logarithmic cap, nor suppress any subsequently chosen affine
boundary marking.

\item \emph{Solvability for every recorded fibre.}
For an $I_n$ fibre the component lattice of total degree
zero modulo the image of $J$ is $\mathbb Z/n$; its quotient
map is $c\mapsto\sum_i i c_i\pmod n$.
This follows by the integral successive-difference reduction
of the cycle Laplacian, whose nonzero Smith factors are
$1,\ldots,1,n$. For $n=1$ the lattice is zero; for $n=2$
the two distinct edges are both retained.
Write $q$ for the component of $Q$, $g=\gcd(n,q)$ and
$t=n/g$. A $q$-translation-invariant multidegree has one
coefficient $u_j$ on each of its $g$ orbits. Its total degree
is zero exactly when $\sum_j u_j=0$, and its component
class is $t\sum_j j u_j\pmod n$.
Thus the subgroup represented by invariant multidegrees
is exactly $t\mathbb Z/n$. The class of $e_p-e_0$ lies in it
if and only if $t\mid p$, equivalently $n\mid pq$.
This proves integral solvability of \eqref{maps:eq:D.VC3} for every
multiplicative fibre.

\item The only other types in the twenty-one profiles are
$II,III,IV,III^*,IV^*,I_1^*,I_3^*$.
The respective component-group orders are $1,2,3,2,3,4,4$.
For these finite lists, the exact marked intersection
matrices and component permutations are supplied in
\path{code/III/root_vertical_corrections.py}.
Their multiplicity kernels and permutation invariance are
checked there before any equation is solved. The integral
Smith calculation for every $p,q$ in each component group
gives the same criterion
\[
              \operatorname{ord}(\Phi_f)\mid pq.           \tag{D.VC6}\label{maps:eq:D.VC6}
\]
For $III$ and $IV$ this is also the same cycle-matrix
calculation just given. The remaining star cases involve
only $4$, $9$, $16$ and $16$ pairs respectively; their
full solutions and all integer kernel generators are
retained in the output. This is an exhaustive finite
proof for those types, not a sampling of section multiples.

\item For every recorded profile the component of $G$ is a unit
in each nontrivial component group, and the least common
multiple of their orders is $d_0$. These finite marked
data are checked in the same executable against the
twenty-one source records. Hence \eqref{maps:eq:D.VC6} for $P=aG,Q=bG$
at all fibres is equivalent to $d_0\mid ab$.
Choose a Smith solution independently at each fibre and
sum the corresponding component divisors. This constructs
the required global $D$ and proves sufficiency.
All integers $a,b$ are covered because their component
classes depend only on their residues modulo $d_0$.
The supplemental check enumerates all $5696$ such residue
pairs, using $227$ distinct local component pairs, and
finds no discrepancy with \eqref{maps:eq:D.VC1}.

\item \emph{A concrete correction and the remaining topology.}
At an $I_4$ fibre with $P$ and $Q$ both meeting component
two, the original multidegree is $(-1,0,1,0)$.
Taking $D=-C_0$ changes it to $(1,-1,1,-1)$, which is
invariant under translation by two. Neither section is
narrow. The full solution set is obtained by adding
integer multiples of $C_0+C_2$ and $C_1+C_3$, as the
actual Smith kernel confirms.
Globally the profiles with such additional residue pairs
are precisely OS45, OS46, OS50 and OS55, with the listed
additive variants retained. The new initial data cannot
be discarded merely because they violate the earlier OR
condition. Their original caps carry the actual corrected
multidegree \eqref{maps:eq:D.VC2}, which can be nonprimitive and can
be permuted nontrivially by $Q$. Computing those marked
cap maps is a separate remaining part of D.2.

\item \emph{The sixteen new all-Mumford examples.}
The actual OS45/default data are
\[
 \begin{gathered}
 (I_8,I_1,I_1,I_1,I_1),\qquad d_0=8,\qquad
 \overline G_{I_8}=3\pmod8,\\
 g_i=(3,0),(0,0),(2,-1),(1,-2),(0,0),\\
 (B_0,\ldots,B_4)=\begin{pmatrix}1&8\\0&1\end{pmatrix},
 \begin{pmatrix}1&1\\0&1\end{pmatrix},
 \begin{pmatrix}3&4\\-1&-1\end{pmatrix},
 \begin{pmatrix}3&1\\-4&-1\end{pmatrix},
 \begin{pmatrix}0&1\\-1&2\end{pmatrix}.
 \end{gathered}                                             \tag{D.VC9}\label{maps:eq:D.VC9}
\]
At $I_8$ one may take the following literal coefficient vectors:
\[
\begin{array}{c|c|c}
 (a,b)&(D_{C_0},\ldots,D_{C_7})&c_D\\ \hline
 (2,4)&(-1,-1,-1,0,0,0,0,0)&(0,0,1,-1,0,0,1,-1)\\
 (4,2)&(-1,0,0,1,1,1,0,0)&(1,-1,1,-1,1,-1,1,-1).
\end{array}                                                \tag{D.VC10}\label{maps:eq:D.VC10}
\]
They are invariant respectively under translation by four and by six;
no correction is needed at the four irreducible fibres. Both height
pairings equal one. Use the punctured marking \eqref{maps:eq:D.VC8}, one weight one
at a selected $I_1$, weight zero elsewhere, and exactly the affine
boundary data in the statement.

\item Each positive-weight cap is the primitive index-one triangular cap
of \ref{maps:src:D.2.43}. Every other cap is the rank-one Mumford cap of its actual
matrix. Thus all local filling kernels in \eqref{maps:eq:D.GI3} and \eqref{maps:eq:D.GI5} apply
without a narrowness assumption: their proofs concern these toric
monodromies and their marked primitive angular discs. The common
piece and every ordinary boundary map are the same integral bar
models as in \ref{maps:src:D.2.41}. No newly corrected original cap is used.

\item For completeness the entire finite calculation is retained in
\path{code/III/fourfolds_corrected_Mumford_examples.py} and its JSON output.
The input range is exactly the two pairs, four zero positions,
and two clutch signs specified above. In all sixteen cases the
unsaturated normal closure of the toric vanishing vectors is the
whole vertical lattice $\mathbb Z\langle e_1,e_2,\delta\rangle$.
All singular peripheral loops are killed by their Mumford fillings;
the remaining ordinary relation kills $\gamma$ with coefficient
$\pm1$. This proves full simple connectivity, not only $H_1=0$.

\item The degree-two bar presentation has $23$ generators and $70$
relations. Its exact Smith form has $22$ unit factors and last
factor $b$. The additional relative free rank in \eqref{maps:eq:D.GI6} is zero:
there are five component classes and the finite-filled exterior
has first Betti number five. The only Euler contribution is the
index-one positive cap, so $\chi=2$. Integral duality and universal
coefficients now give every group in \eqref{maps:eq:D.VC7}.
These actual matrices, vertical corrections, and boundary markings
are retained for review. Invariant orbit sums added to $D$ are not
silently identified with these marked filled spaces; their induced
boundary rigidifications must also be transported.
\end{enumerate}
\end{proof}

The independent certificate
\path{code/fourfolds_vertical_correction_audit.py}
checks all $227$ local pairs with unimodular Smith transformations,
all $5696$ global residues, and the $82$ actual elliptic
monodromy/section-cocycle orders against the recorded component
maps. It also records all sixteen additional positive height-one
section pairs across the full profile list. Those latter sixteen
pairs are initial data, distinct from the sixteen fully filled
OS45 examples in \eqref{maps:eq:D.VC7}.


\begin{theorem}[invariant-Chern original additive caps and their integral extension]\label{maps:src:D.2.51}
\StatementRecord{maps:src:D.2.51}{proved}{Historical events: Claude Fable 5.2 (Anthropic), effort high, 16 November 2026: carried as background without fault found. GPT-6 Astra (OpenAI), ChatGPT Pro mode, effort Pro, S6\_09\_PAPERS, 17 November 2026: targeted proof check. Also checked in this writing. Pre-continuation writing events: GPT-6 Astra (OpenAI), ChatGPT Pro mode, effort Pro, S6\_10\_PAPERS, 18 November 2026: targeted proof check. Continuation metadata events: GPT-6 Astra (OpenAI), ChatGPT Pro mode, effort Pro, S6\_10\_PAPERS continuation, 19 November 2026: used and rechecked in this run; status-vocabulary reconciliation; metadata reconciliation only; no new mathematical check. \emph{Pre-continuation recorded proof scope.} GPT-6 Astra (OpenAI, ChatGPT Pro mode), effort Pro, 2026-11-17 UTC; Claude Fable 5.2 (Anthropic), effort high, 16 November 2026; also checked in this writing by GPT-6 Astra (OpenAI), ChatGPT Pro mode, effort Pro, 18 November 2026}{the hypotheses, notation and definitions stated here; \cref{maps:src:D.2.50}; \cref{maps:src:D.2.31}; \cref{maps:src:D.2.16}}


Let $K$ be the reduced central fibre of an original additive
Kodaira neighbourhood, of type
$II,III,IV,IV^*,III^*,I_1^*$ or $I_3^*$.
Write its rational components as $C_i$, its primitive
multiplicity vector as $\mu=(\mu_i)$, and let $\sigma$
be the actual component permutation induced by the section $Q$.
Let the original line bundle have invariant component degrees
$c_i$, so that
\[
                  \sigma c=c,\qquad \sum_i\mu_i c_i=0.
                                                        \tag{D.IC1}\label{maps:eq:D.IC1}
\]
These include every corrected line bundle in \ref{maps:src:D.2.50}.
Let the component orbits be $O_1,\ldots,O_\rho$, with
lengths $l_j$ and common degrees $c_j$ on $O_j$.
For $c\ne0$ put
\[
                  g=\gcd_j(c_j),\qquad
                  h=\gcd_j(l_jc_j),
                                                        \tag{D.IC2}\label{maps:eq:D.IC2}
\]
where both gcds are positive. For the resulting original
six-dimensional cap $U$, there is an exact integral cellular
cone, with its actual circle-fibre and regular-area classes,
\[
\begin{gathered}
 C_0=\mathbb Z,\quad C_1=\mathbb Z^2,\quad
 C_2=\mathbb Z^v\oplus\mathbb Zt,\quad
 C_3=\mathbb Z^v\oplus\mathbb Z^v,\quad C_4=\mathbb Z^v,\\
 d_2=\begin{pmatrix}c^t&0\\0&0\end{pmatrix},\qquad
 d_3=\begin{pmatrix}0&\sigma-I\\0&c^t\end{pmatrix},\qquad
 d_4=\begin{pmatrix}I-\sigma\\0\end{pmatrix}.
\end{gathered}                                             \tag{D.IC3}\label{maps:eq:D.IC3}
\]
The degree-one generators are the scalar circle $\delta$
and a positive long circle $\gamma$. In particular
\[
\begin{array}{c|ccccc}
i&0&1&2&3&4\\\hline
H_i(U;\mathbb Z)&\mathbb Z&
\mathbb Z\oplus\mathbb Z/g&
\mathbb Z^{\rho-1}\oplus\mathbb Z/h&
\mathbb Z^{2\rho-1}&\mathbb Z^\rho .
\end{array}                                               \tag{D.IC4}\label{maps:eq:D.IC4}
\]
Moreover $\pi_1(U)=\mathbb Z\oplus\mathbb Z/g$.
When $c=0$, these groups are instead
\[
 (H_0,H_1,H_2,H_3,H_4)=
 (\mathbb Z,\mathbb Z^2,\mathbb Z^{\rho+1},
                         \mathbb Z^{2\rho},\mathbb Z^\rho).
                                                        \tag{D.IC5}\label{maps:eq:D.IC5}
\]
The actual regular-fibre torus $\delta\wedge\gamma$
maps to the class of $t$ in \eqref{maps:eq:D.IC3}, which generates
\emph{all} of the cyclic torsion $\mathbb Z/h$ when
$c\ne0$. The elliptic tori crossed with $\delta$ map
to zero. The complex elliptic area maps to the literal
cycle $(\mu,0)$ in \eqref{maps:eq:D.IC3}.

Let $B$ be the elliptic monodromy, $q$ its marked
$Q$-cocycle, and $z=(B-I)^{-1}q$. Put
\[
        q_Q=\gcd_j(l_j\mu_j)
             =\operatorname{ord}(\overline Q).
                                                        \tag{D.IC6}\label{maps:eq:D.IC6}
\]
The actual logarithm $z$ has exact denominator $q_Q$.
For $c\ne0$, the full regular-fibre image in $H_2(U)$
is $\mathbb Z\oplus\mathbb Z/h$ with free cokernel
of rank $\rho-2$. For $c=0$ it is a primitive
$\mathbb Z^2$ with free cokernel of rank $\rho-1$.
In both cases this is already the full boundary image.
Consequently the full boundary kernel projects onto
all of $\operatorname{Fix}(A)$ in the Wang sequence;
there is no omitted finite normal index.

If $h=1$, these facts determine the complete marked
local kernels. In the regular basis
$(e_1,e_2,\delta,\gamma)$ the actual fibre kernel is
\[
 \ker\bigl(q_Q(1,0,z_2,0,-z_1,0):\mathbb Z^6\to\mathbb Z\bigr),
                                                        \tag{D.IC7}\label{maps:eq:D.IC7}
\]
where the six coordinates have their usual increasing
exterior order. Both literal fixed-normal tori bound.
The cap kills every vertical fundamental-group
coordinate and retains the primitive $\gamma$.
This includes the canonical new corrected star degrees
\[
 \begin{array}{c|c}
 I_1^*&(1,1,0,0,-1,0)\\
 I_3^*&(1,1,0,0,-1,0,0,0)
 \end{array}
                                                        \tag{D.IC8}\label{maps:eq:D.IC8}
\]
for component classes $P=Q=2$ in $\mathbb Z/4$.
Their actual permutation has two paired outer orbits
and fixed inner components, and $g=h=1$.

For $h>1$, \eqref{maps:eq:D.IC3}--\eqref{maps:eq:D.IC6} do not by themselves specify
the torsion corrections of the two mixed
elliptic--gamma tori and the second fixed-normal torus.
The theorem does not replace those marked columns by
their saturation. In particular it does not supply an
unrestricted global gluing algorithm for those caps.

\end{theorem}


\begin{proof}
\begin{enumerate}[label=\arabic*.,leftmargin=*]
\item \emph{The actual original cap and the circle bundle.}
The reduced additive fibre is simply connected and
is homotopy equivalent to a wedge of its $v$ rational
component spheres. The tangency in type $III$ and the
triple point in type $IV$ do not add a graph loop;
the same assertion follows directly by taking the
union of the rational components and their actual
connected intersections. The other listed fibres
have a tree of rational components.
The elliptic neighbourhood retracts onto this fibre.
The unit original linearization makes its compact
cap the mapping torus of the principal scalar-circle
bundle $P_c\to K$, with the lift of $Q$.
This is the same actual original filling used in
\ref{maps:src:D.2.31}--\ref{maps:src:D.2.32}, now retaining the full degree vector.

\item A principal circle bundle over this wedge is determined
by its $v$ Chern evaluations. Its cellular complex has
one zero-cell, the scalar one-cell, $v$ two-cells with
boundary degrees $c_i$, and $v$ three-cells with zero
boundary. The lift of $Q$ acts by the component
permutation on the two- and three-cells. There is no
additional homotopy class of a gauge transformation:
$[K,S^1]=H^1(K;\mathbb Z)=0$. More explicitly, after
choosing a fixed point of the induced component-tree
map and a fibre point there, the component maps are
homotopic to their orientation-preserving permutations.
The invariant vector $c$ lifts this homotopy to the
circle bundle. A comparison unit has a continuous
logarithm on $K$, so its lift is homotopic through
bundle maps to the identity. A constant scalar at
the chosen point is removed by a scalar isotopy.
This fixes the scalar circle and its product with
the mapping-torus direction throughout the comparison.
The signs in \eqref{maps:eq:D.IC3} use the literal product orientation
$t=\delta\wedge\gamma$: the boundary of a two-cell
crossed with the long interval contributes $+c_i t$.
This is the usual mapping cone after changing the signs
of its shifted summands. Thus \eqref{maps:eq:D.IC3} is an actual model, not
merely a rational model with the same permutation.

\item The circle-bundle homotopy sequence gives
$\pi_1(P_c)=\mathbb Z/g$ for $c\ne0$.
The return fixes the scalar generator; the mapping-torus
extension by the infinite cyclic group therefore splits.
This proves the fundamental-group statement and $H_1$.

\item \emph{The integral extension.}
In $C_2/\operatorname{im}d_3$, successive cells in an
orbit have the relations
\[
                  e_{i+1}-e_i+c_jt=0.
                                                        \tag{D.IC9}\label{maps:eq:D.IC9}
\]
Eliminating them leaves one free $e_j$ for each orbit,
and the remaining relations are $l_jc_jt=0$.
Hence this quotient is
$\mathbb Z^\rho\oplus\mathbb Z/h$.
Its remaining boundary sends $e_j$ to $c_j\delta$
and sends $t$ to zero. Its kernel is therefore
$\mathbb Z^{\rho-1}\oplus\mathbb Z/h$.
This proves the full nonsplit extension in $H_2$.
For example, two interchanged components of degree
$2$ contribute a relation $4t=0$, not two unrelated
order-two summands. The fibre product torus
$\delta\wedge\gamma$ is the cone generator $t$ itself.

\item The degree-three cone has the direct sum of the
coinvariants of the component three-cells, of rank
$\rho$, and the invariant component two-cycles with
Chern evaluation zero, of rank $\rho-1$.
Both groups are free. Degree four is the invariant
component three-cell group of rank $\rho$.
The same calculation with $c=0$ gives \eqref{maps:eq:D.IC5}.

\item There is also an intrinsic check of the extension.
Let $Y$ be the mapping torus of $K$ and $Q$.
The cap is the principal $S^1_\delta$ bundle over $Y$.
The Euler class evaluates as $c_j$ on the orbit
representatives in $H_2(Y)$, and as $l_jc_j$ on
invariant orbit sums in $H_3(Y)$ paired with the
positive $\gamma$. The integral Gysin sequence
therefore injects $\mathbb Z/h$ into $H_2(U)$ by
$\gamma\mapsto\delta\wedge\gamma$.
The quotient is the free kernel of the first Chern
functional on $H_2(Y)$. This verifies both the
extension and the actual generator without a
choice of an abstract splitting.

\item \emph{The regular area and the full boundary image.}
In the single circle bundle $P_c\to K$ the map
$H_2(P_c)\to H_2(K)$ is injective, with image
$\ker(c^t)$. The regular elliptic area projects to
$\mu$, the multiplicity cycle of the central fibre,
and hence has its unique lift $(\mu,0)$.
The elliptic loops crossed with the scalar circle
vanish because $H_1(K)=0$.
After passing to $Y$, this area projects to the
orbit vector $(l_j\mu_j)$. Its content is $q_Q$.
The equality with the order of the component of
$Q$ follows directly from the marked component
permutations and multiplicities for the seven listed
types; the retained exact certificate checks every
component class. The actual section-cocycle model
has the same order in $\operatorname{coker}(B-I)$,
so the logarithm $z$ has exactly this denominator.

\item To identify the rational mixed area coefficients,
pass to the $q_Q$-fold long-circle cover.
The section $q_QQ$ is narrow; its marked local
logarithm straightens the translation through the
original elliptic neighbourhood, as in \ref{maps:src:D.2.16} and
\ref{maps:src:D.2.31}. Its corrected long direction is
$q_Q(\gamma-z)$. An elliptic loop crossed with this
direction is zero in the base mapping torus because
the elliptic loop bounds in the Kodaira neighbourhood.
Thus the two mixed free coefficients are $z_2$ and
$-z_1$ times the elliptic-area coefficient.
The integers $q_Q,q_Qz_1,q_Qz_2$ have gcd one.
Consequently these actual fibre tori generate the
primitive free orbit vector
$(l_j\mu_j)/q_Q$.
The torus $\delta\wedge\gamma$ already generates
all cap torsion. This proves the claimed fibre
image and its free cokernel.

\item Finally $H_2(U,\partial U)=H^4(U)=\mathbb Z^\rho$,
since $H_3(U)$ is free. The boundary first Betti
number is three, while that of $U$ is one for
$c\ne0$ and two for $c=0$. The boundary maps onto
$H_1(U)$: it contains the scalar and long generators.
The relative sequence therefore makes the full
boundary degree-two image have rational rank one
or two, respectively. Those are already the ranks
of the primitive fibre images just proved. Their
cokernels are free, so no further finite enlargement
is possible. The images are equal integrally.
The Wang projection of the full kernel is therefore
all of the fixed lattice.

\item For completeness, the scalar fixed-normal torus
bounds the scalar-circle bundle over the zero-section
disc. The other fixed-normal torus has zero projection
to $H_2(Y)$: after the same narrow cover and marked
straightening, it is the boundary of the zero-section
disc crossed with the long circle. Since $H_2(Y)$
is free, it was zero before the cover as well.
Its image in $H_2(U)$ is consequently torsion.
If $h=1$ both normal images vanish, and the full
fibre map is the primitive integral area functional
in \eqref{maps:eq:D.IC7}. This proves the complete marked assertion
in that subsector.
\end{enumerate}
\end{proof}

The exact implementation is
\path{code/III/manifolds_invariant_chern_caps.py}.
It constructs the full cone, retains the actual
$\delta\gamma$ and multiplicity cycles, checks their
orders and cokernels, and verifies the component-order
and logarithm-denominator markings against every
recorded additive site. Its finite checks corroborate
the all-integer elimination in \eqref{maps:eq:D.IC9}; they are not
bounds on the permitted Chern coefficients.


\begin{theorem}[all height-integral initial pairs with corrected original additive caps]\label{maps:src:D.2.52}
\StatementRecord{maps:src:D.2.52}{proved}{Historical events: Claude Fable 5.2 (Anthropic), effort high, 16 November 2026: carried as background without fault found. GPT-6 Astra (OpenAI), ChatGPT Pro mode, effort Pro, S6\_09\_PAPERS, 17 November 2026: targeted proof check. Also checked in this writing. Pre-continuation writing events: GPT-6 Astra (OpenAI), ChatGPT Pro mode, effort Pro, S6\_10\_PAPERS, 18 November 2026: targeted proof check. Continuation metadata events: GPT-6 Astra (OpenAI), ChatGPT Pro mode, effort Pro, S6\_10\_PAPERS continuation, 19 November 2026: used and rechecked in this run; status-vocabulary reconciliation; metadata reconciliation only; no new mathematical check. \emph{Pre-continuation recorded proof scope.} GPT-6 Astra (OpenAI, ChatGPT Pro mode), effort Pro, 2026-11-17 UTC; Claude Fable 5.2 (Anthropic), effort high, 16 November 2026; also checked in this writing by GPT-6 Astra (OpenAI), ChatGPT Pro mode, effort Pro, 18 November 2026}{the hypotheses, notation and definitions stated here; \cref{maps:src:D.2.50}; \cref{maps:eq:D.IC2}; \cref{maps:src:D.2.46}; \cref{maps:eq:D.IC8}; \cref{maps:eq:D.VC3}; \cref{maps:eq:D.IC7}; \cref{maps:eq:D.GI2}; \cref{maps:src:D.2.41}; \cref{maps:src:D.2.43}; \cref{maps:src:D.2.51}; \cref{maps:src:D.2.31}; \cref{maps:src:D.2.33}; \cref{maps:eq:D.GI6}; \cref{maps:eq:D.CI1}}


For any of the twenty-one recorded Mordell--Weil rank-one
profiles, take $P=aG,Q=bG$ with $a,b>0$ and
$d_0\mid ab$. Choose any nonnegative whole-fibre
linearization weights of total $ab/d_0$, supported
at multiplicative or additional regular fibres.
Use the actual corrected line bundle of \ref{maps:src:D.2.50}.
The following filling choices have an exact all-order
fundamental-group presentation and complete integral
homology algorithm:
\begin{enumerate}[label=(\roman*)]
\item every marked same-fan rank-one or triangular
Mumford cap at a multiplicative or regular site;
\item every free finite-good additive logarithmic cap,
with all its affine powers and all five integral Thom
columns retained;
\item every original additive cap whose corrected
multidegree has $h=1$ in \eqref{maps:eq:D.IC2}, and every unchanged
zero-degree original additive cap;
\item every original multiplicative cap with its
uncorrected degree vector, as in \ref{maps:src:D.2.46};
\item arbitrarily many regular logarithmic caps and
any ordinary integral clutch.
\end{enumerate}
The odd-star fibres are filled originally in this
algorithm. In particular, the recorded Smith
corrections from \ref{maps:src:D.2.50} admit every positive pair
with integral height: all newly non-narrow additive
positions are the component pair $(2,2)$ at an odd
star, and their corrected degrees are \eqref{maps:eq:D.IC8}.
Newly non-narrow multiplicative positions may all
use the marked Mumford filling.
Thus this theorem removes the narrowness-OR restriction
on the initial sections in \ref{maps:src:D.2.46}.
It does not identify original multiplicative caps
with their nonzero vertical corrections, nor the
unresolved marked $h>1$ additive caps.

The exact implementation is
\path{code/III/manifolds_corrected_initial.py}.
It accepts the same unbounded data as \ref{maps:src:D.2.46} and
optional literal component-divisor vectors. A supplied
vector is checked against the integral equation
\eqref{maps:eq:D.VC3}; no rational correction is substituted.
For an original additive site it computes the
orbit lengths and $h$, and applies the complete
marked kernel \eqref{maps:eq:D.IC7} only when $h=1$.
A nonzero whole-fibre change at a zero-degree
original cap is also retained as outside this
implementation until its relative meridian marking
is supplied.

\end{theorem}


\begin{proof}
\begin{enumerate}[label=\arabic*.,leftmargin=*]
\item \ref{maps:src:D.2.50} constructs the actual initial line bundle
and proves that its punctured monodromy is exactly
\eqref{maps:eq:D.GI2}, without the narrowness-OR assumption.
All entries are integral and their ordered product
is the identity. The common punctured torus bundle
therefore has exactly the same integral fibre-and-graph
bar model as in \ref{maps:src:D.2.41} and \ref{maps:src:D.2.46}, including its
quadratic comparison and literal last-boundary word.

\item Every selected Mumford cap depends on the actual
marked unipotent period matrix and the selected fan.
The rank-one calculation and \ref{maps:src:D.2.43} therefore apply
to these newly corrected inputs verbatim: the full
nonsaturated angular-area kernel, primitive normal
discs and all component meridians remain as stated.
No extension of the vertical divisor over a different
original cap is being inferred from this fact.
The finite-good and regular logarithmic caps use
the actual affine quotient action, so their freeness
criterion, norm columns and fifth Thom generator
are unchanged as well.

\item At an original additive site with $h=1$, \ref{maps:src:D.2.51}
supplies every required integral marking. Its
fundamental group is generated by the primitive
long circle; all three vertical directions are
killed. The fibre degree-two kernel is the actual
primitive area kernel \eqref{maps:eq:D.IC7}, and both literal
fixed-normal tori bound. Its component contribution
to $H_4$ is the number $\rho$ of component orbits,
not the number of components before the permutation.
The unchanged zero-degree case and the uncorrected
original multiplicative case are already proved
in \ref{maps:src:D.2.31}, \ref{maps:src:D.2.33} and \ref{maps:src:D.2.46}.

\item The implementation therefore adds these exact
local columns to the common bar presentation.
Its fundamental-group presentation retains the
commuting fibre generators, the conjugation matrices,
the actual ordered peripheral product, each finite
power relation, and each nonfinite killed vector.
The normal closure of the vertical vectors is
computed by repeatedly adjoining their conjugates
and taking integral Hermite normal form until the
lattice stabilizes. The ascending-chain condition
on submodules of $\mathbb Z^3$ proves termination.
No saturation or assumed abelian quotient replaces
this presentation.

\item For degree-two homology, form the full integer
relation matrix $M$ and retain its Smith cokernel.
If $E$ denotes the common piece after the finite
caps have been filled, the remaining relative
free rank is exactly
\[
           \sum_i b_4(U_i)-b_1(E)+b_1(X).
                                                        \tag{D.CI1}\label{maps:eq:D.CI1}
\]
The proof is the same integral pair sequence as
\eqref{maps:eq:D.GI6}: each $H_2(U_i,\partial U_i)=H^4(U_i)$
is free, and its actual boundary columns are
already included. The cap torsion permitted by
\ref{maps:src:D.2.43} remains in its nonsaturated fibre kernel;
the new original additive caps have no degree-two
torsion because $h=1$. The quotient in \eqref{maps:eq:D.CI1}
is free, so it splits as an abelian-group extension.
This gives $H_2(X)$ with all torsion.
The Euler number is twice the sum of the indices
of the positive rank-two caps. The exact $H_1$,
this Euler number, integral duality and universal
coefficients give every remaining group, including
the Euler-zero cases.

\item Finally, the canonical correction list in \ref{maps:src:D.2.50}
has $45$ additional residue pairs over the full
profile list. At additive sites the component
orders are prime except for the two odd stars.
The criterion $\operatorname{ord}(\Phi)\mid pq$
therefore forces $p=0$ or $q=0$ except for
$p=q=2$ in $\mathbb Z/4$.
The displayed canonical correction at those
stars has a fixed inner coefficient $-1$,
so $h=1$. This proves that the algorithm applies
to every positive height-integral initial pair
with these retained corrections and the stated
filling choices. The conclusion is unbounded in
section multiples, weights, logarithmic orders
and every permitted affine twist.
\end{enumerate}
\end{proof}

The independent replay script
\path{code/III/manifolds_corrected_initial_audit.py}
compares the full matrices against the earlier
\ref{maps:src:D.2.46} implementation and the sixteen independently
computed \ref{maps:src:D.2.50} examples. It also runs every one
of the $45$ new residue pairs, with all listed
multiplicative zero positions, both available
additive filling patterns, and additional
higher-weight and regular-log inputs.
Those finite executions check the implementation;
the actual local models and integral pair argument
above prove the stated all-parameter scope.


\begin{theorem}[complete primitive cubic and fundamental-class invariant of the weighted OS45 example]\label{maps:src:D.2.53}
\StatementRecord{maps:src:D.2.53}{proved}{Historical events: Claude Fable 5.2 (Anthropic), effort high, 16 November 2026: carried as background without fault found. GPT-6 Astra (OpenAI), ChatGPT Pro mode, effort Pro, S6\_09\_PAPERS, 17 November 2026: targeted proof check. Also checked in this writing. Pre-continuation writing events: GPT-6 Astra (OpenAI), ChatGPT Pro mode, effort Pro, S6\_10\_PAPERS, 18 November 2026: targeted proof check. Continuation metadata events: GPT-6 Astra (OpenAI), ChatGPT Pro mode, effort Pro, S6\_10\_PAPERS continuation, 19 November 2026: used and rechecked in this run; status-vocabulary reconciliation; metadata reconciliation only; no new mathematical check. \emph{Pre-continuation recorded proof scope.} GPT-6 Astra (OpenAI, ChatGPT Pro mode), effort Pro, 2026-11-17 UTC; Claude Fable 5.2 (Anthropic), effort high, 16 November 2026; also checked in this writing by GPT-6 Astra (OpenAI), ChatGPT Pro mode, effort Pro, 18 November 2026}{the hypotheses, notation and definitions stated here; \cref{maps:eq:D.PC11}; \cref{maps:eq:D.PC12}; \cref{maps:eq:D.OF1}; \cref{maps:eq:D.OF14}; \cref{maps:eq:D.OF17}; \cref{maps:eq:D.OF18}}


Take the marked OS45/default rational elliptic surface, the initial
sections \(P=2G,Q=8G\), whole-fibre weights
\((2,0,0,0,0)\) supported at \(I_8\), all multiplicative Mumford
caps, and the ordinary clutch \(e=1\). The resulting simply
connected manifold \(X\) has
\[
 H_2(X;\mathbb Z)=\mathbb Z^{15}\oplus\mathbb Z/8,\qquad
 H_3(X;\mathbb Z)=\mathbb Z/8,\qquad
 c_1(X)=c_2(X)=\lambda(X)=0.                              \tag{D.OS1}\label{maps:eq:D.OS1}
\]
Its sixteen normalized cusp components are indexed by
\((r,s)\in\mathbb Z/8\times\mathbb Z/2\), with \(i=2r+s\).
Their divisor classes generate an index-two sublattice in
\(H^2(X;\mathbb Z)\). The full primitive lattice is generated
by \(x_0,\ldots,x_{13},y\), with
\[
             2y=\sum_{r=0}^7x_{r,1},\qquad
             \sum_{r,s}x_{r,s}=0.                        \tag{D.OS2}\label{maps:eq:D.OS2}
\]
The complete integral cubic is \eqref{maps:eq:D.PC11}--\eqref{maps:eq:D.PC12} below.
There is a primitive cyclic coordinate \(u\in H^2(X;\mathbb Z/8)\)
with the full boundary marking \eqref{maps:eq:D.OF1}--\eqref{maps:eq:D.OF14}. In the splitting
fixed by \(u\) and this primitive basis,
\[
 H_6(K(H_2(X),2);\mathbb Z)
   =\mathbb Z^{680}\oplus(\mathbb Z/8)^{120}
                       \oplus(\mathbb Z/16)^{15}
                       \oplus\mathbb Z/8,                \tag{D.OS3}\label{maps:eq:D.OS3}
\]
and the complete fundamental-class invariant is specified by
the integral cubic and \eqref{maps:eq:D.OF17}--\eqref{maps:eq:D.OF18}.
In particular
\[
 \begin{gathered}
 y^2=0,\qquad
 \langle\mathcal P_2(u)x_i,[X]\rangle=7\pmod{16}
                   \ (0\le i<14),\\
 \langle\mathcal P_2(u)y,[X]\rangle=0,\qquad
 \langle u^3,[X]\rangle=0.
 \end{gathered}                                           \tag{D.OS4}\label{maps:eq:D.OS4}
\]
This theorem does not assert a completed smooth classification:
the refined characteristic class is still one of \(0,4\)
in the cyclic order-eight summand of \(H^4(X;\mathbb Z)\).
The full finite-coefficient invariant here does not choose
between those two possibilities.

\end{theorem}
\begin{proof}
\begin{enumerate}[label=\arabic*.,leftmargin=*]
\item \textbf{The primitive component lattice and cubic in the sixteen-component OS45 example}
We first record a saturation calculation which uses the actual component
meridians. Let $X$ be simply connected, let $N$ be a single rank-two
Mumford cusp with $D$ components, and put $E=X\setminus\operatorname{int}N$.
Assume $b_2(X)=D-1$, $H_3(N;\mathbb Z)$ is free, and
\[
 H_1(E;\mathbb Z)=\mathbb Z\oplus A,\qquad
                       \mu_i=(1,a_i),                      \tag{D.PC1}\label{maps:eq:D.PC1}
\]
where $A$ is finite and $\mu_i$ is the positively oriented meridian of
component $i$. The component classes $x_i$ then generate a sublattice
\[
 L=\mathbb Z^D/\mathbb Z(1,\ldots,1)
                       \subset H^2(X;\mathbb Z)
\]
whose full primitive extension, inside $L\otimes\mathbb Q$, is
\[
 H^2(X;\mathbb Z)
   =L+\left\{\sum_i\widetilde{\chi(a_i)}x_i:
                         \chi\in\operatorname{Hom}(A,\mathbb Q/\mathbb Z)
        \right\}.                                        \tag{D.PC2}\label{maps:eq:D.PC2}
\]
A tilde means any rational lift of a value in $\mathbb Q/\mathbb Z$;
changing it adds an element of $L$. In particular the component
lattice has index $|A|$.

\item Poincare--Lefschetz duality identifies
$H^2(X,E;\mathbb Z)=H_4(N;\mathbb Z)=\mathbb Z^D$,
with its positive component Thom basis, and
$H^3(X,E;\mathbb Z)=H_3(N;\mathbb Z)$ is free. The connecting
map from $H^1(E;\mathbb Z)$ evaluates on the meridians \eqref{maps:eq:D.PC1},
so its image is precisely $\mathbb Z(1,\ldots,1)$.
The relative exact sequence therefore injects $L$ into $H^2(X)$.
Since $X$ is simply connected, the latter group is free; its rank is
$D-1$ by hypothesis. Its quotient by $L$ is finite and embeds in
$H^2(E)$, whose torsion is
$\operatorname{Ext}(H_1(E),\mathbb Z)=\operatorname{Hom}(A,\mathbb Q/\mathbb Z)$.
All this torsion lifts to $H^2(X)$ because the next relative group is
free. Thus the quotient has exactly that order.

\item Here is also its precise embedding, rather than only its index.
For a character $\chi$ of order $m$, extend it by zero on the free
summand of $H_1(E)$. Regard it as a class in
$H^1(E;\mathbb Z/m)$ by multiplying its values by $m$.
Its relative coboundary has component coordinates
$m\widetilde{\chi(a_i)}\pmod m$. Choose the displayed integers
as an integral relative lift. Its image in $H^2(X;\mathbb Z/m)$
is zero by exactness. Hence its integral image is divisible by $m$
in $H^2(X;\mathbb Z)$, giving the class in \eqref{maps:eq:D.PC2}.
The coefficient exact sequences identify its exterior restriction
with the corresponding torsion Bockstein, up to the harmless common
connecting-map sign. Since these characters give the full torsion
of $H^2(E)$, they give the full primitive extension. This proves
\eqref{maps:eq:D.PC2}, including every finite denominator.

\item Now take the exact marked data
\[
 \mathrm{OS45/default},\quad P=2G,\quad Q=8G,\quad
 (w_0,w_1,w_2,w_3,w_4)=(2,0,0,0,0),\quad e=1,              \tag{D.PC3}\label{maps:eq:D.PC3}
\]
with all multiplicative fibres filled by their marked Mumford caps.
\ref{maps:src:D.2.46} gives
\[
 H_2(X)=\mathbb Z^{15}\oplus\mathbb Z/8,\qquad
 H_3(X)=\mathbb Z/8,\qquad \pi_1(X)=1.
                                                               \tag{D.PC4}\label{maps:eq:D.PC4}
\]
The elliptic monodromy at $I_8$ is
$\left(\begin{smallmatrix}1&8\\0&1\end{smallmatrix}\right)$
and its section cocycle is $(3,0)$. Formula \eqref{maps:eq:D.GI2} therefore gives
\[
 A_0=\begin{pmatrix}
 1&8&0&24\\0&1&0&0\\0&-6&1&-16\\0&0&0&1
 \end{pmatrix},\qquad
 L_{\rm per}=\left\langle(8,-6),(24,-16)\right\rangle
                 =8\mathbb Z\oplus2\mathbb Z.             \tag{D.PC5}\label{maps:eq:D.PC5}
\]
The short basis is $(e_1,\delta)$ and the long basis is $(e_2,\gamma)$.
This is an equality of sublattices, not merely a comparison of
indices. The determinant is sixteen and the Smith factors are two
and eight. Label the normalized components by
\[
             (r,s)\in\mathbb Z/8\times\mathbb Z/2,
             \qquad i=2r+s.                              \tag{D.PC6}\label{maps:eq:D.PC6}
\]

\item The actual cusp complement satisfies
\[
 H_1(E)=\mathbb Z\oplus\mathbb Z/2,
                      \qquad \mu_{r,s}=(1,s\bmod2).      \tag{D.PC7}\label{maps:eq:D.PC7}
\]
To verify the marking, retain all peripheral generators and the
ordinary clutch in the exact degree-one presentation. Its Smith
form has seven unit factors, one factor two, and one free generator.
The torsion character in the original nine coordinates is
\[
                  (0,-2,1,2,0,0,2,-2,2)\pmod2,
                                                               \tag{D.PC8}\label{maps:eq:D.PC8}
\]
which is precisely $\delta^*$ modulo two. The free character
evaluates to one on the cusp meridian. A height-one ray at
$(r,s)$ has literal meridian $\theta+r e_1+s\delta$; it
therefore gives exactly \eqref{maps:eq:D.PC7}. The full original presentation,
Smith change of basis and all sixteen meridian images are saved
in the executable below.

\item Applying \eqref{maps:eq:D.PC2}, there is a unique integral class $y$ with
\[
                    2y=\sum_{r=0}^{7}x_{r,1}.             \tag{D.PC9}\label{maps:eq:D.PC9}
\]
The component lattice has exact index two. In particular it would
be incorrect to declare any fifteen of the component classes a
primitive basis. A convenient actual basis is
\[
 x_0,x_1,\ldots,x_{13},y,\qquad
 x_{14}=-2y-\sum_{\substack{i<14\\i\ \mathrm{even}}}x_i,
 \qquad
 x_{15}=2y-\sum_{\substack{i<14\\i\ \mathrm{odd}}}x_i.
                                                               \tag{D.PC10}\label{maps:eq:D.PC10}
\]
The integral matrix of the sixteen component classes in this basis
has Smith factors $1^{14},2$ and kernel generated by the all-one
vector. These are the asserted primitive markings.

\item Here is the complete cubic on that basis. For component classes,
\ref{maps:src:D.2.47} specializes to the following symmetric integral trilinear
form. With distinct vertices as appropriate,
\[
 \begin{aligned}
 T(x_i,x_i,x_i)&=6,\\
 T(x_i,x_i,x_j)&=
 \begin{cases}
 -2,&i=(r,s),\ j=(r,1-s),\\
 -1,&i=(r,s),\ j=(r\pm1,t),\quad t=0,1,\\
 0,&\text{otherwise},
 \end{cases}\\
 T(x_i,x_j,x_k)&=
 \begin{cases}
 1,&\{i,j,k\}\text{ is three of the four vertices in}\\[-2pt]
   &\hspace{12pt}\{(r,0),(r,1),(r+1,0),(r+1,1)\}
                   \text{ for some }r,\\
 0,&\text{otherwise}.
 \end{cases}
 \end{aligned}                                             \tag{D.PC11}\label{maps:eq:D.PC11}
\]
All horizontal indices are cyclic modulo eight. Repeated geometric
edges and triangle cells are retained in this formula.
The remaining primitive entries are
\[
 \begin{aligned}
 T(y,x_{r,s},x_{r',s'})
 &=(-1)^s\left(
       \mathbf1_{s'=s,\ r'=r+1}
      +\mathbf1_{s'=s,\ r'=r-1}
      -2\mathbf1_{(r',s')=(r,s)}\right),\\
 T(y,y,x_i)&=0,\qquad T(y,y,y)=0.
 \end{aligned}                                             \tag{D.PC12}\label{maps:eq:D.PC12}
\]
These formulas determine every entry on the primitive basis
\eqref{maps:eq:D.PC10}, and every resulting entry is integral.

For a geometric verification of the half-class entries, on a
normalized $dP_6$ take the cyclic boundary classes
\[
 E_1,\ H-E_1-E_2,\ E_2,\ H-E_2-E_3,\ E_3,\ H-E_1-E_3
                                                               \tag{D.PC13}\label{maps:eq:D.PC13}
\]
in the ray order of \ref{maps:src:D.2.47}. The four nonhorizontal boundary classes
sum to $2f$, where $f=H-E_1$ and $f^2=0$. Restricting \eqref{maps:eq:D.PC9}
therefore gives
\[
                         y|_{P_{r,s}}=(-1)^s f.            \tag{D.PC14}\label{maps:eq:D.PC14}
\]
This division is valid in the torsion-free group $H^2(dP_6;\mathbb Z)$.
The fibre $f$ has degree one on each horizontal boundary curve
and zero on the other four. The Cartier restriction formula in
\ref{maps:src:D.2.47} now gives the first line of \eqref{maps:eq:D.PC12}. Squaring \eqref{maps:eq:D.PC14}
gives $y^2x_i=0$; multiplying \eqref{maps:eq:D.PC9} by $y^2$ gives $y^3=0$.
Thus no rational cubic has been silently substituted for the
integral form.

Finally $c_1(X)=0$ and both $c_2$ and $p_1$ pair to zero with
every class in the primitive free basis. For $c_1$, the normalized
component tangent calculation is zero and the component restriction
is rationally injective by the connected toric graph Laplacian;
$H^2(X)$ is torsion-free, so this proves integral vanishing.
For $c_2,p_1$, \ref{maps:src:D.2.47} gives zero on every component. The component
classes have finite index, and the evaluations are integers, so they
also vanish on $y$. These statements concern the free characteristic
pairings; they do not determine the torsion class of $c_2$, the
refined class, or the full finite-coefficient fundamental-class
invariant.

\path{code/III/fourfolds_OS45_primitive_cubic.py}
reconstructs the actual complement before using its torsion
character, checks all sixteen meridians, computes the exact
index-two primitive lattice, and verifies all $15^3=3375$ primitive
triples and all $16^3=4096$ recovered component triples. The
normalization restriction \eqref{maps:eq:D.PC14}, the actual period basis, and
all sparse integral entries are retained in its JSON output.
The preliminary component-only calculation is retained separately
in \path{code/III/fourfolds_OS45_component_ring.py}.

\paragraph{Integral vanishing of the critical class.}
The primitive-lattice argument above gives the literal relation
\(\sum x_{r,s}=0\). Therefore the target line of \(df\) is
topologically trivial. All fibres in this example are reduced
semistable fibres. The exact complex of \ref{maps:src:D.2.44} consequently
applies without any additive correction: its abstract kernel
has the nowhere-zero scalar section, and its quotient line is
topologically trivial because \(c_1(X)=0\). The local unit Bott
clutch gives
\[
                     c_2(X)=\operatorname{PD}(C_{\rm crit}).
                                                               \tag{D.OS5}\label{maps:eq:D.OS5}
\]
Here every actual complex double curve occurs once.

We retain the actual integral specialization of this complete cycle.
In the ordered bivector basis
\((e_1e_2,e_1\delta,e_1\gamma,e_2\delta,e_2\gamma,\delta\gamma)\),
the five fibre preimages, in their actual site order, are
\[
 \begin{array}{c|rrrrrr}
 I_8\text{ weighted cusp}&-16&0&6&8&0&2\\
 I_1&0&0&0&0&0&1\\
 I_1&0&8&0&0&0&1\\
 I_1&0&0&0&8&0&1\\
 I_1&0&0&0&0&0&1
 \end{array}.                                             \tag{D.OS6}\label{maps:eq:D.OS6}
\]
For the first row, the actual local period cover has matrix
\(L=\left(\begin{smallmatrix}8&24\\-6&-16\end{smallmatrix}\right)\)
and degree sixteen. Transfer the minimal cusp's literal
three-curve preimage
\(a_1b_1+b_1a_2+a_2b_2\), expressing \(b_1,b_2\) by \(L^{-1}\)
in the actual long basis, and multiply by sixteen. This gives
the first integral row of \eqref{maps:eq:D.OS6}, without division in homology
or an assumed covering of the global filled manifolds.
Each other row is the determinant-one transverse torus of the
corresponding primitive node, as in \ref{maps:src:D.2.44}.
Their sum is
\[
                        z=(-16,8,6,16,0,6).              \tag{D.OS7}\label{maps:eq:D.OS7}
\]
In the actual twenty-three-row relation matrix \(M\) of
\ref{maps:src:D.2.46}, augment this vector by seventeen zeros. With columns
numbered from zero, a literal integer preimage is the column
combination with nonzero coefficients
\[
 \begin{array}{c|rrrrrrrr}
 j&6&8&10&15&16&18&19&20\\\hline
 c_j&-6&2&2&-28&2&-2&-52&-8\\
 \end{array}
 \qquad
 \begin{array}{c|rrrrrrr}
 j&26&38&41&51&56&61&66\\\hline
 c_j&2&10&10&10&10&10&10 .
 \end{array}                                             \tag{D.OS8}\label{maps:eq:D.OS8}
\]
Direct multiplication gives \(Mc=(z,0)\).
All columns of \(M\) are the actual core prism and cap Thom
relations, so this proves that the literal critical cycle is
integrally nullhomologous. Equations \eqref{maps:eq:D.OS5}--\eqref{maps:eq:D.OS8} prove
the absolute Chern and spin-characteristic vanishing in
\eqref{maps:eq:D.OS1}. They do not by themselves compute linking with an
odd middle-dimensional generator: local specialization chains
must be retained for that separate refined calculation.

\textbf{The actual torsion coordinate and complete fundamental-class
invariant of the sixteen-component OS45 manifold}
We use exactly the datum \eqref{maps:eq:D.PC3}, its period marking \eqref{maps:eq:D.PC5}, and the
primitive free basis \eqref{maps:eq:D.PC10}. Put \(k=\mathbb Z/8\).
There is a primitive torsion coordinate
\(u\in H^2(X;k)\) whose complete cusp support is
\[
 \Psi(r,s)=4\left\lfloor r/8\right\rfloor
       +3(r\bmod8)\left\lfloor s/2\right\rfloor
                                      \quad\hbox{in }k.       \tag{D.OF1}\label{maps:eq:D.OF1}
\]
Here and below \(r\bmod8\) means the representative between zero and
seven. The assertion includes the entire boundary cohomology class,
including its coefficient-extension part. We prove this before using
any support-function evaluations.

\paragraph{The specified core cochain.}
Use the actual degree-two core model of \ref{maps:src:D.2.46}. Its fibre cells are
ordered
\[
 e_1e_2,\ e_1\delta,\ e_1\gamma,\ e_2\delta,\
                                      e_2\gamma,\ \delta\gamma.
\]
For the six boundary monodromies, including the ordinary identity cap,
let
\[
 D=(A_0-I,\ldots,A_4-I),\qquad K=\ker_{\mathbb Z}D.
\]
Choose the integral cycle basis and the integral left inverse \(K^+\)
retained by the exact program. The ordinary degree-two cycles are the
six fibre cells followed by this graph-cycle basis. The actual
relation matrix \(M\) has twenty-three rows. If
\(UMV\) is its stored Smith form, the unique nonunit factor is eight.
Let \(\varphi\) be the negative of the corresponding row of \(U\).
Then
\[
 \varphi M=0\pmod8,\qquad
 \varphi|_{\Lambda^2\Lambda}=(1,0,0,0,0,0).                \tag{D.OF2}\label{maps:eq:D.OF2}
\]
Its value on the stored Smith torsion generator is \(-1\).
This is already a primitive functional on the order-eight image of
the core in \(H_2(X)\).

We specify the cochain, rather than merely its homology functional.
The inclusion \(K\subset\mathbb Z^{20}\) is primitive because it is
the kernel of an integral map. Extend the graph row by
\(\varphi_{\rm graph}K^+\), and use the fibre row in \eqref{maps:eq:D.OF2}.
This gives a literal degree-two cochain on the core mapping-cylinder
model. The degree-three boundary columns are the actual exterior
monodromy relations, so \eqref{maps:eq:D.OF2} says that it is closed modulo eight.
The fixed bar-prism comparisons in \ref{maps:src:D.2.46} identify this cochain with
one on the actual torus bundle. In particular its restriction is not
left unspecified by the possible \(2\)-torsion in \(H_1\) of the core.

At the rank-two cusp the actual peripheral inclusion has graph
matrix \(J=K^+J_{\rm raw}\). The retained cochain gives
\[
 \varphi_{\rm graph}J=(0,52,0,148)
       \quad\hbox{on }(e_1,e_2,\delta,\gamma).             \tag{D.OF3}\label{maps:eq:D.OF3}
\]
The quadratic correction on each of these four basis loops is zero
at this cusp. Thus the two long-strip constants are
\[
                      (c_1,c_2)=(4,4)\quad\hbox{in }k.   \tag{D.OF4}\label{maps:eq:D.OF4}
\]
The full Smith row, the chosen left inverse, and the actual
peripheral matrices are reproducible by
\path{code/III/maps_weighted_os45.py}. In particular \eqref{maps:eq:D.OF3} is not
obtained by restricting only the regular-fibre alternating form.

\paragraph{A full boundary comparison for finite-coefficient supports.}
We give the comparison explicitly, since this cusp has torsion in
its boundary \(H_1\). Let \(\widetilde N\) be the common universal
cover of the triangular Mumford cusp and let
\(\widetilde M\) be its boundary. Its normalized components are
indexed by \(v\in\mathbb Z^2\). The boundary cover retracts to the
three-torus generated by the two short angular loops and the positive
\(q\)-meridian \(\theta\). \ref{maps:src:D.2.43} gives
\[
 H_1(\widetilde N;\mathbb Z)=0,\quad
 H_2(\widetilde N;\mathbb Z)=I^2,\quad
 \operatorname{im}\{H_2(\widetilde M;\mathbb Z)\to
                         H_2(\widetilde N;\mathbb Z)\}=0. \tag{D.OF5}\label{maps:eq:D.OF5}
\]
Indeed the three boundary generators are the
angular area and the two angular-normal tori, all with the actual
bounding sweeps described there.

The relative component Thom cochains identify
\[
 H^2(\widetilde N,\widetilde M;k)
                  =\operatorname{Map}(\mathbb Z^2,k).
\]
One can obtain this directly from locally finite
component chains: a degree-four locally finite cycle is an arbitrary
family of fundamental cycles of the normalized components, and
degree-four intersection strata do not occur. Equivalently,
Poincare--Lefschetz duality uses Borel--Moore degree-four homology.
No compact-support convention replaces the displayed product of
coefficient groups by a direct sum.

The positive meridian of component \(v\) is
\(\theta+v_1e_1+v_2\delta\). Therefore the relative connecting map
from \(H^1(\widetilde M;k)\) sends a covector \((c,m_1,m_2)\) to
the affine function \(c+m_1v_1+m_2v_2\). By \eqref{maps:eq:D.OF5}, the relative
sequence is exactly
\[
 0\longrightarrow\operatorname{Aff}(\mathbb Z^2,k)
 \longrightarrow\operatorname{Map}(\mathbb Z^2,k)
 \longrightarrow H^2(\widetilde N;k)\longrightarrow0.     \tag{D.OF6}\label{maps:eq:D.OF6}
\]
Here \(\operatorname{Aff}\) consists of these three-parameter
affine functions. This also identifies the quotient with
\(\operatorname{Hom}(I^2,k)\), consistently with the actual
normalization computation of \ref{maps:src:D.2.43}.

Let the actual period lattice \(L\) act by translation on the
argument \(v\). If a function \(\Psi\) satisfies
\[
                  \Psi(v+l)-\Psi(v)=a_l(v)
                       \in\operatorname{Aff},           \tag{D.OF7}\label{maps:eq:D.OF7}
\]
its absolute class on \(\widetilde N\) is \(L\)-invariant.
It descends to \(N_L\): in total degree two the covering
cochain double complex has \(H^1(\widetilde N;k)=0\), and
the only possible obstruction lies in \(H^3(L;k)=0\).
The choices of descent form a torsor for \(H^2(L;k)=k\).

There is a useful precise version of the boundary part of this
argument. The boundary group is the split extension
\[
       \pi_1(M_L)=\mathbb Z^3\rtimes L,                    \tag{D.OF8}\label{maps:eq:D.OF8}
\]
where the normal \(\mathbb Z^3\) is generated by the two shorts
and \(\theta\), and the two long loops give the specified splitting.
Its dual module \(H^1(\widetilde M;k)\) is exactly
\(\operatorname{Aff}\): evaluating a covector on
\(\theta+v\) produces its affine function, and translating \(v\)
gives the deck action.
For boundary classes with zero restriction to \(\widetilde M\),
the degree-two filtration is the short exact sequence
\[
 0\longrightarrow H^2(L;k)\longrightarrow
 \ker\{H^2(M_L;k)\to H^2(\widetilde M;k)\}
 \longrightarrow H^1(L;\operatorname{Aff})\longrightarrow0.
                                                               \tag{D.OF9}\label{maps:eq:D.OF9}
\]
The splitting of \eqref{maps:eq:D.OF8} kills the degree-two transgression;
there is no degree-three group cohomology of \(L\).
These two observations also prove \eqref{maps:eq:D.OF9} directly from the
two-dimensional free resolution of \(L=\mathbb Z^2\).
Restriction to the actual long two-torus splits its leftmost term.

The image in the rightmost group of the boundary value of a
descended support is represented by the full affine cocycle
\(l\mapsto a_l\) in \eqref{maps:eq:D.OF7}. This follows either by chasing
\eqref{maps:eq:D.OF6} in that resolution, or by choosing the relative Thom
cochain \(\Psi\): its deck difference is the boundary one-cocycle
whose value on \(\theta+v\) is precisely \(a_l(v)\).
Thus the gradients, constants, and one long-torus value determine
the entire boundary class. This statement retains every
\(\operatorname{Ext}(H_1(M_L),k)\) ambiguity.

In our positive fan convention a core cochain with short-long
matrix \(E\), zero angular area and zero short-normal strip values
has affine cocycle
\[
 a_{l_j}(v)=-E(v,b_j)+c_j,\qquad j=1,2,                  \tag{D.OF10}\label{maps:eq:D.OF10}
\]
where \(b_1=e_2,b_2=\gamma\) and \(l_j=(A-I)b_j\).
For a direct verification, translate the marked \(\theta,b_j\)
strip from component zero to component \(v\). Its change is the
oriented short-long torus with value \(-E(v,b_j)\).
The initial strip contributes \(c_j\). There is no additional
angular-area term, since that coefficient is zero. The literal
basis-loop prism correction is the one already included in
\eqref{maps:eq:D.OF3}. This gives \eqref{maps:eq:D.OF10} at the cochain level, including
its constant term.

For \eqref{maps:eq:D.OF2}--\eqref{maps:eq:D.OF4}, \(E=\left(\begin{smallmatrix}1&0\\0&0
\end{smallmatrix}\right)\). The required affine functions are
therefore
\[
             a_{(8,-6)}=-r+4,\qquad a_{(24,-16)}=4.
                                                               \tag{D.OF11}\label{maps:eq:D.OF11}
\]
The support \eqref{maps:eq:D.OF1} has exactly these differences.
This is an identity for all integers, not a bounded numerical
test: write \(r=8a+r_0,s=2b+s_0\) with \(0\le r_0<8\) and
\(0\le s_0<2\), and substitute. The two differences are
\(4-9r_0\) and \(12-24r_0\), respectively, which reduce to
\eqref{maps:eq:D.OF11}.
Choose the remaining descent parameter so its value on the
actual long torus \(\langle e_2,\gamma\rangle\) is zero,
as it is for \eqref{maps:eq:D.OF2}.
Equations \eqref{maps:eq:D.OF9}--\eqref{maps:eq:D.OF11} now prove equality on the full
cusp boundary, including both its finite-coefficient and
integral-torsion contributions.

For an independent check of the short-long marking, retain the
actual annulus map from \ref{maps:src:D.2.43}, not only its augmentation.
If \(x^py^q-1=(x-1)S_x+(y-1)S_y\), the angular vector \((a,b)\)
closes to the \(I^2\)-class with coefficients
\[
             (aS_x,\ bS_x+aS_y,\ bS_y)                   \tag{D.OF12}\label{maps:eq:D.OF12}
\]
on \(((x-1)^2,(x-1)(y-1),(y-1)^2)\).
Its literal Picard columns are
\[
 x^{-1}y^{-2}E_2+y^{-1}E_3,\quad
       -x^{-1}y^{-1}E_2,\quad
       x^{-1}E_1+x^{-2}y^{-1}E_2.                        \tag{D.OF13}\label{maps:eq:D.OF13}
\]
In the present positive component labels \(x\) translates by
\((-1,0)\), and \(y\) by \((0,1)\); the path and angular exponents
are changed simultaneously. Substituting \eqref{maps:eq:D.OF1} into all the
shifted columns gives the integer values
\[
                         \begin{pmatrix}9&24\\0&0\end{pmatrix}
                  \equiv E\pmod8.                       \tag{D.OF14}\label{maps:eq:D.OF14}
\]
The path dependence in \(S_x,S_y\) is exactly an
Hilbert--Burch relation, so this checks the actual annulus
marking rather than a rational replacement.

\paragraph{Extension through the other caps.}
Each of the four other multiplicative fibres has primitive
rank-one monodromy. Its boundary \(H_1\) is free; so is the
ordinary cap boundary \(H_1\).
The regular-fibre specialization is onto the degree-two group
of each such cap: in a primitive nodal marking this is the
specialization from \(T^2\times T^2\) to
\((S^2\vee S^1)\times T^2\); for the ordinary cap it is the
identity on the central \(T^4\).
Their actual degree-two boundary kernels
are precisely the columns used in \(M\), including the
ordinary cap's affine-meridian column.
Equation \eqref{maps:eq:D.OF2} therefore makes the specified core cochain
extend across each such cap. The free boundary \(H_1\)
ensures that agreement of the resulting integral-cycle
functional is agreement of the full finite-coefficient
cohomology class. These choices and the cusp comparison
glue to the asserted \(u\in H^2(X;k)\).
Its torsion value is the unit \(-1\), unchanged from \eqref{maps:eq:D.OF2};
we choose the opposite torsion generator to normalize it to one.
Choose free lifts of the homology basis dual to \eqref{maps:eq:D.PC10}
which are annihilated by \(u\). This fixes the splitting
\[
                         H_2(X)=\mathbb Z^{15}\oplus\mathbb Z/8
                                                               \tag{D.OF15}\label{maps:eq:D.OF15}
\]
used for the fundamental-class coordinates.

\paragraph{The complete coordinate values.}
Let \(R_0,\ldots,R_5\) be the six positive fan rays in cyclic
order and define the actual wall values
\[
 d_j(v)=\Psi(v)+\Psi(v+R_j)
                    -\Psi(v+R_{j-1})-\Psi(v+R_{j+1}).
\]
An integral representative for \(u|_{P_v}\) is
\[
 \eta_v=(d_0+d_1+d_2)H-d_0E_1-d_2E_2-d_4E_3.             \tag{D.OF16}\label{maps:eq:D.OF16}
\]
Opposite wall restrictions agree modulo eight. Changes of integral
representative by eight times a class preserve its square modulo
sixteen. The full exact component table is retained in
\path{records/maps_weighted_os45_mu.json}.

For component labels \(i,j\), all mixed ordinary coordinates are
specified without a long matrix display by
\[
 \begin{aligned}
 \langle ux_i^2,[X]\rangle
                    &=-\sum_{a=0}^{5}d_a(v_i)\pmod8,\\
 \langle ux_ix_j,[X]\rangle
                    &=\sum_{\{a:\ v_i+R_a=v_j\ {\rm mod}\ L\}}
                                       d_a(v_i)\pmod8\quad(i\ne j),\\
 \langle ux_i y,[X]\rangle
                    &=(-1)^{s_i}\eta_i\cdot(H-E_1)\pmod8.
 \end{aligned}                                             \tag{D.OF17}\label{maps:eq:D.OF17}
\]
The last value is four for component labels \(0,1,14,15\)
and zero for the other twelve. These formulas use actual
Cartier restrictions and retain the repeated vertical edges.
On the primitive free basis \(x_0,\ldots,x_{13},y\), the remaining
power coordinates are
\[
 \begin{split}
 \langle\mathcal P_2(u)x_i,[X]\rangle&=7\pmod{16}
                                             \quad(0\le i<14),\\
 \langle\mathcal P_2(u)y,[X]\rangle&=0\pmod{16},\\
 \langle u^3,[X]\rangle&=0\pmod8,\qquad
 \langle uy^2,[X]\rangle=0\pmod8.
 \end{split}                                               \tag{D.OF18}\label{maps:eq:D.OF18}
\]
The first line is \(\eta_i^2\) on the normalized component.
The two omitted component labels have square fifteen modulo
sixteen, consistently with their primitive-basis relations.
The remaining three evaluations, including the full
order-sixteen value which cannot be obtained by division,
are proved in the following localization input.
\textbf{The remaining finite-coefficient evaluations for OS45}
Let $u\in H^2(X;\mathbb Z/8)$ be the actual primitive torsion
coordinate constructed by the complete boundary comparison, with
cusp support
\[
 \Psi(r,s)=4\left\lfloor\frac r8\right\rfloor
       +3(r\bmod8)\left\lfloor\frac s2\right\rfloor
                           \quad\pmod8.                  \tag{D.OL1}\label{maps:eq:D.OL1}
\]
Here $r\bmod8$ is the representative from zero to seven. This
subsection uses that comparison on the whole cusp boundary;
matching only the regular-fibre bilinear form would not suffice.
The actual affine changes are
\[
 \Delta_{(8,-6)}\Psi=-r+4,\qquad
 \Delta_{(24,-16)}\Psi=4\quad\pmod8.                     \tag{D.OL2}\label{maps:eq:D.OL2}
\]
In particular both have zero slope in the scalar $s$ direction.
Let $y$ be the primitive class in \eqref{maps:eq:D.PC9}.
We prove the three full evaluations
\[
 \langle u^3,[X]\rangle=0\pmod8,\qquad
 \langle\mathcal P_2(u)y,[X]\rangle=0\pmod{16},\qquad
 \langle uy^2,[X]\rangle=0\pmod8.                         \tag{D.OL3}\label{maps:eq:D.OL3}
\]
The second formula includes the order-two part which cannot be
recovered by halving the sum of the component evaluations.

The scalar circle action is semifree. Its fixed locus consists of
the sixteen vertical toric curves $C_{r,s}$ joining vertices
$(r,s)$ and $(r,s+1)$, with $(r,s)$ indexed as in \eqref{maps:eq:D.PC6}.
Each has normal bundle $\mathcal O(-1)\oplus\mathcal O(-1)$
with weights $1,-1$. In every other triangular toric chart the
nonzero weights are also primitive. At each of the four zero-free
$I_1$ caps the primitive angular collapsing vector has first
coordinate one in its elliptic/scalar plane; the scalar circle
therefore projects primitively to the transverse torus and is
free. The ordinary torus cap is free as well. Thus there are no
unrecorded finite-stabilizer strata in this localization.

Since $X$ is simply connected, both $u$ and $y$ have circle-equivariant
lifts. Indeed the only outgoing differentials from ordinary degree
two in the Borel spectral sequence have targets involving
$H^1(X)$ or $H^3(BS^1)$, both zero. Let $x$ be the positive
generator of $H^2(C_{r,s};\mathbb Z)$ and let $t$ be the generator
of $H^2(BS^1;\mathbb Z)$. Write
\[
 \widehat u|_{C_{r,s}}=d_{r,s}x+\lambda_{r,s}t,
 \qquad
 \widehat y|_{C_{r,s}}=\mu_{r,s}t.                         \tag{D.OL4}\label{maps:eq:D.OL4}
\]
There is no $x$ term in the second expression: \eqref{maps:eq:D.PC14} and
$f\cdot C_{r,s}=0$ give $y(C_{r,s})=0$.

For clarity the ordinary degrees and the fibre characters are
computed separately. The degree of the support on a vertical wall is
\[
 d_{r,s}=\Psi(r,s)+\Psi(r,s+1)
              -\Psi(r+1,s)-\Psi(r-1,s+1)\pmod8 .         \tag{D.OL5}\label{maps:eq:D.OL5}
\]
Choosing representatives from zero to seven gives
\[
\begin{array}{c|cc}
 r&d_{r,0}&d_{r,1}\\\hline
 0,7&4&7\\
 1,\ldots,6&0&3.
\end{array}                                                \tag{D.OL6}\label{maps:eq:D.OL6}
\]
The scalar fibre character supplied by the support is its vertical
slope, up to a single common sign. Thus one may take
\[
                 \lambda_{r,0}=0,\qquad
                 \lambda_{r,1}=3r .                      \tag{D.OL7}\label{maps:eq:D.OL7}
\]
This is periodic under the actual period lattice because of
\eqref{maps:eq:D.OL2}. To justify the common character, rather than making a
separate choice on each normalized component, compare along the
actual toric orbit curves. Their degrees determine the differences
of the characters at their fixed endpoints. Horizontal and diagonal
orbit curves connect the sixteen fixed curves in a connected graph:
they join $(r,s)$ to $(r+1,s-1)$ and to $(r-1,s)$ respectively.
The support realizes all these differences. Consequently any global
equivariant lift has the characters in \eqref{maps:eq:D.OL7}, with a common sign
convention and one common additive character $L\in\mathbb Z/8$.
No global integral line lifting $u$ is assumed.

The integral component divisors have their canonical equivariant
linearizations. Along $C_{r,s}$ the two adjacent component normal
weights are $-1,+1$ in the lower/upper vertex order. Thus the weight
of the odd-component sum in \eqref{maps:eq:D.PC9} is $+1$ for $s=0$ and $-1$
for $s=1$. The difference
$2\widehat y-\sum_r\widehat x_{r,1}$ is a global multiple $ct$.
The existence of integral fixed-point weights forces $c$ odd;
changing the lift of $y$ changes $c$ by an even integer. Normalize
$c=-1$. Then
\[
                       \mu_{r,0}=0,\qquad
                       \mu_{r,1}=-1.                     \tag{D.OL8}\label{maps:eq:D.OL8}
\]
Every other lift adds one common integer $V$.

Each restriction of $\widehat u$ in \eqref{maps:eq:D.OL4} lifts integrally on
$C_{r,s}\times BS^1$. Thomas's cochain formula therefore gives
$\mathcal P_2(\widehat u)=(d x+\lambda t)^2$ with the full
modulus sixteen \cite[equations (2.1), (3.4), pp.267--268]{ffThomas56}.
The ordinary cube uses modulus eight. The equivariant normal Euler
class is
\[
                           (-x+t)(-x-t)=-t^2.
\]
Localization is valid over these coefficient rings because the
complement of the fixed locus is free and all normal weights are
units. Taking the coefficient of $x$ after division by $-t^2$ gives
\[
 \begin{aligned}
 \langle u^3,[X]\rangle&=-3\sum_{r,s}d_{r,s}\lambda_{r,s}^2
                                                        &&\pmod8,\\
 \langle\mathcal P_2(u)y,[X]\rangle
       &=-2\sum_{r,s}d_{r,s}\lambda_{r,s}\mu_{r,s}
                                                        &&\pmod{16},\\
 \langle uy^2,[X]\rangle&=-\sum_{r,s}d_{r,s}\mu_{r,s}^2
                                                        &&\pmod8.
 \end{aligned}                                             \tag{D.OL9}\label{maps:eq:D.OL9}
\]
The degree-zero base ring injects into its localization, so these
are identities before inverting $t$ as well.

The representatives \eqref{maps:eq:D.OL6}--\eqref{maps:eq:D.OL8} satisfy the exact integer sums
\[
 \sum d_{r,s}=40,\qquad \sum d_{r,s}\lambda_{r,s}=336,
 \qquad \sum d_{r,s}\lambda_{r,s}^2=5544,
 \qquad \sum_r d_{r,1}=32.                               \tag{D.OL10}\label{maps:eq:D.OL10}
\]
The three values in \eqref{maps:eq:D.OL9} are consequently $-16632$, $672$,
and $-32$, proving \eqref{maps:eq:D.OL3} with their full moduli.
For completeness, under the replacements
$\lambda\mapsto L+\epsilon\lambda$ and $\mu\mapsto V+\mu$,
where $\epsilon=\pm1$, the exact values are
\[
 \begin{aligned}
 -120L^2-2016\epsilon L-16632,\\
 -80LV+64L-672\epsilon V+672\epsilon,\\
 -40V^2+64V-32.
 \end{aligned}                                             \tag{D.OL11}\label{maps:eq:D.OL11}
\]
Every coefficient is divisible by eight, sixteen, and eight,
respectively. Thus the results are independent of all global
linearization choices, not only of the displayed normalization.

The same support gives all component mixed evaluations. For a
vertex $v$, let $R_j$ be the six cyclic rays of \ref{maps:src:D.2.47} and put
\[
 d_j(v)=\Psi(v)+\Psi(v+R_j)
                  -\Psi(v+R_{j-1})-\Psi(v+R_{j+1}).
\]
An integral representative of its component restriction is
\[
 \eta_v=(d_0+d_1+d_2)H-d_0E_1-d_2E_2-d_4E_3.
                                                               \tag{D.OL12}\label{maps:eq:D.OL12}
\]
Changing the representative by eight times an integral class
changes neither its square modulo sixteen nor any of its needed
pairings modulo eight. In the primitive basis of \eqref{maps:eq:D.PC10},
\[
 \langle\mathcal P_2(u)x_i,[X]\rangle=7\pmod{16}
                     \quad(0\le i<14),
\]
while the two omitted components $i=14,15$ have value fifteen.
Also $\langle ux_i y,[X]\rangle$ equals four for $i=0,1,14,15$
and zero for all the other component labels. The remaining
$\langle ux_i x_j,[X]\rangle$ are exactly
$\eta_i\cdot\nu_i^*D_j\pmod8$, with the explicit Cartier
restriction in \ref{maps:src:D.2.47}. Together with \eqref{maps:eq:D.OL3} these determine the
finite-coordinate evaluations once the support comparison is
established.

The exact program
\path{code/III/fourfolds_OS45_localization.py}
retains every wall value and every $\eta_v$, verifies the actual
quasiperiods on all sixteen residues, and proves the six
all-character polynomial divisibility identities \eqref{maps:eq:D.OL11}.
It explicitly records the required geometric support comparison
as a separate hypothesis.

In particular \(y^2=0\) integrally. Its free part is zero by
\eqref{maps:eq:D.PC12}, and its torsion part is detected injectively by the
primitive \(u\); the last equality of \eqref{maps:eq:D.OF18} kills that part.

Finally \ref{maps:src:D.2.48}'s primary coefficient-group calculation and
the integral Kunneth decomposition give
\[
 H_6(K(\mathbb Z^{15}\oplus\mathbb Z/8,2);\mathbb Z)
       =\mathbb Z^{680}\oplus(\mathbb Z/8)^{120}
                         \oplus(\mathbb Z/16)^{15}
                         \oplus\mathbb Z/8.              \tag{D.OF19}\label{maps:eq:D.OF19}
\]
The free coordinates are exactly the primitive integral cubic
\eqref{maps:eq:D.PC11}--\eqref{maps:eq:D.PC12}. Equations \eqref{maps:eq:D.OF17}--\eqref{maps:eq:D.OF18} give all
remaining coordinates: the first mixed group is detected by
\(ux_ix_j\), the second by the full Pontryagin square, and the
last by \(u^3\). Hence this determines the complete
fundamental-class invariant \(\mu\), not merely its image
under ordinary triple cup products.
The refined class remains a separate \(0\) versus \(4\)
calculation until the actual odd middle-dimensional generator
and its reference framing have been evaluated.

\path{code/III/maps_weighted_os45_mu.py} checks the exact Smith
cochain, both full affine strip constants, all sixteen wall
restrictions, the full shifted annulus map, every primitive
mixed entry and the complete moduli in \eqref{maps:eq:D.OF18}.
The independent fourfolds programs check all primitive cubic
entries and all-character localization identities.
\end{enumerate}
\end{proof}
\subsection{Moving the zero of the linearization}
\begin{theorem}[the global moving-zero family]\label{maps:src:D.3.1}
\StatementRecord{maps:src:D.3.1}{proved}{Historical events: Claude Fable 5.2 (Anthropic), effort high, 16 November 2026: carried as background without fault found. GPT-6 Astra (OpenAI), ChatGPT Pro mode, effort Pro, S6\_09\_PAPERS, 17 November 2026: targeted proof check. Also checked in this writing. Pre-continuation writing events: GPT-6 Astra (OpenAI), ChatGPT Pro mode, effort Pro, S6\_10\_PAPERS, 18 November 2026: targeted proof check. Continuation metadata events: GPT-6 Astra (OpenAI), ChatGPT Pro mode, effort Pro, S6\_10\_PAPERS continuation, 19 November 2026: used and rechecked in this run; status-vocabulary reconciliation; metadata reconciliation only; no new mathematical check. \emph{Pre-continuation recorded proof scope.} GPT-6 Astra (OpenAI, ChatGPT Pro mode), effort Pro, 2026-11-17 UTC; Claude Fable 5.2 (Anthropic), effort high, 16 November 2026; also checked in this writing by GPT-6 Astra (OpenAI), ChatGPT Pro mode, effort Pro, 18 November 2026}{the hypotheses, notation and definitions stated here; \cite{manEngel26}}


Let $S\to\mathbb P^1$ be Engel's rational elliptic surface, with
fibres $IV^*,III,I_1$ at $0,1,\infty$, respectively. Put
$L=\mathcal O_S(P-O)$, $T=t_{6P}$ and
$J=L\otimes T^*L^{-1}\simeq\pi^*\mathcal O_{\mathbb P^1}(1)$,
with the direction of $T$ fixed in \cref{push:src:A.1.2}. We write $X^{\log}(u)$ for the logarithmically transformed sphere, and $X(u)$ for the untwisted fibre of \cref{push:src:A.1.2} and \cref{push:src:A.2.2}. Choose sections
$\phi_0,\phi_1$ of $\mathcal O(1)$ with divisors $\infty,0$,
respectively, normalized so that near infinity
$\phi_0=q\phi_1$, $q=1/t$. Set
$\phi_s=\pi^*(\phi_0-s\phi_1)$.
For every sufficiently small fixed $u\ne0$ there is a proper flat
holomorphic family $\mathcal Y_u\to\Delta_s$, with a holomorphic
map to $\mathbb P^1$, having the following properties. We perform
the same logarithmic transformations of multiplicities $3,4$ and
the same admissible sphere twist data as for the chosen $X^{\log}(u)$.
The central threefold $Y_{u,0}$ has exactly one ordinary double point
and is the contraction of the diagonal $(-1,-1)$-curve of $X^{\log}(u)$.
Every fibre $Y_{u,s}$ for $s\ne0$ is smooth. Its local change at the
contracted point is
\[
 xy=u(zv-s),\qquad q=zv.
\]
This is the conifold transition described in Engel's Remark 8.5.
It is transverse to the $u\to0$ degeneration of \cref{push:src:A.1.2}.
\end{theorem}


\begin{proof}
\begin{enumerate}[label=\arabic*.,leftmargin=*]
\item The equality defining $J$ and its normalization are checked in
\cref{push:src:A.1.2}, using Engel's Lemma~2.1 and Corollary~2.2 \cite{manEngel26}. We retain
that equality, but allow the section of $J$ to vary. Here are the
charts and their identifications. Choose Hermitian norms and
$0<r<1/4$. On the compact surface $S$, after restricting both
parameters there is a uniform inequality
$|u|\|\phi_s(x)\|<r^2$. Take a core
\[
 C=\{(x,a,s):a\in L_x,\ r<\|a\|<r^{-1}\}
\]
and a node chart
\[
 N=\{(x,a,b,s):a\in L_x,\ b\in L^{-1}_{Tx},\
           a\otimes b=u\phi_s(x),\ \|a\|,\|b\|<2r\}.
\]
On the collar $r<\|a\|<2r$, identify $N$ with $C$ by
$(x,a,b,s)\mapsto(x,a,s)$. On $r<\|b\|<2r$, identify it by
$(x,a,b,s)\mapsto(Tx,b^{-1},s)$. These maps and their inverses are
holomorphic; on the first collar $b=u\phi_s/a$, and on the second
$a=u\phi_s b^{-1}$. The two node collars are disjoint by the
uniform inequality. A change of local line frames multiplies both
sides of the node equation by the same unit, so the construction is
intrinsic.

\item The quotient is Hausdorff. At each end of an attaching collar one
chart contains the limiting circle and the other omits it; there is
no pair of distinct unpaired limits of the gluing graph. The two
node collars are separated, so no new chain of identifications
changes this conclusion. For properness, the closed core
$3r/2\leq\|a\|\leq2/(3r)$ and the closed node subset
$\|a\|,\|b\|\leq7r/4$ cover the quotient. Over any compact set
of parameters their preimages are compact because $S$ is compact.
This proves properness, also for the joint map to
$\mathbb P^1\times\Delta_s$.

\item On the core, flatness over $s$ is immediate. On the node chart, the
local rings are hypersurface rings over the smooth surface $S$;
the equation $ab=u\phi_s$ is not divisible by $s-s_0$, so
multiplication by $s-s_0$ is injective. A coherent module over the
one-dimensional regular local parameter ring is flat precisely
when it is torsion free; hence the family is flat. Alternatively,
for the joint map to $(t,s)$ the fibres all have pure dimension
two, and the Cohen--Macaulay regular-sequence criterion gives
flatness over the smooth two-dimensional base.

\item Away from the zero of $\phi_s$, the collar relation gives exactly
$L^\times/\langle u\phi_s\rangle$, with the contracting generator
covering $T^{-1}$. The uniform norm inequality ensures that every
orbit meets a core or node representative, so this is an actual
biholomorphic comparison. Near $0$ and $1$ the ratio
$\phi_s/\phi_0$ is a holomorphic unit, uniformly for small $s$.
A local logarithm changes the scalar period by a holomorphic
function. Thus the same good-reduction torus cover and the same
finite affine translation data extend holomorphically in $s$.
Their finite cyclic actions stay free, by the unchanged primitive
invariant translation coordinate. We replace the two corresponding
pieces by these smooth logarithmic-transform models and glue on
the punctured overlaps by the just-described period comparison.
This is precisely the original finite gluing when $s=0$.
Properness and flatness persist, since the replacements are proper
families with the same compact annular boundaries and are locally
finite free quotients of smooth families.

\item At the node of $S_\infty$ write $q=zv$. A unit change in a line
coordinate gives $xy=u(zv-s)$. The relative gradient in
$(x,y,z,v)$ is $(y,x,-uv,-uz)$, so the only possible critical point
is the origin, and it lies on the fibre exactly when $s=0$.
At a smooth point of a zero fibre of $\phi_s$, its zero is simple,
so $xy=u(q-s)$ is smooth. These observations exhaust the singular
points: the other nodes are absent because the elliptic surface
has only its stated three exceptional fibres, and the two finite
ones were replaced by the smooth free quotients. At $s=0$ the
quadratic form $xy-uzv$ is nondegenerate and gives one ordinary
double point. Its small resolution is the diagonal subdivision of \cref{push:src:A.2.1} and \cref{push:src:A.2.2}. In one chart write $x=az$, $uv=ay$; in the other write
$z=bx$, $y=buv$, with $b=a^{-1}$. The two normal coordinates
transform by multiplication by $a$, giving
$\mathcal O(-1)\oplus\mathcal O(-1)$ over the exceptional
$\mathbb P^1$. There is one such diagonal modulo the deck periods.
The comparison with \cref{push:src:A.1.2} and \cref{push:src:A.2.2} identifies its small resolution with the
chosen $X^{\log}(u)$.

\item Before the finite logarithmic replacements, both constructions sit in
the global core/node plumbing with the explicit two-parameter equation
$xy=u(zv-s)$. The checkpoint follows the slice $s=0$ and varies
$u$ towards zero; the present construction fixes $u\ne0$ and varies
$s$. At $u=0$ the node equation is $xy=0$ and recovers the same
untwisted pinched projective bundle for every $s$. The finite
logarithmic replacements used for the sphere have only been constructed
here for $u\ne0$; no extension of those replacements across $u=0$
is asserted. The common cusp chart and the untwisted global plumbing
are the precise comparison with \cref{push:src:A.1.2} and \cref{push:src:A.2.2}.
\end{enumerate}
\end{proof}

Engel's general zero/pole construction is stated in \cite[Remark~8.5]{manEngel26}; his Remark~8.1 concerns the separate degeneration as $u$ tends to zero. The full moving-section plumbing and its geometric comparisons are proved above.

The complete credited statement and re-derivation are in \cref{maps:src:D.3.2}.
\begin{theorem}[the complex invariants of the smoothing]\label{maps:src:D.3.3}
\StatementRecord{maps:src:D.3.3}{proved}{Historical events: Claude Fable 5.2 (Anthropic), effort high, 16 November 2026: carried as background without fault found. GPT-6 Astra (OpenAI), ChatGPT Pro mode, effort Pro, S6\_09\_PAPERS, 17 November 2026: targeted proof check. Also checked in this writing. Pre-continuation writing events: GPT-6 Astra (OpenAI), ChatGPT Pro mode, effort Pro, S6\_10\_PAPERS, 18 November 2026: targeted proof check. Continuation metadata events: GPT-6 Astra (OpenAI), ChatGPT Pro mode, effort Pro, S6\_10\_PAPERS continuation, 19 November 2026: used and rechecked in this run; status-vocabulary reconciliation; metadata reconciliation only; no new mathematical check. \emph{Pre-continuation recorded proof scope.} GPT-6 Astra (OpenAI, ChatGPT Pro mode), effort Pro, 2026-11-17 UTC; Claude Fable 5.2 (Anthropic), effort high, 16 November 2026; also checked in this writing by GPT-6 Astra (OpenAI), ChatGPT Pro mode, effort Pro, 18 November 2026}{the hypotheses, notation and definitions stated here; the explicit constructions and calculations in the proof}


For $u\ne0$ sufficiently small and $0<|s|\ll1$, let $g:Y_{u,s}\to B$
be the above map and $S'_2=(g^{-1}(1))_{\mathrm{red}}$. Then
\[
 K_{Y_{u,s}}\simeq g^*\mathcal O_B(-1)\otimes\mathcal O_{Y_{u,s}}(2S'_2),
 \quad a(Y_{u,s})=1,
 \quad (h^{0,0},h^{0,1},h^{0,2},h^{0,3})=(1,1,0,0).
\]
In particular $K^{\otimes12}\simeq g^*\mathcal O_B(-6)$ and the
Kodaira dimension is $-\infty$.
\end{theorem}


\begin{proof}
\begin{enumerate}[label=\arabic*.,leftmargin=*]
\item We first compute the canonical bundle by actual transition forms.
On the core over a surface chart, the canonical form is
$\Omega_S\wedge da/a$. Changing a line frame adds a base one-form
to $da/a$, which wedges to zero with $\Omega_S$. On the second
plumbing collar the transition covers $T$ and changes the line
coordinate by the scalar $u\phi_s$; its logarithmic derivative
again contributes only a base form. Translation on the elliptic
surface preserves its invariant relative differential. Thus the
canonical transition functions are those of $K_S=\pi^*\mathcal O(-1)$.
On $xy=u(q-s)$ the residue expression
$dq\wedge dr\wedge dx/x$ becomes, using the equation,
a nonzero constant multiple of $dy\wedge dr\wedge dx$; it extends
across the node-chart loci where a line coordinate vanishes.
The expression over $q=zv$ has the same residue interpretation.
For $s\ne0$ it is regular wherever the total threefold is smooth.
Consequently moving $\phi$ introduces no new canonical divisor.

\item The two finite modifications give the same precise corrections as
in the original threefold. Indeed, for good reduction at $IV^*$
write $t=\zeta^3$, $x=\zeta^4X$, $y=\zeta^6Y$; then
$dt\wedge dx/y$ has order $2-2=0$ in $\zeta$. At $III$ write
$t-1=\zeta^4$, $x=\zeta^2X$, $y=\zeta^3Y$; its order is
$3-1=2$. Wedging with the logarithmic line coordinate gives the
threefold form, and the free quotient is etale as a map of total
spaces. These orders therefore descend respectively to zero and
$2S'_2$. This proves the asserted canonical formula. Since
$4S'_2=g^*(1)$, its twelfth power is $g^*\mathcal O(-6)$.
Connected fibres imply $g_*\mathcal O_Y=\mathcal O_B$, so none of
its positive powers has a section; if a power of $K$ had a section,
its twelfth power would contradict this fact.

\item For the algebraic dimension, use a marked regular fibre with
period matrix $\left(\begin{smallmatrix}6\mu&\tau&1&0\\
\beta_s&\mu&0&1\end{smallmatrix}\right)$.
Moving the zero changes only the scalar period $\beta_s$. Its
monodromy around the new zero $q=s$ is
$\beta_s\mapsto\beta_s\pm1$, while $\tau,\mu$ are unchanged.
If constants $a,b,c,d,e$ satisfy
$a+b\tau+c\mu+d\beta_s+e(6\mu^2-\tau\beta_s)=0$
identically, continuation around that zero gives $d-e\tau=0$.
The elliptic period $\tau$ is nonconstant, so $d=e=0$.
Continuation around infinity gives $b=0$ because
$\tau\mapsto\tau-1$ and $\mu$ is unchanged. Finally $\mu$ is
nonconstant: a constant satisfying
$\mu\circ g_1=(1-\mu)/\tau$ would force a nonconstant $\tau$ to
be constant, or give $0=1$. Hence $c=a=0$.
The type-$(1,1)$ equation for an integral alternating form now
shows that, outside a countable union of discrete subsets of the
base, the Neron--Severi group is generated by
$q=u\wedge w+6\gamma\wedge\delta$. Its associated Hermitian form
is indefinite: it is nondegenerate, its signature varies
continuously on the connected smooth-fibre locus, and at $s=0$
it has signature $(1,1)$ by the explicit period inequality.
It therefore supplies neither an ample class nor a nonzero
square-zero semipositive class. The algebraic reduction of a
complex torus is an abelian quotient; an abelian surface quotient
would supply the former class, and an elliptic quotient the latter.
Thus a very general fibre has algebraic dimension zero.
Every meromorphic function on $Y$ is constant on these fibres.
Its vertical differential consequently vanishes identically on
the regular locus, by the identity theorem, and it descends to a
meromorphic function on $B$. At an exceptional value this descent
extends meromorphically by restriction to a transverse holomorphic
disc and the bounded-pole criterion. Therefore
$\mathbb C(Y)=g^*\mathbb C(B)$ and $a(Y)=1$.

\item It remains to compute the structure-sheaf cohomology. The small
resolution $X^{\log}(u)\to Y_{u,0}$ has
$\pi_*\mathcal O_X=\mathcal O_{Y_{u,0}}$ and
$R^i\pi_*\mathcal O_X=0$ for $i>0$. For the latter assertion,
formal neighbourhoods of the exceptional curve have successive
quotients $\operatorname{Sym}^n(\mathcal O(1)^{\oplus2})$;
all their positive cohomology groups vanish, and formal functions
and coherence give the stated vanishing. Thus
$H^q(Y_{u,0},\mathcal O)=H^q(X^{\log}(u),\mathcal O)$, whose dimensions
are $(1,1,0,0)$, as computed by the supplied direct-image formulas.
Upper semicontinuity gives $h^{0,2}=h^{0,3}=0$ nearby. In a proper
flat analytic family the holomorphic Euler characteristic is
constant, as follows from the finite complex representing the
proper derived direct image over a small parameter disc. Its
value here is zero; connectedness gives $h^{0,0}=1$, and hence
$h^{0,1}=1$.
\end{enumerate}
\end{proof}

The complete credited statement and re-derivation are in \cref{maps:src:D.3.4}.

\begin{theorem}[Credited re-derivation: the diffeomorphism type of every cover]\label{maps:src:D.1.11}
\StatementRecord{maps:src:D.1.11}{proved}{Historical events: Claude Fable 5.2 (Anthropic), effort high, 16 November 2026: carried as background without fault found. GPT-6 Astra (OpenAI), ChatGPT Pro mode, effort Pro, S6\_09\_PAPERS, 17 November 2026: targeted proof check. Also checked in this writing. Pre-continuation writing events: GPT-6 Astra (OpenAI), ChatGPT Pro mode, effort Pro, S6\_10\_PAPERS, 18 November 2026: targeted proof check. Continuation metadata events: GPT-6 Astra (OpenAI), ChatGPT Pro mode, effort Pro, S6\_10\_PAPERS continuation, 19 November 2026: used and rechecked in this run; status-vocabulary reconciliation; metadata reconciliation only; no new mathematical check. \emph{Pre-continuation recorded proof scope.} GPT-6 Astra (OpenAI, ChatGPT Pro mode), effort Pro, 2026-11-17 UTC; Claude Fable 5.2 (Anthropic), effort high, 16 November 2026; also checked in this writing by GPT-6 Astra (OpenAI), ChatGPT Pro mode, effort Pro, 18 November 2026}{the hypotheses, notation and definitions stated here; \cref{maps:src:D.1.9}; \cite{manZhubr75}; \cref{maps:src:D.1.4}; \cref{maps:src:D.1.8}; \cref{maps:src:D.1.6}; \cref{maps:src:D.1.10}; \cite{manReiser24}}


For each $d\geq1$ prime to six let $M_d$ denote the oriented
simply connected spin six-manifold with the following complete
Zhubr data: $A=\mathbb Z^{d-1}\oplus\mathbb Z/d$, $b_3=0$,
refined Pontryagin class zero, integral cubic zero, and finite
cyclic cubic as in \ref{maps:src:D.1.9} in the basis
$x_0,\ldots,x_{d-2},u$. These data specify a unique oriented
diffeomorphism type. One has
\[
 \widetilde X_p\cong M_{|p|},\qquad M_1=S^6.
\]
For $d>1$ this is a classification by a complete published invariant;
no assertion that $M_d$ is a torsion-free connected sum is made.
\end{theorem}


\begin{proof}
\begin{enumerate}[label=\arabic*.,leftmargin=*]
\item We apply Zhubr's Theorem~1 \cite{manZhubr75}, with the definitions
on pp.~794--795. Its invariant quadruple consists of $H_2$, the
middle Betti number, the refined class $p$ and
$\mu=f_*[M]\in H_6(K(H_2,2);\mathbb Z)$, modulo isomorphism.
\ref{maps:src:D.1.4} gives the first two data. \ref{maps:src:D.1.8} gives $p_1=0$ and the refined
class satisfies $4p=\operatorname{PD}(p_1)$. Multiplication by four
is injective on $A$, so $p=0$. \ref{maps:src:D.1.6} and \ref{maps:src:D.1.9} give precisely all
the evaluations of $\mu$ enumerated in \ref{maps:src:D.1.10}. Thus the entire
quadruple is specified. Existence is supplied by the actual covers
already constructed, so its admissibility conditions require no
separate realization assumption. The injective classification in
Zhubr's theorem gives uniqueness and the displayed diffeomorphism.
For $d=1$ the same data are those of $S^6$. This also explains why
integral cup products alone, without \ref{maps:src:D.1.9}, would not have sufficed;
compare the torsion discussion in \cite{manReiser24} before Remark~4.7.
\end{enumerate}
\end{proof}
